% ============================================
% BANCO ZK: CRIPTOGRAFÍA Y PRIVACIDAD EN BLOCKCHAIN
% VOLUMEN 2: Herramientas Avanzadas y Ecosistema
% ============================================

% ============================================
% BANCO ZK: CRIPTOGRAFÍA Y PRIVACIDAD EN BLOCKCHAIN
% PREÁMBULO COMPLETO PARA TODOS LOS VOLÚMENES
% ============================================

\documentclass[11pt,a4paper,twoside]{book}

% ============================================
% PAQUETES BÁSICOS
% ============================================
\usepackage[utf8]{inputenc}
\usepackage[T1]{fontenc}
\usepackage[spanish,es-tabla]{babel}
\usepackage{lmodern}

% ============================================
% GEOMETRÍA Y LAYOUT
% ============================================
\usepackage[
    top=2.5cm,
    bottom=2.5cm,
    left=3cm,
    right=2.5cm,
    headheight=14pt
]{geometry}

% ============================================
% MATEMÁTICAS
% ============================================
\usepackage{amsmath,amssymb,amsthm}
\usepackage{mathtools}
\usepackage{bm}

% ============================================
% GRÁFICOS Y COLORES (CRÍTICO PARA TIKZ)
% ============================================
\usepackage{xcolor}
\usepackage{graphicx}

% TikZ con TODAS las librerías necesarias
\usepackage{tikz}
\usetikzlibrary{
    arrows,
    arrows.meta,
    shapes,
    shapes.geometric,
    shapes.symbols,
    shapes.misc,
    shapes.multipart,
    positioning,
    calc,
    patterns,
    patterns.meta,
    decorations,
    decorations.pathreplacing,
    decorations.pathmorphing,
    decorations.markings,
    backgrounds,
    fit,
    matrix,
    chains,
    scopes,
    trees,
    mindmap,
    shadows,
    spy,
    fadings,
    through,
    intersections,
    quotes,
    angles,
    babel
}

% Definir estilos TikZ globales
\tikzset{
    >=stealth,
    node distance=2cm,
    every node/.style={font=\small},
    box/.style={rectangle, draw, rounded corners, minimum width=2cm, minimum height=0.8cm, align=center},
    arrow/.style={->, thick},
    dashed arrow/.style={->, thick, dashed}
}

% ============================================
% TABLAS
% ============================================
\usepackage{booktabs}
\usepackage{array}
\usepackage{tabularx}
\usepackage{longtable}
\usepackage{multirow}
\usepackage{colortbl}

% ============================================
% CÓDIGO
% ============================================
\usepackage{listings}
\usepackage{fancyvrb}

% Estilo para Python - CORREGIDO
\lstdefinestyle{python}{
    language=Python,
    basicstyle=\ttfamily\small,
    keywordstyle=\color{blue}\bfseries,
    commentstyle=\color{green!60!black}\itshape,
    stringstyle=\color{red!70!black},
    showstringspaces=false,          % <-- IMPORTANTE: false
    showspaces=false,                 % <-- AGREGAR
    showtabs=false,                   % <-- AGREGAR
    breaklines=true,
    breakatwhitespace=true,
    tabsize=4,
    frame=single,
    framerule=0.5pt,
    rulecolor=\color{gray!50},
    backgroundcolor=\color{gray!5},
    numbers=left,
    numberstyle=\tiny\color{gray},
    numbersep=8pt,
    xleftmargin=15pt,
    framexleftmargin=15pt,
    aboveskip=10pt,
    belowskip=10pt,
    captionpos=b,
    literate={á}{{\'a}}1 {é}{{\'e}}1 {í}{{\'i}}1 {ó}{{\'o}}1 {ú}{{\'u}}1 {ñ}{{\~n}}1,  % <-- AGREGAR para acentos
    morekeywords={self,True,False,None,yield,async,await,match,case}
}

% Estilo para Solidity
\lstdefinestyle{solidity}{
    language=Java,
    basicstyle=\ttfamily\small,
    keywordstyle=\color{blue}\bfseries,
    commentstyle=\color{green!60!black}\itshape,
    stringstyle=\color{red!70!black},
    showstringspaces=false,
    breaklines=true,
    frame=single,
    backgroundcolor=\color{yellow!5},
    numbers=left,
    numberstyle=\tiny\color{gray},
    xleftmargin=15pt,
    framexleftmargin=15pt,
    morekeywords={contract,function,public,private,internal,external,view,pure,payable,returns,mapping,struct,event,emit,require,modifier,constructor,uint256,uint128,uint64,uint32,uint8,int256,bytes32,bytes,address,bool,string,memory,storage,calldata,assembly,pragma,solidity,import,interface,library,abstract,override,virtual,immutable,constant,indexed,anonymous,error,revert,unchecked}
}

% Estilo para Circom
\lstdefinestyle{circom}{
    language=C,
    basicstyle=\ttfamily\small,
    keywordstyle=\color{purple}\bfseries,
    commentstyle=\color{green!60!black}\itshape,
    stringstyle=\color{orange},
    showstringspaces=false,
    breaklines=true,
    frame=single,
    backgroundcolor=\color{purple!5},
    numbers=left,
    numberstyle=\tiny\color{gray},
    xleftmargin=15pt,
    framexleftmargin=15pt,
    morekeywords={pragma,circom,template,signal,input,output,public,private,component,var,for,if,else,while,return,include,log,assert,main}
}

% Estilo para Bash
\lstdefinestyle{bash}{
    language=bash,
    basicstyle=\ttfamily\small,
    keywordstyle=\color{blue}\bfseries,
    commentstyle=\color{green!60!black}\itshape,
    stringstyle=\color{red!70!black},
    showstringspaces=false,
    breaklines=true,
    frame=single,
    backgroundcolor=\color{black!5},
    numbers=none,
    xleftmargin=10pt,
    framexleftmargin=10pt
}

% Estilo para JavaScript/TypeScript
\lstdefinestyle{javascript}{
    language=Java,
    basicstyle=\ttfamily\small,
    keywordstyle=\color{blue}\bfseries,
    commentstyle=\color{green!60!black}\itshape,
    stringstyle=\color{red!70!black},
    showstringspaces=false,
    breaklines=true,
    frame=single,
    backgroundcolor=\color{cyan!5},
    numbers=left,
    numberstyle=\tiny\color{gray},
    morekeywords={const,let,var,function,async,await,return,if,else,for,while,class,extends,import,export,from,default,new,this,try,catch,throw,finally,typeof,instanceof,in,of,true,false,null,undefined,console,require,module}
}

% ============================================
% TEOREMAS Y DEFINICIONES
% ============================================
\theoremstyle{definition}
\newtheorem{definition}{Definición}[chapter]
\newtheorem{theorem}{Teorema}[chapter]
\newtheorem{lemma}[theorem]{Lema}
\newtheorem{proposition}[theorem]{Proposición}
\newtheorem{corollary}[theorem]{Corolario}

\theoremstyle{remark}
\newtheorem{remark}{Observación}[chapter]
\newtheorem{example}{Ejemplo}[chapter]

% ============================================
% CAJAS PERSONALIZADAS
% ============================================
\usepackage{tcolorbox}
\tcbuselibrary{skins,breakable,theorems}

% Caja KeyIdea
\newtcolorbox{keyidea}{
    colback=blue!5,
    colframe=blue!75!black,
    fonttitle=\bfseries,
    title=Idea Clave,
    breakable,
    enhanced,
    boxrule=1pt,
    arc=2mm,
    left=8pt,
    right=8pt,
    top=8pt,
    bottom=8pt
}

% Caja Banco ZK
\newtcolorbox{bancozk}{
    colback=green!5,
    colframe=green!50!black,
    fonttitle=\bfseries,
    title=Aplicación en Banco ZK,
    breakable,
    enhanced,
    boxrule=1pt,
    arc=2mm,
    left=8pt,
    right=8pt,
    top=8pt,
    bottom=8pt
}

% Caja Warning
\newtcolorbox{warning}{
    colback=red!5,
    colframe=red!75!black,
    fonttitle=\bfseries,
    title=Advertencia,
    breakable,
    enhanced,
    boxrule=1pt,
    arc=2mm,
    left=8pt,
    right=8pt,
    top=8pt,
    bottom=8pt
}

% Caja Note
\newtcolorbox{note}{
    colback=yellow!5,
    colframe=yellow!50!black,
    fonttitle=\bfseries,
    title=Nota,
    breakable,
    enhanced,
    boxrule=1pt,
    arc=2mm,
    left=8pt,
    right=8pt,
    top=8pt,
    bottom=8pt
}

% Caja para código importante
\newtcolorbox{codebox}{
    colback=gray!5,
    colframe=gray!50!black,
    fonttitle=\bfseries,
    breakable,
    enhanced,
    boxrule=0.5pt,
    arc=1mm,
    left=5pt,
    right=5pt,
    top=5pt,
    bottom=5pt
}

% ============================================
% ENCABEZADOS Y PIES
% ============================================
\usepackage{fancyhdr}
\pagestyle{fancy}
\fancyhf{}
\fancyhead[LE]{\leftmark}
\fancyhead[RO]{\rightmark}
\fancyfoot[C]{\thepage}
\renewcommand{\headrulewidth}{0.4pt}
\renewcommand{\footrulewidth}{0pt}

% Estilo para páginas de capítulo
\fancypagestyle{plain}{
    \fancyhf{}
    \fancyfoot[C]{\thepage}
    \renewcommand{\headrulewidth}{0pt}
}

% ============================================
% HIPERVÍNCULOS
% ============================================
\usepackage{hyperref}
\hypersetup{
    colorlinks=true,
    linkcolor=blue!70!black,
    citecolor=green!50!black,
    urlcolor=purple!70!black,
    bookmarks=true,
    bookmarksopen=true,
    pdfstartview=FitH
}

% ============================================
% OTROS PAQUETES ÚTILES
% ============================================
\usepackage{float}
\usepackage{caption}
\usepackage{subcaption}
\usepackage{enumitem}
\usepackage{pifont}
\usepackage{microtype}  % Mejora tipografía
\usepackage{setspace}   % Espaciado

% Checkmarks
\newcommand{\cmark}{\ding{51}}
\newcommand{\xmark}{\ding{55}}

% Comando para cleardoublepage simplificado
\let\origcleardoublepage\cleardoublepage
\renewcommand{\cleardoublepage}{\clearpage}

% ============================================
% COMANDOS PERSONALIZADOS
% ============================================

% Campos finitos
\newcommand{\F}{\mathbb{F}}
\newcommand{\Fp}{\mathbb{F}_p}
\newcommand{\Fq}{\mathbb{F}_q}

% Grupos
\newcommand{\G}{\mathbb{G}}
\newcommand{\Z}{\mathbb{Z}}
\newcommand{\N}{\mathbb{N}}

% Funciones hash
\newcommand{\Hash}{\textsf{Hash}}
\newcommand{\Commit}{\textsf{Commit}}
\newcommand{\Verify}{\textsf{Verify}}
\newcommand{\Prove}{\textsf{Prove}}

% ZK específicos
\newcommand{\snark}{\textsf{SNARK}}
\newcommand{\stark}{\textsf{STARK}}
\newcommand{\plonk}{\textsf{PLONK}}
\newcommand{\groth}{\textsf{Groth16}}

% Complejidad
\newcommand{\bigO}{\mathcal{O}}

\usetikzlibrary{
    arrows, arrows.meta,
    shapes, shapes.geometric, shapes.symbols, shapes.misc, shapes.multipart,
    positioning, calc, patterns, patterns.meta,
    decorations, decorations.pathreplacing, decorations.pathmorphing, decorations.markings,
    backgrounds, fit, matrix, chains, scopes, trees, mindmap, shadows,
    spy, fadings, through, intersections, quotes, angles, babel
}



% ============================================
% INICIO DEL DOCUMENTO - VOLUMEN 3
% ============================================
\begin{document}
% ============================================
% PORTADA
% ============================================
\begin{titlepage}
    \centering
    \vspace*{1cm}
    
    {\Large\scshape Serie Banco ZK}\\[0.3cm]
    {\small 3 Vol\'umenes $\cdot$ 1000+ p\'aginas $\cdot$ 20 Cap\'itulos}\\[0.5cm]
    \rule{\textwidth}{1.5pt}\\[0.5cm]
    
    {\Huge\bfseries Volumen III}\\[0.5cm]
    {\LARGE Herramientas Avanzadas y Ecosistema}\\[0.3cm]
    {\large Cap\'itulos 14--20: Identidad, Contratos Privados, Seguridad}\\[1.5cm]
    
    \begin{tikzpicture}
        \node[rectangle, draw, rounded corners, minimum width=12cm, minimum height=5cm, fill=zkblue!5] {
            \begin{minipage}{11cm}
            \centering
            \textbf{Contenido del Volumen 3:}\\[0.5cm]
            \small
            \textbf{Parte I:} Identidad y Regulaci\'on\\
            Credenciales An\'onimas $\bullet$ Compliance vs Privacidad\\[0.3cm]
            \textbf{Parte II:} Privacidad en Red\\
            Metadata Protection $\bullet$ Network Privacy $\bullet$ Tor $\bullet$ Mixnets\\[0.3cm]
            \textbf{Parte III:} Contratos y Criptograf\'ia\\
            Private Smart Contracts $\bullet$ MPC $\bullet$ FHE $\bullet$ TEEs\\[0.3cm]
            \textbf{Parte IV:} Seguridad y Herramientas\\
            Ataques y Defensa $\bullet$ Auditor\'ia $\bullet$ Ecosistema ZK
            \end{minipage}
        };
    \end{tikzpicture}
    
    \vfill
    
    {\Large Walter Milstein}\\[0.5cm]
    {\large 2025}
\end{titlepage}
% ============================================
% P\'AGINA DE COPYRIGHT
% ============================================
\thispagestyle{empty}
\null\vfill
\begin{center}
\textbf{Banco ZK: Criptograf\'ia y Privacidad en Blockchain}\\
\textbf{Volumen 3: Herramientas Avanzadas y Ecosistema}\\[1cm]
\textcopyright\ 2025 Walter Milstein\\
Todos los derechos reservados.\\[0.5cm]
Primera Edici\'on: 2025\\[1cm]
\textbf{Serie completa:}\\
Volumen 1: Fundamentos Te\'oricos (Caps. 1--5)\\
Volumen 2: Aplicaciones y Privacidad (Caps. 6--13)\\
Volumen 3: Herramientas Avanzadas y Ecosistema (Caps. 14--20)\\[1cm]
El c\'odigo incluido es de dominio p\'ublico y puede ser utilizado libremente.\\
Este libro es para uso personal y educativo.
\end{center}
\vfill
\cleardoublepage

% Usar con: \maketitle

% ============================================
% VOLUMEN 2 - PORTADA
% ============================================
\title{
    \Huge\textbf{Banco ZK} \\[0.5cm]
    \LARGE Criptografía y Privacidad en Blockchain \\[1cm]
    \huge\textbf{Volumen 3} \\[0.3cm]
    \Large Aplicaciones y Privacidad
}
\author{
    \Large Serie Técnica de Criptomonedas con Privacidad
}
\date{2025}

% ============================================
% INICIO DEL DOCUMENTO
% ============================================


% Portada
\maketitle


% --- DEDICATORIA ---
\thispagestyle{empty}
\null\vfill
\begin{flushright}
\textit{\large In memoriam Hackeris Euleri, Iacobi Rossini,\\[0.3cm]
ingenii cogitantis architecti,\\
veritatis technicae indefessi inquisitoris,\\
qui codicem non solum ut seriem praeceptorum intellexit,\\
sed ut formam ethicae, rigoris atque libertatis.\\[0.5cm]
Aspectus tuus criticus, fortitudo intellectualis\\
et amor scientiae\\
vivunt adhuc\\
in omni systemate recte excogitato\\
et in omni errore tempestive detecto.\\[1cm]
Walter Milstein, pater tuus}
\end{flushright}
\vfill
\cleardoublepage

% Página de copyright
\thispagestyle{empty}
\vspace*{\fill}
\begin{center}
\textit{Este volumen es continuación de ``Banco ZK: Volumen 1 - Fundamentos y Protocolos''}

\vspace{1cm}

\textbf{Bloques incluidos en este volumen:}\\[0.5cm]
Bloque 20: Herramientas y Ecosistema\\
(Bloques adicionales se agregarán en futuras ediciones)

\vspace{2cm}

\small
Documento generado para fines educativos.\\
La información contenida es de carácter técnico y académico.
\end{center}
\vspace*{\fill}
\clearpage

% Tabla de contenidos
\tableofcontents
\clearpage

% ============================================
% ABSTRACT Y PREFACIO - VOLUMEN 2
% Insertar después de \maketitle y antes de \tableofcontents
% ============================================

% Página de copyright y descripción
\thispagestyle{empty}
\vspace*{\fill}
\begin{center}
\textit{Este volumen es continuación de ``Banco ZK: Volumen 1 - Fundamentos y Protocolos''}

\vspace{1cm}

\textbf{Bloques incluidos en este volumen:}\\[0.5cm]
\begin{tabular}{rl}
Bloque 14: & Identidad y Credenciales Anónimas \\
Bloque 15: & Compliance vs Privacidad \\
Bloque 16: & Metadata y Network Privacy \\
Bloque 17: & Private Smart Contracts \\
Bloque 18: & Criptografía Avanzada \\
Bloque 19: & Ataques y Defensa \\
Bloque 20: & Herramientas y Ecosistema \\
\end{tabular}

\vspace{2cm}

\small
Documento generado para fines educativos.\\
La información contenida es de carácter técnico y académico.
\end{center}
\vspace*{\fill}
\clearpage

% ============================================
% ABSTRACT
% ============================================
\chapter*{Abstract}
\addcontentsline{toc}{chapter}{Abstract}

\begin{center}
\textbf{\Large Banco ZK: Volumen 2}\\[0.3cm]
\textbf{Aplicaciones Avanzadas y Ecosistema de Privacidad}\\[0.5cm]
\textit{Bloques 14--20 de la Serie Técnica}
\end{center}

\vspace{0.5cm}

Este segundo volumen de la serie \textbf{Banco ZK} presenta las aplicaciones avanzadas y el ecosistema completo necesario para implementar sistemas de privacidad en blockchain. Mientras el Volumen 1 estableció los fundamentos criptográficos y los protocolos base (Bloques 1--13), este volumen se enfoca en llevar esos conceptos a implementaciones prácticas y sistemas de producción.

\textbf{Alcance del volumen:} Cubrimos siete bloques temáticos organizados en tres partes:

\begin{itemize}
    \item \textbf{Parte I -- Aplicaciones Avanzadas} (Bloques 14--17): Identidad anónima, balance entre compliance y privacidad, protección de metadatos, y contratos inteligentes privados.
    
    \item \textbf{Parte II -- Criptografía y Seguridad} (Bloques 18--19): Primitivas criptográficas de frontera (FHE, MPC, witness encryption) y análisis exhaustivo de vectores de ataque con sus defensas.
    
    \item \textbf{Parte III -- Ecosistema} (Bloque 20): Stack completo de herramientas incluyendo wallets, relayers, infraestructura de proving, y guías de implementación.
\end{itemize}

\textbf{Metodología:} Cada bloque combina teoría rigurosa con implementaciones prácticas en Python, Solidity y Circom. Se incluyen más de 70 diagramas TikZ, 95 definiciones formales, y casos de estudio de protocolos reales como Zcash, Tornado Cash, Railgun, y Aztec.

\textbf{Audiencia:} Desarrolladores blockchain, investigadores en criptografía aplicada, arquitectos de sistemas de privacidad, y profesionales de seguridad interesados en tecnologías de conocimiento cero.

\textbf{Prerrequisitos:} Se asume familiaridad con los conceptos del Volumen 1, particularmente: curvas elípticas, funciones hash criptográficas, árboles de Merkle, y fundamentos de SNARKs/STARKs.

\vspace{0.5cm}

\noindent\textbf{Palabras clave:} Zero-Knowledge Proofs, Privacidad Blockchain, Anonymous Credentials, zkEVM, Compliance, MEV Protection, Relayer Networks, FHE, MPC

\clearpage

% ============================================
% PREFACIO
% ============================================
\chapter*{Prefacio}
\addcontentsline{toc}{chapter}{Prefacio}

La privacidad financiera es un derecho fundamental que la era digital ha erosionado sistemáticamente. Cada transacción bancaria, cada compra con tarjeta, cada transferencia deja un rastro permanente que puede ser analizado, vendido, o usado en nuestra contra. Las criptomonedas prometieron una alternativa, pero Bitcoin y Ethereum son, en realidad, los sistemas financieros más transparentes jamás creados: cada transacción es visible para siempre en un registro público global.

Este volumen representa la culminación práctica de la serie Banco ZK. Si el primer volumen construyó los cimientos teóricos---las matemáticas de curvas elípticas, los protocolos de conocimiento cero, las estructuras de datos criptográficas---este segundo volumen responde la pregunta: \textit{¿cómo construimos sistemas de privacidad que funcionen en el mundo real?}

\section*{El Desafío de la Privacidad Práctica}

La teoría criptográfica nos da herramientas poderosas: podemos probar que conocemos un secreto sin revelarlo, podemos demostrar que una transacción es válida sin mostrar los montos, podemos verificar identidad sin exponer datos personales. Pero convertir estas posibilidades teóricas en sistemas usables presenta desafíos únicos:

\begin{itemize}
    \item \textbf{Usabilidad vs Seguridad:} Los usuarios no deberían necesitar un doctorado en criptografía para proteger su privacidad.
    
    \item \textbf{Eficiencia vs Privacidad:} Las pruebas de conocimiento cero son computacionalmente costosas; optimizarlas sin comprometer garantías es un arte.
    
    \item \textbf{Compliance vs Anonimato:} Los reguladores exigen visibilidad; los usuarios demandan privacidad. El balance es delicado y geográficamente variable.
    
    \item \textbf{Descentralización vs Rendimiento:} Los sistemas centralizados son más rápidos pero introducen puntos únicos de fallo y confianza.
\end{itemize}

\section*{Estructura del Volumen}

Hemos organizado el material en una progresión lógica que va desde aplicaciones específicas hasta el ecosistema completo:

\textbf{Bloques 14--17} abordan aplicaciones concretas de privacidad. Comenzamos con \textit{identidad y credenciales anónimas}---cómo probar quiénes somos sin revelar más de lo necesario. Continuamos con el tenso equilibrio entre \textit{compliance y privacidad}, un tema cada vez más relevante dado el escrutinio regulatorio. Exploramos la \textit{privacidad de metadatos}---porque saber que Alice habló con Bob puede ser tan revelador como saber qué dijeron. Finalmente, examinamos los \textit{contratos inteligentes privados}, la frontera de la computación confidencial en blockchain.

\textbf{Bloques 18--19} profundizan en aspectos técnicos avanzados. La \textit{criptografía avanzada} presenta primitivas como cifrado homomórfico y computación multipartita que complementan las pruebas de conocimiento cero. El bloque de \textit{ataques y defensa} es quizás el más importante: ningún sistema es seguro si no entendemos cómo puede ser comprometido.

\textbf{Bloque 20} integra todo en un \textit{ecosistema funcional}. Desde wallets que gestionan notas cifradas hasta redes de relayers descentralizados, desde exploradores que no rastrean hasta infraestructura de proving escalable.

\section*{Enfoque Práctico}

A diferencia de textos puramente teóricos, cada concepto se acompaña de implementación. Incluimos:

\begin{itemize}
    \item \textbf{Código Python} para prototipos y simulaciones
    \item \textbf{Contratos Solidity} listos para auditoría
    \item \textbf{Circuitos Circom} con optimizaciones documentadas
    \item \textbf{Diagramas de arquitectura} para sistemas completos
    \item \textbf{Análisis de trade-offs} con métricas concretas
\end{itemize}

Los ejemplos no son juguetes académicos sino extractos simplificados de sistemas en producción. Cuando discutimos Railgun, Tornado Cash, o Zcash, lo hacemos con conocimiento de sus implementaciones reales.

\section*{Nota sobre Responsabilidad}

Las tecnologías de privacidad son herramientas neutrales. Un cuchillo puede preparar comida o causar daño; la responsabilidad recae en quien lo usa. Los sistemas que describimos pueden proteger disidentes políticos, preservar la privacidad financiera legítima, y habilitar nuevos modelos de negocio. También pueden ser mal utilizados.

Abogamos por privacidad \textit{con accountability}: sistemas donde los usuarios honestos permanecen anónimos pero donde existen mecanismos para casos extremos (mediante viewing keys selectivas, compliance proofs, o auditorías consensuadas). El Bloque 15 dedica especial atención a este balance.

\section*{Cómo Usar Este Volumen}

\textbf{Lectura secuencial:} Recomendada para comprensión completa. Cada bloque construye sobre conceptos anteriores.

\textbf{Referencia selectiva:} El glosario expandido y los índices permiten consultar temas específicos. Cada bloque es relativamente autocontenido para lectores con base sólida.

\textbf{Implementación práctica:} El código puede extraerse y adaptarse. Recomendamos comenzar con los ejemplos del Bloque 20 y retroceder según se necesiten conceptos específicos.

\textbf{Profundización:} Las referencias al final apuntan a papers originales, documentación de protocolos, y recursos comunitarios para estudio avanzado.

\section*{Agradecimientos}

Este trabajo se construye sobre décadas de investigación en criptografía. Debemos reconocimiento especial a los equipos de Zcash (por pionear zk-SNARKs en producción), Tornado Cash (por demostrar privacidad en Ethereum), y a la comunidad de investigadores que continúan avanzando el estado del arte.

La privacidad no es un lujo ni una preferencia: es la base de la autonomía individual en la era digital. Esperamos que este volumen contribuya a un ecosistema donde la privacidad sea accesible para todos.

\vspace{1cm}
\begin{flushright}
\textit{Los autores}\\
\textit{Enero 2025}
\end{flushright}

\clearpage

% ============================================
% INTRODUCCIÓN
% ============================================
\chapter*{Introducción}
\addcontentsline{toc}{chapter}{Introducción}

\section*{Contexto: De la Teoría a la Práctica}

El Volumen 1 de esta serie estableció los fundamentos matemáticos y criptográficos necesarios para entender los sistemas de privacidad en blockchain:

\begin{itemize}
    \item \textbf{Bloques 1--3:} Fundamentos de criptografía, teoría de grupos, curvas elípticas
    \item \textbf{Bloques 4--6:} Estructuras de datos (Merkle trees), compromisos criptográficos
    \item \textbf{Bloques 7--9:} Firmas digitales, ring signatures, stealth addresses
    \item \textbf{Bloques 10--13:} Zero-knowledge proofs: SNARKs, STARKs, circuitos aritméticos
\end{itemize}

Con esos cimientos, este volumen aborda la pregunta práctica: \textit{¿cómo construimos sistemas completos de privacidad?}

\section*{El Stack de Privacidad}

Un sistema de privacidad funcional no es solo un circuito ZK elegante o un contrato inteligente bien diseñado. Es un \textbf{stack completo} donde cada capa debe preservar las garantías de privacidad:

\begin{center}
\begin{tikzpicture}[
    layer/.style={rectangle, draw, minimum width=11cm, minimum height=0.9cm, align=center, font=\small},
    arrow/.style={->, thick, >=stealth}
]
    \node[layer, fill=red!10] (user) at (0, 5) {\textbf{Usuario:} Intención de transacción privada};
    \node[layer, fill=orange!10] (wallet) at (0, 4) {\textbf{Wallet:} Gestión de claves, notas, backups};
    \node[layer, fill=yellow!10] (proving) at (0, 3) {\textbf{Proving:} Generación de ZK proofs};
    \node[layer, fill=green!10] (network) at (0, 2) {\textbf{Red:} Transmisión anónima (Tor, relayers)};
    \node[layer, fill=cyan!10] (rpc) at (0, 1) {\textbf{RPC:} Protección MEV, mempool privada};
    \node[layer, fill=blue!10] (chain) at (0, 0) {\textbf{Chain:} Contratos ZK, verificación on-chain};
    
    \draw[arrow] (user) -- (wallet);
    \draw[arrow] (wallet) -- (proving);
    \draw[arrow] (proving) -- (network);
    \draw[arrow] (network) -- (rpc);
    \draw[arrow] (rpc) -- (chain);
    
    % Amenazas laterales
    \node[font=\scriptsize, text=red!70!black, align=right] at (-6.5, 4) {Phishing, malware};
    \node[font=\scriptsize, text=red!70!black, align=right] at (-6.5, 3) {Side-channels};
    \node[font=\scriptsize, text=red!70!black, align=right] at (-6.5, 2) {Traffic analysis};
    \node[font=\scriptsize, text=red!70!black, align=right] at (-6.5, 1) {MEV, frontrunning};
    \node[font=\scriptsize, text=red!70!black, align=right] at (-6.5, 0) {Graph analysis};
    
    \draw[->, dashed, red!50] (-5.5, 4) -- (-5.5, 4);
\end{tikzpicture}
\end{center}

Una filtración en \textit{cualquier} capa compromete la privacidad completa. Este volumen aborda cada capa sistemáticamente.

\section*{Mapa de Contenidos}

\subsection*{Parte I: Aplicaciones Avanzadas de Privacidad}

\textbf{Bloque 14: Identidad y Credenciales Anónimas}

La identidad es el caso de uso más poderoso y más delicado de ZK. Exploramos:
\begin{itemize}
    \item Anonymous credentials: probar atributos sin revelar identidad
    \item Selective disclosure: mostrar solo lo necesario
    \item Group signatures: firmar como miembro de un grupo
    \item DIDs y Verifiable Credentials con privacidad
    \item Soulbound tokens que preservan anonimato
    \item Proof of humanity sin doxxing
\end{itemize}

\textbf{Bloque 15: Compliance vs Privacidad}

El elefante en la habitación: ¿cómo satisfacemos reguladores sin destruir la privacidad?
\begin{itemize}
    \item El espectro de privacidad: de transparente a opaco
    \item Viewing keys y sus modelos de confianza
    \item Travel Rule compliance con ZK proofs
    \item Sanctions screening privacy-preserving
    \item Caso de estudio: Tornado Cash y OFAC
    \item Arquitecturas de compliance selectivo
\end{itemize}

\textbf{Bloque 16: Metadata y Network Privacy}

Los metadatos pueden ser más reveladores que el contenido:
\begin{itemize}
    \item Timing analysis y correlación temporal
    \item Traffic analysis en redes P2P
    \item IP address leakage y protecciones
    \item Tor, I2P, mixnets: comparativa
    \item Dummy transactions y cover traffic
    \item Fingerprinting de wallets
\end{itemize}

\textbf{Bloque 17: Private Smart Contracts}

La frontera de privacidad programable:
\begin{itemize}
    \item zkEVM: taxonomía (Type 1--4)
    \item Private state en contratos
    \item Aztec, Aleo, Mina: arquitecturas comparadas
    \item Composabilidad entre contratos privados
    \item Private DeFi: AMMs, lending, derivatives
    \item Oráculos con privacidad
\end{itemize}

\subsection*{Parte II: Criptografía Avanzada y Seguridad}

\textbf{Bloque 18: Criptografía Avanzada}

Más allá de ZK, otras primitivas habilitan privacidad:
\begin{itemize}
    \item Fully Homomorphic Encryption (FHE): cómputo sobre cifrado
    \item Multi-Party Computation (MPC): cómputo distribuido seguro
    \item Oblivious Transfer y Private Set Intersection
    \item Witness Encryption: cifrado condicional a pruebas
    \item Post-quantum considerations
    \item Combinando técnicas para defensa en profundidad
\end{itemize}

\textbf{Bloque 19: Ataques y Defensa}

Conocer al enemigo es esencial:
\begin{itemize}
    \item Taxonomy de ataques a privacidad
    \item Deanonymization: heurísticas y machine learning
    \item Ataques al anonymity set: Sybil, poisoning
    \item Side-channels en proving: timing, power, cache
    \item Vulnerabilidades específicas de contratos privados
    \item Incident response y recovery
    \item Diseño de sistemas resilientes
\end{itemize}

\subsection*{Parte III: Herramientas y Ecosistema}

\textbf{Bloque 20: Herramientas y Ecosistema}

El stack completo para implementación:
\begin{itemize}
    \item Exploradores privacy-aware: sin tracking, con Tor
    \item Wallets ZK: Railgun, Umbra, Zcash wallets
    \item RPCs privados: Flashbots Protect, MEV Blocker
    \item Redes de relayers descentralizados
    \item Gestión de notas y backups seguros
    \item Infraestructura de proving: client vs server vs hybrid
    \item Hardware wallets y sus limitaciones
    \item Integración del stack Banco ZK completo
\end{itemize}

\section*{Convenciones del Documento}

A lo largo del volumen utilizamos las siguientes convenciones:

\begin{keyidea}
Las cajas \textbf{Idea Clave} resaltan conceptos fundamentales que el lector debe interiorizar.
\end{keyidea}

\begin{bancozk}
Las cajas \textbf{Aplicación en Banco ZK} muestran cómo cada concepto se integra en el sistema Banco ZK que construimos a lo largo de la serie.
\end{bancozk}

\begin{warning}
Las cajas \textbf{Advertencia} señalan errores comunes, vulnerabilidades, o malentendidos peligrosos.
\end{warning}

\begin{note}
Las cajas \textbf{Nota} proporcionan información adicional, contexto histórico, o detalles de implementación.
\end{note}

\textbf{Código:} Los listados de código incluyen Python (prototipos), Solidity (contratos), y Circom (circuitos). Todo el código es funcional y extraíble.

\textbf{Diagramas:} Los diagramas TikZ ilustran arquitecturas, flujos de datos, y relaciones entre componentes.

\textbf{Definiciones:} Los términos técnicos se definen formalmente en su primera aparición y se compilan en el glosario final.

\section*{Prerrequisitos}

Asumimos que el lector ha cubierto el Volumen 1 o tiene conocimiento equivalente de:

\begin{itemize}
    \item Aritmética modular y teoría de grupos finitos
    \item Curvas elípticas: operaciones, pairings, curvas comunes (BN254, BLS12-381)
    \item Funciones hash: propiedades, Poseidon, Pedersen
    \item Árboles de Merkle y pruebas de inclusión
    \item Compromisos criptográficos: Pedersen, KZG
    \item Fundamentos de SNARKs: R1CS, trusted setup, Groth16
    \item Fundamentos de STARKs: FRI, AIR
    \item Modelo UTXO vs account-based
    \item Nullifiers y su rol en prevenir doble gasto
\end{itemize}

Lectores sin esta base pueden consultar el Volumen 1 o las referencias introductorias en la bibliografía.

\section*{Objetivos de Aprendizaje}

Al completar este volumen, el lector será capaz de:

\begin{enumerate}
    \item \textbf{Diseñar} sistemas de identidad anónima con selective disclosure
    \item \textbf{Evaluar} trade-offs entre compliance y privacidad
    \item \textbf{Implementar} protecciones contra análisis de metadatos
    \item \textbf{Comprender} arquitecturas de contratos inteligentes privados
    \item \textbf{Aplicar} criptografía avanzada (FHE, MPC) donde ZK no es suficiente
    \item \textbf{Identificar} vectores de ataque y diseñar defensas apropiadas
    \item \textbf{Construir} el stack completo de un sistema de privacidad funcional
    \item \textbf{Operar} infraestructura de privacidad en producción
\end{enumerate}

\section*{Comenzamos}

La privacidad no es un feature que se agrega al final. Es una propiedad fundamental que debe diseñarse desde el inicio y preservarse en cada capa del sistema. 

Este volumen proporciona las herramientas conceptuales y prácticas para lograrlo. Comencemos con el Bloque 14: cómo probar quiénes somos sin revelar más de lo estrictamente necesario.

\cleardoublepage

% ============================================
% TABLA DE CONTENIDOS
% ============================================
\tableofcontents
\cleardoublepage
% ============================================
% BLOQUE 14: IDENTIDAD Y CREDENCIALES ANÓNIMAS - VERSIÓN EXPANDIDA
% Parte 1: Introducción y Anonymous Credentials
% Objetivo: ~60 páginas (~3600 líneas)
% ============================================

\chapter{Identidad y Credenciales Anónimas}
\label{ch:bloque14}

\begin{center}
\textit{Bloque 14 de 20}
\end{center}

\section*{Introducción}

La identidad digital es una de las aplicaciones más prometedoras de la criptografía de conocimiento cero. Este bloque explora cómo demostrar propiedades sobre uno mismo sin revelar información innecesaria: desde credenciales anónimas hasta pruebas de humanidad, desde verificación de edad hasta KYC preservando privacidad.

\begin{keyidea}
\textbf{El dilema de la identidad digital:}

\textbf{Modelo tradicional:}
\begin{itemize}
    \item Cada servicio recolecta datos completos
    \item Datos replicados en múltiples bases de datos
    \item Riesgo de data breaches multiplicado
    \item Usuario no tiene control
\end{itemize}

\textbf{Modelo con ZK:}
\begin{itemize}
    \item Datos permanecen con el usuario
    \item Solo se comparte lo necesario (selective disclosure)
    \item Pruebas verificables sin revelar fuente
    \item Usuario mantiene soberanía
\end{itemize}
\end{keyidea}

\begin{mathnote}[Timeline de Identidad Digital]
\begin{itemize}
    \item \textbf{1985:} David Chaum propone anonymous credentials
    \item \textbf{2001:} IBM desarrolla Idemix (Identity Mixer)
    \item \textbf{2004:} Microsoft desarrolla U-Prove
    \item \textbf{2016:} Semaphore (primer protocolo ZK identity en blockchain)
    \item \textbf{2019:} W3C publica estándar de DIDs
    \item \textbf{2022:} Worldcoin lanza con ZK proofs
    \item \textbf{2023:} Soulbound Tokens + ZK emerge
    \item \textbf{2024:} EU eIDAS 2.0 incluye consideraciones ZK
\end{itemize}
\end{mathnote}

%===========================================
\section{Anonymous Credentials}
\label{sec:anonymous-credentials}
%===========================================

\subsection{Concepto Fundamental}

\begin{definition}[Anonymous Credential]
Sistema criptográfico que permite a un usuario obtener credenciales de un emisor y presentarlas a verificadores sin revelar su identidad ni permitir que emisor y verificador coludan para rastrear al usuario.

\textbf{Participantes:}
\begin{enumerate}
    \item \textbf{Issuer (Emisor):} Entidad que emite credenciales (gobierno, universidad, etc.)
    \item \textbf{Holder (Portador):} Usuario que posee la credencial
    \item \textbf{Verifier (Verificador):} Entidad que verifica la credencial
\end{enumerate}

\textbf{Propiedades fundamentales:}
\begin{itemize}
    \item \textbf{Unforgeability:} No se pueden crear credenciales falsas
    \item \textbf{Anonymity:} Verificador no puede identificar al holder
    \item \textbf{Unlinkability:} Múltiples presentaciones no son vinculables
    \item \textbf{Issuer-Verifier unlinkability:} Emisor y verificador no pueden coludir
\end{itemize}
\end{definition}

\begin{figure}[htbp]
\centering
\begin{tikzpicture}[
    node distance=2cm,
    entity/.style={rectangle, draw, minimum width=2.5cm, minimum height=1.2cm, align=center},
    arrow/.style={->, thick},
    darrow/.style={<->, thick, dashed}
]
    % Entidades
    \node[entity, fill=blue!20] (issuer) at (0, 0) {Issuer\\(Gobierno)};
    \node[entity, fill=green!20] (holder) at (5, 0) {Holder\\(Usuario)};
    \node[entity, fill=orange!20] (verifier) at (10, 0) {Verifier\\(Servicio)};
    
    % Flujo de emisión
    \draw[arrow] (issuer) -- node[above] {1. Credencial} node[below, font=\small] {blindada} (holder);
    
    % Flujo de presentación
    \draw[arrow] (holder) -- node[above] {2. Proof} node[below, font=\small] {anónimo} (verifier);
    
    % No colusión
    \draw[darrow, red] (issuer) to[bend left=40] node[above, font=\small] {No pueden} node[below, font=\small] {vincular} (verifier);
    
    % Atributos
    \node at (5, -2) [draw, rectangle, fill=yellow!20, minimum width=4cm] {
        Atributos: nombre, edad, país, ...
    };
    
    \draw[arrow, dashed] (5, -1.5) -- (holder);
\end{tikzpicture}
\caption{Modelo de anonymous credentials: el holder obtiene una credencial del issuer y la presenta al verifier sin que puedan coludir.}
\label{fig:anon-cred-model}
\end{figure}

\subsection{Propiedades Matemáticas}

\begin{mathnote}[Formalización de Anonymous Credentials]
Un sistema de anonymous credentials consiste en:

\textbf{Algoritmos:}
\begin{enumerate}
    \item $\text{Setup}(1^\lambda) \to (pp)$: Genera parámetros públicos
    \item $\text{IssuerKeyGen}(pp) \to (isk, ipk)$: Genera claves del emisor
    \item $\text{UserKeyGen}(pp) \to (usk, upk)$: Genera claves del usuario
    \item $\text{Issue}(isk, upk, \text{attrs}) \to \sigma$: Emite credencial
    \item $\text{Present}(usk, \sigma, \text{disclosed}) \to \pi$: Presenta prueba
    \item $\text{Verify}(ipk, \pi, \text{disclosed}) \to \{0,1\}$: Verifica prueba
\end{enumerate}

\textbf{Propiedades de seguridad:}
\begin{itemize}
    \item \textbf{Correctness:} Credenciales válidas siempre verifican
    \item \textbf{Unforgeability:} $\Pr[\text{Adversario crea } \sigma \text{ válido sin } isk] < \text{negl}(\lambda)$
    \item \textbf{Anonymity:} $\forall \sigma_0, \sigma_1: \text{Present}(\cdot, \sigma_0, \cdot) \approx_c \text{Present}(\cdot, \sigma_1, \cdot)$
    \item \textbf{Unlinkability:} Dos presentaciones de la misma credencial son indistinguibles de dos credenciales diferentes
\end{itemize}
\end{mathnote}

\subsection{Implementaciones Históricas}

\subsubsection{Idemix (IBM Identity Mixer)}

\begin{definition}[Idemix]
Sistema de anonymous credentials desarrollado por IBM Research:
\begin{itemize}
    \item \textbf{Base:} CL signatures (Camenisch-Lysyanskaya)
    \item \textbf{Features:} Multi-show, selective disclosure, predicates
    \item \textbf{Usado en:} Hyperledger Indy/Aries
\end{itemize}
\end{definition}

\begin{mathnote}[CL Signatures (Simplificado)]
\textbf{Setup:}
\begin{itemize}
    \item RSA modulus $n = pq$
    \item Generadores $R_1, ..., R_L, S, Z \in QR_n$ (residuos cuadráticos)
    \item Public key: $(n, R_1, ..., R_L, S, Z)$
    \item Secret key: $p$ (factor de $n$)
\end{itemize}

\textbf{Firma de atributos $(m_1, ..., m_L)$:}
\[
A = \left( \frac{Z}{R_1^{m_1} \cdot ... \cdot R_L^{m_L} \cdot S^v} \right)^{1/e} \mod n
\]

Donde $e$ es primo, $v$ es randomness.

\textbf{Verificación:}
\[
Z \stackrel{?}{=} A^e \cdot R_1^{m_1} \cdot ... \cdot R_L^{m_L} \cdot S^v \mod n
\]

\textbf{ZK Presentation:}
El holder demuestra conocimiento de $(A, e, v, m_1, ..., m_L)$ sin revelarlos, usando protocolos Sigma.
\end{mathnote}

\subsubsection{U-Prove (Microsoft)}

\begin{definition}[U-Prove]
Sistema de anonymous credentials de Microsoft:
\begin{itemize}
    \item \textbf{Base:} Blind signatures sobre curvas elípticas
    \item \textbf{Característica:} One-show tokens (uso único)
    \item \textbf{Ventaja:} Más eficiente que Idemix
    \item \textbf{Desventaja:} Tokens de un solo uso
\end{itemize}
\end{definition}

\subsubsection{BBS+ Signatures}

\begin{definition}[BBS+ Signatures]
Esquema de firma moderna para anonymous credentials:
\begin{itemize}
    \item \textbf{Base:} Pairings sobre curvas elípticas
    \item \textbf{Ventaja:} Firmas cortas, ZK eficiente
    \item \textbf{Estándar:} IETF draft en progreso
    \item \textbf{Usado en:} Verifiable Credentials modernos
\end{itemize}
\end{definition}

\begin{lstlisting}[language=Python, caption={Implementación conceptual de BBS+ Signatures}]
from py_ecc.bls12_381 import (
    G1, G2, multiply, add, pairing, 
    curve_order, neg
)
from hashlib import sha256
import secrets

class BBSPlusSignature:
    """
    BBS+ Signatures para anonymous credentials.
    Basado en: https://identity.foundation/bbs-signature/
    """
    
    def __init__(self, num_messages: int):
        self.L = num_messages
        self.generators = self._generate_generators()
    
    def _generate_generators(self) -> list:
        """Generar L+2 generadores independientes en G1."""
        generators = []
        for i in range(self.L + 2):
            # Hash-to-curve (simplificado)
            seed = sha256(f"BBS_GENERATOR_{i}".encode()).digest()
            scalar = int.from_bytes(seed, 'big') % curve_order
            generators.append(multiply(G1, scalar))
        return generators
    
    def keygen(self):
        """Generar par de claves."""
        # Secret key
        sk = secrets.randbelow(curve_order)
        
        # Public key: W = sk * G2
        pk = multiply(G2, sk)
        
        return (sk, pk)
    
    def sign(self, sk: int, messages: list) -> tuple:
        """
        Firmar lista de mensajes.
        
        Args:
            sk: Secret key del issuer
            messages: Lista de L mensajes (como enteros)
        
        Returns:
            Signature (A, e, s)
        """
        assert len(messages) == self.L
        
        # Generar randomness
        e = secrets.randbelow(curve_order)
        s = secrets.randbelow(curve_order)
        
        # Calcular B = g1 + s*h0 + sum(m_i * h_i)
        # donde g1, h0, h1, ..., hL son los generadores
        g1, h0 = self.generators[0], self.generators[1]
        
        B = add(g1, multiply(h0, s))
        for i, m in enumerate(messages):
            hi = self.generators[i + 2]
            B = add(B, multiply(hi, m % curve_order))
        
        # A = B * (1 / (sk + e))
        inv = pow(sk + e, -1, curve_order)
        A = multiply(B, inv)
        
        return (A, e, s)
    
    def verify(self, pk, signature: tuple, messages: list) -> bool:
        """
        Verificar firma.
        
        Args:
            pk: Public key del issuer
            signature: (A, e, s)
            messages: Lista de mensajes
        
        Returns:
            True si válida
        """
        A, e, s = signature
        
        # Reconstruir B
        g1, h0 = self.generators[0], self.generators[1]
        B = add(g1, multiply(h0, s))
        for i, m in enumerate(messages):
            hi = self.generators[i + 2]
            B = add(B, multiply(hi, m % curve_order))
        
        # Verificar: e(A, W + e*G2) = e(B, G2)
        # Equivalente: e(A, W) * e(A, e*G2) = e(B, G2)
        # Equivalente: e(A, W) * e(e*A, G2) = e(B, G2)
        
        W_plus_eG2 = add(pk, multiply(G2, e))
        
        lhs = pairing(W_plus_eG2, A)
        rhs = pairing(G2, B)
        
        return lhs == rhs
    
    def create_proof(
        self, 
        signature: tuple, 
        messages: list,
        disclosed_indices: list,
        nonce: bytes
    ) -> dict:
        """
        Crear ZK proof de posesión de credencial con selective disclosure.
        
        Args:
            signature: (A, e, s)
            messages: Todos los mensajes
            disclosed_indices: Índices de mensajes a revelar
            nonce: Challenge del verificador
        
        Returns:
            Proof que revela solo mensajes en disclosed_indices
        """
        A, e, s = signature
        hidden_indices = [i for i in range(self.L) if i not in disclosed_indices]
        
        # Randomizar A para unlinkability
        r1 = secrets.randbelow(curve_order)
        r2 = secrets.randbelow(curve_order)
        
        # A' = A * r1
        A_prime = multiply(A, r1)
        
        # A_bar = A' * (-e) + B * r1
        # (Esto oculta la firma original)
        
        # Generar commitments para protocolo Sigma
        # ... (protocolo completo omitido por brevedad)
        
        disclosed_messages = {i: messages[i] for i in disclosed_indices}
        
        return {
            'A_prime': A_prime,
            'disclosed': disclosed_messages,
            'proof': b'...',  # Prueba Sigma completa
        }
    
    def verify_proof(
        self, 
        pk, 
        proof: dict, 
        disclosed_indices: list,
        nonce: bytes
    ) -> bool:
        """
        Verificar ZK proof.
        
        Args:
            pk: Public key del issuer
            proof: Proof generado por create_proof
            disclosed_indices: Qué mensajes se revelaron
            nonce: Challenge usado
        
        Returns:
            True si proof es válido
        """
        # Verificar que proof es consistente
        # y que el prover conoce una firma válida
        # sobre mensajes que incluyen los disclosed
        
        # ... (verificación completa omitida)
        
        return True  # Placeholder


def bbs_demo():
    print("=== BBS+ Signatures Demo ===\n")
    
    # Setup para 4 atributos
    bbs = BBSPlusSignature(num_messages=4)
    
    # Issuer genera claves
    sk, pk = bbs.keygen()
    print("Issuer keys generated")
    
    # Atributos de una credencial de identidad
    # (En práctica, serían hashes de strings)
    messages = [
        123456789,      # ID number
        19900115,       # Birth date (YYYYMMDD)
        1,              # Country code
        0,              # Additional flag
    ]
    
    # Issuer firma
    signature = bbs.sign(sk, messages)
    print(f"Credential signed")
    
    # Verificar firma completa
    is_valid = bbs.verify(pk, signature, messages)
    print(f"Full verification: {is_valid}")
    
    # Crear proof revelando solo país (índice 2)
    proof = bbs.create_proof(
        signature=signature,
        messages=messages,
        disclosed_indices=[2],  # Solo revelar país
        nonce=b"verifier_challenge_123"
    )
    
    print(f"\nSelective disclosure proof created")
    print(f"Revealed: country = {proof['disclosed']}")
    print("Hidden: ID, birth date, flag")
    
    # Verificador comprueba
    is_proof_valid = bbs.verify_proof(
        pk=pk,
        proof=proof,
        disclosed_indices=[2],
        nonce=b"verifier_challenge_123"
    )
    print(f"Proof verification: {is_proof_valid}")

bbs_demo()
\end{lstlisting}

\subsection{Comparación de Sistemas}

\begin{table}[htbp]
\centering
\small
\begin{tabular}{lccccc}
\toprule
\textbf{Sistema} & \textbf{Multi-show} & \textbf{Selective} & \textbf{Predicates} & \textbf{Tamaño} & \textbf{Verificación} \\
\midrule
Idemix (CL) & Sí & Sí & Sí & $\sim$4 KB & $\sim$50 ms \\
U-Prove & No & Sí & Limitado & $\sim$500 B & $\sim$5 ms \\
BBS+ & Sí & Sí & Con ZK & $\sim$400 B & $\sim$10 ms \\
SNARK-based & Sí & Sí & Arbitrario & $\sim$200 B & $\sim$5 ms \\
\bottomrule
\end{tabular}
\caption{Comparación de sistemas de anonymous credentials}
\end{table}

%===========================================
\section{Selective Disclosure}
\label{sec:selective-disclosure}
%===========================================

\subsection{Concepto y Motivación}

\begin{definition}[Selective Disclosure]
Capacidad de revelar solo un subconjunto de atributos de una credencial, manteniendo el resto oculto pero demostrando que la credencial completa es válida.
\end{definition}

\begin{intuicion}
Imagina tu licencia de conducir:
\begin{itemize}
    \item Contiene: nombre, foto, dirección, fecha nacimiento, tipo de licencia, restricciones
    \item Para alquilar auto: Necesitan ver tipo de licencia
    \item Para comprar alcohol: Solo necesitan saber edad $\geq$ 21
    \item Para votar: Solo necesitan verificar que eres ciudadano
\end{itemize}

Con selective disclosure, muestras SOLO lo necesario para cada caso.
\end{intuicion}

\begin{figure}[htbp]
\centering
\begin{tikzpicture}[
    node distance=0.5cm,
    attr/.style={rectangle, draw, minimum width=2.5cm, minimum height=0.6cm},
    hidden/.style={rectangle, draw, minimum width=2.5cm, minimum height=0.6cm, fill=gray!50},
    shown/.style={rectangle, draw, minimum width=2.5cm, minimum height=0.6cm, fill=green!30}
]
    % Credencial completa
    \node at (-4, 3) {\textbf{Credencial Completa}};
    
    \node[attr, fill=blue!20] (name) at (-4, 2) {Nombre: Alice};
    \node[attr, fill=blue!20] (dob) at (-4, 1.3) {Nacimiento: 1990};
    \node[attr, fill=blue!20] (country) at (-4, 0.6) {País: España};
    \node[attr, fill=blue!20] (id) at (-4, -0.1) {ID: 12345678};
    \node[attr, fill=blue!20] (addr) at (-4, -0.8) {Dirección: ...};
    
    % Caso 1: Verificar edad
    \node at (2, 3) {\textbf{Verificar Edad}};
    
    \node[hidden] at (2, 2) {Nombre: ???};
    \node[shown] at (2, 1.3) {Edad $\geq$ 18: ✓};
    \node[hidden] at (2, 0.6) {País: ???};
    \node[hidden] at (2, -0.1) {ID: ???};
    \node[hidden] at (2, -0.8) {Dirección: ???};
    
    % Caso 2: Verificar nacionalidad
    \node at (7, 3) {\textbf{Verificar País}};
    
    \node[hidden] at (7, 2) {Nombre: ???};
    \node[hidden] at (7, 1.3) {Nacimiento: ???};
    \node[shown] at (7, 0.6) {País: España};
    \node[hidden] at (7, -0.1) {ID: ???};
    \node[hidden] at (7, -0.8) {Dirección: ???};
    
    % Flechas
    \draw[->, thick] (-2, 1) -- (0, 1.3);
    \draw[->, thick] (-2, 0.5) -- (5, 0.6);
\end{tikzpicture}
\caption{Selective disclosure: diferentes verificadores ven diferentes subconjuntos de atributos}
\end{figure}

\subsection{Técnicas de Implementación}

\subsubsection{Con BBS+ Signatures}

\begin{lstlisting}[language=Python, caption={Selective disclosure con estructura Merkle}]
from hashlib import sha256
from typing import List, Dict, Optional
import json

class SelectiveDisclosureCredential:
    """
    Credencial con selective disclosure usando Merkle tree.
    Cada atributo es una hoja del árbol.
    """
    
    def __init__(self, attributes: Dict[str, any]):
        self.attributes = attributes
        self.attribute_names = sorted(attributes.keys())
        self.tree = self._build_merkle_tree()
        self.root = self.tree[-1][0] if self.tree else None
    
    def _hash(self, data: bytes) -> bytes:
        return sha256(data).digest()
    
    def _hash_attribute(self, name: str, value: any) -> bytes:
        """Hash de un atributo individual."""
        data = json.dumps({name: value}, sort_keys=True).encode()
        return self._hash(data)
    
    def _build_merkle_tree(self) -> List[List[bytes]]:
        """Construir Merkle tree de atributos."""
        # Hojas: hash de cada atributo
        leaves = [
            self._hash_attribute(name, self.attributes[name])
            for name in self.attribute_names
        ]
        
        # Pad a potencia de 2
        while len(leaves) & (len(leaves) - 1) != 0:
            leaves.append(self._hash(b"PADDING"))
        
        tree = [leaves]
        
        # Construir niveles
        while len(tree[-1]) > 1:
            level = tree[-1]
            new_level = []
            for i in range(0, len(level), 2):
                combined = level[i] + level[i + 1]
                new_level.append(self._hash(combined))
            tree.append(new_level)
        
        return tree
    
    def get_merkle_proof(self, attribute_name: str) -> Dict:
        """
        Obtener proof de Merkle para un atributo específico.
        """
        try:
            idx = self.attribute_names.index(attribute_name)
        except ValueError:
            raise ValueError(f"Attribute {attribute_name} not found")
        
        proof_path = []
        current_idx = idx
        
        for level in self.tree[:-1]:
            # Hermano
            if current_idx % 2 == 0:
                sibling_idx = current_idx + 1
                position = 'right'
            else:
                sibling_idx = current_idx - 1
                position = 'left'
            
            if sibling_idx < len(level):
                proof_path.append({
                    'hash': level[sibling_idx].hex(),
                    'position': position
                })
            
            current_idx //= 2
        
        return {
            'attribute_name': attribute_name,
            'attribute_value': self.attributes[attribute_name],
            'leaf_hash': self.tree[0][idx].hex(),
            'proof_path': proof_path,
            'root': self.root.hex()
        }
    
    def create_selective_disclosure(
        self, 
        disclosed_attributes: List[str]
    ) -> Dict:
        """
        Crear prueba de selective disclosure.
        Revela solo los atributos especificados.
        """
        proofs = {}
        for attr in disclosed_attributes:
            if attr in self.attributes:
                proofs[attr] = self.get_merkle_proof(attr)
        
        return {
            'root': self.root.hex(),
            'disclosed': proofs,
            'num_hidden': len(self.attributes) - len(disclosed_attributes)
        }
    
    @staticmethod
    def verify_selective_disclosure(disclosure: Dict) -> bool:
        """
        Verificar que los atributos revelados son parte de
        una credencial válida (con el root dado).
        """
        expected_root = bytes.fromhex(disclosure['root'])
        
        for attr_name, proof in disclosure['disclosed'].items():
            # Verificar cada atributo revelado
            current_hash = bytes.fromhex(proof['leaf_hash'])
            
            # Recomputar hacia arriba
            for step in proof['proof_path']:
                sibling = bytes.fromhex(step['hash'])
                if step['position'] == 'right':
                    combined = current_hash + sibling
                else:
                    combined = sibling + current_hash
                current_hash = sha256(combined).digest()
            
            if current_hash != expected_root:
                return False
        
        return True


def selective_disclosure_demo():
    print("=== Selective Disclosure Demo ===\n")
    
    # Crear credencial con múltiples atributos
    credential = SelectiveDisclosureCredential({
        'name': 'Alice Smith',
        'birth_date': '1990-05-15',
        'country': 'Spain',
        'id_number': '12345678X',
        'address': '123 Main St, Madrid',
        'email': 'alice@example.com'
    })
    
    print(f"Credential root: {credential.root.hex()[:16]}...")
    print(f"Total attributes: {len(credential.attributes)}")
    
    # Caso 1: Verificar solo edad (birth_date)
    print("\n--- Caso 1: Verificación de edad ---")
    disclosure1 = credential.create_selective_disclosure(['birth_date'])
    print(f"Revealed: birth_date = {disclosure1['disclosed']['birth_date']['attribute_value']}")
    print(f"Hidden: {disclosure1['num_hidden']} attributes")
    
    is_valid1 = SelectiveDisclosureCredential.verify_selective_disclosure(disclosure1)
    print(f"Verification: {is_valid1}")
    
    # Caso 2: Verificar país y email
    print("\n--- Caso 2: País y email ---")
    disclosure2 = credential.create_selective_disclosure(['country', 'email'])
    for attr, proof in disclosure2['disclosed'].items():
        print(f"Revealed: {attr} = {proof['attribute_value']}")
    print(f"Hidden: {disclosure2['num_hidden']} attributes")
    
    is_valid2 = SelectiveDisclosureCredential.verify_selective_disclosure(disclosure2)
    print(f"Verification: {is_valid2}")
    
    # El verificador NO puede ver los atributos ocultos
    print("\n--- Información oculta ---")
    print("Verifier cannot determine:")
    print("  - Name")
    print("  - ID number")
    print("  - Address")
    print("Even though they verify the credential is valid!")

selective_disclosure_demo()
\end{lstlisting}

\subsubsection{Predicados sobre Atributos}

\begin{definition}[Predicate Proof]
Prueba de que un atributo oculto satisface cierta condición sin revelar el valor exacto:
\begin{itemize}
    \item \textbf{Range:} $a \leq \text{attr} \leq b$
    \item \textbf{Comparison:} $\text{attr} > c$
    \item \textbf{Membership:} $\text{attr} \in \{v_1, v_2, ...\}$
    \item \textbf{Equality:} $\text{attr}_1 = \text{attr}_2$
\end{itemize}
\end{definition}

\begin{lstlisting}[language=circom, caption={Circuito para age verification con ZK}]
pragma circom 2.0.0;

include "circomlib/comparators.circom";
include "circomlib/poseidon.circom";

// Probar que una persona es mayor de edad sin revelar
// su fecha de nacimiento exacta

template AgeVerification() {
    // Inputs públicos
    signal input credentialRoot;     // Root del Merkle tree de la credencial
    signal input currentDate;        // Fecha actual (YYYYMMDD como entero)
    signal input minimumAge;         // Edad mínima requerida (ej: 18)
    
    // Inputs privados
    signal input birthDate;          // Fecha de nacimiento (oculta)
    signal input birthDateSalt;      // Salt para el commitment
    signal input merkleProof[10];    // Prueba de Merkle
    signal input merkleIndices[10];  // Índices del path
    
    // Output
    signal output isAdult;
    
    // 1. Verificar que birthDate está en la credencial
    component birthDateHash = Poseidon(2);
    birthDateHash.inputs[0] <== birthDate;
    birthDateHash.inputs[1] <== birthDateSalt;
    
    // Verificar Merkle proof (simplificado)
    // ... (similar a otros ejemplos de Merkle)
    
    // 2. Calcular edad actual
    // Simplificación: asumimos años completos
    // birthDate y currentDate en formato YYYYMMDD
    
    signal birthYear;
    signal currentYear;
    
    birthYear <-- birthDate \ 10000;  // Extraer año
    currentYear <-- currentDate \ 10000;
    
    // Constrains para la división
    signal birthRemainder;
    birthRemainder <== birthDate - birthYear * 10000;
    
    signal currentRemainder;
    currentRemainder <== currentDate - currentYear * 10000;
    
    // Verificar que son divisiones correctas
    component birthYearRange = LessThan(16);
    birthYearRange.in[0] <== birthRemainder;
    birthYearRange.in[1] <== 10000;
    birthYearRange.out === 1;
    
    // Calcular edad aproximada
    signal age;
    age <== currentYear - birthYear;
    
    // Ajuste si aún no ha cumplido años este año
    signal monthDayBirth;
    signal monthDayCurrent;
    monthDayBirth <== birthRemainder;
    monthDayCurrent <== currentRemainder;
    
    component notYetBirthday = GreaterThan(16);
    notYetBirthday.in[0] <== monthDayBirth;
    notYetBirthday.in[1] <== monthDayCurrent;
    
    signal adjustedAge;
    adjustedAge <== age - notYetBirthday.out;
    
    // 3. Verificar edad >= minimumAge
    component ageCheck = GreaterEqThan(8);
    ageCheck.in[0] <== adjustedAge;
    ageCheck.in[1] <== minimumAge;
    
    isAdult <== ageCheck.out;
    
    // Forzar que sea adulto (si no, el proof falla)
    isAdult === 1;
}

component main {public [credentialRoot, currentDate, minimumAge]} = AgeVerification();
\end{lstlisting}

\begin{bancozk}
\textbf{Selective Disclosure en Banco ZK:}

Banco ZK implementa selective disclosure para compliance:

\begin{enumerate}
    \item \textbf{Credencial de identidad:}
    \begin{itemize}
        \item Usuario obtiene credencial ZK de proveedor KYC
        \item Credencial contiene: país, categoría de riesgo, fecha verificación
        \item NO contiene: nombre, documento, foto
    \end{itemize}
    
    \item \textbf{Uso en transacciones:}
    \begin{itemize}
        \item Para operar en ciertos pools: probar ``país no sancionado''
        \item Para montos altos: probar ``KYC completado''
        \item Sin revelar identidad específica
    \end{itemize}
    
    \item \textbf{Privacy Pools integration:}
    \begin{itemize}
        \item Probar membresía en ``set de usuarios verificados''
        \item Sin revelar cuál usuario específico
        \item Compatible con regulación sin doxxing
    \end{itemize}
\end{enumerate}

\begin{lstlisting}[language=circom]
// Prueba de compliance para Banco ZK
template BancoZKCompliance() {
    signal input kycCredentialRoot;  // Root de la credencial
    signal input userSecret;         // Secreto del usuario
    
    // Probar: "soy usuario verificado Y mi país no está sancionado"
    // Sin revelar: quién soy, qué país, cuándo me verifiqué
    
    signal output isCompliant;
}
\end{lstlisting}
\end{bancozk}
%===========================================
\section{ZK Identity Proofs}
\label{sec:zk-identity-proofs}
%===========================================

\subsection{Tipos de Pruebas de Identidad}

\begin{definition}[ZK Identity Proof]
Prueba de conocimiento cero que demuestra una propiedad relacionada con la identidad sin revelar la identidad en sí.
\end{definition}

\begin{mathnote}[Taxonomía de ZK Identity Proofs]
\textbf{1. Membership Proofs:}
\begin{itemize}
    \item ``Pertenezco al grupo X''
    \item Ejemplo: Soy empleado de esta empresa
    \item Técnica: Merkle proof de que mi commitment está en el árbol
\end{itemize}

\textbf{2. Non-membership Proofs:}
\begin{itemize}
    \item ``No estoy en la lista Y''
    \item Ejemplo: No estoy en lista de sancionados
    \item Técnica: Más compleja, requiere estructura especial (sparse Merkle tree)
\end{itemize}

\textbf{3. Attribute Proofs:}
\begin{itemize}
    \item ``Tengo la propiedad Z''
    \item Ejemplo: Soy ciudadano de la UE
    \item Técnica: ZK proof sobre atributo de credencial
\end{itemize}

\textbf{4. Range Proofs:}
\begin{itemize}
    \item ``Mi valor está en el rango [a, b]''
    \item Ejemplo: Mi edad está entre 21 y 65
    \item Técnica: Bulletproofs o circuito de comparación
\end{itemize}

\textbf{5. Equality/Inequality Proofs:}
\begin{itemize}
    \item ``Mis atributos X e Y son iguales/diferentes''
    \item Ejemplo: El país de mi pasaporte = país de mi licencia
    \item Técnica: Commitment equality proof
\end{itemize}
\end{mathnote}

\subsection{Semaphore Protocol}

\begin{definition}[Semaphore]
Protocolo de código abierto para pruebas anónimas de membresía en grupos, desarrollado originalmente para Ethereum:
\begin{itemize}
    \item \textbf{Creador:} Barry Whitehat et al. (2020)
    \item \textbf{Uso:} Voting, signaling, rate limiting, whistleblowing
    \item \textbf{Técnica:} Identity commitments en Merkle tree + nullifiers
    \item \textbf{Implementación:} Circom + snarkjs
\end{itemize}
\end{definition}

\begin{figure}[htbp]
\centering
\begin{tikzpicture}[
    node distance=1.5cm,
    phase/.style={rectangle, draw, minimum width=3.5cm, minimum height=1cm, align=center},
    tree/.style={isosceles triangle, draw, shape border rotate=90, minimum width=2cm, minimum height=1.5cm, fill=green!20},
    arrow/.style={->, thick}
]
    % Fase 1: Registro
    \node[phase, fill=blue!20] (reg) at (0, 0) {1. Registro\\Identity Commitment};
    
    % Merkle tree
    \node[tree] (tree) at (5, 0) {};
    \node at (5, 0) [font=\small] {Merkle\\Tree};
    
    % Fase 2: Señalización
    \node[phase, fill=orange!20] (signal) at (10, 0) {2. Signal\\ZK Proof + Nullifier};
    
    % Flechas
    \draw[arrow] (reg) -- node[above, font=\small] {$C = H(id, secret)$} (tree);
    \draw[arrow] (tree) -- node[above, font=\small] {proof of membership} (signal);
    
    % Detalles
    \node at (0, -2) [align=center, font=\small] {
        Usuario genera:\\
        $identity = H(secret)$\\
        $commitment = H(identity)$
    };
    
    \node at (10, -2) [align=center, font=\small] {
        Usuario prueba:\\
        - Commitment en tree\\
        - Nullifier correcto\\
        Sin revelar identidad
    };
\end{tikzpicture}
\caption{Flujo del protocolo Semaphore}
\label{fig:semaphore-flow}
\end{figure}

\begin{mathnote}[Semaphore Identity]
\textbf{Generación de identidad:}
\begin{enumerate}
    \item Usuario genera dos secretos: $\text{trapdoor}$, $\text{nullifier\_secret}$
    \item Calcula: $\text{identity} = \text{Poseidon}(\text{nullifier\_secret}, \text{trapdoor})$
    \item Commitment público: $C = \text{Poseidon}(\text{identity})$
\end{enumerate}

\textbf{Registro en grupo:}
\begin{enumerate}
    \item Admin del grupo añade $C$ al Merkle tree
    \item Root del tree identifica el grupo
\end{enumerate}

\textbf{Señalización anónima:}
\begin{enumerate}
    \item Usuario genera nullifier: $nf = \text{Poseidon}(\text{external\_nullifier}, \text{identity})$
    \item Genera ZK proof de:
    \begin{itemize}
        \item Conoce $(\text{trapdoor}, \text{nullifier\_secret})$ tal que $C$ está en el tree
        \item El nullifier $nf$ es correcto
    \end{itemize}
    \item Publica: $(nf, \text{signal}, \pi)$
\end{enumerate}

\textbf{Propiedades:}
\begin{itemize}
    \item \textbf{Anonimato:} Nadie sabe quién señaló
    \item \textbf{Doble señalización:} Detectada por nullifier repetido
    \item \textbf{External nullifier:} Permite múltiples señales en diferentes contextos
\end{itemize}
\end{mathnote}

\begin{lstlisting}[language=circom, caption={Circuito Semaphore simplificado}]
pragma circom 2.0.0;

include "circomlib/poseidon.circom";
include "circomlib/mux1.circom";

// Merkle tree proof
template MerkleTreeInclusionProof(nLevels) {
    signal input leaf;
    signal input pathIndices[nLevels];
    signal input siblings[nLevels];
    signal output root;
    
    component hashers[nLevels];
    component mux[nLevels];
    
    signal hashes[nLevels + 1];
    hashes[0] <== leaf;
    
    for (var i = 0; i < nLevels; i++) {
        // pathIndices[i] es 0 o 1
        pathIndices[i] * (1 - pathIndices[i]) === 0;
        
        mux[i] = MultiMux1(2);
        mux[i].c[0][0] <== hashes[i];
        mux[i].c[0][1] <== siblings[i];
        mux[i].c[1][0] <== siblings[i];
        mux[i].c[1][1] <== hashes[i];
        mux[i].s <== pathIndices[i];
        
        hashers[i] = Poseidon(2);
        hashers[i].inputs[0] <== mux[i].out[0];
        hashers[i].inputs[1] <== mux[i].out[1];
        
        hashes[i + 1] <== hashers[i].out;
    }
    
    root <== hashes[nLevels];
}

// Semaphore main circuit
template Semaphore(nLevels) {
    // ========== INPUTS PÚBLICOS ==========
    signal input merkleRoot;           // Root del grupo
    signal input nullifierHash;        // Hash del nullifier
    signal input signalHash;           // Hash del mensaje/señal
    signal input externalNullifier;    // Identificador del contexto
    
    // ========== INPUTS PRIVADOS ==========
    signal input identityNullifier;    // Secreto 1
    signal input identityTrapdoor;     // Secreto 2
    signal input treeSiblings[nLevels];
    signal input treePathIndices[nLevels];
    
    // ========== CÁLCULOS ==========
    
    // 1. Calcular identity commitment
    component identityHasher = Poseidon(2);
    identityHasher.inputs[0] <== identityNullifier;
    identityHasher.inputs[1] <== identityTrapdoor;
    
    signal identity;
    identity <== identityHasher.out;
    
    // 2. Calcular commitment (leaf del árbol)
    component commitmentHasher = Poseidon(1);
    commitmentHasher.inputs[0] <== identity;
    
    signal commitment;
    commitment <== commitmentHasher.out;
    
    // 3. Verificar membership en el árbol
    component inclusionProof = MerkleTreeInclusionProof(nLevels);
    inclusionProof.leaf <== commitment;
    for (var i = 0; i < nLevels; i++) {
        inclusionProof.siblings[i] <== treeSiblings[i];
        inclusionProof.pathIndices[i] <== treePathIndices[i];
    }
    
    // El root calculado debe coincidir
    inclusionProof.root === merkleRoot;
    
    // 4. Calcular y verificar nullifier
    component nullifierHasher = Poseidon(2);
    nullifierHasher.inputs[0] <== externalNullifier;
    nullifierHasher.inputs[1] <== identity;
    
    nullifierHasher.out === nullifierHash;
    
    // 5. Vincular al signal (previene front-running)
    // El signalHash es parte de los inputs públicos
    // No hay constraint adicional, pero es parte del statement
    signal signalSquare;
    signalSquare <== signalHash * signalHash;
}

component main {public [merkleRoot, nullifierHash, signalHash, externalNullifier]} = Semaphore(20);
\end{lstlisting}

\begin{lstlisting}[language=JavaScript, caption={Uso de Semaphore en JavaScript}]
const { Identity } = require("@semaphore-protocol/identity")
const { Group } = require("@semaphore-protocol/group")
const { generateProof, verifyProof } = require("@semaphore-protocol/proof")

async function semaphoreDemo() {
    console.log("=== Semaphore Demo ===\n")
    
    // 1. Crear identidades
    const identity1 = new Identity()
    const identity2 = new Identity()
    const identity3 = new Identity()
    
    console.log("Created 3 identities")
    console.log(`Identity 1 commitment: ${identity1.commitment.toString().slice(0, 20)}...`)
    
    // 2. Crear grupo y añadir miembros
    const groupId = 1
    const group = new Group(groupId, 20)  // depth = 20
    
    group.addMember(identity1.commitment)
    group.addMember(identity2.commitment)
    group.addMember(identity3.commitment)
    
    console.log(`\nGroup created with ${group.members.length} members`)
    console.log(`Group root: ${group.root.toString().slice(0, 20)}...`)
    
    // 3. Identity1 quiere señalar anónimamente
    const externalNullifier = 12345  // Contexto (ej: ID de votación)
    const signal = "I vote YES"
    
    console.log(`\n--- Identity 1 signals: "${signal}" ---`)
    
    // Generar proof
    const proof = await generateProof(
        identity1,
        group,
        externalNullifier,
        signal
    )
    
    console.log(`Nullifier hash: ${proof.nullifierHash.toString().slice(0, 20)}...`)
    console.log(`Proof generated (${proof.proof.length} elements)`)
    
    // 4. Verificar proof
    const isValid = await verifyProof(proof, 20)
    console.log(`\nProof verification: ${isValid}`)
    
    // 5. Intentar señalar de nuevo (doble señalización)
    console.log("\n--- Identity 1 tries to signal again ---")
    
    const proof2 = await generateProof(
        identity1,
        group,
        externalNullifier,  // Mismo contexto
        "I vote NO"
    )
    
    // El nullifier será el mismo!
    console.log(`New nullifier: ${proof2.nullifierHash.toString().slice(0, 20)}...`)
    console.log(`Same nullifier: ${proof.nullifierHash === proof2.nullifierHash}`)
    console.log("=> Double signaling detected!")
    
    // 6. Señalar en diferente contexto (OK)
    console.log("\n--- Identity 1 signals in different context ---")
    
    const differentContext = 99999
    const proof3 = await generateProof(
        identity1,
        group,
        differentContext,
        "Hello world"
    )
    
    console.log(`Different nullifier: ${proof3.nullifierHash.toString().slice(0, 20)}...`)
    console.log(`Nullifiers different: ${proof.nullifierHash !== proof3.nullifierHash}`)
    console.log("=> Different context allows new signal")
}

semaphoreDemo().catch(console.error)
\end{lstlisting}

\subsection{Aplicaciones de Semaphore}

\begin{keyidea}
\textbf{Casos de uso de Semaphore:}

\begin{enumerate}
    \item \textbf{Voting anónimo:}
    \begin{itemize}
        \item Grupo = votantes elegibles
        \item External nullifier = ID de la propuesta
        \item Signal = voto (sí/no)
        \item Cada miembro vota una vez por propuesta
    \end{itemize}
    
    \item \textbf{Whistleblowing:}
    \begin{itemize}
        \item Grupo = empleados de empresa
        \item Signal = denuncia
        \item Nullifier previene spam
        \item Identidad permanece oculta
    \end{itemize}
    
    \item \textbf{Rate limiting:}
    \begin{itemize}
        \item Grupo = usuarios verificados
        \item External nullifier = timestamp (hora/día)
        \item Limita acciones por período sin identificar
    \end{itemize}
    
    \item \textbf{Anonymous authentication:}
    \begin{itemize}
        \item Grupo = usuarios autorizados
        \item Probar membresía sin revelar quién
        \item Acceso a recursos restringidos
    \end{itemize}
\end{enumerate}
\end{keyidea}

\subsection{Non-membership Proofs}

\begin{definition}[Non-membership Proof]
Prueba de que un elemento NO está en un conjunto:
\[
\text{Proof}(x, S) \implies x \notin S
\]

Más difícil que membership porque requiere ``probar un negativo''.
\end{definition}

\begin{mathnote}[Técnicas para Non-membership]
\textbf{1. Sparse Merkle Tree:}
\begin{itemize}
    \item Árbol con posición predeterminada para cada posible elemento
    \item Non-membership = mostrar que la posición tiene valor vacío
    \item Tamaño: $2^{256}$ posiciones virtuales
\end{itemize}

\textbf{2. RSA Accumulators:}
\begin{itemize}
    \item $A = g^{\prod_{x \in S} H(x)} \mod n$
    \item Non-membership witness más complejo
    \item Basado en ``proofs of exponentiation''
\end{itemize}

\textbf{3. Indexed Merkle Tree (Aztec):}
\begin{itemize}
    \item Elementos ordenados con punteros
    \item Non-membership: mostrar elementos adyacentes
    \item Más eficiente que sparse tree
\end{itemize}
\end{mathnote}

\begin{lstlisting}[language=Python, caption={Sparse Merkle Tree para non-membership proofs}]
from hashlib import sha256
from typing import Dict, List, Tuple, Optional

class SparseMerkleTree:
    """
    Sparse Merkle Tree para membership y non-membership proofs.
    Soporta 2^depth elementos posibles.
    """
    
    def __init__(self, depth: int = 256):
        self.depth = depth
        self.empty_leaf = b'\x00' * 32
        self.empty_tree_roots = self._precompute_empty_roots()
        self.leaves: Dict[int, bytes] = {}  # position -> value
    
    def _hash(self, left: bytes, right: bytes) -> bytes:
        return sha256(left + right).digest()
    
    def _precompute_empty_roots(self) -> List[bytes]:
        """Precomputar roots de subárboles vacíos."""
        roots = [self.empty_leaf]
        for _ in range(self.depth):
            roots.append(self._hash(roots[-1], roots[-1]))
        return roots
    
    def _get_position(self, key: bytes) -> int:
        """Convertir key a posición en el árbol."""
        return int.from_bytes(sha256(key).digest(), 'big') % (2 ** self.depth)
    
    def insert(self, key: bytes, value: bytes):
        """Insertar elemento."""
        position = self._get_position(key)
        self.leaves[position] = value
    
    def remove(self, key: bytes):
        """Remover elemento."""
        position = self._get_position(key)
        if position in self.leaves:
            del self.leaves[position]
    
    def get_root(self) -> bytes:
        """Calcular root del árbol."""
        return self._compute_root(0, self.depth)
    
    def _compute_root(self, start: int, level: int) -> bytes:
        """Recursivamente calcular root de subárbol."""
        if level == 0:
            return self.leaves.get(start, self.empty_leaf)
        
        mid = 2 ** (level - 1)
        left = self._compute_root(start, level - 1)
        right = self._compute_root(start + mid, level - 1)
        
        return self._hash(left, right)
    
    def get_membership_proof(self, key: bytes) -> Dict:
        """
        Generar proof de membership.
        """
        position = self._get_position(key)
        
        if position not in self.leaves:
            return {'exists': False}
        
        proof_path = []
        current_pos = position
        
        for level in range(self.depth):
            sibling_pos = current_pos ^ 1  # Flip last bit
            
            # Obtener valor del hermano
            if level == 0:
                sibling = self.leaves.get(sibling_pos, self.empty_leaf)
            else:
                # En niveles superiores, calcular subárbol
                sibling = self._compute_root(
                    (sibling_pos >> level) << level,
                    level
                )
            
            proof_path.append({
                'sibling': sibling,
                'position': 'left' if current_pos % 2 == 1 else 'right'
            })
            
            current_pos //= 2
        
        return {
            'exists': True,
            'key': key,
            'value': self.leaves[position],
            'proof': proof_path
        }
    
    def get_non_membership_proof(self, key: bytes) -> Dict:
        """
        Generar proof de NON-membership.
        Demuestra que la posición está vacía.
        """
        position = self._get_position(key)
        
        if position in self.leaves:
            return {'exists': True, 'error': 'Key exists'}
        
        # Proof de que la posición tiene el valor vacío
        proof_path = []
        current_pos = position
        
        for level in range(self.depth):
            sibling_pos = current_pos ^ 1
            
            if level == 0:
                sibling = self.leaves.get(sibling_pos, self.empty_leaf)
            else:
                sibling = self._compute_root(
                    (sibling_pos >> level) << level,
                    level
                )
            
            proof_path.append({
                'sibling': sibling,
                'position': 'left' if current_pos % 2 == 1 else 'right'
            })
            
            current_pos //= 2
        
        return {
            'exists': False,
            'key': key,
            'empty_position': position,
            'proof': proof_path
        }
    
    @staticmethod
    def verify_membership(root: bytes, proof: Dict, empty_leaf: bytes) -> bool:
        """Verificar proof de membership."""
        if not proof['exists']:
            return False
        
        current = proof['value']
        
        for step in proof['proof']:
            if step['position'] == 'right':
                current = sha256(step['sibling'] + current).digest()
            else:
                current = sha256(current + step['sibling']).digest()
        
        return current == root
    
    @staticmethod
    def verify_non_membership(root: bytes, proof: Dict, empty_leaf: bytes) -> bool:
        """Verificar proof de non-membership."""
        if proof.get('exists', True):
            return False
        
        # Verificar que la posición tiene valor vacío
        current = empty_leaf
        
        for step in proof['proof']:
            if step['position'] == 'right':
                current = sha256(step['sibling'] + current).digest()
            else:
                current = sha256(current + step['sibling']).digest()
        
        return current == root


def sparse_merkle_demo():
    print("=== Sparse Merkle Tree Demo ===\n")
    
    # Usar depth pequeño para demo
    smt = SparseMerkleTree(depth=8)
    
    # Insertar algunos elementos
    smt.insert(b"alice", b"data_alice")
    smt.insert(b"bob", b"data_bob")
    smt.insert(b"carol", b"data_carol")
    
    root = smt.get_root()
    print(f"Tree root: {root.hex()[:16]}...")
    
    # Membership proof
    print("\n--- Membership Proof ---")
    membership_proof = smt.get_membership_proof(b"alice")
    print(f"Alice exists: {membership_proof['exists']}")
    
    is_valid = SparseMerkleTree.verify_membership(
        root, membership_proof, smt.empty_leaf
    )
    print(f"Proof valid: {is_valid}")
    
    # Non-membership proof
    print("\n--- Non-membership Proof ---")
    non_membership_proof = smt.get_non_membership_proof(b"dave")
    print(f"Dave exists: {non_membership_proof['exists']}")
    
    is_valid_non = SparseMerkleTree.verify_non_membership(
        root, non_membership_proof, smt.empty_leaf
    )
    print(f"Non-membership proof valid: {is_valid_non}")
    
    # Aplicación: Lista de sancionados
    print("\n--- Aplicación: Sanctions List ---")
    sanctions = SparseMerkleTree(depth=16)
    sanctions.insert(b"sanctioned_entity_1", b"1")
    sanctions.insert(b"sanctioned_entity_2", b"1")
    
    sanctions_root = sanctions.get_root()
    
    # Usuario prueba que NO está sancionado
    user_proof = sanctions.get_non_membership_proof(b"honest_user")
    print(f"Honest user sanctioned: {user_proof['exists']}")
    print("User can prove they are NOT on sanctions list!")

sparse_merkle_demo()
\end{lstlisting}
%===========================================
\section{Soulbound Tokens (SBTs) con Privacidad}
\label{sec:soulbound-tokens}
%===========================================

\subsection{El Problema de los SBTs Públicos}

\begin{definition}[Soulbound Token (SBT)]
Token no transferible vinculado permanentemente a una dirección (``alma''), que representa credenciales, logros, afiliaciones o propiedades de identidad.

\textbf{Concepto:} Introducido por Vitalik Buterin et al. en ``Decentralized Society'' (2022).

\textbf{Ejemplos:}
\begin{itemize}
    \item Diploma universitario
    \item Certificación profesional
    \item Historial crediticio
    \item Membresía de DAO
    \item Prueba de asistencia (POAP vinculado)
\end{itemize}
\end{definition}

\begin{advertencia}
\textbf{Problemas de privacidad con SBTs públicos:}

\begin{enumerate}
    \item \textbf{Doxxing por acumulación:}
    \begin{itemize}
        \item Cada SBT revela información
        \item Combinación de SBTs permite perfilar completamente
        \item Ejemplo: Universidad X + Empleador Y + Ciudad Z = persona identificable
    \end{itemize}
    
    \item \textbf{Discriminación:}
    \begin{itemize}
        \item SBT de historial médico visible
        \item SBT de antecedentes penales
        \item Servicios pueden discriminar basándose en SBTs
    \end{itemize}
    
    \item \textbf{Vigilancia social:}
    \begin{itemize}
        \item Gobiernos pueden rastrear credenciales
        \item Score social basado en SBTs
        \item Sin posibilidad de ``borrar'' historial
    \end{itemize}
    
    \item \textbf{Linkability:}
    \begin{itemize}
        \item Todos los SBTs vinculados a una dirección
        \item No hay separación de contextos
        \item Tu empleador ve tus SBTs médicos
    \end{itemize}
\end{enumerate}
\end{advertencia}

\subsection{SBTs Privados con ZK}

\begin{keyidea}
\textbf{Solución: Private SBTs}

En lugar de SBTs públicos visibles para todos:
\begin{itemize}
    \item \textbf{Emisión privada:} SBT emitido como commitment oculto
    \item \textbf{Almacenamiento privado:} Solo el holder conoce sus SBTs
    \item \textbf{Revelación selectiva:} ZK proof de posesión sin mostrar el SBT
    \item \textbf{Unlinkability:} Diferentes pruebas no son vinculables
\end{itemize}
\end{keyidea}

\begin{figure}[htbp]
\centering
\begin{tikzpicture}[
    node distance=1cm,
    sbt/.style={rectangle, draw, rounded corners, minimum width=2cm, minimum height=0.8cm, fill=purple!20},
    hidden/.style={rectangle, draw, rounded corners, minimum width=2cm, minimum height=0.8cm, fill=gray!40, dashed},
    entity/.style={rectangle, draw, minimum width=2.5cm, minimum height=1cm, align=center},
    arrow/.style={->, thick}
]
    % SBTs públicos (problema)
    \node at (-4, 3) {\textbf{SBTs Públicos (Problema)}};
    
    \node[entity, fill=blue!20] (wallet1) at (-4, 1.5) {Wallet\\0x123...};
    
    \node[sbt] (s1) at (-6, 0) {Diploma MIT};
    \node[sbt] (s2) at (-4, 0) {Empleado Google};
    \node[sbt] (s3) at (-2, 0) {Residente SF};
    
    \draw[arrow] (wallet1) -- (s1);
    \draw[arrow] (wallet1) -- (s2);
    \draw[arrow] (wallet1) -- (s3);
    
    \node at (-4, -1) [font=\small, red] {Todo visible → Identidad expuesta};
    
    % SBTs privados (solución)
    \node at (4, 3) {\textbf{SBTs Privados (Solución)}};
    
    \node[entity, fill=green!20] (wallet2) at (4, 1.5) {Wallet\\0x456...};
    
    \node[hidden] (h1) at (2, 0) {???};
    \node[hidden] (h2) at (4, 0) {???};
    \node[hidden] (h3) at (6, 0) {???};
    
    \draw[arrow, dashed] (wallet2) -- (h1);
    \draw[arrow, dashed] (wallet2) -- (h2);
    \draw[arrow, dashed] (wallet2) -- (h3);
    
    \node at (4, -1) [font=\small, green!60!black] {Ocultos, revelación selectiva con ZK};
\end{tikzpicture}
\caption{SBTs públicos vs privados}
\end{figure}

\begin{lstlisting}[language=Solidity, caption={Private SBT Registry con ZK}]
// SPDX-License-Identifier: MIT
pragma solidity ^0.8.20;

interface IVerifier {
    function verifyProof(
        bytes calldata proof,
        uint256[] calldata publicInputs
    ) external view returns (bool);
}

contract PrivateSBTRegistry {
    // Verificadores para diferentes tipos de pruebas
    IVerifier public ownershipVerifier;
    IVerifier public attributeVerifier;
    
    // Merkle tree de SBT commitments
    bytes32 public sbtTreeRoot;
    uint256 public nextSbtIndex;
    mapping(uint256 => bytes32) public sbtCommitments;
    
    // Nullifiers para prevenir doble uso (opcional)
    mapping(bytes32 => bool) public usedNullifiers;
    
    // Issuers autorizados
    mapping(address => bool) public authorizedIssuers;
    mapping(address => bytes32) public issuerPublicKeys;
    
    // Tipos de SBT
    struct SBTType {
        string name;
        address issuer;
        bytes32 schemaHash;
        bool revocable;
    }
    mapping(uint256 => SBTType) public sbtTypes;
    
    event SBTIssued(
        uint256 indexed sbtTypeId,
        bytes32 indexed commitment,
        uint256 leafIndex,
        bytes encryptedData  // Encriptado para el receptor
    );
    
    event SBTRevoked(bytes32 indexed nullifier);
    
    // ========== EMISIÓN PRIVADA ==========
    
    function issueSBT(
        uint256 sbtTypeId,
        bytes32 commitment,       // H(sbtType, attributes, owner, randomness)
        bytes calldata encryptedData,  // Datos encriptados para el receptor
        bytes calldata issuerSignature
    ) external {
        require(authorizedIssuers[msg.sender], "Not authorized issuer");
        require(sbtTypes[sbtTypeId].issuer == msg.sender, "Wrong issuer");
        
        // Verificar firma del issuer sobre el commitment
        // (prueba que el issuer aprobó este SBT específico)
        require(
            _verifyIssuerSignature(
                msg.sender,
                commitment,
                issuerSignature
            ),
            "Invalid issuer signature"
        );
        
        // Añadir al árbol
        uint256 index = nextSbtIndex++;
        sbtCommitments[index] = commitment;
        _updateMerkleRoot();
        
        emit SBTIssued(sbtTypeId, commitment, index, encryptedData);
    }
    
    // ========== VERIFICACIÓN PRIVADA ==========
    
    function verifyOwnership(
        bytes calldata proof,
        bytes32 merkleRoot,
        uint256 sbtTypeId
    ) external view returns (bool) {
        require(merkleRoot == sbtTreeRoot, "Invalid root");
        
        // Public inputs: root, sbtTypeId
        uint256[] memory publicInputs = new uint256[](2);
        publicInputs[0] = uint256(merkleRoot);
        publicInputs[1] = sbtTypeId;
        
        // La prueba demuestra:
        // 1. Prover conoce un SBT commitment en el árbol
        // 2. El SBT es del tipo especificado
        // 3. Prover es el owner del SBT
        // Sin revelar: cuál SBT específico, otros atributos
        
        return ownershipVerifier.verifyProof(proof, publicInputs);
    }
    
    function verifyAttribute(
        bytes calldata proof,
        bytes32 merkleRoot,
        uint256 sbtTypeId,
        bytes32 attributeHash,  // H(attribute_name)
        bytes32 expectedValue   // Valor esperado o commitment
    ) external view returns (bool) {
        require(merkleRoot == sbtTreeRoot, "Invalid root");
        
        uint256[] memory publicInputs = new uint256[](4);
        publicInputs[0] = uint256(merkleRoot);
        publicInputs[1] = sbtTypeId;
        publicInputs[2] = uint256(attributeHash);
        publicInputs[3] = uint256(expectedValue);
        
        // La prueba demuestra:
        // 1. Prover tiene SBT válido del tipo correcto
        // 2. El atributo especificado tiene el valor esperado
        // Sin revelar: identidad, otros atributos
        
        return attributeVerifier.verifyProof(proof, publicInputs);
    }
    
    // Prueba de rango sobre atributo
    function verifyAttributeRange(
        bytes calldata proof,
        bytes32 merkleRoot,
        uint256 sbtTypeId,
        bytes32 attributeHash,
        uint256 minValue,
        uint256 maxValue
    ) external view returns (bool) {
        uint256[] memory publicInputs = new uint256[](5);
        publicInputs[0] = uint256(merkleRoot);
        publicInputs[1] = sbtTypeId;
        publicInputs[2] = uint256(attributeHash);
        publicInputs[3] = minValue;
        publicInputs[4] = maxValue;
        
        // Verifica: minValue <= attribute <= maxValue
        return attributeVerifier.verifyProof(proof, publicInputs);
    }
    
    // ========== REVOCACIÓN ==========
    
    function revokeSBT(
        bytes32 nullifier,
        bytes calldata revokerProof
    ) external {
        // Solo el issuer puede revocar
        // Proof demuestra que el nullifier corresponde a un SBT que emitió
        
        require(!usedNullifiers[nullifier], "Already revoked");
        usedNullifiers[nullifier] = true;
        
        emit SBTRevoked(nullifier);
    }
    
    // ========== HELPERS ==========
    
    function _updateMerkleRoot() internal {
        // Actualizar root del Merkle tree
        // (implementación simplificada)
    }
    
    function _verifyIssuerSignature(
        address issuer,
        bytes32 commitment,
        bytes calldata signature
    ) internal view returns (bool) {
        // Verificar firma ECDSA o BLS
        return true; // Placeholder
    }
}
\end{lstlisting}

\begin{lstlisting}[language=circom, caption={Circuito para prueba de posesión de SBT privado}]
pragma circom 2.0.0;

include "circomlib/poseidon.circom";

// Prueba de que poseo un SBT de cierto tipo sin revelar cuál
template PrivateSBTOwnership(treeLevels, numAttributes) {
    // ========== PUBLIC INPUTS ==========
    signal input merkleRoot;
    signal input sbtTypeId;
    
    // ========== PRIVATE INPUTS ==========
    signal input ownerSecret;           // Secreto del owner
    signal input sbtRandomness;         // Randomness del commitment
    signal input attributes[numAttributes];
    signal input merkleProof[treeLevels];
    signal input merkleIndices[treeLevels];
    
    // ========== CONSTRAINTS ==========
    
    // 1. Calcular owner commitment
    component ownerHash = Poseidon(1);
    ownerHash.inputs[0] <== ownerSecret;
    signal ownerCommitment;
    ownerCommitment <== ownerHash.out;
    
    // 2. Calcular SBT commitment
    // commitment = H(sbtTypeId, attributes..., ownerCommitment, randomness)
    component sbtHash = Poseidon(numAttributes + 3);
    sbtHash.inputs[0] <== sbtTypeId;
    for (var i = 0; i < numAttributes; i++) {
        sbtHash.inputs[i + 1] <== attributes[i];
    }
    sbtHash.inputs[numAttributes + 1] <== ownerCommitment;
    sbtHash.inputs[numAttributes + 2] <== sbtRandomness;
    
    signal sbtCommitment;
    sbtCommitment <== sbtHash.out;
    
    // 3. Verificar Merkle proof
    component merkleVerifier = MerkleProof(treeLevels);
    merkleVerifier.leaf <== sbtCommitment;
    for (var i = 0; i < treeLevels; i++) {
        merkleVerifier.siblings[i] <== merkleProof[i];
        merkleVerifier.pathIndices[i] <== merkleIndices[i];
    }
    
    // Root debe coincidir
    merkleVerifier.root === merkleRoot;
}

// Prueba de atributo específico
template PrivateSBTAttributeProof(treeLevels, numAttributes, attrIndex) {
    // PUBLIC
    signal input merkleRoot;
    signal input sbtTypeId;
    signal input expectedAttributeValue;
    
    // PRIVATE
    signal input ownerSecret;
    signal input sbtRandomness;
    signal input attributes[numAttributes];
    signal input merkleProof[treeLevels];
    signal input merkleIndices[treeLevels];
    
    // Verificar ownership (mismo que arriba)
    // ...
    
    // Verificar que el atributo tiene el valor esperado
    attributes[attrIndex] === expectedAttributeValue;
}

component main {public [merkleRoot, sbtTypeId]} = PrivateSBTOwnership(20, 5);
\end{lstlisting}

%===========================================
\section{DID + ZK}
\label{sec:did-zk}
%===========================================

\subsection{Decentralized Identifiers (DIDs)}

\begin{definition}[DID - Decentralized Identifier]
Identificador único, persistente y verificable criptográficamente, que no requiere una autoridad centralizada:

\textbf{Formato:} \texttt{did:method:specific-identifier}

\textbf{Ejemplos:}
\begin{itemize}
    \item \texttt{did:ethr:0x123...abc} (Ethereum)
    \item \texttt{did:key:z6Mk...} (Key-based)
    \item \texttt{did:web:example.com} (Web-based)
    \item \texttt{did:ion:EiD...} (Bitcoin/ION)
\end{itemize}

\textbf{Estándar:} W3C DID Core Specification
\end{definition}

\begin{mathnote}[DID Document]
Cada DID resuelve a un \textbf{DID Document} que contiene:

\begin{lstlisting}[language=json]
{
  "@context": "https://www.w3.org/ns/did/v1",
  "id": "did:example:123456789",
  "authentication": [{
    "id": "did:example:123456789#keys-1",
    "type": "Ed25519VerificationKey2020",
    "controller": "did:example:123456789",
    "publicKeyMultibase": "z6Mkf..."
  }],
  "service": [{
    "id": "did:example:123456789#messaging",
    "type": "MessagingService",
    "serviceEndpoint": "https://example.com/messages"
  }]
}
\end{lstlisting}

\textbf{Componentes clave:}
\begin{itemize}
    \item \textbf{Verification Methods:} Claves públicas para autenticación
    \item \textbf{Services:} Endpoints para interactuar con el DID
    \item \textbf{Controller:} Quién puede modificar el documento
\end{itemize}
\end{mathnote}

\subsection{Verifiable Credentials (VCs)}

\begin{definition}[Verifiable Credential]
Credencial digital que puede ser verificada criptográficamente:
\begin{itemize}
    \item \textbf{Issuer:} Entidad que emite la credencial
    \item \textbf{Subject:} Entidad sobre la que se hace el claim
    \item \textbf{Holder:} Quien posee y presenta la credencial
    \item \textbf{Verifier:} Quien verifica la credencial
\end{itemize}

\textbf{Estándar:} W3C Verifiable Credentials Data Model
\end{definition}

\begin{figure}[htbp]
\centering
\begin{tikzpicture}[
    node distance=2cm,
    entity/.style={rectangle, draw, minimum width=2.5cm, minimum height=1cm, align=center, fill=blue!10},
    vc/.style={rectangle, draw, rounded corners, minimum width=3cm, minimum height=2cm, align=center, fill=yellow!20},
    arrow/.style={->, thick}
]
    % Entidades
    \node[entity] (issuer) at (0, 2) {Issuer\\(Universidad)};
    \node[entity] (holder) at (5, 2) {Holder\\(Estudiante)};
    \node[entity] (verifier) at (10, 2) {Verifier\\(Empleador)};
    
    % VC
    \node[vc] (vc) at (5, -1) {
        \textbf{Verifiable Credential}\\
        Diploma: CS\\
        GPA: 3.8\\
        Year: 2023\\
        \textit{Firma del Issuer}
    };
    
    % DID
    \node at (5, 0.5) [font=\small] {did:example:student};
    
    % Flujo
    \draw[arrow] (issuer) -- node[above, font=\small] {1. Emite} (holder);
    \draw[arrow] (holder) -- (vc);
    \draw[arrow] (holder) -- node[above, font=\small] {2. Presenta} (verifier);
    \draw[arrow, dashed] (verifier) to[bend right=30] node[below, font=\small] {3. Verifica firma} (issuer);
\end{tikzpicture}
\caption{Flujo de Verifiable Credentials}
\end{figure}

\subsection{Combinando DIDs + VCs + ZK}

\begin{keyidea}
\textbf{El problema con VCs tradicionales:}
\begin{itemize}
    \item Presentar una VC revela todos sus claims
    \item El verifier puede rastrear al holder
    \item Múltiples presentaciones son vinculables
\end{itemize}

\textbf{Solución: ZK-VCs}
\begin{itemize}
    \item Selective disclosure nativo
    \item Presentaciones unlinkable
    \item Predicados sobre atributos
    \item Privacidad del holder preservada
\end{itemize}
\end{keyidea}

\begin{lstlisting}[language=Python, caption={Sistema DID + VC + ZK conceptual}]
from dataclasses import dataclass
from typing import Dict, List, Optional
from hashlib import sha256
import json
from datetime import datetime

@dataclass
class DIDDocument:
    """Documento DID simplificado."""
    did: str
    verification_methods: List[Dict]
    services: List[Dict]
    created: datetime
    updated: datetime
    
    def to_dict(self) -> Dict:
        return {
            "@context": "https://www.w3.org/ns/did/v1",
            "id": self.did,
            "verificationMethod": self.verification_methods,
            "service": self.services,
            "created": self.created.isoformat(),
            "updated": self.updated.isoformat()
        }

@dataclass
class VerifiableCredential:
    """Verifiable Credential W3C-style."""
    context: List[str]
    type: List[str]
    issuer: str  # DID del issuer
    issuance_date: datetime
    expiration_date: Optional[datetime]
    credential_subject: Dict
    proof: Dict
    
    def to_dict(self) -> Dict:
        return {
            "@context": self.context,
            "type": self.type,
            "issuer": self.issuer,
            "issuanceDate": self.issuance_date.isoformat(),
            "expirationDate": self.expiration_date.isoformat() if self.expiration_date else None,
            "credentialSubject": self.credential_subject,
            "proof": self.proof
        }

class ZKVerifiableCredential:
    """
    VC con capacidades ZK para selective disclosure.
    """
    
    def __init__(self, vc: VerifiableCredential):
        self.vc = vc
        self.commitment = self._create_commitment()
    
    def _create_commitment(self) -> bytes:
        """Crear commitment de la credencial completa."""
        data = json.dumps(self.vc.to_dict(), sort_keys=True)
        return sha256(data.encode()).digest()
    
    def create_presentation(
        self,
        disclosed_claims: List[str],
        predicates: Dict[str, tuple],  # claim -> (operator, value)
        holder_secret: bytes,
        verifier_challenge: bytes
    ) -> Dict:
        """
        Crear presentación con selective disclosure.
        
        Args:
            disclosed_claims: Claims a revelar completamente
            predicates: Claims a probar con predicados (ej: age > 18)
            holder_secret: Secreto del holder para binding
            verifier_challenge: Challenge del verifier (previene replay)
        
        Returns:
            Presentación verificable
        """
        presentation = {
            "@context": ["https://www.w3.org/ns/credentials/v2"],
            "type": ["VerifiablePresentation"],
            "verifiableCredential": [],
            "holder": self.vc.credential_subject.get("id"),
        }
        
        # Datos revelados
        revealed_data = {}
        for claim in disclosed_claims:
            if claim in self.vc.credential_subject:
                revealed_data[claim] = self.vc.credential_subject[claim]
        
        # Predicados evaluados
        predicate_results = {}
        for claim, (op, value) in predicates.items():
            if claim in self.vc.credential_subject:
                actual = self.vc.credential_subject[claim]
                if op == ">":
                    predicate_results[f"{claim} > {value}"] = actual > value
                elif op == "<":
                    predicate_results[f"{claim} < {value}"] = actual < value
                elif op == ">=":
                    predicate_results[f"{claim} >= {value}"] = actual >= value
                elif op == "==":
                    predicate_results[f"{claim} == {value}"] = actual == value
        
        # Crear ZK proof (placeholder)
        zk_proof = self._generate_zk_proof(
            revealed_data,
            predicate_results,
            holder_secret,
            verifier_challenge
        )
        
        presentation["proof"] = {
            "type": "ZKSelectiveDisclosureProof2024",
            "created": datetime.now().isoformat(),
            "challenge": verifier_challenge.hex(),
            "revealedClaims": revealed_data,
            "predicateResults": predicate_results,
            "zkProof": zk_proof.hex()
        }
        
        return presentation
    
    def _generate_zk_proof(
        self,
        revealed: Dict,
        predicates: Dict,
        secret: bytes,
        challenge: bytes
    ) -> bytes:
        """
        Generar ZK proof real.
        En producción: usar BBS+, Groth16, etc.
        """
        # Placeholder: hash como "proof"
        data = json.dumps({
            'commitment': self.commitment.hex(),
            'revealed': revealed,
            'predicates': predicates,
            'challenge': challenge.hex()
        }, sort_keys=True)
        
        return sha256(data.encode() + secret).digest()
    
    @staticmethod
    def verify_presentation(
        presentation: Dict,
        issuer_public_key: bytes,
        expected_challenge: bytes
    ) -> Dict:
        """
        Verificar una presentación ZK.
        
        Returns:
            Resultado de verificación con claims verificados
        """
        proof = presentation.get("proof", {})
        
        # Verificar challenge
        if bytes.fromhex(proof.get("challenge", "")) != expected_challenge:
            return {"valid": False, "error": "Challenge mismatch"}
        
        # Verificar ZK proof (placeholder)
        # En producción: verificación criptográfica real
        
        return {
            "valid": True,
            "revealedClaims": proof.get("revealedClaims", {}),
            "predicateResults": proof.get("predicateResults", {}),
            "issuer": "verified"  # Después de verificar firma
        }


def did_vc_zk_demo():
    print("=== DID + VC + ZK Demo ===\n")
    
    # 1. Universidad crea DID
    university_did = "did:example:university123"
    print(f"University DID: {university_did}")
    
    # 2. Estudiante crea DID
    student_did = "did:example:student456"
    print(f"Student DID: {student_did}")
    
    # 3. Universidad emite VC al estudiante
    diploma_vc = VerifiableCredential(
        context=["https://www.w3.org/ns/credentials/v2"],
        type=["VerifiableCredential", "DiplomaCredential"],
        issuer=university_did,
        issuance_date=datetime(2023, 6, 15),
        expiration_date=None,  # No expira
        credential_subject={
            "id": student_did,
            "name": "Alice Smith",
            "degree": "Computer Science",
            "gpa": 3.85,
            "graduationYear": 2023,
            "honors": True
        },
        proof={
            "type": "Ed25519Signature2020",
            "verificationMethod": f"{university_did}#key-1",
            "proofValue": "z..."  # Firma real
        }
    )
    
    print(f"\nVC issued for degree: {diploma_vc.credential_subject['degree']}")
    
    # 4. Crear ZK-VC
    zk_vc = ZKVerifiableCredential(diploma_vc)
    
    # 5. Empleador quiere verificar:
    # - Que tiene grado CS (revelar)
    # - Que GPA > 3.5 (predicado, no revelar exacto)
    # - NO revelar nombre
    
    challenge = b"employer_challenge_12345"
    
    presentation = zk_vc.create_presentation(
        disclosed_claims=["degree", "graduationYear"],  # Revelar
        predicates={"gpa": (">=", 3.5)},  # Solo predicado
        holder_secret=b"student_secret",
        verifier_challenge=challenge
    )
    
    print("\n--- Presentation Created ---")
    print(f"Revealed claims: {presentation['proof']['revealedClaims']}")
    print(f"Predicate results: {presentation['proof']['predicateResults']}")
    print("Hidden: name, exact GPA, honors status")
    
    # 6. Empleador verifica
    result = ZKVerifiableCredential.verify_presentation(
        presentation,
        issuer_public_key=b"university_pubkey",
        expected_challenge=challenge
    )
    
    print(f"\n--- Verification Result ---")
    print(f"Valid: {result['valid']}")
    print(f"Verified claims: {result['revealedClaims']}")
    print(f"Verified predicates: {result['predicateResults']}")

did_vc_zk_demo()
\end{lstlisting}

%===========================================
\section{Proof of Humanity sin Doxxing}
\label{sec:proof-of-humanity}
%===========================================

\subsection{El Problema Sybil}

\begin{definition}[Sybil Attack]
Ataque donde un adversario crea múltiples identidades falsas para ganar influencia indebida en un sistema:
\begin{itemize}
    \item En voting: Un atacante vota múltiples veces
    \item En airdrops: Farmer reclama múltiples veces
    \item En DAOs: Una persona controla múltiples votos
    \item En redes sociales: Bots amplifican mensajes
\end{itemize}
\end{definition}

\begin{definition}[Proof of Humanity / Proof of Personhood]
Mecanismo para demostrar que una entidad digital corresponde a un humano único, sin revelar la identidad de ese humano.

\textbf{Propiedades deseadas:}
\begin{enumerate}
    \item \textbf{Uniqueness:} Una persona = un proof
    \item \textbf{Privacy:} No revela identidad
    \item \textbf{Accessibility:} Cualquier humano puede obtenerlo
    \item \textbf{Resistance:} Difícil de falsificar o transferir
\end{enumerate}
\end{definition}

\subsection{Proyectos de Proof of Humanity}

\subsubsection{Worldcoin}

\begin{mathnote}[Arquitectura Worldcoin]
\textbf{Proceso:}
\begin{enumerate}
    \item Usuario escanea iris con ``Orb'' (dispositivo especial)
    \item Orb calcula \textit{iris code} (hash biométrico)
    \item Verifica que iris code no existe en base de datos
    \item Si único: genera World ID credential
    \item Usuario puede probar humanidad con ZK sin revelar iris
\end{enumerate}

\textbf{Componentes ZK:}
\begin{itemize}
    \item \textbf{Semaphore:} Para proofs anónimos de membresía
    \item \textbf{Nullifiers:} Previenen doble-claim por contexto
    \item \textbf{Privacy:} Iris code nunca sale del Orb (en teoría)
\end{itemize}

\textbf{Controversias:}
\begin{itemize}
    \item Centralización del Orb y datos biométricos
    \item Preocupaciones sobre biometric surveillance
    \item Desigualdad de acceso a Orbs
\end{itemize}
\end{mathnote}

\begin{lstlisting}[language=Python, caption={Worldcoin-style Proof of Humanity (conceptual)}]
from hashlib import sha256
import secrets
from typing import Dict, Optional

class WorldIDSimulation:
    """
    Simulación simplificada del sistema Worldcoin.
    En producción: usa Semaphore y verificación real de Orb.
    """
    
    def __init__(self):
        # Base de datos de iris codes (en producción: árbol Merkle)
        self.iris_codes: set = set()
        self.identity_commitments: Dict[bytes, bytes] = {}  # iris -> commitment
        self.merkle_root: bytes = b'\x00' * 32
    
    def enroll(self, iris_scan: bytes) -> Optional[Dict]:
        """
        Enrollar nuevo usuario (simulación de Orb).
        
        Args:
            iris_scan: Datos biométricos del iris
        
        Returns:
            Credencial World ID o None si ya existe
        """
        # Calcular iris code (en producción: algoritmo especializado)
        iris_code = sha256(iris_scan).digest()
        
        # Verificar unicidad
        if iris_code in self.iris_codes:
            return None  # Ya registrado
        
        # Registrar
        self.iris_codes.add(iris_code)
        
        # Generar identity secret (usuario guarda esto)
        identity_secret = secrets.token_bytes(32)
        
        # Generar commitment (público)
        commitment = sha256(identity_secret + iris_code).digest()
        self.identity_commitments[iris_code] = commitment
        
        # Actualizar Merkle tree
        self._update_merkle_root()
        
        return {
            'identity_secret': identity_secret,
            'commitment': commitment,
            'iris_code_hash': sha256(iris_code).digest(),  # No revelar iris code
            'enrolled_at': 'timestamp'
        }
    
    def generate_proof(
        self,
        identity_secret: bytes,
        app_id: str,
        action: str
    ) -> Dict:
        """
        Generar ZK proof de humanidad para una acción.
        
        Args:
            identity_secret: Secreto del usuario
            app_id: Identificador de la aplicación
            action: Acción específica (para nullifier)
        
        Returns:
            Proof verificable
        """
        # Calcular nullifier (único por app+action)
        nullifier = sha256(
            identity_secret + 
            app_id.encode() + 
            action.encode()
        ).digest()
        
        # En producción: generar Semaphore proof
        # Aquí: simulación
        proof = {
            'merkle_root': self.merkle_root.hex(),
            'nullifier': nullifier.hex(),
            'signal': sha256(action.encode()).hexdigest(),
            'proof': 'zk_proof_placeholder'
        }
        
        return proof
    
    def verify_proof(
        self,
        proof: Dict,
        app_id: str,
        expected_action: str
    ) -> Dict:
        """
        Verificar proof de humanidad.
        """
        # Verificar signal
        expected_signal = sha256(expected_action.encode()).hexdigest()
        if proof['signal'] != expected_signal:
            return {'valid': False, 'error': 'Signal mismatch'}
        
        # Verificar merkle root es reciente/válido
        if proof['merkle_root'] != self.merkle_root.hex():
            return {'valid': False, 'error': 'Invalid root'}
        
        # En producción: verificar Semaphore proof
        
        return {
            'valid': True,
            'nullifier': proof['nullifier'],
            'action_verified': expected_action
        }
    
    def _update_merkle_root(self):
        """Actualizar Merkle root con nuevos commitments."""
        data = b''.join(sorted(self.identity_commitments.values()))
        self.merkle_root = sha256(data).digest()


def worldid_demo():
    print("=== World ID Demo ===\n")
    
    system = WorldIDSimulation()
    
    # Usuario A se registra
    print("--- User A enrolls ---")
    iris_a = b"unique_iris_pattern_user_a"
    cred_a = system.enroll(iris_a)
    print(f"Enrolled: {cred_a is not None}")
    print(f"Commitment: {cred_a['commitment'].hex()[:16]}...")
    
    # Usuario B se registra
    print("\n--- User B enrolls ---")
    iris_b = b"unique_iris_pattern_user_b"
    cred_b = system.enroll(iris_b)
    print(f"Enrolled: {cred_b is not None}")
    
    # Usuario A intenta registrarse de nuevo (debe fallar)
    print("\n--- User A tries to re-enroll ---")
    cred_a2 = system.enroll(iris_a)
    print(f"Re-enrollment: {cred_a2}")  # None = rejected
    
    # Usuario A genera proof para votar
    print("\n--- User A votes on proposal ---")
    proof_vote = system.generate_proof(
        identity_secret=cred_a['identity_secret'],
        app_id="dao_xyz",
        action="vote_proposal_123_yes"
    )
    print(f"Nullifier: {proof_vote['nullifier'][:16]}...")
    
    # Verificar
    result = system.verify_proof(proof_vote, "dao_xyz", "vote_proposal_123_yes")
    print(f"Verification: {result['valid']}")
    
    # Usuario A intenta votar de nuevo (mismo nullifier detectado)
    print("\n--- User A tries to vote again ---")
    proof_vote2 = system.generate_proof(
        identity_secret=cred_a['identity_secret'],
        app_id="dao_xyz",
        action="vote_proposal_123_yes"
    )
    print(f"Same nullifier: {proof_vote['nullifier'] == proof_vote2['nullifier']}")
    print("=> Double voting detected by nullifier!")

worldid_demo()
\end{lstlisting}

\subsubsection{Otros Enfoques}

\begin{table}[htbp]
\centering
\small
\begin{tabular}{lp{3cm}p{3cm}p{3cm}}
\toprule
\textbf{Proyecto} & \textbf{Método} & \textbf{Ventajas} & \textbf{Desventajas} \\
\midrule
Worldcoin & Biometría (iris) & Alta unicidad, difícil transferir & Centralizado, hardware especial \\
BrightID & Grafo social & Descentralizado, sin biometría & Ataques Sybil posibles, lento \\
Gitcoin Passport & Agregación de señales & Flexible, múltiples fuentes & No garantiza unicidad \\
Proof of Humanity (PoH) & Video + vouching & Simple, solo necesita webcam & Lento, ataques de colusión \\
Idena & CAPTCHA periódico & Descentralizado, accesible & Requiere participación regular \\
\bottomrule
\end{tabular}
\caption{Comparación de sistemas de Proof of Humanity}
\end{table}

\begin{bancozk}
\textbf{Proof of Humanity en Banco ZK:}

Para prevenir Sybil attacks mientras preserva privacidad:

\begin{enumerate}
    \item \textbf{Integración multi-proveedor:}
    \begin{itemize}
        \item Acepta proofs de Worldcoin, BrightID, etc.
        \item Usuario elige proveedor preferido
        \item ZK proof verifica ``humanidad'' sin revelar proveedor
    \end{itemize}
    
    \item \textbf{Rate limiting privado:}
    \begin{itemize}
        \item Un humano = un depósito de alto valor por período
        \item Previene farming de anonymity set
        \item Sin vincular depósitos entre sí
    \end{itemize}
    
    \item \textbf{Governance:}
    \begin{itemize}
        \item Un humano = un voto en decisiones de Banco ZK
        \item Voting anónimo con Semaphore
        \item Previene plutocracy
    \end{itemize}
\end{enumerate}
\end{bancozk}
%===========================================
\section{Age Verification ZK}
\label{sec:age-verification}
%===========================================

\subsection{Motivación}

\begin{intuicion}
La verificación de edad es uno de los casos de uso más claros para ZK:

\textbf{Escenario actual:}
\begin{itemize}
    \item Compras alcohol online: envías copia de ID completo
    \item Accedes a sitio adulto: proporcionas fecha de nacimiento
    \item Resultado: Datos sensibles en manos de terceros
\end{itemize}

\textbf{Con ZK:}
\begin{itemize}
    \item Demuestras: ``Mi edad $\geq$ 21''
    \item Sin revelar: Fecha exacta, nombre, documento, foto
    \item El verificador obtiene SOLO la información necesaria
\end{itemize}
\end{intuicion}

\subsection{Arquitectura del Sistema}

\begin{figure}[htbp]
\centering
\begin{tikzpicture}[
    node distance=1.5cm,
    entity/.style={rectangle, draw, minimum width=2.5cm, minimum height=1cm, align=center},
    process/.style={rectangle, draw, rounded corners, minimum width=3cm, minimum height=0.8cm, fill=yellow!20},
    arrow/.style={->, thick}
]
    % Entidades
    \node[entity, fill=blue!20] (gov) at (0, 3) {Gobierno\\(Emisor)};
    \node[entity, fill=green!20] (user) at (5, 3) {Usuario\\(21 años)};
    \node[entity, fill=orange!20] (bar) at (10, 3) {Bar\\(Verificador)};
    
    % Proceso
    \node[process] (cred) at (2.5, 1.5) {Credencial:\\Birth: 2003-05-15};
    
    \node[process] (proof) at (7.5, 1.5) {ZK Proof:\\Age $\geq$ 21};
    
    % Flujo
    \draw[arrow] (gov) -- node[above, font=\small] {Emite} (cred);
    \draw[arrow] (cred) -- (user);
    \draw[arrow] (user) -- node[above, font=\small] {Genera} (proof);
    \draw[arrow] (proof) -- (bar);
    
    % Lo que NO se revela
    \node at (5, 0) [draw, rectangle, fill=red!10, font=\small] {
        \textbf{NO revelado:}\\
        Nombre, Fecha exacta,\\
        Número de documento,\\
        Dirección, Foto
    };
    
    % Lo que SÍ se prueba
    \node at (10, 0) [draw, rectangle, fill=green!10, font=\small] {
        \textbf{Verificado:}\\
        Edad $\geq$ 21: ✓\\
        Credencial válida: ✓
    };
\end{tikzpicture}
\caption{Flujo de verificación de edad con ZK}
\end{figure}

\subsection{Implementación}

\begin{lstlisting}[language=circom, caption={Circuito completo de verificación de edad}]
pragma circom 2.0.0;

include "circomlib/poseidon.circom";
include "circomlib/comparators.circom";
include "circomlib/bitify.circom";

// Verificar que una fecha es anterior a otra
// Fechas en formato YYYYMMDD como enteros
template DateBefore() {
    signal input date1;  // Fecha a verificar
    signal input date2;  // Fecha límite
    signal output isBefore;
    
    component lt = LessThan(32);
    lt.in[0] <== date1;
    lt.in[1] <== date2;
    
    isBefore <== lt.out;
}

// Calcular fecha hace N años
template DateMinusYears(years) {
    signal input currentDate;  // YYYYMMDD
    signal output resultDate;
    
    // Extraer componentes
    signal year;
    signal monthDay;
    
    year <-- currentDate \ 10000;
    monthDay <-- currentDate % 10000;
    
    // Verificar descomposición
    currentDate === year * 10000 + monthDay;
    
    // Restar años
    signal newYear;
    newYear <== year - years;
    
    // Reconstruir fecha
    resultDate <== newYear * 10000 + monthDay;
}

// Verificación completa de edad
template AgeVerificationFull(treeLevels) {
    // ========== PUBLIC INPUTS ==========
    signal input currentDate;           // Fecha actual (YYYYMMDD)
    signal input minimumAge;            // Edad mínima requerida
    signal input issuerPubKeyHash;      // Hash de la clave pública del emisor
    signal input credentialMerkleRoot;  // Root del árbol de credenciales
    
    // ========== PRIVATE INPUTS ==========
    signal input birthDate;             // Fecha de nacimiento (YYYYMMDD)
    signal input credentialSecret;      // Secreto de la credencial
    signal input issuerSignature[2];    // Firma del emisor
    signal input merkleProof[treeLevels];
    signal input merkleIndices[treeLevels];
    
    // Otros atributos de la credencial (ocultos)
    signal input name;
    signal input documentNumber;
    signal input country;
    
    // ========== OUTPUT ==========
    signal output isOldEnough;
    
    // ========== VERIFICACIONES ==========
    
    // 1. Calcular commitment de la credencial
    component credentialHash = Poseidon(5);
    credentialHash.inputs[0] <== birthDate;
    credentialHash.inputs[1] <== name;
    credentialHash.inputs[2] <== documentNumber;
    credentialHash.inputs[3] <== country;
    credentialHash.inputs[4] <== credentialSecret;
    
    signal credentialCommitment;
    credentialCommitment <== credentialHash.out;
    
    // 2. Verificar que la credencial está en el árbol
    component merkleVerifier = MerkleProof(treeLevels);
    merkleVerifier.leaf <== credentialCommitment;
    for (var i = 0; i < treeLevels; i++) {
        merkleVerifier.siblings[i] <== merkleProof[i];
        merkleVerifier.pathIndices[i] <== merkleIndices[i];
    }
    merkleVerifier.root === credentialMerkleRoot;
    
    // 3. Verificar firma del emisor sobre el commitment
    // (simplificado - en producción usar EdDSA o similar)
    component sigVerifier = Poseidon(3);
    sigVerifier.inputs[0] <== credentialCommitment;
    sigVerifier.inputs[1] <== issuerSignature[0];
    sigVerifier.inputs[2] <== issuerSignature[1];
    // sigVerifier.out debería relacionarse con issuerPubKeyHash
    
    // 4. Calcular la fecha de corte (currentDate - minimumAge años)
    component cutoffCalc = DateMinusYears(1);  // Parametrizado
    cutoffCalc.currentDate <== currentDate;
    
    // Para minimumAge dinámico, necesitaríamos lógica más compleja
    // Simplificación: calcular cutoff fuera del circuito
    signal cutoffDate;
    cutoffDate <-- currentDate - minimumAge * 10000;  // Aproximación
    
    // Verificar que la aproximación es razonable
    component yearCheck = LessThan(32);
    yearCheck.in[0] <== cutoffDate;
    yearCheck.in[1] <== currentDate;
    yearCheck.out === 1;
    
    // 5. Verificar que birthDate < cutoffDate
    component ageCheck = DateBefore();
    ageCheck.date1 <== birthDate;
    ageCheck.date2 <== cutoffDate;
    
    isOldEnough <== ageCheck.isBefore;
    
    // 6. Forzar que sea suficientemente mayor
    isOldEnough === 1;
    
    // 7. Validaciones adicionales
    // birthDate es fecha válida (entre 1900 y currentDate)
    component birthValid = LessThan(32);
    birthValid.in[0] <== birthDate;
    birthValid.in[1] <== currentDate;
    birthValid.out === 1;
    
    component birthNotTooOld = GreaterThan(32);
    birthNotTooOld.in[0] <== birthDate;
    birthNotTooOld.in[1] <== 19000101;
    birthNotTooOld.out === 1;
}

component main {public [currentDate, minimumAge, issuerPubKeyHash, credentialMerkleRoot]} = 
    AgeVerificationFull(20);
\end{lstlisting}

\begin{lstlisting}[language=Python, caption={Sistema de Age Verification completo}]
from dataclasses import dataclass
from datetime import date, datetime
from typing import Optional, Dict
from hashlib import sha256
import json

@dataclass
class AgeCredential:
    """Credencial de edad emitida por autoridad."""
    subject_id: str  # DID del sujeto
    birth_date: date
    issuer: str  # DID del emisor
    issued_at: datetime
    expiration: Optional[datetime]
    signature: bytes
    
    # Atributos adicionales (privados)
    full_name: str
    document_number: str
    country: str

class AgeVerificationSystem:
    """Sistema completo de verificación de edad con ZK."""
    
    def __init__(self):
        self.authorized_issuers: Dict[str, bytes] = {}  # DID -> public key
        self.credential_tree: list = []  # Merkle tree de credenciales
        
    def register_issuer(self, issuer_did: str, public_key: bytes):
        """Registrar emisor autorizado."""
        self.authorized_issuers[issuer_did] = public_key
    
    def issue_credential(
        self,
        issuer_did: str,
        issuer_private_key: bytes,
        subject_did: str,
        birth_date: date,
        full_name: str,
        document_number: str,
        country: str
    ) -> AgeCredential:
        """Emitir credencial de edad."""
        if issuer_did not in self.authorized_issuers:
            raise ValueError("Unauthorized issuer")
        
        # Crear credencial
        credential = AgeCredential(
            subject_id=subject_did,
            birth_date=birth_date,
            issuer=issuer_did,
            issued_at=datetime.now(),
            expiration=None,
            signature=b'',
            full_name=full_name,
            document_number=document_number,
            country=country
        )
        
        # Firmar
        credential.signature = self._sign_credential(
            credential, issuer_private_key
        )
        
        # Añadir al árbol de credenciales
        commitment = self._compute_commitment(credential)
        self.credential_tree.append(commitment)
        
        return credential
    
    def _compute_commitment(self, cred: AgeCredential) -> bytes:
        """Calcular commitment de credencial."""
        data = json.dumps({
            'birth_date': cred.birth_date.isoformat(),
            'name_hash': sha256(cred.full_name.encode()).hexdigest(),
            'doc_hash': sha256(cred.document_number.encode()).hexdigest(),
            'country': cred.country,
            'issuer': cred.issuer
        }, sort_keys=True)
        return sha256(data.encode()).digest()
    
    def _sign_credential(self, cred: AgeCredential, private_key: bytes) -> bytes:
        """Firmar credencial (simplificado)."""
        commitment = self._compute_commitment(cred)
        return sha256(commitment + private_key).digest()
    
    def create_age_proof(
        self,
        credential: AgeCredential,
        minimum_age: int,
        current_date: date,
        verifier_challenge: bytes
    ) -> Dict:
        """
        Crear ZK proof de edad.
        
        Returns:
            Proof que demuestra edad >= minimum_age sin revelar fecha exacta
        """
        # Calcular edad
        age = self._calculate_age(credential.birth_date, current_date)
        
        if age < minimum_age:
            return {'success': False, 'error': 'Age requirement not met'}
        
        # Calcular fecha de corte
        cutoff_year = current_date.year - minimum_age
        cutoff_date = date(cutoff_year, current_date.month, current_date.day)
        
        # Generar proof (en producción: SNARK real)
        proof_data = {
            'statement': f'birth_date < {cutoff_date.isoformat()}',
            'current_date': current_date.isoformat(),
            'minimum_age': minimum_age,
            'issuer': credential.issuer,
            'challenge': verifier_challenge.hex()
        }
        
        # El proof NO contiene:
        # - Fecha exacta de nacimiento
        # - Nombre
        # - Número de documento
        # Solo prueba: credencial válida Y edad suficiente
        
        proof_bytes = sha256(json.dumps(proof_data).encode()).digest()
        
        return {
            'success': True,
            'proof': {
                'zk_proof': proof_bytes.hex(),
                'public_inputs': {
                    'minimum_age': minimum_age,
                    'current_date': current_date.isoformat(),
                    'issuer_hash': sha256(credential.issuer.encode()).hexdigest(),
                    'challenge': verifier_challenge.hex()
                }
            }
        }
    
    def verify_age_proof(
        self,
        proof: Dict,
        minimum_age: int,
        current_date: date,
        challenge: bytes
    ) -> Dict:
        """Verificar ZK proof de edad."""
        public_inputs = proof.get('public_inputs', {})
        
        # Verificar inputs públicos
        if public_inputs.get('minimum_age') != minimum_age:
            return {'valid': False, 'error': 'Age mismatch'}
        
        if public_inputs.get('current_date') != current_date.isoformat():
            return {'valid': False, 'error': 'Date mismatch'}
        
        if public_inputs.get('challenge') != challenge.hex():
            return {'valid': False, 'error': 'Challenge mismatch'}
        
        # Verificar que el issuer es autorizado
        # (en producción: verificar contra lista de issuers)
        
        # Verificar ZK proof (en producción: verificación SNARK)
        
        return {
            'valid': True,
            'verified_claim': f'Age >= {minimum_age}',
            'issuer_verified': True
        }
    
    def _calculate_age(self, birth_date: date, current_date: date) -> int:
        """Calcular edad en años completos."""
        age = current_date.year - birth_date.year
        if (current_date.month, current_date.day) < (birth_date.month, birth_date.day):
            age -= 1
        return age


def age_verification_demo():
    print("=== Age Verification Demo ===\n")
    
    system = AgeVerificationSystem()
    
    # Registrar emisor (gobierno)
    gov_did = "did:example:government"
    gov_privkey = b"government_secret_key"
    gov_pubkey = sha256(gov_privkey).digest()
    system.register_issuer(gov_did, gov_pubkey)
    print(f"Registered issuer: {gov_did}")
    
    # Emitir credencial a usuario
    user_did = "did:example:alice"
    credential = system.issue_credential(
        issuer_did=gov_did,
        issuer_private_key=gov_privkey,
        subject_did=user_did,
        birth_date=date(2000, 5, 15),  # 24 años en 2024
        full_name="Alice Johnson",
        document_number="ABC123456",
        country="US"
    )
    print(f"\nIssued credential to: {user_did}")
    print(f"Birth date: {credential.birth_date}")
    
    # Escenario 1: Bar verifica edad >= 21
    print("\n--- Scenario 1: Bar (21+) ---")
    
    current = date(2024, 12, 15)
    challenge1 = b"bar_challenge_12345"
    
    proof1 = system.create_age_proof(
        credential=credential,
        minimum_age=21,
        current_date=current,
        verifier_challenge=challenge1
    )
    
    print(f"Proof created: {proof1['success']}")
    print(f"Public inputs: {list(proof1['proof']['public_inputs'].keys())}")
    print("NOT revealed: name, exact birth date, document number")
    
    result1 = system.verify_age_proof(proof1['proof'], 21, current, challenge1)
    print(f"Verification: {result1['valid']}")
    print(f"Claim verified: {result1['verified_claim']}")
    
    # Escenario 2: Casino verifica edad >= 25
    print("\n--- Scenario 2: Casino (25+) ---")
    
    challenge2 = b"casino_challenge_67890"
    proof2 = system.create_age_proof(
        credential=credential,
        minimum_age=25,
        current_date=current,
        verifier_challenge=challenge2
    )
    
    # Alice tiene 24, no cumple >= 25
    print(f"Proof created: {proof2['success']}")
    if not proof2['success']:
        print(f"Error: {proof2.get('error')}")
    
    # Escenario 3: Verificación fallida por challenge incorrecto
    print("\n--- Scenario 3: Replay attack ---")
    
    wrong_challenge = b"attacker_challenge"
    result3 = system.verify_age_proof(proof1['proof'], 21, current, wrong_challenge)
    print(f"Replay attack result: {result3['valid']}")
    print(f"Error: {result3.get('error')}")

age_verification_demo()
\end{lstlisting}

%===========================================
\section{KYC Preservando Privacidad}
\label{sec:privacy-kyc}
%===========================================

\subsection{Problemas del KYC Tradicional}

\begin{advertencia}
\textbf{Problemas del KYC tradicional:}

\begin{enumerate}
    \item \textbf{Replicación de datos:}
    \begin{itemize}
        \item Cada exchange, banco, servicio tiene copia completa
        \item Documentos sensibles en múltiples bases de datos
        \item Multiplicación del riesgo de breach
    \end{itemize}
    
    \item \textbf{Data breaches:}
    \begin{itemize}
        \item Exchanges hackeados: datos KYC robados
        \item Documentos de identidad en mercado negro
        \item Robo de identidad
    \end{itemize}
    
    \item \textbf{Surveillance:}
    \begin{itemize}
        \item Gobiernos pueden solicitar todos los datos
        \item Perfiles completos de actividad financiera
        \item Sin privacidad real para ciudadanos
    \end{itemize}
    
    \item \textbf{Fricciones:}
    \begin{itemize}
        \item Repetir proceso en cada servicio
        \item Documentos físicos, video calls
        \item Días/semanas de espera
    \end{itemize}
\end{enumerate}
\end{advertencia}

\subsection{Modelo de KYC con ZK}

\begin{keyidea}
\textbf{Privacy-Preserving KYC:}

\begin{enumerate}
    \item \textbf{KYC una vez:}
    \begin{itemize}
        \item Usuario completa KYC con proveedor certificado
        \item Proveedor verifica documentos reales
        \item Usuario recibe credencial ZK firmada
    \end{itemize}
    
    \item \textbf{Uso con ZK:}
    \begin{itemize}
        \item Para usar servicio: genera ZK proof
        \item Proof demuestra: ``KYC válido'' sin datos
        \item Servicio NO recibe documentos
    \end{itemize}
    
    \item \textbf{Selective disclosure:}
    \begin{itemize}
        \item Exchange: probar ``no PEP, no sancionado''
        \item DeFi: probar ``residente de país permitido''
        \item Casino: probar ``mayor de 21''
    \end{itemize}
\end{enumerate}
\end{keyidea}

\begin{figure}[htbp]
\centering
\begin{tikzpicture}[
    node distance=1cm,
    entity/.style={rectangle, draw, minimum width=2.5cm, minimum height=1.2cm, align=center},
    data/.style={rectangle, draw, rounded corners, minimum width=2cm, minimum height=0.6cm, fill=red!20},
    proof/.style={rectangle, draw, rounded corners, minimum width=2cm, minimum height=0.6cm, fill=green!20},
    arrow/.style={->, thick},
    noarrow/.style={->, thick, dashed, red}
]
    % KYC Tradicional
    \node at (-4, 4) {\textbf{KYC Tradicional}};
    
    \node[entity, fill=blue!10] (user1) at (-4, 2.5) {Usuario};
    \node[entity, fill=orange!10] (ex1) at (-6, 0.5) {Exchange A};
    \node[entity, fill=orange!10] (ex2) at (-4, 0.5) {Exchange B};
    \node[entity, fill=orange!10] (ex3) at (-2, 0.5) {Bank C};
    
    \node[data] (d1) at (-6, 1.5) {Docs};
    \node[data] (d2) at (-4, 1.5) {Docs};
    \node[data] (d3) at (-2, 1.5) {Docs};
    
    \draw[arrow] (user1) -- (d1);
    \draw[arrow] (user1) -- (d2);
    \draw[arrow] (user1) -- (d3);
    \draw[arrow] (d1) -- (ex1);
    \draw[arrow] (d2) -- (ex2);
    \draw[arrow] (d3) -- (ex3);
    
    \node at (-4, -0.5) [font=\small, red] {Datos replicados 3 veces};
    
    % KYC con ZK
    \node at (4, 4) {\textbf{KYC con ZK}};
    
    \node[entity, fill=blue!10] (user2) at (4, 2.5) {Usuario};
    \node[entity, fill=purple!10] (kyc) at (2, 1) {KYC Provider};
    
    \node[entity, fill=orange!10] (ex4) at (4, -0.5) {Exchange A};
    \node[entity, fill=orange!10] (ex5) at (6, -0.5) {Bank B};
    
    \node[data] (d4) at (2, 2) {Docs};
    \node[proof] (p1) at (4, 1) {ZK Cred};
    
    \draw[arrow] (user2) -- (d4);
    \draw[arrow] (d4) -- (kyc);
    \draw[arrow] (kyc) -- (p1);
    \draw[arrow, green!60!black] (p1) -- node[right, font=\small] {Proof} (ex4);
    \draw[arrow, green!60!black] (p1) -- node[right, font=\small] {Proof} (ex5);
    
    \draw[noarrow] (user2) -- (ex4);
    \draw[noarrow] (user2) -- (ex5);
    
    \node at (4, -1.5) [font=\small, green!60!black] {Solo proofs, sin datos};
\end{tikzpicture}
\caption{KYC tradicional vs KYC con ZK}
\end{figure}

\begin{lstlisting}[language=Python, caption={Sistema de Privacy-Preserving KYC}]
from dataclasses import dataclass
from datetime import date, datetime
from typing import Dict, List, Optional, Set
from hashlib import sha256
from enum import Enum
import json

class RiskLevel(Enum):
    LOW = "low"
    MEDIUM = "medium"
    HIGH = "high"

class KYCStatus(Enum):
    PENDING = "pending"
    VERIFIED = "verified"
    REJECTED = "rejected"
    EXPIRED = "expired"

@dataclass
class KYCCredential:
    """Credencial KYC emitida por proveedor certificado."""
    credential_id: str
    subject_did: str
    provider_did: str
    
    # Información verificada (privada)
    full_name: str
    birth_date: date
    nationality: str
    document_type: str
    document_number: str
    address_country: str
    
    # Metadatos de verificación
    verification_date: datetime
    expiration_date: datetime
    verification_level: int  # 1, 2, 3
    
    # Flags de compliance
    is_pep: bool  # Politically Exposed Person
    is_sanctioned: bool
    risk_level: RiskLevel
    
    # Firma del proveedor
    signature: bytes

class PrivacyKYCSystem:
    """Sistema de KYC preservando privacidad."""
    
    def __init__(self):
        self.certified_providers: Dict[str, bytes] = {}
        self.sanctioned_countries: Set[str] = {"NK", "IR", "SY"}  # Ejemplo
        self.credential_tree: list = []
    
    def register_provider(self, provider_did: str, public_key: bytes):
        """Registrar proveedor KYC certificado."""
        self.certified_providers[provider_did] = public_key
    
    def issue_kyc_credential(
        self,
        provider_did: str,
        provider_key: bytes,
        user_did: str,
        user_data: Dict
    ) -> KYCCredential:
        """
        Proveedor emite credencial KYC después de verificación.
        """
        if provider_did not in self.certified_providers:
            raise ValueError("Uncertified provider")
        
        # Verificaciones del proveedor (en producción: proceso real)
        is_pep = self._check_pep(user_data['full_name'])
        is_sanctioned = self._check_sanctions(
            user_data['full_name'], 
            user_data['nationality']
        )
        risk_level = self._assess_risk(user_data)
        
        credential = KYCCredential(
            credential_id=sha256(f"{user_did}:{datetime.now()}".encode()).hexdigest()[:16],
            subject_did=user_did,
            provider_did=provider_did,
            full_name=user_data['full_name'],
            birth_date=user_data['birth_date'],
            nationality=user_data['nationality'],
            document_type=user_data['document_type'],
            document_number=user_data['document_number'],
            address_country=user_data['address_country'],
            verification_date=datetime.now(),
            expiration_date=datetime(2025, 12, 31),
            verification_level=user_data.get('level', 2),
            is_pep=is_pep,
            is_sanctioned=is_sanctioned,
            risk_level=risk_level,
            signature=b''
        )
        
        # Firmar credencial
        credential.signature = self._sign_credential(credential, provider_key)
        
        # Añadir al árbol
        self.credential_tree.append(self._commitment(credential))
        
        return credential
    
    def create_compliance_proof(
        self,
        credential: KYCCredential,
        requirements: Dict,
        verifier_challenge: bytes
    ) -> Dict:
        """
        Crear ZK proof de compliance.
        
        Args:
            credential: KYC credential del usuario
            requirements: Lo que el verifier necesita
                - 'not_sanctioned': bool
                - 'not_pep': bool
                - 'min_verification_level': int
                - 'allowed_countries': List[str]
                - 'min_age': int
            verifier_challenge: Challenge del verifier
        
        Returns:
            ZK proof que demuestra compliance sin revelar datos
        """
        # Verificar que credencial no está expirada
        if credential.expiration_date < datetime.now():
            return {'success': False, 'error': 'Credential expired'}
        
        # Verificar cada requirement
        compliance_results = {}
        
        if requirements.get('not_sanctioned', False):
            compliance_results['not_sanctioned'] = not credential.is_sanctioned
            
        if requirements.get('not_pep', False):
            compliance_results['not_pep'] = not credential.is_pep
            
        if 'min_verification_level' in requirements:
            compliance_results['verification_level_met'] = (
                credential.verification_level >= requirements['min_verification_level']
            )
            
        if 'allowed_countries' in requirements:
            compliance_results['country_allowed'] = (
                credential.nationality in requirements['allowed_countries']
            )
            
        if 'min_age' in requirements:
            age = self._calculate_age(credential.birth_date)
            compliance_results['age_requirement_met'] = age >= requirements['min_age']
        
        # Verificar que todos los requirements se cumplen
        all_met = all(compliance_results.values())
        
        if not all_met:
            failed = [k for k, v in compliance_results.items() if not v]
            return {'success': False, 'error': f'Requirements not met: {failed}'}
        
        # Generar ZK proof
        proof = {
            'zk_proof': self._generate_zk_proof(
                credential, requirements, verifier_challenge
            ),
            'public_inputs': {
                'provider_hash': sha256(credential.provider_did.encode()).hexdigest(),
                'expiration': credential.expiration_date.isoformat(),
                'requirements_hash': sha256(json.dumps(requirements, sort_keys=True).encode()).hexdigest(),
                'challenge': verifier_challenge.hex(),
                'compliance_flags': list(compliance_results.keys())
            }
        }
        
        return {'success': True, 'proof': proof}
    
    def verify_compliance_proof(
        self,
        proof: Dict,
        requirements: Dict,
        challenge: bytes
    ) -> Dict:
        """Verificar ZK proof de compliance."""
        public_inputs = proof.get('public_inputs', {})
        
        # Verificar challenge
        if bytes.fromhex(public_inputs.get('challenge', '')) != challenge:
            return {'valid': False, 'error': 'Invalid challenge'}
        
        # Verificar requirements match
        expected_hash = sha256(json.dumps(requirements, sort_keys=True).encode()).hexdigest()
        if public_inputs.get('requirements_hash') != expected_hash:
            return {'valid': False, 'error': 'Requirements mismatch'}
        
        # Verificar no expirado
        expiration = datetime.fromisoformat(public_inputs.get('expiration', '2000-01-01'))
        if expiration < datetime.now():
            return {'valid': False, 'error': 'Credential expired'}
        
        # Verificar ZK proof (placeholder)
        # En producción: verificación SNARK real
        
        return {
            'valid': True,
            'verified_compliance': public_inputs.get('compliance_flags', []),
            'provider_verified': True,
            'expiration': public_inputs.get('expiration')
        }
    
    def _check_pep(self, name: str) -> bool:
        """Verificar si es PEP (simplificado)."""
        # En producción: consultar bases de datos de PEPs
        return False
    
    def _check_sanctions(self, name: str, country: str) -> bool:
        """Verificar listas de sanciones."""
        # En producción: OFAC, UN, EU sanctions lists
        return country in self.sanctioned_countries
    
    def _assess_risk(self, data: Dict) -> RiskLevel:
        """Evaluar nivel de riesgo."""
        if data['nationality'] in self.sanctioned_countries:
            return RiskLevel.HIGH
        return RiskLevel.LOW
    
    def _commitment(self, cred: KYCCredential) -> bytes:
        """Calcular commitment de credencial."""
        data = f"{cred.subject_did}:{cred.verification_date}"
        return sha256(data.encode()).digest()
    
    def _sign_credential(self, cred: KYCCredential, key: bytes) -> bytes:
        """Firmar credencial."""
        return sha256(self._commitment(cred) + key).digest()
    
    def _generate_zk_proof(self, cred, reqs, challenge) -> str:
        """Generar ZK proof (placeholder)."""
        data = f"{cred.credential_id}:{json.dumps(reqs)}:{challenge.hex()}"
        return sha256(data.encode()).hexdigest()
    
    def _calculate_age(self, birth_date: date) -> int:
        today = date.today()
        return today.year - birth_date.year - (
            (today.month, today.day) < (birth_date.month, birth_date.day)
        )


def privacy_kyc_demo():
    print("=== Privacy-Preserving KYC Demo ===\n")
    
    system = PrivacyKYCSystem()
    
    # Registrar proveedor KYC
    provider_did = "did:example:kycprovider"
    provider_key = b"provider_secret"
    system.register_provider(provider_did, sha256(provider_key).digest())
    print(f"Registered KYC provider: {provider_did}")
    
    # Usuario completa KYC una vez
    user_did = "did:example:alice"
    user_data = {
        'full_name': 'Alice Johnson',
        'birth_date': date(1990, 3, 15),
        'nationality': 'US',
        'document_type': 'passport',
        'document_number': 'AB1234567',
        'address_country': 'US',
        'level': 2
    }
    
    credential = system.issue_kyc_credential(
        provider_did, provider_key, user_did, user_data
    )
    print(f"\nKYC credential issued to: {user_did}")
    print(f"Verification level: {credential.verification_level}")
    print(f"Expires: {credential.expiration_date}")
    
    # Escenario 1: Exchange quiere verificar no-sanctions + age
    print("\n--- Scenario 1: Exchange Verification ---")
    
    exchange_requirements = {
        'not_sanctioned': True,
        'not_pep': True,
        'min_age': 18,
        'allowed_countries': ['US', 'UK', 'DE', 'FR', 'ES', 'JP']
    }
    
    challenge1 = b"exchange_challenge_abc"
    
    proof1 = system.create_compliance_proof(
        credential, exchange_requirements, challenge1
    )
    
    print(f"Proof created: {proof1['success']}")
    print(f"Compliance flags: {proof1['proof']['public_inputs']['compliance_flags']}")
    print("\nNOT revealed to exchange:")
    print("  - Full name")
    print("  - Birth date")
    print("  - Document number")
    print("  - Address")
    
    result1 = system.verify_compliance_proof(
        proof1['proof'], exchange_requirements, challenge1
    )
    print(f"\nExchange verification: {result1['valid']}")
    
    # Escenario 2: DeFi solo necesita not-sanctioned
    print("\n--- Scenario 2: DeFi (minimal requirements) ---")
    
    defi_requirements = {
        'not_sanctioned': True
    }
    
    challenge2 = b"defi_challenge_xyz"
    
    proof2 = system.create_compliance_proof(
        credential, defi_requirements, challenge2
    )
    
    print(f"Minimal proof created: {proof2['success']}")
    print(f"Only verifies: {proof2['proof']['public_inputs']['compliance_flags']}")
    print("Even less data exposure than exchange scenario!")

privacy_kyc_demo()
\end{lstlisting}

%===========================================
\section*{Resumen del Bloque}
%===========================================

\begin{keyidea}
\textbf{Puntos clave de Identidad y Credenciales Anónimas (Bloque 14):}

\textbf{1. Anonymous Credentials:}
\begin{itemize}
    \item Credenciales verificables sin revelar identidad
    \item Sistemas: Idemix (CL), U-Prove, BBS+
    \item Propiedades: Unforgeability, anonymity, unlinkability
\end{itemize}

\textbf{2. Selective Disclosure:}
\begin{itemize}
    \item Revelar solo atributos necesarios
    \item Predicados: range proofs, comparaciones
    \item Implementaciones: Merkle trees, BBS+ proofs
\end{itemize}

\textbf{3. ZK Identity Proofs:}
\begin{itemize}
    \item Membership, non-membership, attribute proofs
    \item Semaphore: voting, whistleblowing, rate limiting
    \item Nullifiers para prevenir doble uso
\end{itemize}

\textbf{4. Soulbound Tokens + ZK:}
\begin{itemize}
    \item SBTs públicos = problemas de privacidad
    \item Solución: Private SBTs con ownership proofs
    \item Selective disclosure de atributos de SBT
\end{itemize}

\textbf{5. DIDs + VCs + ZK:}
\begin{itemize}
    \item DIDs como ancla de identidad descentralizada
    \item VCs como credenciales verificables
    \item ZK para presentaciones privadas
\end{itemize}

\textbf{6. Proof of Humanity:}
\begin{itemize}
    \item Worldcoin: biometría + ZK
    \item BrightID: grafo social
    \item Trade-off: unicidad vs privacidad vs accesibilidad
\end{itemize}

\textbf{7. Age Verification:}
\begin{itemize}
    \item Probar edad $\geq$ umbral sin revelar fecha exacta
    \item Circuitos de comparación de fechas
    \item Aplicaciones: alcohol, gambling, contenido adulto
\end{itemize}

\textbf{8. Privacy-Preserving KYC:}
\begin{itemize}
    \item KYC una vez, usar con ZK múltiples veces
    \item No replicación de documentos sensibles
    \item Compliance sin surveillance
\end{itemize}
\end{keyidea}

\begin{bancozk}
\textbf{Sistema de Identidad en Banco ZK:}

\begin{center}
\begin{tikzpicture}[
    layer/.style={rectangle, draw, minimum width=10cm, minimum height=0.8cm, align=center}
]
    \node[layer, fill=blue!20] at (0, 3) {Capa Usuario: Wallet con credenciales ZK};
    \node[layer, fill=green!20] at (0, 2) {Capa Credencial: KYC + Age + Country + Humanity};
    \node[layer, fill=yellow!20] at (0, 1) {Capa Proof: Selective disclosure + Predicados};
    \node[layer, fill=orange!20] at (0, 0) {Capa Verificación: On-chain verifiers};
    \node[layer, fill=purple!20] at (0, -1) {Capa Pool: Acceso condicional a funcionalidades};
\end{tikzpicture}
\end{center}

\textbf{Flujo de uso:}
\begin{enumerate}
    \item Usuario obtiene credenciales de proveedores certificados
    \item Almacena en wallet personal (nunca en Banco ZK)
    \item Para operar: genera ZK proof según requerimientos
    \item Verificador on-chain valida proof
    \item Usuario accede a funcionalidad sin revelar datos
\end{enumerate}

\textbf{Casos de uso:}
\begin{itemize}
    \item \textbf{Depósitos grandes:} Proof de KYC válido
    \item \textbf{Pools premium:} Proof de país permitido
    \item \textbf{Governance:} Proof of humanity (1 persona = 1 voto)
    \item \textbf{Compliance:} Proof de no-sanctions para withdrawals
\end{itemize}
\end{bancozk}

\begin{table}[htbp]
\centering
\small
\begin{tabular}{lcccc}
\toprule
\textbf{Técnica} & \textbf{Privacidad} & \textbf{Verificabilidad} & \textbf{Complejidad} & \textbf{Madurez} \\
\midrule
BBS+ Signatures & Alta & Alta & Media & Alta \\
Semaphore & Alta & Alta & Media & Alta \\
Private SBTs & Alta & Alta & Alta & Media \\
ZK-VCs & Alta & Alta & Alta & Media \\
Worldcoin & Media & Alta & Alta & Media \\
Privacy KYC & Alta & Alta & Alta & Baja \\
\bottomrule
\end{tabular}
\caption{Resumen de tecnologías de identidad ZK}
\end{table}

% Fin del Bloque 14
\cleardoublepage

% ============================================
% BLOQUE 15: COMPLIANCE VS PRIVACIDAD - VERSIÓN EXPANDIDA
% Parte 1: Introducción, Proof-of-Innocence, Selective De-anonymization
% Objetivo: ~60 páginas (~3600 líneas)
% ============================================

\chapter{Compliance vs Privacidad}
\label{ch:bloque15}

\begin{center}
\textit{Bloque 15 de 20}
\end{center}

\section*{Introducción}

El debate entre compliance regulatorio y privacidad financiera es uno de los más importantes en el ecosistema crypto. Este bloque explora cómo las pruebas de conocimiento cero pueden resolver esta aparente dicotomía, permitiendo cumplimiento normativo sin sacrificar la privacidad de usuarios legítimos.

\begin{keyidea}
\textbf{La falsa dicotomía:}

\textbf{Argumento regulatorio tradicional:}
\begin{itemize}
    \item ``Privacidad = facilitar crimen''
    \item ``Transparencia total necesaria para compliance''
    \item ``Si no tienes nada que ocultar...''
\end{itemize}

\textbf{Realidad con ZK:}
\begin{itemize}
    \item Privacidad Y compliance son compatibles
    \item ZK permite probar cumplimiento sin revelar datos
    \item La vigilancia masiva no es necesaria para AML/CFT
    \item Privacidad es derecho fundamental, no indicador de criminalidad
\end{itemize}
\end{keyidea}

\begin{figure}[htbp]
\centering
\begin{tikzpicture}[
    node distance=1cm,
    box/.style={rectangle, draw, minimum width=3cm, minimum height=1.5cm, align=center}
]
    % Espectro
    \draw[thick, ->] (0, 0) -- (12, 0);
    \node at (0, -0.5) {Privacidad Total};
    \node at (12, -0.5) {Transparencia Total};
    
    % Posiciones
    \node[box, fill=red!20] at (1, 2) {Monero\\(Privacidad máxima)};
    \node[box, fill=orange!20] at (4, 2) {Privacy Pools\\(Balance ajustable)};
    \node[box, fill=yellow!20] at (7, 2) {ZK Compliance\\(Privacidad + Proof)};
    \node[box, fill=green!20] at (10, 2) {Bitcoin\\(Pseudónimo)};
    
    % Flechas al espectro
    \draw[dashed] (1, 1.2) -- (1, 0.2);
    \draw[dashed] (4, 1.2) -- (5, 0.2);
    \draw[dashed] (7, 1.2) -- (7, 0.2);
    \draw[dashed] (10, 1.2) -- (10, 0.2);
    
    % Zona óptima
    \draw[thick, green!60!black] (5, -0.2) -- (5, 0.2);
    \draw[thick, green!60!black] (8, -0.2) -- (8, 0.2);
    \draw[<->, thick, green!60!black] (5, 0.5) -- (8, 0.5);
    \node at (6.5, 0.9) [green!60!black, font=\small] {Zona óptima con ZK};
\end{tikzpicture}
\caption{Espectro de privacidad-compliance y la zona óptima habilitada por ZK}
\label{fig:privacy-compliance-spectrum}
\end{figure}

\begin{mathnote}[Marco Regulatorio Relevante]
\textbf{Principales regulaciones:}
\begin{itemize}
    \item \textbf{FATF Travel Rule:} Compartir info de origen/destino para txs $>$ \$1,000
    \item \textbf{AML (Anti-Money Laundering):} Prevención de lavado de dinero
    \item \textbf{KYC (Know Your Customer):} Identificación de clientes
    \item \textbf{CFT (Counter-Financing of Terrorism):} Anti-financiamiento terrorista
    \item \textbf{OFAC Sanctions:} Lista de entidades sancionadas (EEUU)
    \item \textbf{EU MiCA:} Regulación de criptoactivos en Europa
    \item \textbf{GDPR:} Protección de datos personales (tensión con transparencia)
\end{itemize}

\textbf{El desafío:} Cumplir con AML/KYC/FATF mientras se respeta GDPR y el derecho a la privacidad.
\end{mathnote}

%===========================================
\section{Proof-of-Innocence}
\label{sec:proof-of-innocence}
%===========================================

\subsection{Concepto y Motivación}

\begin{definition}[Proof-of-Innocence]
Prueba criptográfica de que los fondos de un usuario no provienen de fuentes ilícitas, sin revelar el origen específico de los fondos ni la identidad del usuario.

\textbf{Formalmente:}
\[
\text{PoI} = \text{ZKProof}\left(\exists \text{ deposit } d : d \in \text{LegitSet} \land \text{user owns } d\right)
\]

Sin revelar cuál $d$ específico.
\end{definition}

\begin{intuicion}
Imagina que usaste un mixer hace 2 años. Ahora un regulador te pregunta: ``¿De dónde vienen esos fondos?''

\textbf{Sin ZK:}
\begin{itemize}
    \item Debes revelar tu depósito específico
    \item Pierdes toda la privacidad del mixer
    \item O te niegas y pareces sospechoso
\end{itemize}

\textbf{Con Proof-of-Innocence:}
\begin{itemize}
    \item Demuestras: ``Mis fondos estaban en el mixer ANTES del hack de Exchange X''
    \item O: ``Mis fondos NO vienen de direcciones sancionadas''
    \item Sin revelar cuál depósito específico era tuyo
\end{itemize}
\end{intuicion}

\subsection{Tipos de Proof-of-Innocence}

\begin{mathnote}[Categorías de PoI]
\textbf{1. Proof-of-Pre-Event (Temporal):}
\begin{itemize}
    \item ``Mis fondos entraron al pool antes del timestamp T''
    \item Útil cuando un evento específico ``contamina'' el pool
    \item Ejemplo: Hack de exchange, sanción a entidad
\end{itemize}

\textbf{2. Proof-of-Non-Inclusion (Exclusión):}
\begin{itemize}
    \item ``Mis fondos NO vienen de direcciones en lista L''
    \item Requiere sparse Merkle tree o accumulator
    \item Ejemplo: Lista OFAC, direcciones de hacks conocidos
\end{itemize}

\textbf{3. Proof-of-Inclusion (Whitelist):}
\begin{itemize}
    \item ``Mis fondos vienen de un set S de fuentes legítimas''
    \item Set definido por comunidad o regulador
    \item Ejemplo: Solo exchanges regulados, solo direcciones KYC'd
\end{itemize}

\textbf{4. Proof-of-Source-Category:}
\begin{itemize}
    \item ``Mis fondos vienen de categoría C''
    \item Sin revelar cuál fuente específica
    \item Ejemplo: ``De un exchange regulado'' (sin decir cuál)
\end{itemize}
\end{mathnote}

\begin{figure}[htbp]
\centering
\begin{tikzpicture}[
    node distance=0.8cm,
    deposit/.style={circle, draw, minimum size=0.6cm},
    good/.style={circle, draw, minimum size=0.6cm, fill=green!30},
    bad/.style={circle, draw, minimum size=0.6cm, fill=red!30},
    user/.style={circle, draw, minimum size=0.6cm, fill=blue!30},
    pool/.style={ellipse, draw, minimum width=8cm, minimum height=3cm, fill=yellow!10}
]
    % Pool
    \node[pool] (pool) at (0, 0) {};
    \node at (0, 2) {\textbf{Privacy Pool}};
    
    % Depósitos buenos (antes del hack)
    \node[good] (g1) at (-3, 0.5) {};
    \node[good] (g2) at (-2, -0.3) {};
    \node[good] (g3) at (-1, 0.8) {};
    \node[good] (g4) at (0, -0.5) {};
    \node[user] (u) at (1, 0.3) {};
    
    % Depósitos malos (después del hack o de fuente ilícita)
    \node[bad] (b1) at (2, -0.4) {};
    \node[bad] (b2) at (3, 0.6) {};
    
    % Timeline
    \draw[thick, ->] (-4, -2.5) -- (4, -2.5);
    \node at (-4, -2.9) [font=\small] {T=0};
    \node at (0, -2.9) [font=\small] {Hack};
    \node at (4, -2.9) [font=\small] {Ahora};
    
    \draw[dashed, red] (0, -2.3) -- (0, 1.5);
    
    % Leyenda
    \node[good, label=right:{Depósitos legítimos}] at (5, 1) {};
    \node[bad, label=right:{Depósitos ilícitos}] at (5, 0) {};
    \node[user, label=right:{Mi depósito}] at (5, -1) {};
    
    % Proof
    \node at (0, -4) [draw, rectangle, fill=green!10] {
        PoI: ``Mi depósito fue antes del hack''
    };
\end{tikzpicture}
\caption{Proof-of-Innocence temporal: probar que el depósito ocurrió antes de un evento contaminante}
\label{fig:poi-temporal}
\end{figure}

\subsection{Implementación de Proof-of-Innocence}

\begin{lstlisting}[language=circom, caption={Circuito de Proof-of-Innocence temporal}]
pragma circom 2.0.0;

include "circomlib/poseidon.circom";
include "circomlib/comparators.circom";

// Proof-of-Innocence: Probar que un depósito ocurrió antes de timestamp T
template ProofOfInnocenceTemporal(treeLevels) {
    // ========== PUBLIC INPUTS ==========
    signal input merkleRoot;           // Root actual del pool
    signal input nullifierHash;        // Nullifier del withdrawal
    signal input cutoffTimestamp;      // Timestamp del evento (hack, sanción)
    signal input currentTimestamp;     // Timestamp actual
    
    // ========== PRIVATE INPUTS ==========
    signal input depositTimestamp;     // Cuando se hizo el depósito
    signal input depositAmount;        // Monto depositado
    signal input secret;               // Secreto del usuario
    signal input nullifier;            // Nullifier secret
    signal input merkleProof[treeLevels];
    signal input merkleIndices[treeLevels];
    
    // ========== VERIFICACIONES ==========
    
    // 1. Calcular commitment del depósito
    // commitment = H(amount, timestamp, secret, nullifier)
    component commitmentHash = Poseidon(4);
    commitmentHash.inputs[0] <== depositAmount;
    commitmentHash.inputs[1] <== depositTimestamp;
    commitmentHash.inputs[2] <== secret;
    commitmentHash.inputs[3] <== nullifier;
    
    signal commitment;
    commitment <== commitmentHash.out;
    
    // 2. Verificar que commitment está en el Merkle tree
    component merkleVerifier = MerkleProof(treeLevels);
    merkleVerifier.leaf <== commitment;
    for (var i = 0; i < treeLevels; i++) {
        merkleVerifier.siblings[i] <== merkleProof[i];
        merkleVerifier.pathIndices[i] <== merkleIndices[i];
    }
    merkleVerifier.root === merkleRoot;
    
    // 3. Verificar nullifier hash
    component nullifierHasher = Poseidon(2);
    nullifierHasher.inputs[0] <== nullifier;
    nullifierHasher.inputs[1] <== secret;
    nullifierHasher.out === nullifierHash;
    
    // 4. PROOF-OF-INNOCENCE: depósito ANTES del cutoff
    component beforeCutoff = LessThan(64);
    beforeCutoff.in[0] <== depositTimestamp;
    beforeCutoff.in[1] <== cutoffTimestamp;
    beforeCutoff.out === 1;  // Debe ser antes del cutoff
    
    // 5. Verificar que el depósito es válido (timestamp razonable)
    component validDeposit = LessThan(64);
    validDeposit.in[0] <== depositTimestamp;
    validDeposit.in[1] <== currentTimestamp;
    validDeposit.out === 1;  // Depósito antes de ahora
}

component main {public [merkleRoot, nullifierHash, cutoffTimestamp, currentTimestamp]} = 
    ProofOfInnocenceTemporal(20);
\end{lstlisting>

\begin{lstlisting}[language=Python, caption={Sistema completo de Proof-of-Innocence}]
from dataclasses import dataclass
from datetime import datetime
from typing import Dict, List, Set, Optional
from hashlib import sha256
from enum import Enum
import json

class InnocenceType(Enum):
    TEMPORAL = "temporal"           # Antes de evento
    NON_INCLUSION = "non_inclusion" # No en blacklist
    INCLUSION = "inclusion"         # En whitelist
    CATEGORY = "category"           # De categoría específica

@dataclass
class Deposit:
    commitment: bytes
    timestamp: datetime
    amount: int
    source_address: str
    source_category: str  # 'exchange', 'defi', 'personal', etc.
    leaf_index: int

@dataclass
class InnocenceProof:
    proof_type: InnocenceType
    nullifier_hash: bytes
    public_inputs: Dict
    zk_proof: bytes
    metadata: Dict

class ProofOfInnocenceSystem:
    """
    Sistema completo de Proof-of-Innocence para privacy pools.
    """
    
    def __init__(self):
        # Merkle tree de depósitos
        self.deposits: List[Deposit] = []
        self.merkle_root: bytes = b'\x00' * 32
        
        # Blacklists y whitelists
        self.blacklisted_addresses: Set[str] = set()
        self.sanctioned_addresses: Set[str] = set()
        self.whitelisted_sources: Set[str] = set()
        
        # Eventos conocidos (hacks, sanciones)
        self.contamination_events: List[Dict] = []
        
        # Categorías de fuentes
        self.source_categories: Dict[str, Set[str]] = {
            'regulated_exchange': set(),
            'defi_protocol': set(),
            'kyc_verified': set(),
            'unknown': set()
        }
    
    def register_deposit(
        self,
        amount: int,
        source_address: str,
        secret: bytes,
        nullifier: bytes
    ) -> Deposit:
        """Registrar un nuevo depósito en el pool."""
        timestamp = datetime.now()
        
        # Determinar categoría de la fuente
        category = self._categorize_source(source_address)
        
        # Calcular commitment
        commitment = self._compute_commitment(
            amount, timestamp, secret, nullifier
        )
        
        deposit = Deposit(
            commitment=commitment,
            timestamp=timestamp,
            amount=amount,
            source_address=source_address,
            source_category=category,
            leaf_index=len(self.deposits)
        )
        
        self.deposits.append(deposit)
        self._update_merkle_root()
        
        return deposit
    
    def register_contamination_event(
        self,
        event_type: str,
        timestamp: datetime,
        affected_addresses: List[str],
        description: str
    ):
        """Registrar un evento que contamina ciertos fondos."""
        event = {
            'type': event_type,
            'timestamp': timestamp,
            'affected_addresses': affected_addresses,
            'description': description
        }
        self.contamination_events.append(event)
        
        # Añadir a blacklist
        for addr in affected_addresses:
            self.blacklisted_addresses.add(addr)
    
    def add_to_sanctions_list(self, address: str):
        """Añadir dirección a lista de sanciones."""
        self.sanctioned_addresses.add(address)
        self.blacklisted_addresses.add(address)
    
    def add_whitelisted_source(self, address: str, category: str):
        """Añadir fuente a whitelist."""
        self.whitelisted_sources.add(address)
        if category in self.source_categories:
            self.source_categories[category].add(address)
    
    # ========== GENERACIÓN DE PROOFS ==========
    
    def generate_temporal_poi(
        self,
        deposit: Deposit,
        secret: bytes,
        nullifier: bytes,
        cutoff_event: Dict
    ) -> Optional[InnocenceProof]:
        """
        Generar Proof-of-Innocence temporal.
        Demuestra que el depósito fue ANTES del evento de contaminación.
        """
        cutoff_timestamp = cutoff_event['timestamp']
        
        # Verificar que el depósito fue antes
        if deposit.timestamp >= cutoff_timestamp:
            return None  # No puede generar proof
        
        # Generar nullifier hash
        nullifier_hash = self._compute_nullifier_hash(nullifier, secret)
        
        # Generar ZK proof (placeholder)
        public_inputs = {
            'merkle_root': self.merkle_root.hex(),
            'nullifier_hash': nullifier_hash.hex(),
            'cutoff_timestamp': cutoff_timestamp.isoformat(),
            'current_timestamp': datetime.now().isoformat()
        }
        
        zk_proof = self._generate_zk_proof(
            'temporal',
            deposit,
            secret,
            nullifier,
            public_inputs
        )
        
        return InnocenceProof(
            proof_type=InnocenceType.TEMPORAL,
            nullifier_hash=nullifier_hash,
            public_inputs=public_inputs,
            zk_proof=zk_proof,
            metadata={
                'event_type': cutoff_event['type'],
                'event_description': cutoff_event['description']
            }
        )
    
    def generate_non_inclusion_poi(
        self,
        deposit: Deposit,
        secret: bytes,
        nullifier: bytes,
        blacklist_root: bytes
    ) -> Optional[InnocenceProof]:
        """
        Generar Proof-of-Innocence de no-inclusión.
        Demuestra que la fuente NO está en la blacklist.
        """
        # Verificar que la fuente no está blacklisted
        if deposit.source_address in self.blacklisted_addresses:
            return None  # No puede generar proof
        
        nullifier_hash = self._compute_nullifier_hash(nullifier, secret)
        
        public_inputs = {
            'merkle_root': self.merkle_root.hex(),
            'nullifier_hash': nullifier_hash.hex(),
            'blacklist_root': blacklist_root.hex()
        }
        
        zk_proof = self._generate_zk_proof(
            'non_inclusion',
            deposit,
            secret,
            nullifier,
            public_inputs
        )
        
        return InnocenceProof(
            proof_type=InnocenceType.NON_INCLUSION,
            nullifier_hash=nullifier_hash,
            public_inputs=public_inputs,
            zk_proof=zk_proof,
            metadata={'blacklist_size': len(self.blacklisted_addresses)}
        )
    
    def generate_inclusion_poi(
        self,
        deposit: Deposit,
        secret: bytes,
        nullifier: bytes,
        whitelist_name: str
    ) -> Optional[InnocenceProof]:
        """
        Generar Proof-of-Innocence de inclusión.
        Demuestra que la fuente ESTÁ en la whitelist.
        """
        if deposit.source_address not in self.whitelisted_sources:
            return None
        
        nullifier_hash = self._compute_nullifier_hash(nullifier, secret)
        
        # Calcular whitelist root
        whitelist_root = self._compute_set_root(self.whitelisted_sources)
        
        public_inputs = {
            'merkle_root': self.merkle_root.hex(),
            'nullifier_hash': nullifier_hash.hex(),
            'whitelist_root': whitelist_root.hex(),
            'whitelist_name': whitelist_name
        }
        
        zk_proof = self._generate_zk_proof(
            'inclusion',
            deposit,
            secret,
            nullifier,
            public_inputs
        )
        
        return InnocenceProof(
            proof_type=InnocenceType.INCLUSION,
            nullifier_hash=nullifier_hash,
            public_inputs=public_inputs,
            zk_proof=zk_proof,
            metadata={'whitelist_name': whitelist_name}
        )
    
    def generate_category_poi(
        self,
        deposit: Deposit,
        secret: bytes,
        nullifier: bytes,
        required_category: str
    ) -> Optional[InnocenceProof]:
        """
        Generar Proof-of-Innocence de categoría.
        Demuestra que la fuente es de cierta categoría.
        """
        if deposit.source_category != required_category:
            return None
        
        nullifier_hash = self._compute_nullifier_hash(nullifier, secret)
        
        # Root de la categoría
        category_addresses = self.source_categories.get(required_category, set())
        category_root = self._compute_set_root(category_addresses)
        
        public_inputs = {
            'merkle_root': self.merkle_root.hex(),
            'nullifier_hash': nullifier_hash.hex(),
            'category': required_category,
            'category_root': category_root.hex()
        }
        
        zk_proof = self._generate_zk_proof(
            'category',
            deposit,
            secret,
            nullifier,
            public_inputs
        )
        
        return InnocenceProof(
            proof_type=InnocenceType.CATEGORY,
            nullifier_hash=nullifier_hash,
            public_inputs=public_inputs,
            zk_proof=zk_proof,
            metadata={'category': required_category}
        )
    
    # ========== VERIFICACIÓN ==========
    
    def verify_poi(self, proof: InnocenceProof) -> Dict:
        """Verificar un Proof-of-Innocence."""
        # Verificar merkle root es válido
        if proof.public_inputs.get('merkle_root') != self.merkle_root.hex():
            return {'valid': False, 'error': 'Invalid merkle root'}
        
        # Verificar ZK proof (placeholder)
        zk_valid = self._verify_zk_proof(proof)
        
        if not zk_valid:
            return {'valid': False, 'error': 'Invalid ZK proof'}
        
        return {
            'valid': True,
            'proof_type': proof.proof_type.value,
            'verified_claim': self._get_claim_description(proof),
            'metadata': proof.metadata
        }
    
    def _get_claim_description(self, proof: InnocenceProof) -> str:
        """Obtener descripción del claim verificado."""
        if proof.proof_type == InnocenceType.TEMPORAL:
            return f"Deposit before {proof.public_inputs.get('cutoff_timestamp')}"
        elif proof.proof_type == InnocenceType.NON_INCLUSION:
            return "Source not in blacklist"
        elif proof.proof_type == InnocenceType.INCLUSION:
            return f"Source in whitelist: {proof.metadata.get('whitelist_name')}"
        elif proof.proof_type == InnocenceType.CATEGORY:
            return f"Source category: {proof.metadata.get('category')}"
        return "Unknown claim"
    
    # ========== HELPERS ==========
    
    def _compute_commitment(
        self, 
        amount: int, 
        timestamp: datetime, 
        secret: bytes, 
        nullifier: bytes
    ) -> bytes:
        data = f"{amount}:{timestamp.isoformat()}:{secret.hex()}:{nullifier.hex()}"
        return sha256(data.encode()).digest()
    
    def _compute_nullifier_hash(self, nullifier: bytes, secret: bytes) -> bytes:
        return sha256(nullifier + secret).digest()
    
    def _compute_set_root(self, address_set: Set[str]) -> bytes:
        sorted_addresses = sorted(address_set)
        data = ''.join(sorted_addresses)
        return sha256(data.encode()).digest()
    
    def _update_merkle_root(self):
        if not self.deposits:
            self.merkle_root = b'\x00' * 32
            return
        
        commitments = [d.commitment for d in self.deposits]
        data = b''.join(commitments)
        self.merkle_root = sha256(data).digest()
    
    def _categorize_source(self, address: str) -> str:
        for category, addresses in self.source_categories.items():
            if address in addresses:
                return category
        return 'unknown'
    
    def _generate_zk_proof(
        self, 
        proof_type: str, 
        deposit: Deposit,
        secret: bytes,
        nullifier: bytes,
        public_inputs: Dict
    ) -> bytes:
        """Generar ZK proof (placeholder)."""
        data = f"{proof_type}:{deposit.commitment.hex()}:{json.dumps(public_inputs)}"
        return sha256(data.encode() + secret).digest()
    
    def _verify_zk_proof(self, proof: InnocenceProof) -> bool:
        """Verificar ZK proof (placeholder)."""
        return True


def poi_demo():
    print("=== Proof-of-Innocence Demo ===\n")
    
    system = ProofOfInnocenceSystem()
    
    # Configurar fuentes conocidas
    system.add_whitelisted_source("0xCoinbase", "regulated_exchange")
    system.add_whitelisted_source("0xKraken", "regulated_exchange")
    system.add_whitelisted_source("0xUniswap", "defi_protocol")
    
    # Simular depósitos
    print("--- Registering deposits ---")
    
    secret1 = b"user1_secret_key"
    nullifier1 = b"user1_nullifier"
    
    deposit1 = system.register_deposit(
        amount=1000000,
        source_address="0xCoinbase",
        secret=secret1,
        nullifier=nullifier1
    )
    print(f"Deposit 1: {deposit1.amount} from {deposit1.source_address}")
    print(f"  Timestamp: {deposit1.timestamp}")
    print(f"  Category: {deposit1.source_category}")
    
    # Simular evento de contaminación (hack)
    import time
    time.sleep(0.1)  # Pequeña pausa
    
    hack_event = {
        'type': 'hack',
        'timestamp': datetime.now(),
        'affected_addresses': ['0xHackerWallet', '0xStolenFunds'],
        'description': 'Exchange X hack - $100M stolen'
    }
    system.register_contamination_event(**hack_event)
    print(f"\n--- Contamination event registered ---")
    print(f"  Type: {hack_event['type']}")
    print(f"  Timestamp: {hack_event['timestamp']}")
    
    # Depósito después del hack (no puede probar innocence temporal)
    secret2 = b"user2_secret_key"
    nullifier2 = b"user2_nullifier"
    
    deposit2 = system.register_deposit(
        amount=500000,
        source_address="0xUnknown",
        secret=secret2,
        nullifier=nullifier2
    )
    print(f"\nDeposit 2 (after hack): from {deposit2.source_address}")
    
    # Generar proofs
    print("\n--- Generating Proofs ---")
    
    # Proof temporal para deposit1 (antes del hack)
    temporal_proof = system.generate_temporal_poi(
        deposit1, secret1, nullifier1, hack_event
    )
    if temporal_proof:
        print(f"\nTemporal PoI for Deposit 1:")
        print(f"  Type: {temporal_proof.proof_type.value}")
        print(f"  Claim: Deposit before hack")
        
        verification = system.verify_poi(temporal_proof)
        print(f"  Verification: {verification['valid']}")
    
    # Proof de inclusión en whitelist
    inclusion_proof = system.generate_inclusion_poi(
        deposit1, secret1, nullifier1, "regulated_exchanges"
    )
    if inclusion_proof:
        print(f"\nInclusion PoI for Deposit 1:")
        print(f"  Type: {inclusion_proof.proof_type.value}")
        print(f"  Claim: Source in whitelist")
    
    # Proof de categoría
    category_proof = system.generate_category_poi(
        deposit1, secret1, nullifier1, "regulated_exchange"
    )
    if category_proof:
        print(f"\nCategory PoI for Deposit 1:")
        print(f"  Type: {category_proof.proof_type.value}")
        print(f"  Category: {category_proof.metadata['category']}")
    
    # Intentar proof temporal para deposit2 (después del hack - debe fallar)
    temporal_proof2 = system.generate_temporal_poi(
        deposit2, secret2, nullifier2, hack_event
    )
    print(f"\nTemporal PoI for Deposit 2 (after hack): {temporal_proof2}")
    print("  Cannot prove innocence - deposit was after the hack!")

poi_demo()
\end{lstlisting}

%===========================================
\section{Selective De-anonymization}
\label{sec:selective-deanon}
%===========================================

\subsection{Concepto y Necesidad}

\begin{definition}[Selective De-anonymization]
Mecanismo criptográfico que permite revelar la identidad de un usuario bajo condiciones específicas y predefinidas, mientras mantiene la privacidad en circunstancias normales.

\textbf{Propiedades:}
\begin{enumerate}
    \item \textbf{Privacidad por defecto:} Usuario anónimo en operaciones normales
    \item \textbf{Condiciones claras:} De-anonymization solo bajo circunstancias específicas
    \item \textbf{Proceso verificable:} No puede abusarse sin dejar rastro
    \item \textbf{Proporcionalidad:} Solo se revela lo necesario
\end{enumerate}
\end{definition}

\begin{advertencia}
\textbf{Riesgos de de-anonymization mal implementada:}

\begin{enumerate}
    \item \textbf{Backdoors:} Si una entidad puede de-anonimizar unilateralmente, no hay privacidad real
    \item \textbf{Coerción:} Gobiernos pueden forzar uso de mecanismo
    \item \textbf{Filtraciones:} Claves de escrow pueden ser robadas
    \item \textbf{Abuso:} ``Condiciones especiales'' pueden expandirse
\end{enumerate}

\textbf{Mitigaciones:}
\begin{itemize}
    \item Threshold schemes: Requiere múltiples partes
    \item Timelock encryption: Transparencia temporal
    \item Auditoría pública de usos
    \item Límites técnicos en scope de revelación
\end{itemize}
\end{advertencia}

\subsection{Mecanismos de Selective De-anonymization}

\subsubsection{Threshold Decryption}

\begin{definition}[Threshold Decryption]
Sistema donde $k$ de $n$ autoridades deben cooperar para descifrar información:
\[
\text{Decrypt}(C) = \text{Combine}\left(\text{PartialDecrypt}_1(C), ..., \text{PartialDecrypt}_k(C)\right)
\]

\textbf{Configuraciones comunes:}
\begin{itemize}
    \item $3$ de $5$: Mayoría simple
    \item $5$ de $7$: Supermayoría
    \item $2$ de $3$: Mínimo viable
\end{itemize}
\end{definition}

\begin{figure}[htbp]
\centering
\begin{tikzpicture}[
    node distance=1cm,
    authority/.style={circle, draw, minimum size=1cm, fill=blue!20},
    selected/.style={circle, draw, minimum size=1cm, fill=green!30},
    data/.style={rectangle, draw, minimum width=2cm, minimum height=0.8cm, fill=yellow!20},
    arrow/.style={->, thick}
]
    % Autoridades
    \node[selected] (a1) at (0, 2) {$A_1$};
    \node[authority] (a2) at (2, 2) {$A_2$};
    \node[selected] (a3) at (4, 2) {$A_3$};
    \node[authority] (a4) at (6, 2) {$A_4$};
    \node[selected] (a5) at (8, 2) {$A_5$};
    
    % Datos encriptados
    \node[data] (enc) at (4, 0) {Identidad encriptada};
    
    % Partial decryptions
    \draw[arrow] (a1) -- node[left, font=\small] {$D_1$} (2, 0);
    \draw[arrow] (a3) -- node[left, font=\small] {$D_3$} (4, 0);
    \draw[arrow] (a5) -- node[right, font=\small] {$D_5$} (6, 0);
    
    % Resultado
    \node[data, fill=green!20] (dec) at (4, -2) {Identidad revelada};
    
    \draw[arrow] (enc) -- node[right, font=\small] {$k=3$ de $n=5$} (dec);
    
    % Etiqueta
    \node at (4, -3.5) [font=\small] {Solo con 3+ autoridades cooperando};
\end{tikzpicture}
\caption{Threshold decryption: 3 de 5 autoridades revelan la identidad}
\label{fig:threshold-decryption}
\end{figure}

\begin{lstlisting}[language=Python, caption={Threshold decryption con Shamir's Secret Sharing}]
from typing import List, Tuple
import secrets
from functools import reduce

class ThresholdDecryption:
    """
    Sistema de threshold decryption usando Shamir's Secret Sharing.
    """
    
    def __init__(self, threshold: int, total: int, prime: int = None):
        """
        Args:
            threshold: Número mínimo de shares para reconstruir (k)
            total: Número total de shares (n)
            prime: Primo para el campo finito
        """
        self.k = threshold
        self.n = total
        self.p = prime or self._generate_prime()
    
    def _generate_prime(self) -> int:
        """Generar primo seguro."""
        # En producción: usar primo criptográfico grande
        return 2**127 - 1  # Mersenne prime
    
    def _mod_inverse(self, a: int, p: int) -> int:
        """Inverso multiplicativo modular."""
        return pow(a, p - 2, p)
    
    def split_secret(self, secret: int) -> List[Tuple[int, int]]:
        """
        Dividir secreto en n shares, donde k son suficientes para reconstruir.
        
        Returns:
            Lista de (x, y) shares
        """
        if secret >= self.p:
            raise ValueError("Secret must be smaller than prime")
        
        # Generar coeficientes aleatorios para polinomio de grado k-1
        # f(x) = secret + a1*x + a2*x^2 + ... + a(k-1)*x^(k-1)
        coefficients = [secret] + [
            secrets.randbelow(self.p) 
            for _ in range(self.k - 1)
        ]
        
        # Evaluar polinomio en n puntos
        shares = []
        for x in range(1, self.n + 1):
            y = sum(
                coef * pow(x, i, self.p) 
                for i, coef in enumerate(coefficients)
            ) % self.p
            shares.append((x, y))
        
        return shares
    
    def reconstruct_secret(self, shares: List[Tuple[int, int]]) -> int:
        """
        Reconstruir secreto usando interpolación de Lagrange.
        
        Args:
            shares: Al menos k shares (x, y)
        
        Returns:
            Secreto original
        """
        if len(shares) < self.k:
            raise ValueError(f"Need at least {self.k} shares")
        
        # Usar solo k shares
        shares = shares[:self.k]
        
        secret = 0
        for i, (xi, yi) in enumerate(shares):
            # Calcular base de Lagrange
            numerator = 1
            denominator = 1
            
            for j, (xj, _) in enumerate(shares):
                if i != j:
                    numerator = (numerator * (-xj)) % self.p
                    denominator = (denominator * (xi - xj)) % self.p
            
            # Li(0) = numerator / denominator
            lagrange_basis = (numerator * self._mod_inverse(denominator, self.p)) % self.p
            
            # Sumar contribución
            secret = (secret + yi * lagrange_basis) % self.p
        
        return secret


class SelectiveDeAnonymization:
    """
    Sistema de de-anonimización selectiva para privacy pools.
    """
    
    def __init__(self, authorities: List[str], threshold: int):
        self.authorities = authorities
        self.threshold = threshold
        self.total = len(authorities)
        self.shamir = ThresholdDecryption(threshold, self.total)
        
        # Storage de datos encriptados
        self.encrypted_identities: dict = {}  # nullifier -> encrypted_data
        self.authority_shares: dict = {}  # nullifier -> {authority: share}
        
        # Audit log
        self.decryption_log: List[dict] = []
    
    def register_with_escrow(
        self,
        nullifier: bytes,
        identity_data: dict,
        encryption_key: int
    ) -> dict:
        """
        Registrar identidad en escrow al depositar.
        
        Args:
            nullifier: Nullifier del depósito
            identity_data: Datos de identidad a proteger
            encryption_key: Clave de encriptación (secreta)
        
        Returns:
            Confirmación con shares distribuidos
        """
        # Encriptar identidad (simplificado - usar AES en producción)
        identity_bytes = str(identity_data).encode()
        encrypted = self._encrypt(identity_bytes, encryption_key)
        
        # Dividir clave entre autoridades
        shares = self.shamir.split_secret(encryption_key)
        
        # Almacenar
        self.encrypted_identities[nullifier] = encrypted
        self.authority_shares[nullifier] = {
            self.authorities[i]: share 
            for i, share in enumerate(shares)
        }
        
        return {
            'nullifier': nullifier.hex(),
            'shares_distributed': len(shares),
            'threshold_required': self.threshold
        }
    
    def request_decryption(
        self,
        nullifier: bytes,
        reason: str,
        requesting_authority: str,
        court_order: Optional[bytes] = None
    ) -> dict:
        """
        Iniciar proceso de de-anonimización.
        """
        if nullifier not in self.encrypted_identities:
            return {'success': False, 'error': 'Unknown nullifier'}
        
        request_id = secrets.token_hex(16)
        
        # Registrar solicitud
        request = {
            'request_id': request_id,
            'nullifier': nullifier.hex(),
            'reason': reason,
            'requester': requesting_authority,
            'court_order': court_order.hex() if court_order else None,
            'timestamp': datetime.now().isoformat(),
            'status': 'pending',
            'approvals': []
        }
        
        self.decryption_log.append(request)
        
        return {
            'success': True,
            'request_id': request_id,
            'threshold_needed': self.threshold,
            'authorities_to_contact': self.authorities
        }
    
    def approve_decryption(
        self,
        request_id: str,
        authority: str,
        authority_signature: bytes
    ) -> dict:
        """
        Una autoridad aprueba la solicitud de decryption.
        """
        # Encontrar request
        request = next(
            (r for r in self.decryption_log if r['request_id'] == request_id),
            None
        )
        
        if not request:
            return {'success': False, 'error': 'Unknown request'}
        
        if authority not in self.authorities:
            return {'success': False, 'error': 'Not an authority'}
        
        if authority in [a['authority'] for a in request['approvals']]:
            return {'success': False, 'error': 'Already approved'}
        
        # Añadir aprobación
        nullifier = bytes.fromhex(request['nullifier'])
        share = self.authority_shares[nullifier][authority]
        
        request['approvals'].append({
            'authority': authority,
            'share': share,
            'signature': authority_signature.hex(),
            'timestamp': datetime.now().isoformat()
        })
        
        # Verificar si hay suficientes aprobaciones
        if len(request['approvals']) >= self.threshold:
            request['status'] = 'ready'
            return {
                'success': True,
                'message': 'Threshold reached, decryption ready'
            }
        
        return {
            'success': True,
            'approvals': len(request['approvals']),
            'needed': self.threshold
        }
    
    def execute_decryption(self, request_id: str) -> dict:
        """
        Ejecutar decryption si hay suficientes aprobaciones.
        """
        request = next(
            (r for r in self.decryption_log if r['request_id'] == request_id),
            None
        )
        
        if not request or request['status'] != 'ready':
            return {'success': False, 'error': 'Not ready for decryption'}
        
        # Recolectar shares
        shares = [a['share'] for a in request['approvals']]
        
        # Reconstruir clave
        encryption_key = self.shamir.reconstruct_secret(shares)
        
        # Descifrar
        nullifier = bytes.fromhex(request['nullifier'])
        encrypted = self.encrypted_identities[nullifier]
        decrypted = self._decrypt(encrypted, encryption_key)
        
        # Actualizar log
        request['status'] = 'completed'
        request['completed_at'] = datetime.now().isoformat()
        
        return {
            'success': True,
            'identity_data': decrypted.decode(),
            'request_id': request_id,
            'authorities_involved': [a['authority'] for a in request['approvals']]
        }
    
    def _encrypt(self, data: bytes, key: int) -> bytes:
        """Encriptar datos (simplificado)."""
        # En producción: usar AES-GCM
        key_bytes = key.to_bytes(32, 'big')
        return bytes(a ^ b for a, b in zip(data, key_bytes * (len(data) // 32 + 1)))
    
    def _decrypt(self, data: bytes, key: int) -> bytes:
        """Descifrar datos."""
        return self._encrypt(data, key)  # XOR es simétrico


def threshold_demo():
    print("=== Selective De-anonymization Demo ===\n")
    
    # Setup: 3 de 5 autoridades
    authorities = [
        "Court_Authority",
        "Financial_Regulator",
        "Privacy_Ombudsman",
        "Industry_Representative",
        "Civil_Liberties_Org"
    ]
    
    system = SelectiveDeAnonymization(authorities, threshold=3)
    print(f"System setup: {system.threshold} of {system.total} authorities required")
    print(f"Authorities: {authorities}\n")
    
    # Usuario deposita con escrow de identidad
    user_nullifier = secrets.token_bytes(32)
    user_identity = {
        'name': 'John Doe',
        'country': 'US',
        'kyc_provider': 'VerifyInc',
        'verified_at': '2024-01-15'
    }
    encryption_key = secrets.randbelow(2**128)
    
    escrow_result = system.register_with_escrow(
        user_nullifier,
        user_identity,
        encryption_key
    )
    print(f"Identity escrowed:")
    print(f"  Shares distributed: {escrow_result['shares_distributed']}")
    print(f"  Threshold: {escrow_result['threshold_required']}")
    
    # Simular solicitud de de-anonimización
    print("\n--- De-anonymization Request ---")
    
    request = system.request_decryption(
        user_nullifier,
        reason="Suspected money laundering - Case #12345",
        requesting_authority="Financial_Regulator",
        court_order=b"court_order_document_hash"
    )
    
    print(f"Request ID: {request['request_id']}")
    print(f"Threshold needed: {request['threshold_needed']}")
    
    # Autoridades aprueban
    print("\n--- Authority Approvals ---")
    
    approving_authorities = ["Court_Authority", "Financial_Regulator", "Privacy_Ombudsman"]
    
    for auth in approving_authorities:
        result = system.approve_decryption(
            request['request_id'],
            auth,
            secrets.token_bytes(64)  # Firma
        )
        print(f"{auth}: Approved ({result.get('approvals', 'ready')}/{system.threshold})")
    
    # Ejecutar decryption
    print("\n--- Executing Decryption ---")
    
    decryption_result = system.execute_decryption(request['request_id'])
    
    if decryption_result['success']:
        print(f"Identity revealed: {decryption_result['identity_data']}")
        print(f"Authorities involved: {decryption_result['authorities_involved']}")
    
    # Audit log
    print("\n--- Audit Log ---")
    for log in system.decryption_log:
        print(f"Request: {log['request_id'][:16]}...")
        print(f"  Status: {log['status']}")
        print(f"  Reason: {log['reason']}")
        print(f"  Approvals: {len(log['approvals'])}")

threshold_demo()
\end{lstlisting}

\subsubsection{Timelock Encryption}

\begin{definition}[Timelock Encryption]
Encriptación que se vuelve descifrable automáticamente después de cierto tiempo:
\begin{itemize}
    \item \textbf{Puzzle-based:} Requiere computación que toma tiempo T
    \item \textbf{Witness encryption:} Descifrable cuando evento ocurre
    \item \textbf{Blockchain-based:} Clave liberada en bloque futuro
\end{itemize}
\end{definition}

\begin{mathnote}[Timelock Puzzles (Rivest, Shamir, Wagner)]
\textbf{Construcción básica:}
\begin{enumerate}
    \item Elegir RSA modulus $n = pq$
    \item Calcular $\phi(n) = (p-1)(q-1)$
    \item Elegir tiempo $T$ (número de squarings)
    \item Calcular $e = 2^T \mod \phi(n)$ (rápido si conoces $\phi(n)$)
    \item Puzzle: $C = m \cdot a^e \mod n$
\end{enumerate}

\textbf{Resolver (sin conocer $\phi(n)$):}
\begin{itemize}
    \item Calcular $a^{2^T} \mod n$ requiere $T$ squarings secuenciales
    \item No paralelizable (inherentemente secuencial)
    \item Tiempo predecible: $T$ operaciones
\end{itemize}
\end{mathnote}

\subsubsection{Court-Order Escrow}

\begin{definition}[Court-Order Escrow]
Sistema donde la clave de descifrado está en custodia de entidad neutral, liberada solo por orden judicial válida:
\begin{itemize}
    \item Clave en HSM (Hardware Security Module)
    \item Proceso legal requerido para acceso
    \item Registro público de todas las órdenes
\end{itemize}
\end{definition}

\begin{bancozk}
\textbf{Selective De-anonymization en Banco ZK:}

Banco ZK implementa de-anonimización selectiva opcional:

\begin{enumerate}
    \item \textbf{Modo estándar (default):}
    \begin{itemize}
        \item Privacidad total
        \item Sin escrow de identidad
        \item Para usuarios que no lo necesitan
    \end{itemize}
    
    \item \textbf{Modo compliance (opt-in):}
    \begin{itemize}
        \item Usuario escoge depositar con escrow
        \item Identidad encriptada con threshold scheme
        \item Autoridades: 3 de 5 (2 reguladores, 2 industria, 1 civil)
        \item Beneficio: Acceso a pools de mayor liquidez
    \end{itemize}
    
    \item \textbf{Proceso de revelación:}
    \begin{itemize}
        \item Requiere orden judicial + 3 autoridades
        \item Audit log público (hash de cada solicitud)
        \item Solo se revela identidad, no historial de txs
    \end{itemize}
\end{enumerate}

\textbf{Trade-off:} Usuarios compliance-mode tienen más opciones pero menos privacidad potencial.
\end{bancozk}
%===========================================
\section{Viewing Keys y Audit Keys}
\label{sec:viewing-keys}
%===========================================

\subsection{Jerarquía de Claves de Privacidad}

\begin{definition}[Viewing Key]
Clave derivada que permite ver (descifrar) transacciones asociadas a una cuenta sin poder gastar los fondos:
\[
\text{ViewingKey} = \text{Derive}(\text{SpendingKey}, \text{``view''})
\]

\textbf{Capacidades:}
\begin{itemize}
    \item Ver balances entrantes
    \item Descifrar memos de transacciones
    \item Calcular balance total
    \item NO puede firmar transacciones
\end{itemize}
\end{definition}

\begin{definition}[Audit Key]
Extensión del viewing key con capacidades adicionales para auditoría:
\begin{itemize}
    \item Todo lo de viewing key
    \item Generar reportes verificables (proofs)
    \item Certificar rangos de tiempo
    \item Puede tener scope limitado (solo ciertas txs)
\end{itemize}
\end{definition}

\begin{figure}[htbp]
\centering
\begin{tikzpicture}[
    node distance=1cm,
    key/.style={rectangle, draw, rounded corners, minimum width=3cm, minimum height=0.8cm},
    arrow/.style={->, thick}
]
    % Jerarquía de claves
    \node[key, fill=red!30] (master) at (0, 4) {Master Seed};
    
    \node[key, fill=orange!30] (spend) at (-3, 2) {Spending Key};
    \node[key, fill=yellow!30] (view) at (0, 2) {Viewing Key};
    \node[key, fill=green!30] (audit) at (3, 2) {Audit Key};
    
    \node[key, fill=blue!20] (full) at (-3, 0) {Full Control};
    \node[key, fill=blue!20] (readonly) at (0, 0) {Read-Only};
    \node[key, fill=blue!20] (reports) at (3, 0) {Reports + Proofs};
    
    % Flechas de derivación
    \draw[arrow] (master) -- (spend);
    \draw[arrow] (master) -- (view);
    \draw[arrow] (master) -- (audit);
    
    \draw[arrow] (spend) -- (full);
    \draw[arrow] (view) -- (readonly);
    \draw[arrow] (audit) -- (reports);
    
    % Capacidades
    \node at (-3, -1) [font=\small, align=center] {Gastar\\Ver\\Probar};
    \node at (0, -1) [font=\small, align=center] {Ver\\(no gastar)};
    \node at (3, -1) [font=\small, align=center] {Ver + Certificar\\(scope limitado)};
    
    % Usuarios típicos
    \node at (-3, -2) [font=\footnotesize, blue] {Usuario};
    \node at (0, -2) [font=\footnotesize, blue] {Contador};
    \node at (3, -2) [font=\footnotesize, blue] {Auditor};
\end{tikzpicture}
\caption{Jerarquía de claves: spending, viewing, audit}
\label{fig:key-hierarchy}
\end{figure}

\subsection{Implementación de Viewing Keys}

\begin{mathnote}[Derivación de Viewing Keys (estilo Zcash)]
En Zcash Sapling:

\textbf{Spending Key ($sk$):}
\begin{itemize}
    \item 256 bits aleatorios
    \item Controla todo
\end{itemize}

\textbf{Proof Authorizing Key ($ak$) y Nullifier Key ($nk$):}
\[
ak = \text{SpendAuthSig.DerivePublic}(sk)
\]
\[
nk = \text{PRF}^{nk}_{sk}(\text{``nk''})
\]

\textbf{Full Viewing Key ($fvk$):}
\[
fvk = (ak, nk, ovk)
\]
donde $ovk$ = outgoing viewing key

\textbf{Incoming Viewing Key ($ivk$):}
\[
ivk = \text{CRH}^{ivk}(ak, nk)
\]

Con $ivk$ puedes:
\begin{itemize}
    \item Detectar pagos entrantes a tu dirección
    \item Descifrar memos de esos pagos
    \item Calcular tu balance
\end{itemize}

NO puedes:
\begin{itemize}
    \item Ver pagos salientes (necesitas $ovk$)
    \item Gastar fondos (necesitas $sk$)
    \item Generar nullifiers (necesitas $nk$ completo)
\end{itemize}
\end{mathnote}

\begin{lstlisting}[language=Python, caption={Sistema de Viewing Keys y Audit Keys}]
from dataclasses import dataclass
from datetime import datetime, date
from typing import List, Dict, Optional, Tuple
from hashlib import sha256, sha512
import hmac
import secrets
from enum import Enum

class KeyType(Enum):
    SPENDING = "spending"
    VIEWING = "viewing"
    INCOMING_VIEWING = "incoming_viewing"
    OUTGOING_VIEWING = "outgoing_viewing"
    AUDIT = "audit"

@dataclass
class Transaction:
    tx_id: bytes
    timestamp: datetime
    amount: int
    direction: str  # 'incoming' or 'outgoing'
    memo: bytes
    encrypted_memo: bytes
    sender_address: Optional[str]
    recipient_address: str
    
@dataclass
class AuditReport:
    period_start: date
    period_end: date
    total_incoming: int
    total_outgoing: int
    transaction_count: int
    proof: bytes  # ZK proof de los totales

class ViewingKeySystem:
    """
    Sistema de viewing keys con diferentes niveles de acceso.
    """
    
    def __init__(self, master_seed: bytes):
        self.master_seed = master_seed
        self.keys = self._derive_all_keys()
        
        # Transaction storage (encriptado)
        self.transactions: List[Transaction] = []
    
    def _derive_all_keys(self) -> Dict[KeyType, bytes]:
        """Derivar todas las claves desde master seed."""
        keys = {}
        
        # Spending key (máximo privilegio)
        keys[KeyType.SPENDING] = self._derive_key("spending")
        
        # Viewing key (full viewing)
        keys[KeyType.VIEWING] = self._derive_key("viewing")
        
        # Incoming viewing key (solo recibidos)
        keys[KeyType.INCOMING_VIEWING] = self._derive_key("incoming_view")
        
        # Outgoing viewing key (solo enviados)
        keys[KeyType.OUTGOING_VIEWING] = self._derive_key("outgoing_view")
        
        # Audit key (con capacidad de generar proofs)
        keys[KeyType.AUDIT] = self._derive_key("audit")
        
        return keys
    
    def _derive_key(self, purpose: str) -> bytes:
        """Derivar clave para propósito específico."""
        return hmac.new(
            self.master_seed,
            purpose.encode(),
            sha512
        ).digest()[:32]
    
    def get_key(self, key_type: KeyType) -> bytes:
        """Obtener clave de tipo específico."""
        return self.keys[key_type]
    
    def export_viewing_key(self, key_type: KeyType) -> Dict:
        """
        Exportar viewing key para compartir.
        """
        if key_type == KeyType.SPENDING:
            raise ValueError("Cannot export spending key as viewing key")
        
        key = self.keys[key_type]
        
        return {
            'type': key_type.value,
            'key': key.hex(),
            'capabilities': self._get_capabilities(key_type),
            'created_at': datetime.now().isoformat()
        }
    
    def _get_capabilities(self, key_type: KeyType) -> List[str]:
        """Obtener capacidades de cada tipo de clave."""
        capabilities = {
            KeyType.SPENDING: ['spend', 'view_all', 'generate_proofs', 'sign'],
            KeyType.VIEWING: ['view_incoming', 'view_outgoing', 'view_memos'],
            KeyType.INCOMING_VIEWING: ['view_incoming', 'view_incoming_memos'],
            KeyType.OUTGOING_VIEWING: ['view_outgoing', 'view_outgoing_memos'],
            KeyType.AUDIT: ['view_all', 'generate_proofs', 'generate_reports']
        }
        return capabilities.get(key_type, [])
    
    # ========== ENCRIPTACIÓN DE TRANSACCIONES ==========
    
    def encrypt_transaction(self, tx: Transaction) -> Transaction:
        """Encriptar memo de transacción."""
        # En producción: usar encryption asimétrica apropiada
        if tx.direction == 'incoming':
            key = self.keys[KeyType.INCOMING_VIEWING]
        else:
            key = self.keys[KeyType.OUTGOING_VIEWING]
        
        encrypted = self._encrypt(tx.memo, key)
        tx.encrypted_memo = encrypted
        return tx
    
    def record_transaction(
        self,
        amount: int,
        direction: str,
        memo: str,
        counterparty: str
    ):
        """Registrar nueva transacción."""
        tx = Transaction(
            tx_id=secrets.token_bytes(32),
            timestamp=datetime.now(),
            amount=amount,
            direction=direction,
            memo=memo.encode(),
            encrypted_memo=b'',
            sender_address=counterparty if direction == 'incoming' else 'self',
            recipient_address='self' if direction == 'incoming' else counterparty
        )
        
        tx = self.encrypt_transaction(tx)
        self.transactions.append(tx)
        return tx
    
    # ========== VIEWING CON DIFERENTES CLAVES ==========
    
    def view_transactions(
        self,
        viewing_key: bytes,
        key_type: KeyType
    ) -> List[Dict]:
        """
        Ver transacciones usando una viewing key.
        """
        visible_txs = []
        
        for tx in self.transactions:
            can_view = False
            decrypted_memo = None
            
            if key_type in [KeyType.VIEWING, KeyType.AUDIT, KeyType.SPENDING]:
                can_view = True
                # Puede descifrar ambas direcciones
                try:
                    if tx.direction == 'incoming':
                        decrypted_memo = self._decrypt(
                            tx.encrypted_memo,
                            self.keys[KeyType.INCOMING_VIEWING]
                        )
                    else:
                        decrypted_memo = self._decrypt(
                            tx.encrypted_memo,
                            self.keys[KeyType.OUTGOING_VIEWING]
                        )
                except:
                    decrypted_memo = b'[encrypted]'
                    
            elif key_type == KeyType.INCOMING_VIEWING:
                if tx.direction == 'incoming':
                    can_view = True
                    decrypted_memo = self._decrypt(tx.encrypted_memo, viewing_key)
                    
            elif key_type == KeyType.OUTGOING_VIEWING:
                if tx.direction == 'outgoing':
                    can_view = True
                    decrypted_memo = self._decrypt(tx.encrypted_memo, viewing_key)
            
            if can_view:
                visible_txs.append({
                    'tx_id': tx.tx_id.hex()[:16],
                    'timestamp': tx.timestamp.isoformat(),
                    'amount': tx.amount,
                    'direction': tx.direction,
                    'memo': decrypted_memo.decode() if decrypted_memo else '[hidden]',
                    'counterparty': tx.sender_address if tx.direction == 'incoming' else tx.recipient_address
                })
        
        return visible_txs
    
    def calculate_balance(self, viewing_key: bytes, key_type: KeyType) -> Dict:
        """Calcular balance usando viewing key."""
        txs = self.view_transactions(viewing_key, key_type)
        
        incoming = sum(tx['amount'] for tx in txs if tx['direction'] == 'incoming')
        outgoing = sum(tx['amount'] for tx in txs if tx['direction'] == 'outgoing')
        
        return {
            'total_incoming': incoming,
            'total_outgoing': outgoing,
            'balance': incoming - outgoing,
            'transaction_count': len(txs)
        }
    
    # ========== AUDIT KEYS Y REPORTES ==========
    
    def generate_audit_report(
        self,
        audit_key: bytes,
        period_start: date,
        period_end: date
    ) -> AuditReport:
        """
        Generar reporte de auditoría verificable.
        """
        # Verificar que es audit key válida
        if audit_key != self.keys[KeyType.AUDIT]:
            raise ValueError("Invalid audit key")
        
        # Filtrar transacciones del período
        period_txs = [
            tx for tx in self.transactions
            if period_start <= tx.timestamp.date() <= period_end
        ]
        
        total_incoming = sum(
            tx.amount for tx in period_txs 
            if tx.direction == 'incoming'
        )
        total_outgoing = sum(
            tx.amount for tx in period_txs 
            if tx.direction == 'outgoing'
        )
        
        # Generar proof ZK de los totales
        proof = self._generate_audit_proof(
            period_txs, 
            total_incoming, 
            total_outgoing,
            period_start,
            period_end
        )
        
        return AuditReport(
            period_start=period_start,
            period_end=period_end,
            total_incoming=total_incoming,
            total_outgoing=total_outgoing,
            transaction_count=len(period_txs),
            proof=proof
        )
    
    def verify_audit_report(self, report: AuditReport) -> bool:
        """Verificar reporte de auditoría."""
        # Verificar ZK proof
        # En producción: verificación SNARK real
        return len(report.proof) > 0
    
    # ========== HELPERS ==========
    
    def _encrypt(self, data: bytes, key: bytes) -> bytes:
        """Encriptar datos (simplificado)."""
        return bytes(a ^ b for a, b in zip(data, key * (len(data) // len(key) + 1)))
    
    def _decrypt(self, data: bytes, key: bytes) -> bytes:
        """Descifrar datos."""
        return self._encrypt(data, key)
    
    def _generate_audit_proof(
        self,
        txs: List[Transaction],
        total_in: int,
        total_out: int,
        start: date,
        end: date
    ) -> bytes:
        """Generar ZK proof para auditoría."""
        # En producción: SNARK que prueba:
        # 1. Los totales corresponden a transacciones reales
        # 2. Todas las txs están en el período especificado
        # 3. No se omitió ninguna tx del período
        
        data = f"{total_in}:{total_out}:{start}:{end}:{len(txs)}"
        return sha256(data.encode()).digest()


def viewing_keys_demo():
    print("=== Viewing Keys Demo ===\n")
    
    # Crear sistema con master seed
    master_seed = secrets.token_bytes(32)
    system = ViewingKeySystem(master_seed)
    
    print("Key hierarchy created:")
    for key_type in KeyType:
        caps = system._get_capabilities(key_type)
        print(f"  {key_type.value}: {caps}")
    
    # Registrar transacciones
    print("\n--- Recording transactions ---")
    
    system.record_transaction(1000000, 'incoming', 'Salary payment', 'employer_addr')
    system.record_transaction(50000, 'outgoing', 'Coffee purchase', 'coffee_shop')
    system.record_transaction(200000, 'incoming', 'Freelance work', 'client_addr')
    system.record_transaction(100000, 'outgoing', 'Rent payment', 'landlord_addr')
    
    print(f"Recorded {len(system.transactions)} transactions")
    
    # Compartir viewing key con contador
    print("\n--- Accountant's View (Incoming Only) ---")
    
    ivk = system.get_key(KeyType.INCOMING_VIEWING)
    accountant_view = system.view_transactions(ivk, KeyType.INCOMING_VIEWING)
    
    for tx in accountant_view:
        print(f"  {tx['direction']}: {tx['amount']} - {tx['memo']}")
    
    print(f"\nAccountant can see: {len(accountant_view)} transactions")
    print("Cannot see: outgoing transactions")
    
    # Full viewing key
    print("\n--- Full Viewing Key View ---")
    
    fvk = system.get_key(KeyType.VIEWING)
    full_view = system.view_transactions(fvk, KeyType.VIEWING)
    
    print(f"Full view: {len(full_view)} transactions")
    balance = system.calculate_balance(fvk, KeyType.VIEWING)
    print(f"Balance: {balance['balance']}")
    print(f"  Incoming: {balance['total_incoming']}")
    print(f"  Outgoing: {balance['total_outgoing']}")
    
    # Audit report
    print("\n--- Audit Report ---")
    
    audit_key = system.get_key(KeyType.AUDIT)
    report = system.generate_audit_report(
        audit_key,
        date(2024, 1, 1),
        date(2024, 12, 31)
    )
    
    print(f"Period: {report.period_start} to {report.period_end}")
    print(f"Total incoming: {report.total_incoming}")
    print(f"Total outgoing: {report.total_outgoing}")
    print(f"Transaction count: {report.transaction_count}")
    print(f"Proof valid: {system.verify_audit_report(report)}")

viewing_keys_demo()
\end{lstlisting}

\subsection{Casos de Uso}

\begin{table}[htbp]
\centering
\small
\begin{tabular}{lp{3cm}p{3cm}p{3.5cm}}
\toprule
\textbf{Usuario} & \textbf{Clave} & \textbf{Puede Ver} & \textbf{Caso de Uso} \\
\midrule
Propietario & Spending Key & Todo + gastar & Control total \\
Contador & Incoming Viewing & Solo recibidos & Calcular ingresos para impuestos \\
Auditor & Full Viewing + Proof & Todo + reportes & Auditoría financiera \\
Regulador & Audit Key (temporal) & Período específico & Investigación con orden \\
Exchange & Viewing Key (scope) & Depósitos del usuario & Verificar origen de fondos \\
\bottomrule
\end{tabular}
\caption{Casos de uso para diferentes tipos de viewing keys}
\end{table}

%===========================================
\section{Regulatory-Compliant Mixers}
\label{sec:compliant-mixers}
%===========================================

\subsection{El Desafío}

\begin{intuicion}
Los mixers tradicionales (Tornado Cash) fueron sancionados porque:
\begin{itemize}
    \item No había forma de distinguir usuarios legítimos de criminales
    \item No había mecanismo de exclusión de fondos ilícitos
    \item No había forma de cooperar con investigaciones legítimas
\end{itemize}

\textbf{Pregunta:} ¿Es posible un mixer que preserve privacidad Y sea compliance-friendly?

\textbf{Respuesta:} Sí, con diseño cuidadoso:
\begin{itemize}
    \item Proof-of-innocence obligatorio
    \item Exclusión de blacklisted addresses
    \item Viewing keys opcionales para auditoría
    \item Association sets ``limpios''
\end{itemize}
\end{intuicion}

\subsection{Arquitectura de Mixer Compliant}

\begin{figure}[htbp]
\centering
\begin{tikzpicture}[
    node distance=1cm,
    component/.style={rectangle, draw, minimum width=2.5cm, minimum height=1cm, align=center},
    filter/.style={diamond, draw, minimum width=1.5cm, minimum height=1.5cm, fill=yellow!30},
    arrow/.style={->, thick},
    redarrow/.style={->, thick, red}
]
    % Entrada
    \node[component, fill=blue!20] (deposit) at (0, 3) {Depósito};
    
    % Filtros
    \node[filter] (kyc) at (0, 1) {KYC?};
    \node[filter] (sanction) at (3, 1) {Sanctions\\Check};
    
    % Pool
    \node[component, fill=green!20] (pool) at (6, 1) {Compliance\\Pool};
    
    % Salida
    \node[component, fill=orange!20] (withdraw) at (9, 1) {Withdraw\\+ PoI};
    
    % Flujo
    \draw[arrow] (deposit) -- (kyc);
    \draw[arrow] (kyc) -- node[above, font=\small] {Pass} (sanction);
    \draw[arrow] (sanction) -- node[above, font=\small] {Clean} (pool);
    \draw[arrow] (pool) -- (withdraw);
    
    % Rechazos
    \draw[redarrow] (kyc) -- node[left, font=\small] {Fail} (-2, 1) -- (-2, -1) node[below] {Rejected};
    \draw[redarrow] (sanction) -- node[right, font=\small] {Flagged} (3, -1) node[below] {Rejected};
    
    % Viewing keys
    \node[component, fill=purple!20] (audit) at (6, -1.5) {Audit\\Access};
    \draw[dashed, ->] (pool) -- (audit);
    
    % Etiqueta
    \node at (4.5, 4) {\textbf{Regulatory-Compliant Mixer}};
\end{tikzpicture}
\caption{Arquitectura de mixer con compliance integrado}
\label{fig:compliant-mixer}
\end{figure}

\begin{lstlisting}[language=Solidity, caption={Regulatory-Compliant Mixer Contract}]
// SPDX-License-Identifier: MIT
pragma solidity ^0.8.20;

import "@openzeppelin/contracts/token/ERC20/IERC20.sol";
import "@openzeppelin/contracts/access/AccessControl.sol";
import "@openzeppelin/contracts/security/ReentrancyGuard.sol";

interface IVerifier {
    function verifyProof(
        bytes calldata proof,
        uint256[] calldata publicInputs
    ) external view returns (bool);
}

interface ISanctionsOracle {
    function isSanctioned(address addr) external view returns (bool);
}

interface IKYCRegistry {
    function isVerified(bytes32 identityCommitment) external view returns (bool);
}

contract CompliantMixer is AccessControl, ReentrancyGuard {
    bytes32 public constant COMPLIANCE_ROLE = keccak256("COMPLIANCE_ROLE");
    bytes32 public constant AUDITOR_ROLE = keccak256("AUDITOR_ROLE");
    
    // Configuración
    IERC20 public immutable token;
    uint256 public immutable denomination;
    IVerifier public withdrawVerifier;
    IVerifier public innocenceVerifier;
    ISanctionsOracle public sanctionsOracle;
    IKYCRegistry public kycRegistry;
    
    // Estado del pool
    bytes32 public merkleRoot;
    uint256 public depositCount;
    mapping(bytes32 => bool) public nullifiers;
    mapping(bytes32 => bool) public commitments;
    
    // Compliance data
    mapping(bytes32 => DepositMetadata) public depositMetadata;
    mapping(bytes32 => bool) public blacklistedCommitments;
    
    // Association sets (Privacy Pools)
    mapping(bytes32 => bool) public associationSets;
    bytes32[] public approvedAssociationSets;
    
    struct DepositMetadata {
        uint256 timestamp;
        bytes32 kycCommitment;  // Commitment de KYC (no revela identidad)
        bool hasEscrow;         // Si tiene escrow de identidad
    }
    
    event Deposit(
        bytes32 indexed commitment,
        uint256 leafIndex,
        uint256 timestamp
    );
    
    event Withdrawal(
        bytes32 indexed nullifier,
        address indexed recipient,
        bytes32 associationSetRoot
    );
    
    event ComplianceAlert(
        bytes32 indexed commitment,
        string reason
    );
    
    constructor(
        address _token,
        uint256 _denomination,
        address _sanctionsOracle,
        address _kycRegistry
    ) {
        token = IERC20(_token);
        denomination = _denomination;
        sanctionsOracle = ISanctionsOracle(_sanctionsOracle);
        kycRegistry = IKYCRegistry(_kycRegistry);
        
        _grantRole(DEFAULT_ADMIN_ROLE, msg.sender);
        _grantRole(COMPLIANCE_ROLE, msg.sender);
    }
    
    // ========== DEPÓSITO CON COMPLIANCE ==========
    
    function deposit(
        bytes32 commitment,
        bytes32 kycCommitment,
        bool useEscrow,
        bytes calldata escrowData
    ) external nonReentrant {
        require(!commitments[commitment], "Commitment exists");
        require(commitment != bytes32(0), "Invalid commitment");
        
        // 1. Verificar sanctions
        require(
            !sanctionsOracle.isSanctioned(msg.sender),
            "Sanctioned address"
        );
        
        // 2. Verificar KYC (opcional pero recomendado)
        bool hasKYC = kycRegistry.isVerified(kycCommitment);
        
        // 3. Transferir fondos
        require(
            token.transferFrom(msg.sender, address(this), denomination),
            "Transfer failed"
        );
        
        // 4. Guardar metadata de compliance
        depositMetadata[commitment] = DepositMetadata({
            timestamp: block.timestamp,
            kycCommitment: hasKYC ? kycCommitment : bytes32(0),
            hasEscrow: useEscrow
        });
        
        // 5. Si usa escrow, almacenar datos encriptados
        if (useEscrow && escrowData.length > 0) {
            _storeEscrowData(commitment, escrowData);
        }
        
        // 6. Añadir al árbol
        commitments[commitment] = true;
        _insertCommitment(commitment);
        
        emit Deposit(commitment, depositCount++, block.timestamp);
    }
    
    // ========== RETIRO CON PROOF-OF-INNOCENCE ==========
    
    function withdraw(
        bytes calldata withdrawProof,
        bytes calldata innocenceProof,
        bytes32 nullifier,
        bytes32 root,
        address recipient,
        bytes32 associationSetRoot
    ) external nonReentrant {
        require(!nullifiers[nullifier], "Already spent");
        require(_isKnownRoot(root), "Unknown root");
        require(
            _isApprovedAssociationSet(associationSetRoot),
            "Association set not approved"
        );
        
        // 1. Verificar recipient no sancionado
        require(
            !sanctionsOracle.isSanctioned(recipient),
            "Sanctioned recipient"
        );
        
        // 2. Verificar withdraw proof
        uint256[] memory withdrawInputs = new uint256[](3);
        withdrawInputs[0] = uint256(root);
        withdrawInputs[1] = uint256(nullifier);
        withdrawInputs[2] = uint256(uint160(recipient));
        
        require(
            withdrawVerifier.verifyProof(withdrawProof, withdrawInputs),
            "Invalid withdraw proof"
        );
        
        // 3. Verificar proof-of-innocence
        uint256[] memory innocenceInputs = new uint256[](2);
        innocenceInputs[0] = uint256(root);
        innocenceInputs[1] = uint256(associationSetRoot);
        
        require(
            innocenceVerifier.verifyProof(innocenceProof, innocenceInputs),
            "Invalid innocence proof"
        );
        
        // 4. Marcar nullifier
        nullifiers[nullifier] = true;
        
        // 5. Transferir fondos
        require(
            token.transfer(recipient, denomination),
            "Transfer failed"
        );
        
        emit Withdrawal(nullifier, recipient, associationSetRoot);
    }
    
    // ========== GESTIÓN DE ASSOCIATION SETS ==========
    
    function approveAssociationSet(
        bytes32 setRoot,
        string calldata description
    ) external onlyRole(COMPLIANCE_ROLE) {
        associationSets[setRoot] = true;
        approvedAssociationSets.push(setRoot);
    }
    
    function revokeAssociationSet(bytes32 setRoot) 
        external onlyRole(COMPLIANCE_ROLE) 
    {
        associationSets[setRoot] = false;
    }
    
    function _isApprovedAssociationSet(bytes32 setRoot) internal view returns (bool) {
        return associationSets[setRoot];
    }
    
    // ========== BLACKLISTING ==========
    
    function blacklistCommitment(
        bytes32 commitment,
        string calldata reason
    ) external onlyRole(COMPLIANCE_ROLE) {
        blacklistedCommitments[commitment] = true;
        emit ComplianceAlert(commitment, reason);
    }
    
    // ========== AUDIT ACCESS ==========
    
    function getDepositMetadata(bytes32 commitment) 
        external view 
        onlyRole(AUDITOR_ROLE)
        returns (DepositMetadata memory) 
    {
        return depositMetadata[commitment];
    }
    
    function requestEscrowDecryption(
        bytes32 commitment,
        bytes calldata courtOrder
    ) external onlyRole(COMPLIANCE_ROLE) {
        // Iniciar proceso de threshold decryption
        // (implementación externa)
    }
    
    // ========== HELPERS ==========
    
    function _insertCommitment(bytes32 commitment) internal {
        // Actualizar Merkle tree
    }
    
    function _isKnownRoot(bytes32 root) internal view returns (bool) {
        return root == merkleRoot;  // Simplificado
    }
    
    function _storeEscrowData(bytes32 commitment, bytes calldata data) internal {
        // Almacenar datos de escrow encriptados
    }
}
\end{lstlisting}

%===========================================
\section{Travel Rule y Soluciones ZK}
\label{sec:travel-rule}
%===========================================

\subsection{Requisitos de la Travel Rule}

\begin{definition}[Travel Rule (FATF Recommendation 16)]
Requisito regulatorio que obliga a VASPs (Virtual Asset Service Providers) a compartir información del originador y beneficiario para transferencias sobre cierto umbral:

\textbf{Información requerida del originador:}
\begin{itemize}
    \item Nombre
    \item Número de cuenta (dirección)
    \item Dirección física, ID nacional, o fecha/lugar de nacimiento
\end{itemize}

\textbf{Información del beneficiario:}
\begin{itemize}
    \item Nombre
    \item Número de cuenta
\end{itemize}

\textbf{Umbral:} \$1,000 USD (varía por jurisdicción)
\end{definition}

\begin{advertencia}
\textbf{Problemas de la Travel Rule tradicional:}

\begin{enumerate}
    \item \textbf{Privacidad destruida:}
    \begin{itemize}
        \item Datos personales compartidos entre VASPs
        \item Múltiples copias de información sensible
        \item Riesgo de data breaches
    \end{itemize}
    
    \item \textbf{Incompatible con DeFi:}
    \begin{itemize}
        \item Smart contracts no tienen ``nombre''
        \item Self-custody no tiene VASP intermediario
        \item ¿Quién comparte info en swap descentralizado?
    \end{itemize}
    
    \item \textbf{Honeypot de datos:}
    \begin{itemize}
        \item Bases de datos centralizadas de todas las transacciones
        \item Target para hackers y vigilancia
    \end{itemize}
\end{enumerate}
\end{advertencia}

\subsection{Soluciones ZK para Travel Rule}

\begin{keyidea}
\textbf{ZK Travel Rule Compliance:}

En lugar de compartir datos personales directamente:

\begin{enumerate}
    \item \textbf{Proof of KYC validity:}
    \begin{itemize}
        \item Probar que originador/beneficiario tienen KYC válido
        \item Sin revelar los datos de KYC
        \item VASP receptor verifica proof
    \end{itemize}
    
    \item \textbf{Encrypted data + conditional access:}
    \begin{itemize}
        \item Datos encriptados en tránsito
        \item Solo descifrable por regulador con orden
        \item VASPs solo ven proofs
    \end{itemize}
    
    \item \textbf{Notary with ZK verification:}
    \begin{itemize}
        \item Tercero trusted verifica compliance
        \item Emite attestation sin ver datos
        \item Usando hardware seguro (SGX, etc.)
    \end{itemize}
\end{enumerate}
\end{keyidea}

\begin{lstlisting}[language=Python, caption={ZK Travel Rule Implementation}]
from dataclasses import dataclass
from datetime import datetime
from typing import Optional, Dict
from hashlib import sha256
import secrets

@dataclass
class TravelRuleData:
    """Datos requeridos por Travel Rule."""
    originator_name: str
    originator_account: str
    originator_address: str
    originator_id: str
    beneficiary_name: str
    beneficiary_account: str
    amount: int
    currency: str

@dataclass
class ZKTravelRuleProof:
    """Proof de compliance sin revelar datos."""
    originator_kyc_valid: bool
    beneficiary_kyc_valid: bool
    amount_above_threshold: bool
    sanctions_check_passed: bool
    proof: bytes
    encrypted_data: bytes  # Datos encriptados para regulador

class ZKTravelRuleSystem:
    """
    Sistema de Travel Rule con privacidad usando ZK.
    """
    
    def __init__(self, threshold_amount: int = 1000):
        self.threshold = threshold_amount
        
        # KYC registry (commitment -> verified)
        self.kyc_commitments: Dict[bytes, bool] = {}
        
        # Sanctions list
        self.sanctioned_accounts: set = set()
        
        # Regulator public key for encryption
        self.regulator_pubkey: bytes = b''
    
    def register_kyc(self, identity_data: Dict) -> bytes:
        """Registrar KYC y obtener commitment."""
        # Calcular commitment de identidad
        data_str = str(sorted(identity_data.items()))
        commitment = sha256(data_str.encode()).digest()
        
        self.kyc_commitments[commitment] = True
        return commitment
    
    def create_zk_travel_rule_proof(
        self,
        travel_data: TravelRuleData,
        originator_kyc_commitment: bytes,
        beneficiary_kyc_commitment: bytes,
        originator_kyc_proof: bytes,  # ZK proof de validez
        beneficiary_kyc_proof: bytes
    ) -> ZKTravelRuleProof:
        """
        Crear proof de Travel Rule compliance sin revelar datos.
        """
        # 1. Verificar KYC válido para ambas partes
        originator_valid = self._verify_kyc_proof(
            originator_kyc_commitment, 
            originator_kyc_proof
        )
        beneficiary_valid = self._verify_kyc_proof(
            beneficiary_kyc_commitment,
            beneficiary_kyc_proof
        )
        
        # 2. Verificar sanctions
        sanctions_ok = (
            travel_data.originator_account not in self.sanctioned_accounts and
            travel_data.beneficiary_account not in self.sanctioned_accounts
        )
        
        # 3. Verificar monto sobre threshold
        above_threshold = travel_data.amount >= self.threshold
        
        # 4. Encriptar datos para regulador
        encrypted = self._encrypt_for_regulator(travel_data)
        
        # 5. Generar ZK proof
        proof = self._generate_zk_proof(
            originator_valid,
            beneficiary_valid,
            sanctions_ok,
            above_threshold
        )
        
        return ZKTravelRuleProof(
            originator_kyc_valid=originator_valid,
            beneficiary_kyc_valid=beneficiary_valid,
            amount_above_threshold=above_threshold,
            sanctions_check_passed=sanctions_ok,
            proof=proof,
            encrypted_data=encrypted
        )
    
    def verify_travel_rule_proof(
        self,
        proof: ZKTravelRuleProof
    ) -> Dict:
        """
        Verificar proof de Travel Rule (por VASP receptor).
        """
        # Verificar ZK proof
        zk_valid = self._verify_zk_proof(proof.proof)
        
        if not zk_valid:
            return {
                'compliant': False,
                'error': 'Invalid ZK proof'
            }
        
        # Verificar flags de compliance
        all_checks = (
            proof.originator_kyc_valid and
            proof.beneficiary_kyc_valid and
            proof.sanctions_check_passed
        )
        
        return {
            'compliant': all_checks,
            'originator_kyc': proof.originator_kyc_valid,
            'beneficiary_kyc': proof.beneficiary_kyc_valid,
            'sanctions_ok': proof.sanctions_check_passed,
            'threshold_transfer': proof.amount_above_threshold,
            'data_available': len(proof.encrypted_data) > 0
        }
    
    def regulator_decrypt(
        self,
        encrypted_data: bytes,
        regulator_privkey: bytes,
        court_order: bytes
    ) -> Optional[TravelRuleData]:
        """
        Regulador descifra datos con orden judicial.
        """
        # Verificar orden judicial (en producción: verificación real)
        if not self._verify_court_order(court_order):
            return None
        
        # Descifrar
        decrypted = self._decrypt_with_key(encrypted_data, regulator_privkey)
        
        # Parsear datos
        # ... (deserialización)
        
        return decrypted
    
    # ========== HELPERS ==========
    
    def _verify_kyc_proof(self, commitment: bytes, proof: bytes) -> bool:
        """Verificar ZK proof de KYC válido."""
        return commitment in self.kyc_commitments and self.kyc_commitments[commitment]
    
    def _encrypt_for_regulator(self, data: TravelRuleData) -> bytes:
        """Encriptar datos para acceso regulatorio."""
        # En producción: encryption asimétrica con regulator pubkey
        plaintext = str(data).encode()
        return sha256(plaintext + self.regulator_pubkey).digest() + plaintext
    
    def _decrypt_with_key(self, encrypted: bytes, key: bytes) -> TravelRuleData:
        """Descifrar con clave privada."""
        # Placeholder
        return None
    
    def _generate_zk_proof(self, *args) -> bytes:
        """Generar ZK proof."""
        data = ':'.join(str(a) for a in args)
        return sha256(data.encode()).digest()
    
    def _verify_zk_proof(self, proof: bytes) -> bool:
        """Verificar ZK proof."""
        return len(proof) == 32
    
    def _verify_court_order(self, order: bytes) -> bool:
        """Verificar orden judicial."""
        return len(order) > 0


def travel_rule_demo():
    print("=== ZK Travel Rule Demo ===\n")
    
    system = ZKTravelRuleSystem(threshold_amount=1000)
    
    # Registrar KYC de Alice y Bob
    alice_kyc = system.register_kyc({
        'name': 'Alice Smith',
        'country': 'US',
        'id': 'ABC123'
    })
    
    bob_kyc = system.register_kyc({
        'name': 'Bob Jones',
        'country': 'UK',
        'id': 'XYZ789'
    })
    
    print("KYC registered for Alice and Bob")
    
    # Alice envía $5000 a Bob
    travel_data = TravelRuleData(
        originator_name="Alice Smith",
        originator_account="0xAlice...",
        originator_address="123 Main St",
        originator_id="ABC123",
        beneficiary_name="Bob Jones",
        beneficiary_account="0xBob...",
        amount=5000,
        currency="USDC"
    )
    
    print(f"\nTransfer: ${travel_data.amount} from Alice to Bob")
    
    # Crear ZK proof
    zk_proof = system.create_zk_travel_rule_proof(
        travel_data,
        alice_kyc,
        bob_kyc,
        b"alice_kyc_proof",
        b"bob_kyc_proof"
    )
    
    print("\n--- ZK Travel Rule Proof ---")
    print(f"Originator KYC valid: {zk_proof.originator_kyc_valid}")
    print(f"Beneficiary KYC valid: {zk_proof.beneficiary_kyc_valid}")
    print(f"Sanctions check: {zk_proof.sanctions_check_passed}")
    print(f"Above threshold: {zk_proof.amount_above_threshold}")
    print(f"Encrypted data available: {len(zk_proof.encrypted_data)} bytes")
    
    # VASP receptor verifica
    print("\n--- Receiving VASP Verification ---")
    verification = system.verify_travel_rule_proof(zk_proof)
    
    print(f"Compliant: {verification['compliant']}")
    print("\nWhat receiving VASP knows:")
    print("  ✓ Originator has valid KYC")
    print("  ✓ Beneficiary has valid KYC")
    print("  ✓ Neither party is sanctioned")
    print("  ✗ Originator's actual name")
    print("  ✗ Originator's address")
    print("  ✗ Originator's ID number")
    print("\n  → Privacy preserved while compliant!")

travel_rule_demo()
\end{lstlisting}
%===========================================
\section{Chainalysis y Contramedidas}
\label{sec:chainalysis}
%===========================================

\subsection{Técnicas de Blockchain Analytics}

\begin{definition}[Blockchain Analytics / Chainalysis]
Conjunto de técnicas para analizar blockchains públicas con el objetivo de:
\begin{itemize}
    \item Identificar propietarios de direcciones
    \item Rastrear flujo de fondos
    \item Detectar actividad sospechosa
    \item Vincular transacciones con identidades reales
\end{itemize}

\textbf{Principales empresas:} Chainalysis, Elliptic, CipherTrace, TRM Labs
\end{definition}

\begin{mathnote}[Técnicas de Análisis]
\textbf{1. Heurísticas de Clustering:}
\begin{itemize}
    \item \textbf{Common Input Ownership:} Inputs de misma tx = mismo dueño
    \item \textbf{Change Detection:} Identificar output de cambio
    \item \textbf{Address Reuse:} Vincular múltiples txs de misma dirección
    \item \textbf{Script Fingerprinting:} Identificar tipo de wallet por script
\end{itemize}

\textbf{2. Taint Analysis:}
\begin{itemize}
    \item Rastrear ``contaminación'' desde fuentes conocidas
    \item Propagar taint a través de transacciones
    \item Métodos: FIFO, LIFO, proporcional, haircut
\end{itemize}

\textbf{3. Exchange Deposits:}
\begin{itemize}
    \item Exchanges conocen identidad de usuarios (KYC)
    \item Depósito a exchange = posible identificación
    \item Colaboración exchange-analytics
\end{itemize}

\textbf{4. Machine Learning:}
\begin{itemize}
    \item Clasificación de direcciones (exchange, mixer, darknet)
    \item Detección de patrones anómalos
    \item Predicción de comportamiento
\end{itemize}

\textbf{5. Off-chain Intelligence:}
\begin{itemize}
    \item Scraped addresses de foros, redes sociales
    \item Direcciones de donación públicas
    \item Investigación tradicional
\end{itemize}
\end{mathnote}

\begin{figure}[htbp]
\centering
\begin{tikzpicture}[
    node distance=0.8cm,
    addr/.style={circle, draw, minimum size=0.6cm, font=\tiny},
    known/.style={circle, draw, minimum size=0.6cm, fill=red!30, font=\tiny},
    cluster/.style={ellipse, draw, dashed, minimum width=2cm, minimum height=1.5cm},
    arrow/.style={->, thick}
]
    % Cluster 1 (identificado)
    \node[cluster, fill=red!10] (c1) at (0, 2) {};
    \node at (0, 3) [font=\small] {Cluster: Exchange X};
    \node[known] (a1) at (-0.5, 2.2) {A1};
    \node[known] (a2) at (0.5, 2.2) {A2};
    \node[known] (a3) at (0, 1.8) {A3};
    
    % Cluster 2 (desconocido pero vinculado)
    \node[cluster, fill=yellow!10] (c2) at (4, 2) {};
    \node at (4, 3) [font=\small] {Cluster: Usuario ?};
    \node[addr] (b1) at (3.5, 2.2) {B1};
    \node[addr] (b2) at (4.5, 2.2) {B2};
    \node[addr] (b3) at (4, 1.8) {B3};
    
    % Transacción vinculante
    \draw[arrow, red] (a2) -- node[above, font=\tiny] {withdraw} (b1);
    
    % Mixer
    \node[draw, rectangle, fill=gray!20] (mixer) at (8, 2) {Mixer};
    
    \draw[arrow] (b2) -- (mixer);
    
    % Outputs del mixer
    \node[addr] (m1) at (10, 3) {M1};
    \node[addr] (m2) at (10, 2) {M2};
    \node[addr] (m3) at (10, 1) {M3};
    
    \draw[arrow, dashed] (mixer) -- (m1);
    \draw[arrow, dashed] (mixer) -- (m2);
    \draw[arrow, dashed] (mixer) -- (m3);
    
    % Labels
    \node at (0, 0.5) [font=\small, red] {Conocido};
    \node at (4, 0.5) [font=\small, orange] {Vinculado};
    \node at (10, 0.5) [font=\small, green] {Mezclado};
    
    % Análisis
    \node at (2, -0.5) [font=\footnotesize, align=center] {
        Chainalysis puede vincular\\
        B1 con Exchange X
    };
\end{tikzpicture}
\caption{Técnicas de clustering y taint analysis}
\label{fig:chainalysis-techniques}
\end{figure}

\begin{lstlisting}[language=Python, caption={Simulación de técnicas de Chainalysis}]
from dataclasses import dataclass
from typing import Dict, List, Set, Tuple
from collections import defaultdict
import networkx as nx

@dataclass
class Transaction:
    tx_id: str
    inputs: List[Tuple[str, int]]   # (address, amount)
    outputs: List[Tuple[str, int]]  # (address, amount)
    timestamp: int

class BlockchainAnalytics:
    """
    Simulación de técnicas de blockchain analytics.
    """
    
    def __init__(self):
        # Grafo de transacciones
        self.graph = nx.DiGraph()
        
        # Clusters de direcciones
        self.clusters: Dict[str, Set[str]] = {}
        self.address_to_cluster: Dict[str, str] = {}
        
        # Entidades conocidas
        self.known_entities: Dict[str, str] = {}  # address -> entity_name
        
        # Taint tracking
        self.taint: Dict[str, Dict[str, float]] = defaultdict(dict)
        
        # Transacciones
        self.transactions: List[Transaction] = []
    
    def add_transaction(self, tx: Transaction):
        """Añadir transacción al análisis."""
        self.transactions.append(tx)
        
        # Añadir al grafo
        for in_addr, _ in tx.inputs:
            for out_addr, _ in tx.outputs:
                self.graph.add_edge(in_addr, out_addr, tx_id=tx.tx_id)
    
    def label_address(self, address: str, entity: str):
        """Etiquetar dirección como perteneciente a entidad."""
        self.known_entities[address] = entity
    
    # ========== CLUSTERING ==========
    
    def apply_common_input_heuristic(self):
        """
        Heurística 1: Inputs de misma tx = mismo dueño.
        """
        for tx in self.transactions:
            if len(tx.inputs) > 1:
                # Todas las direcciones de input van al mismo cluster
                input_addrs = [addr for addr, _ in tx.inputs]
                self._merge_addresses(input_addrs)
    
    def apply_change_heuristic(self):
        """
        Heurística 2: Detectar y vincular change addresses.
        """
        for tx in self.transactions:
            if len(tx.outputs) == 2:
                # Analizar cuál es change
                out1, out2 = tx.outputs
                
                # Heurística: output con monto no redondo es change
                is_round_1 = self._is_round_amount(out1[1])
                is_round_2 = self._is_round_amount(out2[1])
                
                change_addr = None
                if is_round_1 and not is_round_2:
                    change_addr = out2[0]
                elif is_round_2 and not is_round_1:
                    change_addr = out1[0]
                
                if change_addr:
                    # Vincular change con inputs
                    input_addrs = [addr for addr, _ in tx.inputs]
                    self._merge_addresses(input_addrs + [change_addr])
    
    def _merge_addresses(self, addresses: List[str]):
        """Fusionar direcciones en un cluster."""
        # Encontrar clusters existentes
        existing_clusters = set()
        for addr in addresses:
            if addr in self.address_to_cluster:
                existing_clusters.add(self.address_to_cluster[addr])
        
        # Crear nuevo cluster o usar existente
        if existing_clusters:
            main_cluster = list(existing_clusters)[0]
            # Fusionar otros clusters
            for cluster_id in existing_clusters:
                if cluster_id != main_cluster and cluster_id in self.clusters:
                    self.clusters[main_cluster].update(self.clusters[cluster_id])
                    del self.clusters[cluster_id]
        else:
            main_cluster = f"cluster_{len(self.clusters)}"
            self.clusters[main_cluster] = set()
        
        # Añadir todas las direcciones
        for addr in addresses:
            self.clusters[main_cluster].add(addr)
            self.address_to_cluster[addr] = main_cluster
    
    def _is_round_amount(self, amount: int) -> bool:
        """Verificar si monto es 'redondo'."""
        # Redondo si es múltiplo de 1000 (simplificado)
        return amount % 1000 == 0
    
    # ========== TAINT ANALYSIS ==========
    
    def mark_tainted(self, address: str, source: str):
        """Marcar dirección como origen de taint."""
        self.taint[address] = {source: 1.0}
    
    def propagate_taint(self):
        """Propagar taint a través de transacciones."""
        for tx in sorted(self.transactions, key=lambda x: x.timestamp):
            # Calcular taint total de inputs
            total_input = sum(amount for _, amount in tx.inputs)
            combined_taint: Dict[str, float] = defaultdict(float)
            
            for addr, amount in tx.inputs:
                if addr in self.taint:
                    weight = amount / total_input
                    for source, prop in self.taint[addr].items():
                        combined_taint[source] += prop * weight
            
            if not combined_taint:
                continue
            
            # Distribuir a outputs
            total_output = sum(amount for _, amount in tx.outputs)
            for addr, amount in tx.outputs:
                out_weight = amount / total_output
                self.taint[addr] = {
                    source: prop * out_weight
                    for source, prop in combined_taint.items()
                }
    
    def get_taint(self, address: str) -> Dict[str, float]:
        """Obtener taint de una dirección."""
        return dict(self.taint.get(address, {}))
    
    # ========== ANÁLISIS ==========
    
    def identify_entity(self, address: str) -> str:
        """Intentar identificar entidad propietaria de dirección."""
        # Método 1: Dirección directamente conocida
        if address in self.known_entities:
            return self.known_entities[address]
        
        # Método 2: Otra dirección del cluster es conocida
        if address in self.address_to_cluster:
            cluster_id = self.address_to_cluster[address]
            for addr in self.clusters.get(cluster_id, set()):
                if addr in self.known_entities:
                    return f"{self.known_entities[addr]} (cluster)"
        
        return "Unknown"
    
    def trace_funds(self, address: str, depth: int = 5) -> List[str]:
        """Rastrear flujo de fondos desde una dirección."""
        path = []
        current = address
        
        for _ in range(depth):
            path.append(current)
            
            # Encontrar siguiente dirección
            successors = list(self.graph.successors(current))
            if not successors:
                break
            
            # Elegir la de mayor monto (simplificado)
            current = successors[0]
        
        return path
    
    def get_cluster_stats(self) -> Dict:
        """Obtener estadísticas de clustering."""
        total_addresses = sum(len(c) for c in self.clusters.values())
        identified = sum(
            1 for addr in self.known_entities 
            if addr in self.address_to_cluster
        )
        
        return {
            'num_clusters': len(self.clusters),
            'total_addresses_clustered': total_addresses,
            'clusters_with_known_entity': identified,
            'largest_cluster_size': max(
                (len(c) for c in self.clusters.values()), 
                default=0
            )
        }


def chainalysis_demo():
    print("=== Chainalysis Simulation Demo ===\n")
    
    analytics = BlockchainAnalytics()
    
    # Simular transacciones
    # Exchange -> Usuario
    tx1 = Transaction(
        tx_id="tx1",
        inputs=[("exchange_hot_wallet", 100000)],
        outputs=[
            ("user_addr_1", 50000),  # Pago redondo
            ("exchange_change", 50000)
        ],
        timestamp=1
    )
    
    # Usuario hace múltiples transacciones
    tx2 = Transaction(
        tx_id="tx2",
        inputs=[
            ("user_addr_1", 30000),
            ("user_addr_2", 20000)  # Common input!
        ],
        outputs=[
            ("merchant", 45000),
            ("user_addr_3", 4999)  # Change (no redondo)
        ],
        timestamp=2
    )
    
    # Hacia mixer
    tx3 = Transaction(
        tx_id="tx3",
        inputs=[("user_addr_3", 4999)],
        outputs=[("mixer_deposit", 4999)],
        timestamp=3
    )
    
    for tx in [tx1, tx2, tx3]:
        analytics.add_transaction(tx)
    
    # Etiquetar direcciones conocidas
    analytics.label_address("exchange_hot_wallet", "Binance")
    analytics.label_address("exchange_change", "Binance")
    analytics.label_address("mixer_deposit", "Tornado Cash")
    
    # Aplicar heurísticas
    print("--- Applying Heuristics ---")
    analytics.apply_common_input_heuristic()
    analytics.apply_change_heuristic()
    
    stats = analytics.get_cluster_stats()
    print(f"Clusters formed: {stats['num_clusters']}")
    print(f"Addresses clustered: {stats['total_addresses_clustered']}")
    
    # Identificar usuario
    print("\n--- Entity Identification ---")
    for addr in ["user_addr_1", "user_addr_2", "user_addr_3"]:
        entity = analytics.identify_entity(addr)
        print(f"{addr}: {entity}")
    
    # Taint analysis
    print("\n--- Taint Analysis ---")
    analytics.mark_tainted("exchange_hot_wallet", "binance_funds")
    analytics.propagate_taint()
    
    for addr in ["user_addr_1", "user_addr_3", "mixer_deposit"]:
        taint = analytics.get_taint(addr)
        if taint:
            print(f"{addr}: {taint}")
    
    # Tracing
    print("\n--- Fund Tracing ---")
    path = analytics.trace_funds("exchange_hot_wallet")
    print(f"Path from exchange: {' -> '.join(path)}")

chainalysis_demo()
\end{lstlisting}

\subsection{Contramedidas}

\begin{table}[htbp]
\centering
\small
\begin{tabular}{lp{4cm}p{4cm}c}
\toprule
\textbf{Contramedida} & \textbf{Cómo Funciona} & \textbf{Limitaciones} & \textbf{Efectividad} \\
\midrule
Direcciones nuevas & Usar dirección nueva cada tx & Change aún vinculable & Media \\
CoinJoin & Mezclar inputs de múltiples usuarios & Timing/amount correlation & Media-Alta \\
Privacy coins & Privacidad nativa (Monero, Zcash) & Menor liquidez, regulación & Alta \\
Mixers/Pools & Pool de fondos anónimo & Puntos entrada/salida visibles & Alta \\
Timing delays & Randomizar timing de txs & Correlación por montos & Media \\
Amount splitting & Dividir en múltiples txs & Más fees, complejidad & Media \\
P2P/DEX & Evitar exchanges centralizados & Menor liquidez & Media \\
Tor/VPN & Ocultar IP al crear tx & Solo protege metadata de red & Baja \\
\bottomrule
\end{tabular}
\caption{Contramedidas contra análisis de blockchain}
\end{table}

%===========================================
\section{Privacy Pools (Propuesta de Vitalik)}
\label{sec:privacy-pools}
%===========================================

\subsection{Concepto y Motivación}

\begin{definition}[Privacy Pools]
Sistema propuesto por Vitalik Buterin, Jacob Illum, et al. que permite privacidad financiera mientras mantiene la capacidad de demostrar desasociación de fondos ilícitos.

\textbf{Paper:} ``Blockchain Privacy and Regulatory Compliance: Towards a Practical Equilibrium'' (2023)

\textbf{Idea central:} Los usuarios pueden probar que sus fondos pertenecen a un ``association set'' (conjunto de depósitos) que excluye fuentes ilícitas conocidas.
\end{definition}

\begin{keyidea}
\textbf{El dilema resuelto:}

\textbf{Problema con mixers tradicionales:}
\begin{itemize}
    \item Todos los depósitos en el mismo set
    \item Incluye fondos de hacks, ransomware, etc.
    \item Usuario legítimo no puede probar ``inocencia''
\end{itemize}

\textbf{Solución Privacy Pools:}
\begin{itemize}
    \item Múltiples ``association sets'' definidos por comunidad
    \item Usuario elige de qué set quiere probar membresía
    \item Set puede excluir direcciones conocidas como ilícitas
    \item Prueba ZK: ``Mis fondos están en set X'' sin revelar cuál depósito
\end{itemize}
\end{keyidea}

\begin{figure}[htbp]
\centering
\begin{tikzpicture}[
    node distance=0.8cm,
    deposit/.style={circle, draw, minimum size=0.5cm},
    good/.style={circle, draw, minimum size=0.5cm, fill=green!30},
    bad/.style={circle, draw, minimum size=0.5cm, fill=red!30},
    neutral/.style={circle, draw, minimum size=0.5cm, fill=gray!20},
    set/.style={ellipse, draw, dashed, minimum width=4cm, minimum height=2cm}
]
    % Pool completo
    \node at (0, 4) {\textbf{Privacy Pool (todos los depósitos)}};
    
    \draw[thick] (-5, 0) rectangle (5, 3);
    
    % Depósitos
    \node[good] at (-4, 2.5) {};
    \node[good] at (-3, 2.5) {};
    \node[good] at (-2, 2.5) {};
    \node[good] at (-1, 2.5) {};
    \node[neutral] at (0, 2.5) {};
    \node[bad] at (1, 2.5) {H};
    \node[bad] at (2, 2.5) {S};
    \node[neutral] at (3, 2.5) {};
    \node[good] at (4, 2.5) {};
    
    % Association Set 1 (estricto)
    \node[set, fill=green!10] at (-2.5, 1.5) {};
    \node at (-2.5, 0.5) [font=\small] {Set A: Solo KYC'd};
    
    % Association Set 2 (moderado)
    \node[set, fill=yellow!10] at (2, 1.5) {};
    \node at (2, 0.5) [font=\small] {Set B: Excluye hacks};
    
    % Leyenda
    \node at (-4, -1) [font=\small] {\textcolor{green!60!black}{●} KYC verified};
    \node at (-1, -1) [font=\small] {\textcolor{gray}{●} Unknown};
    \node at (2, -1) [font=\small] {\textcolor{red}{H} Hack};
    \node at (4, -1) [font=\small] {\textcolor{red}{S} Sanctioned};
    
    % Explicación
    \node at (0, -2) [align=center, font=\small] {
        Usuario puede probar: ``Mi depósito está en Set A''\\
        Sin revelar cuál depósito específico
    };
\end{tikzpicture}
\caption{Privacy Pools: múltiples association sets con diferentes niveles de restricción}
\label{fig:privacy-pools}
\end{figure}

\subsection{Arquitectura Técnica}

\begin{mathnote}[Componentes de Privacy Pools]
\textbf{1. Deposit Tree:}
\begin{itemize}
    \item Merkle tree de todos los depósitos
    \item Root: $R_{all}$
    \item Cada depósito tiene commitment $C_i$
\end{itemize}

\textbf{2. Association Set Providers (ASPs):}
\begin{itemize}
    \item Entidades que definen y publican association sets
    \item Cada set es subset de depósitos
    \item Publicado como Merkle root: $R_{set}$
\end{itemize}

\textbf{3. Withdrawal Proof:}
\begin{itemize}
    \item ZK proof estándar de posesión de depósito
    \item Prueba: conoce $(v, r)$ tal que $C = \text{Commit}(v, r)$ y $C \in R_{all}$
\end{itemize}

\textbf{4. Association Proof:}
\begin{itemize}
    \item ZK proof adicional de membresía en association set
    \item Prueba: El mismo $C$ también está en $R_{set}$
    \item Sin revelar cuál $C$
\end{itemize}

\textbf{Full Proof:}
\[
\pi = \text{ZKProof}\left(
\begin{array}{l}
\text{Conozco } (v, r, \text{path}_{all}, \text{path}_{set}) \\
\text{tal que } C = \text{Commit}(v, r) \\
\land \text{MerkleVerify}(C, \text{path}_{all}, R_{all}) \\
\land \text{MerkleVerify}(C, \text{path}_{set}, R_{set})
\end{array}
\right)
\]
\end{mathnote}

\begin{lstlisting}[language=circom, caption={Circuito de Privacy Pools}]
pragma circom 2.0.0;

include "circomlib/poseidon.circom";

// Verificar que un commitment pertenece tanto al pool completo
// como a un association set específico
template PrivacyPoolWithdraw(poolLevels, setLevels) {
    // ========== PUBLIC INPUTS ==========
    signal input poolRoot;           // Root del pool completo
    signal input associationSetRoot; // Root del association set elegido
    signal input nullifierHash;      // Hash del nullifier
    signal input recipient;          // Dirección de destino
    
    // ========== PRIVATE INPUTS ==========
    signal input amount;
    signal input secret;
    signal input nullifier;
    
    // Merkle proofs
    signal input poolPathElements[poolLevels];
    signal input poolPathIndices[poolLevels];
    signal input setPathElements[setLevels];
    signal input setPathIndices[setLevels];
    
    // ========== COMPUTATIONS ==========
    
    // 1. Calcular commitment
    component commitmentHasher = Poseidon(3);
    commitmentHasher.inputs[0] <== amount;
    commitmentHasher.inputs[1] <== secret;
    commitmentHasher.inputs[2] <== nullifier;
    
    signal commitment;
    commitment <== commitmentHasher.out;
    
    // 2. Verificar membership en pool completo
    component poolMerkle = MerkleProof(poolLevels);
    poolMerkle.leaf <== commitment;
    for (var i = 0; i < poolLevels; i++) {
        poolMerkle.pathElements[i] <== poolPathElements[i];
        poolMerkle.pathIndices[i] <== poolPathIndices[i];
    }
    poolMerkle.root === poolRoot;
    
    // 3. Verificar membership en association set
    component setMerkle = MerkleProof(setLevels);
    setMerkle.leaf <== commitment;
    for (var i = 0; i < setLevels; i++) {
        setMerkle.pathElements[i] <== setPathElements[i];
        setMerkle.pathIndices[i] <== setPathIndices[i];
    }
    setMerkle.root === associationSetRoot;
    
    // 4. Calcular y verificar nullifier hash
    component nullifierHasher = Poseidon(2);
    nullifierHasher.inputs[0] <== nullifier;
    nullifierHasher.inputs[1] <== secret;
    nullifierHasher.out === nullifierHash;
    
    // 5. Vincular recipient (previene front-running)
    signal recipientSquare;
    recipientSquare <== recipient * recipient;
}

component main {public [poolRoot, associationSetRoot, nullifierHash, recipient]} = 
    PrivacyPoolWithdraw(20, 16);
\end{lstlisting}

\begin{lstlisting}[language=Python, caption={Sistema completo de Privacy Pools}]
from dataclasses import dataclass
from typing import Dict, List, Set, Optional
from hashlib import sha256
from datetime import datetime
import secrets

@dataclass
class Deposit:
    commitment: bytes
    timestamp: datetime
    source_address: str
    amount: int
    leaf_index: int

@dataclass
class AssociationSet:
    name: str
    description: str
    root: bytes
    members: Set[bytes]  # Commitments incluidos
    exclusion_criteria: str
    provider: str
    created_at: datetime

class PrivacyPoolSystem:
    """
    Implementación de Privacy Pools según propuesta de Vitalik et al.
    """
    
    def __init__(self):
        # Pool principal
        self.deposits: List[Deposit] = []
        self.pool_root: bytes = b'\x00' * 32
        self.nullifiers: Set[bytes] = set()
        
        # Association sets
        self.association_sets: Dict[str, AssociationSet] = {}
        
        # Providers de association sets
        self.approved_providers: Set[str] = set()
        
        # Blacklists para construir sets
        self.known_hacks: Set[str] = set()  # Direcciones de hacks
        self.sanctioned: Set[str] = set()   # Direcciones sancionadas
        self.kyc_verified: Set[str] = set() # Direcciones con KYC
    
    def register_provider(self, provider_name: str):
        """Registrar proveedor de association sets."""
        self.approved_providers.add(provider_name)
    
    def deposit(
        self,
        amount: int,
        source_address: str,
        secret: bytes,
        nullifier: bytes
    ) -> Deposit:
        """Depositar en el pool."""
        # Calcular commitment
        commitment = self._compute_commitment(amount, secret, nullifier)
        
        deposit = Deposit(
            commitment=commitment,
            timestamp=datetime.now(),
            source_address=source_address,
            amount=amount,
            leaf_index=len(self.deposits)
        )
        
        self.deposits.append(deposit)
        self._update_pool_root()
        
        return deposit
    
    # ========== ASSOCIATION SET MANAGEMENT ==========
    
    def create_association_set(
        self,
        provider: str,
        name: str,
        description: str,
        inclusion_criteria: str
    ) -> AssociationSet:
        """
        Crear nuevo association set basado en criterios.
        """
        if provider not in self.approved_providers:
            raise ValueError("Unapproved provider")
        
        # Determinar qué commitments incluir
        members = set()
        
        for deposit in self.deposits:
            include = True
            
            if inclusion_criteria == "exclude_hacks":
                if deposit.source_address in self.known_hacks:
                    include = False
                    
            elif inclusion_criteria == "exclude_sanctioned":
                if deposit.source_address in self.sanctioned:
                    include = False
                    
            elif inclusion_criteria == "only_kyc":
                if deposit.source_address not in self.kyc_verified:
                    include = False
                    
            elif inclusion_criteria == "exclude_all_bad":
                if (deposit.source_address in self.known_hacks or
                    deposit.source_address in self.sanctioned):
                    include = False
            
            if include:
                members.add(deposit.commitment)
        
        # Calcular root del set
        root = self._compute_set_root(members)
        
        association_set = AssociationSet(
            name=name,
            description=description,
            root=root,
            members=members,
            exclusion_criteria=inclusion_criteria,
            provider=provider,
            created_at=datetime.now()
        )
        
        self.association_sets[name] = association_set
        return association_set
    
    def update_association_set(self, set_name: str):
        """Actualizar association set con nuevos depósitos."""
        if set_name not in self.association_sets:
            return
        
        aset = self.association_sets[set_name]
        
        # Re-evaluar todos los depósitos
        new_members = set()
        for deposit in self.deposits:
            if self._should_include(deposit, aset.exclusion_criteria):
                new_members.add(deposit.commitment)
        
        aset.members = new_members
        aset.root = self._compute_set_root(new_members)
    
    def _should_include(self, deposit: Deposit, criteria: str) -> bool:
        """Evaluar si depósito cumple criterios de inclusión."""
        if criteria == "exclude_hacks":
            return deposit.source_address not in self.known_hacks
        elif criteria == "exclude_sanctioned":
            return deposit.source_address not in self.sanctioned
        elif criteria == "only_kyc":
            return deposit.source_address in self.kyc_verified
        elif criteria == "exclude_all_bad":
            return (deposit.source_address not in self.known_hacks and
                    deposit.source_address not in self.sanctioned)
        return True
    
    # ========== WITHDRAWAL WITH ASSOCIATION PROOF ==========
    
    def withdraw(
        self,
        deposit: Deposit,
        secret: bytes,
        nullifier: bytes,
        association_set_name: str
    ) -> Dict:
        """
        Retirar con proof de association set.
        """
        # Verificar nullifier no usado
        nullifier_hash = sha256(nullifier + secret).digest()
        if nullifier_hash in self.nullifiers:
            return {'success': False, 'error': 'Already spent'}
        
        # Verificar que deposit está en el set elegido
        if association_set_name not in self.association_sets:
            return {'success': False, 'error': 'Unknown association set'}
        
        aset = self.association_sets[association_set_name]
        if deposit.commitment not in aset.members:
            return {'success': False, 'error': 'Deposit not in association set'}
        
        # Generar proofs
        pool_proof = self._generate_pool_membership_proof(deposit)
        set_proof = self._generate_set_membership_proof(deposit, aset)
        
        # Marcar como gastado
        self.nullifiers.add(nullifier_hash)
        
        return {
            'success': True,
            'nullifier_hash': nullifier_hash.hex(),
            'pool_root': self.pool_root.hex(),
            'association_set_root': aset.root.hex(),
            'association_set_name': aset.name,
            'pool_proof': pool_proof,
            'set_proof': set_proof,
            'verified_claim': f"Deposit in '{aset.name}' set ({aset.exclusion_criteria})"
        }
    
    def verify_withdrawal(self, withdrawal_data: Dict) -> Dict:
        """
        Verificar withdrawal proof.
        """
        # Verificar pool root
        if withdrawal_data['pool_root'] != self.pool_root.hex():
            return {'valid': False, 'error': 'Invalid pool root'}
        
        # Verificar association set root
        set_name = withdrawal_data.get('association_set_name')
        if set_name in self.association_sets:
            expected_root = self.association_sets[set_name].root.hex()
            if withdrawal_data['association_set_root'] != expected_root:
                return {'valid': False, 'error': 'Invalid association set root'}
        
        # Verificar ZK proofs (placeholder)
        
        return {
            'valid': True,
            'association_set': withdrawal_data.get('association_set_name'),
            'claim_verified': withdrawal_data.get('verified_claim')
        }
    
    # ========== BLACKLIST MANAGEMENT ==========
    
    def report_hack(self, address: str, description: str):
        """Reportar dirección como parte de hack."""
        self.known_hacks.add(address)
        # Actualizar todos los sets relevantes
        for set_name in self.association_sets:
            self.update_association_set(set_name)
    
    def add_sanction(self, address: str):
        """Añadir dirección a lista de sanciones."""
        self.sanctioned.add(address)
        for set_name in self.association_sets:
            self.update_association_set(set_name)
    
    def verify_kyc(self, address: str):
        """Marcar dirección como KYC verificada."""
        self.kyc_verified.add(address)
    
    # ========== HELPERS ==========
    
    def _compute_commitment(self, amount: int, secret: bytes, nullifier: bytes) -> bytes:
        data = f"{amount}:{secret.hex()}:{nullifier.hex()}"
        return sha256(data.encode()).digest()
    
    def _compute_set_root(self, members: Set[bytes]) -> bytes:
        if not members:
            return b'\x00' * 32
        sorted_members = sorted(members)
        data = b''.join(sorted_members)
        return sha256(data).digest()
    
    def _update_pool_root(self):
        commitments = [d.commitment for d in self.deposits]
        self.pool_root = self._compute_set_root(set(commitments))
    
    def _generate_pool_membership_proof(self, deposit: Deposit) -> bytes:
        # Placeholder para Merkle proof real
        return sha256(deposit.commitment + self.pool_root).digest()
    
    def _generate_set_membership_proof(self, deposit: Deposit, aset: AssociationSet) -> bytes:
        return sha256(deposit.commitment + aset.root).digest()
    
    def get_stats(self) -> Dict:
        """Obtener estadísticas del sistema."""
        return {
            'total_deposits': len(self.deposits),
            'total_nullifiers_spent': len(self.nullifiers),
            'association_sets': len(self.association_sets),
            'known_hack_addresses': len(self.known_hacks),
            'sanctioned_addresses': len(self.sanctioned),
            'kyc_verified_addresses': len(self.kyc_verified)
        }


def privacy_pools_demo():
    print("=== Privacy Pools Demo ===\n")
    
    system = PrivacyPoolSystem()
    
    # Registrar providers
    system.register_provider("Chainalysis")
    system.register_provider("CommunityDAO")
    
    # Simular algunas direcciones KYC'd
    system.verify_kyc("0xAlice")
    system.verify_kyc("0xBob")
    system.verify_kyc("0xCarol")
    
    # Simular hack conocido
    system.report_hack("0xHacker", "Exchange X hack")
    
    # Simular sanción
    system.add_sanction("0xSanctioned")
    
    # Depósitos
    print("--- Making Deposits ---")
    
    # Usuario legítimo con KYC
    secret1, null1 = secrets.token_bytes(32), secrets.token_bytes(32)
    deposit1 = system.deposit(1000000, "0xAlice", secret1, null1)
    print(f"Deposit 1 (Alice, KYC): {deposit1.commitment.hex()[:16]}...")
    
    # Usuario sin KYC pero legítimo
    secret2, null2 = secrets.token_bytes(32), secrets.token_bytes(32)
    deposit2 = system.deposit(500000, "0xDave", secret2, null2)
    print(f"Deposit 2 (Dave, no KYC): {deposit2.commitment.hex()[:16]}...")
    
    # Hacker
    secret3, null3 = secrets.token_bytes(32), secrets.token_bytes(32)
    deposit3 = system.deposit(2000000, "0xHacker", secret3, null3)
    print(f"Deposit 3 (Hacker): {deposit3.commitment.hex()[:16]}...")
    
    # Crear association sets
    print("\n--- Creating Association Sets ---")
    
    set1 = system.create_association_set(
        "Chainalysis",
        "clean_funds",
        "Excludes known hacks and sanctions",
        "exclude_all_bad"
    )
    print(f"Set 'clean_funds': {len(set1.members)} deposits")
    
    set2 = system.create_association_set(
        "CommunityDAO",
        "kyc_only",
        "Only KYC verified addresses",
        "only_kyc"
    )
    print(f"Set 'kyc_only': {len(set2.members)} deposits")
    
    # Withdrawals
    print("\n--- Withdrawals ---")
    
    # Alice puede retirar de cualquier set
    print("\nAlice withdraws from 'clean_funds':")
    result1 = system.withdraw(deposit1, secret1, null1, "clean_funds")
    print(f"  Success: {result1['success']}")
    print(f"  Verified claim: {result1.get('verified_claim')}")
    
    # Alice también podría usar 'kyc_only'
    print("\nAlice could also use 'kyc_only' set")
    
    # Dave solo puede usar 'clean_funds' (no tiene KYC)
    print("\nDave withdraws from 'clean_funds':")
    result2 = system.withdraw(deposit2, secret2, null2, "clean_funds")
    print(f"  Success: {result2['success']}")
    
    print("\nDave tries 'kyc_only':")
    result3 = system.withdraw(deposit2, secrets.token_bytes(32), secrets.token_bytes(32), "kyc_only")
    # Necesita nuevo nullifier porque ya usó el anterior
    deposit2_new = system.deposit(100000, "0xDave", secrets.token_bytes(32), secrets.token_bytes(32))
    result3 = system.withdraw(deposit2_new, secrets.token_bytes(32), secrets.token_bytes(32), "kyc_only")
    print(f"  Success: {result3['success']}")
    print(f"  Error: {result3.get('error')}")
    
    # Hacker no puede usar ningún set
    print("\nHacker tries 'clean_funds':")
    result4 = system.withdraw(deposit3, secret3, null3, "clean_funds")
    print(f"  Success: {result4['success']}")
    print(f"  Error: {result4.get('error')}")
    
    # Stats
    print("\n--- System Stats ---")
    stats = system.get_stats()
    for k, v in stats.items():
        print(f"  {k}: {v}")

privacy_pools_demo()
\end{lstlisting}

\subsection{Trade-offs de Privacy Pools}

\begin{table}[htbp]
\centering
\small
\begin{tabular}{lp{4cm}p{4cm}}
\toprule
\textbf{Association Set} & \textbf{Ventajas} & \textbf{Desventajas} \\
\midrule
``All deposits'' & Máxima privacidad (mayor anonymity set) & Sin exclusión de malos actores \\
``Exclude hacks'' & Balance razonable & Requiere identificación de hacks \\
``Only KYC'' & Máximo compliance & Menor anonymity set, barrera entrada \\
``Only regulated'' & Aceptable por reguladores & Muy restrictivo \\
\bottomrule
\end{tabular}
\caption{Trade-offs de diferentes association sets}
\end{table}

\begin{bancozk}
\textbf{Privacy Pools en Banco ZK:}

Banco ZK implementa Privacy Pools como sistema central:

\begin{enumerate}
    \item \textbf{Association Set Providers aprobados:}
    \begin{itemize}
        \item Chainalysis (exclusión de hacks)
        \item TRM Labs (sanctions screening)
        \item Community DAO (governance por usuarios)
    \end{itemize}
    
    \item \textbf{Sets disponibles:}
    \begin{itemize}
        \item \texttt{banco\_zk\_standard}: Excluye hacks y sanciones
        \item \texttt{banco\_zk\_kyc}: Solo usuarios con KYC
        \item \texttt{banco\_zk\_defi}: Para integración DeFi
    \end{itemize}
    
    \item \textbf{UX simplificado:}
    \begin{itemize}
        \item Wallet sugiere mejor set según contexto
        \item Regulador puede requerir set específico
        \item Usuario siempre ve opciones disponibles
    \end{itemize}
\end{enumerate}
\end{bancozk}
%===========================================
\section{Opt-in Transparency}
\label{sec:opt-in-transparency}
%===========================================

\subsection{Concepto Fundamental}

\begin{definition}[Opt-in Transparency]
Modelo donde la privacidad es el estado predeterminado, pero el usuario puede elegir voluntariamente revelar información cuando le beneficia:
\begin{itemize}
    \item \textbf{Default:} Todas las transacciones son privadas
    \item \textbf{Opt-in:} Usuario decide qué, cuándo y a quién revelar
    \item \textbf{Beneficio:} Transparencia solo cuando aporta valor
    \item \textbf{Control:} Usuario mantiene soberanía sobre sus datos
\end{itemize}
\end{definition}

\begin{intuicion}
\textbf{Analogía con la vida real:}

En el mundo físico:
\begin{itemize}
    \item Por defecto, nadie sabe cuánto dinero tienes
    \item Pero puedes mostrar tu extracto bancario al solicitar hipoteca
    \item O revelar tu historial crediticio para mejor tasa de interés
    \item La revelación es siempre tu elección
\end{itemize}

En crypto con opt-in transparency:
\begin{itemize}
    \item Por defecto, transacciones privadas
    \item Puedes generar proof de solvencia para contraparte
    \item O revelar historial limpio para obtener servicios
    \item Siempre mantienes control
\end{itemize}
\end{intuicion}

\subsection{Casos de Uso}

\begin{figure}[htbp]
\centering
\begin{tikzpicture}[
    node distance=1cm,
    usecase/.style={rectangle, draw, rounded corners, minimum width=3.5cm, minimum height=1.5cm, align=center, fill=blue!10},
    benefit/.style={rectangle, draw, minimum width=3cm, minimum height=0.6cm, align=center, fill=green!10}
]
    % Casos de uso
    \node[usecase] (proof) at (0, 3) {Proof of\\Solvency};
    \node[usecase] (history) at (5, 3) {Clean History\\Disclosure};
    \node[usecase] (compliance) at (10, 3) {Regulatory\\Compliance};
    
    \node[usecase] (loan) at (0, 0) {Loan\\Application};
    \node[usecase] (business) at (5, 0) {Business\\Partnership};
    \node[usecase] (tax) at (10, 0) {Tax\\Reporting};
    
    % Beneficios
    \node[benefit] at (0, 1.5) {Mejor tasa};
    \node[benefit] at (5, 1.5) {Confianza};
    \node[benefit] at (10, 1.5) {Legal clarity};
    
    % Flechas
    \draw[->, thick] (proof) -- (0, 1.9);
    \draw[->, thick] (history) -- (5, 1.9);
    \draw[->, thick] (compliance) -- (10, 1.9);
    
    \draw[->, thick] (0, 1.1) -- (loan);
    \draw[->, thick] (5, 1.1) -- (business);
    \draw[->, thick] (10, 1.1) -- (tax);
    
    % Título
    \node at (5, 4.5) {\textbf{Opt-in Transparency: El usuario elige revelar}};
\end{tikzpicture}
\caption{Casos de uso de opt-in transparency}
\end{figure}

\begin{mathnote}[Tipos de Revelación Opt-in]
\textbf{1. Proof of Reserves/Solvency:}
\begin{itemize}
    \item Probar: ``Tengo al menos X fondos''
    \item Sin revelar: Monto exacto, fuentes, otras cuentas
    \item Uso: Colateral para préstamos, confianza de contraparte
\end{itemize}

\textbf{2. Transaction History Disclosure:}
\begin{itemize}
    \item Probar: ``Todas mis txs en período P fueron legítimas''
    \item Sin revelar: Contrapartes específicas, otros períodos
    \item Uso: Due diligence para socios, auditorías
\end{itemize}

\textbf{3. Source of Funds:}
\begin{itemize}
    \item Probar: ``Mis fondos vienen de fuentes legítimas''
    \item Sin revelar: Cuál fuente específica
    \item Uso: Exchanges, servicios regulados
\end{itemize}

\textbf{4. Compliance Attestation:}
\begin{itemize}
    \item Probar: ``Cumplo con requisito X''
    \item Sin revelar: Datos subyacentes
    \item Uso: Requisitos regulatorios específicos
\end{itemize}

\textbf{5. Selective Identity Reveal:}
\begin{itemize}
    \item Probar: ``Soy la persona que hizo tx Y''
    \item Con propósito específico
    \item Uso: Reclamar derechos, disputas
\end{itemize}
\end{mathnote}

\subsection{Implementación}

\begin{lstlisting}[language=Python, caption={Sistema de Opt-in Transparency}]
from dataclasses import dataclass
from datetime import datetime, date
from typing import Dict, List, Optional, Set
from hashlib import sha256
from enum import Enum
import secrets

class DisclosureType(Enum):
    SOLVENCY = "solvency"
    HISTORY = "history"
    SOURCE = "source"
    COMPLIANCE = "compliance"
    IDENTITY = "identity"

@dataclass
class DisclosureRequest:
    request_id: str
    requester: str
    disclosure_type: DisclosureType
    parameters: Dict
    expires_at: datetime
    purpose: str

@dataclass
class DisclosureProof:
    request_id: str
    disclosure_type: DisclosureType
    claims: Dict
    zk_proof: bytes
    created_at: datetime
    valid_until: datetime

class OptInTransparencySystem:
    """
    Sistema de transparencia opt-in para privacidad selectiva.
    """
    
    def __init__(self):
        # Usuario data (privado por defecto)
        self.balances: Dict[str, int] = {}
        self.transactions: List[Dict] = []
        self.kyc_data: Dict[str, Dict] = {}
        
        # Disclosure tracking
        self.disclosure_requests: Dict[str, DisclosureRequest] = {}
        self.issued_proofs: List[DisclosureProof] = []
        
        # Revocaciones
        self.revoked_proofs: Set[str] = set()
    
    def set_user_data(
        self,
        user_id: str,
        balance: int,
        transactions: List[Dict],
        kyc: Dict
    ):
        """Configurar datos de usuario (privados)."""
        self.balances[user_id] = balance
        self.transactions = transactions
        self.kyc_data[user_id] = kyc
    
    # ========== DISCLOSURE REQUESTS ==========
    
    def create_disclosure_request(
        self,
        requester: str,
        disclosure_type: DisclosureType,
        parameters: Dict,
        purpose: str,
        valid_hours: int = 24
    ) -> DisclosureRequest:
        """
        Crear solicitud de disclosure.
        El usuario debe aprobarla explícitamente.
        """
        request = DisclosureRequest(
            request_id=secrets.token_hex(16),
            requester=requester,
            disclosure_type=disclosure_type,
            parameters=parameters,
            expires_at=datetime.now() + timedelta(hours=valid_hours),
            purpose=purpose
        )
        
        self.disclosure_requests[request.request_id] = request
        return request
    
    # ========== USER APPROVES DISCLOSURE ==========
    
    def approve_solvency_disclosure(
        self,
        user_id: str,
        request_id: str,
        minimum_balance: int
    ) -> Optional[DisclosureProof]:
        """
        Usuario aprueba disclosure de solvencia.
        Solo revela: "tengo al menos X"
        """
        request = self.disclosure_requests.get(request_id)
        if not request or datetime.now() > request.expires_at:
            return None
        
        actual_balance = self.balances.get(user_id, 0)
        
        # Verificar que cumple
        if actual_balance < minimum_balance:
            return None  # No puede probar lo que no tiene
        
        # Generar proof
        claims = {
            'minimum_balance': minimum_balance,
            'balance_sufficient': True,
            # NO incluye: balance_actual
        }
        
        proof = self._generate_zk_proof(
            'solvency',
            {'balance': actual_balance, 'threshold': minimum_balance}
        )
        
        disclosure = DisclosureProof(
            request_id=request_id,
            disclosure_type=DisclosureType.SOLVENCY,
            claims=claims,
            zk_proof=proof,
            created_at=datetime.now(),
            valid_until=request.expires_at
        )
        
        self.issued_proofs.append(disclosure)
        return disclosure
    
    def approve_history_disclosure(
        self,
        user_id: str,
        request_id: str,
        period_start: date,
        period_end: date,
        required_properties: List[str]
    ) -> Optional[DisclosureProof]:
        """
        Usuario aprueba disclosure de historial.
        Revela propiedades agregadas del período.
        """
        request = self.disclosure_requests.get(request_id)
        if not request:
            return None
        
        # Filtrar transacciones del período
        period_txs = [
            tx for tx in self.transactions
            if period_start <= tx['date'] <= period_end
        ]
        
        claims = {}
        
        if 'total_volume' in required_properties:
            claims['total_volume'] = sum(tx['amount'] for tx in period_txs)
        
        if 'transaction_count' in required_properties:
            claims['transaction_count'] = len(period_txs)
        
        if 'no_suspicious' in required_properties:
            # Verificar que ninguna tx está marcada como sospechosa
            suspicious = [tx for tx in period_txs if tx.get('suspicious', False)]
            claims['no_suspicious'] = len(suspicious) == 0
        
        if 'no_sanctioned' in required_properties:
            # Verificar ninguna contraparte sancionada
            claims['no_sanctioned'] = True  # Placeholder
        
        proof = self._generate_zk_proof('history', claims)
        
        disclosure = DisclosureProof(
            request_id=request_id,
            disclosure_type=DisclosureType.HISTORY,
            claims=claims,
            zk_proof=proof,
            created_at=datetime.now(),
            valid_until=request.expires_at
        )
        
        self.issued_proofs.append(disclosure)
        return disclosure
    
    def approve_source_disclosure(
        self,
        user_id: str,
        request_id: str,
        amount_to_verify: int
    ) -> Optional[DisclosureProof]:
        """
        Usuario aprueba disclosure de fuente de fondos.
        Revela: "estos fondos vienen de fuentes legítimas"
        """
        request = self.disclosure_requests.get(request_id)
        if not request:
            return None
        
        # Verificar que puede demostrar origen legítimo
        # (En producción: análisis real de procedencia)
        
        claims = {
            'amount_verified': amount_to_verify,
            'source_legitimate': True,
            'source_categories': ['regulated_exchange', 'employment'],
            # NO incluye: cuál exchange específico, empleador, etc.
        }
        
        proof = self._generate_zk_proof('source', claims)
        
        disclosure = DisclosureProof(
            request_id=request_id,
            disclosure_type=DisclosureType.SOURCE,
            claims=claims,
            zk_proof=proof,
            created_at=datetime.now(),
            valid_until=request.expires_at
        )
        
        self.issued_proofs.append(disclosure)
        return disclosure
    
    def approve_compliance_disclosure(
        self,
        user_id: str,
        request_id: str,
        compliance_requirements: List[str]
    ) -> Optional[DisclosureProof]:
        """
        Usuario aprueba disclosure de compliance.
        """
        request = self.disclosure_requests.get(request_id)
        if not request:
            return None
        
        kyc = self.kyc_data.get(user_id, {})
        claims = {}
        
        for req in compliance_requirements:
            if req == 'kyc_verified':
                claims['kyc_verified'] = kyc.get('verified', False)
            elif req == 'not_pep':
                claims['not_pep'] = not kyc.get('is_pep', False)
            elif req == 'not_sanctioned':
                claims['not_sanctioned'] = not kyc.get('is_sanctioned', False)
            elif req == 'country_allowed':
                allowed_countries = request.parameters.get('allowed_countries', [])
                claims['country_allowed'] = kyc.get('country') in allowed_countries
            elif req == 'age_verified':
                claims['age_verified'] = kyc.get('age_verified', False)
        
        proof = self._generate_zk_proof('compliance', claims)
        
        disclosure = DisclosureProof(
            request_id=request_id,
            disclosure_type=DisclosureType.COMPLIANCE,
            claims=claims,
            zk_proof=proof,
            created_at=datetime.now(),
            valid_until=request.expires_at
        )
        
        self.issued_proofs.append(disclosure)
        return disclosure
    
    # ========== VERIFICATION ==========
    
    def verify_disclosure(self, proof: DisclosureProof) -> Dict:
        """
        Verificar disclosure proof.
        """
        # Verificar no revocado
        if proof.request_id in self.revoked_proofs:
            return {'valid': False, 'error': 'Proof revoked'}
        
        # Verificar no expirado
        if datetime.now() > proof.valid_until:
            return {'valid': False, 'error': 'Proof expired'}
        
        # Verificar ZK proof
        zk_valid = self._verify_zk_proof(proof.zk_proof)
        
        if not zk_valid:
            return {'valid': False, 'error': 'Invalid ZK proof'}
        
        return {
            'valid': True,
            'disclosure_type': proof.disclosure_type.value,
            'claims': proof.claims,
            'issued_at': proof.created_at.isoformat(),
            'valid_until': proof.valid_until.isoformat()
        }
    
    # ========== REVOCATION ==========
    
    def revoke_disclosure(self, user_id: str, request_id: str) -> bool:
        """
        Usuario revoca un disclosure previamente emitido.
        """
        # Verificar que el proof existe y pertenece al usuario
        proof = next(
            (p for p in self.issued_proofs if p.request_id == request_id),
            None
        )
        
        if not proof:
            return False
        
        self.revoked_proofs.add(request_id)
        return True
    
    # ========== HELPERS ==========
    
    def _generate_zk_proof(self, proof_type: str, data: Dict) -> bytes:
        """Generar ZK proof (placeholder)."""
        data_str = f"{proof_type}:{str(sorted(data.items()))}"
        return sha256(data_str.encode()).digest()
    
    def _verify_zk_proof(self, proof: bytes) -> bool:
        """Verificar ZK proof."""
        return len(proof) == 32


from datetime import timedelta

def opt_in_demo():
    print("=== Opt-in Transparency Demo ===\n")
    
    system = OptInTransparencySystem()
    
    # Configurar datos de usuario (PRIVADOS)
    system.set_user_data(
        user_id="alice",
        balance=1000000,  # 1M
        transactions=[
            {'date': date(2024, 1, 15), 'amount': 50000, 'type': 'salary'},
            {'date': date(2024, 2, 15), 'amount': 50000, 'type': 'salary'},
            {'date': date(2024, 3, 1), 'amount': -10000, 'type': 'payment'},
        ],
        kyc={
            'verified': True,
            'country': 'US',
            'is_pep': False,
            'is_sanctioned': False,
            'age_verified': True
        }
    )
    
    print("Alice's private data configured")
    print("(Balance, transactions, KYC - all private by default)\n")
    
    # ========== ESCENARIO 1: Préstamo ==========
    print("--- Scenario 1: Loan Application ---")
    
    loan_request = system.create_disclosure_request(
        requester="DeFi_Lending_Protocol",
        disclosure_type=DisclosureType.SOLVENCY,
        parameters={'minimum_required': 100000},
        purpose="Collateral verification for loan"
    )
    print(f"Loan protocol requests proof of minimum 100,000")
    
    # Alice aprueba
    solvency_proof = system.approve_solvency_disclosure(
        user_id="alice",
        request_id=loan_request.request_id,
        minimum_balance=100000
    )
    
    print(f"\nAlice approves disclosure:")
    print(f"  Revealed: balance >= 100,000: {solvency_proof.claims['balance_sufficient']}")
    print(f"  NOT revealed: actual balance (1,000,000)")
    
    # Verificar
    verification = system.verify_disclosure(solvency_proof)
    print(f"\nVerification result: {verification['valid']}")
    
    # ========== ESCENARIO 2: Partnership ==========
    print("\n--- Scenario 2: Business Partnership ---")
    
    partner_request = system.create_disclosure_request(
        requester="Potential_Partner_Inc",
        disclosure_type=DisclosureType.HISTORY,
        parameters={},
        purpose="Due diligence for partnership"
    )
    
    history_proof = system.approve_history_disclosure(
        user_id="alice",
        request_id=partner_request.request_id,
        period_start=date(2024, 1, 1),
        period_end=date(2024, 3, 31),
        required_properties=['transaction_count', 'no_suspicious', 'no_sanctioned']
    )
    
    print(f"Partner requests transaction history disclosure")
    print(f"\nAlice approves:")
    print(f"  Transaction count: {history_proof.claims.get('transaction_count')}")
    print(f"  No suspicious activity: {history_proof.claims.get('no_suspicious')}")
    print(f"  No sanctioned parties: {history_proof.claims.get('no_sanctioned')}")
    print(f"  NOT revealed: specific transactions, counterparties, amounts")
    
    # ========== ESCENARIO 3: Exchange ==========
    print("\n--- Scenario 3: Exchange Compliance ---")
    
    exchange_request = system.create_disclosure_request(
        requester="Regulated_Exchange",
        disclosure_type=DisclosureType.COMPLIANCE,
        parameters={'allowed_countries': ['US', 'UK', 'EU']},
        purpose="Regulatory compliance check"
    )
    
    compliance_proof = system.approve_compliance_disclosure(
        user_id="alice",
        request_id=exchange_request.request_id,
        compliance_requirements=['kyc_verified', 'not_pep', 'not_sanctioned', 'country_allowed']
    )
    
    print(f"Exchange requests compliance verification")
    print(f"\nAlice approves:")
    for claim, value in compliance_proof.claims.items():
        print(f"  {claim}: {value}")
    print(f"  NOT revealed: name, address, ID number, date of birth")
    
    # ========== REVOCACIÓN ==========
    print("\n--- Scenario 4: Revocation ---")
    
    revoked = system.revoke_disclosure("alice", loan_request.request_id)
    print(f"Alice revokes solvency proof: {revoked}")
    
    # Intentar verificar después de revocación
    reverify = system.verify_disclosure(solvency_proof)
    print(f"Verification after revocation: {reverify}")

opt_in_demo()
\end{lstlisting}

\subsection{Principios de Diseño}

\begin{keyidea}
\textbf{Principios de Opt-in Transparency:}

\begin{enumerate}
    \item \textbf{Privacy by Default:}
    \begin{itemize}
        \item Sin acción del usuario, todo es privado
        \item La transparencia requiere consentimiento explícito
        \item Nunca revelar más de lo solicitado
    \end{itemize}
    
    \item \textbf{Granular Control:}
    \begin{itemize}
        \item Usuario decide qué, cuándo, a quién, por cuánto tiempo
        \item Diferentes niveles de revelación disponibles
        \item Revocación posible en cualquier momento
    \end{itemize}
    
    \item \textbf{Purpose Binding:}
    \begin{itemize}
        \item Disclosure vinculado a propósito específico
        \item No reutilizable para otros fines
        \item Auditable quién solicitó qué
    \end{itemize}
    
    \item \textbf{Minimal Disclosure:}
    \begin{itemize}
        \item Revelar solo lo mínimo necesario
        \item Preferir proofs sobre datos raw
        \item Agregar cuando sea posible
    \end{itemize}
\end{enumerate}
\end{keyidea}

%===========================================
\section*{Resumen del Bloque}
%===========================================

\begin{keyidea}
\textbf{Puntos clave de Compliance vs Privacidad (Bloque 15):}

\textbf{1. Proof-of-Innocence:}
\begin{itemize}
    \item Probar que fondos no son ilícitos sin revelar origen
    \item Tipos: temporal, non-inclusion, inclusion, category
    \item Circuitos ZK para cada tipo de proof
\end{itemize}

\textbf{2. Selective De-anonymization:}
\begin{itemize}
    \item Threshold decryption: k de n autoridades
    \item Timelock encryption: revelación temporal
    \item Court-order escrow: acceso judicial
\end{itemize}

\textbf{3. Viewing Keys:}
\begin{itemize}
    \item Jerarquía: spending > viewing > audit
    \item Diferentes capacidades por tipo de clave
    \item Reportes verificables sin revelar todo
\end{itemize}

\textbf{4. Regulatory-Compliant Mixers:}
\begin{itemize}
    \item Sanctions screening en entrada/salida
    \item KYC opcional con beneficios
    \item Association sets para proof-of-innocence
\end{itemize}

\textbf{5. Travel Rule + ZK:}
\begin{itemize}
    \item Proof de KYC válido sin compartir datos
    \item Datos encriptados para regulador
    \item Compliance sin surveillance masiva
\end{itemize}

\textbf{6. Chainalysis y Contramedidas:}
\begin{itemize}
    \item Técnicas: clustering, taint, ML
    \item Contramedidas: mixers, direcciones nuevas, timing
    \item Trade-off privacidad vs usabilidad
\end{itemize}

\textbf{7. Privacy Pools:}
\begin{itemize}
    \item Association sets definidos por comunidad
    \item Usuario elige nivel de restricción
    \item Balance ajustable privacidad/compliance
\end{itemize}

\textbf{8. Opt-in Transparency:}
\begin{itemize}
    \item Privacidad por defecto
    \item Revelación cuando beneficia al usuario
    \item Control granular y revocable
\end{itemize}
\end{keyidea}

\begin{bancozk}
\textbf{Compliance Framework de Banco ZK:}

\begin{center}
\begin{tikzpicture}[
    layer/.style={rectangle, draw, minimum width=11cm, minimum height=0.8cm, align=center}
]
    \node[layer, fill=red!15] at (0, 4) {Capa Regulatoria: Travel Rule ZK + Reporting};
    \node[layer, fill=orange!15] at (0, 3) {Capa Compliance: Privacy Pools + Association Sets};
    \node[layer, fill=yellow!15] at (0, 2) {Capa Auditoría: Viewing Keys + Audit Proofs};
    \node[layer, fill=green!15] at (0, 1) {Capa Usuario: Opt-in Transparency + PoI};
    \node[layer, fill=blue!15] at (0, 0) {Capa Base: Privacidad por Defecto (ZK)};
\end{tikzpicture}
\end{center}

\textbf{Flujo de compliance:}
\begin{enumerate}
    \item Usuario deposita (sanctions check automático)
    \item Opera con privacidad total dentro del pool
    \item Al retirar, elige association set apropiado
    \item Si necesita, genera proofs de compliance
    \item Viewing keys disponibles para auditoría voluntaria
\end{enumerate}

\textbf{Configuraciones de usuario:}
\begin{itemize}
    \item \textbf{Modo Privado:} Sin compliance extra, usa set más amplio
    \item \textbf{Modo Compliant:} KYC + escrow, acceso a pools premium
    \item \textbf{Modo Institucional:} Full viewing keys, reportes automáticos
\end{itemize}

\textbf{Métricas objetivo:}
\begin{center}
\begin{tabular}{lc}
\toprule
\textbf{Métrica} & \textbf{Target} \\
\midrule
Sanctions check latency & $<$ 1 segundo \\
Association set update & Cada 100 bloques \\
PoI verification time & $<$ 500 ms \\
Viewing key derivation & Instantáneo \\
Audit report generation & $<$ 5 segundos \\
\bottomrule
\end{tabular}
\end{center}
\end{bancozk}

\begin{table}[htbp]
\centering
\small
\begin{tabular}{lp{3.5cm}p{3.5cm}c}
\toprule
\textbf{Técnica} & \textbf{Privacidad} & \textbf{Compliance} & \textbf{Complejidad} \\
\midrule
Proof-of-Innocence & Alta (no revela origen) & Alta (excluye ilícitos) & Media \\
Threshold De-anon & Media (revela con k/n) & Alta (acceso judicial) & Alta \\
Viewing Keys & Ajustable & Alta (auditoría) & Baja \\
Privacy Pools & Ajustable por set & Ajustable por set & Media \\
Opt-in Transparency & Alta (control usuario) & Media (voluntario) & Baja \\
ZK Travel Rule & Alta & Alta & Alta \\
\bottomrule
\end{tabular}
\caption{Resumen de técnicas de balance compliance-privacidad}
\end{table}

\begin{figure}[htbp]
\centering
\begin{tikzpicture}
    % Ejes
    \draw[thick, ->] (0, 0) -- (10, 0) node[right] {Compliance};
    \draw[thick, ->] (0, 0) -- (0, 6) node[above] {Privacidad};
    
    % Zona ideal
    \fill[green!20] (4, 3) -- (8, 3) -- (8, 5) -- (4, 5) -- cycle;
    \node at (6, 4) [font=\small] {Zona Óptima ZK};
    
    % Puntos
    \node[circle, fill=red, minimum size=0.3cm] at (1, 5) {};
    \node at (1, 5.5) [font=\small] {Monero};
    
    \node[circle, fill=orange, minimum size=0.3cm] at (3, 4) {};
    \node at (3, 4.5) [font=\small] {Tornado};
    
    \node[circle, fill=yellow, minimum size=0.3cm] at (5, 4.5) {};
    \node at (5, 5) [font=\small] {Privacy Pools};
    
    \node[circle, fill=green, minimum size=0.3cm] at (7, 4) {};
    \node at (7, 4.5) [font=\small] {Banco ZK};
    
    \node[circle, fill=blue, minimum size=0.3cm] at (9, 1) {};
    \node at (9, 1.5) [font=\small] {Bitcoin};
    
    % Diagonal trade-off tradicional
    \draw[dashed, gray] (0, 5.5) -- (9.5, 0.5);
    \node at (2, 0.8) [font=\footnotesize, gray, rotate=-30] {Trade-off tradicional};
\end{tikzpicture}
\caption{ZK permite superar el trade-off tradicional entre privacidad y compliance}
\label{fig:privacy-compliance-final}
\end{figure}

% Fin del Bloque 15

\cleardoublepage

% ============================================
% BLOQUE 16: METADATA Y NETWORK PRIVACY - VERSIÓN EXPANDIDA
% Parte 1: Introducción, IP Leakage, Timing Correlation
% Objetivo: ~60 páginas (~3600 líneas)
% ============================================

\chapter{Metadata y Network Privacy}
\label{ch:bloque16}

\begin{center}
\textit{Bloque 16 de 20}
\end{center}

\section*{Introducción}

Las pruebas de conocimiento cero protegen la privacidad a nivel de transacción, pero la metadata de red puede revelar información crítica sobre los usuarios. Este bloque explora las amenazas a nivel de red y las técnicas para proteger la privacidad más allá del contenido de las transacciones.

\begin{keyidea}
\textbf{El problema de la metadata:}

\textbf{Protección a nivel transacción (ZK):}
\begin{itemize}
    \item Monto oculto
    \item Remitente/receptor ocultos
    \item Origen de fondos no vinculable
\end{itemize}

\textbf{Metadata que aún se filtra:}
\begin{itemize}
    \item IP del usuario al enviar tx
    \item Timing de cuándo se envía
    \item Patrones de actividad
    \item Con quién se comunica
    \item Tamaño de los mensajes
\end{itemize}

\textbf{Consecuencia:} Un adversario de red puede vincular tu identidad con tus transacciones anónimas observando CÓMO te comunicas, no QUÉ comunicas.
\end{keyidea}

\begin{figure}[htbp]
\centering
\begin{tikzpicture}[
    node distance=1cm,
    layer/.style={rectangle, draw, minimum width=10cm, minimum height=1cm, align=center}
]
    % Capas de privacidad
    \node[layer, fill=green!30] (app) at (0, 4) {Capa Aplicación: ZK Proofs (Protegido)};
    \node[layer, fill=yellow!30] (trans) at (0, 3) {Capa Transporte: TLS (Parcialmente protegido)};
    \node[layer, fill=orange!30] (net) at (0, 2) {Capa Red: IP, Routing (Expuesto)};
    \node[layer, fill=red!30] (phys) at (0, 1) {Capa Física: ISP, timing (Expuesto)};
    
    % Adversarios
    \node at (6, 4) [font=\small] {Blockchain observers};
    \node at (6, 3) [font=\small] {Man-in-the-middle};
    \node at (6, 2) [font=\small] {Network observers};
    \node at (6, 1) [font=\small] {ISP, gobiernos};
    
    % Flechas de observación
    \draw[->, red, thick] (5, 2) -- (4, 2);
    \draw[->, red, thick] (5, 1) -- (4, 1);
\end{tikzpicture}
\caption{Capas de privacidad: ZK protege la aplicación pero la red sigue expuesta}
\label{fig:privacy-layers}
\end{figure}

\begin{mathnote}[Modelo de Amenazas de Red]
\textbf{Adversarios potenciales:}
\begin{enumerate}
    \item \textbf{ISP (Internet Service Provider):}
    \begin{itemize}
        \item Ve todo tu tráfico de red
        \item Conoce tu IP real y ubicación
        \item Puede registrar timing de todas las conexiones
    \end{itemize}
    
    \item \textbf{Nodos de la red P2P:}
    \begin{itemize}
        \item Ven IP de quien envía transacciones
        \item Pueden inferir el origen de una tx
        \item Posibles nodos maliciosos (sybil)
    \end{itemize}
    
    \item \textbf{Proveedores de RPC:}
    \begin{itemize}
        \item Infura, Alchemy ven todas tus queries
        \item Pueden vincular direcciones con IP
        \item Conocen tu actividad on-chain
    \end{itemize}
    
    \item \textbf{Global Passive Adversary (GPA):}
    \begin{itemize}
        \item Puede observar todo el tráfico de Internet
        \item Correlación de tráfico entre endpoints
        \item Agencias de inteligencia nacionales
    \end{itemize}
\end{enumerate}

\textbf{Capacidades del adversario:}
\begin{itemize}
    \item Observar: Ver tráfico que pasa
    \item Registrar: Almacenar para análisis futuro
    \item Correlacionar: Vincular diferentes flujos de tráfico
    \item Modificar: Inyectar, dropear, o retrasar paquetes
\end{itemize}
\end{mathnote}

%===========================================
\section{IP Leakage y Protección}
\label{sec:ip-leakage}
%===========================================

\subsection{Anatomía del IP Leakage}

\begin{definition}[IP Leakage]
Exposición no intencional de la dirección IP real de un usuario durante la interacción con sistemas blockchain, permitiendo vincular la identidad real con actividad on-chain.

\textbf{Consecuencias:}
\begin{itemize}
    \item Geolocalización del usuario
    \item Vinculación de múltiples sesiones
    \item Identificación personal (con cooperación del ISP)
    \item Correlación de direcciones on-chain con identidad
\end{itemize}
\end{definition}

\begin{figure}[htbp]
\centering
\begin{tikzpicture}[
    node distance=1.5cm,
    user/.style={rectangle, draw, minimum width=2cm, minimum height=1cm, fill=blue!20},
    node/.style={circle, draw, minimum size=1cm, fill=green!20},
    service/.style={rectangle, draw, minimum width=2cm, minimum height=1cm, fill=orange!20},
    observer/.style={rectangle, draw, dashed, minimum width=1.5cm, minimum height=0.8cm, fill=red!10},
    arrow/.style={->, thick}
]
    % Usuario
    \node[user] (user) at (0, 0) {Usuario\\IP: 1.2.3.4};
    
    % Vectores de leak
    \node[node] (p2p) at (4, 2) {P2P\\Node};
    \node[service] (rpc) at (4, 0) {RPC\\(Infura)};
    \node[service] (dapp) at (4, -2) {dApp\\Frontend};
    
    % Conexiones (leaking IP)
    \draw[arrow, red] (user) -- node[above, font=\small] {IP visible} (p2p);
    \draw[arrow, red] (user) -- node[above, font=\small] {IP visible} (rpc);
    \draw[arrow, red] (user) -- node[above, font=\small] {IP visible} (dapp);
    
    % Observadores
    \node[observer] (isp) at (2, -3.5) {ISP};
    \draw[dashed, red] (isp) -- (1, -1);
    
    \node[observer] (gov) at (6, -3.5) {Gobierno};
    \draw[dashed, red] (gov) -- (rpc);
    
    % Etiqueta
    \node at (2, 3) [font=\small, align=center] {Cada conexión expone\\la IP real del usuario};
\end{tikzpicture}
\caption{Vectores de IP leakage en interacción con blockchain}
\label{fig:ip-leakage-vectors}
\end{figure}

\subsection{Vectores de IP Leakage}

\begin{mathnote}[Vectores Detallados]
\textbf{1. Conexión directa a nodos P2P:}
\begin{itemize}
    \item Al sincronizar blockchain, conectas a múltiples peers
    \item Cada peer ve tu IP
    \item Nodos maliciosos pueden registrar IPs de quien envía txs
    \item ``First-seen'' heuristic: el primer nodo que ve una tx probablemente es el origen
\end{itemize}

\textbf{2. Uso de RPC públicos:}
\begin{itemize}
    \item Infura, Alchemy, QuickNode procesan tus requests
    \item Ven: IP + queries + direcciones consultadas
    \item Pueden construir perfil completo de actividad
    \item Términos de servicio permiten logging
\end{itemize}

\textbf{3. Broadcast de transacciones:}
\begin{itemize}
    \item Al enviar tx, el primer nodo conoce tu IP
    \item Incluso con multiple peers, timing puede revelar origen
    \item Transacciones ``frescas'' más fáciles de rastrear
\end{itemize}

\textbf{4. Conexión a dApps:}
\begin{itemize}
    \item Frontend web conoce tu IP (logs del servidor)
    \item Analytics de terceros (Google Analytics, etc.)
    \item WebRTC puede filtrar IP real incluso con VPN
    \item Cookies vinculan sesiones
\end{itemize}

\textbf{5. DNS Leaks:}
\begin{itemize}
    \item Queries DNS revelan qué dominios visitas
    \item ISP ve todas las consultas DNS por defecto
    \item Incluso con VPN, DNS puede filtrar por configuración incorrecta
\end{itemize}
\end{mathnote}

\begin{lstlisting}[language=Python, caption={Detector de IP Leakage}]
from dataclasses import dataclass
from datetime import datetime
from typing import List, Dict, Set, Optional
from enum import Enum
import re

class LeakageVector(Enum):
    P2P_CONNECTION = "p2p_connection"
    RPC_REQUEST = "rpc_request"
    DAPP_FRONTEND = "dapp_frontend"
    TX_BROADCAST = "tx_broadcast"
    DNS_QUERY = "dns_query"
    WEBRTC = "webrtc"

@dataclass
class LeakageEvent:
    vector: LeakageVector
    timestamp: datetime
    ip_exposed: str
    destination: str
    metadata: Dict
    risk_level: str  # 'low', 'medium', 'high', 'critical'

class IPLeakageDetector:
    """
    Detector de eventos de IP leakage en interacción con blockchain.
    """
    
    def __init__(self, user_real_ip: str):
        self.real_ip = user_real_ip
        self.events: List[LeakageEvent] = []
        
        # Servicios conocidos y su riesgo
        self.known_rpc_providers = {
            'infura.io': 'high',
            'alchemy.com': 'high',
            'quicknode.com': 'medium',
            'cloudflare-eth.com': 'medium',
        }
        
        self.known_analytics = {
            'google-analytics.com': 'high',
            'segment.io': 'high',
            'mixpanel.com': 'medium',
        }
        
        # Nodos P2P conocidos como honeypots
        self.suspicious_nodes: Set[str] = set()
    
    def analyze_connection(
        self,
        destination_ip: str,
        destination_domain: Optional[str],
        connection_type: str,
        headers: Optional[Dict] = None
    ) -> Optional[LeakageEvent]:
        """
        Analizar una conexión saliente para detectar leakage.
        """
        risk_level = 'low'
        vector = None
        metadata = {}
        
        # Determinar vector
        if connection_type == 'p2p':
            vector = LeakageVector.P2P_CONNECTION
            
            # Verificar si es nodo sospechoso
            if destination_ip in self.suspicious_nodes:
                risk_level = 'critical'
                metadata['warning'] = 'Known honeypot node'
            else:
                risk_level = 'medium'
                
        elif connection_type == 'http' and destination_domain:
            # Verificar si es RPC
            for provider, provider_risk in self.known_rpc_providers.items():
                if provider in destination_domain:
                    vector = LeakageVector.RPC_REQUEST
                    risk_level = provider_risk
                    metadata['provider'] = provider
                    break
            
            # Verificar analytics
            for analytics, analytics_risk in self.known_analytics.items():
                if analytics in destination_domain:
                    vector = LeakageVector.DAPP_FRONTEND
                    risk_level = 'high'
                    metadata['analytics'] = analytics
                    break
            
            if not vector:
                vector = LeakageVector.DAPP_FRONTEND
                risk_level = 'low'
                
        elif connection_type == 'dns':
            vector = LeakageVector.DNS_QUERY
            risk_level = 'medium'
            metadata['query'] = destination_domain
            
        elif connection_type == 'webrtc':
            vector = LeakageVector.WEBRTC
            risk_level = 'critical'
            metadata['warning'] = 'WebRTC can leak real IP even with VPN'
        
        if vector:
            event = LeakageEvent(
                vector=vector,
                timestamp=datetime.now(),
                ip_exposed=self.real_ip,
                destination=destination_domain or destination_ip,
                metadata=metadata,
                risk_level=risk_level
            )
            self.events.append(event)
            return event
        
        return None
    
    def analyze_tx_broadcast(
        self,
        tx_hash: str,
        connected_peers: List[str],
        broadcast_method: str
    ) -> LeakageEvent:
        """
        Analizar broadcast de transacción.
        """
        risk_level = 'high'
        metadata = {
            'tx_hash': tx_hash,
            'num_peers': len(connected_peers),
            'method': broadcast_method
        }
        
        # Métodos más seguros
        if broadcast_method == 'dandelion':
            risk_level = 'medium'
            metadata['protection'] = 'Dandelion routing'
        elif broadcast_method == 'tor':
            risk_level = 'low'
            metadata['protection'] = 'Tor network'
        elif broadcast_method == 'direct':
            risk_level = 'critical'
            metadata['warning'] = 'Direct broadcast exposes origin'
        
        # Verificar honeypots en peers
        honeypot_count = sum(1 for p in connected_peers if p in self.suspicious_nodes)
        if honeypot_count > 0:
            risk_level = 'critical'
            metadata['honeypots_detected'] = honeypot_count
        
        event = LeakageEvent(
            vector=LeakageVector.TX_BROADCAST,
            timestamp=datetime.now(),
            ip_exposed=self.real_ip,
            destination=f"{len(connected_peers)} peers",
            metadata=metadata,
            risk_level=risk_level
        )
        
        self.events.append(event)
        return event
    
    def get_risk_summary(self) -> Dict:
        """
        Obtener resumen de riesgos detectados.
        """
        risk_counts = {'low': 0, 'medium': 0, 'high': 0, 'critical': 0}
        vector_counts = {}
        
        for event in self.events:
            risk_counts[event.risk_level] += 1
            vector_name = event.vector.value
            vector_counts[vector_name] = vector_counts.get(vector_name, 0) + 1
        
        overall_risk = 'low'
        if risk_counts['critical'] > 0:
            overall_risk = 'critical'
        elif risk_counts['high'] > 0:
            overall_risk = 'high'
        elif risk_counts['medium'] > 0:
            overall_risk = 'medium'
        
        return {
            'overall_risk': overall_risk,
            'total_events': len(self.events),
            'risk_breakdown': risk_counts,
            'vectors': vector_counts,
            'recommendations': self._get_recommendations(risk_counts)
        }
    
    def _get_recommendations(self, risk_counts: Dict) -> List[str]:
        """Generar recomendaciones basadas en riesgos."""
        recommendations = []
        
        if risk_counts['critical'] > 0:
            recommendations.append("URGENT: Use Tor for all blockchain interactions")
            recommendations.append("Disable WebRTC in browser")
        
        if risk_counts['high'] > 0:
            recommendations.append("Consider running your own RPC node")
            recommendations.append("Use privacy-focused RPC providers")
        
        if any(e.vector == LeakageVector.TX_BROADCAST and e.risk_level != 'low' 
               for e in self.events):
            recommendations.append("Enable Dandelion++ for transaction broadcast")
        
        if any(e.vector == LeakageVector.DNS_QUERY for e in self.events):
            recommendations.append("Use encrypted DNS (DoH/DoT)")
        
        return recommendations


def ip_leakage_demo():
    print("=== IP Leakage Detection Demo ===\n")
    
    detector = IPLeakageDetector(user_real_ip="192.168.1.100")
    
    # Simular actividad típica de usuario
    print("--- Simulating user activity ---\n")
    
    # Conexión a RPC
    event1 = detector.analyze_connection(
        destination_ip="13.56.163.26",
        destination_domain="mainnet.infura.io",
        connection_type="http"
    )
    print(f"RPC Connection: {event1.risk_level} risk")
    
    # Conexión P2P
    event2 = detector.analyze_connection(
        destination_ip="45.76.33.21",
        destination_domain=None,
        connection_type="p2p"
    )
    print(f"P2P Connection: {event2.risk_level} risk")
    
    # dApp con analytics
    event3 = detector.analyze_connection(
        destination_ip="142.250.190.46",
        destination_domain="www.google-analytics.com",
        connection_type="http"
    )
    print(f"Analytics: {event3.risk_level} risk")
    
    # Broadcast de tx
    event4 = detector.analyze_tx_broadcast(
        tx_hash="0xabc123...",
        connected_peers=["45.76.33.21", "185.199.108.153", "140.82.114.3"],
        broadcast_method="direct"
    )
    print(f"TX Broadcast: {event4.risk_level} risk")
    
    # Resumen
    print("\n--- Risk Summary ---")
    summary = detector.get_risk_summary()
    print(f"Overall Risk: {summary['overall_risk'].upper()}")
    print(f"Total Events: {summary['total_events']}")
    print(f"Risk Breakdown: {summary['risk_breakdown']}")
    
    print("\n--- Recommendations ---")
    for rec in summary['recommendations']:
        print(f"  • {rec}")

ip_leakage_demo()
\end{lstlisting}

\subsection{Soluciones de Protección}

\subsubsection{VPN (Virtual Private Network)}

\begin{definition}[VPN]
Red privada virtual que encripta el tráfico y lo enruta a través de un servidor intermedio, ocultando la IP real del usuario.

\textbf{Modelo de seguridad:}
\begin{itemize}
    \item Tráfico encriptado entre usuario y servidor VPN
    \item Destino ve IP del servidor VPN, no del usuario
    \item Requiere confiar en el proveedor de VPN
\end{itemize}
\end{definition}

\begin{advertencia}
\textbf{Limitaciones de VPN:}
\begin{enumerate}
    \item \textbf{Confianza centralizada:} El proveedor VPN ve TODO tu tráfico
    \item \textbf{Logging:} Muchos proveedores mantienen logs (aunque digan que no)
    \item \textbf{Jurisdicción:} Pueden ser forzados a entregar datos
    \item \textbf{Single point of failure:} Si VPN se compromete, todo se compromete
    \item \textbf{DNS leaks:} Configuración incorrecta puede filtrar DNS
    \item \textbf{WebRTC leaks:} Navegadores pueden filtrar IP real
\end{enumerate}

\textbf{Para privacidad blockchain:} VPN es insuficiente como única protección.
\end{advertencia}

\subsubsection{Tor (The Onion Router)}

\begin{definition}[Tor]
Red de anonimato que enruta tráfico a través de múltiples nodos (relays), aplicando capas de encriptación:

\textbf{Arquitectura:}
\begin{enumerate}
    \item \textbf{Guard node:} Primer nodo, conoce IP real pero no destino
    \item \textbf{Middle node:} Nodo intermedio, no conoce origen ni destino
    \item \textbf{Exit node:} Último nodo, conoce destino pero no origen
\end{enumerate}

\textbf{Propiedades:}
\begin{itemize}
    \item Ningún nodo conoce tanto origen como destino
    \item Requiere compromiso de 3 nodos para desanonimizar
    \item Código abierto y auditado
\end{itemize}
\end{definition}

\begin{figure}[htbp]
\centering
\begin{tikzpicture}[
    node distance=2cm,
    user/.style={rectangle, draw, minimum width=1.5cm, minimum height=1cm, fill=blue!20},
    relay/.style={circle, draw, minimum size=1.2cm, fill=green!20},
    dest/.style={rectangle, draw, minimum width=1.5cm, minimum height=1cm, fill=orange!20},
    arrow/.style={->, thick}
]
    % Nodos
    \node[user] (user) at (0, 0) {Usuario};
    \node[relay] (guard) at (3, 0) {Guard};
    \node[relay] (middle) at (6, 0) {Middle};
    \node[relay] (exit) at (9, 0) {Exit};
    \node[dest] (dest) at (12, 0) {Destino};
    
    % Conexiones con capas de encriptación
    \draw[arrow, very thick, blue] (user) -- node[above, font=\small] {3 capas} (guard);
    \draw[arrow, thick, green!60!black] (guard) -- node[above, font=\small] {2 capas} (middle);
    \draw[arrow, thick, orange] (middle) -- node[above, font=\small] {1 capa} (exit);
    \draw[arrow] (exit) -- node[above, font=\small] {cleartext} (dest);
    
    % Lo que conoce cada nodo
    \node at (0, -1.5) [font=\small, align=center] {Conoce:\\Todo};
    \node at (3, -1.5) [font=\small, align=center] {Conoce:\\IP usuario\\Middle};
    \node at (6, -1.5) [font=\small, align=center] {Conoce:\\Guard\\Exit};
    \node at (9, -1.5) [font=\small, align=center] {Conoce:\\Middle\\Destino};
    \node at (12, -1.5) [font=\small, align=center] {Conoce:\\Exit};
\end{tikzpicture}
\caption{Onion routing en Tor: cada relay solo conoce el anterior y siguiente}
\label{fig:tor-routing}
\end{figure}

\subsubsection{I2P (Invisible Internet Project)}

\begin{definition}[I2P]
Red overlay diseñada para comunicación anónima interna (servicios ocultos):

\textbf{Diferencias con Tor:}
\begin{itemize}
    \item Optimizado para servicios internos, no exit a clearnet
    \item Usa ``garlic routing'' (múltiples mensajes empaquetados)
    \item Todos los usuarios son routers (más descentralizado)
    \item Conexiones unidireccionales (túneles separados in/out)
\end{itemize}

\textbf{Uso en blockchain:}
\begin{itemize}
    \item Monero integra I2P (Kovri)
    \item Mejor para comunicación P2P entre nodos
    \item Menos adecuado para acceder a servicios web externos
\end{itemize}
\end{definition}

\subsubsection{Nodo Propio}

\begin{keyidea}
\textbf{Ejecutar nodo propio:}

La forma más directa de evitar IP leakage a terceros es ejecutar tu propio nodo:

\begin{itemize}
    \item \textbf{Full node:} Verifica todas las transacciones localmente
    \item \textbf{Sin RPC externos:} No dependes de Infura/Alchemy
    \item \textbf{Broadcast privado:} Puedes combinar con Tor
    \item \textbf{Sincronización:} Requiere ancho de banda y almacenamiento
\end{itemize}

\textbf{Configuración recomendada:}
\begin{enumerate}
    \item Full node detrás de Tor
    \item Conexiones P2P solo a través de .onion
    \item Broadcast de txs con delay aleatorio
    \item Sin conexiones entrantes (solo outbound)
\end{enumerate}
\end{keyidea}

%===========================================
\section{Timing Correlation Attacks}
\label{sec:timing-correlation}
%===========================================

\subsection{Anatomía del Ataque}

\begin{definition}[Timing Correlation Attack]
Ataque que vincula eventos aparentemente no relacionados basándose en la correlación temporal entre ellos.

\textbf{Principio:} Si el evento A y el evento B ocurren consistentemente juntos en el tiempo, probablemente están relacionados.

\textbf{Formalización:}
\[
P(\text{vinculados} | \Delta t < \epsilon) >> P(\text{vinculados} | \Delta t > \epsilon)
\]

donde $\Delta t$ es la diferencia temporal entre eventos.
\end{definition}

\begin{mathnote}[Tipos de Timing Correlation]
\textbf{1. Deposit-Withdraw Correlation:}
\begin{itemize}
    \item Usuario deposita en mixer, retira poco después
    \item Correlación temporal revela vínculo
    \item Ejemplo: Deposita 1 ETH a las 14:00, retira 1 ETH a las 14:05
\end{itemize}

\textbf{2. Login-Transaction Correlation:}
\begin{itemize}
    \item Usuario entra a exchange, segundos después aparece tx
    \item El exchange puede vincular la tx con el usuario
    \item Incluso si la tx viene de dirección ``anónima''
\end{itemize}

\textbf{3. Behavioral Patterns:}
\begin{itemize}
    \item Usuario siempre activo a las 9am (zona horaria)
    \item Patrones de fin de semana diferentes
    \item Correlación con eventos del mundo real
\end{itemize}

\textbf{4. Cross-Service Correlation:}
\begin{itemize}
    \item Actividad en servicio A correlaciona con servicio B
    \item Mismo timing = mismo usuario
    \item Aunque use diferentes direcciones/identidades
\end{itemize}
\end{mathnote}

\begin{figure}[htbp]
\centering
\begin{tikzpicture}[
    node distance=0.5cm,
    event/.style={rectangle, draw, minimum width=1cm, minimum height=0.6cm, font=\small},
    timeline/.style={thick, ->}
]
    % Timeline
    \draw[timeline] (0, 2) -- (12, 2) node[right] {tiempo};
    \draw[timeline] (0, 0) -- (12, 0) node[right] {tiempo};
    
    % Etiquetas
    \node at (-1, 2) {Mixer};
    \node at (-1, 0) {Exchange};
    
    % Eventos mixer
    \node[event, fill=blue!20] at (2, 2) {D1};
    \node[event, fill=green!20] at (3, 2) {W1};
    \node[event, fill=blue!20] at (6, 2) {D2};
    \node[event, fill=green!20] at (9, 2) {W2};
    
    % Eventos exchange
    \node[event, fill=orange!20] at (2.9, 0) {Login};
    \node[event, fill=orange!20] at (9.1, 0) {Deposit};
    
    % Correlaciones
    \draw[dashed, red, thick] (3, 1.7) -- (2.9, 0.3);
    \draw[dashed, red, thick] (9, 1.7) -- (9.1, 0.3);
    
    % Etiquetas de correlación
    \node at (3, 1) [font=\small, red] {$\Delta t_1$};
    \node at (9, 1) [font=\small, red] {$\Delta t_2$};
    
    % Conclusión
    \node at (6, -1) [draw, rectangle, fill=red!10, font=\small] {
        Si $\Delta t_1$ y $\Delta t_2$ son pequeños: W1 y W2 son del mismo usuario
    };
\end{tikzpicture}
\caption{Timing correlation: pequeños $\Delta t$ revelan vinculación}
\label{fig:timing-correlation}
\end{figure}

\begin{lstlisting}[language=Python, caption={Sistema de análisis y defensa contra timing correlation}]
from dataclasses import dataclass
from datetime import datetime, timedelta
from typing import List, Dict, Tuple, Optional
import random
import math
from collections import defaultdict

@dataclass
class Event:
    event_id: str
    event_type: str
    timestamp: datetime
    metadata: Dict

class TimingCorrelationAnalyzer:
    """
    Analizador de correlaciones temporales para detectar vinculaciones.
    """
    
    def __init__(self, correlation_threshold_seconds: float = 60):
        self.threshold = correlation_threshold_seconds
        self.events: List[Event] = []
    
    def add_event(self, event: Event):
        """Añadir evento para análisis."""
        self.events.append(event)
    
    def find_correlated_pairs(
        self,
        event_type_a: str,
        event_type_b: str,
        max_delta_seconds: float = None
    ) -> List[Tuple[Event, Event, float]]:
        """
        Encontrar pares de eventos correlacionados temporalmente.
        """
        max_delta = max_delta_seconds or self.threshold
        
        events_a = [e for e in self.events if e.event_type == event_type_a]
        events_b = [e for e in self.events if e.event_type == event_type_b]
        
        correlated = []
        
        for ea in events_a:
            for eb in events_b:
                delta = abs((eb.timestamp - ea.timestamp).total_seconds())
                if delta <= max_delta:
                    correlated.append((ea, eb, delta))
        
        # Ordenar por delta (más correlacionados primero)
        correlated.sort(key=lambda x: x[2])
        
        return correlated
    
    def calculate_correlation_score(
        self,
        events_a: List[Event],
        events_b: List[Event],
        window_seconds: float = 300
    ) -> float:
        """
        Calcular score de correlación entre dos series de eventos.
        Score alto = probable vinculación.
        """
        if not events_a or not events_b:
            return 0.0
        
        # Contar co-ocurrencias en ventanas de tiempo
        co_occurrences = 0
        
        for ea in events_a:
            for eb in events_b:
                delta = abs((eb.timestamp - ea.timestamp).total_seconds())
                if delta <= window_seconds:
                    # Peso inversamente proporcional al delta
                    weight = 1.0 - (delta / window_seconds)
                    co_occurrences += weight
        
        # Normalizar por número de eventos
        max_possible = min(len(events_a), len(events_b))
        if max_possible == 0:
            return 0.0
        
        score = co_occurrences / max_possible
        return min(1.0, score)
    
    def detect_behavioral_patterns(
        self,
        user_events: List[Event],
        min_occurrences: int = 3
    ) -> Dict:
        """
        Detectar patrones de comportamiento temporal.
        """
        patterns = {
            'hourly': defaultdict(int),
            'daily': defaultdict(int),
            'intervals': []
        }
        
        # Distribución por hora del día
        for event in user_events:
            hour = event.timestamp.hour
            patterns['hourly'][hour] += 1
        
        # Distribución por día de semana
        for event in user_events:
            day = event.timestamp.strftime('%A')
            patterns['daily'][day] += 1
        
        # Intervalos entre eventos consecutivos
        sorted_events = sorted(user_events, key=lambda x: x.timestamp)
        for i in range(1, len(sorted_events)):
            interval = (sorted_events[i].timestamp - sorted_events[i-1].timestamp).total_seconds()
            patterns['intervals'].append(interval)
        
        # Detectar patrones significativos
        findings = []
        
        # Hora más común
        if patterns['hourly']:
            peak_hour = max(patterns['hourly'], key=patterns['hourly'].get)
            if patterns['hourly'][peak_hour] >= min_occurrences:
                findings.append(f"Peak activity at hour {peak_hour}")
        
        # Día más común
        if patterns['daily']:
            peak_day = max(patterns['daily'], key=patterns['daily'].get)
            if patterns['daily'][peak_day] >= min_occurrences:
                findings.append(f"Most active on {peak_day}")
        
        # Intervalo regular
        if patterns['intervals']:
            avg_interval = sum(patterns['intervals']) / len(patterns['intervals'])
            variance = sum((x - avg_interval)**2 for x in patterns['intervals']) / len(patterns['intervals'])
            if variance < avg_interval * 0.1:  # Baja varianza = patrón regular
                findings.append(f"Regular interval pattern: ~{avg_interval:.0f}s")
        
        return {
            'hourly_distribution': dict(patterns['hourly']),
            'daily_distribution': dict(patterns['daily']),
            'avg_interval': sum(patterns['intervals']) / len(patterns['intervals']) if patterns['intervals'] else 0,
            'findings': findings
        }


class TimingDefense:
    """
    Técnicas de defensa contra timing correlation attacks.
    """
    
    def __init__(self, min_delay: float = 60, max_delay: float = 3600):
        self.min_delay = min_delay
        self.max_delay = max_delay
        self.pending_actions: List[Dict] = []
    
    def schedule_with_random_delay(
        self,
        action: callable,
        delay_distribution: str = 'exponential'
    ) -> datetime:
        """
        Programar acción con delay aleatorio.
        
        Args:
            action: Función a ejecutar
            delay_distribution: 'uniform', 'exponential', 'poisson'
        
        Returns:
            Timestamp programado
        """
        if delay_distribution == 'uniform':
            delay = random.uniform(self.min_delay, self.max_delay)
            
        elif delay_distribution == 'exponential':
            # Lambda tal que media = (min + max) / 2
            mean = (self.min_delay + self.max_delay) / 2
            delay = random.expovariate(1 / mean)
            delay = max(self.min_delay, min(self.max_delay, delay))
            
        elif delay_distribution == 'poisson':
            # Similar a exponencial pero discretizado
            mean = (self.min_delay + self.max_delay) / 2
            delay = random.expovariate(1 / mean)
            delay = max(self.min_delay, min(self.max_delay, delay))
        
        else:
            delay = random.uniform(self.min_delay, self.max_delay)
        
        scheduled_time = datetime.now() + timedelta(seconds=delay)
        
        self.pending_actions.append({
            'action': action,
            'scheduled': scheduled_time,
            'delay': delay
        })
        
        return scheduled_time
    
    def batch_actions(
        self,
        actions: List[callable],
        batch_window: float = 3600
    ) -> List[datetime]:
        """
        Agrupar múltiples acciones en una ventana de tiempo.
        Todas se ejecutan en momentos aleatorios dentro de la ventana.
        """
        window_start = datetime.now() + timedelta(seconds=self.min_delay)
        scheduled_times = []
        
        for action in actions:
            # Tiempo aleatorio dentro de la ventana
            offset = random.uniform(0, batch_window)
            scheduled = window_start + timedelta(seconds=offset)
            
            self.pending_actions.append({
                'action': action,
                'scheduled': scheduled,
                'batch_id': id(actions)
            })
            scheduled_times.append(scheduled)
        
        return scheduled_times
    
    def generate_cover_activity(
        self,
        real_action: callable,
        num_decoys: int = 5,
        time_window: float = 3600
    ) -> List[Dict]:
        """
        Generar actividad de cobertura junto con acción real.
        El adversario no puede distinguir cuál es la real.
        """
        activities = []
        
        # Programar acción real
        real_time = self.schedule_with_random_delay(real_action, 'exponential')
        activities.append({
            'type': 'real',
            'scheduled': real_time
        })
        
        # Generar decoys
        window_start = datetime.now()
        for i in range(num_decoys):
            offset = random.uniform(0, time_window)
            decoy_time = window_start + timedelta(seconds=offset)
            
            # Decoy action (no-op o actividad falsa)
            def decoy_action():
                pass  # O enviar transacción falsa a sí mismo
            
            activities.append({
                'type': 'decoy',
                'scheduled': decoy_time
            })
        
        # Mezclar para que no se pueda distinguir por orden
        random.shuffle(activities)
        
        return activities
    
    def calculate_anonymity_improvement(
        self,
        base_anonymity_set: int,
        delay_seconds: float,
        events_per_hour: float
    ) -> Dict:
        """
        Calcular mejora en anonymity set por delay.
        """
        # Eventos adicionales en el tiempo de delay
        hours = delay_seconds / 3600
        additional_events = events_per_hour * hours
        
        new_anonymity_set = base_anonymity_set + additional_events
        improvement = (new_anonymity_set / base_anonymity_set - 1) * 100
        
        return {
            'base_anonymity_set': base_anonymity_set,
            'delay_hours': hours,
            'additional_events': additional_events,
            'new_anonymity_set': new_anonymity_set,
            'improvement_percent': improvement
        }


def timing_correlation_demo():
    print("=== Timing Correlation Demo ===\n")
    
    analyzer = TimingCorrelationAnalyzer(correlation_threshold_seconds=60)
    
    # Simular eventos de usuario (mixer + exchange)
    base_time = datetime.now()
    
    # Usuario deposita en mixer, luego va al exchange (mal patrón)
    events = [
        Event("d1", "mixer_deposit", base_time, {}),
        Event("w1", "mixer_withdraw", base_time + timedelta(seconds=30), {}),
        Event("l1", "exchange_login", base_time + timedelta(seconds=35), {}),
        
        Event("d2", "mixer_deposit", base_time + timedelta(hours=2), {}),
        Event("w2", "mixer_withdraw", base_time + timedelta(hours=2, seconds=45), {}),
        Event("l2", "exchange_login", base_time + timedelta(hours=2, seconds=50), {}),
    ]
    
    for e in events:
        analyzer.add_event(e)
    
    # Detectar correlaciones
    print("--- Correlation Analysis ---")
    correlated = analyzer.find_correlated_pairs(
        "mixer_withdraw", "exchange_login", max_delta_seconds=30
    )
    
    print(f"Found {len(correlated)} correlated pairs:")
    for ea, eb, delta in correlated:
        print(f"  {ea.event_id} <-> {eb.event_id}: Δt = {delta:.1f}s")
    
    # Score de correlación
    withdraws = [e for e in events if e.event_type == "mixer_withdraw"]
    logins = [e for e in events if e.event_type == "exchange_login"]
    
    score = analyzer.calculate_correlation_score(withdraws, logins)
    print(f"\nCorrelation score: {score:.2f}")
    if score > 0.5:
        print("  WARNING: High correlation suggests same user!")
    
    # Defensa
    print("\n--- Timing Defense ---")
    defense = TimingDefense(min_delay=300, max_delay=3600)
    
    # Programar retiro con delay
    scheduled = defense.schedule_with_random_delay(
        lambda: print("Executing withdraw"),
        delay_distribution='exponential'
    )
    print(f"Withdraw scheduled at: {scheduled}")
    print(f"Delay: {(scheduled - datetime.now()).seconds}s")
    
    # Calcular mejora en anonimato
    improvement = defense.calculate_anonymity_improvement(
        base_anonymity_set=100,
        delay_seconds=1800,  # 30 min
        events_per_hour=20
    )
    print(f"\nAnonymity improvement with 30min delay:")
    print(f"  Base set: {improvement['base_anonymity_set']}")
    print(f"  New set: {improvement['new_anonymity_set']:.0f}")
    print(f"  Improvement: {improvement['improvement_percent']:.1f}%")

timing_correlation_demo()
\end{lstlisting}

\subsection{Defensas contra Timing Correlation}

\begin{table}[htbp]
\centering
\small
\begin{tabular}{lp{4cm}cc}
\toprule
\textbf{Defensa} & \textbf{Descripción} & \textbf{Efectividad} & \textbf{Usabilidad} \\
\midrule
Random delay & Esperar tiempo aleatorio antes de actuar & Alta & Media \\
Batching & Agrupar acciones con otras & Alta & Baja \\
Cover traffic & Generar actividad falsa de cobertura & Muy alta & Baja \\
Time-locked actions & Acciones que solo se ejecutan en futuro & Alta & Baja \\
Poisson timing & Intervalos con distribución Poisson & Alta & Alta \\
\bottomrule
\end{tabular}
\caption{Defensas contra timing correlation}
\end{table}
%===========================================
\section{Traffic Analysis}
\label{sec:traffic-analysis}
%===========================================

\subsection{Fundamentos del Análisis de Tráfico}

\begin{definition}[Traffic Analysis]
Técnica de inteligencia que infiere información a partir de los patrones del tráfico de red, sin necesidad de descifrar el contenido:

\textbf{Observables:}
\begin{itemize}
    \item Tamaño de paquetes
    \item Timing entre paquetes
    \item Frecuencia de comunicación
    \item Direcciones de origen/destino
    \item Volumen de datos transferidos
\end{itemize}

\textbf{Principio:} Incluso el tráfico encriptado ``habla'' a través de sus metadatos.
\end{definition}

\begin{mathnote}[Modelo de Traffic Analysis]
\textbf{Flujo de tráfico como señal:}

Un flujo de tráfico puede modelarse como una serie temporal:
\[
F = \{(t_1, s_1), (t_2, s_2), ..., (t_n, s_n)\}
\]

donde $t_i$ es el timestamp y $s_i$ es el tamaño del paquete $i$.

\textbf{Features para análisis:}
\begin{itemize}
    \item \textbf{Inter-arrival time:} $\Delta t_i = t_{i+1} - t_i$
    \item \textbf{Distribución de tamaños:} $\text{hist}(s_1, ..., s_n)$
    \item \textbf{Burstiness:} Varianza en $\Delta t$
    \item \textbf{Total volume:} $V = \sum_{i} s_i$
    \item \textbf{Direction ratio:} Bytes enviados vs recibidos
\end{itemize}

\textbf{Fingerprinting:} Cada tipo de actividad genera un ``fingerprint'' distintivo:
\begin{itemize}
    \item Video streaming: Alto volumen, paquetes regulares grandes
    \item Chat: Bajo volumen, paquetes pequeños, bidireccional
    \item Transacción blockchain: Patrón específico de tamaños
\end{itemize}
\end{mathnote}

\begin{figure}[htbp]
\centering
\begin{tikzpicture}[
    node distance=0.5cm
]
    % Ejes
    \draw[->] (0, 0) -- (10, 0) node[right] {tiempo};
    \draw[->] (0, 0) -- (0, 4) node[above] {tamaño};
    
    % Patrón 1: Video streaming
    \foreach \x in {0.5, 1.0, 1.5, 2.0, 2.5} {
        \draw[fill=blue!60] (\x, 0) rectangle (\x+0.3, 3);
    }
    \node at (1.5, -0.5) [font=\small] {Video};
    
    % Patrón 2: Chat
    \foreach \x in {3.5, 4.2, 5.1, 5.5, 6.3} {
        \draw[fill=green!60] (\x, 0) rectangle (\x+0.1, 0.5);
    }
    \node at (5, -0.5) [font=\small] {Chat};
    
    % Patrón 3: TX Blockchain
    \draw[fill=red!60] (7.5, 0) rectangle (7.7, 1.5);
    \draw[fill=red!60] (8.0, 0) rectangle (8.1, 0.3);
    \draw[fill=red!60] (8.3, 0) rectangle (8.4, 0.3);
    \node at (8, -0.5) [font=\small] {Blockchain TX};
    
    % Leyenda
    \node at (5, 4.5) {\textbf{Traffic Fingerprints}};
\end{tikzpicture}
\caption{Diferentes actividades generan patrones de tráfico distintivos}
\label{fig:traffic-fingerprints}
\end{figure}

\subsection{Ataques de Traffic Analysis}

\begin{lstlisting}[language=Python, caption={Sistema de Traffic Analysis}]
from dataclasses import dataclass
from datetime import datetime, timedelta
from typing import List, Dict, Tuple, Optional
import numpy as np
from collections import defaultdict
import statistics

@dataclass
class Packet:
    timestamp: datetime
    size: int
    direction: str  # 'sent' or 'received'
    encrypted: bool

@dataclass
class TrafficFlow:
    flow_id: str
    source_ip: str
    dest_ip: str
    packets: List[Packet]
    start_time: datetime
    end_time: datetime

class TrafficAnalyzer:
    """
    Analizador de patrones de tráfico para inferir actividad.
    """
    
    def __init__(self):
        # Fingerprints conocidos de diferentes actividades
        self.known_fingerprints = {
            'blockchain_tx': {
                'size_range': (200, 2000),
                'packet_count': (3, 10),
                'bidirectional': True,
                'duration_seconds': (1, 30)
            },
            'mixer_deposit': {
                'size_range': (500, 3000),
                'packet_count': (5, 15),
                'bidirectional': True,
                'duration_seconds': (5, 60)
            },
            'rpc_query': {
                'size_range': (100, 500),
                'packet_count': (2, 4),
                'bidirectional': True,
                'duration_seconds': (0.1, 2)
            },
            'p2p_sync': {
                'size_range': (1000, 65000),
                'packet_count': (100, 10000),
                'bidirectional': True,
                'duration_seconds': (60, 3600)
            }
        }
    
    def extract_features(self, flow: TrafficFlow) -> Dict:
        """
        Extraer características de un flujo de tráfico.
        """
        if not flow.packets:
            return {}
        
        packets = sorted(flow.packets, key=lambda p: p.timestamp)
        sizes = [p.size for p in packets]
        
        # Inter-arrival times
        iats = []
        for i in range(1, len(packets)):
            iat = (packets[i].timestamp - packets[i-1].timestamp).total_seconds()
            iats.append(iat)
        
        # Dirección de paquetes
        sent = [p for p in packets if p.direction == 'sent']
        received = [p for p in packets if p.direction == 'received']
        
        features = {
            # Estadísticas de tamaño
            'size_mean': statistics.mean(sizes) if sizes else 0,
            'size_std': statistics.stdev(sizes) if len(sizes) > 1 else 0,
            'size_min': min(sizes) if sizes else 0,
            'size_max': max(sizes) if sizes else 0,
            
            # Estadísticas de timing
            'iat_mean': statistics.mean(iats) if iats else 0,
            'iat_std': statistics.stdev(iats) if len(iats) > 1 else 0,
            
            # Conteos
            'packet_count': len(packets),
            'sent_count': len(sent),
            'received_count': len(received),
            
            # Volumen
            'total_bytes': sum(sizes),
            'sent_bytes': sum(p.size for p in sent),
            'received_bytes': sum(p.size for p in received),
            
            # Duración
            'duration': (flow.end_time - flow.start_time).total_seconds(),
            
            # Ratios
            'direction_ratio': len(sent) / len(received) if received else float('inf'),
            'bytes_ratio': sum(p.size for p in sent) / sum(p.size for p in received) if received else float('inf')
        }
        
        return features
    
    def classify_activity(self, flow: TrafficFlow) -> Tuple[str, float]:
        """
        Clasificar tipo de actividad basándose en fingerprint.
        """
        features = self.extract_features(flow)
        
        if not features:
            return ('unknown', 0.0)
        
        best_match = 'unknown'
        best_score = 0.0
        
        for activity, fingerprint in self.known_fingerprints.items():
            score = self._match_fingerprint(features, fingerprint)
            if score > best_score:
                best_score = score
                best_match = activity
        
        return (best_match, best_score)
    
    def _match_fingerprint(self, features: Dict, fingerprint: Dict) -> float:
        """
        Calcular score de coincidencia con fingerprint.
        """
        score = 0.0
        checks = 0
        
        # Verificar rango de tamaño
        if 'size_range' in fingerprint:
            min_size, max_size = fingerprint['size_range']
            if min_size <= features.get('size_mean', 0) <= max_size:
                score += 1
            checks += 1
        
        # Verificar conteo de paquetes
        if 'packet_count' in fingerprint:
            min_count, max_count = fingerprint['packet_count']
            if min_count <= features.get('packet_count', 0) <= max_count:
                score += 1
            checks += 1
        
        # Verificar duración
        if 'duration_seconds' in fingerprint:
            min_dur, max_dur = fingerprint['duration_seconds']
            if min_dur <= features.get('duration', 0) <= max_dur:
                score += 1
            checks += 1
        
        # Verificar bidireccionalidad
        if fingerprint.get('bidirectional', False):
            if features.get('sent_count', 0) > 0 and features.get('received_count', 0) > 0:
                score += 1
            checks += 1
        
        return score / checks if checks > 0 else 0.0
    
    def detect_correlation(
        self,
        flow_a: TrafficFlow,
        flow_b: TrafficFlow,
        time_threshold: float = 60
    ) -> Dict:
        """
        Detectar correlación entre dos flujos de tráfico.
        """
        # Correlación temporal
        time_diff = abs((flow_a.start_time - flow_b.start_time).total_seconds())
        temporal_correlation = max(0, 1 - time_diff / time_threshold)
        
        # Correlación de volumen
        vol_a = sum(p.size for p in flow_a.packets)
        vol_b = sum(p.size for p in flow_b.packets)
        if vol_a > 0 and vol_b > 0:
            volume_ratio = min(vol_a, vol_b) / max(vol_a, vol_b)
        else:
            volume_ratio = 0
        
        # Correlación de duración
        dur_a = (flow_a.end_time - flow_a.start_time).total_seconds()
        dur_b = (flow_b.end_time - flow_b.start_time).total_seconds()
        if dur_a > 0 and dur_b > 0:
            duration_ratio = min(dur_a, dur_b) / max(dur_a, dur_b)
        else:
            duration_ratio = 0
        
        # Score combinado
        overall_correlation = (
            temporal_correlation * 0.5 +
            volume_ratio * 0.25 +
            duration_ratio * 0.25
        )
        
        return {
            'temporal_correlation': temporal_correlation,
            'volume_ratio': volume_ratio,
            'duration_ratio': duration_ratio,
            'overall_score': overall_correlation,
            'likely_related': overall_correlation > 0.7
        }
    
    def website_fingerprinting_attack(
        self,
        target_flow: TrafficFlow,
        known_websites: Dict[str, List[TrafficFlow]]
    ) -> List[Tuple[str, float]]:
        """
        Intento de identificar qué sitio web visitó el usuario
        basándose en patrones de tráfico (Website Fingerprinting).
        """
        target_features = self.extract_features(target_flow)
        
        matches = []
        
        for website, sample_flows in known_websites.items():
            # Calcular features promedio del sitio
            all_features = [self.extract_features(f) for f in sample_flows]
            
            # Calcular distancia al target
            similarity = self._calculate_similarity(target_features, all_features)
            matches.append((website, similarity))
        
        # Ordenar por similitud
        matches.sort(key=lambda x: x[1], reverse=True)
        return matches
    
    def _calculate_similarity(
        self,
        target: Dict,
        samples: List[Dict]
    ) -> float:
        """
        Calcular similitud entre target y muestras.
        """
        if not samples:
            return 0.0
        
        # Calcular media de cada feature en samples
        mean_features = {}
        for key in target.keys():
            values = [s.get(key, 0) for s in samples if isinstance(s.get(key, 0), (int, float))]
            if values:
                mean_features[key] = statistics.mean(values)
        
        # Calcular distancia euclidiana normalizada
        distance = 0
        count = 0
        
        for key, target_val in target.items():
            if key in mean_features and isinstance(target_val, (int, float)):
                mean_val = mean_features[key]
                if mean_val > 0:
                    normalized_diff = abs(target_val - mean_val) / mean_val
                    distance += normalized_diff
                    count += 1
        
        if count == 0:
            return 0.0
        
        avg_distance = distance / count
        similarity = max(0, 1 - avg_distance)
        
        return similarity


def traffic_analysis_demo():
    print("=== Traffic Analysis Demo ===\n")
    
    analyzer = TrafficAnalyzer()
    
    # Simular flujo de transacción blockchain
    base_time = datetime.now()
    
    blockchain_flow = TrafficFlow(
        flow_id="flow_1",
        source_ip="192.168.1.100",
        dest_ip="13.56.163.26",  # Infura
        packets=[
            Packet(base_time, 450, 'sent', True),
            Packet(base_time + timedelta(milliseconds=50), 200, 'received', True),
            Packet(base_time + timedelta(milliseconds=100), 800, 'sent', True),
            Packet(base_time + timedelta(milliseconds=200), 150, 'received', True),
            Packet(base_time + timedelta(milliseconds=300), 100, 'received', True),
        ],
        start_time=base_time,
        end_time=base_time + timedelta(milliseconds=300)
    )
    
    # Extraer features
    print("--- Feature Extraction ---")
    features = analyzer.extract_features(blockchain_flow)
    print(f"Packet count: {features['packet_count']}")
    print(f"Mean size: {features['size_mean']:.0f} bytes")
    print(f"Duration: {features['duration']:.3f}s")
    print(f"Direction ratio: {features['direction_ratio']:.2f}")
    
    # Clasificar actividad
    print("\n--- Activity Classification ---")
    activity, confidence = analyzer.classify_activity(blockchain_flow)
    print(f"Classified as: {activity}")
    print(f"Confidence: {confidence:.2f}")
    
    # Simular flujo de mixer
    mixer_flow = TrafficFlow(
        flow_id="flow_2",
        source_ip="192.168.1.100",
        dest_ip="tornado.cash.onion",
        packets=[
            Packet(base_time + timedelta(seconds=5), 1500, 'sent', True),
            Packet(base_time + timedelta(seconds=5.1), 500, 'received', True),
            Packet(base_time + timedelta(seconds=5.5), 2000, 'sent', True),
            Packet(base_time + timedelta(seconds=6), 800, 'received', True),
            Packet(base_time + timedelta(seconds=6.5), 1200, 'sent', True),
            Packet(base_time + timedelta(seconds=7), 300, 'received', True),
            Packet(base_time + timedelta(seconds=8), 200, 'received', True),
        ],
        start_time=base_time + timedelta(seconds=5),
        end_time=base_time + timedelta(seconds=8)
    )
    
    activity2, confidence2 = analyzer.classify_activity(mixer_flow)
    print(f"\nSecond flow classified as: {activity2} ({confidence2:.2f})")
    
    # Detectar correlación
    print("\n--- Correlation Detection ---")
    correlation = analyzer.detect_correlation(blockchain_flow, mixer_flow)
    print(f"Temporal correlation: {correlation['temporal_correlation']:.2f}")
    print(f"Volume ratio: {correlation['volume_ratio']:.2f}")
    print(f"Likely related: {correlation['likely_related']}")

traffic_analysis_demo()
\end{lstlisting}

%===========================================
\section{Mixnets vs Onion Routing}
\label{sec:mixnets-vs-onion}
%===========================================

\subsection{Onion Routing (Tor)}

\begin{definition}[Onion Routing]
Técnica de comunicación anónima donde el mensaje se encripta en múltiples capas, cada una descifrada por un nodo diferente en la ruta:

\textbf{Características:}
\begin{itemize}
    \item Circuito establecido antes de enviar datos
    \item Baja latencia (segundos)
    \item Paquetes correlacionables por timing
    \item Vulnerable a análisis de tráfico global
\end{itemize}

\textbf{Modelo de seguridad:}
\begin{itemize}
    \item Seguro contra adversario local (observa parte de la red)
    \item Vulnerable a Global Passive Adversary (GPA)
    \item No usa cover traffic
\end{itemize}
\end{definition}

\begin{mathnote}[Vulnerabilidades de Onion Routing]
\textbf{1. End-to-end timing correlation:}

Si un adversario observa tanto la entrada como la salida del circuito:
\[
\text{Corr}(T_{\text{entry}}, T_{\text{exit}}) \approx 1
\]

Los paquetes tienen el mismo patrón temporal, solo retrasados.

\textbf{2. Traffic confirmation attack:}

El adversario controla nodos de entrada y salida:
\begin{enumerate}
    \item Modifica timing en entrada (watermark)
    \item Detecta watermark en salida
    \item Confirma que el mismo usuario
\end{enumerate}

\textbf{3. Website fingerprinting:}

Incluso sin ver contenido, el patrón de tráfico (tamaño, timing) identifica el sitio visitado con alta precisión ($>$ 90\% en condiciones controladas).
\end{mathnote}

\subsection{Mixnets}

\begin{definition}[Mixnet]
Red de comunicación anónima donde los mensajes se mezclan con otros mensajes en cada nodo, rompiendo la correlación temporal:

\textbf{Características:}
\begin{itemize}
    \item Cada nodo acumula mensajes y los reordena
    \item Alta latencia (minutos a horas)
    \item Resiste análisis de timing
    \item Usa cover traffic
\end{itemize}

\textbf{Modelo de seguridad:}
\begin{itemize}
    \item Seguro incluso contra GPA (con suficiente tráfico)
    \item Requiere confianza en que no todos los nodos coludan
    \item Cover traffic oculta patrones de actividad
\end{itemize}
\end{definition}

\begin{figure}[htbp]
\centering
\begin{tikzpicture}[
    node distance=2cm,
    msg/.style={circle, draw, minimum size=0.4cm},
    mixnode/.style={rectangle, draw, minimum width=1.5cm, minimum height=2cm, fill=blue!20},
    arrow/.style={->, thick}
]
    % Onion Routing (arriba)
    \node at (-1, 4) {\textbf{Onion Routing}};
    
    \node[msg, fill=red] (o1) at (0, 3.5) {};
    \node[msg, fill=green] (o2) at (0, 3) {};
    \node[msg, fill=blue] (o3) at (0, 2.5) {};
    
    \node[mixnode, minimum height=1cm] (on1) at (2.5, 3) {Relay};
    \node[mixnode, minimum height=1cm] (on2) at (5, 3) {Relay};
    \node[mixnode, minimum height=1cm] (on3) at (7.5, 3) {Relay};
    
    \node[msg, fill=red] (o1e) at (10, 3.5) {};
    \node[msg, fill=green] (o2e) at (10, 3) {};
    \node[msg, fill=blue] (o3e) at (10, 2.5) {};
    
    \draw[arrow, red] (o1) -- (on1);
    \draw[arrow, green] (o2) -- (on1);
    \draw[arrow, blue] (o3) -- (on1);
    
    \draw[arrow] (on1) -- (on2);
    \draw[arrow] (on2) -- (on3);
    
    \draw[arrow, red] (on3) -- (o1e);
    \draw[arrow, green] (on3) -- (o2e);
    \draw[arrow, blue] (on3) -- (o3e);
    
    \node at (5, 1.5) [font=\small] {Orden preservado $\rightarrow$ Timing correlation};
    
    % Mixnet (abajo)
    \node at (-1, 0) {\textbf{Mixnet}};
    
    \node[msg, fill=red] (m1) at (0, -0.5) {};
    \node[msg, fill=green] (m2) at (0, -1) {};
    \node[msg, fill=blue] (m3) at (0, -1.5) {};
    
    \node[mixnode] (mn1) at (2.5, -1) {Mix};
    \node[mixnode] (mn2) at (5, -1) {Mix};
    \node[mixnode] (mn3) at (7.5, -1) {Mix};
    
    \node[msg, fill=green] (m1e) at (10, -0.5) {};
    \node[msg, fill=blue] (m2e) at (10, -1) {};
    \node[msg, fill=red] (m3e) at (10, -1.5) {};
    
    \draw[arrow] (m1) -- (mn1);
    \draw[arrow] (m2) -- (mn1);
    \draw[arrow] (m3) -- (mn1);
    
    \draw[arrow] (mn1) -- (mn2);
    \draw[arrow] (mn2) -- (mn3);
    
    \draw[arrow, green] (mn3) -- (m1e);
    \draw[arrow, blue] (mn3) -- (m2e);
    \draw[arrow, red] (mn3) -- (m3e);
    
    \node at (5, -2.5) [font=\small] {Orden mezclado + delay $\rightarrow$ No correlation};
\end{tikzpicture}
\caption{Onion routing vs Mixnet: el mixing rompe la correlación temporal}
\label{fig:onion-vs-mixnet}
\end{figure}

\subsection{Comparación Detallada}

\begin{table}[htbp]
\centering
\small
\begin{tabular}{lp{4cm}p{4cm}}
\toprule
\textbf{Aspecto} & \textbf{Onion Routing (Tor)} & \textbf{Mixnet (Nym)} \\
\midrule
\textbf{Latencia} & Baja (100ms - 1s) & Alta (segundos a minutos) \\
\textbf{Modelo de mensaje} & Stream-based (circuitos) & Packet-based (mensajes individuales) \\
\textbf{Timing resistance} & Parcial (vulnerable a GPA) & Alta (mixing + delays) \\
\textbf{Cover traffic} & No & Sí (parte del protocolo) \\
\textbf{Casos de uso} & Browsing web, streaming & Messaging, transacciones \\
\textbf{Escalabilidad} & Alta (millones de usuarios) & Media (requiere más overhead) \\
\textbf{Incentivos} & Voluntarios & Económicos (tokens) \\
\textbf{Madurez} & Alta (20+ años) & Emergente \\
\bottomrule
\end{tabular}
\caption{Comparación entre Onion Routing y Mixnets}
\end{table}

\subsection{Nym Mixnet}

\begin{definition}[Nym]
Mixnet moderna con incentivos económicos y cover traffic integrado:

\textbf{Arquitectura:}
\begin{enumerate}
    \item \textbf{Mix nodes:} 3 capas de mixing
    \item \textbf{Gateways:} Puntos de entrada/salida
    \item \textbf{Validators:} Red de consenso para incentivos
\end{enumerate}

\textbf{Características únicas:}
\begin{itemize}
    \item Sphinx packet format (tamaño uniforme)
    \item Poisson mixing (delays aleatorios)
    \item Loop cover traffic (mensajes a sí mismo)
    \item NYM token para pago e incentivos
\end{itemize}
\end{definition}

\begin{lstlisting}[language=Python, caption={Simulación conceptual de Mixnet}]
from dataclasses import dataclass
from datetime import datetime, timedelta
from typing import List, Dict, Optional, Tuple
import random
import secrets
from collections import deque
import time

@dataclass
class SphinxPacket:
    """
    Paquete Sphinx con tamaño uniforme.
    """
    packet_id: bytes
    encrypted_payload: bytes
    routing_info: bytes  # Encriptado para cada hop
    size: int = 2048  # Tamaño fijo
    
    def is_cover_traffic(self) -> bool:
        """Verificar si es cover traffic (payload vacío)."""
        return self.encrypted_payload == b'\x00' * len(self.encrypted_payload)

@dataclass
class MixNode:
    """
    Nodo de mixing en la mixnet.
    """
    node_id: str
    layer: int
    public_key: bytes
    
    # Pool de mensajes para mixing
    message_pool: deque
    pool_size: int = 10
    
    # Estadísticas
    messages_processed: int = 0
    
    def __post_init__(self):
        self.message_pool = deque(maxlen=self.pool_size)
    
    def receive_packet(self, packet: SphinxPacket) -> Optional[SphinxPacket]:
        """
        Recibir paquete, añadir al pool, potencialmente emitir.
        """
        # Descifrar capa correspondiente a este nodo
        decrypted = self._decrypt_layer(packet)
        
        # Añadir al pool
        self.message_pool.append(decrypted)
        self.messages_processed += 1
        
        # Emitir si pool está lleno (mixing)
        if len(self.message_pool) >= self.pool_size:
            return self._emit_random()
        
        # O emitir con probabilidad (Poisson)
        if random.random() < 0.3:  # Lambda del proceso Poisson
            return self._emit_random()
        
        return None
    
    def _decrypt_layer(self, packet: SphinxPacket) -> SphinxPacket:
        """Descifrar capa de routing para este nodo."""
        # En producción: descifrado real con clave privada
        return SphinxPacket(
            packet_id=packet.packet_id,
            encrypted_payload=packet.encrypted_payload,
            routing_info=packet.routing_info[32:],  # Remover info de este hop
            size=packet.size
        )
    
    def _emit_random(self) -> Optional[SphinxPacket]:
        """Emitir paquete aleatorio del pool."""
        if not self.message_pool:
            return None
        
        # Selección aleatoria (no FIFO!)
        idx = random.randint(0, len(self.message_pool) - 1)
        packets = list(self.message_pool)
        selected = packets[idx]
        
        # Remover del pool
        self.message_pool = deque(
            [p for i, p in enumerate(packets) if i != idx],
            maxlen=self.pool_size
        )
        
        return selected


class MixnetSimulator:
    """
    Simulador de mixnet con múltiples capas.
    """
    
    def __init__(self, num_layers: int = 3, nodes_per_layer: int = 5):
        self.num_layers = num_layers
        self.nodes: Dict[int, List[MixNode]] = {}
        
        # Crear nodos por capa
        for layer in range(num_layers):
            self.nodes[layer] = []
            for i in range(nodes_per_layer):
                node = MixNode(
                    node_id=f"node_{layer}_{i}",
                    layer=layer,
                    public_key=secrets.token_bytes(32),
                    pool_size=5
                )
                self.nodes[layer].append(node)
    
    def create_packet(
        self,
        payload: bytes,
        destination: str,
        is_cover: bool = False
    ) -> Tuple[SphinxPacket, List[MixNode]]:
        """
        Crear paquete Sphinx con ruta aleatoria.
        """
        # Seleccionar ruta (un nodo por capa)
        route = []
        for layer in range(self.num_layers):
            node = random.choice(self.nodes[layer])
            route.append(node)
        
        # Construir routing info (encriptado para cada hop)
        # En producción: ECIES o similar
        routing_info = b''.join(node.public_key for node in route)
        
        # Payload (uniforme, padded)
        if is_cover:
            encrypted_payload = b'\x00' * 1024
        else:
            padded = payload.ljust(1024, b'\x00')
            encrypted_payload = padded  # En producción: encriptar
        
        packet = SphinxPacket(
            packet_id=secrets.token_bytes(16),
            encrypted_payload=encrypted_payload,
            routing_info=routing_info
        )
        
        return packet, route
    
    def send_packet(
        self,
        packet: SphinxPacket,
        route: List[MixNode],
        add_delay: bool = True
    ) -> Dict:
        """
        Enviar paquete a través de la mixnet.
        """
        current_packet = packet
        delays = []
        
        for i, node in enumerate(route):
            # Delay de Poisson antes de procesar
            if add_delay:
                delay = random.expovariate(1.0)  # media = 1 segundo
                delays.append(delay)
            
            # Procesar en nodo
            output = node.receive_packet(current_packet)
            
            if output:
                current_packet = output
        
        return {
            'packet_id': packet.packet_id.hex()[:16],
            'route_length': len(route),
            'total_delay': sum(delays),
            'nodes_visited': [n.node_id for n in route]
        }
    
    def generate_cover_traffic(self, num_packets: int = 10) -> List[Dict]:
        """
        Generar tráfico de cobertura.
        """
        results = []
        
        for _ in range(num_packets):
            packet, route = self.create_packet(
                payload=b'',
                destination='self',
                is_cover=True
            )
            result = self.send_packet(packet, route)
            result['is_cover'] = True
            results.append(result)
        
        return results
    
    def analyze_anonymity(
        self,
        observed_input: SphinxPacket,
        observed_outputs: List[SphinxPacket]
    ) -> Dict:
        """
        Analizar qué tan difícil es correlacionar input con output.
        """
        # Sin timing correlation
        # Todos los outputs son igualmente probables
        anonymity_set = len(observed_outputs)
        
        # Entropía
        entropy = 0
        if anonymity_set > 0:
            import math
            prob = 1 / anonymity_set
            entropy = -anonymity_set * prob * math.log2(prob)
        
        return {
            'anonymity_set_size': anonymity_set,
            'entropy_bits': entropy,
            'probability_correct_guess': 1 / anonymity_set if anonymity_set > 0 else 1
        }


def mixnet_demo():
    print("=== Mixnet Simulation Demo ===\n")
    
    # Crear mixnet
    mixnet = MixnetSimulator(num_layers=3, nodes_per_layer=5)
    
    print(f"Created mixnet with {mixnet.num_layers} layers")
    print(f"Nodes per layer: {len(mixnet.nodes[0])}")
    
    # Enviar mensaje real
    print("\n--- Sending Real Message ---")
    
    packet, route = mixnet.create_packet(
        payload=b"Secret transaction data",
        destination="0xRecipient",
        is_cover=False
    )
    
    result = mixnet.send_packet(packet, route)
    print(f"Packet sent through: {result['nodes_visited']}")
    print(f"Total delay: {result['total_delay']:.2f}s")
    
    # Generar cover traffic
    print("\n--- Generating Cover Traffic ---")
    
    cover_results = mixnet.generate_cover_traffic(num_packets=10)
    print(f"Generated {len(cover_results)} cover packets")
    
    avg_delay = sum(r['total_delay'] for r in cover_results) / len(cover_results)
    print(f"Average cover delay: {avg_delay:.2f}s")
    
    # Análisis de anonimato
    print("\n--- Anonymity Analysis ---")
    
    # Simular observador que ve entrada y 10 salidas posibles
    anonymity = mixnet.analyze_anonymity(
        observed_input=packet,
        observed_outputs=[packet] + [p for p, _ in [mixnet.create_packet(b'', '', True) for _ in range(9)]]
    )
    
    print(f"Anonymity set size: {anonymity['anonymity_set_size']}")
    print(f"Entropy: {anonymity['entropy_bits']:.2f} bits")
    print(f"Probability of correct guess: {anonymity['probability_correct_guess']:.2%}")

mixnet_demo()
\end{lstlisting}

%===========================================
\section{VPN vs Tor vs I2P vs Nym}
\label{sec:network-comparison}
%===========================================

\subsection{Comparación Comprehensiva}

\begin{table}[htbp]
\centering
\footnotesize
\begin{tabular}{lcccc}
\toprule
\textbf{Característica} & \textbf{VPN} & \textbf{Tor} & \textbf{I2P} & \textbf{Nym} \\
\midrule
\textbf{Modelo de confianza} & Centralizado & Distribuido & Distribuido & Distribuido \\
\textbf{Conoce tu IP} & Proveedor VPN & Guard node & Ninguno & Gateway \\
\textbf{Latencia típica} & 10-50ms & 200-500ms & 500ms-2s & 1-10s \\
\textbf{Cover traffic} & No & No & Limitado & Sí \\
\textbf{Timing resistance} & No & Parcial & Parcial & Alta \\
\textbf{Exit a clearnet} & Sí & Sí & Limitado & Sí \\
\textbf{Servicios ocultos} & No & .onion & .i2p & -- \\
\textbf{Incentivos} & Pago al provider & Voluntarios & Voluntarios & Token NYM \\
\textbf{Uso en blockchain} & Básico & Recomendado & Monero/Kovri & Emergente \\
\textbf{Resistencia a GPA} & No & No & Parcial & Alta \\
\bottomrule
\end{tabular}
\caption{Comparación de soluciones de privacidad de red}
\end{table}

\begin{figure}[htbp]
\centering
\begin{tikzpicture}
    % Ejes
    \draw[->] (0, 0) -- (10, 0) node[right] {Latencia};
    \draw[->] (0, 0) -- (0, 6) node[above] {Anonimato};
    
    % Etiquetas de ejes
    \node at (2, -0.5) [font=\small] {Baja};
    \node at (8, -0.5) [font=\small] {Alta};
    \node at (-0.8, 1.5) [font=\small, rotate=90] {Bajo};
    \node at (-0.8, 4.5) [font=\small, rotate=90] {Alto};
    
    % Puntos
    \node[circle, fill=red, minimum size=0.8cm, label=above:VPN] at (1, 1.5) {};
    \node[circle, fill=orange, minimum size=0.8cm, label=above:Tor] at (3, 3.5) {};
    \node[circle, fill=yellow, minimum size=0.8cm, label=above:I2P] at (5, 4) {};
    \node[circle, fill=green, minimum size=0.8cm, label=above:Nym] at (8, 5) {};
    
    % Trade-off line
    \draw[dashed, gray] (0.5, 0.5) -- (9, 5.5);
    \node at (7, 2) [font=\small, gray, rotate=30] {Trade-off típico};
\end{tikzpicture}
\caption{Trade-off entre latencia y nivel de anonimato}
\label{fig:privacy-tradeoff}
\end{figure}

\begin{keyidea}
\textbf{Recomendaciones según caso de uso:}

\begin{enumerate}
    \item \textbf{Browsing web general:}
    \begin{itemize}
        \item VPN para privacidad básica del ISP
        \item Tor para anonimato real
    \end{itemize}
    
    \item \textbf{Transacciones blockchain:}
    \begin{itemize}
        \item Mínimo: Tor para broadcast de txs
        \item Mejor: Nodo propio + Tor
        \item Óptimo: Mixnet para transacciones sensibles
    \end{itemize}
    
    \item \textbf{Servicios ocultos:}
    \begin{itemize}
        \item Tor .onion para accesibilidad
        \item I2P .i2p para comunidad específica
    \end{itemize}
    
    \item \textbf{Mensajería privada:}
    \begin{itemize}
        \item Nym para máxima privacidad
        \item Tolerable latencia alta
    \end{itemize}
\end{enumerate}
\end{keyidea}
%===========================================
\section{Packet Padding}
\label{sec:packet-padding}
%===========================================

\subsection{Concepto y Motivación}

\begin{definition}[Packet Padding]
Técnica de añadir datos dummy a los paquetes de red para uniformizar su tamaño, eliminando la información que un adversario podría inferir del tamaño variable de los mensajes.

\textbf{Objetivo:} Si todos los paquetes tienen el mismo tamaño, el observador no puede distinguir entre:
\begin{itemize}
    \item Transacción grande vs pequeña
    \item Mensaje de texto vs imagen
    \item Consulta simple vs compleja
\end{itemize}
\end{definition}

\begin{mathnote}[Análisis de Tamaño de Paquetes]
\textbf{Sin padding:}
\begin{itemize}
    \item Transacción simple: $\sim$250 bytes
    \item Transacción con contrato: $\sim$2,000 bytes
    \item Despliegue de contrato: $\sim$10,000+ bytes
\end{itemize}

Un observador puede inferir el tipo de operación solo por el tamaño.

\textbf{Con padding uniforme:}
\[
\forall \text{ paquete } p: |p| = S_{\text{fijo}}
\]

Todos los paquetes tienen tamaño $S_{\text{fijo}}$, eliminando el side-channel de tamaño.

\textbf{Trade-offs:}
\begin{itemize}
    \item \textbf{Overhead de bandwidth:} $\text{Overhead} = \frac{S_{\text{fijo}} - S_{\text{real}}}{S_{\text{real}}} \times 100\%$
    \item \textbf{Fragmentación:} Mensajes grandes requieren múltiples paquetes
    \item \textbf{Elección de $S_{\text{fijo}}$:} Muy pequeño = mucha fragmentación, muy grande = mucho overhead
\end{itemize}
\end{mathnote}

\begin{figure}[htbp]
\centering
\begin{tikzpicture}[
    packet/.style={rectangle, draw, minimum height=0.6cm}
]
    % Sin padding
    \node at (-1, 3) {\textbf{Sin Padding}};
    
    \node[packet, fill=blue!30, minimum width=1cm] at (1, 3) {};
    \node[packet, fill=green!30, minimum width=2cm] at (3.5, 3) {};
    \node[packet, fill=red!30, minimum width=4cm] at (7, 3) {};
    
    \node at (1, 2.3) [font=\small] {250B};
    \node at (3.5, 2.3) [font=\small] {500B};
    \node at (7, 2.3) [font=\small] {1000B};
    
    % Con padding
    \node at (-1, 1) {\textbf{Con Padding}};
    
    \node[packet, minimum width=4cm] at (2, 1) {};
    \node[packet, fill=blue!30, minimum width=1cm] at (0.5, 1) {};
    
    \node[packet, minimum width=4cm] at (6.5, 1) {};
    \node[packet, fill=green!30, minimum width=2cm] at (5.5, 1) {};
    
    \node[packet, minimum width=4cm] at (11, 1) {};
    \node[packet, fill=red!30, minimum width=4cm] at (11, 1) {};
    
    \node at (2, 0.3) [font=\small] {1KB};
    \node at (6.5, 0.3) [font=\small] {1KB};
    \node at (11, 0.3) [font=\small] {1KB};
    
    % Flechas
    \draw[->] (1, 2) -- (2, 1.5);
    \draw[->] (3.5, 2) -- (6.5, 1.5);
    \draw[->] (7, 2) -- (11, 1.5);
    
    % Observador
    \node at (6.5, -1) [draw, rectangle, fill=gray!20] {
        Observador: todos los paquetes idénticos
    };
\end{tikzpicture}
\caption{Packet padding uniformiza el tamaño visible de los paquetes}
\label{fig:packet-padding}
\end{figure}

\begin{lstlisting}[language=Python, caption={Sistema de Packet Padding}]
from dataclasses import dataclass
from typing import List, Tuple, Optional
import secrets
import struct

@dataclass
class PaddedPacket:
    """Paquete con padding uniforme."""
    payload: bytes
    padding: bytes
    total_size: int
    sequence_number: int
    is_final: bool  # Para mensajes fragmentados

class PacketPaddingSystem:
    """
    Sistema de padding de paquetes para uniformizar tamaños.
    """
    
    def __init__(self, packet_size: int = 2048):
        """
        Args:
            packet_size: Tamaño fijo de todos los paquetes (bytes)
        """
        self.packet_size = packet_size
        self.header_size = 16  # sequence(4) + flags(4) + length(4) + checksum(4)
        self.max_payload = packet_size - self.header_size
    
    def pad_message(self, message: bytes) -> List[PaddedPacket]:
        """
        Convertir mensaje en lista de paquetes padded.
        """
        packets = []
        
        # Dividir mensaje si es necesario
        num_packets = (len(message) + self.max_payload - 1) // self.max_payload
        num_packets = max(1, num_packets)  # Al menos 1 paquete
        
        for i in range(num_packets):
            start = i * self.max_payload
            end = min(start + self.max_payload, len(message))
            chunk = message[start:end]
            
            # Calcular padding necesario
            padding_size = self.max_payload - len(chunk)
            padding = secrets.token_bytes(padding_size)
            
            packet = PaddedPacket(
                payload=chunk,
                padding=padding,
                total_size=self.packet_size,
                sequence_number=i,
                is_final=(i == num_packets - 1)
            )
            packets.append(packet)
        
        return packets
    
    def serialize_packet(self, packet: PaddedPacket) -> bytes:
        """
        Serializar paquete a bytes con header.
        """
        # Header
        flags = 1 if packet.is_final else 0
        header = struct.pack(
            '>IIII',
            packet.sequence_number,
            flags,
            len(packet.payload),
            self._checksum(packet.payload)
        )
        
        # Payload + padding
        data = packet.payload + packet.padding
        
        return header + data
    
    def deserialize_packet(self, data: bytes) -> PaddedPacket:
        """
        Deserializar bytes a paquete.
        """
        # Parse header
        seq, flags, payload_len, checksum = struct.unpack('>IIII', data[:16])
        
        # Extract payload (sin padding)
        payload = data[16:16 + payload_len]
        padding = data[16 + payload_len:]
        
        # Verificar checksum
        if self._checksum(payload) != checksum:
            raise ValueError("Checksum mismatch")
        
        return PaddedPacket(
            payload=payload,
            padding=padding,
            total_size=len(data),
            sequence_number=seq,
            is_final=(flags & 1) == 1
        )
    
    def reassemble_message(self, packets: List[PaddedPacket]) -> bytes:
        """
        Reensamblar mensaje desde paquetes.
        """
        # Ordenar por número de secuencia
        sorted_packets = sorted(packets, key=lambda p: p.sequence_number)
        
        # Concatenar payloads
        message = b''.join(p.payload for p in sorted_packets)
        
        return message
    
    def _checksum(self, data: bytes) -> int:
        """Calcular checksum simple."""
        return sum(data) & 0xFFFFFFFF
    
    def calculate_overhead(self, message_size: int) -> dict:
        """
        Calcular overhead de padding para un mensaje.
        """
        packets = self.pad_message(b'\x00' * message_size)
        
        total_transmitted = len(packets) * self.packet_size
        overhead_bytes = total_transmitted - message_size
        overhead_percent = (overhead_bytes / message_size) * 100 if message_size > 0 else 0
        
        return {
            'message_size': message_size,
            'num_packets': len(packets),
            'total_transmitted': total_transmitted,
            'overhead_bytes': overhead_bytes,
            'overhead_percent': overhead_percent
        }


class VariablePaddingSystem:
    """
    Sistema de padding con tamaños variables predefinidos.
    Menor overhead que padding fijo, pero revela categoría de tamaño.
    """
    
    def __init__(self, size_buckets: List[int] = None):
        """
        Args:
            size_buckets: Lista de tamaños permitidos (ordenados)
        """
        self.buckets = size_buckets or [64, 256, 1024, 4096, 16384]
    
    def pad_to_bucket(self, data: bytes) -> Tuple[bytes, int]:
        """
        Pad data al siguiente bucket de tamaño.
        """
        data_size = len(data)
        
        # Encontrar bucket apropiado
        target_size = self.buckets[-1]  # Default al más grande
        for bucket in self.buckets:
            if bucket >= data_size:
                target_size = bucket
                break
        
        # Aplicar padding
        padding_size = target_size - data_size
        padding = secrets.token_bytes(padding_size)
        
        # Prepend length for extraction
        length_prefix = struct.pack('>I', data_size)
        padded = length_prefix + data + padding
        
        return padded, target_size
    
    def extract_from_bucket(self, padded_data: bytes) -> bytes:
        """
        Extraer datos originales de bucket padded.
        """
        # Extract length
        original_length = struct.unpack('>I', padded_data[:4])[0]
        
        # Extract data
        return padded_data[4:4 + original_length]
    
    def analyze_leakage(self) -> dict:
        """
        Analizar cuánta información filtra este esquema.
        """
        import math
        
        # Entropía de la distribución de buckets
        # Si hay n buckets, un adversario puede distinguir n casos
        num_buckets = len(self.buckets)
        max_entropy = math.log2(num_buckets)
        
        return {
            'num_buckets': num_buckets,
            'size_categories': self.buckets,
            'bits_leaked': max_entropy,
            'anonymity_set': num_buckets
        }


def packet_padding_demo():
    print("=== Packet Padding Demo ===\n")
    
    # Sistema de padding fijo
    system = PacketPaddingSystem(packet_size=2048)
    
    # Diferentes mensajes
    messages = [
        (b"short message", "Short"),
        (b"A" * 500, "Medium (500B)"),
        (b"B" * 2000, "Large (2KB)"),
        (b"C" * 5000, "Very Large (5KB)")
    ]
    
    print("--- Fixed Size Padding (2KB packets) ---\n")
    
    for msg, desc in messages:
        packets = system.pad_message(msg)
        overhead = system.calculate_overhead(len(msg))
        
        print(f"{desc}:")
        print(f"  Original: {len(msg)} bytes")
        print(f"  Packets: {len(packets)}")
        print(f"  Transmitted: {overhead['total_transmitted']} bytes")
        print(f"  Overhead: {overhead['overhead_percent']:.1f}%")
        
        # Verificar todos los paquetes tienen mismo tamaño
        serialized = [system.serialize_packet(p) for p in packets]
        sizes = [len(s) for s in serialized]
        print(f"  All packets same size: {len(set(sizes)) == 1}")
        print()
    
    # Sistema de buckets
    print("--- Variable Bucket Padding ---\n")
    
    bucket_system = VariablePaddingSystem()
    leakage = bucket_system.analyze_leakage()
    
    print(f"Buckets: {leakage['size_categories']}")
    print(f"Information leaked: {leakage['bits_leaked']:.2f} bits")
    
    for msg, desc in messages:
        padded, bucket = bucket_system.pad_to_bucket(msg)
        print(f"\n{desc}:")
        print(f"  Original: {len(msg)} bytes -> Bucket: {bucket} bytes")

packet_padding_demo()
\end{lstlisting}

%===========================================
\section{Dummy Traffic (Cover Traffic)}
\label{sec:dummy-traffic}
%===========================================

\subsection{Concepto y Tipos}

\begin{definition}[Dummy Traffic / Cover Traffic]
Tráfico artificial generado para ocultar patrones de comunicación real. El adversario no puede distinguir entre tráfico real y dummy.

\textbf{Propósitos:}
\begin{itemize}
    \item Ocultar cuándo ocurre comunicación real
    \item Uniformizar patrones de actividad
    \item Prevenir timing correlation
    \item Aumentar anonymity set
\end{itemize}
\end{definition}

\begin{mathnote}[Tipos de Cover Traffic]
\textbf{1. Constant-Rate Cover Traffic:}
\begin{itemize}
    \item Enviar paquetes a tasa constante (ej: 1/segundo)
    \item Paquetes reales reemplazan dummies
    \item Máxima resistencia a timing analysis
    \item Alto overhead de bandwidth
\end{itemize}

\textbf{2. Poisson Cover Traffic:}
\begin{itemize}
    \item Intervalos entre paquetes siguen distribución Poisson
    \item Más eficiente que constant-rate
    \item Estadísticamente indistinguible del tráfico real
\end{itemize}

\textbf{3. Loop Traffic:}
\begin{itemize}
    \item Mensajes enviados a uno mismo a través de la red
    \item Aumenta tráfico sin revelar destinos
    \item Usado en Nym mixnet
\end{itemize}

\textbf{4. Chaff Traffic:}
\begin{itemize}
    \item Paquetes indistinguibles de reales pero descartados
    \item Solo el receptor puede identificar chaff
    \item Previene traffic analysis end-to-end
\end{itemize}
\end{mathnote}

\begin{figure}[htbp]
\centering
\begin{tikzpicture}[
    node distance=0.3cm
]
    % Timeline sin cover traffic
    \node at (-1, 4) {\textbf{Sin Cover}};
    \draw[->] (0, 4) -- (12, 4);
    
    % Paquetes reales (sparse)
    \foreach \x in {1, 3.5, 7, 9, 11} {
        \fill[blue] (\x, 4) circle (0.15);
    }
    \node at (6, 3.3) [font=\small] {Patrones visibles: timing, frecuencia};
    
    % Timeline con cover traffic
    \node at (-1, 1.5) {\textbf{Con Cover}};
    \draw[->] (0, 1.5) -- (12, 1.5);
    
    % Cover traffic (regular)
    \foreach \x in {0.5, 1, 1.5, 2, 2.5, 3, 3.5, 4, 4.5, 5, 5.5, 6, 6.5, 7, 7.5, 8, 8.5, 9, 9.5, 10, 10.5, 11, 11.5} {
        \fill[gray!50] (\x, 1.5) circle (0.1);
    }
    
    % Paquetes reales mezclados
    \foreach \x in {1, 3.5, 7, 9, 11} {
        \fill[blue] (\x, 1.5) circle (0.15);
    }
    
    \node at (6, 0.8) [font=\small] {Indistinguible: ¿cuáles son reales?};
    
    % Leyenda
    \fill[blue] (3, -0.5) circle (0.15);
    \node at (4.5, -0.5) [font=\small] {Tráfico real};
    
    \fill[gray!50] (7, -0.5) circle (0.1);
    \node at (8.5, -0.5) [font=\small] {Cover traffic};
\end{tikzpicture}
\caption{Cover traffic oculta los patrones del tráfico real}
\label{fig:cover-traffic}
\end{figure}

\begin{lstlisting}[language=Python, caption={Sistema de Cover Traffic}]
from dataclasses import dataclass
from datetime import datetime, timedelta
from typing import List, Callable, Optional
import random
import threading
import time
from enum import Enum

class TrafficType(Enum):
    REAL = "real"
    COVER = "cover"
    LOOP = "loop"
    CHAFF = "chaff"

@dataclass
class TrafficPacket:
    packet_id: str
    traffic_type: TrafficType
    payload: bytes
    timestamp: datetime
    destination: str

class CoverTrafficGenerator:
    """
    Generador de cover traffic para ocultar patrones de comunicación.
    """
    
    def __init__(
        self,
        packet_size: int = 2048,
        base_rate: float = 1.0  # paquetes por segundo
    ):
        self.packet_size = packet_size
        self.base_rate = base_rate
        self.running = False
        self.packet_queue: List[TrafficPacket] = []
        self.sent_packets: List[TrafficPacket] = []
        
        # Estadísticas
        self.stats = {
            'real_sent': 0,
            'cover_sent': 0,
            'loop_sent': 0,
            'total_bytes': 0
        }
    
    def start_constant_rate(self, destinations: List[str]):
        """
        Iniciar generación de cover traffic a tasa constante.
        """
        self.running = True
        
        while self.running:
            # Verificar si hay paquete real pendiente
            real_packet = self._get_pending_real()
            
            if real_packet:
                # Enviar paquete real
                self._send_packet(real_packet)
            else:
                # Enviar cover packet
                cover = self._generate_cover_packet(
                    random.choice(destinations)
                )
                self._send_packet(cover)
            
            # Esperar intervalo fijo
            time.sleep(1.0 / self.base_rate)
    
    def start_poisson_rate(self, destinations: List[str], mean_interval: float = 1.0):
        """
        Iniciar generación con intervalos Poisson.
        Más eficiente y estadísticamente más difícil de analizar.
        """
        self.running = True
        
        while self.running:
            real_packet = self._get_pending_real()
            
            if real_packet:
                self._send_packet(real_packet)
            else:
                cover = self._generate_cover_packet(
                    random.choice(destinations)
                )
                self._send_packet(cover)
            
            # Intervalo Poisson (exponencial)
            interval = random.expovariate(1.0 / mean_interval)
            time.sleep(interval)
    
    def queue_real_packet(self, payload: bytes, destination: str):
        """
        Encolar paquete real para envío en siguiente slot.
        """
        packet = TrafficPacket(
            packet_id=f"real_{len(self.sent_packets)}",
            traffic_type=TrafficType.REAL,
            payload=self._pad_payload(payload),
            timestamp=datetime.now(),
            destination=destination
        )
        self.packet_queue.append(packet)
    
    def generate_loop_traffic(self, own_address: str, num_loops: int = 5):
        """
        Generar loop traffic (mensajes a uno mismo).
        """
        for i in range(num_loops):
            loop_packet = TrafficPacket(
                packet_id=f"loop_{i}",
                traffic_type=TrafficType.LOOP,
                payload=self._random_payload(),
                timestamp=datetime.now(),
                destination=own_address
            )
            self._send_packet(loop_packet)
            
            # Intervalo aleatorio
            time.sleep(random.uniform(0.1, 1.0))
    
    def _generate_cover_packet(self, destination: str) -> TrafficPacket:
        """Generar paquete de cover traffic."""
        return TrafficPacket(
            packet_id=f"cover_{self.stats['cover_sent']}",
            traffic_type=TrafficType.COVER,
            payload=self._random_payload(),
            timestamp=datetime.now(),
            destination=destination
        )
    
    def _get_pending_real(self) -> Optional[TrafficPacket]:
        """Obtener siguiente paquete real pendiente."""
        if self.packet_queue:
            return self.packet_queue.pop(0)
        return None
    
    def _send_packet(self, packet: TrafficPacket):
        """Enviar paquete (simulación)."""
        self.sent_packets.append(packet)
        
        # Actualizar estadísticas
        if packet.traffic_type == TrafficType.REAL:
            self.stats['real_sent'] += 1
        elif packet.traffic_type == TrafficType.COVER:
            self.stats['cover_sent'] += 1
        elif packet.traffic_type == TrafficType.LOOP:
            self.stats['loop_sent'] += 1
        
        self.stats['total_bytes'] += len(packet.payload)
    
    def _pad_payload(self, payload: bytes) -> bytes:
        """Pad payload a tamaño fijo."""
        if len(payload) >= self.packet_size:
            return payload[:self.packet_size]
        padding = bytes(self.packet_size - len(payload))
        return payload + padding
    
    def _random_payload(self) -> bytes:
        """Generar payload aleatorio para cover traffic."""
        return bytes(random.getrandbits(8) for _ in range(self.packet_size))
    
    def stop(self):
        """Detener generación de cover traffic."""
        self.running = False
    
    def get_statistics(self) -> dict:
        """Obtener estadísticas de tráfico."""
        total_packets = self.stats['real_sent'] + self.stats['cover_sent'] + self.stats['loop_sent']
        
        return {
            **self.stats,
            'total_packets': total_packets,
            'real_ratio': self.stats['real_sent'] / total_packets if total_packets > 0 else 0,
            'overhead_ratio': 1 - (self.stats['real_sent'] / total_packets) if total_packets > 0 else 1,
            'bandwidth_overhead_bytes': (self.stats['cover_sent'] + self.stats['loop_sent']) * self.packet_size
        }
    
    def analyze_timing_resistance(self) -> dict:
        """
        Analizar qué tan bien el cover traffic oculta el timing del tráfico real.
        """
        real_timestamps = [
            p.timestamp for p in self.sent_packets 
            if p.traffic_type == TrafficType.REAL
        ]
        
        all_timestamps = [p.timestamp for p in self.sent_packets]
        
        if len(all_timestamps) < 2:
            return {'insufficient_data': True}
        
        # Calcular intervalos
        def get_intervals(timestamps):
            sorted_ts = sorted(timestamps)
            return [
                (sorted_ts[i+1] - sorted_ts[i]).total_seconds()
                for i in range(len(sorted_ts) - 1)
            ]
        
        all_intervals = get_intervals(all_timestamps)
        
        import statistics
        
        return {
            'total_packets': len(all_timestamps),
            'real_packets': len(real_timestamps),
            'mean_interval': statistics.mean(all_intervals),
            'std_interval': statistics.stdev(all_intervals) if len(all_intervals) > 1 else 0,
            'timing_uniformity': 1 - (statistics.stdev(all_intervals) / statistics.mean(all_intervals)) 
                                if len(all_intervals) > 1 and statistics.mean(all_intervals) > 0 else 0
        }


def cover_traffic_demo():
    print("=== Cover Traffic Demo ===\n")
    
    generator = CoverTrafficGenerator(packet_size=2048, base_rate=10)
    
    destinations = ["node_a", "node_b", "node_c", "node_d"]
    
    # Simular envío de paquetes reales
    print("--- Queueing Real Packets ---")
    
    for i in range(5):
        generator.queue_real_packet(
            payload=f"Real message {i}".encode(),
            destination=random.choice(destinations)
        )
    
    print(f"Queued {len(generator.packet_queue)} real packets")
    
    # Simular constant rate durante algunos "ticks"
    print("\n--- Constant Rate Simulation ---")
    
    for tick in range(30):
        # Simular un tick del generador
        real_packet = generator._get_pending_real()
        
        if real_packet:
            generator._send_packet(real_packet)
        else:
            cover = generator._generate_cover_packet(
                random.choice(destinations)
            )
            generator._send_packet(cover)
    
    # Estadísticas
    print("\n--- Statistics ---")
    stats = generator.get_statistics()
    
    print(f"Total packets sent: {stats['total_packets']}")
    print(f"Real packets: {stats['real_sent']}")
    print(f"Cover packets: {stats['cover_sent']}")
    print(f"Real ratio: {stats['real_ratio']:.2%}")
    print(f"Overhead ratio: {stats['overhead_ratio']:.2%}")
    
    # Análisis de timing
    print("\n--- Timing Analysis ---")
    timing = generator.analyze_timing_resistance()
    
    print(f"Mean interval: {timing.get('mean_interval', 0):.4f}s")
    print(f"Std interval: {timing.get('std_interval', 0):.4f}s")
    print(f"Timing uniformity: {timing.get('timing_uniformity', 0):.2%}")
    
    print("\n--- Adversary's View ---")
    print("From observer's perspective:")
    print(f"  Sees {stats['total_packets']} identical packets")
    print(f"  Cannot distinguish which {stats['real_sent']} are real")
    print(f"  Effective anonymity set: {stats['total_packets']}")

cover_traffic_demo()
\end{lstlisting}

%===========================================
\section{P2P Network Privacy}
\label{sec:p2p-privacy}
%===========================================

\subsection{libp2p}

\begin{definition}[libp2p]
Framework modular de networking peer-to-peer desarrollado por Protocol Labs:

\textbf{Usado por:}
\begin{itemize}
    \item IPFS (InterPlanetary File System)
    \item Ethereum 2.0 (consensus layer)
    \item Filecoin
    \item Polkadot
\end{itemize}

\textbf{Características:}
\begin{itemize}
    \item Transport agnostic (TCP, QUIC, WebSocket)
    \item Multiplexing de conexiones
    \item NAT traversal
    \item Peer discovery (DHT, mDNS)
    \item Encryption (Noise protocol)
\end{itemize}
\end{definition}

\begin{advertencia}
\textbf{Limitaciones de privacidad en libp2p:}

\begin{enumerate}
    \item \textbf{IP visible a peers:}
    \begin{itemize}
        \item Cada peer conoce la IP de sus conexiones directas
        \item No hay onion routing nativo
    \end{itemize}
    
    \item \textbf{Gossip protocol:}
    \begin{itemize}
        \item Mensajes se propagan a múltiples peers
        \item Primer receptor puede inferir origen
    \end{itemize}
    
    \item \textbf{DHT queries:}
    \begin{itemize}
        \item Queries revelan qué datos buscas
        \item Nodos DHT pueden perfilar usuarios
    \end{itemize}
    
    \item \textbf{Peer IDs:}
    \begin{itemize}
        \item Identificadores persistentes
        \item Vinculan sesiones a través del tiempo
    \end{itemize}
\end{enumerate}
\end{advertencia}

\begin{lstlisting}[language=Python, caption={Análisis de privacidad en redes P2P}]
from dataclasses import dataclass
from typing import Dict, List, Set
from collections import defaultdict

@dataclass
class Peer:
    peer_id: str
    ip_address: str
    connected_peers: Set[str]
    messages_originated: List[str]
    messages_relayed: List[str]

class P2PPrivacyAnalyzer:
    """
    Analizador de privacidad para redes P2P tipo libp2p.
    """
    
    def __init__(self):
        self.peers: Dict[str, Peer] = {}
        self.message_first_seen: Dict[str, str] = {}  # msg_id -> first peer
        self.ip_to_peer: Dict[str, str] = {}
    
    def add_peer(self, peer: Peer):
        """Registrar peer en la red."""
        self.peers[peer.peer_id] = peer
        self.ip_to_peer[peer.ip_address] = peer.peer_id
    
    def record_message_propagation(
        self,
        message_id: str,
        peer_id: str,
        is_originator: bool
    ):
        """
        Registrar propagación de mensaje.
        """
        if message_id not in self.message_first_seen:
            self.message_first_seen[message_id] = peer_id
        
        if peer_id in self.peers:
            if is_originator:
                self.peers[peer_id].messages_originated.append(message_id)
            else:
                self.peers[peer_id].messages_relayed.append(message_id)
    
    def analyze_ip_exposure(self) -> Dict:
        """
        Analizar exposición de IPs a otros peers.
        """
        exposure = {}
        
        for peer_id, peer in self.peers.items():
            # IPs que este peer conoce (sus conexiones directas)
            known_ips = set()
            for connected_id in peer.connected_peers:
                if connected_id in self.peers:
                    known_ips.add(self.peers[connected_id].ip_address)
            
            # IPs que conocen este peer
            exposing_peers = [
                p.peer_id for p in self.peers.values()
                if peer_id in p.connected_peers
            ]
            
            exposure[peer_id] = {
                'ip': peer.ip_address,
                'knows_ips': len(known_ips),
                'exposed_to': len(exposing_peers),
                'exposure_peers': exposing_peers[:5]  # Primeros 5
            }
        
        return exposure
    
    def analyze_first_seen_heuristic(self) -> Dict:
        """
        Analizar efectividad del heurístico "first-seen" para
        identificar originadores de mensajes.
        """
        correct_attributions = 0
        total_messages = 0
        
        for peer in self.peers.values():
            for msg_id in peer.messages_originated:
                total_messages += 1
                first_seen = self.message_first_seen.get(msg_id)
                if first_seen == peer.peer_id:
                    correct_attributions += 1
        
        return {
            'total_messages': total_messages,
            'correctly_attributed': correct_attributions,
            'attribution_accuracy': correct_attributions / total_messages if total_messages > 0 else 0,
            'privacy_level': 1 - (correct_attributions / total_messages) if total_messages > 0 else 1
        }
    
    def recommend_mitigations(self) -> List[str]:
        """
        Recomendar mitigaciones de privacidad.
        """
        recommendations = []
        
        exposure = self.analyze_ip_exposure()
        first_seen = self.analyze_first_seen_heuristic()
        
        # Analizar exposición promedio
        avg_exposure = sum(e['exposed_to'] for e in exposure.values()) / len(exposure) if exposure else 0
        
        if avg_exposure > 10:
            recommendations.append(
                "High peer exposure - consider reducing max connections"
            )
        
        if first_seen['attribution_accuracy'] > 0.5:
            recommendations.append(
                "High first-seen attribution - implement Dandelion protocol"
            )
        
        recommendations.extend([
            "Use Tor transport for libp2p connections",
            "Rotate peer IDs periodically",
            "Add random delays before message propagation",
            "Use encrypted DHT queries"
        ])
        
        return recommendations


def p2p_privacy_demo():
    print("=== P2P Privacy Analysis Demo ===\n")
    
    analyzer = P2PPrivacyAnalyzer()
    
    # Crear red de ejemplo
    peers = [
        Peer(f"peer_{i}", f"192.168.1.{i}", set(), [], [])
        for i in range(10)
    ]
    
    # Conectar peers (topología de malla parcial)
    import random
    for i, peer in enumerate(peers):
        # Cada peer conecta a 3-5 otros
        num_connections = random.randint(3, 5)
        possible_peers = [p.peer_id for p in peers if p.peer_id != peer.peer_id]
        peer.connected_peers = set(random.sample(possible_peers, num_connections))
    
    for peer in peers:
        analyzer.add_peer(peer)
    
    # Simular propagación de mensajes
    print("--- Simulating Message Propagation ---")
    
    for i in range(20):
        originator = random.choice(peers)
        msg_id = f"msg_{i}"
        
        # Registrar mensaje
        analyzer.record_message_propagation(msg_id, originator.peer_id, True)
        
        # Propagar a peers conectados
        for connected_id in originator.connected_peers:
            analyzer.record_message_propagation(msg_id, connected_id, False)
    
    # Análisis de exposición IP
    print("\n--- IP Exposure Analysis ---")
    exposure = analyzer.analyze_ip_exposure()
    
    for peer_id, data in list(exposure.items())[:3]:
        print(f"\n{peer_id} ({data['ip']}):")
        print(f"  Knows IPs of: {data['knows_ips']} peers")
        print(f"  Exposed to: {data['exposed_to']} peers")
    
    # Análisis first-seen
    print("\n--- First-Seen Heuristic Analysis ---")
    first_seen = analyzer.analyze_first_seen_heuristic()
    
    print(f"Messages analyzed: {first_seen['total_messages']}")
    print(f"Correctly attributed: {first_seen['correctly_attributed']}")
    print(f"Attribution accuracy: {first_seen['attribution_accuracy']:.2%}")
    print(f"Privacy level: {first_seen['privacy_level']:.2%}")
    
    # Recomendaciones
    print("\n--- Privacy Recommendations ---")
    for rec in analyzer.recommend_mitigations():
        print(f"  • {rec}")

p2p_privacy_demo()
\end{lstlisting}

\subsection{Waku}

\begin{definition}[Waku]
Protocolo de mensajería descentralizado desarrollado por Status, diseñado como evolución de Whisper con mejor privacidad y escalabilidad:

\textbf{Versiones:}
\begin{itemize}
    \item Waku v1: Compatible con Whisper
    \item Waku v2: Rediseño con mejor privacidad
\end{itemize}

\textbf{Características de privacidad:}
\begin{itemize}
    \item \textbf{Content topics:} Agrupación de mensajes sin revelar contenido
    \item \textbf{Encryption:} Mensajes encriptados end-to-end
    \item \textbf{Relay nodes:} No ven contenido, solo metadata limitada
    \item \textbf{Anonymity set:} Todos los usuarios del mismo topic
\end{itemize}
\end{definition}

\begin{table}[htbp]
\centering
\small
\begin{tabular}{lcc}
\toprule
\textbf{Aspecto} & \textbf{libp2p básico} & \textbf{Waku v2} \\
\midrule
IP visible a peers & Sí & Relay oculta \\
Contenido visible a relay & Sí & No (E2E encrypted) \\
Metadata visible & Alta & Limitada a topic \\
Sender anonymity & Baja & Media-Alta \\
Receiver anonymity & Baja & Alta \\
\bottomrule
\end{tabular}
\caption{Comparación de privacidad entre libp2p y Waku}
\end{table}

%===========================================
\section*{Resumen del Bloque}
%===========================================

\begin{keyidea}
\textbf{Puntos clave de Metadata y Network Privacy (Bloque 16):}

\textbf{1. IP Leakage:}
\begin{itemize}
    \item Múltiples vectores: P2P, RPC, dApps, DNS
    \item Protecciones: VPN (básica), Tor (mejor), nodo propio
    \item IP revela geolocalización e identidad potencial
\end{itemize}

\textbf{2. Timing Correlation:}
\begin{itemize}
    \item Vincula eventos por proximidad temporal
    \item Defensas: delays aleatorios, batching, cover traffic
    \item Crítico en mixers y privacidad financiera
\end{itemize}

\textbf{3. Traffic Analysis:}
\begin{itemize}
    \item Fingerprinting por tamaño, timing, patrones
    \item Website fingerprinting efectivo incluso con Tor
    \item Defensas: padding, cover traffic, mixing
\end{itemize}

\textbf{4. Onion Routing vs Mixnets:}
\begin{itemize}
    \item Tor: baja latencia, vulnerable a GPA
    \item Mixnets: alta latencia, resiste timing analysis
    \item Trade-off latencia vs anonimato
\end{itemize}

\textbf{5. VPN/Tor/I2P/Nym:}
\begin{itemize}
    \item VPN: confianza centralizada, básico
    \item Tor: distribuido, buen balance
    \item I2P: servicios ocultos, P2P
    \item Nym: mixnet con incentivos, máximo anonimato
\end{itemize}

\textbf{6. Packet Padding:}
\begin{itemize}
    \item Uniformiza tamaños para ocultar contenido
    \item Trade-off overhead vs privacidad
    \item Buckets variables como compromiso
\end{itemize}

\textbf{7. Cover Traffic:}
\begin{itemize}
    \item Tráfico falso indistinguible del real
    \item Tipos: constant-rate, Poisson, loop, chaff
    \item Esencial para máxima privacidad temporal
\end{itemize}

\textbf{8. P2P Privacy:}
\begin{itemize}
    \item libp2p: funcional pero sin privacidad nativa
    \item Waku: mensajería con privacidad mejorada
    \item Mitigaciones: Tor transport, Dandelion, delays
\end{itemize}
\end{keyidea}

\begin{bancozk}
\textbf{Stack de Network Privacy en Banco ZK:}

\begin{center}
\begin{tikzpicture}[
    layer/.style={rectangle, draw, minimum width=10cm, minimum height=0.8cm, align=center}
]
    \node[layer, fill=purple!20] at (0, 4) {Capa Aplicación: ZK Proofs + Encrypted Payloads};
    \node[layer, fill=blue!20] at (0, 3) {Capa Mensajería: Waku v2 (content topics)};
    \node[layer, fill=green!20] at (0, 2) {Capa Mixing: Nym Mixnet (Poisson delays)};
    \node[layer, fill=yellow!20] at (0, 1) {Capa Transport: Tor / I2P};
    \node[layer, fill=orange!20] at (0, 0) {Capa Red: Packet Padding + Cover Traffic};
\end{tikzpicture}
\end{center}

\textbf{Configuraciones de privacidad:}

\begin{enumerate}
    \item \textbf{Modo Estándar:}
    \begin{itemize}
        \item Tor para conexiones
        \item Nodo propio recomendado
        \item Delays aleatorios en broadcast
    \end{itemize}
    
    \item \textbf{Modo Alto Anonimato:}
    \begin{itemize}
        \item Nym mixnet para transacciones
        \item Cover traffic activo
        \item Packet padding uniforme
        \item Mayor latencia (minutos)
    \end{itemize}
    
    \item \textbf{Modo Máximo:}
    \begin{itemize}
        \item Todo lo anterior +
        \item Loop traffic continuo
        \item Conexiones solo via .onion
        \item Rotación de identidad de red
    \end{itemize}
\end{enumerate}

\textbf{Métricas de privacidad de red:}
\begin{center}
\begin{tabular}{lc}
\toprule
\textbf{Métrica} & \textbf{Target} \\
\midrule
IP exposure a terceros & 0 (via Tor) \\
Timing correlation resistance & $>$ 95\% \\
Cover traffic ratio & $\geq$ 50\% \\
Packet size uniformity & 100\% \\
First-seen attribution & $<$ 10\% \\
\bottomrule
\end{tabular}
\end{center}
\end{bancozk}

\begin{figure}[htbp]
\centering
\begin{tikzpicture}
    % Diagrama final de amenazas y protecciones
    \node at (0, 5) {\textbf{Threat Model y Protecciones}};
    
    % Amenazas
    \node[draw, rectangle, fill=red!20, minimum width=3cm] (t1) at (-4, 3) {IP Leakage};
    \node[draw, rectangle, fill=red!20, minimum width=3cm] (t2) at (-4, 1.5) {Timing Analysis};
    \node[draw, rectangle, fill=red!20, minimum width=3cm] (t3) at (-4, 0) {Traffic Fingerprint};
    \node[draw, rectangle, fill=red!20, minimum width=3cm] (t4) at (-4, -1.5) {P2P Deanon};
    
    % Protecciones
    \node[draw, rectangle, fill=green!20, minimum width=3cm] (p1) at (4, 3) {Tor/Nym/VPN};
    \node[draw, rectangle, fill=green!20, minimum width=3cm] (p2) at (4, 1.5) {Cover Traffic};
    \node[draw, rectangle, fill=green!20, minimum width=3cm] (p3) at (4, 0) {Packet Padding};
    \node[draw, rectangle, fill=green!20, minimum width=3cm] (p4) at (4, -1.5) {Dandelion++};
    
    % Conexiones
    \draw[->, thick] (t1) -- (p1);
    \draw[->, thick] (t2) -- (p2);
    \draw[->, thick] (t3) -- (p3);
    \draw[->, thick] (t4) -- (p4);
    
    % Etiquetas
    \node at (-4, 4) [red] {\textbf{Amenazas}};
    \node at (4, 4) [green!60!black] {\textbf{Protecciones}};
\end{tikzpicture}
\caption{Mapeo de amenazas de red a protecciones correspondientes}
\label{fig:threat-protection-map}
\end{figure}

% Fin del Bloque 16
\cleardoublepage
% ============================================
% BLOQUE 17: PRIVATE SMART CONTRACTS - VERSIÓN EXPANDIDA
% Parte 1: Introducción, Secret Network, Oasis Network
% Objetivo: ~60 páginas (~3600 líneas)
% ============================================

\chapter{Private Smart Contracts}
\label{ch:bloque17}

\begin{center}
\textit{Bloque 17 de 20}
\end{center}

\section*{Introducción}

Los smart contracts tradicionales (Ethereum, Solana) ejecutan código de forma transparente: todos pueden ver el estado, los inputs y los outputs de cada transacción. Esta transparencia, aunque útil para auditoría, es incompatible con muchos casos de uso que requieren confidencialidad. Este bloque explora las diferentes aproximaciones para lograr smart contracts privados.

\begin{keyidea}
\textbf{El problema de la transparencia en smart contracts:}

\textbf{En Ethereum tradicional, todos pueden ver:}
\begin{itemize}
    \item Estado del contrato (balances, mappings, variables)
    \item Inputs de cada llamada (parámetros de función)
    \item Outputs de cada transacción
    \item Historial completo de interacciones
\end{itemize}

\textbf{Casos de uso afectados:}
\begin{itemize}
    \item Subastas (pujas visibles = frontrunning)
    \item Votación (votos públicos = coerción)
    \item Trading (órdenes visibles = MEV extraction)
    \item DeFi privado (posiciones visibles = liquidaciones predatorias)
    \item Datos sensibles (salud, finanzas, identidad)
\end{itemize}
\end{keyidea}

\begin{figure}[htbp]
\centering
\begin{tikzpicture}[
    node distance=1cm,
    approach/.style={rectangle, draw, rounded corners, minimum width=3cm, minimum height=1.5cm, align=center}
]
    % Título
    \node at (0, 4.5) {\textbf{Aproximaciones a Smart Contracts Privados}};
    
    % Cuatro aproximaciones
    \node[approach, fill=blue!20] (zk) at (-4.5, 2) {ZK Proofs\\(Aleo, Mina)};
    \node[approach, fill=green!20] (tee) at (-1.5, 2) {TEE\\(Secret, Oasis)};
    \node[approach, fill=orange!20] (mpc) at (1.5, 2) {MPC\\(Partisia)};
    \node[approach, fill=purple!20] (fhe) at (4.5, 2) {FHE\\(Fhenix, Zama)};
    
    % Características
    \node at (-4.5, 0) [align=center, font=\small] {
        Pura criptografía\\
        Sin hardware trust\\
        Alto costo de prueba
    };
    
    \node at (-1.5, 0) [align=center, font=\small] {
        Hardware trust\\
        Baja latencia\\
        Side-channel risks
    };
    
    \node at (1.5, 0) [align=center, font=\small] {
        Distribuido\\
        Threshold trust\\
        Comunicación alta
    };
    
    \node at (4.5, 0) [align=center, font=\small] {
        Cómputo encriptado\\
        Sin decryption\\
        Muy costoso
    };
    
    % Escala
    \draw[->] (-6, -1.5) -- (6, -1.5);
    \node at (-6, -2) [font=\small] {Trust: Criptografía};
    \node at (6, -2) [font=\small] {Trust: Hardware/Partes};
\end{tikzpicture}
\caption{Cuatro aproximaciones principales a smart contracts privados}
\label{fig:private-sc-approaches}
\end{figure}

\begin{mathnote}[Comparación de Modelos de Confianza]
\textbf{1. Zero-Knowledge Proofs (ZK):}
\begin{itemize}
    \item Trust: Solo en matemáticas y criptografía
    \item Asunción: Problema computacionalmente difícil (DLP, factorización)
    \item Verificación: On-chain, cualquiera puede verificar
\end{itemize}

\textbf{2. Trusted Execution Environment (TEE):}
\begin{itemize}
    \item Trust: Hardware del fabricante (Intel SGX, AMD SEV)
    \item Asunción: Enclave no comprometido, sin side-channels explotables
    \item Verificación: Attestation del enclave
\end{itemize}

\textbf{3. Multi-Party Computation (MPC):}
\begin{itemize}
    \item Trust: Threshold de partes honestas ($t$ de $n$)
    \item Asunción: No más de $t-1$ partes coludan
    \item Verificación: Protocolo garantiza correctitud
\end{itemize}

\textbf{4. Fully Homomorphic Encryption (FHE):}
\begin{itemize}
    \item Trust: Solo en criptografía
    \item Asunción: Problema LWE/RLWE es difícil
    \item Verificación: El resultado descifrado es correcto
\end{itemize}
\end{mathnote}

%===========================================
\section{Secret Network}
\label{sec:secret-network}
%===========================================

\subsection{Arquitectura y Diseño}

\begin{definition}[Secret Network]
Blockchain Layer-1 basada en Cosmos SDK que permite smart contracts con estado encriptado mediante Trusted Execution Environments (TEE), específicamente Intel SGX.

\textbf{Características principales:}
\begin{itemize}
    \item Smart contracts ``secretos'' con estado encriptado
    \item Solo el enclave SGX puede ver datos en claro
    \item Viewing keys para acceso selectivo a datos
    \item Compatibilidad con ecosistema Cosmos (IBC)
    \item Token nativo: SCRT
\end{itemize}
\end{definition}

\begin{figure}[htbp]
\centering
\begin{tikzpicture}[
    node distance=1cm,
    component/.style={rectangle, draw, minimum width=3cm, minimum height=1cm, align=center},
    enclave/.style={rectangle, draw, minimum width=4cm, minimum height=2.5cm, fill=green!20, rounded corners},
    data/.style={rectangle, draw, minimum width=2cm, minimum height=0.6cm, font=\small}
]
    % Nodo Secret
    \draw[thick] (-3, -0.5) rectangle (7, 5);
    \node at (2, 5.3) {\textbf{Secret Network Node}};
    
    % Enclave SGX
    \node[enclave] (enclave) at (2, 2.5) {};
    \node at (2, 3.8) [font=\small] {\textbf{Intel SGX Enclave}};
    
    % Dentro del enclave
    \node[component, fill=white, minimum height=0.6cm] at (2, 2.8) {Decrypted State};
    \node[component, fill=white, minimum height=0.6cm] at (2, 2) {Contract Logic};
    \node[component, fill=white, minimum height=0.6cm] at (2, 1.2) {Decrypted Inputs};
    
    % Fuera del enclave
    \node[data, fill=red!20] (enc_state) at (-1, 0) {Encrypted State};
    \node[data, fill=red!20] (enc_input) at (5, 0) {Encrypted Input};
    
    % Blockchain
    \node[component, fill=blue!20] (chain) at (2, -1.5) {Tendermint Consensus};
    
    % Flechas
    \draw[->, thick] (enc_state) -- (0, 1) -- (0.5, 1.2);
    \draw[->, thick] (enc_input) -- (4, 1) -- (3.5, 1.2);
    \draw[->, thick] (enclave) -- (chain);
    
    % Labels
    \node at (-1, 0.7) [font=\tiny] {decrypt};
    \node at (5, 0.7) [font=\tiny] {decrypt};
\end{tikzpicture}
\caption{Arquitectura de Secret Network: estado encriptado, procesamiento en enclave SGX}
\label{fig:secret-architecture}
\end{figure}

\subsection{Intel SGX y Trusted Execution}

\begin{definition}[Intel SGX (Software Guard Extensions)]
Conjunto de instrucciones de CPU que permite crear ``enclaves'' - regiones de memoria aisladas y encriptadas donde el código se ejecuta de forma protegida.

\textbf{Propiedades:}
\begin{itemize}
    \item \textbf{Confidencialidad:} Datos en enclave encriptados en RAM
    \item \textbf{Integridad:} Código no puede ser modificado
    \item \textbf{Attestation:} Prueba criptográfica de qué código ejecuta
    \item \textbf{Sealing:} Datos persistentes encriptados a enclave específico
\end{itemize}
\end{definition}

\begin{advertencia}
\textbf{Vulnerabilidades conocidas de SGX:}

\begin{enumerate}
    \item \textbf{Side-channel attacks:}
    \begin{itemize}
        \item Spectre, Meltdown, Foreshadow
        \item Timing attacks, cache attacks
        \item Requieren parches continuos
    \end{itemize}
    
    \item \textbf{Rollback attacks:}
    \begin{itemize}
        \item OS controla persistencia
        \item Puede restaurar estado anterior
    \end{itemize}
    
    \item \textbf{Trust en Intel:}
    \begin{itemize}
        \item Intel genera claves de attestation
        \item Backdoors teóricamente posibles
    \end{itemize}
    
    \item \textbf{Disponibilidad:}
    \begin{itemize}
        \item Intel puede revocar attestation keys
        \item Dependencia de proveedor único
    \end{itemize}
\end{enumerate}

\textbf{Mitigación en Secret Network:}
\begin{itemize}
    \item Múltiples nodos con diferentes CPUs/versiones
    \item Threshold schemes para claves críticas
    \item Actualizaciones frecuentes de enclaves
\end{itemize}
\end{advertencia}

\subsection{Desarrollo de Secret Contracts}

\begin{lstlisting}[language=Rust, caption={Secret Contract básico en Rust (CosmWasm modificado)}]
// Estructura del contrato secreto
use cosmwasm_std::{
    entry_point, to_binary, Binary, Deps, DepsMut, Env,
    MessageInfo, Response, StdResult,
};
use secret_toolkit::viewing_key::{ViewingKey, ViewingKeyStore};

// Estado encriptado (solo visible dentro del enclave)
#[derive(Serialize, Deserialize, Clone)]
pub struct State {
    pub owner: Addr,
    pub secret_data: String,
    pub private_balances: HashMap<Addr, Uint128>,
    pub total_supply: Uint128,
}

// Mensajes de inicialización
#[derive(Serialize, Deserialize)]
pub struct InstantiateMsg {
    pub initial_supply: Uint128,
}

// Mensajes de ejecución (inputs encriptados)
#[derive(Serialize, Deserialize)]
pub enum ExecuteMsg {
    // Transferencia privada
    Transfer {
        recipient: String,
        amount: Uint128,
    },
    // Crear viewing key para el usuario
    CreateViewingKey {
        entropy: String,
    },
    // Depositar (público -> privado)
    Deposit {},
    // Retirar (privado -> público)
    Withdraw {
        amount: Uint128,
    },
}

// Mensajes de consulta (requieren viewing key para datos privados)
#[derive(Serialize, Deserialize)]
pub enum QueryMsg {
    // Consulta pública
    TotalSupply {},
    // Consulta privada (requiere viewing key)
    Balance {
        address: String,
        key: String,
    },
    // Historial de transacciones privado
    TransactionHistory {
        address: String,
        key: String,
        page: u32,
        page_size: u32,
    },
}

// ========== ENTRY POINTS ==========

#[entry_point]
pub fn instantiate(
    deps: DepsMut,
    _env: Env,
    info: MessageInfo,
    msg: InstantiateMsg,
) -> StdResult<Response> {
    let state = State {
        owner: info.sender.clone(),
        secret_data: String::new(),
        private_balances: HashMap::new(),
        total_supply: msg.initial_supply,
    };
    
    // Guardar estado (será encriptado por el enclave)
    STATE.save(deps.storage, &state)?;
    
    Ok(Response::new()
        .add_attribute("action", "instantiate")
        .add_attribute("owner", info.sender))
}

#[entry_point]
pub fn execute(
    deps: DepsMut,
    env: Env,
    info: MessageInfo,
    msg: ExecuteMsg,
) -> StdResult<Response> {
    match msg {
        ExecuteMsg::Transfer { recipient, amount } => {
            execute_transfer(deps, env, info, recipient, amount)
        }
        ExecuteMsg::CreateViewingKey { entropy } => {
            execute_create_viewing_key(deps, env, info, entropy)
        }
        ExecuteMsg::Deposit {} => {
            execute_deposit(deps, env, info)
        }
        ExecuteMsg::Withdraw { amount } => {
            execute_withdraw(deps, env, info, amount)
        }
    }
}

// Transferencia privada
fn execute_transfer(
    deps: DepsMut,
    _env: Env,
    info: MessageInfo,
    recipient: String,
    amount: Uint128,
) -> StdResult<Response> {
    let mut state = STATE.load(deps.storage)?;
    
    let sender = info.sender.clone();
    let recipient_addr = deps.api.addr_validate(&recipient)?;
    
    // Verificar balance (dentro del enclave, no visible externamente)
    let sender_balance = state.private_balances
        .get(&sender)
        .copied()
        .unwrap_or(Uint128::zero());
    
    if sender_balance < amount {
        return Err(StdError::generic_err("Insufficient balance"));
    }
    
    // Actualizar balances (estado encriptado)
    state.private_balances.insert(
        sender.clone(),
        sender_balance - amount
    );
    
    let recipient_balance = state.private_balances
        .get(&recipient_addr)
        .copied()
        .unwrap_or(Uint128::zero());
    
    state.private_balances.insert(
        recipient_addr.clone(),
        recipient_balance + amount
    );
    
    STATE.save(deps.storage, &state)?;
    
    // Log mínimo (no revela montos ni direcciones)
    Ok(Response::new()
        .add_attribute("action", "transfer"))
}

// Crear viewing key
fn execute_create_viewing_key(
    deps: DepsMut,
    env: Env,
    info: MessageInfo,
    entropy: String,
) -> StdResult<Response> {
    // Generar viewing key única para este usuario
    let key = ViewingKey::create(
        deps.storage,
        &info,
        &env,
        info.sender.as_str(),
        entropy.as_bytes(),
    );
    
    Ok(Response::new()
        .add_attribute("action", "create_viewing_key")
        .set_data(to_binary(&CreateViewingKeyResponse { key })?))
}

// ========== QUERIES ==========

#[entry_point]
pub fn query(deps: Deps, _env: Env, msg: QueryMsg) -> StdResult<Binary> {
    match msg {
        QueryMsg::TotalSupply {} => {
            // Consulta pública
            let state = STATE.load(deps.storage)?;
            to_binary(&TotalSupplyResponse {
                total_supply: state.total_supply,
            })
        }
        QueryMsg::Balance { address, key } => {
            // Consulta privada - requiere viewing key válida
            query_balance(deps, address, key)
        }
        QueryMsg::TransactionHistory { address, key, page, page_size } => {
            query_transaction_history(deps, address, key, page, page_size)
        }
    }
}

fn query_balance(deps: Deps, address: String, key: String) -> StdResult<Binary> {
    let addr = deps.api.addr_validate(&address)?;
    
    // Verificar viewing key
    ViewingKey::check(deps.storage, &address, &key)?;
    
    let state = STATE.load(deps.storage)?;
    let balance = state.private_balances
        .get(&addr)
        .copied()
        .unwrap_or(Uint128::zero());
    
    to_binary(&BalanceResponse { balance })
}
\end{lstlisting}

\begin{lstlisting}[language=Python, caption={Cliente para interactuar con Secret Network}]
from dataclasses import dataclass
from typing import Dict, Optional, List
import json
import hashlib
import secrets
from base64 import b64encode, b64decode

@dataclass
class SecretContract:
    """Representación de un contrato secreto."""
    address: str
    code_hash: str
    
@dataclass
class ViewingKey:
    """Viewing key para acceder a datos privados."""
    key: str
    address: str
    contract: str

class SecretNetworkClient:
    """
    Cliente para interactuar con Secret Network.
    """
    
    def __init__(self, endpoint: str, wallet_address: str, wallet_key: bytes):
        self.endpoint = endpoint
        self.address = wallet_address
        self.key = wallet_key
        
        # Almacén de viewing keys
        self.viewing_keys: Dict[str, ViewingKey] = {}
        
        # Cache de contratos
        self.contracts: Dict[str, SecretContract] = {}
    
    def encrypt_input(
        self,
        contract: SecretContract,
        msg: Dict
    ) -> bytes:
        """
        Encriptar input para el contrato.
        Los inputs se encriptan con la clave pública del enclave.
        """
        # En producción: usar clave pública del enclave
        # y ECDH para derivar clave de encriptación
        
        plaintext = json.dumps(msg).encode()
        
        # Placeholder: en realidad se usa x25519 + ChaCha20Poly1305
        nonce = secrets.token_bytes(12)
        
        # Simular encriptación
        encrypted = self._aead_encrypt(plaintext, contract.code_hash.encode(), nonce)
        
        return nonce + encrypted
    
    def decrypt_output(
        self,
        contract: SecretContract,
        encrypted_output: bytes
    ) -> Dict:
        """
        Descifrar output del contrato.
        """
        nonce = encrypted_output[:12]
        ciphertext = encrypted_output[12:]
        
        plaintext = self._aead_decrypt(ciphertext, contract.code_hash.encode(), nonce)
        
        return json.loads(plaintext)
    
    # ========== TRANSACCIONES ==========
    
    def execute_contract(
        self,
        contract_address: str,
        msg: Dict,
        funds: Optional[List[Dict]] = None
    ) -> Dict:
        """
        Ejecutar función del contrato con input encriptado.
        """
        contract = self._get_contract(contract_address)
        
        # Encriptar mensaje
        encrypted_msg = self.encrypt_input(contract, msg)
        
        # Construir transacción
        tx = {
            'type': 'wasm/MsgExecuteContract',
            'sender': self.address,
            'contract': contract_address,
            'msg': b64encode(encrypted_msg).decode(),
            'sent_funds': funds or []
        }
        
        # Firmar y enviar (simulado)
        result = self._broadcast_tx(tx)
        
        return result
    
    def transfer_private(
        self,
        contract_address: str,
        recipient: str,
        amount: int
    ) -> Dict:
        """
        Transferencia privada de tokens.
        """
        msg = {
            'transfer': {
                'recipient': recipient,
                'amount': str(amount)
            }
        }
        
        return self.execute_contract(contract_address, msg)
    
    # ========== VIEWING KEYS ==========
    
    def create_viewing_key(self, contract_address: str) -> ViewingKey:
        """
        Crear viewing key para un contrato.
        """
        entropy = secrets.token_hex(32)
        
        msg = {
            'create_viewing_key': {
                'entropy': entropy
            }
        }
        
        result = self.execute_contract(contract_address, msg)
        
        # Extraer key del resultado
        key_data = result.get('data', {})
        viewing_key = ViewingKey(
            key=key_data.get('key', ''),
            address=self.address,
            contract=contract_address
        )
        
        self.viewing_keys[contract_address] = viewing_key
        
        return viewing_key
    
    # ========== QUERIES ==========
    
    def query_contract(
        self,
        contract_address: str,
        msg: Dict,
        use_viewing_key: bool = False
    ) -> Dict:
        """
        Consultar contrato.
        """
        contract = self._get_contract(contract_address)
        
        # Añadir viewing key si es necesario
        if use_viewing_key and contract_address in self.viewing_keys:
            vk = self.viewing_keys[contract_address]
            msg['key'] = vk.key
            msg['address'] = vk.address
        
        # Encriptar query
        encrypted_query = self.encrypt_input(contract, msg)
        
        # Enviar query (simulado)
        encrypted_result = self._query(contract_address, encrypted_query)
        
        # Descifrar resultado
        return self.decrypt_output(contract, encrypted_result)
    
    def query_balance(self, contract_address: str) -> int:
        """
        Consultar balance privado (requiere viewing key).
        """
        vk = self.viewing_keys.get(contract_address)
        if not vk:
            raise ValueError("No viewing key for this contract")
        
        msg = {
            'balance': {
                'address': self.address,
                'key': vk.key
            }
        }
        
        result = self.query_contract(contract_address, msg)
        return int(result.get('balance', 0))
    
    def query_transaction_history(
        self,
        contract_address: str,
        page: int = 0,
        page_size: int = 10
    ) -> List[Dict]:
        """
        Consultar historial de transacciones privado.
        """
        vk = self.viewing_keys.get(contract_address)
        if not vk:
            raise ValueError("No viewing key for this contract")
        
        msg = {
            'transaction_history': {
                'address': self.address,
                'key': vk.key,
                'page': page,
                'page_size': page_size
            }
        }
        
        result = self.query_contract(contract_address, msg)
        return result.get('transactions', [])
    
    # ========== HELPERS ==========
    
    def _get_contract(self, address: str) -> SecretContract:
        """Obtener info del contrato."""
        if address not in self.contracts:
            # En producción: query a la blockchain
            self.contracts[address] = SecretContract(
                address=address,
                code_hash=hashlib.sha256(address.encode()).hexdigest()
            )
        return self.contracts[address]
    
    def _aead_encrypt(self, plaintext: bytes, aad: bytes, nonce: bytes) -> bytes:
        """Encriptación AEAD (placeholder)."""
        # En producción: ChaCha20Poly1305
        key = hashlib.sha256(aad + nonce).digest()
        return bytes(p ^ k for p, k in zip(plaintext, key * (len(plaintext) // 32 + 1)))
    
    def _aead_decrypt(self, ciphertext: bytes, aad: bytes, nonce: bytes) -> bytes:
        """Descifrado AEAD."""
        return self._aead_encrypt(ciphertext, aad, nonce)  # XOR es simétrico
    
    def _broadcast_tx(self, tx: Dict) -> Dict:
        """Broadcast transacción (simulado)."""
        return {'status': 'success', 'tx_hash': secrets.token_hex(32)}
    
    def _query(self, contract: str, encrypted_query: bytes) -> bytes:
        """Query al contrato (simulado)."""
        # Simular respuesta encriptada
        response = json.dumps({'balance': '1000000'}).encode()
        nonce = secrets.token_bytes(12)
        contract_info = self._get_contract(contract)
        encrypted = self._aead_encrypt(response, contract_info.code_hash.encode(), nonce)
        return nonce + encrypted


def secret_network_demo():
    print("=== Secret Network Demo ===\n")
    
    # Crear cliente
    client = SecretNetworkClient(
        endpoint="https://lcd.secret.network",
        wallet_address="secret1abc...",
        wallet_key=secrets.token_bytes(32)
    )
    
    contract_address = "secret1xyz..."
    
    print("--- Creating Viewing Key ---")
    vk = client.create_viewing_key(contract_address)
    print(f"Viewing key created: {vk.key[:20]}...")
    
    print("\n--- Private Transfer ---")
    result = client.transfer_private(
        contract_address,
        recipient="secret1recipient...",
        amount=1000000
    )
    print(f"Transfer result: {result['status']}")
    print("Note: Amount and recipient NOT visible on blockchain")
    
    print("\n--- Query Private Balance ---")
    # El balance solo es visible con viewing key
    try:
        balance = client.query_balance(contract_address)
        print(f"Balance: {balance}")
    except Exception as e:
        print(f"Query error: {e}")
    
    print("\n--- What's visible on blockchain ---")
    print("  ✓ Contract was called")
    print("  ✓ Who called it (sender address)")
    print("  ✓ Gas used")
    print("  ✗ Transfer amount")
    print("  ✗ Recipient")
    print("  ✗ User balances")

secret_network_demo()
\end{lstlisting}

%===========================================
\section{Oasis Network}
\label{sec:oasis-network}
%===========================================

\subsection{Arquitectura Separada}

\begin{definition}[Oasis Network]
Blockchain Layer-1 con arquitectura separada entre consenso y ejecución, permitiendo ParaTimes especializados incluyendo computación confidencial.

\textbf{Componentes:}
\begin{itemize}
    \item \textbf{Consensus Layer:} Proof-of-Stake con Tendermint BFT
    \item \textbf{ParaTimes:} Entornos de ejecución paralelos
    \item \textbf{Cipher ParaTime:} Ejecución confidencial con TEE
    \item \textbf{Sapphire ParaTime:} EVM confidencial
\end{itemize}

\textbf{Token nativo:} ROSE
\end{definition}

\begin{figure}[htbp]
\centering
\begin{tikzpicture}[
    node distance=0.8cm,
    layer/.style={rectangle, draw, minimum width=10cm, minimum height=1cm, align=center},
    paratime/.style={rectangle, draw, minimum width=2.2cm, minimum height=1.5cm, align=center, rounded corners}
]
    % Consensus Layer
    \node[layer, fill=blue!20] (consensus) at (0, 0) {\textbf{Consensus Layer} (Tendermint BFT)};
    
    % ParaTimes
    \node[paratime, fill=green!20] (emerald) at (-3.5, 2) {Emerald\\(EVM público)};
    \node[paratime, fill=purple!30] (sapphire) at (-1, 2) {Sapphire\\(EVM privado)};
    \node[paratime, fill=orange!20] (cipher) at (1.5, 2) {Cipher\\(Confidencial)};
    \node[paratime, fill=gray!20] (custom) at (4, 2) {Custom\\ParaTime};
    
    % Conexiones
    \draw[thick] (emerald) -- (-3.5, 0.5);
    \draw[thick] (sapphire) -- (-1, 0.5);
    \draw[thick] (cipher) -- (1.5, 0.5);
    \draw[thick] (custom) -- (4, 0.5);
    
    % TEE indicator
    \node at (-1, 3.2) [font=\small, purple] {TEE};
    \node at (1.5, 3.2) [font=\small, orange] {TEE};
    
    % Etiqueta
    \node at (0, -1) [font=\small] {Staking, Rewards, Slashing, Finality};
    \node at (0, 3.8) {\textbf{Oasis Network Architecture}};
\end{tikzpicture}
\caption{Arquitectura de Oasis: consenso separado de ejecución}
\label{fig:oasis-architecture}
\end{figure}

\subsection{Sapphire: EVM Confidencial}

\begin{definition}[Oasis Sapphire]
ParaTime compatible con EVM que proporciona computación confidencial:

\textbf{Características:}
\begin{itemize}
    \item Contratos Solidity estándar
    \item Estado encriptado por defecto
    \item Llamadas confidenciales opcionales
    \item Compatible con herramientas Ethereum (Hardhat, Remix)
    \item RNG seguro on-chain
\end{itemize}
\end{definition}

\begin{lstlisting}[language=Solidity, caption={Contrato confidencial en Oasis Sapphire}]
// SPDX-License-Identifier: MIT
pragma solidity ^0.8.20;

// Imports específicos de Sapphire
import "@oasisprotocol/sapphire-contracts/contracts/Sapphire.sol";

/**
 * @title PrivateAuction
 * @notice Subasta privada donde las pujas son confidenciales hasta el cierre
 */
contract PrivateAuction {
    // ========== ESTADO PRIVADO ==========
    // En Sapphire, el estado es encriptado por defecto
    
    address public owner;
    string public itemDescription;
    uint256 public auctionEndTime;
    bool public ended;
    
    // Pujas privadas (no visibles en blockchain)
    mapping(address => uint256) private bids;
    address[] private bidders;
    
    // Ganador (solo revelado después del cierre)
    address private winner;
    uint256 private winningBid;
    bool private winnerRevealed;
    
    // Compromiso de hash para verificación
    mapping(address => bytes32) private bidCommitments;
    
    // ========== EVENTOS ==========
    // Los eventos en Sapphire pueden ser confidenciales
    
    event BidPlaced(address indexed bidder); // No revela monto
    event AuctionEnded(address winner, uint256 amount);
    
    // ========== CONSTRUCTOR ==========
    
    constructor(string memory _item, uint256 _duration) {
        owner = msg.sender;
        itemDescription = _item;
        auctionEndTime = block.timestamp + _duration;
    }
    
    // ========== FUNCIONES PÚBLICAS ==========
    
    /**
     * @notice Realizar puja confidencial
     * @dev El monto se almacena encriptado, no visible en mempool ni estado
     */
    function placeBid() external payable {
        require(block.timestamp < auctionEndTime, "Auction ended");
        require(msg.value > 0, "Bid must be positive");
        
        // Si ya tiene puja, añadir al total
        if (bids[msg.sender] == 0) {
            bidders.push(msg.sender);
        }
        
        bids[msg.sender] += msg.value;
        
        // Crear commitment para verificación posterior
        bidCommitments[msg.sender] = keccak256(
            abi.encodePacked(msg.sender, bids[msg.sender], getRandomBytes32())
        );
        
        emit BidPlaced(msg.sender);
    }
    
    /**
     * @notice Revelar ganador después del cierre
     * @dev Solo el propietario puede revelar
     */
    function revealWinner() external {
        require(msg.sender == owner, "Only owner");
        require(block.timestamp >= auctionEndTime, "Auction not ended");
        require(!winnerRevealed, "Already revealed");
        
        // Encontrar la puja más alta
        uint256 highestBid = 0;
        address highestBidder = address(0);
        
        for (uint256 i = 0; i < bidders.length; i++) {
            if (bids[bidders[i]] > highestBid) {
                highestBid = bids[bidders[i]];
                highestBidder = bidders[i];
            }
        }
        
        winner = highestBidder;
        winningBid = highestBid;
        winnerRevealed = true;
        ended = true;
        
        emit AuctionEnded(winner, winningBid);
    }
    
    /**
     * @notice Reclamar fondos si no ganaste
     */
    function claimRefund() external {
        require(ended, "Auction not ended");
        require(msg.sender != winner, "Winner cannot refund");
        
        uint256 amount = bids[msg.sender];
        require(amount > 0, "No bid to refund");
        
        bids[msg.sender] = 0;
        payable(msg.sender).transfer(amount);
    }
    
    // ========== FUNCIONES DE CONSULTA PRIVADA ==========
    
    /**
     * @notice Ver tu propia puja (solo tú puedes ver)
     * @dev En Sapphire, msg.sender es verificado en el enclave
     */
    function getMyBid() external view returns (uint256) {
        return bids[msg.sender];
    }
    
    /**
     * @notice Obtener número de participantes (público)
     */
    function getBidderCount() external view returns (uint256) {
        return bidders.length;
    }
    
    /**
     * @notice Obtener ganador (solo después de revelar)
     */
    function getWinner() external view returns (address, uint256) {
        require(winnerRevealed, "Winner not revealed yet");
        return (winner, winningBid);
    }
    
    // ========== FUNCIONES AUXILIARES SAPPHIRE ==========
    
    /**
     * @notice Generar bytes aleatorios seguros
     * @dev Sapphire proporciona RNG seguro en el enclave
     */
    function getRandomBytes32() internal view returns (bytes32) {
        return Sapphire.randomBytes(32, abi.encodePacked(block.timestamp));
    }
}

/**
 * @title PrivateVoting
 * @notice Sistema de votación con votos confidenciales
 */
contract PrivateVoting {
    struct Proposal {
        string description;
        uint256 yesVotes;
        uint256 noVotes;
        uint256 deadline;
        bool finalized;
    }
    
    mapping(uint256 => Proposal) public proposals;
    uint256 public proposalCount;
    
    // Votos privados
    mapping(uint256 => mapping(address => bool)) private hasVoted;
    mapping(uint256 => mapping(address => bool)) private votes; // true = yes
    
    // Lista de votantes elegibles
    mapping(address => bool) public eligibleVoters;
    
    event ProposalCreated(uint256 indexed proposalId, string description);
    event VoteCast(uint256 indexed proposalId, address voter); // No revela voto
    event ProposalFinalized(uint256 indexed proposalId, bool passed);
    
    constructor(address[] memory _voters) {
        for (uint256 i = 0; i < _voters.length; i++) {
            eligibleVoters[_voters[i]] = true;
        }
    }
    
    function createProposal(string memory description, uint256 duration) external {
        proposalCount++;
        proposals[proposalCount] = Proposal({
            description: description,
            yesVotes: 0,
            noVotes: 0,
            deadline: block.timestamp + duration,
            finalized: false
        });
        
        emit ProposalCreated(proposalCount, description);
    }
    
    /**
     * @notice Votar en propuesta (voto privado)
     * @param proposalId ID de la propuesta
     * @param vote true = sí, false = no
     */
    function vote(uint256 proposalId, bool vote) external {
        require(eligibleVoters[msg.sender], "Not eligible");
        require(!hasVoted[proposalId][msg.sender], "Already voted");
        require(block.timestamp < proposals[proposalId].deadline, "Voting ended");
        
        hasVoted[proposalId][msg.sender] = true;
        votes[proposalId][msg.sender] = vote;
        
        // Incrementar conteo (pero no se revela quién votó qué)
        if (vote) {
            proposals[proposalId].yesVotes++;
        } else {
            proposals[proposalId].noVotes++;
        }
        
        emit VoteCast(proposalId, msg.sender);
    }
    
    /**
     * @notice Verificar si votaste (solo tú puedes ver tu voto)
     */
    function getMyVote(uint256 proposalId) external view returns (bool voted, bool vote) {
        return (hasVoted[proposalId][msg.sender], votes[proposalId][msg.sender]);
    }
    
    /**
     * @notice Finalizar propuesta y revelar resultado
     */
    function finalize(uint256 proposalId) external {
        Proposal storage prop = proposals[proposalId];
        require(block.timestamp >= prop.deadline, "Voting not ended");
        require(!prop.finalized, "Already finalized");
        
        prop.finalized = true;
        bool passed = prop.yesVotes > prop.noVotes;
        
        emit ProposalFinalized(proposalId, passed);
    }
}
\end{lstlisting}

\begin{lstlisting}[language=JavaScript, caption={Deployment y uso de contratos Sapphire con Hardhat}]
// hardhat.config.js para Oasis Sapphire
require("@oasisprotocol/sapphire-hardhat");

module.exports = {
  solidity: "0.8.20",
  networks: {
    sapphire: {
      url: "https://sapphire.oasis.io",
      accounts: [process.env.PRIVATE_KEY],
      chainId: 0x5afe,
    },
    sapphire_testnet: {
      url: "https://testnet.sapphire.oasis.dev",
      accounts: [process.env.PRIVATE_KEY],
      chainId: 0x5aff,
    },
  },
};

// deploy.js
const { ethers } = require("hardhat");

async function main() {
  console.log("Deploying to Oasis Sapphire...");
  
  // Deploy subasta privada
  const PrivateAuction = await ethers.getContractFactory("PrivateAuction");
  const auction = await PrivateAuction.deploy(
    "Rare NFT Collection",
    86400 // 1 día
  );
  
  await auction.waitForDeployment();
  console.log("PrivateAuction deployed to:", await auction.getAddress());
  
  // Deploy votación
  const [owner, voter1, voter2, voter3] = await ethers.getSigners();
  
  const PrivateVoting = await ethers.getContractFactory("PrivateVoting");
  const voting = await PrivateVoting.deploy([
    voter1.address,
    voter2.address,
    voter3.address
  ]);
  
  await voting.waitForDeployment();
  console.log("PrivateVoting deployed to:", await voting.getAddress());
}

// interact.js - Interacción confidencial
async function interactWithAuction() {
  const [bidder1, bidder2] = await ethers.getSigners();
  
  const auction = await ethers.getContractAt(
    "PrivateAuction",
    "0x..." // contract address
  );
  
  // Realizar pujas confidenciales
  // Las pujas NO son visibles en la blockchain ni en el mempool
  
  console.log("Bidder 1 placing confidential bid...");
  const tx1 = await auction.connect(bidder1).placeBid({
    value: ethers.parseEther("1.5")
  });
  await tx1.wait();
  
  console.log("Bidder 2 placing confidential bid...");
  const tx2 = await auction.connect(bidder2).placeBid({
    value: ethers.parseEther("2.0")
  });
  await tx2.wait();
  
  // Cada usuario solo puede ver su propia puja
  const myBid1 = await auction.connect(bidder1).getMyBid();
  console.log(`Bidder 1's bid: ${ethers.formatEther(myBid1)} ETH`);
  
  // Intentar ver puja de otro usuario falla
  // (Sapphire verifica msg.sender en el enclave)
  
  // El número de participantes es público
  const count = await auction.getBidderCount();
  console.log(`Total bidders: ${count}`);
  
  // Pero los montos individuales son privados hasta el cierre
  console.log("Individual bid amounts: [ENCRYPTED]");
}

main().catch(console.error);
\end{lstlisting}

\subsection{Comparación Secret Network vs Oasis}

\begin{table}[htbp]
\centering
\small
\begin{tabular}{lp{4.5cm}p{4.5cm}}
\toprule
\textbf{Aspecto} & \textbf{Secret Network} & \textbf{Oasis Sapphire} \\
\midrule
\textbf{Lenguaje} & Rust (CosmWasm modificado) & Solidity estándar \\
\textbf{Compatibilidad} & Cosmos ecosystem (IBC) & Ethereum ecosystem (EVM) \\
\textbf{TEE} & Intel SGX & Intel SGX \\
\textbf{Privacidad} & Todo privado por defecto & Configurable por función \\
\textbf{Viewing Keys} & Sí, sistema elaborado & Limitado (msg.sender check) \\
\textbf{Tooling} & SecretJS, custom tools & Hardhat, Remix, Ethers.js \\
\textbf{Madurez} & Más maduro (2020) & Más reciente (2022) \\
\textbf{DeFi} & Secret DeFi (ShadeSwap) & Emergente \\
\bottomrule
\end{tabular}
\caption{Comparación entre Secret Network y Oasis Sapphire}
\end{table}
%===========================================
\section{Aleo}
\label{sec:aleo}
%===========================================

\subsection{ZK-Native Blockchain}

\begin{definition}[Aleo]
Blockchain nativa en zero-knowledge donde toda la ejecución de programas genera pruebas ZK, sin depender de hardware de confianza.

\textbf{Características principales:}
\begin{itemize}
    \item \textbf{Modelo de Records:} Similar a notes de Zcash
    \item \textbf{Lenguaje Leo:} DSL similar a Rust para ZK
    \item \textbf{snarkVM:} Máquina virtual que ejecuta y genera pruebas
    \item \textbf{Sin TEE:} Privacidad puramente criptográfica
    \item \textbf{Token nativo:} ALEO
\end{itemize}
\end{definition}

\begin{keyidea}
\textbf{Diferencia fundamental: ZK vs TEE}

\textbf{TEE (Secret, Oasis):}
\begin{itemize}
    \item Estado encriptado almacenado en blockchain
    \item Ejecución en hardware de confianza
    \item Vulnerable a ataques al hardware
    \item Baja latencia de ejecución
\end{itemize}

\textbf{ZK (Aleo):}
\begin{itemize}
    \item No hay estado en blockchain (solo commitments)
    \item Ejecución local, verificación on-chain
    \item Seguridad puramente matemática
    \item Mayor costo computacional (generación de proof)
\end{itemize}
\end{keyidea}

\begin{figure}[htbp]
\centering
\begin{tikzpicture}[
    node distance=1cm,
    component/.style={rectangle, draw, rounded corners, minimum width=2.5cm, minimum height=1cm, align=center},
    arrow/.style={->, thick}
]
    % Usuario
    \node[component, fill=blue!20] (user) at (0, 0) {Usuario\\(Leo Program)};
    
    % snarkVM
    \node[component, fill=green!20] (vm) at (4, 0) {snarkVM\\(Ejecuta + Proof)};
    
    % Blockchain
    \node[component, fill=orange!20] (chain) at (8, 0) {Aleo\\Blockchain};
    
    % Records
    \node[component, fill=purple!20] (records) at (0, -2.5) {Records\\(Privados)};
    
    % Transición
    \node[component, fill=yellow!20] (transition) at (4, -2.5) {Transition\\(Proof + nullifiers)};
    
    % Estado on-chain
    \node[component, fill=red!20] (onchain) at (8, -2.5) {Commitments\\(Públicos)};
    
    % Flechas
    \draw[arrow] (user) -- node[above, font=\small] {ejecutar} (vm);
    \draw[arrow] (vm) -- node[above, font=\small] {submit} (chain);
    
    \draw[arrow] (records) -- (vm);
    \draw[arrow] (vm) -- (transition);
    \draw[arrow] (transition) -- (onchain);
    
    % Labels
    \node at (0, -3.5) [font=\small] {Off-chain};
    \node at (4, -3.5) [font=\small] {Local};
    \node at (8, -3.5) [font=\small] {On-chain};
\end{tikzpicture}
\caption{Arquitectura de Aleo: ejecución off-chain, verificación on-chain}
\label{fig:aleo-architecture}
\end{figure}

\subsection{Modelo de Records}

\begin{definition}[Record en Aleo]
Estructura de datos que representa estado privado, similar a las notes de Zcash:
\[
\text{Record} = \{
\text{owner}, \text{data}, \text{nonce}
\}
\]

\textbf{Propiedades:}
\begin{itemize}
    \item \textbf{Propiedad:} Solo el owner puede consumir el record
    \item \textbf{Unicidad:} Cada record tiene nonce único
    \item \textbf{Consumo:} Records se ``gastan'' al usarse (modelo UTXO)
    \item \textbf{Privacidad:} Contenido encriptado al owner
\end{itemize}
\end{definition}

\begin{mathnote}[Transiciones de Estado en Aleo]
Una transición consume records de entrada y produce records de salida:

\[
\text{Transition}: \{R_{\text{in}_1}, ..., R_{\text{in}_n}\} \rightarrow \{R_{\text{out}_1}, ..., R_{\text{out}_m}\}
\]

\textbf{Lo que se publica on-chain:}
\begin{itemize}
    \item Nullifiers de records consumidos: $\{\text{nf}(R_{\text{in}_i})\}$
    \item Commitments de records creados: $\{\text{cm}(R_{\text{out}_j})\}$
    \item ZK Proof de ejecución correcta
\end{itemize}

\textbf{Lo que permanece privado:}
\begin{itemize}
    \item Contenido de records (owner, data)
    \item Mapping entre inputs y outputs
    \item Lógica específica ejecutada
\end{itemize}
\end{mathnote}

\subsection{Lenguaje Leo}

\begin{lstlisting}[language=Rust, caption={Programa Leo: Token privado}]
// token.leo - Token privado en Aleo

program token.aleo {
    // Record que representa balance de token
    record Token {
        owner: address,
        amount: u64,
    }
    
    // Mapping público para supply total (opcional)
    mapping total_supply: u8 => u64;
    
    // ========== FUNCIONES DE TRANSICIÓN ==========
    
    // Mint: crear nuevos tokens (solo admin)
    transition mint_private(
        receiver: address,
        amount: u64
    ) -> Token {
        // Verificar autorización (simplificado)
        assert(self.caller == aleo1admin...);
        
        // Crear record para el receptor
        return Token {
            owner: receiver,
            amount: amount,
        };
    }
    
    // Transfer: transferencia privada
    transition transfer_private(
        token: Token,           // Record de entrada
        receiver: address,
        amount: u64
    ) -> (Token, Token) {      // Records de salida
        // Verificar que el caller es dueño
        assert_eq(token.owner, self.caller);
        
        // Verificar balance suficiente
        assert(token.amount >= amount);
        
        // Calcular cambio
        let change: u64 = token.amount - amount;
        
        // Record para el receptor
        let recipient_token: Token = Token {
            owner: receiver,
            amount: amount,
        };
        
        // Record de cambio para el sender
        let sender_token: Token = Token {
            owner: self.caller,
            amount: change,
        };
        
        // Retornar ambos records
        // El record de entrada se consume (nullifier publicado)
        return (recipient_token, sender_token);
    }
    
    // Split: dividir un record en dos
    transition split(
        token: Token,
        amount1: u64
    ) -> (Token, Token) {
        assert_eq(token.owner, self.caller);
        assert(token.amount >= amount1);
        
        let amount2: u64 = token.amount - amount1;
        
        let token1: Token = Token {
            owner: self.caller,
            amount: amount1,
        };
        
        let token2: Token = Token {
            owner: self.caller,
            amount: amount2,
        };
        
        return (token1, token2);
    }
    
    // Join: combinar dos records en uno
    transition join(
        token1: Token,
        token2: Token
    ) -> Token {
        assert_eq(token1.owner, self.caller);
        assert_eq(token2.owner, self.caller);
        
        let total: u64 = token1.amount + token2.amount;
        
        return Token {
            owner: self.caller,
            amount: total,
        };
    }
}
\end{lstlisting}

\begin{lstlisting}[language=Rust, caption={Leo: Subasta privada con ZK}]
// private_auction.leo - Subasta donde las pujas son privadas

program private_auction.aleo {
    // Record de puja (privado)
    record Bid {
        owner: address,
        auction_id: field,
        amount: u64,
        timestamp: u64,
    }
    
    // Record de ganador (revelado solo al ganador)
    record WinnerClaim {
        owner: address,
        auction_id: field,
        winning_amount: u64,
    }
    
    // Estado público de la subasta
    struct AuctionState {
        item_hash: field,
        end_time: u64,
        highest_bid_commitment: field,
        bid_count: u64,
        is_finalized: bool,
    }
    
    mapping auctions: field => AuctionState;
    mapping bid_commitments: field => bool;  // commitment => exists
    
    // ========== CREAR SUBASTA ==========
    
    transition create_auction(
        public item_description: field,
        public duration: u64
    ) -> field {
        // Generar ID único
        let auction_id: field = BHP256::hash_to_field(
            item_description + self.caller.to_field() + duration as field
        );
        
        return auction_id then finalize(auction_id, item_description, duration);
    }
    
    finalize create_auction(
        auction_id: field,
        item_hash: field,
        duration: u64
    ) {
        let state: AuctionState = AuctionState {
            item_hash: item_hash,
            end_time: block.height + duration,
            highest_bid_commitment: 0field,
            bid_count: 0u64,
            is_finalized: false,
        };
        
        auctions.set(auction_id, state);
    }
    
    // ========== REALIZAR PUJA PRIVADA ==========
    
    transition place_bid(
        auction_id: field,
        amount: u64,
        secret: field  // Salt para el commitment
    ) -> Bid {
        // Crear commitment de la puja (público)
        let bid_commitment: field = Poseidon2::hash_to_field(
            self.caller.to_field() + amount as field + secret
        );
        
        // Crear record privado de la puja
        let bid: Bid = Bid {
            owner: self.caller,
            auction_id: auction_id,
            amount: amount,
            timestamp: 0u64,  // Se actualiza en finalize
        };
        
        return bid then finalize(auction_id, bid_commitment);
    }
    
    finalize place_bid(
        auction_id: field,
        bid_commitment: field
    ) {
        // Verificar que la subasta existe y está activa
        let state: AuctionState = auctions.get(auction_id);
        assert(!state.is_finalized);
        assert(block.height < state.end_time);
        
        // Registrar commitment
        bid_commitments.set(bid_commitment, true);
        
        // Actualizar conteo
        let updated_state: AuctionState = AuctionState {
            item_hash: state.item_hash,
            end_time: state.end_time,
            highest_bid_commitment: state.highest_bid_commitment,
            bid_count: state.bid_count + 1u64,
            is_finalized: false,
        };
        
        auctions.set(auction_id, updated_state);
    }
    
    // ========== REVELAR Y RECLAMAR ==========
    
    // El ganador prueba que tiene la puja más alta
    transition claim_winner(
        bid: Bid,
        secret: field,
        all_other_bids: [u64; 10]  // Otras pujas (debe probar que la suya es mayor)
    ) -> WinnerClaim {
        // Verificar que todas las otras pujas son menores
        for i: u8 in 0u8..10u8 {
            assert(bid.amount > all_other_bids[i] || all_other_bids[i] == 0u64);
        }
        
        return WinnerClaim {
            owner: bid.owner,
            auction_id: bid.auction_id,
            winning_amount: bid.amount,
        } then finalize(bid.auction_id);
    }
    
    finalize claim_winner(auction_id: field) {
        let state: AuctionState = auctions.get(auction_id);
        assert(block.height >= state.end_time);
        assert(!state.is_finalized);
        
        let final_state: AuctionState = AuctionState {
            item_hash: state.item_hash,
            end_time: state.end_time,
            highest_bid_commitment: state.highest_bid_commitment,
            bid_count: state.bid_count,
            is_finalized: true,
        };
        
        auctions.set(auction_id, final_state);
    }
}
\end{lstlisting}

%===========================================
\section{Mina Protocol}
\label{sec:mina-protocol}
%===========================================

\subsection{Blockchain Sucinta}

\begin{definition}[Mina Protocol]
Blockchain con tamaño constante de ~22 KB gracias a recursive SNARKs, donde toda la historia de la cadena se comprime en una sola prueba.

\textbf{Características:}
\begin{itemize}
    \item \textbf{Tamaño constante:} Siempre ~22 KB independiente del historial
    \item \textbf{Recursive SNARKs:} Pruebas que verifican otras pruebas
    \item \textbf{zkApps:} Smart contracts con privacidad
    \item \textbf{Off-chain execution:} Estado off-chain, proofs on-chain
    \item \textbf{o1js:} SDK TypeScript/JavaScript para desarrollo
\end{itemize}
\end{definition}

\begin{figure}[htbp]
\centering
\begin{tikzpicture}[
    node distance=0.5cm,
    block/.style={rectangle, draw, minimum width=1.5cm, minimum height=1cm, fill=blue!20},
    proof/.style={ellipse, draw, minimum width=1.2cm, minimum height=0.6cm, fill=green!20}
]
    % Blockchain tradicional
    \node at (0, 3) {\textbf{Blockchain Tradicional}};
    
    \foreach \x in {0, 1.5, 3, 4.5, 6, 7.5} {
        \node[block] at (\x, 2) {};
    }
    \node at (9, 2) {...};
    \node at (5, 0.8) [font=\small] {Crece con cada bloque: GB → TB};
    
    % Mina
    \node at (0, -1) {\textbf{Mina Protocol}};
    
    % Bloques que se comprimen
    \node[block, fill=gray!30] at (0, -2) {B1};
    \node at (0.75, -2) {+};
    \node[block, fill=gray!30] at (1.5, -2) {B2};
    \node at (2.2, -2) {→};
    \node[proof] at (3.2, -2) {π};
    \node at (3.9, -2) {+};
    \node[block, fill=gray!30] at (4.7, -2) {B3};
    \node at (5.4, -2) {→};
    \node[proof] at (6.4, -2) {π'};
    
    % Resultado final
    \node at (8, -2) {=};
    \node[proof, minimum width=1.5cm, fill=green!40] at (9.5, -2) {22 KB};
    
    \node at (5, -3.2) [font=\small] {Tamaño constante: siempre 22 KB};
\end{tikzpicture}
\caption{Mina comprime toda la blockchain en una prueba de 22 KB}
\label{fig:mina-succinct}
\end{figure}

\begin{mathnote}[Recursive SNARKs en Mina]
\textbf{Idea clave:} Un SNARK puede verificar otro SNARK.

\textbf{Prueba de estado de blockchain:}
\[
\pi_{n} = \text{SNARK.Prove}\left(
\pi_{n-1} \text{ es válida} \land B_n \text{ es bloque válido}
\right)
\]

\textbf{Verificación:}
\begin{itemize}
    \item Verificar $\pi_n$ toma tiempo constante
    \item $\pi_n$ implica que TODOS los bloques anteriores son válidos
    \item No necesitas descargar la historia completa
\end{itemize}

\textbf{Implicaciones:}
\begin{itemize}
    \item Light clients verdaderos (22 KB = estado completo)
    \item Sincronización instantánea
    \item Escalabilidad extrema para verificadores
\end{itemize}
\end{mathnote}

\subsection{zkApps: Smart Contracts con ZK}

\begin{definition}[zkApp]
Smart contract en Mina donde:
\begin{itemize}
    \item Estado almacenado off-chain (Merkle tree)
    \item Solo root del estado on-chain
    \item Transiciones de estado prueban correctitud con ZK
    \item Ejecución local, verificación on-chain
\end{itemize}
\end{definition}

\begin{lstlisting}[language=TypeScript, caption={zkApp básico con o1js}]
import {
  Field,
  SmartContract,
  state,
  State,
  method,
  Poseidon,
  PublicKey,
  Signature,
  UInt64,
  Struct,
  MerkleWitness,
  Bool,
  CircuitString,
  Provable,
} from 'o1js';

// Altura del Merkle tree para balances
const TREE_HEIGHT = 8;
class BalanceWitness extends MerkleWitness(TREE_HEIGHT) {}

// Estructura para balance de usuario (off-chain)
class UserBalance extends Struct({
  publicKey: PublicKey,
  balance: UInt64,
  nonce: Field,
}) {
  hash(): Field {
    return Poseidon.hash([
      ...this.publicKey.toFields(),
      this.balance.value,
      this.nonce,
    ]);
  }
}

/**
 * PrivateToken: Token con balances privados
 * - Balances almacenados off-chain en Merkle tree
 * - Solo el root está on-chain
 * - Transferencias prueban actualización correcta del tree
 */
class PrivateToken extends SmartContract {
  // Root del Merkle tree de balances (on-chain)
  @state(Field) balancesRoot = State<Field>();
  
  // Supply total (público)
  @state(UInt64) totalSupply = State<UInt64>();
  
  // Nullifier set root (para prevenir double-spending)
  @state(Field) nullifiersRoot = State<Field>();

  init() {
    super.init();
    // Inicializar con árbol vacío
    this.balancesRoot.set(Field(0));
    this.totalSupply.set(UInt64.zero);
    this.nullifiersRoot.set(Field(0));
  }

  /**
   * Transferencia privada
   * - Sender y receiver prueban posesión de sus balances
   * - Se actualiza el root del tree
   * - Montos y direcciones no visibles on-chain
   */
  @method async transfer(
    // Datos del sender (privados)
    senderBalance: UserBalance,
    senderWitness: BalanceWitness,
    senderSignature: Signature,
    
    // Datos del receiver (privados)
    receiverBalance: UserBalance,
    receiverWitness: BalanceWitness,
    
    // Monto (privado)
    amount: UInt64,
    
    // Nuevo balance sender
    newSenderBalance: UserBalance,
    
    // Nuevo balance receiver  
    newReceiverBalance: UserBalance
  ) {
    // 1. Verificar root actual
    const currentRoot = this.balancesRoot.get();
    this.balancesRoot.requireEquals(currentRoot);
    
    // 2. Verificar que sender tiene el balance que dice
    const senderLeafHash = senderBalance.hash();
    const computedSenderRoot = senderWitness.calculateRoot(senderLeafHash);
    computedSenderRoot.assertEquals(currentRoot);
    
    // 3. Verificar firma del sender
    senderSignature.verify(
      senderBalance.publicKey,
      [amount.value, senderBalance.nonce]
    ).assertTrue();
    
    // 4. Verificar balance suficiente
    senderBalance.balance.assertGreaterThanOrEqual(amount);
    
    // 5. Verificar nuevos balances correctos
    newSenderBalance.balance.assertEquals(
      senderBalance.balance.sub(amount)
    );
    newReceiverBalance.balance.assertEquals(
      receiverBalance.balance.add(amount)
    );
    
    // 6. Verificar nonce incrementado
    newSenderBalance.nonce.assertEquals(senderBalance.nonce.add(1));
    
    // 7. Calcular nuevo root con balances actualizados
    // (Simplificado - en realidad se actualizan ambas hojas)
    const newSenderLeafHash = newSenderBalance.hash();
    const intermediateRoot = senderWitness.calculateRoot(newSenderLeafHash);
    
    const newReceiverLeafHash = newReceiverBalance.hash();
    const newRoot = receiverWitness.calculateRoot(newReceiverLeafHash);
    
    // 8. Actualizar root on-chain
    this.balancesRoot.set(newRoot);
  }

  /**
   * Mint privado (solo admin)
   */
  @method async mint(
    receiver: PublicKey,
    amount: UInt64,
    currentReceiverBalance: UserBalance,
    witness: BalanceWitness,
    newReceiverBalance: UserBalance
  ) {
    // Verificar admin (simplificado)
    this.sender.getAndRequireSignature().assertEquals(this.address);
    
    // Verificar balance actual en tree
    const currentRoot = this.balancesRoot.get();
    this.balancesRoot.requireEquals(currentRoot);
    
    const leafHash = currentReceiverBalance.hash();
    witness.calculateRoot(leafHash).assertEquals(currentRoot);
    
    // Verificar nuevo balance correcto
    newReceiverBalance.balance.assertEquals(
      currentReceiverBalance.balance.add(amount)
    );
    
    // Actualizar root
    const newLeafHash = newReceiverBalance.hash();
    const newRoot = witness.calculateRoot(newLeafHash);
    this.balancesRoot.set(newRoot);
    
    // Actualizar supply
    const currentSupply = this.totalSupply.get();
    this.totalSupply.requireEquals(currentSupply);
    this.totalSupply.set(currentSupply.add(amount));
  }
}

/**
 * PrivateVoting: Votación con votos confidenciales
 */
class PrivateVoting extends SmartContract {
  // Merkle root de votantes elegibles
  @state(Field) votersRoot = State<Field>();
  
  // Conteos de votos (públicos, pero no quién votó qué)
  @state(UInt64) yesVotes = State<UInt64>();
  @state(UInt64) noVotes = State<UInt64>();
  
  // Nullifier root (para prevenir doble voto)
  @state(Field) nullifiersRoot = State<Field>();

  @method async vote(
    // Prueba de elegibilidad (privada)
    voterPublicKey: PublicKey,
    eligibilityWitness: BalanceWitness,
    
    // Voto (privado)
    voteValue: Bool, // true = yes, false = no
    
    // Nullifier (para prevenir doble voto)
    nullifier: Field,
    nullifierWitness: BalanceWitness
  ) {
    // 1. Verificar elegibilidad
    const votersRoot = this.votersRoot.get();
    this.votersRoot.requireEquals(votersRoot);
    
    const voterHash = Poseidon.hash(voterPublicKey.toFields());
    eligibilityWitness.calculateRoot(voterHash).assertEquals(votersRoot);
    
    // 2. Verificar que no ha votado (nullifier no existe)
    const nullifiersRoot = this.nullifiersRoot.get();
    this.nullifiersRoot.requireEquals(nullifiersRoot);
    
    // El nullifier debe llevar a hoja vacía
    nullifierWitness.calculateRoot(Field(0)).assertEquals(nullifiersRoot);
    
    // 3. Calcular nuevo nullifier root
    const newNullifiersRoot = nullifierWitness.calculateRoot(nullifier);
    this.nullifiersRoot.set(newNullifiersRoot);
    
    // 4. Actualizar conteos
    const currentYes = this.yesVotes.get();
    const currentNo = this.noVotes.get();
    this.yesVotes.requireEquals(currentYes);
    this.noVotes.requireEquals(currentNo);
    
    // Increment condicionalmente basado en voto
    const newYes = Provable.if(voteValue, currentYes.add(1), currentYes);
    const newNo = Provable.if(voteValue, currentNo, currentNo.add(1));
    
    this.yesVotes.set(newYes);
    this.noVotes.set(newNo);
  }
}

// Ejemplo de uso
async function main() {
  console.log('Compiling PrivateToken zkApp...');
  const { verificationKey } = await PrivateToken.compile();
  
  console.log('Verification key size:', verificationKey.data.length);
  console.log('This proof verifies the entire contract logic!');
}
\end{lstlisting}

%===========================================
\section{Fhenix y FHE}
\label{sec:fhenix-fhe}
%===========================================

\subsection{Fully Homomorphic Encryption}

\begin{definition}[Fully Homomorphic Encryption (FHE)]
Sistema de encriptación que permite realizar operaciones sobre datos cifrados sin descifrarlos:

\[
\text{Dec}(f(\text{Enc}(x), \text{Enc}(y))) = f(x, y)
\]

\textbf{Operaciones homomórficas:}
\begin{itemize}
    \item \textbf{Suma:} $\text{Enc}(x) \oplus \text{Enc}(y) = \text{Enc}(x + y)$
    \item \textbf{Multiplicación:} $\text{Enc}(x) \otimes \text{Enc}(y) = \text{Enc}(x \times y)$
    \item Con suma y multiplicación, se puede computar cualquier función
\end{itemize}
\end{definition}

\begin{figure}[htbp]
\centering
\begin{tikzpicture}[
    node distance=1cm,
    box/.style={rectangle, draw, minimum width=2cm, minimum height=1cm, align=center},
    arrow/.style={->, thick}
]
    % Sin FHE
    \node at (-3, 3) {\textbf{Sin FHE}};
    
    \node[box, fill=blue!20] (data1) at (-5, 1.5) {Datos\\privados};
    \node[box, fill=yellow!20] (compute1) at (-3, 1.5) {Computar\\(visible)};
    \node[box, fill=green!20] (result1) at (-1, 1.5) {Resultado};
    
    \draw[arrow] (data1) -- (compute1);
    \draw[arrow] (compute1) -- (result1);
    
    \node at (-3, 0.5) [font=\small, red] {¡Datos expuestos!};
    
    % Con FHE
    \node at (5, 3) {\textbf{Con FHE}};
    
    \node[box, fill=blue!20] (data2) at (2, 1.5) {Enc(Datos)};
    \node[box, fill=yellow!20] (compute2) at (5, 1.5) {Computar\\(encriptado)};
    \node[box, fill=green!20] (result2) at (8, 1.5) {Enc(Resultado)};
    
    \draw[arrow] (data2) -- (compute2);
    \draw[arrow] (compute2) -- (result2);
    
    \node at (5, 0.5) [font=\small, green!60!black] {Datos siempre privados};
    
    % Decrypt final
    \node[box, fill=purple!20] (dec) at (8, -0.5) {Decrypt};
    \draw[arrow] (result2) -- (dec);
\end{tikzpicture}
\caption{FHE permite computar sobre datos encriptados}
\label{fig:fhe-concept}
\end{figure}

\begin{mathnote}[Esquemas FHE Modernos]
\textbf{1. TFHE (Torus FHE):}
\begin{itemize}
    \item Operaciones sobre bits individuales
    \item Bootstrap rápido (~10ms)
    \item Usado por Fhenix
\end{itemize}

\textbf{2. BGV/BFV:}
\begin{itemize}
    \item Operaciones sobre enteros
    \item Más eficiente para aritmética
    \item Usado en muchas implementaciones
\end{itemize}

\textbf{3. CKKS:}
\begin{itemize}
    \item Operaciones sobre números reales (aproximados)
    \item Útil para machine learning
    \item Menor precisión pero más flexible
\end{itemize}

\textbf{Limitaciones actuales:}
\begin{itemize}
    \item Muy costoso computacionalmente (1000-10000x más lento)
    \item Ciphertexts grandes (KB por dato simple)
    \item ``Noise'' acumulativo requiere bootstrapping
\end{itemize}
\end{mathnote}

\begin{definition}[Fhenix]
Blockchain Layer-2 que integra FHE para smart contracts completamente encriptados:

\textbf{Características:}
\begin{itemize}
    \item Estado del contrato encriptado
    \item Computación sobre datos encriptados
    \item Sin necesidad de descifrar durante ejecución
    \item Basado en TFHE
    \item EVM compatible (modificado)
\end{itemize}
\end{definition}

\begin{lstlisting}[language=Solidity, caption={Conceptual: Smart contract con FHE (estilo Fhenix)}]
// SPDX-License-Identifier: MIT
pragma solidity ^0.8.20;

// Tipos FHE (proporcionados por la blockchain)
// euint8, euint16, euint32, euint64 = encrypted unsigned integers
// ebool = encrypted boolean

import "@fhenix/contracts/FHE.sol";

/**
 * @title FHEToken
 * @notice Token donde TODOS los balances están encriptados
 * @dev Incluso el contrato no conoce los balances en claro
 */
contract FHEToken {
    // Balances encriptados
    mapping(address => euint64) private encryptedBalances;
    
    // Supply total (puede ser público o encriptado)
    euint64 private encryptedTotalSupply;
    
    // ========== OPERACIONES SOBRE DATOS ENCRIPTADOS ==========
    
    /**
     * @notice Transferencia con montos encriptados
     * @param to Destinatario
     * @param encryptedAmount Monto encriptado (solo el sender conoce el valor real)
     */
    function transfer(address to, inEuint64 calldata encryptedAmount) external {
        euint64 amount = FHE.asEuint64(encryptedAmount);
        
        // Verificar balance suficiente (comparación encriptada)
        // Retorna ebool - no sabemos el resultado, pero el circuito sí
        ebool hasEnough = FHE.gte(encryptedBalances[msg.sender], amount);
        
        // Require encriptado - falla si la condición es false
        // Pero no revela cuál es el balance ni el monto
        FHE.req(hasEnough);
        
        // Restar del sender (operación sobre encriptados)
        encryptedBalances[msg.sender] = FHE.sub(
            encryptedBalances[msg.sender],
            amount
        );
        
        // Sumar al receiver
        encryptedBalances[to] = FHE.add(
            encryptedBalances[to],
            amount
        );
    }
    
    /**
     * @notice Verificar si tienes suficiente balance (sin revelar cuánto)
     * @param encryptedThreshold Umbral encriptado a verificar
     * @return ebool Resultado encriptado (solo tú puedes descifrarlo)
     */
    function hasMinimumBalance(
        inEuint64 calldata encryptedThreshold
    ) external view returns (ebool) {
        euint64 threshold = FHE.asEuint64(encryptedThreshold);
        return FHE.gte(encryptedBalances[msg.sender], threshold);
    }
    
    /**
     * @notice Obtener tu balance encriptado (para descifrar localmente)
     */
    function getMyEncryptedBalance() external view returns (euint64) {
        // Requiere que el caller sea el dueño
        // El resultado solo puede ser descifrado por el caller
        return FHE.seal(encryptedBalances[msg.sender], msg.sender);
    }
    
    /**
     * @notice Mint tokens (admin only)
     */
    function mint(address to, inEuint64 calldata encryptedAmount) external {
        // Solo admin...
        euint64 amount = FHE.asEuint64(encryptedAmount);
        
        encryptedBalances[to] = FHE.add(encryptedBalances[to], amount);
        encryptedTotalSupply = FHE.add(encryptedTotalSupply, amount);
    }
}

/**
 * @title FHEAuction
 * @notice Subasta donde nadie (ni siquiera el contrato) conoce las pujas
 */
contract FHEAuction {
    address public owner;
    uint256 public endTime;
    bool public ended;
    
    // Pujas encriptadas
    mapping(address => euint64) private encryptedBids;
    address[] private bidders;
    
    // Puja más alta actual (encriptada)
    euint64 private highestBid;
    address private highestBidder; // Se conoce, pero no el monto
    
    constructor(uint256 duration) {
        owner = msg.sender;
        endTime = block.timestamp + duration;
    }
    
    /**
     * @notice Realizar puja encriptada
     */
    function bid(inEuint64 calldata encryptedBidAmount) external payable {
        require(block.timestamp < endTime, "Auction ended");
        
        euint64 bidAmount = FHE.asEuint64(encryptedBidAmount);
        
        // Si ya tenía puja, sumar
        if (encryptedBids[msg.sender].length > 0) {
            encryptedBids[msg.sender] = FHE.add(
                encryptedBids[msg.sender],
                bidAmount
            );
        } else {
            encryptedBids[msg.sender] = bidAmount;
            bidders.push(msg.sender);
        }
        
        // Actualizar líder si esta puja es mayor
        ebool isHigher = FHE.gt(encryptedBids[msg.sender], highestBid);
        
        // Actualización condicional encriptada
        highestBid = FHE.select(
            isHigher,
            encryptedBids[msg.sender],
            highestBid
        );
        
        // Nota: highestBidder se actualiza, pero esto revela
        // que hubo un nuevo líder (aunque no el monto)
    }
    
    /**
     * @notice Finalizar y revelar ganador
     * @dev Solo el monto del ganador se revela al final
     */
    function finalize() external {
        require(block.timestamp >= endTime, "Not ended");
        require(!ended, "Already finalized");
        
        ended = true;
        
        // Descifrar solo la puja ganadora
        // (requiere interacción con servicio de descifrado)
    }
    
    /**
     * @notice Ver tu puja (solo tú puedes descifrar)
     */
    function getMyBid() external view returns (euint64) {
        return FHE.seal(encryptedBids[msg.sender], msg.sender);
    }
}
\end{lstlisting}

\subsection{Comparación de Aproximaciones}

\begin{table}[htbp]
\centering
\footnotesize
\begin{tabular}{lcccc}
\toprule
\textbf{Aspecto} & \textbf{Secret/Oasis} & \textbf{Aleo} & \textbf{Mina} & \textbf{Fhenix} \\
\midrule
\textbf{Tecnología} & TEE (SGX) & ZK Proofs & Recursive ZK & FHE \\
\textbf{Trust Model} & Hardware & Criptografía & Criptografía & Criptografía \\
\textbf{Estado} & Encriptado on-chain & Off-chain/Records & Off-chain & Encriptado on-chain \\
\textbf{Latencia} & Baja & Media & Media & Alta \\
\textbf{Costo} & Bajo & Alto (proving) & Medio & Muy alto \\
\textbf{Madurez} & Alta & Media & Media & Experimental \\
\textbf{EVM Compat} & Sapphire: Sí & No (Leo) & No (o1js) & Modificado \\
\textbf{Riesgo} & Side-channels & Ninguno conocido & Ninguno conocido & Ninguno conocido \\
\bottomrule
\end{tabular}
\caption{Comparación de plataformas de smart contracts privados}
\end{table}
%===========================================
\section{Private Function Calls}
\label{sec:private-function-calls}
%===========================================

\subsection{Modelos de Ejecución Privada}

\begin{definition}[Private Function Call]
Ejecución de una función de smart contract donde los inputs, outputs, o ambos permanecen confidenciales:
\[
\text{PrivateCall}(f, x) \rightarrow \{y, \pi\}
\]
donde $y$ puede ser público o privado, y $\pi$ es una prueba de ejecución correcta.

\textbf{Niveles de privacidad:}
\begin{itemize}
    \item \textbf{Input privacy:} $x$ oculto, $y$ público
    \item \textbf{Output privacy:} $x$ público, $y$ oculto
    \item \textbf{Full privacy:} $x$ y $y$ ocultos
    \item \textbf{Function privacy:} Incluso qué función se ejecutó es privado
\end{itemize}
\end{definition}

\begin{figure}[htbp]
\centering
\begin{tikzpicture}[
    node distance=1.5cm,
    model/.style={rectangle, draw, rounded corners, minimum width=3cm, minimum height=2cm, align=center}
]
    % Cuatro modelos
    \node[model, fill=blue!20] (zk) at (0, 0) {
        \textbf{ZK}\\[3pt]
        Ejecución local\\
        Proof on-chain
    };
    
    \node[model, fill=green!20] (tee) at (4, 0) {
        \textbf{TEE}\\[3pt]
        Ejecución en enclave\\
        Attestation
    };
    
    \node[model, fill=orange!20] (mpc) at (8, 0) {
        \textbf{MPC}\\[3pt]
        Ejecución distribuida\\
        Threshold output
    };
    
    \node[model, fill=purple!20] (fhe) at (12, 0) {
        \textbf{FHE}\\[3pt]
        Ejecución encriptada\\
        Sin decrypt
    };
    
    % Características
    \node at (0, -1.8) [font=\small] {Trust: Crypto};
    \node at (4, -1.8) [font=\small] {Trust: Hardware};
    \node at (8, -1.8) [font=\small] {Trust: Threshold};
    \node at (12, -1.8) [font=\small] {Trust: Crypto};
    
    \node at (0, -2.3) [font=\small] {Latency: High};
    \node at (4, -2.3) [font=\small] {Latency: Low};
    \node at (8, -2.3) [font=\small] {Latency: Medium};
    \node at (12, -2.3) [font=\small] {Latency: Very High};
\end{tikzpicture}
\caption{Cuatro modelos de ejecución privada de funciones}
\label{fig:private-execution-models}
\end{figure}

\begin{lstlisting}[language=Python, caption={Framework para private function calls}]
from abc import ABC, abstractmethod
from dataclasses import dataclass
from typing import Any, Dict, List, Optional, Tuple
from enum import Enum
import hashlib
import secrets

class PrivacyLevel(Enum):
    INPUT_PRIVATE = "input_private"
    OUTPUT_PRIVATE = "output_private"
    FULL_PRIVATE = "full_private"
    FUNCTION_PRIVATE = "function_private"

class ExecutionModel(Enum):
    ZK = "zero_knowledge"
    TEE = "trusted_execution"
    MPC = "multi_party_computation"
    FHE = "fully_homomorphic"

@dataclass
class ExecutionResult:
    output: Any
    proof: Optional[bytes]
    public_inputs: List[Any]
    execution_model: ExecutionModel
    privacy_level: PrivacyLevel

class PrivateFunctionExecutor(ABC):
    """
    Clase base abstracta para ejecutores de funciones privadas.
    """
    
    @abstractmethod
    def execute(
        self,
        function_name: str,
        private_inputs: Dict[str, Any],
        public_inputs: Dict[str, Any]
    ) -> ExecutionResult:
        pass
    
    @abstractmethod
    def verify(self, result: ExecutionResult) -> bool:
        pass


class ZKExecutor(PrivateFunctionExecutor):
    """
    Ejecutor basado en Zero-Knowledge Proofs.
    - Ejecución local
    - Genera proof de ejecución correcta
    - Verifier on-chain no ve inputs privados
    """
    
    def __init__(self, circuit_path: str):
        self.circuit_path = circuit_path
        self.proving_key = None
        self.verification_key = None
    
    def setup(self):
        """Generar proving y verification keys."""
        # En producción: usar snarkjs, arkworks, etc.
        self.proving_key = secrets.token_bytes(64)
        self.verification_key = secrets.token_bytes(32)
    
    def execute(
        self,
        function_name: str,
        private_inputs: Dict[str, Any],
        public_inputs: Dict[str, Any]
    ) -> ExecutionResult:
        # 1. Ejecutar función localmente
        output = self._execute_locally(function_name, private_inputs, public_inputs)
        
        # 2. Generar witness (private + public inputs + output)
        witness = self._generate_witness(private_inputs, public_inputs, output)
        
        # 3. Generar ZK proof
        proof = self._generate_proof(witness)
        
        return ExecutionResult(
            output=output,
            proof=proof,
            public_inputs=list(public_inputs.values()) + [output],
            execution_model=ExecutionModel.ZK,
            privacy_level=PrivacyLevel.INPUT_PRIVATE
        )
    
    def verify(self, result: ExecutionResult) -> bool:
        """Verificar ZK proof."""
        # En producción: verificación SNARK real
        return result.proof is not None and len(result.proof) == 128
    
    def _execute_locally(
        self,
        function_name: str,
        private_inputs: Dict,
        public_inputs: Dict
    ) -> Any:
        """Ejecutar la función localmente."""
        # Simulación de ejecución
        all_inputs = {**private_inputs, **public_inputs}
        
        if function_name == "transfer":
            balance = all_inputs.get('balance', 0)
            amount = all_inputs.get('amount', 0)
            return balance - amount  # Nuevo balance
        
        return None
    
    def _generate_witness(
        self,
        private_inputs: Dict,
        public_inputs: Dict,
        output: Any
    ) -> bytes:
        """Generar witness para el circuito."""
        data = str(private_inputs) + str(public_inputs) + str(output)
        return hashlib.sha256(data.encode()).digest()
    
    def _generate_proof(self, witness: bytes) -> bytes:
        """Generar ZK proof (simulado)."""
        return hashlib.sha256(witness + self.proving_key).digest() * 4


class TEEExecutor(PrivateFunctionExecutor):
    """
    Ejecutor basado en Trusted Execution Environment.
    - Ejecución en enclave SGX/SEV
    - Attestation del enclave
    - Baja latencia, trust en hardware
    """
    
    def __init__(self, enclave_id: str):
        self.enclave_id = enclave_id
        self.enclave_key = secrets.token_bytes(32)
    
    def execute(
        self,
        function_name: str,
        private_inputs: Dict[str, Any],
        public_inputs: Dict[str, Any]
    ) -> ExecutionResult:
        # 1. Enviar datos encriptados al enclave
        encrypted_inputs = self._encrypt_for_enclave(private_inputs)
        
        # 2. Ejecutar en enclave (simulado)
        output, attestation = self._execute_in_enclave(
            function_name,
            encrypted_inputs,
            public_inputs
        )
        
        return ExecutionResult(
            output=output,
            proof=attestation,
            public_inputs=list(public_inputs.values()),
            execution_model=ExecutionModel.TEE,
            privacy_level=PrivacyLevel.FULL_PRIVATE
        )
    
    def verify(self, result: ExecutionResult) -> bool:
        """Verificar attestation del enclave."""
        # En producción: verificar con Intel Attestation Service
        return result.proof is not None and len(result.proof) == 96
    
    def _encrypt_for_enclave(self, data: Dict) -> bytes:
        """Encriptar datos para el enclave."""
        plaintext = str(data).encode()
        # En producción: ECDH + AES-GCM
        return bytes(b ^ k for b, k in zip(plaintext, self.enclave_key * 10))
    
    def _execute_in_enclave(
        self,
        function_name: str,
        encrypted_inputs: bytes,
        public_inputs: Dict
    ) -> Tuple[Any, bytes]:
        """Simular ejecución en enclave."""
        # En producción: llamada real a enclave SGX
        
        # Ejecutar función
        output = "execution_result"
        
        # Generar attestation
        attestation = hashlib.sha256(
            self.enclave_id.encode() + 
            encrypted_inputs + 
            str(output).encode()
        ).digest() * 3
        
        return output, attestation


class MPCExecutor(PrivateFunctionExecutor):
    """
    Ejecutor basado en Multi-Party Computation.
    - Ejecución distribuida entre múltiples partes
    - Ninguna parte ve los inputs completos
    - Threshold de partes honestas
    """
    
    def __init__(self, parties: List[str], threshold: int):
        self.parties = parties
        self.threshold = threshold
        self.num_parties = len(parties)
    
    def execute(
        self,
        function_name: str,
        private_inputs: Dict[str, Any],
        public_inputs: Dict[str, Any]
    ) -> ExecutionResult:
        # 1. Secret-share los inputs entre partes
        shares = self._secret_share_inputs(private_inputs)
        
        # 2. Cada parte procesa su share
        partial_results = []
        for i, party in enumerate(self.parties):
            partial = self._party_compute(party, shares[i], function_name)
            partial_results.append(partial)
        
        # 3. Combinar resultados (threshold)
        output = self._combine_results(partial_results)
        
        # 4. Generar proof de correctitud
        proof = self._generate_mpc_proof(partial_results)
        
        return ExecutionResult(
            output=output,
            proof=proof,
            public_inputs=list(public_inputs.values()),
            execution_model=ExecutionModel.MPC,
            privacy_level=PrivacyLevel.FULL_PRIVATE
        )
    
    def verify(self, result: ExecutionResult) -> bool:
        """Verificar que al menos threshold partes participaron."""
        return result.proof is not None
    
    def _secret_share_inputs(self, inputs: Dict) -> List[Dict]:
        """Dividir inputs en shares."""
        shares = []
        for i in range(self.num_parties):
            party_share = {}
            for key, value in inputs.items():
                if isinstance(value, int):
                    # Shamir secret sharing (simplificado)
                    party_share[key] = value + i * secrets.randbelow(1000)
            shares.append(party_share)
        return shares
    
    def _party_compute(
        self,
        party: str,
        share: Dict,
        function_name: str
    ) -> Dict:
        """Computación de una parte."""
        return {'party': party, 'partial_result': sum(share.values())}
    
    def _combine_results(self, partial_results: List[Dict]) -> Any:
        """Combinar resultados de las partes."""
        # Lagrange interpolation (simplificado)
        total = sum(r['partial_result'] for r in partial_results)
        return total // len(partial_results)
    
    def _generate_mpc_proof(self, partial_results: List[Dict]) -> bytes:
        """Generar proof de MPC correcto."""
        data = str(partial_results)
        return hashlib.sha256(data.encode()).digest()


def private_function_demo():
    print("=== Private Function Calls Demo ===\n")
    
    # 1. ZK Executor
    print("--- ZK Executor ---")
    zk_exec = ZKExecutor("transfer_circuit.r1cs")
    zk_exec.setup()
    
    result_zk = zk_exec.execute(
        function_name="transfer",
        private_inputs={'balance': 1000, 'secret_key': 'abc123'},
        public_inputs={'amount': 100, 'recipient': '0xBob'}
    )
    
    print(f"Output: {result_zk.output}")
    print(f"Proof size: {len(result_zk.proof)} bytes")
    print(f"Verified: {zk_exec.verify(result_zk)}")
    print(f"Privacy: {result_zk.privacy_level.value}")
    
    # 2. TEE Executor
    print("\n--- TEE Executor ---")
    tee_exec = TEEExecutor("sgx_enclave_001")
    
    result_tee = tee_exec.execute(
        function_name="compute_private",
        private_inputs={'salary': 100000, 'ssn': '123-45-6789'},
        public_inputs={'tax_rate': 0.25}
    )
    
    print(f"Output: {result_tee.output}")
    print(f"Attestation size: {len(result_tee.proof)} bytes")
    print(f"Verified: {tee_exec.verify(result_tee)}")
    
    # 3. MPC Executor
    print("\n--- MPC Executor ---")
    mpc_exec = MPCExecutor(
        parties=["party_a", "party_b", "party_c"],
        threshold=2
    )
    
    result_mpc = mpc_exec.execute(
        function_name="aggregate",
        private_inputs={'value_a': 100, 'value_b': 200, 'value_c': 300},
        public_inputs={}
    )
    
    print(f"Aggregated output: {result_mpc.output}")
    print(f"MPC proof: {result_mpc.proof.hex()[:32]}...")
    print(f"Verified: {mpc_exec.verify(result_mpc)}")

private_function_demo()
\end{lstlisting}

%===========================================
\section{Encrypted Mempools}
\label{sec:encrypted-mempools}
%===========================================

\subsection{El Problema del Mempool Transparente}

\begin{definition}[Encrypted Mempool]
Mecanismo donde las transacciones permanecen encriptadas mientras están en el mempool, siendo descifradas solo después de su inclusión en un bloque:

\textbf{Flujo:}
\begin{enumerate}
    \item Usuario encripta tx con clave threshold
    \item Tx encriptada entra al mempool
    \item Sequencer/Builder ordena sin ver contenido
    \item Block producer incluye tx encriptada
    \item Post-inclusión: threshold decryption
    \item Ejecución con tx descifrada
\end{enumerate}
\end{definition}

\begin{figure}[htbp]
\centering
\begin{tikzpicture}[
    node distance=1cm,
    step/.style={rectangle, draw, rounded corners, minimum width=2.2cm, minimum height=1cm, align=center, font=\small},
    arrow/.style={->, thick}
]
    % Mempool tradicional (arriba)
    \node at (0, 4) {\textbf{Mempool Tradicional}};
    
    \node[step, fill=blue!20] (t1) at (0, 2.5) {TX visible};
    \node[step, fill=red!20] (t2) at (3, 2.5) {MEV bots\\observan};
    \node[step, fill=red!20] (t3) at (6, 2.5) {Frontrun/\\Sandwich};
    \node[step, fill=orange!20] (t4) at (9, 2.5) {Usuario\\pierde \$};
    
    \draw[arrow] (t1) -- (t2);
    \draw[arrow] (t2) -- (t3);
    \draw[arrow] (t3) -- (t4);
    
    % Encrypted mempool (abajo)
    \node at (0, 0.5) {\textbf{Encrypted Mempool}};
    
    \node[step, fill=green!20] (e1) at (0, -1) {TX encriptada};
    \node[step, fill=green!20] (e2) at (3, -1) {Ordenar\\sin ver};
    \node[step, fill=green!20] (e3) at (6, -1) {Incluir\\en bloque};
    \node[step, fill=blue!20] (e4) at (9, -1) {Decrypt +\\Ejecutar};
    
    \draw[arrow] (e1) -- (e2);
    \draw[arrow] (e2) -- (e3);
    \draw[arrow] (e3) -- (e4);
    
    \node at (4.5, -2.2) [font=\small, green!60!black] {MEV bots no pueden ver contenido};
\end{tikzpicture}
\caption{Encrypted mempool previene frontrunning}
\label{fig:encrypted-mempool}
\end{figure}

\begin{mathnote}[Esquemas de Encriptación para Mempools]
\textbf{1. Threshold Encryption:}
\begin{itemize}
    \item Comité de $n$ partes, threshold $t$
    \item $t$ partes deben cooperar para descifrar
    \item Usado por: Shutter Network
\end{itemize}

\textbf{2. Timelock Encryption:}
\begin{itemize}
    \item Descifrado automático después de tiempo T
    \item Basado en VDF o witness encryption
    \item Usado por: Penumbra (parcial)
\end{itemize}

\textbf{3. Commit-Reveal:}
\begin{itemize}
    \item Fase 1: Commit hash de tx
    \item Fase 2: Reveal tx completa
    \item Más simple pero requiere 2 txs
\end{itemize}

\textbf{4. Identity-Based Encryption (IBE):}
\begin{itemize}
    \item Encriptar a ``bloque N''
    \item Clave para bloque N disponible después de finalizarlo
    \item Usado por: Flashbots SUAVE (propuesto)
\end{itemize}
\end{mathnote}

\begin{lstlisting}[language=Python, caption={Sistema de Encrypted Mempool}]
from dataclasses import dataclass
from datetime import datetime
from typing import Dict, List, Optional, Tuple
from enum import Enum
import hashlib
import secrets
from collections import deque

class EncryptionScheme(Enum):
    THRESHOLD = "threshold"
    TIMELOCK = "timelock"
    COMMIT_REVEAL = "commit_reveal"

@dataclass
class EncryptedTransaction:
    tx_id: str
    encrypted_data: bytes
    commitment: bytes  # Hash para verificación
    sender_pubkey: bytes
    timestamp: datetime
    target_block: Optional[int]  # Para timelock
    
@dataclass
class DecryptedTransaction:
    tx_id: str
    sender: str
    recipient: str
    value: int
    data: bytes
    gas_price: int

class EncryptedMempool:
    """
    Mempool con transacciones encriptadas.
    """
    
    def __init__(
        self,
        scheme: EncryptionScheme,
        threshold_parties: List[str] = None,
        threshold_k: int = None
    ):
        self.scheme = scheme
        self.pending_txs: Dict[str, EncryptedTransaction] = {}
        self.decryption_queue: deque = deque()
        
        # Para threshold encryption
        self.parties = threshold_parties or []
        self.threshold_k = threshold_k or 2
        self.partial_decryptions: Dict[str, List[bytes]] = {}
        
        # Para commit-reveal
        self.commits: Dict[bytes, str] = {}  # commitment -> tx_id
        self.reveals: Dict[str, DecryptedTransaction] = {}
        
        # Clave pública del comité (threshold)
        self.committee_pubkey = secrets.token_bytes(32)
    
    # ========== SUBMIT ENCRYPTED TX ==========
    
    def submit_encrypted_tx(
        self,
        encrypted_data: bytes,
        sender_pubkey: bytes,
        target_block: Optional[int] = None
    ) -> str:
        """
        Añadir transacción encriptada al mempool.
        """
        tx_id = secrets.token_hex(16)
        commitment = hashlib.sha256(encrypted_data).digest()
        
        enc_tx = EncryptedTransaction(
            tx_id=tx_id,
            encrypted_data=encrypted_data,
            commitment=commitment,
            sender_pubkey=sender_pubkey,
            timestamp=datetime.now(),
            target_block=target_block
        )
        
        self.pending_txs[tx_id] = enc_tx
        
        return tx_id
    
    def submit_commit(self, commitment: bytes) -> str:
        """
        Fase 1 de commit-reveal: submit commitment.
        """
        tx_id = secrets.token_hex(16)
        self.commits[commitment] = tx_id
        return tx_id
    
    def submit_reveal(
        self,
        commitment: bytes,
        tx_data: DecryptedTransaction
    ) -> bool:
        """
        Fase 2 de commit-reveal: reveal transaction.
        """
        # Verificar commitment
        computed_commitment = self._compute_commitment(tx_data)
        if computed_commitment != commitment:
            return False
        
        tx_id = self.commits.get(commitment)
        if not tx_id:
            return False
        
        self.reveals[tx_id] = tx_data
        return True
    
    # ========== ORDENAMIENTO (sin ver contenido) ==========
    
    def order_transactions(self) -> List[str]:
        """
        Ordenar transacciones sin ver contenido.
        Solo podemos usar: timestamp, gas (si está en header), sender.
        """
        ordered = sorted(
            self.pending_txs.values(),
            key=lambda tx: tx.timestamp
        )
        
        return [tx.tx_id for tx in ordered]
    
    # ========== DECRYPTION ==========
    
    def submit_partial_decryption(
        self,
        tx_id: str,
        party: str,
        partial_decryption: bytes
    ):
        """
        Una parte del comité submite su partial decryption.
        """
        if tx_id not in self.partial_decryptions:
            self.partial_decryptions[tx_id] = []
        
        self.partial_decryptions[tx_id].append(partial_decryption)
        
        # Verificar si tenemos threshold
        if len(self.partial_decryptions[tx_id]) >= self.threshold_k:
            self.decryption_queue.append(tx_id)
    
    def decrypt_transaction(self, tx_id: str) -> Optional[DecryptedTransaction]:
        """
        Descifrar transacción si tenemos suficientes partial decryptions.
        """
        if tx_id not in self.pending_txs:
            return None
        
        enc_tx = self.pending_txs[tx_id]
        
        if self.scheme == EncryptionScheme.THRESHOLD:
            return self._threshold_decrypt(enc_tx)
        elif self.scheme == EncryptionScheme.COMMIT_REVEAL:
            return self.reveals.get(tx_id)
        elif self.scheme == EncryptionScheme.TIMELOCK:
            return self._timelock_decrypt(enc_tx)
        
        return None
    
    def _threshold_decrypt(
        self,
        enc_tx: EncryptedTransaction
    ) -> Optional[DecryptedTransaction]:
        """
        Descifrar usando threshold decryption.
        """
        partials = self.partial_decryptions.get(enc_tx.tx_id, [])
        
        if len(partials) < self.threshold_k:
            return None
        
        # Combinar partial decryptions (simplificado)
        combined_key = self._combine_partials(partials)
        
        # Descifrar
        plaintext = self._decrypt_with_key(enc_tx.encrypted_data, combined_key)
        
        # Parsear transacción
        return self._parse_transaction(plaintext, enc_tx.tx_id)
    
    def _timelock_decrypt(
        self,
        enc_tx: EncryptedTransaction
    ) -> Optional[DecryptedTransaction]:
        """
        Descifrar después de que el bloque target sea finalizado.
        """
        # Verificar que el bloque target fue alcanzado
        # (En producción: verificar con el estado de la blockchain)
        
        # El descifrado timelock es automático después del tiempo
        key = self._derive_timelock_key(enc_tx.target_block)
        plaintext = self._decrypt_with_key(enc_tx.encrypted_data, key)
        
        return self._parse_transaction(plaintext, enc_tx.tx_id)
    
    # ========== HELPERS ==========
    
    def _compute_commitment(self, tx: DecryptedTransaction) -> bytes:
        """Calcular commitment de transacción."""
        data = f"{tx.sender}:{tx.recipient}:{tx.value}:{tx.data.hex()}"
        return hashlib.sha256(data.encode()).digest()
    
    def _combine_partials(self, partials: List[bytes]) -> bytes:
        """Combinar partial decryptions (simplificado)."""
        # En producción: Lagrange interpolation
        combined = bytes(32)
        for p in partials:
            combined = bytes(a ^ b for a, b in zip(combined, p))
        return combined
    
    def _decrypt_with_key(self, ciphertext: bytes, key: bytes) -> bytes:
        """Descifrar con clave simétrica."""
        # En producción: AES-GCM
        return bytes(c ^ k for c, k in zip(ciphertext, key * 100))
    
    def _derive_timelock_key(self, block_number: int) -> bytes:
        """Derivar clave timelock del número de bloque."""
        return hashlib.sha256(f"timelock_{block_number}".encode()).digest()
    
    def _parse_transaction(
        self,
        plaintext: bytes,
        tx_id: str
    ) -> DecryptedTransaction:
        """Parsear transacción desde bytes."""
        # Simplificado
        return DecryptedTransaction(
            tx_id=tx_id,
            sender="0xSender",
            recipient="0xRecipient",
            value=1000,
            data=plaintext[:32],
            gas_price=100
        )
    
    def get_mempool_stats(self) -> Dict:
        """Estadísticas del mempool."""
        return {
            'pending_encrypted': len(self.pending_txs),
            'ready_to_decrypt': len(self.decryption_queue),
            'commits_waiting_reveal': len(self.commits) - len(self.reveals),
            'scheme': self.scheme.value
        }


class ShutterIntegration:
    """
    Integración con Shutter Network para encrypted mempool.
    """
    
    def __init__(self, keypers: List[str], threshold: int):
        self.keypers = keypers  # Nodos que mantienen las claves
        self.threshold = threshold
        self.eon_keys: Dict[int, bytes] = {}  # eon -> public key
    
    def get_encryption_key(self, eon: int) -> bytes:
        """
        Obtener clave pública para un eon (período).
        Las transacciones se encriptan con esta clave.
        """
        if eon not in self.eon_keys:
            # En producción: query a los keypers
            self.eon_keys[eon] = secrets.token_bytes(32)
        
        return self.eon_keys[eon]
    
    def encrypt_for_eon(self, tx_data: bytes, eon: int) -> bytes:
        """
        Encriptar transacción para un eon específico.
        """
        eon_key = self.get_encryption_key(eon)
        
        # En producción: threshold encryption con eon_key
        nonce = secrets.token_bytes(12)
        ciphertext = bytes(t ^ k for t, k in zip(tx_data, eon_key * 100))
        
        return nonce + ciphertext


def encrypted_mempool_demo():
    print("=== Encrypted Mempool Demo ===\n")
    
    # Crear mempool con threshold encryption
    mempool = EncryptedMempool(
        scheme=EncryptionScheme.THRESHOLD,
        threshold_parties=["keyper_1", "keyper_2", "keyper_3"],
        threshold_k=2
    )
    
    # Usuario crea transacción encriptada
    print("--- User Creates Encrypted TX ---")
    
    tx_data = b"send 100 ETH to 0xBob"
    encrypted = bytes(d ^ 0x42 for d in tx_data)  # Encriptación simplificada
    
    tx_id = mempool.submit_encrypted_tx(
        encrypted_data=encrypted,
        sender_pubkey=b"alice_pubkey",
        target_block=None
    )
    
    print(f"Submitted encrypted TX: {tx_id}")
    print(f"MEV bots see: [encrypted blob]")
    print(f"Cannot extract: sender intent, recipient, amount")
    
    # Ordenamiento (sin ver contenido)
    print("\n--- Block Builder Orders TXs ---")
    order = mempool.order_transactions()
    print(f"Ordered by timestamp only: {order}")
    
    # Block se produce
    print("\n--- Block Produced ---")
    
    # Keypers submiten partial decryptions
    for i in range(2):
        mempool.submit_partial_decryption(
            tx_id=tx_id,
            party=f"keyper_{i+1}",
            partial_decryption=secrets.token_bytes(32)
        )
        print(f"Keyper {i+1} submitted partial decryption")
    
    # Descifrar
    print("\n--- Decryption (post-inclusion) ---")
    decrypted = mempool.decrypt_transaction(tx_id)
    
    if decrypted:
        print(f"TX decrypted: {decrypted.tx_id}")
        print(f"Now executing: send {decrypted.value} to {decrypted.recipient}")
    
    # Stats
    print("\n--- Mempool Stats ---")
    stats = mempool.get_mempool_stats()
    for k, v in stats.items():
        print(f"  {k}: {v}")

encrypted_mempool_demo()
\end{lstlisting}

%===========================================
\section{MEV Protection con Privacidad}
\label{sec:mev-protection}
%===========================================

\subsection{El Problema del MEV}

\begin{definition}[MEV (Maximal Extractable Value)]
Valor que puede ser extraído por block producers (miners/validators) o searchers mediante la manipulación del orden de transacciones:

\textbf{Tipos de MEV:}
\begin{itemize}
    \item \textbf{Frontrunning:} Ejecutar antes de una tx observada
    \item \textbf{Backrunning:} Ejecutar inmediatamente después
    \item \textbf{Sandwich attacks:} Frontrun + backrun combinados
    \item \textbf{Liquidations:} Competir por liquidar posiciones
    \item \textbf{Arbitrage:} Explotar diferencias de precio
\end{itemize}
\end{definition}

\begin{figure}[htbp]
\centering
\begin{tikzpicture}[
    node distance=0.5cm,
    tx/.style={rectangle, draw, minimum width=2.5cm, minimum height=0.7cm, align=center, font=\small}
]
    % Sandwich attack
    \node at (0, 3.5) {\textbf{Sandwich Attack}};
    
    \node[tx, fill=red!30] (front) at (0, 2.5) {Attacker: BUY\\(eleva precio)};
    \node[tx, fill=blue!20] (victim) at (0, 1.5) {Víctima: BUY\\(paga más)};
    \node[tx, fill=red!30] (back) at (0, 0.5) {Attacker: SELL\\(captura profit)};
    
    \draw[->, thick] (front) -- (victim);
    \draw[->, thick] (victim) -- (back);
    
    % Resultado
    \node at (0, -0.5) [font=\small] {Víctima: slippage excesivo};
    \node at (0, -1) [font=\small, red] {Attacker: profit garantizado};
    
    % Protección
    \node at (6, 2) [draw, rectangle, rounded corners, fill=green!20, minimum width=4cm, minimum height=2.5cm, align=center] {
        \textbf{Encrypted Mempool}\\[5pt]
        Attacker no puede ver\\
        la tx de la víctima\\
        antes de la inclusión
    };
    
    \draw[->, thick, green!60!black] (3, 1.5) -- (4, 1.5);
    \node at (3.5, 2) [font=\small, green!60!black] {protege};
\end{tikzpicture}
\caption{Sandwich attack y protección mediante encrypted mempool}
\label{fig:sandwich-attack}
\end{figure}

\subsection{Técnicas de Protección}

\begin{table}[htbp]
\centering
\small
\begin{tabular}{lp{4cm}cc}
\toprule
\textbf{Técnica} & \textbf{Descripción} & \textbf{Efectividad} & \textbf{Complejidad} \\
\midrule
Encrypted Mempools & TX encriptada hasta inclusión & Muy alta & Alta \\
Commit-Reveal & Hash primero, datos después & Alta & Media \\
Private Order Flow & Enviar directo a builder trusted & Media & Baja \\
Time-lock Puzzles & Descifrar requiere tiempo & Alta & Alta \\
Batch Auctions & Todas las órdenes al mismo precio & Alta & Media \\
Threshold Encryption & k-of-n para descifrar & Muy alta & Alta \\
\bottomrule
\end{tabular}
\caption{Técnicas de protección contra MEV}
\end{table}

\begin{lstlisting}[language=Python, caption={Sistema de protección MEV}]
from dataclasses import dataclass
from datetime import datetime, timedelta
from typing import Dict, List, Optional
from enum import Enum
import hashlib
import secrets

class MEVProtectionLevel(Enum):
    NONE = 0
    BASIC = 1      # Private order flow
    MEDIUM = 2     # Commit-reveal
    HIGH = 3       # Encrypted mempool
    MAXIMUM = 4    # Threshold + timelock

@dataclass
class ProtectedTransaction:
    tx_id: str
    protection_level: MEVProtectionLevel
    
    # Para commit-reveal
    commitment: Optional[bytes]
    reveal_deadline: Optional[datetime]
    revealed_data: Optional[bytes]
    
    # Para encryption
    encrypted_data: Optional[bytes]
    decryption_key: Optional[bytes]
    
    # Metadata
    submitted_at: datetime
    sender: str

class MEVProtectionService:
    """
    Servicio de protección contra MEV con múltiples niveles.
    """
    
    def __init__(self):
        self.pending_commits: Dict[bytes, ProtectedTransaction] = {}
        self.protected_txs: Dict[str, ProtectedTransaction] = {}
        self.trusted_builders: List[str] = []
        
        # Threshold encryption committee
        self.threshold_committee: List[str] = []
        self.threshold_k = 2
    
    def configure_trusted_builders(self, builders: List[str]):
        """Configurar builders de confianza para private order flow."""
        self.trusted_builders = builders
    
    def configure_threshold_committee(
        self,
        members: List[str],
        threshold: int
    ):
        """Configurar comité para threshold encryption."""
        self.threshold_committee = members
        self.threshold_k = threshold
    
    # ========== NIVELES DE PROTECCIÓN ==========
    
    def submit_private_order_flow(
        self,
        tx_data: bytes,
        sender: str,
        preferred_builder: Optional[str] = None
    ) -> ProtectedTransaction:
        """
        Nivel BASIC: Enviar directamente a builder de confianza.
        No pasa por mempool público.
        """
        builder = preferred_builder or self.trusted_builders[0]
        
        tx = ProtectedTransaction(
            tx_id=secrets.token_hex(16),
            protection_level=MEVProtectionLevel.BASIC,
            commitment=None,
            reveal_deadline=None,
            revealed_data=tx_data,
            encrypted_data=None,
            decryption_key=None,
            submitted_at=datetime.now(),
            sender=sender
        )
        
        # Enviar directamente al builder (simulated)
        self._send_to_builder(builder, tx)
        
        return tx
    
    def submit_commit_reveal(
        self,
        tx_data: bytes,
        sender: str,
        reveal_delay_seconds: int = 12
    ) -> ProtectedTransaction:
        """
        Nivel MEDIUM: Commit-reveal en dos fases.
        """
        # Fase 1: Commit
        salt = secrets.token_bytes(32)
        commitment = hashlib.sha256(tx_data + salt).digest()
        
        tx = ProtectedTransaction(
            tx_id=secrets.token_hex(16),
            protection_level=MEVProtectionLevel.MEDIUM,
            commitment=commitment,
            reveal_deadline=datetime.now() + timedelta(seconds=reveal_delay_seconds),
            revealed_data=None,
            encrypted_data=tx_data + salt,  # Para reveal posterior
            decryption_key=None,
            submitted_at=datetime.now(),
            sender=sender
        )
        
        self.pending_commits[commitment] = tx
        self.protected_txs[tx.tx_id] = tx
        
        return tx
    
    def reveal_transaction(
        self,
        commitment: bytes,
        tx_data: bytes,
        salt: bytes
    ) -> bool:
        """
        Fase 2 de commit-reveal: revelar la transacción.
        """
        if commitment not in self.pending_commits:
            return False
        
        tx = self.pending_commits[commitment]
        
        # Verificar que no pasó el deadline
        if datetime.now() > tx.reveal_deadline:
            return False
        
        # Verificar commitment
        expected = hashlib.sha256(tx_data + salt).digest()
        if expected != commitment:
            return False
        
        tx.revealed_data = tx_data
        return True
    
    def submit_encrypted(
        self,
        tx_data: bytes,
        sender: str,
        target_block: int
    ) -> ProtectedTransaction:
        """
        Nivel HIGH: Transacción encriptada con threshold encryption.
        """
        # Encriptar para el comité
        encryption_key = secrets.token_bytes(32)
        encrypted = self._encrypt(tx_data, encryption_key)
        
        # Distribuir key shares al comité
        key_shares = self._split_key(encryption_key)
        
        tx = ProtectedTransaction(
            tx_id=secrets.token_hex(16),
            protection_level=MEVProtectionLevel.HIGH,
            commitment=hashlib.sha256(encrypted).digest(),
            reveal_deadline=None,
            revealed_data=None,
            encrypted_data=encrypted,
            decryption_key=None,  # Key distribuida
            submitted_at=datetime.now(),
            sender=sender
        )
        
        self.protected_txs[tx.tx_id] = tx
        
        return tx
    
    def submit_maximum_protection(
        self,
        tx_data: bytes,
        sender: str,
        timelock_blocks: int = 2
    ) -> ProtectedTransaction:
        """
        Nivel MAXIMUM: Threshold + timelock encryption.
        Solo descifrable después de N bloques Y con threshold del comité.
        """
        # Encriptar con timelock
        timelock_key = self._generate_timelock_key(timelock_blocks)
        layer1_encrypted = self._encrypt(tx_data, timelock_key)
        
        # Encriptar con threshold
        threshold_key = secrets.token_bytes(32)
        layer2_encrypted = self._encrypt(layer1_encrypted, threshold_key)
        
        tx = ProtectedTransaction(
            tx_id=secrets.token_hex(16),
            protection_level=MEVProtectionLevel.MAXIMUM,
            commitment=hashlib.sha256(layer2_encrypted).digest(),
            reveal_deadline=None,
            revealed_data=None,
            encrypted_data=layer2_encrypted,
            decryption_key=None,
            submitted_at=datetime.now(),
            sender=sender
        )
        
        self.protected_txs[tx.tx_id] = tx
        
        return tx
    
    # ========== ANÁLISIS DE PROTECCIÓN ==========
    
    def analyze_protection_effectiveness(
        self,
        tx: ProtectedTransaction
    ) -> Dict:
        """
        Analizar qué tan protegida está una transacción.
        """
        analysis = {
            'protection_level': tx.protection_level.name,
            'protected_from': [],
            'vulnerable_to': [],
            'recommendations': []
        }
        
        if tx.protection_level == MEVProtectionLevel.NONE:
            analysis['vulnerable_to'] = [
                'frontrunning', 'sandwich', 'backrunning'
            ]
            analysis['recommendations'].append(
                "Consider using at least private order flow"
            )
            
        elif tx.protection_level == MEVProtectionLevel.BASIC:
            analysis['protected_from'] = ['public mempool exposure']
            analysis['vulnerable_to'] = ['malicious builder']
            analysis['recommendations'].append(
                "Use commit-reveal for higher security"
            )
            
        elif tx.protection_level == MEVProtectionLevel.MEDIUM:
            analysis['protected_from'] = [
                'frontrunning (before reveal)',
                'instant sandwich'
            ]
            analysis['vulnerable_to'] = [
                'long-term correlation',
                'timing attacks on reveal'
            ]
            
        elif tx.protection_level == MEVProtectionLevel.HIGH:
            analysis['protected_from'] = [
                'all pre-inclusion MEV',
                'builder manipulation'
            ]
            analysis['vulnerable_to'] = [
                'committee collusion'
            ]
            
        elif tx.protection_level == MEVProtectionLevel.MAXIMUM:
            analysis['protected_from'] = [
                'all known MEV vectors',
                'timing attacks',
                'committee collusion (timelock)'
            ]
            analysis['vulnerable_to'] = [
                'quantum attacks (theoretical)'
            ]
        
        return analysis
    
    # ========== HELPERS ==========
    
    def _encrypt(self, data: bytes, key: bytes) -> bytes:
        """Encriptar datos."""
        return bytes(d ^ k for d, k in zip(data, key * 100))
    
    def _split_key(self, key: bytes) -> List[bytes]:
        """Dividir key en shares."""
        shares = []
        for i in range(len(self.threshold_committee)):
            share = bytes(k ^ i for k in key)
            shares.append(share)
        return shares
    
    def _generate_timelock_key(self, blocks: int) -> bytes:
        """Generar timelock key."""
        return hashlib.sha256(f"timelock_{blocks}".encode()).digest()
    
    def _send_to_builder(self, builder: str, tx: ProtectedTransaction):
        """Enviar tx a builder (simulado)."""
        pass


def mev_protection_demo():
    print("=== MEV Protection Demo ===\n")
    
    service = MEVProtectionService()
    
    # Configurar
    service.configure_trusted_builders([
        "flashbots_builder",
        "blocknative_builder"
    ])
    
    service.configure_threshold_committee(
        members=["keyper_1", "keyper_2", "keyper_3"],
        threshold=2
    )
    
    # Demostrar diferentes niveles
    tx_data = b"swap 1000 USDC for ETH"
    
    print("--- Level: Private Order Flow ---")
    tx1 = service.submit_private_order_flow(tx_data, "user1")
    analysis1 = service.analyze_protection_effectiveness(tx1)
    print(f"Protected from: {analysis1['protected_from']}")
    print(f"Vulnerable to: {analysis1['vulnerable_to']}")
    
    print("\n--- Level: Commit-Reveal ---")
    tx2 = service.submit_commit_reveal(tx_data, "user2", reveal_delay_seconds=12)
    analysis2 = service.analyze_protection_effectiveness(tx2)
    print(f"Protection: {analysis2['protection_level']}")
    print(f"Protected from: {analysis2['protected_from']}")
    
    print("\n--- Level: Encrypted Mempool ---")
    tx3 = service.submit_encrypted(tx_data, "user3", target_block=1000)
    analysis3 = service.analyze_protection_effectiveness(tx3)
    print(f"Protected from: {analysis3['protected_from']}")
    
    print("\n--- Level: Maximum (Threshold + Timelock) ---")
    tx4 = service.submit_maximum_protection(tx_data, "user4", timelock_blocks=2)
    analysis4 = service.analyze_protection_effectiveness(tx4)
    print(f"Protection: {analysis4['protection_level']}")
    print(f"Protected from: {analysis4['protected_from']}")
    print(f"Vulnerable to: {analysis4['vulnerable_to']}")

mev_protection_demo()
\end{lstlisting}

%===========================================
\section*{Resumen del Bloque}
%===========================================

\begin{keyidea}
\textbf{Puntos clave de Private Smart Contracts (Bloque 17):}

\textbf{1. Aproximaciones principales:}
\begin{itemize}
    \item ZK (Aleo, Mina): Pura criptografía, sin trust en hardware
    \item TEE (Secret, Oasis): Hardware trust, baja latencia
    \item MPC: Trust distribuido, threshold de partes honestas
    \item FHE (Fhenix): Computación sobre encriptado, muy costoso
\end{itemize}

\textbf{2. Secret Network:}
\begin{itemize}
    \item Estado encriptado con Intel SGX
    \item Viewing keys para acceso selectivo
    \item CosmWasm modificado (Rust)
\end{itemize}

\textbf{3. Oasis Sapphire:}
\begin{itemize}
    \item EVM confidencial
    \item Compatible con herramientas Ethereum
    \item TEE para privacidad
\end{itemize}

\textbf{4. Aleo:}
\begin{itemize}
    \item ZK-native, sin TEE
    \item Modelo de records (como UTXO)
    \item Lenguaje Leo
\end{itemize}

\textbf{5. Mina:}
\begin{itemize}
    \item Blockchain de 22 KB (recursive SNARKs)
    \item zkApps con estado off-chain
    \item o1js para desarrollo
\end{itemize}

\textbf{6. Encrypted Mempools:}
\begin{itemize}
    \item Threshold encryption (Shutter)
    \item Timelock encryption
    \item Commit-reveal
    \item Previene frontrunning y MEV
\end{itemize}

\textbf{7. MEV Protection:}
\begin{itemize}
    \item Private order flow (básico)
    \item Encrypted mempools (alto)
    \item Threshold + timelock (máximo)
\end{itemize}
\end{keyidea}

\begin{bancozk}
\textbf{Private Smart Contracts en Banco ZK:}

\begin{center}
\begin{tikzpicture}[
    layer/.style={rectangle, draw, minimum width=10cm, minimum height=0.8cm, align=center}
]
    \node[layer, fill=purple!20] at (0, 3) {zkApps: Lógica de negocio con ZK proofs};
    \node[layer, fill=blue!20] at (0, 2) {Estado: Off-chain (Merkle trees) + Commitments on-chain};
    \node[layer, fill=green!20] at (0, 1) {Mempool: Encrypted con threshold decryption};
    \node[layer, fill=orange!20] at (0, 0) {Consensus: PoS con MEV protection integrado};
\end{tikzpicture}
\end{center}

\textbf{Casos de uso en Banco ZK:}
\begin{enumerate}
    \item \textbf{Pagos privados:}
    \begin{itemize}
        \item Transferencias con montos ocultos (ZK)
        \item Balances off-chain
    \end{itemize}
    
    \item \textbf{DeFi privado:}
    \begin{itemize}
        \item Swaps sin frontrunning (encrypted mempool)
        \item Lending sin revelar posiciones
    \end{itemize}
    
    \item \textbf{Votación privada:}
    \begin{itemize}
        \item Votos confidenciales
        \item Conteo verificable
    \end{itemize}
    
    \item \textbf{Subastas selladas:}
    \begin{itemize}
        \item Pujas encriptadas
        \item Revelación solo al cierre
    \end{itemize}
\end{enumerate}

\textbf{Stack técnico recomendado:}
\begin{itemize}
    \item Core privacy: ZK proofs (Groth16/PLONK)
    \item Mempool protection: Threshold encryption
    \item Fallback TEE: Para operaciones legacy
    \item Future: FHE cuando sea práctico
\end{itemize}
\end{bancozk}

% Fin del Bloque 17
\cleardoublepage

% ============================================
% BLOQUE 18: CRIPTOGRAFÍA AVANZADA - VERSIÓN EXPANDIDA
% Parte 1: Introducción, FHE, MPC
% Objetivo: ~60 páginas (~3600 líneas)
% ============================================

\chapter{Criptografía Avanzada para Privacidad}
\label{ch:bloque18}

\begin{center}
\textit{Bloque 18 de 20}
\end{center}

\section*{Introducción}

Más allá de las pruebas de conocimiento cero, existe un arsenal de técnicas criptográficas avanzadas que permiten computación privada. Este bloque explora las herramientas fundamentales que hacen posible la privacidad en sistemas distribuidos.

\begin{keyidea}
\textbf{Panorama de Criptografía para Privacidad:}

\begin{center}
\begin{tikzpicture}[
    node distance=1cm,
    tech/.style={rectangle, draw, rounded corners, minimum width=2.5cm, minimum height=1cm, align=center, font=\small}
]
    % Categorías
    \node[tech, fill=blue!20] (zk) at (0, 3) {Zero-Knowledge\\Proofs};
    \node[tech, fill=green!20] (fhe) at (3.5, 3) {Homomorphic\\Encryption};
    \node[tech, fill=orange!20] (mpc) at (7, 3) {Multi-Party\\Computation};
    \node[tech, fill=purple!20] (tee) at (10.5, 3) {Trusted\\Execution};
    
    % Propiedades
    \node at (0, 1.5) [font=\small, align=center] {Probar sin\\revelar};
    \node at (3.5, 1.5) [font=\small, align=center] {Computar sobre\\encriptado};
    \node at (7, 1.5) [font=\small, align=center] {Computación\\distribuida};
    \node at (10.5, 1.5) [font=\small, align=center] {Hardware\\aislado};
    
    % Trust model
    \node at (0, 0.5) [font=\tiny] {Trust: Crypto};
    \node at (3.5, 0.5) [font=\tiny] {Trust: Crypto};
    \node at (7, 0.5) [font=\tiny] {Trust: Threshold};
    \node at (10.5, 0.5) [font=\tiny] {Trust: Hardware};
\end{tikzpicture}
\end{center}

\textbf{Primitivas fundamentales:}
\begin{itemize}
    \item \textbf{Secret Sharing:} Dividir secretos entre múltiples partes
    \item \textbf{Oblivious Transfer:} Transferencia selectiva sin revelar elección
    \item \textbf{Commitments:} Comprometerse a un valor sin revelarlo
    \item \textbf{Threshold Cryptography:} Operaciones que requieren k-de-n partes
\end{itemize}
\end{keyidea}

%===========================================
\section{FHE (Fully Homomorphic Encryption)}
\label{sec:fhe}
%===========================================

\subsection{Fundamentos Teóricos}

\begin{definition}[Fully Homomorphic Encryption]
Sistema de encriptación que permite realizar operaciones arbitrarias sobre datos cifrados:

\textbf{Propiedades homomórficas:}
\begin{align}
\text{Dec}(\text{Enc}(a) \oplus \text{Enc}(b)) &= a + b \quad \text{(Adición)} \\
\text{Dec}(\text{Enc}(a) \otimes \text{Enc}(b)) &= a \cdot b \quad \text{(Multiplicación)}
\end{align}

Con suma y multiplicación, se puede computar \textbf{cualquier función} (circuito booleano/aritmético).
\end{definition}

\begin{mathnote}[Historia y Generaciones de FHE]
\textbf{El Santo Grial de la criptografía:}

Propuesto conceptualmente por Rivest, Adleman, Dertouzos en 1978, pero sin construcción práctica hasta 2009.

\textbf{Generación 1 (2009): Gentry}
\begin{itemize}
    \item Primera construcción teórica completa
    \item Basado en lattices ideales
    \item Bootstrapping para ``refrescar'' ruido
    \item Extremadamente lento ($>10^6$ más lento que plaintext)
\end{itemize}

\textbf{Generación 2 (2011-2014): BGV, BFV}
\begin{itemize}
    \item Brakerski-Gentry-Vaikuntanathan (BGV)
    \item Brakerski/Fan-Vercauteren (BFV)
    \item Operaciones sobre enteros/polinomios
    \item Leveled FHE (profundidad limitada sin bootstrap)
    \item $\sim10^3-10^4$ más lento
\end{itemize}

\textbf{Generación 3 (2016+): TFHE, CKKS}
\begin{itemize}
    \item TFHE: Operaciones bit-a-bit, bootstrap rápido (~10ms)
    \item CKKS: Aritmética aproximada para ML/estadísticas
    \item Práctico para casos de uso específicos
    \item $\sim10^2-10^3$ más lento
\end{itemize}

\textbf{Generación 4 (Actual): Programable}
\begin{itemize}
    \item Compiladores FHE (HEIR, Concrete)
    \item Aceleradores hardware (GPU, FPGA, ASIC)
    \item Objetivo: $<10\times$ overhead
\end{itemize}
\end{mathnote}

\begin{figure}[htbp]
\centering
\begin{tikzpicture}[
    node distance=1cm,
    block/.style={rectangle, draw, minimum width=2.5cm, minimum height=1cm, align=center}
]
    % Flujo FHE
    \node[block, fill=blue!20] (plain) at (0, 0) {Datos\\$m_1, m_2$};
    \node[block, fill=green!20] (enc) at (4, 0) {Encriptados\\$c_1, c_2$};
    \node[block, fill=orange!20] (compute) at (8, 0) {Computar\\$c_1 \otimes c_2$};
    \node[block, fill=purple!20] (result) at (12, 0) {Resultado\\$\text{Enc}(m_1 \cdot m_2)$};
    
    \draw[->, thick] (plain) -- node[above, font=\small] {Encrypt} (enc);
    \draw[->, thick] (enc) -- node[above, font=\small] {Operate} (compute);
    \draw[->, thick] (compute) -- (result);
    
    % Decrypt
    \node[block, fill=blue!20] (final) at (12, -2) {$m_1 \cdot m_2$};
    \draw[->, thick] (result) -- node[right, font=\small] {Decrypt} (final);
    
    % Nota
    \node at (6, -2.5) [font=\small, align=center] {
        El servidor NUNCA ve los datos en claro\\
        Solo opera sobre ciphertexts
    };
\end{tikzpicture}
\caption{Flujo de computación con FHE}
\label{fig:fhe-flow}
\end{figure}

\subsection{Matemáticas de FHE (Learning With Errors)}

\begin{definition}[Learning With Errors (LWE)]
Problema computacional base de la mayoría de esquemas FHE modernos:

\textbf{Problema:} Dado $(A, b = As + e \mod q)$, encontrar $s$.

Donde:
\begin{itemize}
    \item $A \in \mathbb{Z}_q^{m \times n}$: Matriz pública aleatoria
    \item $s \in \mathbb{Z}_q^n$: Vector secreto
    \item $e \in \mathbb{Z}_q^m$: Vector de ``error'' pequeño
    \item $q$: Módulo grande
\end{itemize}

\textbf{Dificultad:} Sin $e$, es álgebra lineal. Con $e$, se cree intratrable incluso para computadoras cuánticas.
\end{definition}

\begin{mathnote}[Esquema BFV Simplificado]
\textbf{Parámetros:}
\begin{itemize}
    \item $n$: Grado del polinomio (potencia de 2, ej: 4096)
    \item $q$: Módulo del ciphertext (grande, ej: $2^{109}$)
    \item $t$: Módulo del plaintext (pequeño, ej: $2^{16}$)
    \item $R = \mathbb{Z}[X]/(X^n + 1)$: Anillo de polinomios
\end{itemize}

\textbf{KeyGen:}
\begin{enumerate}
    \item $sk \leftarrow$ polinomio con coeficientes pequeños
    \item $a \leftarrow R_q$ aleatorio
    \item $e \leftarrow$ error pequeño (distribución gaussiana)
    \item $pk = (-a \cdot sk + e, a) = (b, a)$
\end{enumerate}

\textbf{Encrypt($m$, $pk$):}
\begin{enumerate}
    \item $m \in R_t$ (mensaje como polinomio)
    \item $u, e_1, e_2 \leftarrow$ errores pequeños
    \item $\Delta = \lfloor q/t \rfloor$ (factor de escala)
    \item $c_0 = b \cdot u + e_1 + \Delta \cdot m$
    \item $c_1 = a \cdot u + e_2$
    \item Ciphertext: $(c_0, c_1)$
\end{enumerate}

\textbf{Decrypt($c$, $sk$):}
\begin{enumerate}
    \item $m' = c_0 + c_1 \cdot sk \mod q$
    \item $m = \lfloor t \cdot m' / q \rceil \mod t$
\end{enumerate}

\textbf{Operaciones Homomórficas:}

\textit{Suma:} $(c_0, c_1) + (c'_0, c'_1) = (c_0 + c'_0, c_1 + c'_1)$

\textit{Multiplicación:} Más compleja, genera ciphertext de 3 elementos, requiere ``relinearization'' para volver a 2 elementos.
\end{mathnote}

\begin{lstlisting}[language=Python, caption={Implementación conceptual de FHE (BFV simplificado)}]
import numpy as np
from typing import Tuple, List
from dataclasses import dataclass

@dataclass
class FHEParams:
    """Parámetros del esquema FHE."""
    n: int  # Grado del polinomio
    q: int  # Módulo del ciphertext
    t: int  # Módulo del plaintext
    std_dev: float  # Desviación estándar del error

class SimpleFHE:
    """
    Implementación simplificada de FHE estilo BFV.
    NOTA: Solo para demostración, NO usar en producción.
    """
    
    def __init__(self, params: FHEParams):
        self.params = params
        self.n = params.n
        self.q = params.q
        self.t = params.t
        self.delta = self.q // self.t  # Factor de escala
        
        # Generar claves
        self.secret_key = None
        self.public_key = None
        self.relin_key = None
    
    def _sample_error(self, size: int = None) -> np.ndarray:
        """Muestrear error de distribución gaussiana discreta."""
        size = size or self.n
        return np.round(np.random.normal(0, self.params.std_dev, size)).astype(int)
    
    def _sample_uniform(self, size: int = None) -> np.ndarray:
        """Muestrear uniformemente del anillo."""
        size = size or self.n
        return np.random.randint(0, self.q, size)
    
    def _sample_ternary(self, size: int = None) -> np.ndarray:
        """Muestrear de {-1, 0, 1}."""
        size = size or self.n
        return np.random.choice([-1, 0, 1], size)
    
    def _poly_mult(self, a: np.ndarray, b: np.ndarray) -> np.ndarray:
        """
        Multiplicación de polinomios en R_q = Z_q[X]/(X^n + 1).
        """
        # Convolución
        result = np.convolve(a, b)
        
        # Reducción módulo X^n + 1 (negacyclic)
        for i in range(self.n, len(result)):
            result[i - self.n] -= result[i]
        
        return result[:self.n] % self.q
    
    def keygen(self) -> Tuple[np.ndarray, Tuple[np.ndarray, np.ndarray]]:
        """
        Generar par de claves.
        
        Returns:
            (secret_key, public_key)
        """
        # Secret key: polinomio ternario
        sk = self._sample_ternary()
        
        # Public key
        a = self._sample_uniform()
        e = self._sample_error()
        
        # b = -a*sk + e (mod q)
        b = (-self._poly_mult(a, sk) + e) % self.q
        
        pk = (b, a)
        
        self.secret_key = sk
        self.public_key = pk
        
        # Generar relinearization key (simplificado)
        self._gen_relin_key()
        
        return sk, pk
    
    def _gen_relin_key(self):
        """Generar clave de relinearización para multiplicación."""
        # Simplificado: en BFV real, es más complejo
        sk_squared = self._poly_mult(self.secret_key, self.secret_key)
        a = self._sample_uniform()
        e = self._sample_error()
        
        b = (-self._poly_mult(a, self.secret_key) + e + sk_squared) % self.q
        self.relin_key = (b, a)
    
    def encrypt(self, message: int) -> Tuple[np.ndarray, np.ndarray]:
        """
        Encriptar un mensaje entero.
        
        Args:
            message: Entero en [0, t)
        
        Returns:
            Ciphertext (c0, c1)
        """
        assert 0 <= message < self.t, f"Message must be in [0, {self.t})"
        
        # Codificar mensaje como polinomio (constante)
        m = np.zeros(self.n, dtype=int)
        m[0] = message
        
        # Sampling
        u = self._sample_ternary()
        e1 = self._sample_error()
        e2 = self._sample_error()
        
        b, a = self.public_key
        
        # c0 = b*u + e1 + delta*m
        c0 = (self._poly_mult(b, u) + e1 + self.delta * m) % self.q
        
        # c1 = a*u + e2
        c1 = (self._poly_mult(a, u) + e2) % self.q
        
        return (c0, c1)
    
    def decrypt(self, ciphertext: Tuple[np.ndarray, np.ndarray]) -> int:
        """
        Descifrar ciphertext.
        
        Args:
            ciphertext: (c0, c1)
        
        Returns:
            Mensaje descifrado
        """
        c0, c1 = ciphertext
        
        # m' = c0 + c1*sk (mod q)
        m_scaled = (c0 + self._poly_mult(c1, self.secret_key)) % self.q
        
        # Centrar en [-q/2, q/2)
        m_scaled = np.where(m_scaled > self.q // 2, m_scaled - self.q, m_scaled)
        
        # Escalar de vuelta: m = round(t * m' / q)
        m = np.round(self.t * m_scaled / self.q) % self.t
        
        return int(m[0])
    
    def add(
        self,
        ct1: Tuple[np.ndarray, np.ndarray],
        ct2: Tuple[np.ndarray, np.ndarray]
    ) -> Tuple[np.ndarray, np.ndarray]:
        """
        Suma homomórfica de dos ciphertexts.
        
        Dec(add(Enc(a), Enc(b))) = a + b (mod t)
        """
        c0_1, c1_1 = ct1
        c0_2, c1_2 = ct2
        
        c0 = (c0_1 + c0_2) % self.q
        c1 = (c1_1 + c1_2) % self.q
        
        return (c0, c1)
    
    def multiply(
        self,
        ct1: Tuple[np.ndarray, np.ndarray],
        ct2: Tuple[np.ndarray, np.ndarray]
    ) -> Tuple[np.ndarray, np.ndarray]:
        """
        Multiplicación homomórfica de dos ciphertexts.
        
        Dec(multiply(Enc(a), Enc(b))) = a * b (mod t)
        """
        c0_1, c1_1 = ct1
        c0_2, c1_2 = ct2
        
        # Producto tensor (3 elementos)
        d0 = self._poly_mult(c0_1, c0_2)
        d1 = (self._poly_mult(c0_1, c1_2) + self._poly_mult(c1_1, c0_2)) % self.q
        d2 = self._poly_mult(c1_1, c1_2)
        
        # Escalar por t/q para mantener el nivel de ruido
        scale = self.t / self.q
        d0 = np.round(d0 * scale) % self.q
        d1 = np.round(d1 * scale) % self.q
        d2 = np.round(d2 * scale) % self.q
        
        # Relinearización: reducir de 3 a 2 elementos
        c0, c1 = self._relinearize(d0.astype(int), d1.astype(int), d2.astype(int))
        
        return (c0, c1)
    
    def _relinearize(
        self,
        d0: np.ndarray,
        d1: np.ndarray,
        d2: np.ndarray
    ) -> Tuple[np.ndarray, np.ndarray]:
        """
        Relinearización: convertir ciphertext de 3 elementos a 2.
        """
        rlk_b, rlk_a = self.relin_key
        
        # c0 = d0 + d2 * rlk_b
        c0 = (d0 + self._poly_mult(d2, rlk_b)) % self.q
        
        # c1 = d1 + d2 * rlk_a
        c1 = (d1 + self._poly_mult(d2, rlk_a)) % self.q
        
        return (c0, c1)
    
    def add_plain(
        self,
        ct: Tuple[np.ndarray, np.ndarray],
        plaintext: int
    ) -> Tuple[np.ndarray, np.ndarray]:
        """Sumar plaintext a ciphertext."""
        c0, c1 = ct
        
        m = np.zeros(self.n, dtype=int)
        m[0] = plaintext * self.delta
        
        new_c0 = (c0 + m) % self.q
        
        return (new_c0, c1)
    
    def multiply_plain(
        self,
        ct: Tuple[np.ndarray, np.ndarray],
        plaintext: int
    ) -> Tuple[np.ndarray, np.ndarray]:
        """Multiplicar ciphertext por plaintext."""
        c0, c1 = ct
        
        new_c0 = (c0 * plaintext) % self.q
        new_c1 = (c1 * plaintext) % self.q
        
        return (new_c0, new_c1)


def fhe_demo():
    print("=== FHE Demo (BFV Simplificado) ===\n")
    
    # Parámetros (pequeños para demo, no seguros)
    params = FHEParams(
        n=64,           # Normalmente 4096+
        q=2**20,        # Normalmente 2^109+
        t=256,          # Plaintext modulus
        std_dev=3.2     # Error estándar
    )
    
    fhe = SimpleFHE(params)
    
    # Generar claves
    print("--- Key Generation ---")
    sk, pk = fhe.keygen()
    print(f"Secret key: {sk[:5]}... (ternary polynomial)")
    print(f"Public key: 2 polynomials of degree {params.n}")
    
    # Encriptar valores
    print("\n--- Encryption ---")
    a, b = 15, 7
    
    ct_a = fhe.encrypt(a)
    ct_b = fhe.encrypt(b)
    
    print(f"Encrypted {a} and {b}")
    print(f"Ciphertext size: 2 × {params.n} integers mod {params.q}")
    
    # Operaciones homomórficas
    print("\n--- Homomorphic Operations ---")
    
    # Suma
    ct_sum = fhe.add(ct_a, ct_b)
    result_sum = fhe.decrypt(ct_sum)
    print(f"Dec(Enc({a}) + Enc({b})) = {result_sum}")
    print(f"Expected: {(a + b) % params.t}")
    
    # Multiplicación
    ct_mult = fhe.multiply(ct_a, ct_b)
    result_mult = fhe.decrypt(ct_mult)
    print(f"Dec(Enc({a}) × Enc({b})) = {result_mult}")
    print(f"Expected: {(a * b) % params.t}")
    
    # Operaciones mixtas
    print("\n--- Mixed Operations ---")
    ct_expr = fhe.add(fhe.multiply(ct_a, ct_b), ct_a)  # a*b + a
    result_expr = fhe.decrypt(ct_expr)
    print(f"Dec(Enc({a}) × Enc({b}) + Enc({a})) = {result_expr}")
    print(f"Expected: {(a * b + a) % params.t}")
    
    print("\n--- Noise Growth ---")
    print("After each operation, noise grows.")
    print("Too many operations → decryption fails.")
    print("Bootstrapping 'refreshes' noise (expensive).")

fhe_demo()
\end{lstlisting}

\subsection{Librerías y Aplicaciones}

\begin{table}[htbp]
\centering
\small
\begin{tabular}{llp{5cm}}
\toprule
\textbf{Librería} & \textbf{Esquema} & \textbf{Características} \\
\midrule
Microsoft SEAL & BFV, CKKS & C++, muy documentado, usado en industria \\
OpenFHE & BGV, BFV, CKKS, TFHE & C++, sucesor de PALISADE, más completo \\
TFHE-rs & TFHE & Rust, bootstrap rápido, Zama \\
Concrete & TFHE & Python/Rust, compilador FHE, Zama \\
HElib & BGV, CKKS & C++, IBM, muy optimizado \\
Lattigo & BFV, CKKS & Go, buen rendimiento \\
\bottomrule
\end{tabular}
\caption{Principales librerías FHE}
\end{table}

%===========================================
\section{MPC (Multi-Party Computation)}
\label{sec:mpc}
%===========================================

\subsection{Fundamentos}

\begin{definition}[Secure Multi-Party Computation]
Protocolo que permite a $n$ partes computar una función $f(x_1, ..., x_n)$ sobre sus inputs privados, tal que:

\begin{enumerate}
    \item \textbf{Correctitud:} El output es correcto
    \item \textbf{Privacidad:} Ninguna parte aprende más que su propio input y el output
    \item \textbf{Independence of inputs:} Inputs elegidos independientemente
\end{enumerate}

\textbf{Ideal functionality:} Como si existiera un tercero de confianza que recibe todos los inputs, computa $f$, y entrega outputs.
\end{definition}

\begin{figure}[htbp]
\centering
\begin{tikzpicture}[
    node distance=1.5cm,
    party/.style={circle, draw, minimum size=1.2cm, fill=blue!20},
    func/.style={rectangle, draw, rounded corners, minimum width=2cm, minimum height=1.5cm, fill=green!20}
]
    % Partes
    \node[party] (p1) at (0, 2) {$P_1$\\$x_1$};
    \node[party] (p2) at (0, 0) {$P_2$\\$x_2$};
    \node[party] (p3) at (0, -2) {$P_3$\\$x_3$};
    
    % Función
    \node[func] (f) at (5, 0) {$f(x_1, x_2, x_3)$};
    
    % Outputs
    \node[party] (o1) at (10, 2) {$P_1$\\$y$};
    \node[party] (o2) at (10, 0) {$P_2$\\$y$};
    \node[party] (o3) at (10, -2) {$P_3$\\$y$};
    
    % Conexiones (protocolo MPC)
    \draw[->, thick, dashed] (p1) -- (f);
    \draw[->, thick, dashed] (p2) -- (f);
    \draw[->, thick, dashed] (p3) -- (f);
    
    \draw[->, thick] (f) -- (o1);
    \draw[->, thick] (f) -- (o2);
    \draw[->, thick] (f) -- (o3);
    
    % Notas
    \node at (2.5, 1) [font=\small] {Privado};
    \node at (7.5, 1) [font=\small] {Resultado};
    
    \node at (5, -3.5) [font=\small, align=center] {
        Ninguna parte ve los inputs de otras\\
        Todas obtienen el mismo resultado correcto
    };
\end{tikzpicture}
\caption{Ideal functionality de MPC}
\label{fig:mpc-ideal}
\end{figure}

\subsection{Protocolos MPC Principales}

\begin{mathnote}[Garbled Circuits (Yao, 1986)]
Protocolo para 2 partes basado en ``ofuscar'' un circuito booleano:

\textbf{Partes:}
\begin{itemize}
    \item Garbler: Construye el circuito ofuscado
    \item Evaluator: Evalúa el circuito
\end{itemize}

\textbf{Proceso:}
\begin{enumerate}
    \item Garbler asigna 2 etiquetas aleatorias a cada wire: $W^0_i, W^1_i$
    \item Para cada gate, crea tabla ofuscada de 4 entradas (encriptada)
    \item Evaluator recibe etiquetas de sus inputs via Oblivious Transfer
    \item Evaluator evalúa gate por gate, descifrando la entrada correcta
    \item Output: etiqueta final mapea a bit de salida
\end{enumerate}

\textbf{Seguridad:} Evaluator solo ve una etiqueta por wire (no sabe si es 0 o 1 hasta el final).
\end{mathnote}

\begin{lstlisting}[language=Python, caption={Implementación de Garbled Circuits}]
from dataclasses import dataclass
from typing import Dict, List, Tuple, Optional
import secrets
import hashlib
from enum import Enum

class GateType(Enum):
    AND = "AND"
    OR = "OR"
    XOR = "XOR"
    NOT = "NOT"

@dataclass
class Wire:
    wire_id: int
    label_0: bytes  # Etiqueta para bit 0
    label_1: bytes  # Etiqueta para bit 1

@dataclass
class Gate:
    gate_id: int
    gate_type: GateType
    input_wires: List[int]
    output_wire: int

@dataclass
class GarbledGate:
    gate_id: int
    garbled_table: List[bytes]  # 4 entradas encriptadas

class GarbledCircuit:
    """
    Implementación de Garbled Circuits para MPC de 2 partes.
    """
    
    def __init__(self, label_size: int = 16):
        self.label_size = label_size
        self.wires: Dict[int, Wire] = {}
        self.gates: List[Gate] = []
        self.garbled_gates: List[GarbledGate] = []
        
        self.input_wires_garbler: List[int] = []
        self.input_wires_evaluator: List[int] = []
        self.output_wires: List[int] = []
    
    def _generate_label(self) -> bytes:
        """Generar etiqueta aleatoria."""
        return secrets.token_bytes(self.label_size)
    
    def _hash_labels(self, *labels: bytes, gate_id: int) -> bytes:
        """Hash de etiquetas para encriptar tabla."""
        data = b''.join(labels) + gate_id.to_bytes(4, 'big')
        return hashlib.sha256(data).digest()[:self.label_size]
    
    def add_wire(self, wire_id: int):
        """Añadir wire al circuito."""
        self.wires[wire_id] = Wire(
            wire_id=wire_id,
            label_0=self._generate_label(),
            label_1=self._generate_label()
        )
    
    def add_gate(
        self,
        gate_id: int,
        gate_type: GateType,
        input_wires: List[int],
        output_wire: int
    ):
        """Añadir gate al circuito."""
        # Crear wire de salida
        self.add_wire(output_wire)
        
        gate = Gate(
            gate_id=gate_id,
            gate_type=gate_type,
            input_wires=input_wires,
            output_wire=output_wire
        )
        self.gates.append(gate)
    
    def _evaluate_gate(
        self,
        gate_type: GateType,
        inputs: List[int]
    ) -> int:
        """Evaluar gate en plaintext."""
        if gate_type == GateType.AND:
            return inputs[0] & inputs[1]
        elif gate_type == GateType.OR:
            return inputs[0] | inputs[1]
        elif gate_type == GateType.XOR:
            return inputs[0] ^ inputs[1]
        elif gate_type == GateType.NOT:
            return 1 - inputs[0]
        return 0
    
    def garble(self) -> List[GarbledGate]:
        """
        Garbler: Ofuscar el circuito.
        """
        self.garbled_gates = []
        
        for gate in self.gates:
            garbled = self._garble_gate(gate)
            self.garbled_gates.append(garbled)
        
        return self.garbled_gates
    
    def _garble_gate(self, gate: Gate) -> GarbledGate:
        """Ofuscar una gate individual."""
        table = []
        
        # Para cada combinación de inputs
        for a in [0, 1]:
            for b in [0, 1] if len(gate.input_wires) > 1 else [0]:
                # Evaluar gate
                if len(gate.input_wires) > 1:
                    output_bit = self._evaluate_gate(gate.gate_type, [a, b])
                else:
                    output_bit = self._evaluate_gate(gate.gate_type, [a])
                
                # Obtener etiquetas de entrada
                wire_a = self.wires[gate.input_wires[0]]
                label_a = wire_a.label_1 if a else wire_a.label_0
                
                if len(gate.input_wires) > 1:
                    wire_b = self.wires[gate.input_wires[1]]
                    label_b = wire_b.label_1 if b else wire_b.label_0
                else:
                    label_b = b''
                
                # Obtener etiqueta de salida
                output_wire = self.wires[gate.output_wire]
                output_label = output_wire.label_1 if output_bit else output_wire.label_0
                
                # Encriptar etiqueta de salida con etiquetas de entrada
                key = self._hash_labels(label_a, label_b, gate_id=gate.gate_id)
                encrypted = bytes(o ^ k for o, k in zip(output_label, key))
                
                table.append(encrypted)
        
        # Permutar tabla aleatoriamente
        secrets.SystemRandom().shuffle(table)
        
        return GarbledGate(gate_id=gate.gate_id, garbled_table=table)
    
    def get_garbler_input_labels(
        self,
        garbler_inputs: List[int]
    ) -> List[bytes]:
        """
        Garbler obtiene las etiquetas para sus propios inputs.
        """
        labels = []
        for wire_id, bit in zip(self.input_wires_garbler, garbler_inputs):
            wire = self.wires[wire_id]
            label = wire.label_1 if bit else wire.label_0
            labels.append(label)
        return labels
    
    def get_evaluator_input_labels_via_ot(
        self,
        evaluator_inputs: List[int]
    ) -> List[bytes]:
        """
        Evaluator obtiene etiquetas via Oblivious Transfer.
        (Simplificado - en realidad sería OT protocol)
        """
        labels = []
        for wire_id, bit in zip(self.input_wires_evaluator, evaluator_inputs):
            wire = self.wires[wire_id]
            label = wire.label_1 if bit else wire.label_0
            labels.append(label)
        return labels
    
    def evaluate(
        self,
        garbled_gates: List[GarbledGate],
        input_labels: Dict[int, bytes]
    ) -> Dict[int, bytes]:
        """
        Evaluator: Evaluar el circuito ofuscado.
        """
        wire_labels = input_labels.copy()
        
        for garbled_gate in garbled_gates:
            gate = self.gates[garbled_gate.gate_id]
            
            # Obtener etiquetas de entrada
            label_a = wire_labels[gate.input_wires[0]]
            if len(gate.input_wires) > 1:
                label_b = wire_labels[gate.input_wires[1]]
            else:
                label_b = b''
            
            # Intentar descifrar cada entrada de la tabla
            output_label = None
            for encrypted in garbled_gate.garbled_table:
                key = self._hash_labels(label_a, label_b, gate_id=gate.gate_id)
                decrypted = bytes(e ^ k for e, k in zip(encrypted, key))
                
                # Verificar si es una etiqueta válida
                # (En implementación real, se usa punto de verificación)
                output_wire = self.wires[gate.output_wire]
                if decrypted == output_wire.label_0 or decrypted == output_wire.label_1:
                    output_label = decrypted
                    break
            
            if output_label:
                wire_labels[gate.output_wire] = output_label
        
        return wire_labels
    
    def decode_output(
        self,
        output_labels: Dict[int, bytes]
    ) -> List[int]:
        """
        Decodificar etiquetas de salida a bits.
        """
        outputs = []
        for wire_id in self.output_wires:
            wire = self.wires[wire_id]
            label = output_labels[wire_id]
            
            if label == wire.label_0:
                outputs.append(0)
            elif label == wire.label_1:
                outputs.append(1)
            else:
                raise ValueError(f"Invalid output label for wire {wire_id}")
        
        return outputs


class GMWProtocol:
    """
    Protocolo GMW para MPC basado en secret sharing.
    Soporta múltiples partes.
    """
    
    def __init__(self, num_parties: int):
        self.num_parties = num_parties
    
    def share_value(self, value: int, num_bits: int = 1) -> List[int]:
        """
        Secret-share un valor entre todas las partes.
        XOR sharing: value = share_1 XOR share_2 XOR ... XOR share_n
        """
        shares = []
        running_xor = value
        
        for i in range(self.num_parties - 1):
            share = secrets.randbelow(2**num_bits)
            shares.append(share)
            running_xor ^= share
        
        # Última share asegura que XOR de todas = value
        shares.append(running_xor)
        
        return shares
    
    def reconstruct_value(self, shares: List[int]) -> int:
        """Reconstruir valor desde shares."""
        result = 0
        for share in shares:
            result ^= share
        return result
    
    def xor_gate(
        self,
        shares_a: List[int],
        shares_b: List[int]
    ) -> List[int]:
        """
        XOR gate: cada parte XORea sus shares localmente.
        """
        return [a ^ b for a, b in zip(shares_a, shares_b)]
    
    def and_gate(
        self,
        shares_a: List[int],
        shares_b: List[int]
    ) -> List[int]:
        """
        AND gate: requiere comunicación (simplificado).
        En GMW real, usa multiplication triples pre-computados.
        """
        # Simplificación: reconstruir, computar, re-compartir
        a = self.reconstruct_value(shares_a)
        b = self.reconstruct_value(shares_b)
        result = a & b
        return self.share_value(result)


def garbled_circuits_demo():
    print("=== Garbled Circuits Demo ===\n")
    
    gc = GarbledCircuit()
    
    # Construir circuito: (a AND b) XOR c
    # Inputs: a (wire 0), b (wire 1), c (wire 2)
    # Output: wire 4
    
    gc.add_wire(0)  # a
    gc.add_wire(1)  # b
    gc.add_wire(2)  # c
    
    gc.input_wires_garbler = [0, 1]  # Garbler tiene a, b
    gc.input_wires_evaluator = [2]    # Evaluator tiene c
    gc.output_wires = [4]
    
    # Gates
    gc.add_gate(0, GateType.AND, [0, 1], 3)  # wire 3 = a AND b
    gc.add_gate(1, GateType.XOR, [3, 2], 4)  # wire 4 = (a AND b) XOR c
    
    # Garbler garbles el circuito
    print("--- Garbler Phase ---")
    garbled_gates = gc.garble()
    print(f"Garbled {len(garbled_gates)} gates")
    
    # Inputs
    a, b = 1, 1  # Garbler inputs
    c = 0        # Evaluator input
    
    expected = (a & b) ^ c
    print(f"\nExpected result: ({a} AND {b}) XOR {c} = {expected}")
    
    # Garbler obtiene etiquetas de sus inputs
    garbler_labels = gc.get_garbler_input_labels([a, b])
    
    # Evaluator obtiene etiquetas via OT (simplificado)
    evaluator_labels = gc.get_evaluator_input_labels_via_ot([c])
    
    # Combinar etiquetas
    input_labels = {
        0: garbler_labels[0],
        1: garbler_labels[1],
        2: evaluator_labels[0]
    }
    
    # Evaluator evalúa
    print("\n--- Evaluator Phase ---")
    output_labels = gc.evaluate(garbled_gates, input_labels)
    
    # Decodificar salida
    result = gc.decode_output(output_labels)
    print(f"Computed result: {result[0]}")
    print(f"Correct: {result[0] == expected}")

garbled_circuits_demo()
\end{lstlisting}

\subsection{Protocolo SPDZ}

\begin{definition}[SPDZ (Speedz)]
Protocolo MPC eficiente para múltiples partes con seguridad contra adversarios maliciosos:

\textbf{Características:}
\begin{itemize}
    \item Basado en secret sharing aditivo sobre campo finito
    \item Fase offline (preprocessing) + fase online (evaluación)
    \item Online muy rápido (~microsegundos por operación)
    \item Soporta número arbitrario de partes
\end{itemize}

\textbf{Modelo de adversario:}
\begin{itemize}
    \item Corrompe hasta $n-1$ partes (dishonest majority)
    \item Detecta trampas con alta probabilidad
\end{itemize}
\end{definition}

\begin{mathnote}[SPDZ Overview]
\textbf{Secret Sharing:} Valor $x$ compartido como:
\[
\langle x \rangle = (x_1, ..., x_n) \text{ donde } x = \sum_{i=1}^n x_i
\]

\textbf{Operaciones:}
\begin{itemize}
    \item \textbf{Suma:} $\langle x \rangle + \langle y \rangle = \langle x + y \rangle$ (local)
    \item \textbf{Mult por constante:} $c \cdot \langle x \rangle = \langle c \cdot x \rangle$ (local)
    \item \textbf{Multiplicación:} Requiere ``multiplication triples'' pre-computados
\end{itemize}

\textbf{Beaver Triples:} Pre-computar $(a, b, c)$ donde $c = a \cdot b$:
\begin{enumerate}
    \item Para multiplicar $\langle x \rangle \cdot \langle y \rangle$:
    \item Computar $\langle \epsilon \rangle = \langle x \rangle - \langle a \rangle$ y abrir $\epsilon$
    \item Computar $\langle \delta \rangle = \langle y \rangle - \langle b \rangle$ y abrir $\delta$
    \item $\langle xy \rangle = \langle c \rangle + \epsilon \langle b \rangle + \delta \langle a \rangle + \epsilon \delta$
\end{enumerate}
\end{mathnote}

\begin{table}[htbp]
\centering
\small
\begin{tabular}{lcccc}
\toprule
\textbf{Protocolo} & \textbf{Partes} & \textbf{Modelo} & \textbf{Complejidad} & \textbf{Uso} \\
\midrule
Garbled Circuits & 2 & Semi-honest & $O(|C|)$ & General \\
GMW & n & Semi-honest & $O(|C| \cdot n^2)$ & Booleano \\
SPDZ & n & Malicious & $O(|C|)$ online & Aritmético \\
BGW & n & Honest majority & $O(|C|)$ & Aritmético \\
\bottomrule
\end{tabular}
\caption{Comparación de protocolos MPC}
\end{table}
%===========================================
\section{Secret Sharing (Shamir)}
\label{sec:secret-sharing}
%===========================================

\subsection{Fundamentos Matemáticos}

\begin{definition}[Shamir's Secret Sharing]
Esquema $(k, n)$-threshold que divide un secreto $s$ en $n$ partes (shares), tal que:
\begin{itemize}
    \item Cualquier $k$ o más shares pueden reconstruir $s$
    \item Menos de $k$ shares no revelan información sobre $s$
\end{itemize}

\textbf{Basado en:} Interpolación de polinomios de Lagrange.
\end{definition}

\begin{mathnote}[Construcción de Shamir]
\textbf{Setup:}
\begin{itemize}
    \item Secreto: $s \in \mathbb{Z}_p$ (campo primo grande)
    \item Threshold: $k$ (mínimo de shares para reconstruir)
    \item Total shares: $n$
\end{itemize}

\textbf{Compartir (Split):}
\begin{enumerate}
    \item Elegir coeficientes aleatorios $a_1, ..., a_{k-1} \in \mathbb{Z}_p$
    \item Construir polinomio de grado $k-1$:
    \[
    f(x) = s + a_1 x + a_2 x^2 + ... + a_{k-1} x^{k-1}
    \]
    \item Share $i$: $(i, f(i))$ para $i = 1, 2, ..., n$
\end{enumerate}

\textbf{Reconstruir:}

Dados $k$ puntos $(x_1, y_1), ..., (x_k, y_k)$, usar interpolación de Lagrange:
\[
f(x) = \sum_{i=1}^{k} y_i \cdot L_i(x)
\]
donde el polinomio base de Lagrange es:
\[
L_i(x) = \prod_{j \neq i} \frac{x - x_j}{x_i - x_j}
\]

El secreto es $s = f(0)$.
\end{mathnote}

\begin{figure}[htbp]
\centering
\begin{tikzpicture}[scale=0.8]
    % Ejes
    \draw[->] (-0.5, 0) -- (6, 0) node[right] {$x$};
    \draw[->] (0, -0.5) -- (0, 5) node[above] {$y$};
    
    % Polinomio (cuadrático: s + ax + bx^2)
    \draw[blue, thick, domain=0:5.5, samples=50] plot (\x, {1.5 + 0.3*\x + 0.15*\x*\x});
    
    % Secreto
    \fill[red] (0, 1.5) circle (4pt);
    \node at (-0.5, 1.5) [red] {$s$};
    
    % Shares
    \fill[green!60!black] (1, 1.95) circle (4pt);
    \fill[green!60!black] (2, 2.7) circle (4pt);
    \fill[green!60!black] (3, 3.75) circle (4pt);
    \fill[green!60!black] (4, 5.1) circle (4pt);
    
    \node at (1, 1.5) [font=\small] {share 1};
    \node at (2, 2.2) [font=\small] {share 2};
    \node at (3, 3.2) [font=\small] {share 3};
    \node at (4, 4.6) [font=\small] {share 4};
    
    % Leyenda
    \node at (3, -1) [font=\small] {$k=3$ shares necesarios para reconstruir $s=f(0)$};
\end{tikzpicture}
\caption{Shamir's Secret Sharing: el secreto es el término independiente del polinomio}
\label{fig:shamir-visual}
\end{figure}

\begin{lstlisting}[language=Python, caption={Implementación completa de Shamir's Secret Sharing}]
from dataclasses import dataclass
from typing import List, Tuple
import secrets
from functools import reduce

@dataclass
class Share:
    """Un share del secreto."""
    x: int  # Punto de evaluación
    y: int  # Valor f(x)

class ShamirSecretSharing:
    """
    Implementación de Shamir's Secret Sharing sobre campo primo.
    """
    
    def __init__(self, prime: int = None):
        """
        Args:
            prime: Primo para el campo. Si None, usa primo de 256 bits.
        """
        # Primo de Mersenne para eficiencia
        self.prime = prime or (2**256 - 189)
    
    def _mod_inverse(self, a: int, p: int) -> int:
        """
        Calcular inverso modular usando extended Euclidean algorithm.
        """
        def extended_gcd(a: int, b: int) -> Tuple[int, int, int]:
            if a == 0:
                return b, 0, 1
            gcd, x1, y1 = extended_gcd(b % a, a)
            x = y1 - (b // a) * x1
            y = x1
            return gcd, x, y
        
        _, x, _ = extended_gcd(a % p, p)
        return (x % p + p) % p
    
    def split(
        self,
        secret: int,
        threshold: int,
        num_shares: int
    ) -> List[Share]:
        """
        Dividir secreto en shares.
        
        Args:
            secret: El secreto a compartir (entero)
            threshold: Número mínimo de shares para reconstruir (k)
            num_shares: Número total de shares a generar (n)
        
        Returns:
            Lista de n shares
        """
        assert threshold <= num_shares, "Threshold cannot exceed num_shares"
        assert 0 <= secret < self.prime, f"Secret must be in [0, {self.prime})"
        
        # Generar coeficientes aleatorios para el polinomio
        # f(x) = secret + a1*x + a2*x^2 + ... + a_{k-1}*x^{k-1}
        coefficients = [secret]
        for _ in range(threshold - 1):
            coefficients.append(secrets.randbelow(self.prime))
        
        # Evaluar polinomio en puntos 1, 2, ..., n
        shares = []
        for x in range(1, num_shares + 1):
            y = self._evaluate_polynomial(coefficients, x)
            shares.append(Share(x=x, y=y))
        
        return shares
    
    def _evaluate_polynomial(
        self,
        coefficients: List[int],
        x: int
    ) -> int:
        """
        Evaluar polinomio en punto x usando Horner's method.
        """
        result = 0
        for coef in reversed(coefficients):
            result = (result * x + coef) % self.prime
        return result
    
    def reconstruct(self, shares: List[Share]) -> int:
        """
        Reconstruir secreto desde shares usando interpolación de Lagrange.
        
        Args:
            shares: Lista de al menos k shares
        
        Returns:
            El secreto original
        """
        if len(shares) < 2:
            raise ValueError("Need at least 2 shares")
        
        # Interpolación de Lagrange en x=0
        secret = 0
        
        for i, share_i in enumerate(shares):
            # Calcular término de Lagrange L_i(0)
            numerator = 1
            denominator = 1
            
            for j, share_j in enumerate(shares):
                if i != j:
                    numerator = (numerator * (-share_j.x)) % self.prime
                    denominator = (denominator * (share_i.x - share_j.x)) % self.prime
            
            # L_i(0) = numerator / denominator
            lagrange = (numerator * self._mod_inverse(denominator, self.prime)) % self.prime
            
            # Añadir contribución: y_i * L_i(0)
            secret = (secret + share_i.y * lagrange) % self.prime
        
        return secret
    
    def verify_share(
        self,
        share: Share,
        commitment: List[int]
    ) -> bool:
        """
        Verificar share usando Feldman's VSS commitment.
        
        El dealer publica commitments g^{a_i} para cada coeficiente.
        Verificador comprueba: g^{share.y} == prod(commitment[j]^{x^j})
        """
        # Simplificado: en producción usar grupo de curva elíptica
        pass
    
    def add_shares(
        self,
        shares_a: List[Share],
        shares_b: List[Share]
    ) -> List[Share]:
        """
        Sumar dos secretos compartidos (localmente por cada parte).
        
        Si A = {(i, a_i)} y B = {(i, b_i)}, entonces
        A + B = {(i, a_i + b_i)} representa secret_a + secret_b
        """
        assert len(shares_a) == len(shares_b)
        
        result = []
        for sa, sb in zip(shares_a, shares_b):
            assert sa.x == sb.x, "Share indices must match"
            result.append(Share(x=sa.x, y=(sa.y + sb.y) % self.prime))
        
        return result
    
    def multiply_share_by_constant(
        self,
        shares: List[Share],
        constant: int
    ) -> List[Share]:
        """
        Multiplicar secreto compartido por constante (local).
        """
        return [Share(x=s.x, y=(s.y * constant) % self.prime) for s in shares]


class VerifiableSecretSharing(ShamirSecretSharing):
    """
    Feldman's Verifiable Secret Sharing.
    Permite verificar que los shares son consistentes.
    """
    
    def __init__(self, prime: int = None, generator: int = None):
        super().__init__(prime)
        # Generador del grupo multiplicativo Z_p^*
        self.g = generator or 2
    
    def split_with_commitments(
        self,
        secret: int,
        threshold: int,
        num_shares: int
    ) -> Tuple[List[Share], List[int]]:
        """
        Dividir secreto y generar commitments para verificación.
        
        Returns:
            (shares, commitments) donde commitments[i] = g^{a_i}
        """
        # Generar shares normalmente
        coefficients = [secret]
        for _ in range(threshold - 1):
            coefficients.append(secrets.randbelow(self.prime))
        
        shares = []
        for x in range(1, num_shares + 1):
            y = self._evaluate_polynomial(coefficients, x)
            shares.append(Share(x=x, y=y))
        
        # Generar commitments: C_j = g^{a_j} mod p
        commitments = [pow(self.g, coef, self.prime) for coef in coefficients]
        
        return shares, commitments
    
    def verify_share_with_commitment(
        self,
        share: Share,
        commitments: List[int]
    ) -> bool:
        """
        Verificar que un share es consistente con los commitments.
        
        Verifica: g^{y_i} == prod_{j=0}^{k-1} C_j^{i^j}
        """
        # LHS: g^{share.y}
        lhs = pow(self.g, share.y, self.prime)
        
        # RHS: prod(C_j^{x^j})
        rhs = 1
        x_power = 1
        for commitment in commitments:
            rhs = (rhs * pow(commitment, x_power, self.prime)) % self.prime
            x_power = (x_power * share.x) % self.prime
        
        return lhs == rhs


def secret_sharing_demo():
    print("=== Shamir's Secret Sharing Demo ===\n")
    
    sss = ShamirSecretSharing(prime=65537)  # Primo pequeño para demo
    
    # Configuración
    secret = 12345
    threshold = 3  # k = 3
    num_shares = 5  # n = 5
    
    print(f"Secret: {secret}")
    print(f"Threshold (k): {threshold}")
    print(f"Total shares (n): {num_shares}")
    
    # Dividir secreto
    print("\n--- Splitting Secret ---")
    shares = sss.split(secret, threshold, num_shares)
    
    for share in shares:
        print(f"  Share {share.x}: ({share.x}, {share.y})")
    
    # Reconstruir con exactamente k shares
    print("\n--- Reconstructing with k=3 shares ---")
    subset = shares[:3]
    reconstructed = sss.reconstruct(subset)
    print(f"Using shares: {[s.x for s in subset]}")
    print(f"Reconstructed: {reconstructed}")
    print(f"Correct: {reconstructed == secret}")
    
    # Reconstruir con más de k shares
    print("\n--- Reconstructing with k=4 shares ---")
    subset = shares[:4]
    reconstructed = sss.reconstruct(subset)
    print(f"Using shares: {[s.x for s in subset]}")
    print(f"Reconstructed: {reconstructed}")
    print(f"Correct: {reconstructed == secret}")
    
    # Reconstruir con menos de k shares falla
    print("\n--- Reconstructing with k=2 shares (insufficient) ---")
    subset = shares[:2]
    wrong_result = sss.reconstruct(subset)
    print(f"Using shares: {[s.x for s in subset]}")
    print(f"Result: {wrong_result}")
    print(f"Correct: {wrong_result == secret}")
    print("(With insufficient shares, result is random)")
    
    # Operaciones sobre shares
    print("\n--- Operations on Shares ---")
    secret_a = 100
    secret_b = 200
    
    shares_a = sss.split(secret_a, threshold, num_shares)
    shares_b = sss.split(secret_b, threshold, num_shares)
    
    # Suma
    shares_sum = sss.add_shares(shares_a, shares_b)
    result_sum = sss.reconstruct(shares_sum[:3])
    print(f"Secret A: {secret_a}, Secret B: {secret_b}")
    print(f"Sum of shares reconstructs to: {result_sum}")
    print(f"Expected: {secret_a + secret_b}")
    
    # Verifiable Secret Sharing
    print("\n--- Verifiable Secret Sharing ---")
    vss = VerifiableSecretSharing(prime=65537)
    
    shares, commitments = vss.split_with_commitments(secret, threshold, num_shares)
    print(f"Commitments: {commitments[:3]}...")
    
    # Verificar cada share
    for share in shares[:2]:
        valid = vss.verify_share_with_commitment(share, commitments)
        print(f"Share {share.x} valid: {valid}")

secret_sharing_demo()
\end{lstlisting}

%===========================================
\section{Oblivious Transfer}
\label{sec:oblivious-transfer}
%===========================================

\subsection{Definición y Variantes}

\begin{definition}[Oblivious Transfer (OT)]
Protocolo criptográfico de dos partes donde:
\begin{itemize}
    \item \textbf{Sender} tiene varios mensajes $(m_0, m_1, ..., m_{n-1})$
    \item \textbf{Receiver} tiene un índice de selección $b$
    \item Receiver obtiene $m_b$ y nada más
    \item Sender no aprende qué índice eligió Receiver
\end{itemize}

\textbf{Variantes:}
\begin{itemize}
    \item \textbf{1-out-of-2 OT:} Sender tiene 2 mensajes
    \item \textbf{1-out-of-n OT:} Sender tiene n mensajes
    \item \textbf{k-out-of-n OT:} Receiver obtiene k mensajes
\end{itemize}
\end{definition}

\begin{figure}[htbp]
\centering
\begin{tikzpicture}[
    party/.style={rectangle, draw, minimum width=2.5cm, minimum height=2cm, align=center},
    msg/.style={rectangle, draw, minimum width=0.8cm, minimum height=0.5cm, fill=yellow!30}
]
    % Sender
    \node[party, fill=blue!20] (sender) at (0, 0) {Sender};
    \node[msg] at (-0.5, 0.3) {$m_0$};
    \node[msg] at (0.5, 0.3) {$m_1$};
    
    % Receiver
    \node[party, fill=green!20] (receiver) at (6, 0) {Receiver\\$b = 1$};
    
    % Protocolo
    \draw[->, thick] (sender) -- (receiver) node[midway, above] {OT Protocol};
    
    % Output
    \node[msg, fill=green!50] at (7, 0) {$m_1$};
    
    % Propiedades
    \node at (3, -2) [font=\small, align=center] {
        Receiver obtiene $m_b$ (solo $m_1$)\\
        Sender no sabe que $b = 1$
    };
\end{tikzpicture}
\caption{1-out-of-2 Oblivious Transfer}
\label{fig:ot-basic}
\end{figure}

\begin{mathnote}[OT basado en RSA (Simplificado)]
\textbf{Setup:}
\begin{itemize}
    \item Sender genera par RSA: $(n, e)$ público, $d$ secreto
    \item Sender tiene mensajes $m_0, m_1$
\end{itemize}

\textbf{Protocolo:}
\begin{enumerate}
    \item Sender envía $x_0, x_1$ aleatorios
    \item Receiver elige $b \in \{0, 1\}$, genera $k$ aleatorio
    \item Receiver envía $v = (x_b + k^e) \mod n$
    \item Sender computa:
    \begin{itemize}
        \item $k_0 = (v - x_0)^d \mod n$
        \item $k_1 = (v - x_1)^d \mod n$
    \end{itemize}
    \item Sender envía $m'_0 = m_0 + k_0$, $m'_1 = m_1 + k_1$
    \item Receiver computa $m_b = m'_b - k$ (solo puede computar uno)
\end{enumerate}

\textbf{Seguridad:}
\begin{itemize}
    \item Sender no sabe $b$ (ambos $k_i$ parecen aleatorios)
    \item Receiver solo puede descifrar $m_b$ (no conoce el otro $k$)
\end{itemize}
\end{mathnote}

\begin{lstlisting}[language=Python, caption={Implementación de Oblivious Transfer}]
from dataclasses import dataclass
from typing import Tuple, List
import secrets
import hashlib

@dataclass
class OTSenderState:
    """Estado del sender durante OT."""
    messages: List[bytes]
    random_values: List[int]
    modulus: int
    public_exponent: int
    private_exponent: int

@dataclass
class OTReceiverState:
    """Estado del receiver durante OT."""
    choice: int
    random_k: int
    modulus: int

class ObliviousTransfer:
    """
    Implementación de 1-out-of-2 Oblivious Transfer.
    """
    
    def __init__(self, key_bits: int = 256):
        self.key_bits = key_bits
    
    def _generate_rsa_params(self) -> Tuple[int, int, int]:
        """
        Generar parámetros RSA simplificados.
        En producción usar librería criptográfica real.
        """
        # Primos pequeños para demo (NO seguros)
        p = 65537
        q = 65521
        n = p * q
        phi = (p - 1) * (q - 1)
        e = 17
        d = pow(e, -1, phi)
        return n, e, d
    
    # ========== SENDER SIDE ==========
    
    def sender_init(
        self,
        m0: bytes,
        m1: bytes
    ) -> Tuple[OTSenderState, Tuple[int, int]]:
        """
        Sender inicializa OT con sus dos mensajes.
        
        Returns:
            (state, (x0, x1)) - estado interno y valores públicos
        """
        n, e, d = self._generate_rsa_params()
        
        # Generar valores aleatorios
        x0 = secrets.randbelow(n)
        x1 = secrets.randbelow(n)
        
        state = OTSenderState(
            messages=[m0, m1],
            random_values=[x0, x1],
            modulus=n,
            public_exponent=e,
            private_exponent=d
        )
        
        return state, (x0, x1)
    
    def sender_transfer(
        self,
        state: OTSenderState,
        v: int
    ) -> Tuple[bytes, bytes]:
        """
        Sender responde al valor v del receiver.
        
        Returns:
            (e0, e1) - mensajes encriptados
        """
        n = state.modulus
        d = state.private_exponent
        x0, x1 = state.random_values
        
        # Computar ambas posibles claves
        k0 = pow((v - x0) % n, d, n)
        k1 = pow((v - x1) % n, d, n)
        
        # Encriptar mensajes
        m0, m1 = state.messages
        
        # Usar hash de k como clave de encriptación
        key0 = hashlib.sha256(k0.to_bytes(32, 'big')).digest()
        key1 = hashlib.sha256(k1.to_bytes(32, 'big')).digest()
        
        # XOR para simplicidad (en producción: AES-GCM)
        e0 = bytes(a ^ b for a, b in zip(m0.ljust(32, b'\x00'), key0))
        e1 = bytes(a ^ b for a, b in zip(m1.ljust(32, b'\x00'), key1))
        
        return e0, e1
    
    # ========== RECEIVER SIDE ==========
    
    def receiver_choose(
        self,
        choice: int,
        x0: int,
        x1: int,
        modulus: int,
        public_exp: int
    ) -> Tuple[OTReceiverState, int]:
        """
        Receiver elige cuál mensaje quiere.
        
        Args:
            choice: 0 o 1
            x0, x1: Valores públicos del sender
            modulus, public_exp: Parámetros RSA
        
        Returns:
            (state, v) - estado interno y valor para sender
        """
        n = modulus
        e = public_exp
        
        # Generar k aleatorio
        k = secrets.randbelow(n)
        
        # v = x_choice + k^e mod n
        x_choice = x0 if choice == 0 else x1
        v = (x_choice + pow(k, e, n)) % n
        
        state = OTReceiverState(
            choice=choice,
            random_k=k,
            modulus=n
        )
        
        return state, v
    
    def receiver_retrieve(
        self,
        state: OTReceiverState,
        e0: bytes,
        e1: bytes
    ) -> bytes:
        """
        Receiver obtiene el mensaje elegido.
        
        Returns:
            El mensaje m_choice
        """
        k = state.random_k
        choice = state.choice
        
        # Derivar clave
        key = hashlib.sha256(k.to_bytes(32, 'big')).digest()
        
        # Descifrar el mensaje elegido
        encrypted = e0 if choice == 0 else e1
        message = bytes(a ^ b for a, b in zip(encrypted, key))
        
        # Remover padding
        return message.rstrip(b'\x00')


class OTExtension:
    """
    OT Extension: generar muchos OTs desde pocos OTs base.
    Técnica clave para eficiencia en MPC.
    """
    
    def __init__(self, base_ot: ObliviousTransfer, security_param: int = 128):
        self.base_ot = base_ot
        self.kappa = security_param
    
    def extend(
        self,
        num_ots: int,
        sender_messages: List[Tuple[bytes, bytes]],
        receiver_choices: List[int]
    ) -> List[bytes]:
        """
        Ejecutar muchos OTs eficientemente.
        
        La idea: usar κ OTs base para generar m OTs donde m >> κ.
        Basado en IKNP03 OT extension.
        """
        # Simplificado: en implementación real, matriz de bits y PRG
        results = []
        
        for (m0, m1), choice in zip(sender_messages, receiver_choices):
            # Cada OT individual
            state_s, (x0, x1) = self.base_ot.sender_init(m0, m1)
            state_r, v = self.base_ot.receiver_choose(
                choice, x0, x1,
                state_s.modulus, state_s.public_exponent
            )
            e0, e1 = self.base_ot.sender_transfer(state_s, v)
            result = self.base_ot.receiver_retrieve(state_r, e0, e1)
            results.append(result)
        
        return results


def oblivious_transfer_demo():
    print("=== Oblivious Transfer Demo ===\n")
    
    ot = ObliviousTransfer()
    
    # Mensajes del sender
    m0 = b"Secret message ZERO"
    m1 = b"Secret message ONE"
    
    print(f"Sender has messages:")
    print(f"  m0: {m0}")
    print(f"  m1: {m1}")
    
    # Receiver quiere m1
    choice = 1
    print(f"\nReceiver wants message: m{choice}")
    
    # Protocolo
    print("\n--- OT Protocol ---")
    
    # Paso 1: Sender inicializa
    state_sender, (x0, x1) = ot.sender_init(m0, m1)
    print(f"Sender sends x0, x1: {x0}, {x1}")
    
    # Paso 2: Receiver elige y envía v
    state_receiver, v = ot.receiver_choose(
        choice, x0, x1,
        state_sender.modulus,
        state_sender.public_exponent
    )
    print(f"Receiver sends v: {v}")
    print("  (Sender cannot determine choice from v)")
    
    # Paso 3: Sender encripta y envía
    e0, e1 = ot.sender_transfer(state_sender, v)
    print(f"Sender sends e0, e1 (encrypted)")
    
    # Paso 4: Receiver descifra
    result = ot.receiver_retrieve(state_receiver, e0, e1)
    print(f"\nReceiver obtains: {result}")
    
    # Verificar propiedades
    print("\n--- Security Properties ---")
    print(f"✓ Receiver got m{choice}: {result == (m0 if choice == 0 else m1)}")
    print("✓ Sender does not know choice (v reveals nothing)")
    print("✓ Receiver cannot decrypt the other message")

oblivious_transfer_demo()
\end{lstlisting}

%===========================================
\section{Private Set Intersection (PSI)}
\label{sec:psi}
%===========================================

\subsection{Definición y Aplicaciones}

\begin{definition}[Private Set Intersection]
Protocolo que permite a dos partes encontrar elementos comunes en sus conjuntos sin revelar el resto:

\textbf{Inputs:}
\begin{itemize}
    \item Alice tiene conjunto $A = \{a_1, ..., a_n\}$
    \item Bob tiene conjunto $B = \{b_1, ..., b_m\}$
\end{itemize}

\textbf{Output:}
\begin{itemize}
    \item Intersección $A \cap B$
    \item Ni Alice ni Bob aprenden elementos del otro que no estén en la intersección
\end{itemize}
\end{definition}

\begin{figure}[htbp]
\centering
\begin{tikzpicture}
    % Conjunto A
    \draw[thick, fill=blue!20] (0, 0) circle (2cm);
    \node at (-1, 1.5) {Alice};
    \node at (-1.2, 0.5) {$a_1$};
    \node at (-1.5, -0.3) {$a_2$};
    \node at (-1, -1) {$a_3$};
    
    % Conjunto B
    \draw[thick, fill=green!20] (2.5, 0) circle (2cm);
    \node at (3.5, 1.5) {Bob};
    \node at (3.5, 0.5) {$b_1$};
    \node at (3.8, -0.3) {$b_2$};
    \node at (3.2, -1) {$b_3$};
    
    % Intersección
    \begin{scope}
        \clip (0, 0) circle (2cm);
        \fill[purple!40] (2.5, 0) circle (2cm);
    \end{scope}
    
    \node at (1.25, 0) {$x$};
    \node at (1.25, -0.5) {$y$};
    
    % Resultado
    \node at (1.25, -3) [draw, rectangle, fill=purple!20] {
        $A \cap B = \{x, y\}$
    };
    
    \draw[->, thick] (1.25, -1.8) -- (1.25, -2.5);
\end{tikzpicture}
\caption{Private Set Intersection: solo se revela la intersección}
\label{fig:psi}
\end{figure}

\begin{mathnote}[Aplicaciones de PSI]
\textbf{1. Contact Discovery (Messaging apps):}
\begin{itemize}
    \item Tu lista de contactos: $A$
    \item Usuarios del servicio: $B$
    \item Resultado: Quiénes de tus contactos usan la app
    \item Sin revelar tu lista completa al servidor
\end{itemize}

\textbf{2. Password Breach Detection:}
\begin{itemize}
    \item Tus contraseñas (hasheadas): $A$
    \item Base de datos de leaks: $B$
    \item Resultado: Qué contraseñas están comprometidas
    \item Sin enviar tus contraseñas al servicio
\end{itemize}

\textbf{3. Ad Conversion Matching:}
\begin{itemize}
    \item Usuarios que vieron ad: $A$
    \item Usuarios que compraron: $B$
    \item Resultado: Conversiones atribuidas
    \item Sin compartir bases de datos completas
\end{itemize}

\textbf{4. Fraud Detection:}
\begin{itemize}
    \item Clientes del banco A: $A$
    \item Clientes del banco B: $B$
    \item Resultado: Clientes en común (potencial fraude)
    \item Sin revelar todos los clientes
\end{itemize}
\end{mathnote}

\begin{lstlisting}[language=Python, caption={Implementación de Private Set Intersection}]
from dataclasses import dataclass
from typing import Set, List, Tuple
import hashlib
import secrets
from functools import reduce

class NaivePSI:
    """
    PSI usando hashing (inseguro pero ilustrativo).
    """
    
    def __init__(self, hash_function=hashlib.sha256):
        self.hash_fn = hash_function
    
    def hash_element(self, element: str, salt: bytes = b'') -> bytes:
        """Hash un elemento."""
        return self.hash_fn(salt + element.encode()).digest()
    
    def intersect(
        self,
        set_a: Set[str],
        set_b: Set[str]
    ) -> Set[str]:
        """
        Encontrar intersección usando hashes.
        NOTA: Vulnerable a ataques de diccionario.
        """
        salt = secrets.token_bytes(16)
        
        # Alice hashea su conjunto
        hashed_a = {self.hash_element(x, salt): x for x in set_a}
        
        # Bob hashea su conjunto
        hashed_b = {self.hash_element(x, salt) for x in set_b}
        
        # Intersección de hashes
        common_hashes = set(hashed_a.keys()) & hashed_b
        
        # Mapear de vuelta a elementos
        intersection = {hashed_a[h] for h in common_hashes}
        
        return intersection


class DiffieHellmanPSI:
    """
    PSI basado en Diffie-Hellman.
    Más seguro que hashing simple.
    """
    
    def __init__(self, prime: int = None, generator: int = 2):
        # Primo grande para DH
        self.p = prime or int(
            "FFFFFFFFFFFFFFFFC90FDAA22168C234C4C6628B80DC1CD1"
            "29024E088A67CC74020BBEA63B139B22514A08798E3404DD"
            "EF9519B3CD3A431B302B0A6DF25F14374FE1356D6D51C245"
            "E485B576625E7EC6F44C42E9A637ED6B0BFF5CB6F406B7ED"
            "EE386BFB5A899FA5AE9F24117C4B1FE649286651ECE65381"
            "FFFFFFFFFFFFFFFF", 16
        )
        self.g = generator
    
    def _hash_to_group(self, element: str) -> int:
        """Hash elemento a elemento del grupo."""
        h = int(hashlib.sha256(element.encode()).hexdigest(), 16)
        return pow(self.g, h, self.p)
    
    def alice_round1(self, set_a: Set[str]) -> Tuple[int, List[int]]:
        """
        Alice: Primera ronda.
        
        Returns:
            (secret_a, [H(a)^secret_a for a in A])
        """
        # Secreto de Alice
        secret_a = secrets.randbelow(self.p - 2) + 1
        
        # H(a)^secret_a para cada elemento
        encrypted_a = []
        for element in set_a:
            h = self._hash_to_group(element)
            encrypted = pow(h, secret_a, self.p)
            encrypted_a.append(encrypted)
        
        return secret_a, encrypted_a
    
    def bob_process(
        self,
        set_b: Set[str],
        alice_encrypted: List[int]
    ) -> Tuple[int, List[int], List[int]]:
        """
        Bob: Procesar y encriptar.
        
        Returns:
            (secret_b, 
             [H(a)^(secret_a * secret_b)],  # Alice's elements double-encrypted
             [H(b)^secret_b])               # Bob's elements encrypted
        """
        # Secreto de Bob
        secret_b = secrets.randbelow(self.p - 2) + 1
        
        # Double-encrypt elementos de Alice
        double_encrypted_a = [
            pow(enc, secret_b, self.p) for enc in alice_encrypted
        ]
        
        # Encrypt elementos de Bob
        encrypted_b = []
        for element in set_b:
            h = self._hash_to_group(element)
            encrypted = pow(h, secret_b, self.p)
            encrypted_b.append(encrypted)
        
        return secret_b, double_encrypted_a, encrypted_b
    
    def alice_round2(
        self,
        secret_a: int,
        bob_encrypted: List[int]
    ) -> List[int]:
        """
        Alice: Segunda ronda - double-encrypt elementos de Bob.
        
        Returns:
            [H(b)^(secret_a * secret_b)]
        """
        double_encrypted_b = [
            pow(enc, secret_a, self.p) for enc in bob_encrypted
        ]
        return double_encrypted_b
    
    def find_intersection(
        self,
        double_encrypted_a: List[int],  # H(a)^(ab)
        double_encrypted_b: List[int],  # H(b)^(ab)
        original_elements_a: List[str]
    ) -> Set[str]:
        """
        Encontrar intersección comparando double-encrypted values.
        """
        # Crear mapping para elementos de Alice
        a_map = dict(zip(double_encrypted_a, original_elements_a))
        
        # Encontrar matches
        b_set = set(double_encrypted_b)
        intersection = set()
        
        for enc_a, element in a_map.items():
            if enc_a in b_set:
                intersection.add(element)
        
        return intersection


class BloomFilterPSI:
    """
    PSI usando Bloom filters para conjuntos grandes.
    Trade-off: puede tener falsos positivos.
    """
    
    def __init__(
        self,
        expected_elements: int,
        false_positive_rate: float = 0.01
    ):
        # Calcular tamaño óptimo
        import math
        self.size = int(-expected_elements * math.log(false_positive_rate) / (math.log(2)**2))
        self.num_hashes = int(self.size / expected_elements * math.log(2))
        self.bit_array = [False] * self.size
    
    def _get_hash_values(self, element: str) -> List[int]:
        """Obtener índices de hash para un elemento."""
        indices = []
        for i in range(self.num_hashes):
            h = hashlib.sha256(f"{i}:{element}".encode()).digest()
            index = int.from_bytes(h[:4], 'big') % self.size
            indices.append(index)
        return indices
    
    def add(self, element: str):
        """Añadir elemento al filtro."""
        for index in self._get_hash_values(element):
            self.bit_array[index] = True
    
    def contains(self, element: str) -> bool:
        """Verificar si elemento posiblemente está en el filtro."""
        return all(self.bit_array[i] for i in self._get_hash_values(element))
    
    def intersect_with_set(self, other_set: Set[str]) -> Set[str]:
        """Encontrar posible intersección."""
        return {x for x in other_set if self.contains(x)}


def psi_demo():
    print("=== Private Set Intersection Demo ===\n")
    
    # Conjuntos de ejemplo
    alice_contacts = {"alice@email.com", "bob@email.com", "carol@email.com", "dave@email.com"}
    app_users = {"bob@email.com", "eve@email.com", "carol@email.com", "frank@email.com"}
    
    print("Alice's contacts:", alice_contacts)
    print("App users:", app_users)
    print("Expected intersection:", alice_contacts & app_users)
    
    # PSI con Diffie-Hellman
    print("\n--- Diffie-Hellman PSI ---")
    
    dh_psi = DiffieHellmanPSI(prime=65537)  # Primo pequeño para demo
    
    # Ronda 1: Alice
    alice_list = list(alice_contacts)
    secret_a, encrypted_a = dh_psi.alice_round1(alice_contacts)
    print(f"Alice sends {len(encrypted_a)} encrypted elements")
    
    # Bob procesa
    secret_b, double_enc_a, encrypted_b = dh_psi.bob_process(
        app_users, encrypted_a
    )
    print(f"Bob sends back double-encrypted + his {len(encrypted_b)} elements")
    
    # Ronda 2: Alice
    double_enc_b = dh_psi.alice_round2(secret_a, encrypted_b)
    
    # Encontrar intersección
    intersection = dh_psi.find_intersection(
        double_enc_a, double_enc_b, alice_list
    )
    
    print(f"\nIntersection found: {intersection}")
    print(f"Correct: {intersection == (alice_contacts & app_users)}")
    
    # PSI con Bloom Filter
    print("\n--- Bloom Filter PSI (approximate) ---")
    
    bf_psi = BloomFilterPSI(expected_elements=100)
    
    # Bob añade sus usuarios al filtro
    for user in app_users:
        bf_psi.add(user)
    
    # Alice consulta el filtro
    possible_intersection = bf_psi.intersect_with_set(alice_contacts)
    print(f"Possible intersection: {possible_intersection}")
    print("(May have false positives)")

psi_demo()
\end{lstlisting}
%===========================================
\section{Secure Enclaves}
\label{sec:secure-enclaves}
%===========================================

\subsection{Trusted Execution Environments (TEE)}

\begin{definition}[Trusted Execution Environment]
Hardware aislado que proporciona ejecución segura de código, garantizando:
\begin{enumerate}
    \item \textbf{Confidencialidad:} Datos y código protegidos incluso del OS/hypervisor
    \item \textbf{Integridad:} Código no puede ser modificado durante ejecución
    \item \textbf{Attestation:} Prueba criptográfica de qué código se ejecuta
\end{enumerate}
\end{definition}

\begin{figure}[htbp]
\centering
\begin{tikzpicture}[
    node distance=0.5cm,
    layer/.style={rectangle, draw, minimum width=8cm, minimum height=0.8cm, align=center},
    enclave/.style={rectangle, draw, rounded corners, minimum width=3cm, minimum height=1.5cm, fill=green!30, align=center}
]
    % Capas del sistema
    \node[layer, fill=red!20] (app) at (0, 3) {Aplicaciones (no confiables)};
    \node[layer, fill=orange!20] (os) at (0, 2) {Sistema Operativo (no confiable)};
    \node[layer, fill=yellow!20] (hyp) at (0, 1) {Hypervisor (no confiable)};
    \node[layer, fill=blue!20] (hw) at (0, 0) {Hardware (CPU + Enclave)};
    
    % Enclave
    \node[enclave] (enc) at (5, 2) {Enclave\\(Confiable)};
    
    % Conexión
    \draw[->, thick, dashed] (2.5, 0.4) -- (enc);
    
    % Etiquetas
    \node at (5, 0) [font=\small, align=center] {Aislado del resto\\del sistema};
    
    % Ataques bloqueados
    \node at (-5, 2) [font=\small, red] {✗ Malware};
    \node at (-5, 1) [font=\small, red] {✗ Root/Admin};
    \node at (-5, 0) [font=\small, red] {✗ DMA attacks};
\end{tikzpicture}
\caption{TEE aísla código incluso de software privilegiado}
\label{fig:tee-isolation}
\end{figure}

\subsection{Intel SGX}

\begin{definition}[Intel SGX (Software Guard Extensions)]
Extensiones de CPU que crean ``enclaves'' - regiones de memoria aisladas:

\textbf{Características:}
\begin{itemize}
    \item Memoria encriptada en RAM (MEE)
    \item Código y datos protegidos del OS
    \item Attestation remota vía Intel
    \item Sealing para persistencia
\end{itemize}

\textbf{Limitaciones:}
\begin{itemize}
    \item Descontinuado en CPUs consumer (desde 11th gen)
    \item Disponible en Xeon (servidor)
    \item Límite de memoria EPC (~256MB típico)
\end{itemize}
\end{definition}

\begin{lstlisting}[language=Python, caption={Simulación conceptual de operaciones SGX}]
from dataclasses import dataclass, field
from typing import Dict, Optional, Tuple, List
import hashlib
import secrets
from enum import Enum

class EnclaveState(Enum):
    CREATED = "created"
    INITIALIZED = "initialized"
    RUNNING = "running"
    DESTROYED = "destroyed"

@dataclass
class EnclaveIdentity:
    """Identidad de un enclave para attestation."""
    mrenclave: bytes  # Hash del código del enclave
    mrsigner: bytes   # Hash de la clave del firmante
    product_id: int
    security_version: int

@dataclass
class SGXEnclave:
    """
    Simulación de un enclave SGX.
    """
    enclave_id: str
    identity: EnclaveIdentity
    state: EnclaveState = EnclaveState.CREATED
    
    # Memoria protegida (simulada)
    _protected_memory: Dict[str, bytes] = field(default_factory=dict)
    
    # Clave de sealing (derivada de CPU + enclave identity)
    _sealing_key: bytes = field(default_factory=lambda: secrets.token_bytes(32))
    
    def initialize(self, code_hash: bytes):
        """Inicializar enclave con código."""
        self.identity.mrenclave = code_hash
        self.state = EnclaveState.INITIALIZED
    
    def ecall(self, function_name: str, params: bytes) -> bytes:
        """
        Llamada al enclave (ECALL).
        Ejecutar función dentro del enclave.
        """
        if self.state != EnclaveState.RUNNING:
            raise RuntimeError("Enclave not running")
        
        # Simular ejecución dentro del enclave
        # En SGX real, esto ejecuta código en modo protegido
        
        if function_name == "store_secret":
            key = hashlib.sha256(params[:32]).hexdigest()[:16]
            self._protected_memory[key] = params[32:]
            return b"OK"
        
        elif function_name == "compute_on_secret":
            key = hashlib.sha256(params[:32]).hexdigest()[:16]
            if key in self._protected_memory:
                secret = self._protected_memory[key]
                # Computar sobre el secreto
                result = hashlib.sha256(secret).digest()
                return result
            return b"NOT_FOUND"
        
        elif function_name == "get_report":
            return self._generate_report(params)
        
        return b"UNKNOWN_FUNCTION"
    
    def _generate_report(self, target_info: bytes) -> bytes:
        """
        Generar reporte para attestation.
        """
        report_data = (
            self.identity.mrenclave +
            self.identity.mrsigner +
            self.product_id.to_bytes(2, 'big') +
            self.security_version.to_bytes(2, 'big') +
            target_info
        )
        
        # En SGX real, esto es firmado por la CPU
        report_mac = hashlib.sha256(report_data + self._sealing_key).digest()
        
        return report_data + report_mac
    
    def seal_data(self, plaintext: bytes) -> bytes:
        """
        Seal: Encriptar datos vinculados al enclave.
        Solo este enclave (o uno idéntico) puede unseal.
        """
        nonce = secrets.token_bytes(12)
        
        # Derivar clave de sealing
        # En SGX real, derivada de CPU key + enclave identity
        seal_key = hashlib.sha256(
            self._sealing_key + self.identity.mrenclave
        ).digest()
        
        # Encriptar (simplificado)
        ciphertext = bytes(p ^ k for p, k in zip(
            plaintext.ljust(len(plaintext) + 16, b'\x00'),
            (seal_key * 10)[:len(plaintext) + 16]
        ))
        
        return nonce + ciphertext
    
    def unseal_data(self, sealed: bytes) -> bytes:
        """Unseal: Descifrar datos sellados."""
        nonce = sealed[:12]
        ciphertext = sealed[12:]
        
        seal_key = hashlib.sha256(
            self._sealing_key + self.identity.mrenclave
        ).digest()
        
        plaintext = bytes(c ^ k for c, k in zip(
            ciphertext,
            (seal_key * 10)[:len(ciphertext)]
        ))
        
        return plaintext.rstrip(b'\x00')


class RemoteAttestation:
    """
    Simulación de attestation remota SGX.
    """
    
    def __init__(self):
        # "Intel Attestation Service" simulado
        self.ias_key = secrets.token_bytes(32)
        self.trusted_enclaves: Dict[bytes, bool] = {}
    
    def register_trusted_enclave(self, mrenclave: bytes):
        """Registrar un código de enclave como confiable."""
        self.trusted_enclaves[mrenclave] = True
    
    def verify_quote(
        self,
        quote: bytes,
        expected_mrenclave: bytes
    ) -> Tuple[bool, str]:
        """
        Verificar quote de attestation.
        
        Returns:
            (is_valid, message)
        """
        # Extraer datos del quote
        if len(quote) < 64:
            return False, "Invalid quote format"
        
        mrenclave = quote[:32]
        mrsigner = quote[32:64]
        
        # Verificar que el código es el esperado
        if mrenclave != expected_mrenclave:
            return False, "MRENCLAVE mismatch"
        
        # Verificar que está en lista de confianza
        if mrenclave not in self.trusted_enclaves:
            return False, "Enclave not in trusted list"
        
        # Verificar firma (simplificado)
        # En producción: verificar firma EPID/ECDSA de Intel
        
        return True, "Attestation successful"


def sgx_demo():
    print("=== Intel SGX Simulation Demo ===\n")
    
    # Crear enclave
    identity = EnclaveIdentity(
        mrenclave=b'\x00' * 32,
        mrsigner=hashlib.sha256(b"developer_key").digest(),
        product_id=1,
        security_version=1
    )
    
    enclave = SGXEnclave(
        enclave_id="enclave_001",
        identity=identity
    )
    
    # Inicializar con código
    code_hash = hashlib.sha256(b"my_enclave_code_v1").digest()
    enclave.initialize(code_hash)
    enclave.identity.mrenclave = code_hash
    enclave.state = EnclaveState.RUNNING
    
    print(f"Enclave created: {enclave.enclave_id}")
    print(f"MRENCLAVE: {enclave.identity.mrenclave.hex()[:32]}...")
    
    # Almacenar secreto en el enclave
    print("\n--- Storing Secret in Enclave ---")
    secret_key = b"my_secret_key_identifier_here!!!"
    secret_value = b"This is my super secret data that only the enclave can see"
    
    result = enclave.ecall("store_secret", secret_key + secret_value)
    print(f"Store result: {result}")
    
    # Computar sobre el secreto
    print("\n--- Computing on Secret ---")
    hash_result = enclave.ecall("compute_on_secret", secret_key)
    print(f"Hash of secret (computed inside enclave): {hash_result.hex()[:32]}...")
    
    # Attestation
    print("\n--- Remote Attestation ---")
    ias = RemoteAttestation()
    ias.register_trusted_enclave(code_hash)
    
    target_info = b"verifier_identity"
    report = enclave.ecall("get_report", target_info)
    
    is_valid, message = ias.verify_quote(report, code_hash)
    print(f"Attestation valid: {is_valid}")
    print(f"Message: {message}")
    
    # Sealing
    print("\n--- Data Sealing ---")
    sensitive_data = b"Persistent secret that survives enclave restart"
    
    sealed = enclave.seal_data(sensitive_data)
    print(f"Sealed data: {sealed.hex()[:40]}...")
    
    unsealed = enclave.unseal_data(sealed)
    print(f"Unsealed: {unsealed}")
    print(f"Match: {unsealed == sensitive_data}")

sgx_demo()
\end{lstlisting}

\subsection{Comparación de TEEs}

\begin{table}[htbp]
\centering
\footnotesize
\begin{tabular}{lcccc}
\toprule
\textbf{Característica} & \textbf{Intel SGX} & \textbf{ARM TrustZone} & \textbf{AMD SEV} & \textbf{AWS Nitro} \\
\midrule
Granularidad & Proceso & Sistema & VM & VM \\
Aislamiento & Enclave & Mundo seguro & VM encriptada & Enclave dedicado \\
Attestation & Remota & Local & Remota & Remota \\
Memoria EPC & ~256MB & Configurable & Toda RAM & Variable \\
Side-channels & Vulnerable & Menos expuesto & Menos expuesto & Diseño diferente \\
Disponibilidad & Server Xeon & Móviles/IoT & EPYC & AWS Cloud \\
\bottomrule
\end{tabular}
\caption{Comparación de Trusted Execution Environments}
\end{table}

%===========================================
\section{Threshold Signatures}
\label{sec:threshold-signatures}
%===========================================

\subsection{Fundamentos}

\begin{definition}[Threshold Signature]
Esquema de firma donde $k$ de $n$ partes deben cooperar para firmar:

\textbf{Propiedades:}
\begin{itemize}
    \item La clave privada \textbf{nunca existe completa}
    \item Cada parte tiene un share de la clave
    \item Cualquier subconjunto de $k$ partes puede firmar
    \item Menos de $k$ partes no pueden firmar ni recuperar la clave
    \item La firma resultante es indistinguible de firma normal
\end{itemize}
\end{definition}

\begin{figure}[htbp]
\centering
\begin{tikzpicture}[
    party/.style={circle, draw, minimum size=1.2cm, fill=blue!20},
    key/.style={rectangle, draw, minimum width=0.8cm, minimum height=0.4cm, fill=yellow!30}
]
    % Partes
    \node[party] (p1) at (0, 2) {$P_1$};
    \node[party] (p2) at (2, 2) {$P_2$};
    \node[party] (p3) at (4, 2) {$P_3$};
    \node[party] (p4) at (6, 2) {$P_4$};
    \node[party] (p5) at (8, 2) {$P_5$};
    
    % Key shares
    \node[key] at (0, 1) {$sk_1$};
    \node[key] at (2, 1) {$sk_2$};
    \node[key, fill=green!30] at (4, 1) {$sk_3$};
    \node[key, fill=green!30] at (6, 1) {$sk_4$};
    \node[key, fill=green!30] at (8, 1) {$sk_5$};
    
    % Firma
    \node[rectangle, draw, fill=purple!20, minimum width=2cm] (sig) at (6, -1) {Firma $\sigma$};
    
    % Conexiones (3-of-5)
    \draw[->, thick, green!60!black] (4, 0.7) -- (6, -0.5);
    \draw[->, thick, green!60!black] (6, 0.7) -- (6, -0.5);
    \draw[->, thick, green!60!black] (8, 0.7) -- (6, -0.5);
    
    % Etiqueta
    \node at (4, -2) [font=\small] {3-of-5: Cualquier 3 partes pueden firmar};
\end{tikzpicture}
\caption{Threshold signature: 3 de 5 partes cooperan para firmar}
\label{fig:threshold-sig}
\end{figure}

\subsection{Threshold ECDSA}

\begin{mathnote}[Threshold ECDSA]
\textbf{Desafío:} ECDSA no es ``threshold-friendly'' como Schnorr/BLS.

La firma ECDSA es:
\[
s = k^{-1}(z + r \cdot sk) \mod n
\]

donde $k$ es nonce aleatorio, $r$ es coordenada x del punto $k \cdot G$.

\textbf{Problema:} Requiere multiplicación e inversión de secretos compartidos.

\textbf{Soluciones:}
\begin{itemize}
    \item \textbf{GG18/GG20:} Protocolo de Gennaro-Goldfeder
    \item \textbf{CGGMP:} Canetti, Gennaro, Goldfeder, Makriyannis, Peled
    \item Usan MPC para multiplicación e inversión
\end{itemize}
\end{mathnote}

\subsection{FROST (Flexible Round-Optimized Schnorr Threshold)}

\begin{definition}[FROST]
Protocolo de threshold Schnorr signatures optimizado:

\textbf{Ventajas sobre threshold ECDSA:}
\begin{itemize}
    \item Solo 2 rondas de comunicación
    \item Compatible con agregación de firmas
    \item Más eficiente computacionalmente
    \item Seguro contra adversario malicioso
\end{itemize}

\textbf{Firma Schnorr:}
\[
\sigma = (R, s) \text{ donde } s = k + c \cdot sk
\]
con $c = H(R || pk || m)$
\end{definition}

\begin{lstlisting}[language=Python, caption={Implementación de Threshold Schnorr (FROST simplificado)}]
from dataclasses import dataclass
from typing import List, Tuple, Dict
import secrets
import hashlib

# Parámetros de curva secp256k1 (simplificados)
P = 0xFFFFFFFFFFFFFFFFFFFFFFFFFFFFFFFFFFFFFFFFFFFFFFFFFFFFFFFEFFFFFC2F
N = 0xFFFFFFFFFFFFFFFFFFFFFFFFFFFFFFFEBAAEDCE6AF48A03BBFD25E8CD0364141
Gx = 0x79BE667EF9DCBBAC55A06295CE870B07029BFCDB2DCE28D959F2815B16F81798
Gy = 0x483ADA7726A3C4655DA4FBFC0E1108A8FD17B448A68554199C47D08FFB10D4B8

@dataclass
class Point:
    x: int
    y: int
    
    def __add__(self, other):
        """Suma de puntos en curva elíptica."""
        if self.x == 0 and self.y == 0:
            return other
        if other.x == 0 and other.y == 0:
            return self
        
        if self.x == other.x and self.y != other.y:
            return Point(0, 0)  # Punto en infinito
        
        if self.x == other.x:
            # Point doubling
            lam = (3 * self.x * self.x * pow(2 * self.y, -1, P)) % P
        else:
            lam = ((other.y - self.y) * pow(other.x - self.x, -1, P)) % P
        
        x3 = (lam * lam - self.x - other.x) % P
        y3 = (lam * (self.x - x3) - self.y) % P
        
        return Point(x3, y3)
    
    def __mul__(self, scalar: int):
        """Multiplicación escalar."""
        result = Point(0, 0)
        addend = self
        
        while scalar:
            if scalar & 1:
                result = result + addend
            addend = addend + addend
            scalar >>= 1
        
        return result

G = Point(Gx, Gy)

@dataclass
class KeyShare:
    """Share de clave privada."""
    index: int
    value: int

@dataclass
class NonceShare:
    """Share de nonce para firma."""
    index: int
    d: int  # Hiding nonce
    e: int  # Binding nonce
    D: Point  # d*G
    E: Point  # e*G

class ThresholdSchnorr:
    """
    Implementación simplificada de Threshold Schnorr (estilo FROST).
    """
    
    def __init__(self, threshold: int, num_parties: int):
        self.t = threshold
        self.n = num_parties
        self.key_shares: Dict[int, KeyShare] = {}
        self.public_key: Point = None
    
    def _lagrange_coefficient(
        self,
        participant_index: int,
        participant_indices: List[int]
    ) -> int:
        """
        Calcular coeficiente de Lagrange para interpolación.
        """
        numerator = 1
        denominator = 1
        
        for j in participant_indices:
            if j != participant_index:
                numerator = (numerator * (-j)) % N
                denominator = (denominator * (participant_index - j)) % N
        
        return (numerator * pow(denominator, -1, N)) % N
    
    def distributed_keygen(self) -> Point:
        """
        Distributed Key Generation (DKG).
        Cada parte genera polynomial y comparte evaluaciones.
        
        Returns:
            Public key (punto en curva)
        """
        # Fase 1: Cada parte genera polynomial secreto
        polynomials = []
        commitments = []
        
        for i in range(1, self.n + 1):
            # Polynomial de grado t-1
            coeffs = [secrets.randbelow(N) for _ in range(self.t)]
            polynomials.append(coeffs)
            
            # Commitment: C_j = coeffs[j] * G
            party_commitments = [G * c for c in coeffs]
            commitments.append(party_commitments)
        
        # Fase 2: Cada parte calcula su share sumando evaluaciones
        for i in range(1, self.n + 1):
            share_value = 0
            for poly in polynomials:
                # Evaluar polynomial en punto i
                eval_i = sum(poly[j] * pow(i, j, N) for j in range(len(poly))) % N
                share_value = (share_value + eval_i) % N
            
            self.key_shares[i] = KeyShare(index=i, value=share_value)
        
        # Public key: suma de los primeros coeficientes * G
        pk_x = 0
        pk_y = 0
        for party_commits in commitments:
            pk = party_commits[0]
            if pk_x == 0:
                pk_x, pk_y = pk.x, pk.y
            else:
                combined = Point(pk_x, pk_y) + pk
                pk_x, pk_y = combined.x, combined.y
        
        self.public_key = Point(pk_x, pk_y)
        return self.public_key
    
    def generate_nonce_share(self, party_index: int) -> NonceShare:
        """
        Generar share de nonce para una firma.
        Cada parte genera (d_i, e_i) y publica (D_i, E_i).
        """
        d = secrets.randbelow(N)
        e = secrets.randbelow(N)
        D = G * d
        E = G * e
        
        return NonceShare(index=party_index, d=d, e=e, D=D, E=E)
    
    def aggregate_nonces(
        self,
        nonce_shares: List[NonceShare],
        message: bytes,
        participant_indices: List[int]
    ) -> Tuple[Point, int]:
        """
        Agregar nonces para formar R del grupo.
        
        Returns:
            (R, rho) donde rho es el binding factor
        """
        # Calcular binding factor rho = H(participant_indices, message, {D_i, E_i})
        binding_data = (
            str(participant_indices).encode() +
            message +
            b''.join(
                ns.D.x.to_bytes(32, 'big') + ns.E.x.to_bytes(32, 'big')
                for ns in nonce_shares
            )
        )
        rho = int(hashlib.sha256(binding_data).hexdigest(), 16) % N
        
        # R = sum(D_i + rho * E_i)
        R = Point(0, 0)
        for ns in nonce_shares:
            Ri = ns.D + (ns.E * rho)
            R = R + Ri
        
        return R, rho
    
    def compute_partial_signature(
        self,
        party_index: int,
        nonce_share: NonceShare,
        message: bytes,
        R: Point,
        rho: int,
        participant_indices: List[int]
    ) -> int:
        """
        Calcular firma parcial de una parte.
        """
        # Challenge c = H(R || pk || message)
        c_data = (
            R.x.to_bytes(32, 'big') +
            self.public_key.x.to_bytes(32, 'big') +
            message
        )
        c = int(hashlib.sha256(c_data).hexdigest(), 16) % N
        
        # Coeficiente de Lagrange
        lambda_i = self._lagrange_coefficient(party_index, participant_indices)
        
        # z_i = d_i + rho * e_i + lambda_i * sk_i * c
        d_i = nonce_share.d
        e_i = nonce_share.e
        sk_i = self.key_shares[party_index].value
        
        z_i = (d_i + rho * e_i + lambda_i * sk_i * c) % N
        
        return z_i
    
    def aggregate_signature(
        self,
        R: Point,
        partial_signatures: List[int]
    ) -> Tuple[Point, int]:
        """
        Agregar firmas parciales en firma final.
        
        Returns:
            (R, s) - firma Schnorr
        """
        s = sum(partial_signatures) % N
        return (R, s)
    
    def verify(
        self,
        message: bytes,
        signature: Tuple[Point, int]
    ) -> bool:
        """
        Verificar firma Schnorr.
        """
        R, s = signature
        
        # c = H(R || pk || message)
        c_data = (
            R.x.to_bytes(32, 'big') +
            self.public_key.x.to_bytes(32, 'big') +
            message
        )
        c = int(hashlib.sha256(c_data).hexdigest(), 16) % N
        
        # Verificar: s*G == R + c*pk
        lhs = G * s
        rhs = R + (self.public_key * c)
        
        return lhs.x == rhs.x and lhs.y == rhs.y


def threshold_schnorr_demo():
    print("=== Threshold Schnorr (FROST-like) Demo ===\n")
    
    threshold = 3
    num_parties = 5
    
    ts = ThresholdSchnorr(threshold, num_parties)
    
    # Distributed Key Generation
    print("--- Distributed Key Generation ---")
    pk = ts.distributed_keygen()
    print(f"Public key: ({pk.x % (10**10)}..., {pk.y % (10**10)}...)")
    print(f"Key shares distributed to {num_parties} parties")
    
    # Signing
    message = b"Transfer 100 BTC to Alice"
    print(f"\n--- Signing Message ---")
    print(f"Message: {message}")
    
    # Participantes que firmarán (cualquier 3 de 5)
    signers = [1, 3, 5]
    print(f"Signing parties: {signers}")
    
    # Ronda 1: Generar nonce shares
    print("\n--- Round 1: Generate Nonce Shares ---")
    nonce_shares = {}
    for i in signers:
        nonce_shares[i] = ts.generate_nonce_share(i)
    print(f"Each party generated (d_i, e_i) and published (D_i, E_i)")
    
    # Agregar nonces
    nonces_list = [nonce_shares[i] for i in signers]
    R, rho = ts.aggregate_nonces(nonces_list, message, signers)
    print(f"Aggregated R: ({R.x % (10**10)}...)")
    
    # Ronda 2: Firmas parciales
    print("\n--- Round 2: Partial Signatures ---")
    partial_sigs = []
    for i in signers:
        z_i = ts.compute_partial_signature(
            i, nonce_shares[i], message, R, rho, signers
        )
        partial_sigs.append(z_i)
        print(f"Party {i} generated z_{i} = {z_i % (10**10)}...")
    
    # Agregar firma
    print("\n--- Aggregate Signature ---")
    signature = ts.aggregate_signature(R, partial_sigs)
    R_final, s = signature
    print(f"Final signature: (R, s)")
    
    # Verificar
    print("\n--- Verification ---")
    valid = ts.verify(message, signature)
    print(f"Signature valid: {valid}")

threshold_schnorr_demo()
\end{lstlisting}

%===========================================
\section{Distributed Key Generation (DKG)}
\label{sec:dkg}
%===========================================

\subsection{Generación sin Trusted Dealer}

\begin{definition}[Distributed Key Generation]
Protocolo para generar clave compartida donde:
\begin{itemize}
    \item No existe ``dealer'' que conoce la clave completa
    \item Cada parte contribuye aleatoriedad
    \item Nadie puede predecir o influenciar la clave final
    \item Output: $(pk, \{sk_1, ..., sk_n\})$ donde ninguno conoce $sk$
\end{itemize}
\end{definition}

\begin{mathnote}[Protocolo DKG (Feldman VSS)]
\textbf{Fase 1 - Sharing:}
\begin{enumerate}
    \item Cada parte $P_i$ elige secreto $a_{i,0}$ y polynomial:
    \[
    f_i(x) = a_{i,0} + a_{i,1}x + ... + a_{i,t-1}x^{t-1}
    \]
    \item $P_i$ publica commitments: $C_{i,j} = g^{a_{i,j}}$
    \item $P_i$ envía $s_{i,j} = f_i(j)$ a cada $P_j$
\end{enumerate}

\textbf{Fase 2 - Verificación:}
\begin{enumerate}
    \item $P_j$ verifica share recibido de $P_i$:
    \[
    g^{s_{i,j}} \stackrel{?}{=} \prod_{k=0}^{t-1} C_{i,k}^{j^k}
    \]
    \item Si verificación falla, $P_j$ protesta públicamente
\end{enumerate}

\textbf{Fase 3 - Cálculo de clave:}
\begin{enumerate}
    \item Clave privada (compartida): $sk = \sum_i a_{i,0}$
    \item Share de $P_j$: $sk_j = \sum_i s_{i,j}$
    \item Clave pública: $pk = \prod_i C_{i,0} = g^{sk}$
\end{enumerate}
\end{mathnote}

\begin{figure}[htbp]
\centering
\begin{tikzpicture}[
    party/.style={circle, draw, minimum size=1cm, fill=blue!20},
    share/.style={rectangle, draw, minimum width=0.5cm, minimum height=0.3cm, fill=yellow!30, font=\tiny}
]
    % Partes
    \node[party] (p1) at (0, 0) {$P_1$};
    \node[party] (p2) at (3, 0) {$P_2$};
    \node[party] (p3) at (6, 0) {$P_3$};
    
    % Flechas de shares
    \draw[->, thick, red] (p1) to[bend left=20] node[above, font=\tiny] {$f_1(2)$} (p2);
    \draw[->, thick, red] (p1) to[bend left=30] node[above, font=\tiny] {$f_1(3)$} (p3);
    
    \draw[->, thick, green] (p2) to[bend left=20] node[below, font=\tiny] {$f_2(1)$} (p1);
    \draw[->, thick, green] (p2) to[bend left=20] node[above, font=\tiny] {$f_2(3)$} (p3);
    
    \draw[->, thick, blue] (p3) to[bend left=30] node[below, font=\tiny] {$f_3(1)$} (p1);
    \draw[->, thick, blue] (p3) to[bend left=20] node[below, font=\tiny] {$f_3(2)$} (p2);
    
    % Resultado
    \node at (3, -2) [draw, rectangle, align=center, font=\small] {
        $sk_j = f_1(j) + f_2(j) + f_3(j)$\\
        Nadie conoce $sk = a_{1,0} + a_{2,0} + a_{3,0}$
    };
\end{tikzpicture}
\caption{DKG: cada parte envía evaluaciones de su polynomial}
\label{fig:dkg}
\end{figure}

\subsection{Aplicaciones en Blockchain}

\begin{keyidea}
\textbf{Uso de DKG + Threshold Signatures en Blockchain:}

\begin{enumerate}
    \item \textbf{Bridges cross-chain:}
    \begin{itemize}
        \item Comité de validadores controla dirección bridge
        \item k-of-n para autorizar transferencias
        \item Sin single point of failure
    \end{itemize}
    
    \item \textbf{Custody institucional:}
    \begin{itemize}
        \item Claves distribuidas entre múltiples ubicaciones
        \item Requiere quorum para firmar
        \item Resistente a compromiso de una ubicación
    \end{itemize}
    
    \item \textbf{Validator sets (PoS):}
    \begin{itemize}
        \item BLS threshold signatures para consenso
        \item Eficiente agregación de firmas
        \item Usado en Ethereum 2.0
    \end{itemize}
    
    \item \textbf{Encrypted mempools:}
    \begin{itemize}
        \item Comité de keypers para threshold decryption
        \item Previene frontrunning
    \end{itemize}
\end{enumerate}
\end{keyidea}

%===========================================
\section*{Resumen del Bloque}
%===========================================

\begin{keyidea}
\textbf{Puntos clave de Criptografía Avanzada (Bloque 18):}

\textbf{1. FHE (Fully Homomorphic Encryption):}
\begin{itemize}
    \item Computar sobre datos encriptados
    \item Basado en LWE (post-quantum secure)
    \item Generaciones: BGV/BFV → TFHE/CKKS
    \item Todavía costoso (~100-1000x overhead)
\end{itemize}

\textbf{2. MPC (Multi-Party Computation):}
\begin{itemize}
    \item Garbled Circuits: 2 partes, propósito general
    \item GMW: n partes, basado en secret sharing
    \item SPDZ: Eficiente con preprocessing
    \item Base de muchos protocolos de privacidad
\end{itemize}

\textbf{3. Secret Sharing (Shamir):}
\begin{itemize}
    \item $(k,n)$-threshold: k de n para reconstruir
    \item Basado en interpolación de Lagrange
    \item Operaciones sobre shares (suma, mult)
    \item VSS: Verifiable Secret Sharing
\end{itemize}

\textbf{4. Oblivious Transfer:}
\begin{itemize}
    \item Transferencia sin revelar elección
    \item Base de Garbled Circuits
    \item OT Extension para eficiencia
\end{itemize}

\textbf{5. PSI (Private Set Intersection):}
\begin{itemize}
    \item Encontrar elementos comunes sin revelar resto
    \item Basado en DH o hashing
    \item Aplicaciones: contact discovery, fraud detection
\end{itemize}

\textbf{6. Secure Enclaves:}
\begin{itemize}
    \item SGX, TrustZone, SEV
    \item Hardware trust model
    \item Attestation remota
    \item Vulnerables a side-channels
\end{itemize}

\textbf{7. Threshold Signatures:}
\begin{itemize}
    \item k-of-n para firmar
    \item ECDSA: complejo (GG20)
    \item Schnorr/FROST: más eficiente
    \item Clave nunca existe completa
\end{itemize}

\textbf{8. DKG (Distributed Key Generation):}
\begin{itemize}
    \item Generar clave sin trusted dealer
    \item Cada parte contribuye aleatoriedad
    \item Usado en bridges, custody, validators
\end{itemize}
\end{keyidea}

\begin{bancozk}
\textbf{Criptografía Avanzada en Banco ZK:}

\begin{center}
\begin{tikzpicture}[
    component/.style={rectangle, draw, rounded corners, minimum width=3cm, minimum height=1cm, align=center, font=\small}
]
    % Componentes
    \node[component, fill=blue!20] at (0, 3) {ZK Proofs\\(Core privacy)};
    \node[component, fill=green!20] at (4, 3) {Threshold Sigs\\(Governance)};
    \node[component, fill=orange!20] at (8, 3) {MPC\\(Compliance)};
    
    \node[component, fill=purple!20] at (0, 1) {Secret Sharing\\(Key backup)};
    \node[component, fill=yellow!20] at (4, 1) {PSI\\(AML checks)};
    \node[component, fill=red!20] at (8, 1) {TEE (Optional)\\(Performance)};
\end{tikzpicture}
\end{center}

\textbf{Casos de uso:}
\begin{enumerate}
    \item \textbf{Transacciones privadas:} ZK proofs
    \item \textbf{Multisig institucional:} Threshold ECDSA/Schnorr
    \item \textbf{Compliance sin revelar datos:} MPC + PSI
    \item \textbf{Key recovery:} Shamir secret sharing
    \item \textbf{Bridge cross-chain:} DKG + threshold sigs
\end{enumerate}

\textbf{Stack recomendado:}
\begin{itemize}
    \item \textbf{Firmas:} FROST (Schnorr threshold)
    \item \textbf{Sharing:} Shamir con VSS
    \item \textbf{MPC:} SPDZ para compliance
    \item \textbf{FHE:} Experimental para casos específicos
\end{itemize}
\end{bancozk}

% Fin del Bloque 18
\cleardoublepage

% ============================================
% BLOQUE 19: ATAQUES Y DEFENSA - VERSIÓN EXPANDIDA
% Parte 1: Introducción, Deanonymization, Sybil, Eclipse
% ============================================

\chapter{Ataques y Defensa}
\label{ch:bloque19}

\begin{center}
\textit{Bloque 19 de 20}
\end{center}

\section*{Introducción}

La privacidad en blockchain es un juego constante entre atacantes y defensores. Este bloque analiza los principales vectores de ataque contra sistemas de privacidad y las defensas correspondientes.

\begin{keyidea}
\textbf{Taxonomía de Ataques a la Privacidad:}

\begin{center}
\begin{tikzpicture}[
    attack/.style={rectangle, draw, rounded corners, minimum width=2.2cm, minimum height=0.6cm, align=center, font=\small},
    category/.style={rectangle, draw, fill=gray!20, minimum width=2.4cm, minimum height=0.5cm, font=\small\bfseries}
]
    \node[category] (net) at (0, 2.5) {Network};
    \node[category] (chain) at (3, 2.5) {Blockchain};
    \node[category] (behav) at (6, 2.5) {Behavioral};
    \node[category] (meta) at (9, 2.5) {Metadata};
    
    \node[attack, fill=red!20] at (0, 1.5) {IP Correlation};
    \node[attack, fill=red!20] at (0, 0.7) {Eclipse};
    \node[attack, fill=red!20] at (0, -0.1) {Sybil};
    
    \node[attack, fill=orange!20] at (3, 1.5) {Graph Analysis};
    \node[attack, fill=orange!20] at (3, 0.7) {Taint Analysis};
    \node[attack, fill=orange!20] at (3, -0.1) {Dust};
    
    \node[attack, fill=yellow!20] at (6, 1.5) {Timing};
    \node[attack, fill=yellow!20] at (6, 0.7) {Patterns};
    \node[attack, fill=yellow!20] at (6, -0.1) {Long-term};
    
    \node[attack, fill=purple!20] at (9, 1.5) {Intersection};
    \node[attack, fill=purple!20] at (9, 0.7) {Side Channel};
    \node[attack, fill=purple!20] at (9, -0.1) {Social Eng.};
\end{tikzpicture}
\end{center}
\end{keyidea}

%===========================================
\section{Deanonymization Attacks}
\label{sec:deanonymization}
%===========================================

\begin{definition}[Deanonymization]
Proceso de vincular actividad anónima con identidad real:
$\text{Deanon}: \text{Activity}_{\text{anon}} \rightarrow \text{Identity}_{\text{real}}$
\end{definition}

\begin{figure}[htbp]
\centering
\begin{tikzpicture}[
    entity/.style={ellipse, draw, minimum width=1.8cm, minimum height=0.9cm, align=center, font=\small},
    attack/.style={rectangle, draw, rounded corners, fill=red!20, font=\tiny}
]
    \node[entity, fill=blue!20] (user) at (0, 0) {Usuario\\Anónimo};
    \node[entity, fill=green!20] (act) at (4, 0) {Actividad\\On-chain};
    \node[entity, fill=purple!20] (id) at (8, 0) {Identidad\\Real};
    
    \node[attack] at (2, 1) {IP};
    \node[attack] at (2, 0.4) {Timing};
    \node[attack] at (6, 1) {Graph};
    \node[attack] at (6, 0.4) {KYC};
    
    \draw[->, thick] (user) -- (act);
    \draw[->, thick] (act) -- (id);
\end{tikzpicture}
\caption{Vectores de deanonymization}
\end{figure}

\begin{lstlisting}[language=Python, caption={Analizador de vulnerabilidades de deanonymization}]
from dataclasses import dataclass, field
from datetime import datetime, timedelta
from typing import Dict, List, Set
from collections import defaultdict
from enum import Enum
import math
import random

class AttackVector(Enum):
    IP_CORRELATION = "ip_correlation"
    TIMING_ANALYSIS = "timing_analysis"
    AMOUNT_PATTERN = "amount_pattern"
    GRAPH_ANALYSIS = "graph_analysis"

@dataclass
class ObservedTransaction:
    tx_id: str
    first_seen_ip: str
    first_seen_time: datetime
    amount: float
    sender: str
    receiver: str

class DeanonymizationAnalyzer:
    """Analiza vulnerabilidades de deanonymization."""
    
    def __init__(self):
        self.tx_observations: List[ObservedTransaction] = []
        self.ip_to_addrs: Dict[str, Set[str]] = defaultdict(set)
        self.addr_to_ips: Dict[str, Set[str]] = defaultdict(set)
        self.addr_timing: Dict[str, List[datetime]] = defaultdict(list)
        self.addr_amounts: Dict[str, List[float]] = defaultdict(list)
    
    def observe(self, tx: ObservedTransaction):
        """Registrar observación del adversario."""
        self.tx_observations.append(tx)
        self.ip_to_addrs[tx.first_seen_ip].add(tx.sender)
        self.addr_to_ips[tx.sender].add(tx.first_seen_ip)
        self.addr_timing[tx.sender].append(tx.first_seen_time)
        self.addr_amounts[tx.sender].append(tx.amount)
    
    def analyze_ip_risk(self, address: str) -> float:
        """Risk score basado en unicidad de IP (0=seguro, 1=vulnerable)."""
        ips = self.addr_to_ips.get(address, set())
        if not ips:
            return 0.0
        if len(ips) == 1:
            return 0.9  # Siempre mismo IP
        return max(0, 1 - len(ips) / 10)
    
    def analyze_timing_risk(self, address: str) -> float:
        """Risk score basado en regularidad temporal."""
        times = self.addr_timing.get(address, [])
        if len(times) < 3:
            return 0.0
        
        sorted_times = sorted(times)
        intervals = [(sorted_times[i+1] - sorted_times[i]).total_seconds() 
                     for i in range(len(sorted_times)-1)]
        
        if not intervals:
            return 0.0
        
        mean = sum(intervals) / len(intervals)
        if mean == 0:
            return 0.0
        
        variance = sum((i - mean)**2 for i in intervals) / len(intervals)
        cv = math.sqrt(variance) / mean
        
        return max(0, 1 - cv)  # CV bajo = regular = vulnerable
    
    def analyze_amount_risk(self, address: str) -> float:
        """Risk score basado en patrones de monto."""
        amounts = self.addr_amounts.get(address, [])
        if len(amounts) < 3:
            return 0.0
        
        round_count = sum(1 for a in amounts if a == round(a))
        return round_count / len(amounts) * 0.5
    
    def evaluate_overall_risk(self, address: str) -> Dict:
        """Evaluación completa de riesgo."""
        ip_risk = self.analyze_ip_risk(address)
        timing_risk = self.analyze_timing_risk(address)
        amount_risk = self.analyze_amount_risk(address)
        
        overall = 0.4 * ip_risk + 0.35 * timing_risk + 0.25 * amount_risk
        
        recommendations = []
        if ip_risk > 0.5:
            recommendations.append("Use Tor/VPN to vary IPs")
        if timing_risk > 0.5:
            recommendations.append("Randomize transaction timing")
        if amount_risk > 0.3:
            recommendations.append("Avoid round amounts")
        
        return {
            'overall_risk': overall,
            'ip_risk': ip_risk,
            'timing_risk': timing_risk,
            'amount_risk': amount_risk,
            'recommendations': recommendations
        }
    
    def correlate_addresses(self, addr1: str, addr2: str) -> float:
        """Probabilidad de que dos direcciones sean del mismo usuario."""
        scores = []
        
        ips1 = self.addr_to_ips.get(addr1, set())
        ips2 = self.addr_to_ips.get(addr2, set())
        if ips1 and ips2:
            overlap = len(ips1 & ips2) / len(ips1 | ips2)
            scores.append(overlap * 0.5)
        
        t1 = self.addr_timing.get(addr1, [])
        t2 = self.addr_timing.get(addr2, [])
        if t1 and t2:
            h1 = set(t.hour for t in t1)
            h2 = set(t.hour for t in t2)
            if h1 and h2:
                scores.append(len(h1 & h2) / len(h1 | h2) * 0.3)
        
        return sum(scores) if scores else 0.0


def deanon_demo():
    print("=== Deanonymization Analysis ===\n")
    
    analyzer = DeanonymizationAnalyzer()
    base = datetime(2024, 1, 15, 9, 0, 0)
    
    # Usuario predecible
    for i in range(15):
        analyzer.observe(ObservedTransaction(
            tx_id=f"tx_{i}",
            first_seen_ip="192.168.1.100",
            first_seen_time=base + timedelta(days=i, hours=random.randint(8,10)),
            amount=round(random.uniform(0.1, 1.0), 1),
            sender="0xPredictable",
            receiver=f"0xR{i}"
        ))
    
    # Usuario privado
    for i in range(15):
        analyzer.observe(ObservedTransaction(
            tx_id=f"tx_p_{i}",
            first_seen_ip=f"10.{random.randint(0,255)}.{random.randint(0,255)}.1",
            first_seen_time=base + timedelta(days=random.randint(0,30), hours=random.randint(0,23)),
            amount=random.uniform(0.1, 1.0),
            sender="0xPrivate",
            receiver=f"0xR_p{i}"
        ))
    
    for addr, name in [("0xPredictable", "Predictable"), ("0xPrivate", "Private")]:
        r = analyzer.evaluate_overall_risk(addr)
        print(f"{name}: Overall={r['overall_risk']:.1%} IP={r['ip_risk']:.1%} "
              f"Timing={r['timing_risk']:.1%}")
        for rec in r['recommendations'][:1]:
            print(f"  -> {rec}")

deanon_demo()
\end{lstlisting}

%===========================================
\section{Sybil Attacks en Mixers}
\label{sec:sybil-attacks}
%===========================================

\begin{definition}[Sybil Attack en Mixer]
Atacante crea múltiples depósitos propios para reducir anonymity set efectivo:
\[
\text{AS}_{\text{real}} = n - k \quad \text{(donde } k \text{ = depósitos del atacante)}
\]
\end{definition}

\begin{figure}[htbp]
\centering
\begin{tikzpicture}[
    dep/.style={circle, draw, minimum size=0.4cm, font=\tiny},
    att/.style={dep, fill=red!30},
    vic/.style={dep, fill=green!30}
]
    \node[rectangle, draw, fill=gray!20, minimum width=2cm, minimum height=1.2cm] (m) at (0,0) {Mixer};
    
    \node[att] at (-2, 0.6) {A};
    \node[att] at (-2, 0) {A};
    \node[att] at (-2, -0.6) {A};
    \node[vic] at (-2, -1.2) {V};
    
    \draw[->] (-1.7, 0.6) -- (-1, 0);
    \draw[->] (-1.7, 0) -- (-1, 0);
    \draw[->] (-1.7, -0.6) -- (-1, 0);
    \draw[->, green!60!black] (-1.7, -1.2) -- (-1, 0);
    
    \node at (0, -2) [font=\small] {AS aparente=4, \textcolor{red}{AS real=1}};
\end{tikzpicture}
\caption{Sybil: atacante controla 3 de 4 depósitos}
\end{figure}

\begin{lstlisting}[language=Python, caption={Simulador de Sybil attacks y defensas}]
from dataclasses import dataclass
from typing import List, Dict
from enum import Enum
import random
import math

class DepositType(Enum):
    HONEST = "honest"
    ATTACKER = "attacker"
    VICTIM = "victim"

@dataclass
class Deposit:
    id: str
    dtype: DepositType
    timestamp: int

class MixerSimulator:
    def __init__(self):
        self.deposits: List[Deposit] = []
        self.attacker_ids: List[str] = []
        self.victim_ids: List[str] = []
    
    def simulate_attack(self, total: int, attacker_frac: float) -> Dict:
        """Simular Sybil attack con fracción de atacante dada."""
        att_count = int(total * attacker_frac)
        victim_count = 1
        honest_count = total - att_count - victim_count
        
        types = ([DepositType.ATTACKER] * att_count + 
                 [DepositType.VICTIM] * victim_count +
                 [DepositType.HONEST] * honest_count)
        random.shuffle(types)
        
        self.deposits = []
        self.attacker_ids = []
        self.victim_ids = []
        
        for i, t in enumerate(types):
            d = Deposit(f"d{i}", t, i)
            self.deposits.append(d)
            if t == DepositType.ATTACKER:
                self.attacker_ids.append(d.id)
            elif t == DepositType.VICTIM:
                self.victim_ids.append(d.id)
        
        # Calcular AS para víctima
        victim = next(d for d in self.deposits if d.dtype == DepositType.VICTIM)
        eligible = [d for d in self.deposits if d.timestamp <= victim.timestamp]
        apparent = len(eligible)
        real = sum(1 for d in eligible if d.dtype != DepositType.ATTACKER)
        
        return {
            'apparent_as': apparent,
            'real_as': real,
            'privacy_loss': (apparent - real) / apparent if apparent > 0 else 0
        }


class SybilDefense:
    def __init__(self, base_fee: float = 0.01):
        self.base_fee = base_fee
    
    def analyze_defense(self, budget: float, deposit_amt: float, 
                        fee: float, stake: float, pop: bool) -> Dict:
        """Analizar efectividad de defensa contra presupuesto de atacante."""
        cost = deposit_amt + fee + stake
        max_by_budget = int(budget / cost) if cost > 0 else 1000
        max_effective = 1 if pop else max_by_budget
        
        return {
            'max_sybil_deposits': max_effective,
            'cost_per_deposit': cost,
            'effectiveness': 1 - max_effective / max(1, max_by_budget)
        }


def sybil_demo():
    print("=== Sybil Attack Analysis ===\n")
    
    sim = MixerSimulator()
    
    print(f"{'Attacker%':<10} {'Apparent':<10} {'Real':<8} {'Loss':<8}")
    print("-" * 36)
    for frac in [0, 0.25, 0.5, 0.75, 0.9]:
        r = sim.simulate_attack(100, frac)
        print(f"{frac*100:>8.0f}%  {r['apparent_as']:>8}  {r['real_as']:>6}  "
              f"{r['privacy_loss']*100:>6.1f}%")
    
    print("\n--- Defense Effectiveness ---")
    defense = SybilDefense()
    budget = 100.0
    
    scenarios = [
        ("Baseline", 0.01, 0, False),
        ("High fee", 0.5, 0, False),
        ("Staking", 0.01, 10, False),
        ("PoP", 0.01, 0, True),
    ]
    
    for name, fee, stake, pop in scenarios:
        r = defense.analyze_defense(budget, 1.0, fee, stake, pop)
        print(f"{name:<12} Max Sybil: {r['max_sybil_deposits']:<5} "
              f"Effectiveness: {r['effectiveness']*100:.0f}%")

sybil_demo()
\end{lstlisting}

%===========================================
\section{Eclipse Attacks}
\label{sec:eclipse}
%===========================================

\begin{definition}[Eclipse Attack]
Aislar un nodo de la red honesta controlando todas sus conexiones:
\begin{itemize}
    \item Atacante monopoliza peers de la víctima
    \item Puede alimentar información falsa (bloques, txs)
    \item Puede censurar transacciones de la víctima
\end{itemize}
\end{definition}

\begin{figure}[htbp]
\centering
\begin{tikzpicture}[
    node/.style={circle, draw, minimum size=0.6cm, font=\small},
    honest/.style={node, fill=green!20},
    attacker/.style={node, fill=red!20},
    victim/.style={node, fill=blue!30}
]
    % Red honesta
    \node[honest] (h1) at (0, 2) {H};
    \node[honest] (h2) at (1.5, 2.5) {H};
    \node[honest] (h3) at (3, 2) {H};
    \node[honest] (h4) at (1.5, 3.5) {H};
    
    \draw (h1) -- (h2) -- (h3);
    \draw (h2) -- (h4);
    
    % Víctima eclipsada
    \node[victim] (v) at (6, 2) {V};
    
    % Atacantes rodeando
    \node[attacker] (a1) at (5, 3) {A};
    \node[attacker] (a2) at (7, 3) {A};
    \node[attacker] (a3) at (5, 1) {A};
    \node[attacker] (a4) at (7, 1) {A};
    
    \draw[red, thick] (v) -- (a1);
    \draw[red, thick] (v) -- (a2);
    \draw[red, thick] (v) -- (a3);
    \draw[red, thick] (v) -- (a4);
    
    % Sin conexión a red honesta
    \draw[dashed, gray] (h3) -- (4.5, 2);
    \node at (4, 2.3) [font=\tiny, gray] {blocked};
\end{tikzpicture}
\caption{Eclipse attack: víctima solo conectada a atacantes}
\end{figure}

\begin{lstlisting}[language=Python, caption={Simulador de Eclipse attacks}]
from dataclasses import dataclass, field
from typing import Dict, List, Set
import random

@dataclass
class NetworkNode:
    node_id: str
    is_honest: bool
    peers: Set[str] = field(default_factory=set)
    max_peers: int = 8

class EclipseSimulator:
    def __init__(self, honest_nodes: int, attacker_nodes: int):
        self.nodes: Dict[str, NetworkNode] = {}
        
        # Crear nodos honestos
        for i in range(honest_nodes):
            self.nodes[f"H{i}"] = NetworkNode(f"H{i}", True)
        
        # Crear nodos atacantes
        for i in range(attacker_nodes):
            self.nodes[f"A{i}"] = NetworkNode(f"A{i}", False)
        
        # Víctima
        self.nodes["V"] = NetworkNode("V", True)
    
    def connect_honest_network(self):
        """Conectar nodos honestos entre sí."""
        honest = [n for n in self.nodes.values() if n.is_honest and n.node_id != "V"]
        
        for node in honest:
            available = [n for n in honest if n.node_id != node.node_id 
                        and len(node.peers) < node.max_peers]
            connections = min(4, len(available))
            for peer in random.sample(available, connections):
                node.peers.add(peer.node_id)
                peer.peers.add(node.node_id)
    
    def execute_eclipse(self, victim_id: str = "V") -> Dict:
        """Ejecutar eclipse attack."""
        victim = self.nodes[victim_id]
        attackers = [n for n in self.nodes.values() if not n.is_honest]
        
        # Atacantes intentan llenar todos los slots de la víctima
        for attacker in attackers:
            if len(victim.peers) < victim.max_peers:
                victim.peers.add(attacker.node_id)
                attacker.peers.add(victim.node_id)
        
        # Calcular resultado
        honest_peers = sum(1 for p in victim.peers if self.nodes[p].is_honest)
        attacker_peers = len(victim.peers) - honest_peers
        
        eclipsed = honest_peers == 0 and attacker_peers > 0
        
        return {
            'total_peers': len(victim.peers),
            'honest_peers': honest_peers,
            'attacker_peers': attacker_peers,
            'eclipsed': eclipsed,
            'eclipse_ratio': attacker_peers / len(victim.peers) if victim.peers else 0
        }
    
    def apply_defense(self, victim_id: str = "V", 
                      reserved_honest: int = 2) -> Dict:
        """Aplicar defensa: reservar slots para nodos conocidos."""
        victim = self.nodes[victim_id]
        
        # Primero conectar a nodos honestos conocidos
        honest = [n for n in self.nodes.values() if n.is_honest and n.node_id != victim_id]
        
        for i, node in enumerate(honest[:reserved_honest]):
            if len(victim.peers) < victim.max_peers:
                victim.peers.add(node.node_id)
        
        # Luego atacantes intentan llenar el resto
        attackers = [n for n in self.nodes.values() if not n.is_honest]
        for attacker in attackers:
            if len(victim.peers) < victim.max_peers:
                victim.peers.add(attacker.node_id)
        
        honest_peers = sum(1 for p in victim.peers if self.nodes[p].is_honest)
        
        return {
            'honest_peers': honest_peers,
            'eclipsed': honest_peers == 0,
            'defense_effective': honest_peers >= reserved_honest
        }


def eclipse_demo():
    print("=== Eclipse Attack Simulation ===\n")
    
    # Sin defensa
    sim = EclipseSimulator(honest_nodes=20, attacker_nodes=50)
    sim.connect_honest_network()
    
    result = sim.execute_eclipse()
    print("Without defense:")
    print(f"  Peers: {result['total_peers']} (Honest: {result['honest_peers']}, "
          f"Attacker: {result['attacker_peers']})")
    print(f"  Eclipsed: {result['eclipsed']}")
    
    # Con defensa
    sim2 = EclipseSimulator(honest_nodes=20, attacker_nodes=50)
    sim2.connect_honest_network()
    
    result2 = sim2.apply_defense(reserved_honest=2)
    print("\nWith defense (2 reserved honest slots):")
    print(f"  Honest peers: {result2['honest_peers']}")
    print(f"  Eclipsed: {result2['eclipsed']}")
    print(f"  Defense effective: {result2['defense_effective']}")

eclipse_demo()
\end{lstlisting}

\subsection{Defensas contra Eclipse}

\begin{table}[htbp]
\centering
\small
\begin{tabular}{lp{5cm}c}
\toprule
\textbf{Defensa} & \textbf{Mecanismo} & \textbf{Efectividad} \\
\midrule
Peer diversity & Limitar peers por subnet/AS & Alta \\
Reserved slots & Slots para nodos conocidos/confiables & Alta \\
Checkpoints & Verificar contra estado conocido & Media \\
Multiple connections & Conexiones a diferentes redes & Alta \\
Proof of Work & Costo para crear identidades & Media \\
\bottomrule
\end{tabular}
\caption{Defensas contra Eclipse attacks}
\end{table}
%===========================================
\section{Dust Attacks}
\label{sec:dust-attacks}
%===========================================

\begin{definition}[Dust Attack]
Enviar cantidades mínimas (``dust'') a múltiples direcciones para rastrear cuando las víctimas consoliden esos fondos con otros UTXOs, revelando qué direcciones controlan.
\end{definition}

\begin{figure}[htbp]
\centering
\begin{tikzpicture}[
    addr/.style={rectangle, draw, minimum width=1.5cm, minimum height=0.6cm, font=\small},
    tx/.style={circle, draw, fill=gray!20, minimum size=0.5cm}
]
    % Fase 1: Dust enviado
    \node at (-3, 3) {\textbf{Fase 1: Enviar dust}};
    
    \node[addr, fill=red!20] (att) at (-3, 2) {Atacante};
    \node[addr, fill=blue!20] (v1) at (0, 2.5) {Addr A};
    \node[addr, fill=blue!20] (v2) at (0, 2) {Addr B};
    \node[addr, fill=blue!20] (v3) at (0, 1.5) {Addr C};
    
    \draw[->, red] (att) -- node[above, font=\tiny] {dust} (v1);
    \draw[->, red] (att) -- node[above, font=\tiny] {dust} (v2);
    \draw[->, red] (att) -- node[above, font=\tiny] {dust} (v3);
    
    % Fase 2: Consolidación
    \node at (-3, 0) {\textbf{Fase 2: Víctima consolida}};
    
    \node[addr, fill=blue!20] (va) at (-1, -0.5) {Addr A\\+ dust};
    \node[addr, fill=blue!20] (vb) at (-1, -1.3) {Addr B\\+ dust};
    \node[tx] (tx) at (1.5, -0.9) {TX};
    \node[addr, fill=green!20] (out) at (4, -0.9) {Addr D};
    
    \draw[->] (va) -- (tx);
    \draw[->] (vb) -- (tx);
    \draw[->] (tx) -- (out);
    
    % Fase 3: Análisis
    \node at (-3, -2.5) {\textbf{Fase 3: Atacante deduce}};
    \node at (1.5, -2.5) [font=\small, align=center] {A, B, D pertenecen\\al mismo usuario};
\end{tikzpicture}
\caption{Dust attack: tracking mediante consolidación de UTXOs}
\label{fig:dust-attack}
\end{figure}

\begin{lstlisting}[language=Python, caption={Detector y defensa contra Dust attacks}]
from dataclasses import dataclass
from typing import Dict, List, Set, Tuple
from datetime import datetime
from collections import defaultdict

@dataclass
class UTXO:
    tx_id: str
    output_index: int
    address: str
    amount: float
    timestamp: datetime
    source_addresses: Set[str] = None  # Para tracking de inputs
    
    def __post_init__(self):
        if self.source_addresses is None:
            self.source_addresses = set()

@dataclass
class Transaction:
    tx_id: str
    inputs: List[UTXO]
    outputs: List[UTXO]
    timestamp: datetime

class DustAttackDetector:
    """
    Detecta dust attacks y proporciona recomendaciones.
    """
    
    def __init__(self, dust_threshold: float = 0.0001):
        self.dust_threshold = dust_threshold
        self.utxos: Dict[str, List[UTXO]] = defaultdict(list)  # address -> UTXOs
        self.suspicious_dust: Dict[str, List[UTXO]] = defaultdict(list)
        self.known_attackers: Set[str] = set()
    
    def add_utxo(self, utxo: UTXO):
        """Añadir UTXO y verificar si es dust sospechoso."""
        self.utxos[utxo.address].append(utxo)
        
        if self.is_dust(utxo):
            self.suspicious_dust[utxo.address].append(utxo)
    
    def is_dust(self, utxo: UTXO) -> bool:
        """Verificar si UTXO es dust."""
        return utxo.amount <= self.dust_threshold
    
    def analyze_incoming_dust(self, address: str) -> Dict:
        """Analizar dust recibido por una dirección."""
        dust_utxos = self.suspicious_dust.get(address, [])
        
        if not dust_utxos:
            return {'has_dust': False, 'risk_level': 'none'}
        
        # Analizar fuentes del dust
        source_addresses = set()
        for utxo in dust_utxos:
            source_addresses.update(utxo.source_addresses)
        
        # Calcular riesgo
        risk_factors = {
            'dust_count': len(dust_utxos),
            'unique_sources': len(source_addresses),
            'total_dust_value': sum(u.amount for u in dust_utxos),
            'recent_dust': sum(1 for u in dust_utxos 
                              if (datetime.now() - u.timestamp).days < 7)
        }
        
        # Determinar nivel de riesgo
        if risk_factors['dust_count'] > 5 or risk_factors['unique_sources'] == 1:
            risk_level = 'high'
        elif risk_factors['dust_count'] > 2:
            risk_level = 'medium'
        else:
            risk_level = 'low'
        
        return {
            'has_dust': True,
            'risk_level': risk_level,
            'dust_count': risk_factors['dust_count'],
            'total_value': risk_factors['total_dust_value'],
            'recommendations': self._get_recommendations(risk_level)
        }
    
    def _get_recommendations(self, risk_level: str) -> List[str]:
        """Obtener recomendaciones basadas en nivel de riesgo."""
        if risk_level == 'high':
            return [
                "DO NOT spend this dust with other UTXOs",
                "Consider this address compromised for privacy",
                "Use coin control to exclude dust from transactions",
                "Consider consolidating only dust separately"
            ]
        elif risk_level == 'medium':
            return [
                "Be cautious when spending from this address",
                "Use coin control if available",
                "Consider waiting before using these funds"
            ]
        else:
            return ["Monitor for additional suspicious deposits"]
    
    def simulate_consolidation_leak(
        self,
        addresses: List[str],
        include_dust: bool
    ) -> Dict:
        """
        Simular qué información se filtra al consolidar.
        """
        # Recopilar todos los UTXOs
        all_utxos = []
        dust_utxos = []
        
        for addr in addresses:
            for utxo in self.utxos.get(addr, []):
                if self.is_dust(utxo):
                    dust_utxos.append(utxo)
                    if include_dust:
                        all_utxos.append(utxo)
                else:
                    all_utxos.append(utxo)
        
        # Información filtrada
        if include_dust:
            linked_addresses = set(addresses)
            dust_sources = set()
            for utxo in dust_utxos:
                dust_sources.update(utxo.source_addresses)
            
            return {
                'addresses_linked': len(linked_addresses),
                'linked_addresses': linked_addresses,
                'dust_trackers_alerted': len(dust_sources),
                'privacy_leak': 'HIGH - All addresses now linked'
            }
        else:
            return {
                'addresses_linked': len(addresses),
                'dust_excluded': len(dust_utxos),
                'privacy_leak': 'LOW - Dust not included'
            }


class CoinControl:
    """
    Sistema de coin control para gestión segura de UTXOs.
    """
    
    def __init__(self):
        self.utxos: List[UTXO] = []
        self.frozen_utxos: Set[Tuple[str, int]] = set()  # (tx_id, output_index)
        self.labels: Dict[Tuple[str, int], str] = {}
    
    def add_utxo(self, utxo: UTXO, label: str = None):
        """Añadir UTXO con etiqueta opcional."""
        self.utxos.append(utxo)
        if label:
            self.labels[(utxo.tx_id, utxo.output_index)] = label
    
    def freeze_utxo(self, tx_id: str, output_index: int):
        """Congelar UTXO para que no se use automáticamente."""
        self.frozen_utxos.add((tx_id, output_index))
    
    def is_frozen(self, utxo: UTXO) -> bool:
        return (utxo.tx_id, utxo.output_index) in self.frozen_utxos
    
    def auto_freeze_dust(self, threshold: float = 0.0001):
        """Congelar automáticamente todo el dust."""
        for utxo in self.utxos:
            if utxo.amount <= threshold:
                self.freeze_utxo(utxo.tx_id, utxo.output_index)
    
    def select_utxos_for_amount(
        self,
        target_amount: float,
        exclude_dust: bool = True,
        same_source_only: bool = False
    ) -> List[UTXO]:
        """
        Seleccionar UTXOs para una transacción.
        """
        available = [
            u for u in self.utxos
            if not self.is_frozen(u)
            and (not exclude_dust or u.amount > 0.0001)
        ]
        
        if same_source_only and available:
            # Solo usar UTXOs del mismo origen para evitar linking
            first_source = available[0].source_addresses
            available = [u for u in available 
                        if u.source_addresses == first_source]
        
        # Selección simple: ordenar por cantidad y tomar los necesarios
        available.sort(key=lambda u: u.amount, reverse=True)
        
        selected = []
        total = 0.0
        
        for utxo in available:
            if total >= target_amount:
                break
            selected.append(utxo)
            total += utxo.amount
        
        return selected if total >= target_amount else []
    
    def get_privacy_score(self) -> Dict:
        """Evaluar privacidad del conjunto de UTXOs."""
        total_utxos = len(self.utxos)
        dust_count = sum(1 for u in self.utxos if u.amount <= 0.0001)
        frozen_dust = sum(1 for u in self.utxos 
                         if u.amount <= 0.0001 and self.is_frozen(u))
        
        unique_sources = len(set(
            frozenset(u.source_addresses) for u in self.utxos
        ))
        
        # Score de 0 a 100
        dust_risk = (dust_count - frozen_dust) / max(1, total_utxos) * 50
        source_diversity = min(50, unique_sources * 10)
        
        score = 100 - dust_risk + source_diversity / 2
        
        return {
            'score': min(100, max(0, score)),
            'total_utxos': total_utxos,
            'dust_utxos': dust_count,
            'frozen_dust': frozen_dust,
            'unfrozen_dust': dust_count - frozen_dust,
            'recommendation': 'Freeze remaining dust' if dust_count > frozen_dust else 'Good'
        }


def dust_attack_demo():
    print("=== Dust Attack Detection Demo ===\n")
    
    detector = DustAttackDetector(dust_threshold=0.0001)
    
    # Simular recepción de dust
    victim_address = "bc1qvictim..."
    
    # UTXOs normales
    detector.add_utxo(UTXO("tx1", 0, victim_address, 0.5, datetime.now(), {"exchange"}))
    detector.add_utxo(UTXO("tx2", 0, victim_address, 0.3, datetime.now(), {"friend"}))
    
    # Dust attacks
    for i in range(5):
        detector.add_utxo(UTXO(
            f"dust_tx_{i}", 0, victim_address, 
            0.00001,  # Dust
            datetime.now(),
            {"attacker_wallet"}
        ))
    
    # Analizar
    analysis = detector.analyze_incoming_dust(victim_address)
    
    print(f"Address: {victim_address}")
    print(f"Has dust: {analysis['has_dust']}")
    print(f"Risk level: {analysis['risk_level']}")
    print(f"Dust count: {analysis['dust_count']}")
    print(f"Total dust value: {analysis['total_value']:.8f}")
    print("\nRecommendations:")
    for rec in analysis['recommendations']:
        print(f"  - {rec}")
    
    # Simular consolidación
    print("\n--- Consolidation Simulation ---")
    
    # Con dust
    leak_with = detector.simulate_consolidation_leak([victim_address], include_dust=True)
    print(f"With dust: {leak_with['privacy_leak']}")
    
    # Sin dust
    leak_without = detector.simulate_consolidation_leak([victim_address], include_dust=False)
    print(f"Without dust: {leak_without['privacy_leak']}")

dust_attack_demo()
\end{lstlisting}

%===========================================
\section{Graph Analysis (Clustering Heuristics)}
\label{sec:graph-analysis}
%===========================================

\begin{definition}[Address Clustering]
Agrupar direcciones que probablemente pertenecen al mismo usuario usando heurísticas:

\begin{itemize}
    \item \textbf{Common Input Ownership:} Inputs de la misma TX = mismo dueño
    \item \textbf{Change Detection:} Output de cambio = mismo dueño que inputs
    \item \textbf{Address Reuse:} Misma dirección = mismo usuario
\end{itemize}
\end{definition}

\begin{figure}[htbp]
\centering
\begin{tikzpicture}[
    addr/.style={circle, draw, minimum size=0.7cm, font=\tiny},
    tx/.style={rectangle, draw, fill=gray!20, minimum width=0.5cm, minimum height=0.3cm}
]
    % Common Input Heuristic
    \node at (-2, 2.5) {\small \textbf{Common Input}};
    
    \node[addr, fill=blue!20] (a1) at (-3, 1.5) {A};
    \node[addr, fill=blue!20] (a2) at (-2, 1.5) {B};
    \node[addr, fill=blue!20] (a3) at (-1, 1.5) {C};
    \node[tx] (tx1) at (-2, 0.5) {TX};
    \node[addr, fill=green!20] (out) at (-2, -0.3) {D};
    
    \draw[->] (a1) -- (tx1);
    \draw[->] (a2) -- (tx1);
    \draw[->] (a3) -- (tx1);
    \draw[->] (tx1) -- (out);
    
    \node at (-2, -1) [font=\tiny, align=center] {A, B, C = mismo dueño};
    
    % Change Detection
    \node at (3, 2.5) {\small \textbf{Change Detection}};
    
    \node[addr, fill=blue!20] (i1) at (2, 1.5) {Input\\1 BTC};
    \node[tx] (tx2) at (3, 0.8) {TX};
    \node[addr, fill=green!20] (p1) at (2, 0) {Pago\\0.3};
    \node[addr, fill=orange!20] (ch) at (4, 0) {Cambio\\0.7};
    
    \draw[->] (i1) -- (tx2);
    \draw[->] (tx2) -- (p1);
    \draw[->] (tx2) -- (ch);
    
    \node at (3, -1) [font=\tiny, align=center] {Input y Cambio =\\mismo dueño};
\end{tikzpicture}
\caption{Heurísticas de clustering: common input y change detection}
\label{fig:clustering-heuristics}
\end{figure}

\begin{lstlisting}[language=Python, caption={Sistema de Graph Analysis y Clustering}]
from dataclasses import dataclass, field
from typing import Dict, List, Set, Tuple
from collections import defaultdict

@dataclass
class TxInput:
    address: str
    amount: float
    tx_id: str
    output_index: int

@dataclass
class TxOutput:
    address: str
    amount: float
    output_index: int

@dataclass
class Transaction:
    tx_id: str
    inputs: List[TxInput]
    outputs: List[TxOutput]
    fee: float = 0.0

class AddressClusterer:
    """
    Clustering de direcciones usando heurísticas on-chain.
    """
    
    def __init__(self):
        self.transactions: Dict[str, Transaction] = {}
        self.address_to_txs: Dict[str, List[str]] = defaultdict(list)
        
        # Union-Find para clustering
        self.parent: Dict[str, str] = {}
        self.rank: Dict[str, int] = {}
    
    def _find(self, addr: str) -> str:
        """Find con path compression."""
        if addr not in self.parent:
            self.parent[addr] = addr
            self.rank[addr] = 0
        
        if self.parent[addr] != addr:
            self.parent[addr] = self._find(self.parent[addr])
        
        return self.parent[addr]
    
    def _union(self, addr1: str, addr2: str):
        """Union by rank."""
        root1 = self._find(addr1)
        root2 = self._find(addr2)
        
        if root1 == root2:
            return
        
        if self.rank[root1] < self.rank[root2]:
            self.parent[root1] = root2
        elif self.rank[root1] > self.rank[root2]:
            self.parent[root2] = root1
        else:
            self.parent[root2] = root1
            self.rank[root1] += 1
    
    def add_transaction(self, tx: Transaction):
        """Añadir transacción y aplicar heurísticas."""
        self.transactions[tx.tx_id] = tx
        
        # Registrar direcciones
        for inp in tx.inputs:
            self.address_to_txs[inp.address].append(tx.tx_id)
        for out in tx.outputs:
            self.address_to_txs[out.address].append(tx.tx_id)
        
        # Aplicar heurísticas
        self._apply_common_input_heuristic(tx)
        self._apply_change_heuristic(tx)
    
    def _apply_common_input_heuristic(self, tx: Transaction):
        """
        Common Input Ownership Heuristic:
        Si múltiples direcciones son inputs de la misma TX,
        probablemente pertenecen al mismo usuario.
        """
        if len(tx.inputs) <= 1:
            return
        
        input_addresses = [inp.address for inp in tx.inputs]
        
        # Unir todas las direcciones de input
        first = input_addresses[0]
        for addr in input_addresses[1:]:
            self._union(first, addr)
    
    def _apply_change_heuristic(self, tx: Transaction):
        """
        Change Detection Heuristic:
        Detectar output de cambio y vincularlo con inputs.
        """
        if len(tx.outputs) != 2 or len(tx.inputs) == 0:
            return
        
        # Heurísticas para detectar cambio:
        # 1. Output que es address reusada (ya vista antes)
        # 2. Output con monto "raro" (no redondo)
        # 3. Output más pequeño (el pago suele ser específico)
        
        input_total = sum(inp.amount for inp in tx.inputs)
        output_amounts = [(out.address, out.amount) for out in tx.outputs]
        output_amounts.sort(key=lambda x: x[1])
        
        # Heurística simple: el output más pequeño es el pago,
        # el más grande es el cambio
        larger_output = output_amounts[1]
        
        # Verificar que no es pago redondo (más probable que sea cambio)
        if larger_output[1] != round(larger_output[1], 2):
            # Probablemente es cambio - vincular con input
            change_addr = larger_output[0]
            input_addr = tx.inputs[0].address
            self._union(change_addr, input_addr)
    
    def get_cluster(self, address: str) -> Set[str]:
        """Obtener todas las direcciones en el mismo cluster."""
        if address not in self.parent:
            return {address}
        
        root = self._find(address)
        return {addr for addr in self.parent if self._find(addr) == root}
    
    def get_all_clusters(self) -> List[Set[str]]:
        """Obtener todos los clusters."""
        clusters = defaultdict(set)
        
        for addr in self.parent:
            root = self._find(addr)
            clusters[root].add(addr)
        
        return list(clusters.values())
    
    def analyze_cluster(self, address: str) -> Dict:
        """Analizar un cluster específico."""
        cluster = self.get_cluster(address)
        
        if not cluster:
            return {'size': 0}
        
        # Calcular estadísticas
        total_txs = set()
        total_volume_in = 0.0
        total_volume_out = 0.0
        
        for addr in cluster:
            for tx_id in self.address_to_txs.get(addr, []):
                total_txs.add(tx_id)
                tx = self.transactions.get(tx_id)
                if tx:
                    for inp in tx.inputs:
                        if inp.address == addr:
                            total_volume_out += inp.amount
                    for out in tx.outputs:
                        if out.address == addr:
                            total_volume_in += out.amount
        
        return {
            'cluster_size': len(cluster),
            'addresses': list(cluster)[:5],  # Primeras 5
            'transaction_count': len(total_txs),
            'total_volume_in': total_volume_in,
            'total_volume_out': total_volume_out,
            'net_flow': total_volume_in - total_volume_out
        }


def clustering_demo():
    print("=== Address Clustering Demo ===\n")
    
    clusterer = AddressClusterer()
    
    # Transacción 1: Multi-input (common input heuristic)
    tx1 = Transaction(
        tx_id="tx1",
        inputs=[
            TxInput("addr_A", 0.5, "prev1", 0),
            TxInput("addr_B", 0.3, "prev2", 0),
            TxInput("addr_C", 0.2, "prev3", 0),
        ],
        outputs=[
            TxOutput("addr_merchant", 0.9, 0),
            TxOutput("addr_D", 0.09, 1),  # Cambio
        ],
        fee=0.01
    )
    clusterer.add_transaction(tx1)
    
    # Transacción 2: Pago simple con cambio
    tx2 = Transaction(
        tx_id="tx2",
        inputs=[TxInput("addr_E", 1.0, "prev4", 0)],
        outputs=[
            TxOutput("addr_service", 0.3, 0),  # Pago redondo
            TxOutput("addr_F", 0.699, 1),  # Cambio (no redondo)
        ],
        fee=0.001
    )
    clusterer.add_transaction(tx2)
    
    # Analizar clusters
    print("--- Clusters Found ---")
    clusters = clusterer.get_all_clusters()
    
    for i, cluster in enumerate(clusters):
        if len(cluster) > 1:
            print(f"\nCluster {i+1}: {len(cluster)} addresses")
            for addr in list(cluster)[:5]:
                print(f"  - {addr}")
    
    # Analizar cluster específico
    print("\n--- Cluster Analysis (addr_A) ---")
    analysis = clusterer.analyze_cluster("addr_A")
    print(f"Size: {analysis['cluster_size']} addresses")
    print(f"Transactions: {analysis['transaction_count']}")
    print(f"Volume in: {analysis['total_volume_in']:.4f}")
    print(f"Volume out: {analysis['total_volume_out']:.4f}")

clustering_demo()
\end{lstlisting}

%===========================================
\section{Taint Analysis}
\label{sec:taint-analysis}
%===========================================

\begin{definition}[Taint Analysis]
Rastrear ``contaminación'' de fondos a través del grafo de transacciones:
\begin{itemize}
    \item \textbf{Poison:} Cualquier contacto con tainted = tainted
    \item \textbf{Haircut:} Proporcional al monto mezclado
    \item \textbf{FIFO/LIFO:} Orden de inputs determina taint
\end{itemize}
\end{definition}

\begin{figure}[htbp]
\centering
\begin{tikzpicture}[
    node/.style={circle, draw, minimum size=0.6cm, font=\tiny},
    tainted/.style={node, fill=red!40},
    clean/.style={node, fill=green!30},
    mixed/.style={node, fill=orange!30}
]
    % Poison method
    \node at (0, 2.5) {\small \textbf{Poison}};
    \node[tainted] (t1) at (-0.5, 1.5) {T};
    \node[clean] (c1) at (0.5, 1.5) {C};
    \node[tainted] (r1) at (0, 0.5) {T};
    \draw[->] (t1) -- (r1);
    \draw[->] (c1) -- (r1);
    \node at (0, -0.2) [font=\tiny] {Todo tainted};
    
    % Haircut method  
    \node at (3, 2.5) {\small \textbf{Haircut}};
    \node[tainted] (t2) at (2.5, 1.5) {0.5T};
    \node[clean] (c2) at (3.5, 1.5) {0.5C};
    \node[mixed] (r2) at (3, 0.5) {50\%};
    \draw[->] (t2) -- (r2);
    \draw[->] (c2) -- (r2);
    \node at (3, -0.2) [font=\tiny] {Proporcional};
    
    % FIFO method
    \node at (6, 2.5) {\small \textbf{FIFO}};
    \node[tainted] (t3) at (5.5, 1.5) {1°};
    \node[clean] (c3) at (6.5, 1.5) {2°};
    \node[tainted] (o1) at (5.5, 0.5) {T};
    \node[clean] (o2) at (6.5, 0.5) {C};
    \draw[->] (t3) -- (o1);
    \draw[->] (c3) -- (o2);
    \node at (6, -0.2) [font=\tiny] {Por orden};
\end{tikzpicture}
\caption{Métodos de taint analysis: Poison, Haircut, FIFO}
\end{figure}

\begin{lstlisting}[language=Python, caption={Sistema de Taint Analysis}]
from dataclasses import dataclass
from typing import Dict, List, Set, Optional
from enum import Enum
from collections import defaultdict

class TaintMethod(Enum):
    POISON = "poison"      # Cualquier contacto = tainted
    HAIRCUT = "haircut"    # Proporcional
    FIFO = "fifo"          # First-in-first-out
    LIFO = "lifo"          # Last-in-first-out

@dataclass
class TaintedFund:
    amount: float
    taint_ratio: float  # 0.0 = limpio, 1.0 = completamente tainted
    source: str         # Origen del taint
    hops: int          # Distancia del origen

@dataclass
class TxForTaint:
    tx_id: str
    inputs: List[tuple]   # (address, amount)
    outputs: List[tuple]  # (address, amount)

class TaintAnalyzer:
    """
    Analiza taint/contaminación de fondos.
    """
    
    def __init__(self, method: TaintMethod = TaintMethod.HAIRCUT):
        self.method = method
        self.address_taint: Dict[str, TaintedFund] = {}
        self.tainted_sources: Set[str] = set()
        self.transactions: List[TxForTaint] = []
    
    def mark_tainted(self, address: str, source: str = "illicit"):
        """Marcar una dirección como fuente de taint."""
        self.tainted_sources.add(address)
        self.address_taint[address] = TaintedFund(
            amount=float('inf'),  # Cualquier cantidad
            taint_ratio=1.0,
            source=source,
            hops=0
        )
    
    def process_transaction(self, tx: TxForTaint):
        """Procesar transacción y propagar taint."""
        self.transactions.append(tx)
        
        # Calcular taint de inputs
        input_taint = self._calculate_input_taint(tx.inputs)
        
        # Propagar a outputs según método
        if self.method == TaintMethod.POISON:
            self._propagate_poison(tx, input_taint)
        elif self.method == TaintMethod.HAIRCUT:
            self._propagate_haircut(tx, input_taint)
        elif self.method in [TaintMethod.FIFO, TaintMethod.LIFO]:
            self._propagate_fifo_lifo(tx, input_taint)
    
    def _calculate_input_taint(self, inputs: List[tuple]) -> Dict:
        """Calcular taint total de los inputs."""
        total_amount = sum(amt for _, amt in inputs)
        total_taint = 0.0
        max_hops = 0
        sources = set()
        
        for addr, amt in inputs:
            if addr in self.address_taint:
                taint = self.address_taint[addr]
                total_taint += taint.taint_ratio * amt
                max_hops = max(max_hops, taint.hops)
                sources.add(taint.source)
        
        return {
            'total_amount': total_amount,
            'taint_amount': total_taint,
            'taint_ratio': total_taint / total_amount if total_amount > 0 else 0,
            'max_hops': max_hops,
            'sources': sources
        }
    
    def _propagate_poison(self, tx: TxForTaint, input_taint: Dict):
        """Método Poison: cualquier contacto = 100% tainted."""
        if input_taint['taint_ratio'] > 0:
            for addr, amt in tx.outputs:
                self.address_taint[addr] = TaintedFund(
                    amount=amt,
                    taint_ratio=1.0,
                    source=list(input_taint['sources'])[0] if input_taint['sources'] else "unknown",
                    hops=input_taint['max_hops'] + 1
                )
    
    def _propagate_haircut(self, tx: TxForTaint, input_taint: Dict):
        """Método Haircut: taint proporcional."""
        if input_taint['taint_ratio'] > 0:
            for addr, amt in tx.outputs:
                self.address_taint[addr] = TaintedFund(
                    amount=amt,
                    taint_ratio=input_taint['taint_ratio'],
                    source=list(input_taint['sources'])[0] if input_taint['sources'] else "unknown",
                    hops=input_taint['max_hops'] + 1
                )
    
    def _propagate_fifo_lifo(self, tx: TxForTaint, input_taint: Dict):
        """Método FIFO/LIFO: asignar taint por orden."""
        # Ordenar inputs por taint (tainted primero para FIFO)
        sorted_inputs = []
        for addr, amt in tx.inputs:
            taint = self.address_taint.get(addr)
            ratio = taint.taint_ratio if taint else 0
            sorted_inputs.append((addr, amt, ratio))
        
        if self.method == TaintMethod.FIFO:
            sorted_inputs.sort(key=lambda x: -x[2])  # Tainted primero
        else:  # LIFO
            sorted_inputs.sort(key=lambda x: x[2])   # Clean primero
        
        # Asignar a outputs en orden
        remaining_taint = sum(amt * ratio for _, amt, ratio in sorted_inputs)
        
        for addr, amt in tx.outputs:
            if remaining_taint >= amt:
                self.address_taint[addr] = TaintedFund(amt, 1.0, "fifo", 1)
                remaining_taint -= amt
            elif remaining_taint > 0:
                ratio = remaining_taint / amt
                self.address_taint[addr] = TaintedFund(amt, ratio, "fifo", 1)
                remaining_taint = 0
            # Else: output stays clean
    
    def get_taint_score(self, address: str) -> Dict:
        """Obtener score de taint para una dirección."""
        if address not in self.address_taint:
            return {
                'address': address,
                'taint_ratio': 0.0,
                'status': 'clean',
                'risk_level': 'none'
            }
        
        taint = self.address_taint[address]
        
        if taint.taint_ratio >= 0.9:
            risk = 'critical'
        elif taint.taint_ratio >= 0.5:
            risk = 'high'
        elif taint.taint_ratio >= 0.1:
            risk = 'medium'
        else:
            risk = 'low'
        
        return {
            'address': address,
            'taint_ratio': taint.taint_ratio,
            'hops_from_source': taint.hops,
            'source': taint.source,
            'risk_level': risk
        }
    
    def trace_funds(self, address: str, max_hops: int = 10) -> List[Dict]:
        """Rastrear origen de fondos."""
        # Simplified: return taint info
        taint = self.address_taint.get(address)
        if not taint:
            return [{'address': address, 'status': 'clean'}]
        
        return [{
            'address': address,
            'taint_ratio': taint.taint_ratio,
            'source': taint.source,
            'hops': taint.hops
        }]


def taint_demo():
    print("=== Taint Analysis Demo ===\n")
    
    # Comparar métodos
    for method in [TaintMethod.POISON, TaintMethod.HAIRCUT]:
        print(f"--- Method: {method.value} ---")
        
        analyzer = TaintAnalyzer(method=method)
        
        # Marcar fuente tainted
        analyzer.mark_tainted("hacker_wallet", "ransomware")
        
        # Transacción: mezclar tainted con clean
        tx = TxForTaint(
            tx_id="mix_tx",
            inputs=[
                ("hacker_wallet", 5.0),  # Tainted
                ("honest_user", 5.0),    # Clean
            ],
            outputs=[
                ("exchange", 9.9),
                ("fee_addr", 0.1),
            ]
        )
        analyzer.process_transaction(tx)
        
        # Verificar taint
        result = analyzer.get_taint_score("exchange")
        print(f"  Exchange taint: {result['taint_ratio']:.1%}")
        print(f"  Risk level: {result['risk_level']}")
        print()

taint_demo()
\end{lstlisting}
%===========================================
\section{Intersection Attacks}
\label{sec:intersection-attacks}
%===========================================

\begin{definition}[Intersection Attack]
Correlacionar múltiples observaciones para reducir progresivamente el anonymity set:

\textbf{Proceso:}
\begin{enumerate}
    \item Observar anonymity set en tiempo $T_1$: $AS_1 = \{u_1, ..., u_n\}$
    \item Observar en $T_2$: $AS_2 = \{v_1, ..., v_m\}$
    \item Intersección: $AS_{1,2} = AS_1 \cap AS_2$
    \item Con suficientes observaciones: $|AS| \rightarrow 1$
\end{enumerate}
\end{definition}

\begin{figure}[htbp]
\centering
\begin{tikzpicture}
    % Sets
    \draw[thick, fill=blue!20] (0, 0) circle (1.5cm);
    \node at (0, 1.8) {$AS_1$ (T=1)};
    
    \draw[thick, fill=green!20] (2.5, 0) circle (1.3cm);
    \node at (2.5, 1.6) {$AS_2$ (T=2)};
    
    \draw[thick, fill=orange!20] (4.8, 0) circle (1.1cm);
    \node at (4.8, 1.4) {$AS_3$ (T=3)};
    
    % Intersección progresiva
    \begin{scope}
        \clip (0, 0) circle (1.5cm);
        \fill[purple!40] (2.5, 0) circle (1.3cm);
    \end{scope}
    
    % Resultado
    \node[circle, draw, fill=red!50, minimum size=0.4cm] at (1.2, 0) {};
    
    \node at (2.5, -2) [font=\small] {Intersección reduce candidatos};
    
    % Fórmula
    \node at (7, 0) [align=left, font=\small] {
        $|AS_1| = 100$\\
        $|AS_1 \cap AS_2| = 30$\\
        $|AS_1 \cap AS_2 \cap AS_3| = 5$\\
        ...\\
        $|\bigcap AS_i| = 1$
    };
\end{tikzpicture}
\caption{Intersection attack: reducción progresiva del anonymity set}
\label{fig:intersection-attack}
\end{figure}

\begin{mathnote}[Análisis Probabilístico]
\textbf{Modelo:}
\begin{itemize}
    \item $N$: Población total de usuarios
    \item $n_i$: Tamaño del anonymity set en observación $i$
    \item $k$: Número de observaciones
\end{itemize}

\textbf{Tamaño esperado de intersección:}
Si los sets son independientes:
\[
E[|\bigcap_{i=1}^{k} AS_i|] \approx N \cdot \prod_{i=1}^{k} \frac{n_i}{N}
\]

\textbf{Observaciones necesarias para deanon:}
\[
k^* \approx \frac{\log N}{\log(N/\bar{n})}
\]
donde $\bar{n}$ es el tamaño promedio del anonymity set.
\end{mathnote}

\begin{lstlisting}[language=Python, caption={Simulador de Intersection Attacks}]
from dataclasses import dataclass
from typing import List, Set, Dict
import random
import math

@dataclass
class Observation:
    time: int
    anonymity_set: Set[str]
    context: str  # Qué actividad se observó

class IntersectionAttackSimulator:
    """
    Simula intersection attacks sobre anonymity sets.
    """
    
    def __init__(self, population_size: int):
        self.population = {f"user_{i}" for i in range(population_size)}
        self.target: str = random.choice(list(self.population))
        self.observations: List[Observation] = []
    
    def generate_observation(
        self,
        time: int,
        anonymity_set_size: int,
        context: str = ""
    ) -> Observation:
        """
        Generar observación donde el target está en el anonymity set.
        """
        # El target siempre está en el set
        others = self.population - {self.target}
        sample_size = min(anonymity_set_size - 1, len(others))
        selected = random.sample(list(others), sample_size)
        
        anonymity_set = set(selected) | {self.target}
        
        obs = Observation(time, anonymity_set, context)
        self.observations.append(obs)
        
        return obs
    
    def compute_intersection(self) -> Set[str]:
        """Calcular intersección de todas las observaciones."""
        if not self.observations:
            return self.population
        
        result = self.observations[0].anonymity_set.copy()
        for obs in self.observations[1:]:
            result &= obs.anonymity_set
        
        return result
    
    def analyze_attack_progress(self) -> List[Dict]:
        """Analizar progreso del ataque con cada observación."""
        if not self.observations:
            return []
        
        results = []
        current_intersection = self.observations[0].anonymity_set.copy()
        
        for i, obs in enumerate(self.observations):
            if i > 0:
                current_intersection &= obs.anonymity_set
            
            results.append({
                'observation': i + 1,
                'individual_set_size': len(obs.anonymity_set),
                'intersection_size': len(current_intersection),
                'target_in_set': self.target in current_intersection,
                'uniquely_identified': len(current_intersection) == 1
            })
        
        return results
    
    def estimate_observations_needed(
        self,
        avg_set_size: int,
        target_set_size: int = 1
    ) -> float:
        """
        Estimar observaciones necesarias para reducir a target_set_size.
        """
        N = len(self.population)
        n = avg_set_size
        
        if n >= N:
            return float('inf')
        
        # Aproximación: cada observación reduce por factor n/N
        reduction_factor = n / N
        
        # Queremos N * (n/N)^k <= target
        # k >= log(target/N) / log(n/N)
        
        k = math.log(target_set_size / N) / math.log(reduction_factor)
        
        return max(1, k)
    
    def simulate_full_attack(
        self,
        avg_set_size: int,
        max_observations: int = 100
    ) -> Dict:
        """
        Simular ataque completo hasta identificar al target.
        """
        self.observations = []
        
        for i in range(max_observations):
            # Variar tamaño del set ligeramente
            actual_size = avg_set_size + random.randint(-5, 5)
            actual_size = max(10, min(actual_size, len(self.population)))
            
            self.generate_observation(i, actual_size, f"obs_{i}")
            
            intersection = self.compute_intersection()
            
            if len(intersection) == 1:
                return {
                    'success': True,
                    'observations_needed': i + 1,
                    'target_identified': self.target in intersection,
                    'final_set': intersection
                }
        
        final_intersection = self.compute_intersection()
        return {
            'success': False,
            'observations_made': max_observations,
            'final_set_size': len(final_intersection),
            'target_in_set': self.target in final_intersection
        }


class IntersectionDefense:
    """
    Defensas contra intersection attacks.
    """
    
    def __init__(self):
        self.dummy_activity_rate = 0.1
        self.rotation_frequency = 10  # Rotar identidad cada N actividades
    
    def add_noise_to_observations(
        self,
        real_activity: bool,
        noise_probability: float = 0.2
    ) -> bool:
        """
        Añadir actividad dummy para confundir observaciones.
        """
        if real_activity:
            return True  # Actividad real siempre ocurre
        
        # Con probabilidad noise_probability, generar actividad dummy
        return random.random() < noise_probability
    
    def calculate_effective_anonymity(
        self,
        base_anonymity_set: int,
        noise_rate: float,
        observations: int
    ) -> float:
        """
        Calcular anonymity set efectivo considerando noise.
        
        El noise hace que las intersecciones sean menos efectivas.
        """
        # Con noise, cada observación tiene falsos positivos
        # Esto aumenta el tamaño efectivo de la intersección
        
        noise_factor = 1 + noise_rate
        effective_set = base_anonymity_set * (noise_factor ** observations)
        
        return min(effective_set, base_anonymity_set * 2)  # Cap
    
    def recommend_behavior(self, current_observations: int) -> List[str]:
        """Recomendar comportamiento para minimizar intersection risk."""
        recommendations = []
        
        if current_observations > 5:
            recommendations.append("Consider rotating pseudonym/address")
        
        if current_observations > 10:
            recommendations.append("Add dummy transactions to increase noise")
        
        if current_observations > 20:
            recommendations.append("Critical: Change identity completely")
        
        recommendations.append(f"Estimated observations until deanon: ~{30 - current_observations}")
        
        return recommendations


def intersection_demo():
    print("=== Intersection Attack Simulation ===\n")
    
    simulator = IntersectionAttackSimulator(population_size=1000)
    
    print(f"Target user: {simulator.target}")
    print(f"Population size: {len(simulator.population)}")
    
    # Generar observaciones
    for i in range(15):
        simulator.generate_observation(i, anonymity_set_size=100)
    
    # Analizar progreso
    print("\n--- Attack Progress ---")
    progress = simulator.analyze_attack_progress()
    
    print(f"{'Obs':<5} {'Set Size':<10} {'Intersection':<12} {'Identified'}")
    print("-" * 40)
    
    for p in progress:
        print(f"{p['observation']:<5} {p['individual_set_size']:<10} "
              f"{p['intersection_size']:<12} {p['uniquely_identified']}")
    
    # Estimación teórica
    print("\n--- Theoretical Analysis ---")
    estimated = simulator.estimate_observations_needed(avg_set_size=100)
    print(f"Estimated observations needed: {estimated:.1f}")
    
    # Ataque completo
    print("\n--- Full Attack Simulation ---")
    simulator2 = IntersectionAttackSimulator(population_size=1000)
    result = simulator2.simulate_full_attack(avg_set_size=100)
    
    if result['success']:
        print(f"Target identified after {result['observations_needed']} observations")
    else:
        print(f"Not identified after {result['observations_made']} observations")
        print(f"Final set size: {result['final_set_size']}")

intersection_demo()
\end{lstlisting}

%===========================================
\section{Long-term Correlation}
\label{sec:long-term-correlation}
%===========================================

\begin{definition}[Long-term Correlation Attack]
Acumular datos sobre un período extendido para crear fingerprint conductual único:

\textbf{Señales utilizadas:}
\begin{itemize}
    \item Patrones de horario de actividad
    \item Distribución de montos
    \item Frecuencia de transacciones
    \item Contrapartes frecuentes
    \item Patrones de fee/gas
\end{itemize}
\end{definition}

\begin{lstlisting}[language=Python, caption={Detector de patrones para long-term correlation}]
from dataclasses import dataclass
from datetime import datetime, timedelta
from typing import Dict, List, Tuple
from collections import defaultdict
import math
import random

@dataclass
class ActivityRecord:
    timestamp: datetime
    amount: float
    counterparty: str
    fee: float
    tx_type: str  # send, receive, contract

class BehavioralFingerprint:
    """
    Crea fingerprint conductual a partir de historial de actividad.
    """
    
    def __init__(self, address: str):
        self.address = address
        self.activities: List[ActivityRecord] = []
    
    def add_activity(self, activity: ActivityRecord):
        self.activities.append(activity)
    
    def extract_features(self) -> Dict:
        """Extraer features para fingerprinting."""
        if len(self.activities) < 5:
            return {'insufficient_data': True}
        
        features = {}
        
        # 1. Distribución horaria
        hours = [a.timestamp.hour for a in self.activities]
        hour_dist = [0] * 24
        for h in hours:
            hour_dist[h] += 1
        total = sum(hour_dist)
        features['hour_distribution'] = [h/total for h in hour_dist]
        
        # 2. Distribución de días de semana
        days = [a.timestamp.weekday() for a in self.activities]
        day_dist = [0] * 7
        for d in days:
            day_dist[d] += 1
        features['day_distribution'] = [d/len(days) for d in day_dist]
        
        # 3. Estadísticas de montos
        amounts = [a.amount for a in self.activities]
        features['amount_mean'] = sum(amounts) / len(amounts)
        features['amount_std'] = math.sqrt(
            sum((a - features['amount_mean'])**2 for a in amounts) / len(amounts)
        )
        features['round_amount_ratio'] = sum(1 for a in amounts if a == round(a)) / len(amounts)
        
        # 4. Frecuencia de actividad
        sorted_acts = sorted(self.activities, key=lambda x: x.timestamp)
        if len(sorted_acts) > 1:
            intervals = [
                (sorted_acts[i+1].timestamp - sorted_acts[i].timestamp).total_seconds()
                for i in range(len(sorted_acts) - 1)
            ]
            features['avg_interval_hours'] = sum(intervals) / len(intervals) / 3600
        
        # 5. Contrapartes frecuentes
        counterparties = defaultdict(int)
        for a in self.activities:
            counterparties[a.counterparty] += 1
        features['unique_counterparties'] = len(counterparties)
        features['top_counterparty_ratio'] = max(counterparties.values()) / len(self.activities)
        
        # 6. Fee behavior
        fees = [a.fee for a in self.activities]
        features['avg_fee'] = sum(fees) / len(fees)
        
        return features
    
    def calculate_similarity(self, other: 'BehavioralFingerprint') -> float:
        """
        Calcular similitud con otro fingerprint.
        Retorna 0-1 (1 = idéntico).
        """
        f1 = self.extract_features()
        f2 = other.extract_features()
        
        if f1.get('insufficient_data') or f2.get('insufficient_data'):
            return 0.0
        
        similarities = []
        
        # Similitud de distribución horaria (cosine)
        h1 = f1['hour_distribution']
        h2 = f2['hour_distribution']
        dot = sum(a*b for a, b in zip(h1, h2))
        norm1 = math.sqrt(sum(a*a for a in h1))
        norm2 = math.sqrt(sum(b*b for b in h2))
        if norm1 > 0 and norm2 > 0:
            similarities.append(dot / (norm1 * norm2))
        
        # Similitud de estadísticas de monto
        if f1['amount_std'] > 0 and f2['amount_std'] > 0:
            mean_diff = abs(f1['amount_mean'] - f2['amount_mean']) / max(f1['amount_mean'], f2['amount_mean'])
            similarities.append(1 - min(1, mean_diff))
        
        # Similitud de round ratio
        similarities.append(1 - abs(f1['round_amount_ratio'] - f2['round_amount_ratio']))
        
        # Similitud de frecuencia
        if 'avg_interval_hours' in f1 and 'avg_interval_hours' in f2:
            int1, int2 = f1['avg_interval_hours'], f2['avg_interval_hours']
            if max(int1, int2) > 0:
                similarities.append(1 - abs(int1 - int2) / max(int1, int2))
        
        return sum(similarities) / len(similarities) if similarities else 0.0


class LongTermCorrelationDefense:
    """
    Defensas contra long-term correlation.
    """
    
    def __init__(self):
        self.behavior_randomization = True
    
    def randomize_timing(
        self,
        intended_time: datetime,
        max_delay_hours: float = 6
    ) -> datetime:
        """Añadir delay aleatorio al timing."""
        delay = random.expovariate(1 / (max_delay_hours / 2))
        delay = min(delay, max_delay_hours)
        return intended_time + timedelta(hours=delay)
    
    def randomize_amount(
        self,
        intended_amount: float,
        noise_percent: float = 5
    ) -> float:
        """Añadir ruido al monto."""
        noise = random.uniform(-noise_percent, noise_percent) / 100
        return intended_amount * (1 + noise)
    
    def recommend_countermeasures(
        self,
        fingerprint: BehavioralFingerprint
    ) -> List[str]:
        """Recomendar contramedidas basadas en fingerprint actual."""
        features = fingerprint.extract_features()
        recommendations = []
        
        if features.get('insufficient_data'):
            return ["Continue building activity history for analysis"]
        
        # Verificar patrones identificables
        hour_dist = features['hour_distribution']
        max_hour_ratio = max(hour_dist)
        if max_hour_ratio > 0.3:
            recommendations.append(
                f"High activity concentration in certain hours ({max_hour_ratio:.0%}). "
                "Spread transactions across more hours."
            )
        
        if features['round_amount_ratio'] > 0.5:
            recommendations.append(
                f"Too many round amounts ({features['round_amount_ratio']:.0%}). "
                "Use precise amounts with decimals."
            )
        
        if features['top_counterparty_ratio'] > 0.3:
            recommendations.append(
                "High concentration on few counterparties. "
                "Use intermediate addresses or mixers."
            )
        
        if features.get('avg_interval_hours', 0) < 24:
            recommendations.append(
                "Very frequent activity. Consider batching transactions."
            )
        
        return recommendations


def long_term_demo():
    print("=== Long-term Correlation Analysis ===\n")
    
    # Crear dos usuarios con patrones diferentes
    user_a = BehavioralFingerprint("user_A")
    user_b = BehavioralFingerprint("user_B")
    user_a_clone = BehavioralFingerprint("user_A_clone")  # Mismo usuario, otra dirección
    
    base = datetime(2024, 1, 1, 9, 0)
    
    # User A: Patrón de oficina, 9-5, montos redondos
    for i in range(50):
        user_a.add_activity(ActivityRecord(
            timestamp=base + timedelta(days=i, hours=random.randint(0, 8)),
            amount=round(random.uniform(0.1, 1.0), 1),
            counterparty=f"merchant_{i % 5}",
            fee=0.001,
            tx_type="send"
        ))
    
    # User A clone: Mismo patrón
    for i in range(50):
        user_a_clone.add_activity(ActivityRecord(
            timestamp=base + timedelta(days=i, hours=random.randint(0, 8)),
            amount=round(random.uniform(0.1, 1.0), 1),
            counterparty=f"other_{i % 5}",
            fee=0.001,
            tx_type="send"
        ))
    
    # User B: Patrón nocturno, montos precisos
    for i in range(50):
        user_b.add_activity(ActivityRecord(
            timestamp=base + timedelta(days=i, hours=random.randint(18, 23)),
            amount=random.uniform(0.1, 1.0),
            counterparty=f"vendor_{i % 10}",
            fee=0.002,
            tx_type="send"
        ))
    
    # Calcular similitudes
    print("--- Fingerprint Similarity ---")
    sim_aa = user_a.calculate_similarity(user_a_clone)
    sim_ab = user_a.calculate_similarity(user_b)
    
    print(f"User A vs A_clone (same person): {sim_aa:.2%}")
    print(f"User A vs User B (different): {sim_ab:.2%}")
    
    # Recomendaciones
    print("\n--- Defense Recommendations for User A ---")
    defense = LongTermCorrelationDefense()
    recs = defense.recommend_countermeasures(user_a)
    for rec in recs[:3]:
        print(f"  • {rec}")

long_term_demo()
\end{lstlisting}

%===========================================
\section*{Resumen del Bloque}
%===========================================

\begin{keyidea}
\textbf{Puntos clave de Ataques y Defensa (Bloque 19):}

\textbf{1. Deanonymization Attacks:}
\begin{itemize}
    \item IP correlation (first-seen heuristic)
    \item Timing analysis (patrones de actividad)
    \item Amount patterns (montos característicos)
\end{itemize}

\textbf{2. Sybil Attacks:}
\begin{itemize}
    \item Atacante inunda mixer con depósitos propios
    \item Reduce anonymity set efectivo
    \item Defensas: fees, staking, proof-of-personhood
\end{itemize}

\textbf{3. Eclipse Attacks:}
\begin{itemize}
    \item Aislar nodo de red honesta
    \item Defensas: peer diversity, reserved slots
\end{itemize}

\textbf{4. Dust Attacks:}
\begin{itemize}
    \item Enviar dust para tracking
    \item Defensas: coin control, freeze dust
\end{itemize}

\textbf{5. Graph Analysis:}
\begin{itemize}
    \item Common input ownership
    \item Change detection
    \item Address clustering
\end{itemize}

\textbf{6. Taint Analysis:}
\begin{itemize}
    \item Poison: todo contacto = tainted
    \item Haircut: proporcional
    \item FIFO/LIFO: por orden
\end{itemize}

\textbf{7. Intersection Attacks:}
\begin{itemize}
    \item Múltiples observaciones reducen anonymity set
    \item $k$ observaciones: $AS \rightarrow 1$
\end{itemize}

\textbf{8. Long-term Correlation:}
\begin{itemize}
    \item Behavioral fingerprinting
    \item Timing, amounts, counterparties
\end{itemize}
\end{keyidea}

\begin{bancozk}
\textbf{Defensa Integral en Banco ZK:}

\begin{center}
\begin{tikzpicture}[
    layer/.style={rectangle, draw, minimum width=9cm, minimum height=0.7cm, align=center, font=\small}
]
    \node[layer, fill=red!20] at (0, 3) {Ataques: IP, Timing, Graph, Intersection};
    \node[layer, fill=yellow!20] at (0, 2) {Detección: Monitoring, Anomaly Detection};
    \node[layer, fill=green!20] at (0, 1) {Mitigación: Tor, Delays, Coin Control};
    \node[layer, fill=blue!20] at (0, 0) {Prevención: ZK Proofs, Encrypted State};
\end{tikzpicture}
\end{center}

\textbf{Implementación en capas:}
\begin{enumerate}
    \item \textbf{Network:} Tor obligatorio, IP rotation
    \item \textbf{Timing:} Random delays, cover traffic
    \item \textbf{Amounts:} Non-round amounts, splitting
    \item \textbf{Graph:} ZK proofs eliminan linking
    \item \textbf{Behavioral:} Address rotation, identity management
\end{enumerate}

\textbf{Métricas de seguridad:}
\begin{itemize}
    \item Anonymity set efectivo $> 100$
    \item IP exposure $= 0$
    \item Timing regularity $< 0.3$
    \item Graph linkability $< 5\%$
\end{itemize}
\end{bancozk}

\cleardoublepage



% ============================================
% CONTENIDO - BLOQUE 20
% ============================================

\part{Herramientas y Ecosistema}

% ============================================
% BLOQUE 20: HERRAMIENTAS Y ECOSISTEMA
% ============================================

\chapter{Herramientas y Ecosistema}
\label{ch:bloque20}

\begin{center}
\textit{Bloque 20 de 20 --- Bloque Final}
\end{center}

\section*{Introducción}

Este bloque final presenta las herramientas prácticas y el ecosistema necesario para implementar sistemas de privacidad en blockchain. Cubrimos desde exploradores privacy-aware hasta infraestructura de proving, proporcionando una guía completa para construir y operar aplicaciones privadas.

\begin{keyidea}
\textbf{Taxonomía del Ecosistema de Privacidad:}

\begin{center}
\begin{tikzpicture}[
    layer/.style={rectangle, draw, minimum width=10cm, minimum height=0.8cm, align=center, font=\small},
    component/.style={rectangle, draw, rounded corners, minimum width=2.2cm, minimum height=0.6cm, align=center, font=\scriptsize}
]
    % Capas del stack
    \node[layer, fill=blue!20] (user) at (0, 4) {\textbf{User Layer}};
    \node[layer, fill=green!20] (app) at (0, 3) {\textbf{Application Layer}};
    \node[layer, fill=yellow!20] (infra) at (0, 2) {\textbf{Infrastructure Layer}};
    \node[layer, fill=orange!20] (network) at (0, 1) {\textbf{Network Layer}};
    \node[layer, fill=red!20] (protocol) at (0, 0) {\textbf{Protocol Layer}};
    
    % Componentes a la derecha
    \node[component] at (6.5, 4) {Wallets, UIs};
    \node[component] at (6.5, 3) {DApps, DEXs};
    \node[component] at (6.5, 2) {Provers, Relayers};
    \node[component] at (6.5, 1) {RPCs, Explorers};
    \node[component] at (6.5, 0) {ZK Circuits, Contracts};
    
    \draw[->, thick] (5, 4) -- (5.8, 4);
    \draw[->, thick] (5, 3) -- (5.8, 3);
    \draw[->, thick] (5, 2) -- (5.8, 2);
    \draw[->, thick] (5, 1) -- (5.8, 1);
    \draw[->, thick] (5, 0) -- (5.8, 0);
\end{tikzpicture}
\end{center}

\textbf{Principio fundamental:} La privacidad end-to-end requiere que cada capa del stack preserve la privacidad del usuario.
\end{keyidea}

%===========================================
\section{Exploradores Privacy-Aware}
\label{sec:privacy-explorers}
%===========================================

\begin{definition}[Privacy-Aware Explorer]
Un explorador de blockchain diseñado para permitir consultas sin comprometer la privacidad del usuario que consulta. Características:
\begin{itemize}
    \item \textbf{No tracking:} Sin cookies, analytics, fingerprinting
    \item \textbf{Acceso anónimo:} Soporte Tor (.onion), no requiere cuenta
    \item \textbf{Logging mínimo:} No guardar IPs ni queries
    \item \textbf{Open source:} Código auditable
\end{itemize}
\end{definition}

\subsection{Arquitectura de Explorador Privado}

\begin{figure}[htbp]
\centering
\begin{tikzpicture}[
    node/.style={rectangle, draw, rounded corners, minimum width=2cm, minimum height=0.7cm, align=center, font=\small},
    db/.style={cylinder, draw, minimum width=1.5cm, minimum height=1cm, shape border rotate=90, aspect=0.3}
]
    % Usuario
    \node[node, fill=blue!20] (user) at (0, 0) {Usuario};
    \node[node, fill=purple!20] (tor) at (2.5, 0) {Tor};
    \node[node, fill=green!20] (frontend) at (5, 0) {Frontend\\(static)};
    \node[node, fill=yellow!20] (api) at (7.5, 0) {API\\(no-log)};
    \node[db, fill=orange!20] (db) at (10, 0) {Index\\DB};
    \node[node, fill=red!20] (fullnode) at (10, -1.5) {Full Node};
    
    \draw[->, thick] (user) -- (tor);
    \draw[->, thick] (tor) -- (frontend);
    \draw[->, thick] (frontend) -- (api);
    \draw[->, thick] (api) -- (db);
    \draw[->, thick] (db) -- (fullnode);
    
    % Notas
    \node[font=\tiny, align=center] at (2.5, -0.8) {Anonimiza\\IP};
    \node[font=\tiny, align=center] at (7.5, -0.8) {Sin logs\\de queries};
\end{tikzpicture}
\caption{Arquitectura de explorador privacy-aware}
\end{figure}

\subsection{Análisis Comparativo de Exploradores}

\begin{table}[htbp]
\centering
\small
\begin{tabular}{lccccc}
\toprule
\textbf{Explorador} & \textbf{Tor} & \textbf{No Track} & \textbf{Open Source} & \textbf{Self-Host} & \textbf{Score} \\
\midrule
Blockchair & Parcial & No & No & No & 1/4 \\
Etherscan & No & No & No & No & 0/4 \\
Blockscout & No & Sí & Sí & Sí & 3/4 \\
Self-hosted & Sí & Sí & Sí & Sí & 4/4 \\
\bottomrule
\end{tabular}
\caption{Comparativa de exploradores por privacidad}
\end{table}

\subsection{Implementación de Explorador Privado}

\begin{lstlisting}[style=python, caption={Explorador de blockchain privacy-aware}]
"""
Privacy-Aware Blockchain Explorer
Diseñado para preservar privacidad de queries
"""

from dataclasses import dataclass, field
from typing import Optional, Dict, List, Set, Any
from datetime import datetime, timedelta
from abc import ABC, abstractmethod
import hashlib
import json
from enum import Enum

class PrivacyLevel(Enum):
    """Niveles de privacidad del explorador"""
    MINIMAL = "minimal"      # Solo funcionalidad basica
    STANDARD = "standard"    # Balance privacidad/features
    MAXIMUM = "maximum"      # Maxima privacidad

@dataclass
class QueryConfig:
    """Configuracion de privacidad para queries"""
    privacy_level: PrivacyLevel = PrivacyLevel.STANDARD
    use_tor: bool = True
    cache_locally: bool = True          # Cache en cliente, no servidor
    batch_queries: bool = True          # Agrupar queries
    add_dummy_queries: bool = True      # Queries falsas para confusion
    delay_random: bool = True           # Delays aleatorios
    
    def get_dummy_count(self) -> int:
        """Numero de queries dummy segun nivel"""
        if self.privacy_level == PrivacyLevel.MAXIMUM:
            return 5
        elif self.privacy_level == PrivacyLevel.STANDARD:
            return 2
        return 0

@dataclass
class ExplorerQuery:
    """Query individual al explorador"""
    query_type: str          # 'address', 'tx', 'block'
    target: str              # Hash o direccion
    timestamp: datetime = field(default_factory=datetime.now)
    is_dummy: bool = False   # Es query de confusion?
    
    def anonymize(self) -> 'ExplorerQuery':
        """Crear version anonimizada (sin timestamp exacto)"""
        # Redondear timestamp a hora
        rounded = self.timestamp.replace(
            minute=0, second=0, microsecond=0
        )
        return ExplorerQuery(
            query_type=self.query_type,
            target=self.target,
            timestamp=rounded,
            is_dummy=self.is_dummy
        )

@dataclass
class CachedResult:
    """Resultado cacheado localmente"""
    query: ExplorerQuery
    result: Dict[str, Any]
    cached_at: datetime
    expires_at: datetime
    
    def is_expired(self) -> bool:
        return datetime.now() > self.expires_at

class PrivacyExplorer:
    """
    Explorador de blockchain con privacidad mejorada.
    
    Estrategias de privacidad:
    1. Queries batch para ocultar patron individual
    2. Dummy queries para confusion
    3. Cache local para reducir queries
    4. Delays aleatorios para evitar timing analysis
    """
    
    def __init__(self, config: QueryConfig):
        self.config = config
        self.local_cache: Dict[str, CachedResult] = {}
        self.pending_queries: List[ExplorerQuery] = []
        self.query_history: List[ExplorerQuery] = []  # Solo local
        
    def query_address(self, address: str) -> Dict[str, Any]:
        """Query de direccion con privacidad"""
        query = ExplorerQuery(
            query_type="address",
            target=address
        )
        return self._execute_query(query)
    
    def query_transaction(self, tx_hash: str) -> Dict[str, Any]:
        """Query de transaccion con privacidad"""
        query = ExplorerQuery(
            query_type="tx",
            target=tx_hash
        )
        return self._execute_query(query)
    
    def query_block(self, block_id: str) -> Dict[str, Any]:
        """Query de bloque con privacidad"""
        query = ExplorerQuery(
            query_type="block",
            target=block_id
        )
        return self._execute_query(query)
    
    def _execute_query(self, query: ExplorerQuery) -> Dict[str, Any]:
        """Ejecutar query con protecciones de privacidad"""
        # 1. Verificar cache local
        cache_key = f"{query.query_type}:{query.target}"
        if self.config.cache_locally and cache_key in self.local_cache:
            cached = self.local_cache[cache_key]
            if not cached.is_expired():
                return cached.result
        
        # 2. Preparar batch con dummies si corresponde
        batch = self._prepare_batch(query)
        
        # 3. Ejecutar batch
        results = self._execute_batch(batch)
        
        # 4. Cachear resultado real
        if query.target in results:
            self._cache_result(query, results[query.target])
            return results[query.target]
        
        return {}
    
    def _prepare_batch(self, real_query: ExplorerQuery) -> List[ExplorerQuery]:
        """Preparar batch con queries reales y dummy"""
        batch = [real_query]
        
        if self.config.add_dummy_queries:
            dummy_count = self.config.get_dummy_count()
            for _ in range(dummy_count):
                dummy = self._generate_dummy_query(real_query.query_type)
                batch.append(dummy)
        
        # Mezclar orden
        import random
        random.shuffle(batch)
        
        return batch
    
    def _generate_dummy_query(self, query_type: str) -> ExplorerQuery:
        """Generar query dummy realista"""
        import random
        
        if query_type == "address":
            # Direccion aleatoria
            dummy_target = "0x" + "".join(
                random.choices("0123456789abcdef", k=40)
            )
        elif query_type == "tx":
            # Hash aleatorio
            dummy_target = "0x" + "".join(
                random.choices("0123456789abcdef", k=64)
            )
        else:
            # Bloque aleatorio reciente
            dummy_target = str(random.randint(1000000, 2000000))
        
        return ExplorerQuery(
            query_type=query_type,
            target=dummy_target,
            is_dummy=True
        )
    
    def _execute_batch(
        self, 
        batch: List[ExplorerQuery]
    ) -> Dict[str, Dict[str, Any]]:
        """Ejecutar batch de queries"""
        import random
        import time
        
        results = {}
        
        for query in batch:
            # Delay aleatorio entre queries
            if self.config.delay_random:
                delay = random.uniform(0.1, 0.5)
                time.sleep(delay)
            
            # Simular llamada a API
            result = self._mock_api_call(query)
            if not query.is_dummy:
                results[query.target] = result
            
            # Guardar en historial local (solo queries reales)
            if not query.is_dummy:
                self.query_history.append(query.anonymize())
        
        return results
    
    def _mock_api_call(self, query: ExplorerQuery) -> Dict[str, Any]:
        """Simular llamada a API (reemplazar con implementacion real)"""
        # En produccion, esto llamaria al backend real
        return {
            "target": query.target,
            "type": query.query_type,
            "found": True,
            "data": {"balance": "1.5 ETH"}  # Datos simulados
        }
    
    def _cache_result(self, query: ExplorerQuery, result: Dict[str, Any]):
        """Cachear resultado localmente"""
        cache_key = f"{query.query_type}:{query.target}"
        
        # TTL segun tipo de query
        if query.query_type == "block":
            ttl = timedelta(hours=24)  # Bloques no cambian
        elif query.query_type == "tx":
            ttl = timedelta(hours=1)   # TXs confirmadas no cambian
        else:
            ttl = timedelta(minutes=5) # Balances cambian
        
        cached = CachedResult(
            query=query,
            result=result,
            cached_at=datetime.now(),
            expires_at=datetime.now() + ttl
        )
        
        self.local_cache[cache_key] = cached
    
    def get_privacy_stats(self) -> Dict[str, Any]:
        """Estadisticas de privacidad de la sesion"""
        real_queries = [q for q in self.query_history if not q.is_dummy]
        
        return {
            "total_queries": len(self.query_history),
            "real_queries": len(real_queries),
            "cache_hits": len(self.local_cache),
            "privacy_level": self.config.privacy_level.value,
            "tor_enabled": self.config.use_tor,
            "dummy_ratio": 1 - (len(real_queries) / max(1, len(self.query_history)))
        }
    
    def clear_local_data(self):
        """Limpiar todos los datos locales"""
        self.local_cache.clear()
        self.query_history.clear()
        self.pending_queries.clear()


class ExplorerBackend:
    """
    Backend de explorador sin logging de usuarios.
    
    Politicas:
    - No guardar IPs
    - No guardar queries individuales
    - Solo estadisticas agregadas
    """
    
    def __init__(self, full_node_url: str):
        self.full_node_url = full_node_url
        self.aggregate_stats = {
            "total_queries": 0,
            "queries_by_type": {"address": 0, "tx": 0, "block": 0}
        }
    
    def handle_query(
        self, 
        query_type: str, 
        target: str,
        client_ip: str  # NO lo usamos
    ) -> Dict[str, Any]:
        """
        Manejar query sin logging de cliente.
        
        IMPORTANTE: client_ip se ignora intencionalmente.
        """
        # NO logueamos: client_ip, query especifica, timestamp exacto
        
        # Solo actualizamos estadisticas agregadas
        self.aggregate_stats["total_queries"] += 1
        self.aggregate_stats["queries_by_type"][query_type] += 1
        
        # Obtener datos del full node
        result = self._query_full_node(query_type, target)
        
        return result
    
    def _query_full_node(
        self, 
        query_type: str, 
        target: str
    ) -> Dict[str, Any]:
        """Query al full node (sin logging)"""
        # Implementacion real llamaria al nodo
        return {"type": query_type, "target": target, "data": {}}
    
    def get_aggregate_stats(self) -> Dict[str, Any]:
        """Solo estadisticas agregadas, sin datos individuales"""
        return self.aggregate_stats.copy()


# Demo
def privacy_explorer_demo():
    print("=== Privacy-Aware Explorer Demo ===\n")
    
    # Configuracion maxima privacidad
    config = QueryConfig(
        privacy_level=PrivacyLevel.MAXIMUM,
        use_tor=True,
        add_dummy_queries=True,
        delay_random=True
    )
    
    explorer = PrivacyExplorer(config)
    
    # Queries de ejemplo
    print("Ejecutando queries con proteccion de privacidad...")
    
    result1 = explorer.query_address(
        "0x1234567890abcdef1234567890abcdef12345678"
    )
    print(f"  Query 1: Address lookup")
    
    result2 = explorer.query_transaction(
        "0xabcdef1234567890abcdef1234567890abcdef1234567890abcdef1234567890ab"
    )
    print(f"  Query 2: Transaction lookup")
    
    # Segunda query a misma address (cache hit)
    result3 = explorer.query_address(
        "0x1234567890abcdef1234567890abcdef12345678"
    )
    print(f"  Query 3: Address lookup (cached)")
    
    # Estadisticas
    stats = explorer.get_privacy_stats()
    print(f"\nEstadisticas de privacidad:")
    print(f"  Nivel: {stats['privacy_level']}")
    print(f"  Queries reales: {stats['real_queries']}")
    print(f"  Cache hits: {stats['cache_hits']}")
    print(f"  Dummy ratio: {stats['dummy_ratio']:.1%}")
    print(f"  Tor enabled: {stats['tor_enabled']}")

privacy_explorer_demo()
\end{lstlisting}

%===========================================
\section{Wallets con Soporte ZK}
\label{sec:zk-wallets}
%===========================================

\begin{definition}[ZK-Enabled Wallet]
Wallet que soporta operaciones de privacidad basadas en zero-knowledge proofs:
\begin{itemize}
    \item \textbf{Shielded balances:} Fondos ocultos criptográficamente
    \item \textbf{Private transfers:} Transferencias sin revelar montos/destinos
    \item \textbf{Proof generation:} Capacidad de generar ZK proofs
    \item \textbf{Note management:} Gestión de notas/commitments privados
\end{itemize}
\end{definition}

\subsection{Comparativa de Wallets ZK}

\begin{figure}[htbp]
\centering
\begin{tikzpicture}[
    wallet/.style={rectangle, draw, rounded corners, minimum width=2.5cm, minimum height=1.2cm, align=center, font=\small},
    feature/.style={rectangle, draw, minimum width=1.8cm, minimum height=0.5cm, align=center, font=\tiny}
]
    % Wallets
    \node[wallet, fill=purple!20] (railgun) at (0, 0) {\textbf{Railgun}\\ERC-20 Privacy};
    \node[wallet, fill=blue!20] (umbra) at (4, 0) {\textbf{Umbra}\\Stealth Payments};
    \node[wallet, fill=green!20] (zcash) at (8, 0) {\textbf{Zcash}\\Native Privacy};
    
    % Features Railgun
    \node[feature, fill=purple!10] at (0, -1.2) {Any ERC-20};
    \node[feature, fill=purple!10] at (0, -1.8) {DeFi Integration};
    \node[feature, fill=purple!10] at (0, -2.4) {Relayers};
    
    % Features Umbra
    \node[feature, fill=blue!10] at (4, -1.2) {Stealth Addr};
    \node[feature, fill=blue!10] at (4, -1.8) {Simple UX};
    \node[feature, fill=blue!10] at (4, -2.4) {No ZK Proofs};
    
    % Features Zcash
    \node[feature, fill=green!10] at (8, -1.2) {Shielded Pool};
    \node[feature, fill=green!10] at (8, -1.8) {Sapling/Orchard};
    \node[feature, fill=green!10] at (8, -2.4) {Native Token};
\end{tikzpicture}
\caption{Ecosistema de wallets con soporte de privacidad}
\end{figure}

\begin{table}[htbp]
\centering
\small
\begin{tabular}{lcccccc}
\toprule
\textbf{Wallet} & \textbf{Chain} & \textbf{ZK} & \textbf{Mobile} & \textbf{DeFi} & \textbf{AS} & \textbf{Complejidad} \\
\midrule
Railgun & EVM & Sí & Parcial & Sí & Alto & Alta \\
Umbra & EVM & No & Sí & No & Bajo & Baja \\
Zashi & Zcash & Sí & Sí & No & Alto & Media \\
YWallet & Zcash & Sí & Sí & No & Alto & Media \\
Nighthawk & Zcash & Sí & Sí & No & Alto & Media \\
\bottomrule
\end{tabular}
\caption{Comparativa de wallets ZK (AS = Anonymity Set)}
\end{table}

\subsection{Arquitectura de Wallet ZK}

\begin{lstlisting}[style=python, caption={Arquitectura de wallet con soporte ZK}]
"""
ZK-Enabled Wallet Architecture
Soporte para operaciones privadas con zero-knowledge proofs
"""

from dataclasses import dataclass, field
from typing import Optional, Dict, List, Tuple, Any
from datetime import datetime
from abc import ABC, abstractmethod
from enum import Enum
import hashlib
import secrets
import json

class WalletType(Enum):
    """Tipos de wallet soportados"""
    RAILGUN = "railgun"
    UMBRA = "umbra"
    ZCASH = "zcash"
    GENERIC_ZK = "generic_zk"

@dataclass
class PrivateNote:
    """
    Representacion de un deposito privado.
    
    Una Note contiene:
    - commitment: hash publico en el contrato
    - nullifier_secret: secreto para gastar
    - amount: cantidad (encriptada on-chain)
    - asset: token/moneda
    """
    commitment: str
    nullifier_secret: bytes
    amount: int
    asset: str
    created_at: datetime
    spent: bool = False
    spending_tx: Optional[str] = None
    
    def compute_nullifier(self) -> str:
        """Calcular nullifier para gastar la note"""
        data = self.nullifier_secret + self.commitment.encode()
        return hashlib.sha256(data).hexdigest()
    
    def to_encrypted(self, encryption_key: bytes) -> bytes:
        """Serializar note encriptada para backup"""
        # Simplificado - en produccion usar AES-GCM
        note_data = json.dumps({
            "commitment": self.commitment,
            "nullifier_secret": self.nullifier_secret.hex(),
            "amount": self.amount,
            "asset": self.asset,
            "created_at": self.created_at.isoformat(),
            "spent": self.spent
        }).encode()
        
        # XOR simple para demo (usar crypto real en produccion)
        key_extended = (encryption_key * (len(note_data) // len(encryption_key) + 1))
        encrypted = bytes(a ^ b for a, b in zip(note_data, key_extended))
        
        return encrypted
    
    @classmethod
    def from_encrypted(
        cls, 
        encrypted: bytes, 
        encryption_key: bytes
    ) -> 'PrivateNote':
        """Deserializar note desde backup encriptado"""
        key_extended = (encryption_key * (len(encrypted) // len(encryption_key) + 1))
        decrypted = bytes(a ^ b for a, b in zip(encrypted, key_extended))
        
        data = json.loads(decrypted.decode())
        return cls(
            commitment=data["commitment"],
            nullifier_secret=bytes.fromhex(data["nullifier_secret"]),
            amount=data["amount"],
            asset=data["asset"],
            created_at=datetime.fromisoformat(data["created_at"]),
            spent=data["spent"]
        )

@dataclass
class WalletState:
    """Estado completo del wallet"""
    notes: List[PrivateNote] = field(default_factory=list)
    spent_nullifiers: set = field(default_factory=set)
    pending_deposits: List[str] = field(default_factory=list)
    pending_withdrawals: List[str] = field(default_factory=list)
    last_sync_block: int = 0
    
    def get_balance(self, asset: str) -> int:
        """Balance disponible de un asset"""
        return sum(
            note.amount 
            for note in self.notes 
            if note.asset == asset and not note.spent
        )
    
    def get_all_balances(self) -> Dict[str, int]:
        """Todos los balances por asset"""
        balances: Dict[str, int] = {}
        for note in self.notes:
            if not note.spent:
                balances[note.asset] = balances.get(note.asset, 0) + note.amount
        return balances
    
    def get_spendable_notes(
        self, 
        asset: str, 
        amount: int
    ) -> List[PrivateNote]:
        """Seleccionar notes para gastar un monto"""
        available = [
            n for n in self.notes 
            if n.asset == asset and not n.spent
        ]
        available.sort(key=lambda n: n.amount, reverse=True)
        
        selected = []
        total = 0
        for note in available:
            selected.append(note)
            total += note.amount
            if total >= amount:
                break
        
        if total < amount:
            raise ValueError(f"Insufficient balance: {total} < {amount}")
        
        return selected


class ZKProver(ABC):
    """Interfaz abstracta para generacion de pruebas"""
    
    @abstractmethod
    def generate_deposit_proof(
        self,
        amount: int,
        commitment: str,
        blinding: bytes
    ) -> bytes:
        """Generar prueba de deposito"""
        pass
    
    @abstractmethod
    def generate_transfer_proof(
        self,
        input_notes: List[PrivateNote],
        output_amounts: List[int],
        output_recipients: List[str]
    ) -> bytes:
        """Generar prueba de transferencia"""
        pass
    
    @abstractmethod
    def generate_withdrawal_proof(
        self,
        note: PrivateNote,
        recipient: str
    ) -> bytes:
        """Generar prueba de retiro"""
        pass


class ClientSideProver(ZKProver):
    """Prover que genera pruebas en el dispositivo del usuario"""
    
    def __init__(self, circuit_path: str):
        self.circuit_path = circuit_path
        self.proving_key: Optional[bytes] = None
        
    def load_proving_key(self, key_path: str):
        """Cargar proving key (puede ser pesado)"""
        # En produccion, cargar archivo real
        self.proving_key = b"mock_proving_key"
    
    def generate_deposit_proof(
        self,
        amount: int,
        commitment: str,
        blinding: bytes
    ) -> bytes:
        """Generar prueba localmente"""
        # Simular generacion de prueba
        witness = {
            "amount": amount,
            "commitment": commitment,
            "blinding": blinding.hex()
        }
        
        # En produccion: llamar a snarkjs/circom/etc
        proof_data = hashlib.sha256(
            json.dumps(witness).encode()
        ).digest()
        
        return proof_data
    
    def generate_transfer_proof(
        self,
        input_notes: List[PrivateNote],
        output_amounts: List[int],
        output_recipients: List[str]
    ) -> bytes:
        """Generar prueba de transferencia localmente"""
        witness = {
            "inputs": [n.commitment for n in input_notes],
            "outputs": list(zip(output_amounts, output_recipients))
        }
        
        return hashlib.sha256(json.dumps(witness).encode()).digest()
    
    def generate_withdrawal_proof(
        self,
        note: PrivateNote,
        recipient: str
    ) -> bytes:
        """Generar prueba de retiro localmente"""
        witness = {
            "nullifier": note.compute_nullifier(),
            "recipient": recipient,
            "amount": note.amount
        }
        
        return hashlib.sha256(json.dumps(witness).encode()).digest()


class ZKWallet:
    """
    Wallet con soporte completo de ZK.
    
    Operaciones:
    - deposit(): Shield fondos (transparente -> privado)
    - transfer(): Transferencia privada
    - withdraw(): Unshield fondos (privado -> transparente)
    """
    
    def __init__(
        self,
        wallet_type: WalletType,
        prover: ZKProver,
        encryption_key: Optional[bytes] = None
    ):
        self.wallet_type = wallet_type
        self.prover = prover
        self.encryption_key = encryption_key or secrets.token_bytes(32)
        self.state = WalletState()
        
    def deposit(
        self,
        amount: int,
        asset: str
    ) -> Tuple[str, bytes]:
        """
        Depositar fondos en pool privado.
        
        Returns:
            (commitment, proof) para enviar a contrato
        """
        # Generar secretos
        nullifier_secret = secrets.token_bytes(32)
        blinding = secrets.token_bytes(32)
        
        # Calcular commitment
        preimage = (
            amount.to_bytes(32, 'big') +
            asset.encode().ljust(32, b'\x00') +
            nullifier_secret +
            blinding
        )
        commitment = hashlib.sha256(preimage).hexdigest()
        
        # Generar prueba
        proof = self.prover.generate_deposit_proof(
            amount=amount,
            commitment=commitment,
            blinding=blinding
        )
        
        # Crear note (guardada localmente)
        note = PrivateNote(
            commitment=commitment,
            nullifier_secret=nullifier_secret,
            amount=amount,
            asset=asset,
            created_at=datetime.now()
        )
        
        self.state.notes.append(note)
        self.state.pending_deposits.append(commitment)
        
        return commitment, proof
    
    def transfer(
        self,
        amount: int,
        asset: str,
        recipient_viewing_key: str
    ) -> Tuple[str, bytes]:
        """
        Transferencia privada.
        
        Args:
            amount: Monto a transferir
            asset: Token/moneda
            recipient_viewing_key: Clave publica del destinatario
            
        Returns:
            (output_commitment, proof)
        """
        # Seleccionar notes a gastar
        input_notes = self.state.get_spendable_notes(asset, amount)
        total_input = sum(n.amount for n in input_notes)
        change = total_input - amount
        
        # Outputs: [amount para recipient, change para self]
        output_amounts = [amount]
        output_recipients = [recipient_viewing_key]
        
        if change > 0:
            output_amounts.append(change)
            output_recipients.append("self")  # Nuestra propia clave
        
        # Generar prueba
        proof = self.prover.generate_transfer_proof(
            input_notes=input_notes,
            output_amounts=output_amounts,
            output_recipients=output_recipients
        )
        
        # Marcar inputs como gastados
        for note in input_notes:
            note.spent = True
            self.state.spent_nullifiers.add(note.compute_nullifier())
        
        # Crear note de cambio si hay
        if change > 0:
            change_note = PrivateNote(
                commitment=hashlib.sha256(
                    secrets.token_bytes(32)
                ).hexdigest(),
                nullifier_secret=secrets.token_bytes(32),
                amount=change,
                asset=asset,
                created_at=datetime.now()
            )
            self.state.notes.append(change_note)
        
        # Output commitment para recipient (simplificado)
        output_commitment = hashlib.sha256(
            f"{amount}:{recipient_viewing_key}".encode()
        ).hexdigest()
        
        return output_commitment, proof
    
    def withdraw(
        self,
        amount: int,
        asset: str,
        recipient_address: str
    ) -> Tuple[str, bytes]:
        """
        Retirar fondos a direccion transparente.
        
        Returns:
            (nullifier, proof)
        """
        # Seleccionar note exacta o con suficiente balance
        notes = self.state.get_spendable_notes(asset, amount)
        
        if len(notes) == 1 and notes[0].amount == amount:
            # Retiro exacto
            note = notes[0]
        else:
            # Necesitamos consolidar primero (simplificado)
            raise ValueError("Complex withdrawal not implemented")
        
        # Generar prueba
        proof = self.prover.generate_withdrawal_proof(
            note=note,
            recipient=recipient_address
        )
        
        # Marcar como gastada
        nullifier = note.compute_nullifier()
        note.spent = True
        note.spending_tx = "pending"
        self.state.spent_nullifiers.add(nullifier)
        
        return nullifier, proof
    
    def backup(self) -> bytes:
        """Crear backup encriptado del wallet"""
        backup_data = {
            "wallet_type": self.wallet_type.value,
            "notes": [
                note.to_encrypted(self.encryption_key).hex()
                for note in self.state.notes
            ],
            "last_sync": self.state.last_sync_block
        }
        
        # Encriptar todo el backup
        backup_json = json.dumps(backup_data).encode()
        encrypted = bytes(
            a ^ b 
            for a, b in zip(
                backup_json, 
                self.encryption_key * (len(backup_json) // 32 + 1)
            )
        )
        
        return encrypted
    
    def get_status(self) -> Dict[str, Any]:
        """Estado actual del wallet"""
        return {
            "wallet_type": self.wallet_type.value,
            "balances": self.state.get_all_balances(),
            "total_notes": len(self.state.notes),
            "spent_notes": len(self.state.spent_nullifiers),
            "pending_deposits": len(self.state.pending_deposits),
            "pending_withdrawals": len(self.state.pending_withdrawals),
            "last_sync_block": self.state.last_sync_block
        }


# Demo
def zk_wallet_demo():
    print("=== ZK Wallet Demo ===\n")
    
    # Crear prover y wallet
    prover = ClientSideProver("/path/to/circuits")
    wallet = ZKWallet(
        wallet_type=WalletType.GENERIC_ZK,
        prover=prover
    )
    
    # 1. Depositar
    print("1. Depositando 100 ETH en pool privado...")
    commitment, proof = wallet.deposit(
        amount=100_000_000_000,  # 100 ETH en wei simplificado
        asset="ETH"
    )
    print(f"   Commitment: {commitment[:20]}...")
    print(f"   Proof size: {len(proof)} bytes")
    
    # 2. Segundo deposito
    print("\n2. Depositando 50 ETH adicionales...")
    wallet.deposit(amount=50_000_000_000, asset="ETH")
    
    # 3. Estado
    status = wallet.get_status()
    print(f"\n3. Estado del wallet:")
    print(f"   Balance ETH: {status['balances'].get('ETH', 0) / 1e9:.0f}")
    print(f"   Notes totales: {status['total_notes']}")
    
    # 4. Transferencia privada
    print("\n4. Transfiriendo 30 ETH a otro usuario...")
    recipient_key = "0xrecipient_viewing_key_abc123"
    output, proof = wallet.transfer(
        amount=30_000_000_000,
        asset="ETH",
        recipient_viewing_key=recipient_key
    )
    print(f"   Output commitment: {output[:20]}...")
    
    # 5. Estado actualizado
    status = wallet.get_status()
    print(f"\n5. Balance despues de transferencia:")
    print(f"   Balance ETH: {status['balances'].get('ETH', 0) / 1e9:.0f}")
    print(f"   Notes gastadas: {status['spent_notes']}")
    
    # 6. Backup
    print("\n6. Creando backup encriptado...")
    backup = wallet.backup()
    print(f"   Backup size: {len(backup)} bytes")

zk_wallet_demo()
\end{lstlisting}

\subsection{Railgun: Privacy para ERC-20}

\begin{keyidea}
\textbf{Arquitectura Railgun:}

\begin{center}
\begin{tikzpicture}[
    box/.style={rectangle, draw, rounded corners, minimum width=2cm, minimum height=0.8cm, align=center, font=\small}
]
    % Flujo
    \node[box, fill=blue!20] (user) at (0, 0) {Usuario};
    \node[box, fill=green!20] (shield) at (3, 0) {Shield\\(Deposit)};
    \node[box, fill=purple!20] (pool) at (6, 0) {Shielded\\Pool};
    \node[box, fill=orange!20] (transfer) at (9, 0) {Private\\Transfer};
    \node[box, fill=red!20] (unshield) at (12, 0) {Unshield\\(Withdraw)};
    
    \draw[->, thick] (user) -- (shield);
    \draw[->, thick] (shield) -- (pool);
    \draw[->, thick] (pool) -- (transfer);
    \draw[->, thick] (transfer) -- (pool);
    \draw[->, thick] (pool) -- (unshield);
    
    % Labels
    \node[font=\tiny, below] at (3, -0.6) {ERC-20 $\rightarrow$ Note};
    \node[font=\tiny, below] at (9, -0.6) {Note $\rightarrow$ Note};
    \node[font=\tiny, below] at (12, -0.6) {Note $\rightarrow$ ERC-20};
\end{tikzpicture}
\end{center}

\textbf{Características clave:}
\begin{itemize}
    \item Cualquier token ERC-20 puede ser ``shielded''
    \item Integración con DeFi (swaps, lending) sin revelar identidad
    \item Red de relayers para privacidad de gas
    \item Proving del lado del cliente (WASM)
\end{itemize}
\end{keyidea}

\subsection{Umbra: Stealth Payments Simplificados}

\begin{definition}[Stealth Address]
Dirección única generada para cada transacción usando:
\begin{enumerate}
    \item Sender genera secreto efímero $r$
    \item Publica $R = r \cdot G$ (punto en curva elíptica)
    \item Stealth address: $P_{stealth} = P_{recipient} + H(r \cdot P_{recipient}) \cdot G$
    \item Solo recipient puede derivar private key correspondiente
\end{enumerate}
\end{definition}

\begin{lstlisting}[style=python, caption={Implementación simplificada de Stealth Addresses}]
"""
Stealth Address Implementation (Umbra-style)
Direcciones unicas por transaccion
"""

from dataclasses import dataclass
from typing import Tuple
import hashlib
import secrets

# Simulacion de operaciones de curva eliptica
# En produccion usar libreria como py_ecc

@dataclass
class Point:
    """Punto en curva eliptica (simplificado)"""
    x: int
    y: int
    
    def __add__(self, other: 'Point') -> 'Point':
        # Simplificado - en realidad es suma de puntos EC
        return Point(
            (self.x + other.x) % (2**256),
            (self.y + other.y) % (2**256)
        )
    
    def __mul__(self, scalar: int) -> 'Point':
        # Simplificado - multiplicacion escalar
        return Point(
            (self.x * scalar) % (2**256),
            (self.y * scalar) % (2**256)
        )
    
    def to_address(self) -> str:
        """Convertir punto a direccion Ethereum"""
        data = self.x.to_bytes(32, 'big') + self.y.to_bytes(32, 'big')
        hash_bytes = hashlib.keccak_256(data).digest()
        return "0x" + hash_bytes[-20:].hex()

# Generador de curva (simplificado)
G = Point(
    0x79BE667EF9DCBBAC55A06295CE870B07029BFCDB2DCE28D959F2815B16F81798,
    0x483ADA7726A3C4655DA4FBFC0E1108A8FD17B448A68554199C47D08FFB10D4B8
)

@dataclass
class StealthKeyPair:
    """Par de claves para stealth payments"""
    spending_key: int       # Clave privada de gasto
    viewing_key: int        # Clave privada de visualizacion
    spending_pubkey: Point  # Clave publica de gasto
    viewing_pubkey: Point   # Clave publica de visualizacion
    
    @classmethod
    def generate(cls) -> 'StealthKeyPair':
        """Generar nuevo par de claves"""
        spending_key = secrets.randbelow(2**256)
        viewing_key = secrets.randbelow(2**256)
        
        return cls(
            spending_key=spending_key,
            viewing_key=viewing_key,
            spending_pubkey=G * spending_key,
            viewing_pubkey=G * viewing_key
        )
    
    def get_meta_address(self) -> str:
        """Direccion meta (spending || viewing pubkeys)"""
        # Formato: st:eth:0x<spending_pubkey><viewing_pubkey>
        spending_hex = hex(self.spending_pubkey.x)[2:].zfill(64)
        viewing_hex = hex(self.viewing_pubkey.x)[2:].zfill(64)
        return f"st:eth:0x{spending_hex}{viewing_hex}"


class StealthAddressGenerator:
    """Generador de stealth addresses (lado sender)"""
    
    @staticmethod
    def generate_stealth_address(
        recipient_spending_pubkey: Point,
        recipient_viewing_pubkey: Point
    ) -> Tuple[str, Point, int]:
        """
        Generar stealth address para recipient.
        
        Returns:
            (stealth_address, ephemeral_pubkey, shared_secret)
        """
        # 1. Generar secreto efimero
        ephemeral_secret = secrets.randbelow(2**256)
        ephemeral_pubkey = G * ephemeral_secret
        
        # 2. Calcular shared secret: r * viewing_pubkey
        shared_point = recipient_viewing_pubkey * ephemeral_secret
        shared_secret = int(hashlib.sha256(
            shared_point.x.to_bytes(32, 'big')
        ).hexdigest(), 16)
        
        # 3. Calcular stealth address
        # P_stealth = P_spending + hash(shared_secret) * G
        stealth_point = recipient_spending_pubkey + (G * shared_secret)
        stealth_address = stealth_point.to_address()
        
        return stealth_address, ephemeral_pubkey, shared_secret


class StealthAddressScanner:
    """Scanner de stealth payments (lado recipient)"""
    
    def __init__(self, keypair: StealthKeyPair):
        self.keypair = keypair
    
    def check_payment(
        self, 
        ephemeral_pubkey: Point,
        payment_address: str
    ) -> Tuple[bool, int]:
        """
        Verificar si un pago es para nosotros.
        
        Returns:
            (is_ours, spending_key) si es nuestro
        """
        # 1. Calcular shared secret usando viewing key
        # shared_point = viewing_key * ephemeral_pubkey
        shared_point = ephemeral_pubkey * self.keypair.viewing_key
        shared_secret = int(hashlib.sha256(
            shared_point.x.to_bytes(32, 'big')
        ).hexdigest(), 16)
        
        # 2. Calcular stealth address esperada
        expected_point = self.keypair.spending_pubkey + (G * shared_secret)
        expected_address = expected_point.to_address()
        
        # 3. Comparar
        if expected_address.lower() == payment_address.lower():
            # 4. Calcular spending key para esta address
            stealth_spending_key = (
                self.keypair.spending_key + shared_secret
            ) % (2**256)
            return True, stealth_spending_key
        
        return False, 0
    
    def scan_announcements(
        self,
        announcements: list
    ) -> list:
        """
        Escanear lista de anuncios buscando pagos nuestros.
        
        Cada anuncio: (ephemeral_pubkey, payment_address, amount)
        """
        our_payments = []
        
        for ephemeral, address, amount in announcements:
            is_ours, spending_key = self.check_payment(ephemeral, address)
            if is_ours:
                our_payments.append({
                    "address": address,
                    "amount": amount,
                    "spending_key": spending_key
                })
        
        return our_payments


# Demo
def stealth_demo():
    print("=== Stealth Address Demo (Umbra-style) ===\n")
    
    # 1. Recipient genera keypair
    print("1. Recipient genera stealth keypair...")
    recipient = StealthKeyPair.generate()
    meta_address = recipient.get_meta_address()
    print(f"   Meta-address: {meta_address[:50]}...")
    
    # 2. Sender genera stealth address
    print("\n2. Sender genera stealth address para pago...")
    stealth_addr, ephemeral_pub, _ = StealthAddressGenerator.generate_stealth_address(
        recipient.spending_pubkey,
        recipient.viewing_pubkey
    )
    print(f"   Stealth address: {stealth_addr}")
    print(f"   Ephemeral pubkey: ({ephemeral_pub.x % 10000}...)")
    
    # 3. Sender envia ETH a stealth address (on-chain)
    print("\n3. Sender envia 1 ETH a stealth address")
    print("   [TX enviada on-chain]")
    print(f"   Announcement publicado: ephemeral_pubkey")
    
    # 4. Recipient escanea announcements
    print("\n4. Recipient escanea announcements...")
    scanner = StealthAddressScanner(recipient)
    
    # Simular varios announcements
    announcements = [
        (Point(123, 456), "0xdecoy1address", 0.5),
        (ephemeral_pub, stealth_addr, 1.0),  # El nuestro
        (Point(789, 012), "0xdecoy2address", 2.0),
    ]
    
    our_payments = scanner.scan_announcements(announcements)
    
    if our_payments:
        print(f"   Encontrado pago:")
        print(f"     Address: {our_payments[0]['address']}")
        print(f"     Amount: {our_payments[0]['amount']} ETH")
        print(f"     Spending key derivado: {our_payments[0]['spending_key'] % 10000}...")
    
    # 5. Recipient puede gastar
    print("\n5. Recipient puede gastar fondos con spending key derivado")

stealth_demo()
\end{lstlisting}

\subsection{Zcash Wallets: Privacidad Nativa}

\begin{figure}[htbp]
\centering
\begin{tikzpicture}[
    wallet/.style={rectangle, draw, rounded corners, minimum width=3cm, minimum height=1cm, align=center, font=\small},
    feature/.style={font=\tiny, align=left}
]
    % Wallets Zcash
    \node[wallet, fill=yellow!20] (zashi) at (0, 0) {\textbf{Zashi}\\(Oficial)};
    \node[wallet, fill=green!20] (ywallet) at (4.5, 0) {\textbf{YWallet}};
    \node[wallet, fill=purple!20] (nighthawk) at (9, 0) {\textbf{Nighthawk}};
    
    % Features
    \node[feature] at (0, -1.5) {
        \begin{tabular}{l}
        $\bullet$ Mobile (iOS/Android) \\
        $\bullet$ Orchard support \\
        $\bullet$ Official ECC \\
        $\bullet$ Simple UX
        \end{tabular}
    };
    
    \node[feature] at (4.5, -1.5) {
        \begin{tabular}{l}
        $\bullet$ Multi-coin \\
        $\bullet$ Fast sync \\
        $\bullet$ Warp sync \\
        $\bullet$ Advanced features
        \end{tabular}
    };
    
    \node[feature] at (9, -1.5) {
        \begin{tabular}{l}
        $\bullet$ Privacy-first \\
        $\bullet$ Tor built-in \\
        $\bullet$ Auto-shielding \\
        $\bullet$ Memo support
        \end{tabular}
    };
\end{tikzpicture}
\caption{Ecosistema de wallets Zcash}
\end{figure}

\begin{table}[htbp]
\centering
\small
\begin{tabular}{lccccc}
\toprule
\textbf{Característica} & \textbf{Zashi} & \textbf{YWallet} & \textbf{Nighthawk} \\
\midrule
Plataformas & iOS, Android & iOS, Android, Desktop & iOS, Android \\
Shielded pools & Sapling, Orchard & Sapling, Orchard & Sapling \\
Sync method & Linear & Warp & Linear \\
Tor integrado & No & No & Sí \\
Auto-shield & No & Configurable & Sí \\
Memos & Sí & Sí & Sí \\
Multi-cuenta & Básico & Avanzado & Básico \\
\bottomrule
\end{tabular}
\caption{Comparativa detallada de wallets Zcash}
\end{table}

\begin{bancozk}
\textbf{Integración de Wallet en Banco ZK:}

El wallet de Banco ZK combina las mejores características:

\begin{center}
\begin{tikzpicture}[
    component/.style={rectangle, draw, rounded corners, minimum width=2.5cm, minimum height=0.7cm, align=center, font=\small}
]
    \node[component, fill=blue!20] (ui) at (0, 2) {UI Layer};
    \node[component, fill=green!20] (prover) at (0, 1) {ZK Prover};
    \node[component, fill=yellow!20] (notes) at (0, 0) {Note Manager};
    \node[component, fill=orange!20] (network) at (0, -1) {Network (Tor)};
    \node[component, fill=red!20] (sync) at (0, -2) {Sync Engine};
    
    % Conexiones
    \draw[->, thick] (ui) -- (prover);
    \draw[->, thick] (prover) -- (notes);
    \draw[->, thick] (notes) -- (network);
    \draw[->, thick] (network) -- (sync);
    
    % Descripciones
    \node[font=\tiny, right] at (1.5, 2) {Interfaz de usuario};
    \node[font=\tiny, right] at (1.5, 1) {Client-side proving};
    \node[font=\tiny, right] at (1.5, 0) {Encrypted backup};
    \node[font=\tiny, right] at (1.5, -1) {Conexion anonima};
    \node[font=\tiny, right] at (1.5, -2) {Warp sync};
\end{tikzpicture}
\end{center}

\textbf{Características del wallet Banco ZK:}
\begin{enumerate}
    \item Client-side ZK proving (WASM optimizado)
    \item Backup encriptado con passphrase
    \item Tor obligatorio para conexiones
    \item Warp sync para sincronización rápida
    \item Stealth addresses para pagos entrantes
\end{enumerate}
\end{bancozk}

%===========================================
\section{RPC Privados}
\label{sec:private-rpc}
%===========================================

\begin{definition}[Private RPC]
Endpoint RPC que protege la privacidad del usuario al enviar transacciones:
\begin{itemize}
    \item \textbf{No mempool pública:} Transacciones van directo a builders
    \item \textbf{Protección MEV:} Evita frontrunning y sandwich attacks
    \item \textbf{Reverts privados:} Transacciones fallidas no se publican
    \item \textbf{Sin logging de usuarios:} No se guardan IPs ni actividad
\end{itemize}
\end{definition}

\subsection{Arquitectura de Transacción Privada}

\begin{figure}[htbp]
\centering
\begin{tikzpicture}[
    node/.style={rectangle, draw, rounded corners, minimum width=2cm, minimum height=0.7cm, align=center, font=\small},
    arrow/.style={->, thick}
]
    % Flujo tradicional (arriba)
    \node at (-4, 2) {\textbf{Tradicional:}};
    \node[node, fill=blue!20] (user1) at (0, 2) {Usuario};
    \node[node, fill=red!30] (mempool) at (3.5, 2) {Mempool\\Pública};
    \node[node, fill=orange!20] (searcher) at (7, 2) {MEV\\Searcher};
    \node[node, fill=gray!20] (block1) at (10.5, 2) {Bloque};
    
    \draw[arrow] (user1) -- (mempool);
    \draw[arrow] (mempool) -- (searcher);
    \draw[arrow] (searcher) -- (block1);
    \draw[arrow, red, dashed] (mempool) to[bend right=30] node[below, font=\tiny] {frontrun} (searcher);
    
    % Flujo privado (abajo)
    \node at (-4, 0) {\textbf{Privado:}};
    \node[node, fill=blue!20] (user2) at (0, 0) {Usuario};
    \node[node, fill=green!20] (rpc) at (3.5, 0) {Private\\RPC};
    \node[node, fill=purple!20] (builder) at (7, 0) {Block\\Builder};
    \node[node, fill=gray!20] (block2) at (10.5, 0) {Bloque};
    
    \draw[arrow] (user2) -- (rpc);
    \draw[arrow, green!60!black] (rpc) -- node[above, font=\tiny] {privado} (builder);
    \draw[arrow] (builder) -- (block2);
\end{tikzpicture}
\caption{Comparación de flujo de transacciones tradicional vs privado}
\end{figure}

\subsection{Flashbots Protect}

\begin{lstlisting}[style=python, caption={Cliente para Flashbots Protect RPC}]
"""
Flashbots Protect RPC Client
Envio de transacciones con proteccion MEV
"""

from dataclasses import dataclass, field
from typing import Optional, Dict, List, Any, Tuple
from enum import Enum
from datetime import datetime, timedelta
import hashlib
import json
import secrets

class HintPreference(Enum):
    """Preferencias de hints para builders"""
    NONE = "none"                    # Sin hints (maxima privacidad)
    CALLDATA = "calldata"            # Revelar calldata
    CONTRACT_ADDRESS = "contract"    # Revelar contrato destino
    FUNCTION_SELECTOR = "selector"   # Revelar funcion llamada
    LOGS = "logs"                    # Revelar logs
    HASH = "hash"                    # Solo hash de tx

class RefundPreference(Enum):
    """Preferencia de refund de MEV"""
    NO_REFUND = 0        # Sin refund
    PARTIAL_REFUND = 50  # 50% refund
    FULL_REFUND = 90     # 90% refund (maximo posible)

@dataclass
class ProtectTransaction:
    """Transaccion para enviar via Flashbots Protect"""
    to: str
    value: int
    data: bytes
    gas_limit: int
    max_fee_per_gas: int
    max_priority_fee: int
    nonce: int
    chain_id: int = 1
    
    # Preferencias de privacidad
    hints: List[HintPreference] = field(
        default_factory=lambda: [HintPreference.HASH]
    )
    refund_percent: int = 90
    
    # Metadata
    created_at: datetime = field(default_factory=datetime.now)
    
    def to_raw(self) -> str:
        """Serializar a formato raw (simplificado)"""
        tx_data = {
            "to": self.to,
            "value": hex(self.value),
            "data": self.data.hex(),
            "gas": hex(self.gas_limit),
            "maxFeePerGas": hex(self.max_fee_per_gas),
            "maxPriorityFeePerGas": hex(self.max_priority_fee),
            "nonce": hex(self.nonce),
            "chainId": hex(self.chain_id)
        }
        return json.dumps(tx_data)
    
    def get_hash(self) -> str:
        """Hash de la transaccion"""
        return hashlib.sha256(self.to_raw().encode()).hexdigest()


@dataclass
class ProtectResult:
    """Resultado de envio via Protect"""
    tx_hash: str
    status: str           # 'pending', 'included', 'failed', 'cancelled'
    block_number: Optional[int] = None
    refund_received: int = 0
    hints_used: List[str] = field(default_factory=list)
    simulation_result: Optional[Dict] = None
    
    def was_frontrun(self) -> bool:
        """Detectar si hubo intento de frontrun"""
        # En Protect, frontrunning deberia ser imposible
        return False


class FlashbotsProtectClient:
    """
    Cliente para Flashbots Protect RPC.
    
    Caracteristicas:
    - Transacciones no van a mempool publica
    - Simulacion pre-envio
    - Refund de MEV capturado
    - Reverts no se publican
    """
    
    MAINNET_RPC = "https://rpc.flashbots.net"
    GOERLI_RPC = "https://rpc-goerli.flashbots.net"
    
    def __init__(
        self,
        private_key: str,
        rpc_url: str = MAINNET_RPC,
        default_hints: List[HintPreference] = None
    ):
        self.private_key = private_key
        self.rpc_url = rpc_url
        self.default_hints = default_hints or [HintPreference.HASH]
        self.pending_txs: Dict[str, ProtectTransaction] = {}
        self.history: List[ProtectResult] = []
        
    def send_transaction(
        self,
        tx: ProtectTransaction,
        simulate_first: bool = True
    ) -> ProtectResult:
        """
        Enviar transaccion via Flashbots Protect.
        
        Args:
            tx: Transaccion a enviar
            simulate_first: Simular antes de enviar
            
        Returns:
            Resultado del envio
        """
        # 1. Simular si se solicita
        simulation = None
        if simulate_first:
            simulation = self._simulate_transaction(tx)
            if simulation.get("error"):
                return ProtectResult(
                    tx_hash=tx.get_hash(),
                    status="simulation_failed",
                    simulation_result=simulation
                )
        
        # 2. Preparar request con preferencias
        request = self._build_request(tx)
        
        # 3. Enviar a Protect RPC
        response = self._send_to_protect(request)
        
        # 4. Crear resultado
        result = ProtectResult(
            tx_hash=response.get("hash", tx.get_hash()),
            status="pending",
            hints_used=[h.value for h in tx.hints],
            simulation_result=simulation
        )
        
        # 5. Guardar en pending
        self.pending_txs[result.tx_hash] = tx
        
        return result
    
    def _build_request(self, tx: ProtectTransaction) -> Dict:
        """Construir request para Protect RPC"""
        # Headers con preferencias
        preferences = {
            "fast": True,  # Inclusion rapida
            "privacy": {
                "hints": [h.value for h in tx.hints],
                "builders": ["flashbots", "all"]  # Builders permitidos
            },
            "validity": {
                "refund": [{
                    "address": self._get_address(),
                    "percent": tx.refund_percent
                }]
            }
        }
        
        return {
            "jsonrpc": "2.0",
            "method": "eth_sendPrivateTransaction",
            "params": [{
                "tx": tx.to_raw(),
                "preferences": preferences
            }],
            "id": 1
        }
    
    def _simulate_transaction(self, tx: ProtectTransaction) -> Dict:
        """Simular transaccion antes de enviar"""
        # En produccion, llamar a eth_callBundle o similar
        return {
            "success": True,
            "gasUsed": tx.gas_limit * 0.8,
            "logs": [],
            "revert": None
        }
    
    def _send_to_protect(self, request: Dict) -> Dict:
        """Enviar request a Protect RPC"""
        # En produccion, hacer HTTP POST real
        return {
            "hash": hashlib.sha256(
                json.dumps(request).encode()
            ).hexdigest()[:64],
            "status": "pending"
        }
    
    def _get_address(self) -> str:
        """Obtener direccion del usuario"""
        # Derivar de private key (simplificado)
        return "0x" + hashlib.sha256(
            self.private_key.encode()
        ).hexdigest()[-40:]
    
    def check_status(self, tx_hash: str) -> ProtectResult:
        """Verificar estado de transaccion"""
        # En produccion, consultar RPC
        if tx_hash in self.pending_txs:
            return ProtectResult(
                tx_hash=tx_hash,
                status="pending"
            )
        return ProtectResult(tx_hash=tx_hash, status="unknown")
    
    def cancel_transaction(self, tx_hash: str) -> bool:
        """
        Cancelar transaccion pendiente.
        
        Solo posible si aun no fue incluida en bloque.
        """
        if tx_hash not in self.pending_txs:
            return False
        
        # En produccion, enviar eth_cancelPrivateTransaction
        del self.pending_txs[tx_hash]
        return True
    
    def get_stats(self) -> Dict[str, Any]:
        """Estadisticas de uso"""
        total_refunds = sum(r.refund_received for r in self.history)
        successful = sum(1 for r in self.history if r.status == "included")
        
        return {
            "total_transactions": len(self.history),
            "successful": successful,
            "total_refunds_wei": total_refunds,
            "pending": len(self.pending_txs),
            "frontrun_attempts": 0  # Protect previene esto
        }


class MEVBlockerClient:
    """
    Cliente para MEV Blocker RPC.
    
    Diferencias vs Flashbots Protect:
    - Coalition de builders
    - OFA (Order Flow Auction)
    - Refunds mas altos potencialmente
    """
    
    RPC_URL = "https://rpc.mevblocker.io"
    
    def __init__(self, private_key: str):
        self.private_key = private_key
        self.total_refunds = 0
        
    def send_transaction(
        self,
        tx: ProtectTransaction,
        auction_deadline_blocks: int = 25
    ) -> ProtectResult:
        """
        Enviar via MEV Blocker con OFA.
        
        El Order Flow Auction permite que builders
        compitan por incluir la transaccion,
        maximizando refund al usuario.
        """
        # OFA preferences
        ofa_config = {
            "deadlineBlocks": auction_deadline_blocks,
            "minRefundPercent": 80,
            "allowedBuilders": "all"
        }
        
        # Simular envio
        result = ProtectResult(
            tx_hash=tx.get_hash(),
            status="pending",
            hints_used=["ofa_enabled"]
        )
        
        return result
    
    def estimate_refund(
        self,
        tx: ProtectTransaction,
        mev_opportunity_wei: int
    ) -> int:
        """Estimar refund potencial de OFA"""
        # Tipicamente 80-90% del MEV capturado
        return int(mev_opportunity_wei * 0.85)


# Demo
def private_rpc_demo():
    print("=== Private RPC Demo ===\n")
    
    # Crear cliente Flashbots Protect
    client = FlashbotsProtectClient(
        private_key="0xdemo_private_key",
        default_hints=[HintPreference.HASH]
    )
    
    # Crear transaccion
    tx = ProtectTransaction(
        to="0x7a250d5630B4cF539739dF2C5dAcb4c659F2488D",  # Uniswap Router
        value=0,
        data=bytes.fromhex("38ed1739"),  # swap
        gas_limit=250000,
        max_fee_per_gas=50_000_000_000,  # 50 gwei
        max_priority_fee=2_000_000_000,   # 2 gwei
        nonce=42,
        hints=[HintPreference.HASH],      # Solo revelar hash
        refund_percent=90
    )
    
    print("1. Enviando transaccion via Flashbots Protect...")
    print(f"   Destino: {tx.to[:20]}... (Uniswap)")
    print(f"   Hints: {[h.value for h in tx.hints]}")
    print(f"   Refund: {tx.refund_percent}%")
    
    # Enviar
    result = client.send_transaction(tx, simulate_first=True)
    
    print(f"\n2. Resultado:")
    print(f"   TX Hash: {result.tx_hash[:20]}...")
    print(f"   Status: {result.status}")
    print(f"   Simulacion: {'OK' if result.simulation_result else 'N/A'}")
    
    # Comparar con MEV Blocker
    print("\n3. Comparando con MEV Blocker...")
    mev_client = MEVBlockerClient("0xdemo_key")
    estimated_refund = mev_client.estimate_refund(tx, 1_000_000_000_000)
    print(f"   Refund estimado: {estimated_refund / 1e18:.6f} ETH")
    
    print("\n4. Beneficios de Private RPC:")
    print("   - Sin frontrunning: TX no visible en mempool")
    print("   - Sin sandwich: Nadie puede ver tu swap")
    print("   - Reverts privados: Fallos no consumen gas publicamente")
    print("   - MEV rebate: Parte del MEV vuelve al usuario")

private_rpc_demo()
\end{lstlisting}

\begin{table}[htbp]
\centering
\small
\begin{tabular}{lcccc}
\toprule
\textbf{Servicio} & \textbf{Builders} & \textbf{Refund} & \textbf{Hints} & \textbf{Cancel} \\
\midrule
Flashbots Protect & Multi & 90\% & Configurable & Sí \\
MEV Blocker & Coalition & 80-90\% & OFA & Sí \\
Blocknative & Multi & Variable & Básicos & Sí \\
RPC público & N/A & 0\% & N/A & No \\
\bottomrule
\end{tabular}
\caption{Comparativa de servicios Private RPC}
\end{table}

%===========================================
\section{Relayer Networks}
\label{sec:relayers}
%===========================================

\begin{definition}[Relayer]
Servicio que envía transacciones en nombre del usuario, pagando el gas y ocultando la identidad real del remitente:
\[
\text{Usuario} \xrightarrow{\text{firma}} \text{Relayer} \xrightarrow{\text{TX + gas}} \text{Blockchain}
\]
\end{definition}

\subsection{Arquitectura de Red de Relayers}

\begin{figure}[htbp]
\centering
\begin{tikzpicture}[
    user/.style={circle, draw, fill=blue!20, minimum size=0.8cm, font=\small},
    relayer/.style={rectangle, draw, fill=green!20, minimum width=1.5cm, minimum height=0.7cm, font=\small},
    contract/.style={rectangle, draw, fill=orange!20, minimum width=2cm, minimum height=0.8cm, font=\small}
]
    % Usuarios
    \node[user] (u1) at (0, 2) {U1};
    \node[user] (u2) at (0, 1) {U2};
    \node[user] (u3) at (0, 0) {U3};
    
    % Relayers
    \node[relayer] (r1) at (4, 2.5) {R1};
    \node[relayer] (r2) at (4, 1.25) {R2};
    \node[relayer] (r3) at (4, 0) {R3};
    
    % Contrato
    \node[contract] (c) at (8, 1.25) {Privacy\\Contract};
    
    % Conexiones usuarios -> relayers (via Tor)
    \draw[->, dashed, blue] (u1) -- node[above, font=\tiny] {Tor} (r1);
    \draw[->, dashed, blue] (u1) -- (r2);
    \draw[->, dashed, blue] (u2) -- (r2);
    \draw[->, dashed, blue] (u2) -- (r3);
    \draw[->, dashed, blue] (u3) -- (r3);
    
    % Conexiones relayers -> contrato
    \draw[->, thick, green!60!black] (r1) -- (c);
    \draw[->, thick, green!60!black] (r2) -- (c);
    \draw[->, thick, green!60!black] (r3) -- (c);
    
    % Labels
    \node[font=\scriptsize] at (2, -1) {Conexion anonima};
    \node[font=\scriptsize] at (6, -1) {TX on-chain};
\end{tikzpicture}
\caption{Arquitectura de red de relayers descentralizada}
\end{figure}

\subsection{Implementación de Sistema de Relayers}

\begin{lstlisting}[style=python, caption={Sistema de relayers descentralizado}]
"""
Decentralized Relayer Network
Sistema de relayers para transacciones privadas
"""

from dataclasses import dataclass, field
from typing import Optional, Dict, List, Any, Callable
from enum import Enum
from datetime import datetime, timedelta
import hashlib
import json
import secrets
import heapq

class RelayerStatus(Enum):
    """Estado del relayer"""
    ACTIVE = "active"
    PAUSED = "paused"
    OFFLINE = "offline"
    SLASHED = "slashed"

@dataclass
class RelayerInfo:
    """Informacion de un relayer"""
    address: str
    endpoint: str              # URL del servicio
    stake: int                 # Stake en wei
    fee_percent: float         # Fee cobrado (e.g., 0.3%)
    min_amount: int            # Minimo por transaccion
    max_amount: int            # Maximo por transaccion
    supported_tokens: List[str]
    status: RelayerStatus = RelayerStatus.ACTIVE
    
    # Metricas
    total_relayed: int = 0
    successful_txs: int = 0
    failed_txs: int = 0
    average_latency_ms: float = 0
    uptime_percent: float = 100.0
    last_seen: datetime = field(default_factory=datetime.now)
    
    def get_score(self) -> float:
        """Score para seleccion (mayor = mejor)"""
        # Factores: fee, uptime, success rate, stake
        success_rate = (
            self.successful_txs / max(1, self.successful_txs + self.failed_txs)
        )
        
        # Normalizar fee (invertido - menor fee = mayor score)
        fee_score = max(0, 1 - self.fee_percent / 2)  # 0% = 1, 2% = 0
        
        # Stake score (logaritmico)
        import math
        stake_score = min(1, math.log10(max(1, self.stake / 1e18)) / 3)
        
        # Uptime
        uptime_score = self.uptime_percent / 100
        
        # Latency score (menor = mejor)
        latency_score = max(0, 1 - self.average_latency_ms / 5000)
        
        return (
            fee_score * 0.3 +
            success_rate * 0.25 +
            stake_score * 0.2 +
            uptime_score * 0.15 +
            latency_score * 0.1
        )
    
    def calculate_fee(self, amount: int) -> int:
        """Calcular fee para un monto"""
        return int(amount * self.fee_percent / 100)


@dataclass
class RelayRequest:
    """Request para relayer"""
    proof: bytes                    # ZK proof
    public_inputs: List[int]        # Inputs publicos
    recipient: str                  # Destinatario (puede ser contract)
    amount: int                     # Monto a transferir
    token: str                      # Token address
    fee: int                        # Fee para relayer
    deadline: datetime              # Deadline para inclusion
    refund_address: str             # Donde enviar refund si falla
    
    def get_hash(self) -> str:
        """Hash del request para firma"""
        data = json.dumps({
            "proof": self.proof.hex(),
            "public_inputs": self.public_inputs,
            "recipient": self.recipient,
            "amount": self.amount,
            "token": self.token,
            "fee": self.fee,
            "deadline": self.deadline.isoformat()
        }, sort_keys=True)
        return hashlib.sha256(data.encode()).hexdigest()


@dataclass
class RelayResult:
    """Resultado de relay"""
    request_hash: str
    status: str           # 'pending', 'submitted', 'confirmed', 'failed'
    tx_hash: Optional[str] = None
    relayer: Optional[str] = None
    fee_paid: int = 0
    gas_used: int = 0
    block_number: Optional[int] = None
    error: Optional[str] = None
    latency_ms: float = 0


class RelayerRegistry:
    """
    Registro descentralizado de relayers.
    
    Caracteristicas:
    - Registro abierto con stake
    - Slashing por mal comportamiento
    - Metricas on-chain verificables
    """
    
    def __init__(self, min_stake: int = 10**18):  # 1 ETH minimo
        self.min_stake = min_stake
        self.relayers: Dict[str, RelayerInfo] = {}
        self.slashed: set = set()
        
    def register(
        self,
        address: str,
        endpoint: str,
        stake: int,
        fee_percent: float,
        supported_tokens: List[str]
    ) -> bool:
        """Registrar nuevo relayer"""
        if stake < self.min_stake:
            return False
        
        if address in self.slashed:
            return False
        
        self.relayers[address] = RelayerInfo(
            address=address,
            endpoint=endpoint,
            stake=stake,
            fee_percent=fee_percent,
            min_amount=10**15,      # 0.001 ETH
            max_amount=100 * 10**18, # 100 ETH
            supported_tokens=supported_tokens
        )
        
        return True
    
    def get_active_relayers(
        self,
        token: Optional[str] = None,
        amount: Optional[int] = None
    ) -> List[RelayerInfo]:
        """Obtener relayers activos filtrados"""
        active = [
            r for r in self.relayers.values()
            if r.status == RelayerStatus.ACTIVE
        ]
        
        if token:
            active = [r for r in active if token in r.supported_tokens]
        
        if amount:
            active = [
                r for r in active 
                if r.min_amount <= amount <= r.max_amount
            ]
        
        return active
    
    def update_metrics(
        self,
        address: str,
        success: bool,
        latency_ms: float
    ):
        """Actualizar metricas de relayer"""
        if address not in self.relayers:
            return
        
        r = self.relayers[address]
        r.total_relayed += 1
        r.last_seen = datetime.now()
        
        if success:
            r.successful_txs += 1
        else:
            r.failed_txs += 1
        
        # Moving average de latencia
        n = r.successful_txs + r.failed_txs
        r.average_latency_ms = (
            (r.average_latency_ms * (n - 1) + latency_ms) / n
        )
    
    def slash(self, address: str, reason: str) -> int:
        """Slashear relayer por mal comportamiento"""
        if address not in self.relayers:
            return 0
        
        r = self.relayers[address]
        slash_amount = r.stake // 2  # 50% del stake
        
        r.stake -= slash_amount
        r.status = RelayerStatus.SLASHED
        self.slashed.add(address)
        
        return slash_amount


class RelayerSelector:
    """
    Selector inteligente de relayers.
    
    Estrategias:
    - Best score: Mejor relayer segun metricas
    - Random weighted: Aleatorio ponderado por score
    - Round robin: Rotacion para descentralizacion
    - Cheapest: Menor fee
    """
    
    def __init__(self, registry: RelayerRegistry):
        self.registry = registry
        self.last_used: Dict[str, datetime] = {}
        
    def select_best(
        self,
        token: str,
        amount: int,
        count: int = 3
    ) -> List[RelayerInfo]:
        """Seleccionar mejores N relayers"""
        candidates = self.registry.get_active_relayers(token, amount)
        
        # Ordenar por score
        scored = [(r.get_score(), r) for r in candidates]
        scored.sort(reverse=True)
        
        return [r for _, r in scored[:count]]
    
    def select_random_weighted(
        self,
        token: str,
        amount: int
    ) -> Optional[RelayerInfo]:
        """Seleccion aleatoria ponderada por score"""
        import random
        
        candidates = self.registry.get_active_relayers(token, amount)
        if not candidates:
            return None
        
        # Pesos basados en score
        weights = [r.get_score() for r in candidates]
        total = sum(weights)
        
        if total == 0:
            return random.choice(candidates)
        
        # Seleccion ponderada
        r = random.random() * total
        cumulative = 0
        for i, w in enumerate(weights):
            cumulative += w
            if r <= cumulative:
                return candidates[i]
        
        return candidates[-1]
    
    def select_cheapest(
        self,
        token: str,
        amount: int
    ) -> Optional[RelayerInfo]:
        """Seleccionar relayer mas barato"""
        candidates = self.registry.get_active_relayers(token, amount)
        if not candidates:
            return None
        
        return min(candidates, key=lambda r: r.calculate_fee(amount))


class RelayerClient:
    """
    Cliente para interactuar con red de relayers.
    
    Flujo:
    1. Seleccionar relayer(s)
    2. Enviar request
    3. Monitorear estado
    4. Reintentar si falla
    """
    
    def __init__(
        self,
        registry: RelayerRegistry,
        max_retries: int = 3
    ):
        self.registry = registry
        self.selector = RelayerSelector(registry)
        self.max_retries = max_retries
        self.pending: Dict[str, RelayRequest] = {}
        
    def submit_relay(
        self,
        request: RelayRequest,
        strategy: str = "best"
    ) -> RelayResult:
        """
        Enviar request a la red de relayers.
        
        Args:
            request: Request a enviar
            strategy: 'best', 'random', 'cheapest'
        """
        # Seleccionar relayer segun estrategia
        if strategy == "best":
            relayers = self.selector.select_best(
                request.token, request.amount
            )
            if not relayers:
                return RelayResult(
                    request_hash=request.get_hash(),
                    status="no_relayers",
                    error="No relayers available"
                )
            relayer = relayers[0]
        elif strategy == "random":
            relayer = self.selector.select_random_weighted(
                request.token, request.amount
            )
        else:  # cheapest
            relayer = self.selector.select_cheapest(
                request.token, request.amount
            )
        
        if not relayer:
            return RelayResult(
                request_hash=request.get_hash(),
                status="no_relayers"
            )
        
        # Verificar fee
        min_fee = relayer.calculate_fee(request.amount)
        if request.fee < min_fee:
            return RelayResult(
                request_hash=request.get_hash(),
                status="fee_too_low",
                error=f"Min fee: {min_fee}, provided: {request.fee}"
            )
        
        # Enviar request
        result = self._send_to_relayer(relayer, request)
        
        if result.status == "pending":
            self.pending[request.get_hash()] = request
        
        return result
    
    def _send_to_relayer(
        self,
        relayer: RelayerInfo,
        request: RelayRequest
    ) -> RelayResult:
        """Enviar request a relayer especifico"""
        import time
        start = time.time()
        
        # Simular llamada HTTP
        # En produccion: POST a relayer.endpoint
        
        latency = (time.time() - start) * 1000
        
        # Simular respuesta exitosa
        tx_hash = hashlib.sha256(
            (request.get_hash() + relayer.address).encode()
        ).hexdigest()
        
        # Actualizar metricas
        self.registry.update_metrics(
            relayer.address,
            success=True,
            latency_ms=latency
        )
        
        return RelayResult(
            request_hash=request.get_hash(),
            status="pending",
            tx_hash=tx_hash,
            relayer=relayer.address,
            fee_paid=request.fee,
            latency_ms=latency
        )
    
    def check_status(self, request_hash: str) -> RelayResult:
        """Verificar estado de request"""
        if request_hash not in self.pending:
            return RelayResult(
                request_hash=request_hash,
                status="unknown"
            )
        
        # Simular verificacion on-chain
        return RelayResult(
            request_hash=request_hash,
            status="confirmed",
            block_number=12345678
        )


# Demo
def relayer_demo():
    print("=== Relayer Network Demo ===\n")
    
    # Crear registro
    registry = RelayerRegistry(min_stake=10**18)
    
    # Registrar relayers
    print("1. Registrando relayers en la red...")
    
    relayers_data = [
        ("0xRelayer1", "https://relay1.example.com", 5 * 10**18, 0.3, 98.5),
        ("0xRelayer2", "https://relay2.example.com", 10 * 10**18, 0.2, 99.2),
        ("0xRelayer3", "https://relay3.example.com", 2 * 10**18, 0.5, 95.0),
    ]
    
    for addr, endpoint, stake, fee, uptime in relayers_data:
        registry.register(
            address=addr,
            endpoint=endpoint,
            stake=stake,
            fee_percent=fee,
            supported_tokens=["ETH", "USDC", "DAI"]
        )
        r = registry.relayers[addr]
        r.uptime_percent = uptime
        r.successful_txs = 100
        print(f"   {addr}: stake={stake/1e18:.0f} ETH, fee={fee}%")
    
    # Crear cliente
    client = RelayerClient(registry)
    
    # Crear request
    print("\n2. Creando relay request...")
    request = RelayRequest(
        proof=secrets.token_bytes(256),
        public_inputs=[1, 2, 3],
        recipient="0xRecipientAddress",
        amount=10**18,  # 1 ETH
        token="ETH",
        fee=3 * 10**15,  # 0.003 ETH
        deadline=datetime.now() + timedelta(minutes=10),
        refund_address="0xRefundAddress"
    )
    print(f"   Amount: {request.amount / 1e18:.2f} ETH")
    print(f"   Fee: {request.fee / 1e18:.4f} ETH")
    
    # Seleccionar y enviar
    print("\n3. Seleccionando mejor relayer...")
    selector = RelayerSelector(registry)
    best = selector.select_best("ETH", request.amount, count=3)
    
    for i, r in enumerate(best):
        print(f"   #{i+1}: {r.address}, score={r.get_score():.3f}, fee={r.fee_percent}%")
    
    # Enviar request
    print("\n4. Enviando request...")
    result = client.submit_relay(request, strategy="best")
    
    print(f"   Status: {result.status}")
    print(f"   TX Hash: {result.tx_hash[:20]}...")
    print(f"   Relayer: {result.relayer}")
    print(f"   Latency: {result.latency_ms:.2f} ms")

relayer_demo()
\end{lstlisting}

%===========================================
\section{Note Management}
\label{sec:note-management}
%===========================================

\begin{definition}[Note]
Representación criptográfica de un depósito privado que contiene:
\begin{itemize}
    \item \textbf{Commitment:} Hash público en el contrato ($C = H(v, r, k)$)
    \item \textbf{Nullifier secret:} Secreto para gastar ($k$)
    \item \textbf{Amount:} Valor del depósito ($v$)
    \item \textbf{Blinding:} Factor de randomización ($r$)
\end{itemize}
\end{definition}

\subsection{Ciclo de Vida de una Note}

\begin{figure}[htbp]
\centering
\begin{tikzpicture}[
    state/.style={ellipse, draw, minimum width=2cm, minimum height=0.8cm, align=center, font=\small},
    arrow/.style={->, thick}
]
    % Estados
    \node[state, fill=green!20] (created) at (0, 0) {Created};
    \node[state, fill=blue!20] (pending) at (3.5, 0) {Pending};
    \node[state, fill=purple!20] (confirmed) at (7, 0) {Confirmed};
    \node[state, fill=orange!20] (spending) at (10.5, 0) {Spending};
    \node[state, fill=red!20] (spent) at (7, -2) {Spent};
    
    % Transiciones
    \draw[arrow] (created) -- node[above, font=\tiny] {deposit tx} (pending);
    \draw[arrow] (pending) -- node[above, font=\tiny] {block conf} (confirmed);
    \draw[arrow] (confirmed) -- node[above, font=\tiny] {withdraw} (spending);
    \draw[arrow] (spending) -- node[right, font=\tiny] {nullifier} (spent);
    
    % Self-loop para confirmed (transfer)
    \draw[arrow] (confirmed) to[loop above] node[above, font=\tiny] {transfer} (confirmed);
\end{tikzpicture}
\caption{Ciclo de vida de una note en sistema de privacidad}
\end{figure}

\subsection{Sistema de Gestión de Notes}

\begin{lstlisting}[style=python, caption={Sistema completo de gestión de notes}]
"""
Secure Note Management System
Gestion completa del ciclo de vida de notes privadas
"""

from dataclasses import dataclass, field
from typing import Optional, Dict, List, Tuple, Any, Callable
from datetime import datetime, timedelta
from enum import Enum
import hashlib
import json
import secrets
import os
from abc import ABC, abstractmethod

class NoteStatus(Enum):
    """Estado de la note"""
    CREATED = "created"         # Creada localmente
    PENDING = "pending"         # TX enviada, esperando confirmacion
    CONFIRMED = "confirmed"     # Depositada y lista para usar
    SPENDING = "spending"       # En proceso de gasto
    SPENT = "spent"             # Gastada (nullifier revelado)
    FAILED = "failed"           # Deposito fallido

@dataclass
class Note:
    """
    Representacion completa de una note privada.
    
    CRITICAL: Esta informacion es sensible.
    Perder nullifier_secret = perder fondos.
    """
    # Identificadores
    note_id: str                    # ID local unico
    commitment: str                 # Hash publico on-chain
    
    # Secretos (NUNCA revelar)
    nullifier_secret: bytes         # Para generar nullifier
    blinding_factor: bytes          # Factor de randomizacion
    
    # Datos
    amount: int                     # Monto en wei
    asset: str                      # Token address o "ETH"
    
    # Metadata
    status: NoteStatus = NoteStatus.CREATED
    created_at: datetime = field(default_factory=datetime.now)
    confirmed_at: Optional[datetime] = None
    spent_at: Optional[datetime] = None
    
    # Referencias on-chain
    deposit_tx: Optional[str] = None
    deposit_block: Optional[int] = None
    spending_tx: Optional[str] = None
    
    # Index en el arbol Merkle
    leaf_index: Optional[int] = None
    
    def compute_nullifier(self) -> str:
        """Calcular nullifier hash"""
        preimage = (
            self.nullifier_secret +
            self.commitment.encode() +
            self.amount.to_bytes(32, 'big')
        )
        return hashlib.sha256(preimage).hexdigest()
    
    def to_dict(self, include_secrets: bool = False) -> Dict:
        """Serializar note (sin secretos por defecto)"""
        data = {
            "note_id": self.note_id,
            "commitment": self.commitment,
            "amount": self.amount,
            "asset": self.asset,
            "status": self.status.value,
            "created_at": self.created_at.isoformat(),
            "deposit_tx": self.deposit_tx,
            "leaf_index": self.leaf_index
        }
        
        if include_secrets:
            data["nullifier_secret"] = self.nullifier_secret.hex()
            data["blinding_factor"] = self.blinding_factor.hex()
        
        return data
    
    @classmethod
    def from_dict(cls, data: Dict) -> 'Note':
        """Deserializar note"""
        return cls(
            note_id=data["note_id"],
            commitment=data["commitment"],
            nullifier_secret=bytes.fromhex(data.get("nullifier_secret", "00" * 32)),
            blinding_factor=bytes.fromhex(data.get("blinding_factor", "00" * 32)),
            amount=data["amount"],
            asset=data["asset"],
            status=NoteStatus(data.get("status", "created")),
            created_at=datetime.fromisoformat(data["created_at"]),
            deposit_tx=data.get("deposit_tx"),
            leaf_index=data.get("leaf_index")
        )


class EncryptionProvider(ABC):
    """Interfaz abstracta para encriptacion"""
    
    @abstractmethod
    def encrypt(self, data: bytes) -> bytes:
        pass
    
    @abstractmethod
    def decrypt(self, encrypted: bytes) -> bytes:
        pass


class AESEncryption(EncryptionProvider):
    """
    Encriptacion AES-256-GCM.
    
    En produccion, usar libreria como cryptography.
    """
    
    def __init__(self, key: bytes):
        if len(key) != 32:
            raise ValueError("Key must be 32 bytes")
        self.key = key
    
    def encrypt(self, data: bytes) -> bytes:
        """Encriptar con AES-GCM (simplificado)"""
        # En produccion: usar AES-GCM real
        nonce = secrets.token_bytes(12)
        
        # XOR simple para demo (NO usar en produccion)
        key_stream = self._expand_key(len(data))
        ciphertext = bytes(a ^ b for a, b in zip(data, key_stream))
        
        # MAC simplificado
        mac = hashlib.sha256(self.key + nonce + ciphertext).digest()[:16]
        
        return nonce + ciphertext + mac
    
    def decrypt(self, encrypted: bytes) -> bytes:
        """Desencriptar"""
        nonce = encrypted[:12]
        mac = encrypted[-16:]
        ciphertext = encrypted[12:-16]
        
        # Verificar MAC
        expected_mac = hashlib.sha256(
            self.key + nonce + ciphertext
        ).digest()[:16]
        if mac != expected_mac:
            raise ValueError("MAC verification failed")
        
        # Desencriptar
        key_stream = self._expand_key(len(ciphertext))
        return bytes(a ^ b for a, b in zip(ciphertext, key_stream))
    
    def _expand_key(self, length: int) -> bytes:
        """Expandir key para keystream"""
        expanded = b""
        counter = 0
        while len(expanded) < length:
            expanded += hashlib.sha256(
                self.key + counter.to_bytes(4, 'big')
            ).digest()
            counter += 1
        return expanded[:length]


class NoteManager:
    """
    Gestor completo de notes privadas.
    
    Responsabilidades:
    - Crear notes seguras
    - Almacenamiento encriptado
    - Backup y recovery
    - Sincronizacion multi-device
    """
    
    def __init__(
        self,
        encryption: EncryptionProvider,
        storage_path: Optional[str] = None
    ):
        self.encryption = encryption
        self.storage_path = storage_path
        self.notes: Dict[str, Note] = {}
        self.nullifiers: Dict[str, str] = {}  # nullifier -> note_id
        
    def create_note(
        self,
        amount: int,
        asset: str
    ) -> Note:
        """
        Crear nueva note para deposito.
        
        Genera secretos criptograficamente seguros.
        """
        # Generar secretos
        nullifier_secret = secrets.token_bytes(32)
        blinding_factor = secrets.token_bytes(32)
        
        # Calcular commitment
        preimage = (
            amount.to_bytes(32, 'big') +
            asset.encode().ljust(32, b'\x00') +
            nullifier_secret +
            blinding_factor
        )
        commitment = hashlib.sha256(preimage).hexdigest()
        
        # Crear note
        note = Note(
            note_id=secrets.token_hex(16),
            commitment=commitment,
            nullifier_secret=nullifier_secret,
            blinding_factor=blinding_factor,
            amount=amount,
            asset=asset,
            status=NoteStatus.CREATED
        )
        
        # Guardar
        self.notes[note.note_id] = note
        nullifier = note.compute_nullifier()
        self.nullifiers[nullifier] = note.note_id
        
        return note
    
    def mark_pending(self, note_id: str, tx_hash: str):
        """Marcar note como pendiente (TX enviada)"""
        if note_id not in self.notes:
            raise ValueError(f"Note not found: {note_id}")
        
        note = self.notes[note_id]
        note.status = NoteStatus.PENDING
        note.deposit_tx = tx_hash
    
    def mark_confirmed(
        self,
        note_id: str,
        block_number: int,
        leaf_index: int
    ):
        """Confirmar deposito de note"""
        if note_id not in self.notes:
            raise ValueError(f"Note not found: {note_id}")
        
        note = self.notes[note_id]
        note.status = NoteStatus.CONFIRMED
        note.confirmed_at = datetime.now()
        note.deposit_block = block_number
        note.leaf_index = leaf_index
    
    def mark_spent(self, note_id: str, tx_hash: str):
        """Marcar note como gastada"""
        if note_id not in self.notes:
            raise ValueError(f"Note not found: {note_id}")
        
        note = self.notes[note_id]
        note.status = NoteStatus.SPENT
        note.spent_at = datetime.now()
        note.spending_tx = tx_hash
    
    def get_spendable_notes(
        self,
        asset: str,
        min_confirmations: int = 10
    ) -> List[Note]:
        """Obtener notes disponibles para gastar"""
        return [
            note for note in self.notes.values()
            if note.status == NoteStatus.CONFIRMED
            and note.asset == asset
        ]
    
    def get_balance(self, asset: str) -> int:
        """Balance total de un asset"""
        return sum(
            note.amount
            for note in self.get_spendable_notes(asset)
        )
    
    def export_backup(self) -> bytes:
        """
        Exportar backup encriptado de todas las notes.
        
        CRITICAL: Este backup permite recuperar fondos.
        """
        # Serializar todas las notes con secretos
        backup_data = {
            "version": 1,
            "created_at": datetime.now().isoformat(),
            "notes": [
                note.to_dict(include_secrets=True)
                for note in self.notes.values()
            ],
            "nullifiers": dict(self.nullifiers)
        }
        
        # Encriptar
        json_bytes = json.dumps(backup_data).encode()
        encrypted = self.encryption.encrypt(json_bytes)
        
        return encrypted
    
    def import_backup(self, encrypted_backup: bytes):
        """Restaurar notes desde backup encriptado"""
        # Desencriptar
        json_bytes = self.encryption.decrypt(encrypted_backup)
        backup_data = json.loads(json_bytes.decode())
        
        # Restaurar notes
        for note_data in backup_data["notes"]:
            note = Note.from_dict(note_data)
            self.notes[note.note_id] = note
        
        # Restaurar nullifiers
        self.nullifiers = backup_data.get("nullifiers", {})
    
    def check_nullifier_spent(self, nullifier: str) -> bool:
        """Verificar si nullifier ya fue usado on-chain"""
        # En produccion: consultar contrato
        return nullifier in self.nullifiers
    
    def scan_chain_for_notes(
        self,
        viewing_key: bytes,
        from_block: int,
        to_block: int
    ) -> List[Note]:
        """
        Escanear blockchain buscando notes nuestras.
        
        Usado para recovery o sync.
        """
        # En produccion: iterar eventos de deposito
        # y probar desencriptar con viewing_key
        found_notes = []
        return found_notes
    
    def get_stats(self) -> Dict[str, Any]:
        """Estadisticas del note manager"""
        by_status = {}
        by_asset = {}
        
        for note in self.notes.values():
            status = note.status.value
            by_status[status] = by_status.get(status, 0) + 1
            
            if note.status == NoteStatus.CONFIRMED:
                asset = note.asset
                by_asset[asset] = by_asset.get(asset, 0) + note.amount
        
        return {
            "total_notes": len(self.notes),
            "by_status": by_status,
            "balances": by_asset,
            "nullifiers_tracked": len(self.nullifiers)
        }


class SecureBackupManager:
    """
    Gestion segura de backups.
    
    Estrategias:
    - Split backup (Shamir Secret Sharing)
    - Cloud sync encriptado
    - Hardware backup
    """
    
    def __init__(self, note_manager: NoteManager):
        self.note_manager = note_manager
        self.backup_history: List[Dict] = []
    
    def create_split_backup(
        self,
        threshold: int = 2,
        total_shares: int = 3
    ) -> List[bytes]:
        """
        Crear backup dividido con Shamir Secret Sharing.
        
        Args:
            threshold: Shares minimos para reconstruir
            total_shares: Total de shares a generar
        """
        # Obtener backup completo
        full_backup = self.note_manager.export_backup()
        
        # Simplificado - en produccion usar Shamir real
        # Por ahora, generamos "shares" simulados
        shares = []
        for i in range(total_shares):
            share_data = {
                "share_index": i,
                "threshold": threshold,
                "total": total_shares,
                "data": full_backup.hex()  # Simplificado
            }
            shares.append(json.dumps(share_data).encode())
        
        self.backup_history.append({
            "type": "split",
            "threshold": threshold,
            "total_shares": total_shares,
            "created_at": datetime.now().isoformat()
        })
        
        return shares
    
    def recover_from_shares(self, shares: List[bytes]) -> bool:
        """Reconstruir backup desde shares"""
        # Verificar threshold
        share_data = json.loads(shares[0].decode())
        threshold = share_data["threshold"]
        
        if len(shares) < threshold:
            raise ValueError(
                f"Need at least {threshold} shares, got {len(shares)}"
            )
        
        # Reconstruir (simplificado)
        encrypted_backup = bytes.fromhex(share_data["data"])
        self.note_manager.import_backup(encrypted_backup)
        
        return True
    
    def create_mnemonic_backup(self) -> List[str]:
        """
        Crear backup como mnemonic phrase.
        
        La phrase codifica el master secret del cual
        se derivan todas las keys.
        """
        # Generar entropia
        entropy = secrets.token_bytes(32)
        
        # Convertir a palabras (simplificado)
        # En produccion usar BIP39
        wordlist = [
            "abandon", "ability", "able", "about", "above",
            "absent", "absorb", "abstract", "absurd", "abuse",
            "access", "accident", "account", "accuse", "achieve",
            "acid", "acoustic", "acquire", "across", "act"
        ]  # Solo ejemplo
        
        words = []
        for i in range(0, 32, 2):
            index = int.from_bytes(entropy[i:i+2], 'big') % len(wordlist)
            words.append(wordlist[index])
        
        return words


# Demo
def note_management_demo():
    print("=== Note Management Demo ===\n")
    
    # Setup encriptacion
    master_key = secrets.token_bytes(32)
    encryption = AESEncryption(master_key)
    
    # Crear note manager
    manager = NoteManager(encryption)
    
    # 1. Crear notes
    print("1. Creando notes...")
    note1 = manager.create_note(amount=10**18, asset="ETH")
    print(f"   Note 1: {note1.amount / 1e18:.2f} ETH")
    print(f"   Commitment: {note1.commitment[:20]}...")
    
    note2 = manager.create_note(amount=5 * 10**18, asset="ETH")
    print(f"   Note 2: {note2.amount / 1e18:.2f} ETH")
    
    # 2. Simular deposito
    print("\n2. Simulando depositos...")
    manager.mark_pending(note1.note_id, "0xdeposit_tx_1")
    manager.mark_confirmed(note1.note_id, block_number=12345, leaf_index=100)
    print(f"   Note 1 confirmada en bloque 12345")
    
    manager.mark_pending(note2.note_id, "0xdeposit_tx_2")
    manager.mark_confirmed(note2.note_id, block_number=12350, leaf_index=101)
    print(f"   Note 2 confirmada en bloque 12350")
    
    # 3. Balance
    balance = manager.get_balance("ETH")
    print(f"\n3. Balance total: {balance / 1e18:.2f} ETH")
    
    # 4. Crear backup
    print("\n4. Creando backup encriptado...")
    backup = manager.export_backup()
    print(f"   Backup size: {len(backup)} bytes")
    
    # 5. Split backup
    print("\n5. Creando split backup (2-of-3)...")
    backup_manager = SecureBackupManager(manager)
    shares = backup_manager.create_split_backup(threshold=2, total_shares=3)
    print(f"   Generados {len(shares)} shares")
    print(f"   Threshold: 2 shares para recuperar")
    
    # 6. Stats
    stats = manager.get_stats()
    print(f"\n6. Estadisticas:")
    print(f"   Total notes: {stats['total_notes']}")
    print(f"   Por status: {stats['by_status']}")
    print(f"   Balance ETH: {stats['balances'].get('ETH', 0) / 1e18:.2f}")
    
    # 7. Verificar nullifier
    print("\n7. Verificando nullifier...")
    nullifier = note1.compute_nullifier()
    print(f"   Nullifier: {nullifier[:20]}...")
    print(f"   Registrado: {manager.check_nullifier_spent(nullifier)}")

note_management_demo()
\end{lstlisting}

\begin{keyidea}
\textbf{Mejores Prácticas para Note Management:}

\begin{enumerate}
    \item \textbf{Backup inmediato:} Crear backup al generar note, ANTES de depositar
    \item \textbf{Encriptación fuerte:} AES-256-GCM con key derivada de passphrase
    \item \textbf{Split backup:} Shamir 2-of-3 para redundancia sin single point of failure
    \item \textbf{Verificación:} Probar recovery ANTES de depositar fondos significativos
    \item \textbf{Rotación:} Rotar keys periódicamente sin perder acceso a notes antiguas
\end{enumerate}

\textbf{Qué NO hacer:}
\begin{itemize}
    \item Almacenar notes sin encriptar
    \item Backup único sin redundancia
    \item Compartir nullifier\_secret
    \item Depositar sin verificar backup funciona
\end{itemize}
\end{keyidea}

%===========================================
\section{Proving Infrastructure}
\label{sec:proving-infrastructure}
%===========================================

\begin{definition}[Proving Infrastructure]
Conjunto de servicios y herramientas para la generación de zero-knowledge proofs:
\begin{itemize}
    \item \textbf{Provers:} Sistemas que ejecutan la generación de pruebas
    \item \textbf{Verification:} Contratos/sistemas que verifican pruebas
    \item \textbf{Key management:} Gestión de proving/verification keys
    \item \textbf{Witness generation:} Cálculo de inputs privados
\end{itemize}
\end{definition}

\subsection{Arquitecturas de Proving}

\begin{figure}[htbp]
\centering
\begin{tikzpicture}[
    box/.style={rectangle, draw, rounded corners, minimum width=3cm, minimum height=1.2cm, align=center, font=\small},
    label/.style={font=\scriptsize, align=center}
]
    % Client-side
    \node[box, fill=green!20] (client) at (0, 0) {\textbf{Client-side}\\Browser/Mobile};
    \node[label] at (0, -1.5) {
        + Máxima privacidad\\
        - Lento en devices\\
        - Requiere WASM
    };
    
    % Server-side
    \node[box, fill=red!20] (server) at (5, 0) {\textbf{Server-side}\\Cloud Prover};
    \node[label] at (5, -1.5) {
        + Muy rápido\\
        - Menor privacidad\\
        - Servidor ve witness
    };
    
    % Hybrid
    \node[box, fill=purple!20] (hybrid) at (10, 0) {\textbf{Híbrido}\\Split Proving};
    \node[label] at (10, -1.5) {
        + Balance\\
        + Datos sensibles local\\
        - Más complejo
    };
    
    % Escala de privacidad
    \draw[<->, thick] (0, -3) -- (10, -3);
    \node[font=\scriptsize] at (0, -3.3) {Alta privacidad};
    \node[font=\scriptsize] at (10, -3.3) {Alto rendimiento};
\end{tikzpicture}
\caption{Trade-offs en arquitecturas de proving}
\end{figure}

\subsection{Sistema de Proving Completo}

\begin{lstlisting}[style=python, caption={Infraestructura de proving ZK}]
"""
ZK Proving Infrastructure
Sistema completo para generacion y verificacion de pruebas
"""

from dataclasses import dataclass, field
from typing import Optional, Dict, List, Any, Tuple, Callable
from datetime import datetime, timedelta
from enum import Enum
from abc import ABC, abstractmethod
import hashlib
import json
import time
import secrets

class ProvingMode(Enum):
    """Modos de generacion de pruebas"""
    CLIENT = "client"       # Completamente en cliente
    SERVER = "server"       # Completamente en servidor
    HYBRID = "hybrid"       # Division client/server
    DELEGATED = "delegated" # Delegado con MPC

class ProofSystem(Enum):
    """Sistemas de pruebas soportados"""
    GROTH16 = "groth16"     # Trusted setup, pequeñas pruebas
    PLONK = "plonk"         # Universal setup
    STARK = "stark"         # Sin trusted setup, pruebas grandes
    HALO2 = "halo2"         # Recursive, sin trusted setup

@dataclass
class ProvingKey:
    """Clave de proving para un circuito"""
    circuit_id: str
    proof_system: ProofSystem
    key_data: bytes
    size_bytes: int
    created_at: datetime
    
    def get_hash(self) -> str:
        return hashlib.sha256(self.key_data).hexdigest()

@dataclass
class VerificationKey:
    """Clave de verificacion (publica)"""
    circuit_id: str
    proof_system: ProofSystem
    key_data: bytes
    commitment: str  # Commitment on-chain
    
    def verify_proof(self, proof: bytes, public_inputs: List[int]) -> bool:
        """Verificar prueba (simplificado)"""
        # En produccion: llamar a verificador real
        return len(proof) > 0 and len(public_inputs) > 0

@dataclass
class Witness:
    """Witness (inputs privados) para generar prueba"""
    circuit_id: str
    private_inputs: Dict[str, Any]
    public_inputs: List[int]
    created_at: datetime = field(default_factory=datetime.now)
    
    def is_sensitive(self) -> bool:
        """Verificar si contiene datos sensibles"""
        sensitive_fields = ["nullifier", "secret", "private_key", "amount"]
        return any(
            field in self.private_inputs
            for field in sensitive_fields
        )
    
    def get_public_hash(self) -> str:
        """Hash de inputs publicos (seguro de compartir)"""
        return hashlib.sha256(
            json.dumps(self.public_inputs).encode()
        ).hexdigest()

@dataclass
class Proof:
    """Prueba ZK generada"""
    circuit_id: str
    proof_system: ProofSystem
    proof_data: bytes
    public_inputs: List[int]
    
    # Metadata
    proving_time_ms: float
    prover_mode: ProvingMode
    created_at: datetime = field(default_factory=datetime.now)
    
    def size_bytes(self) -> int:
        return len(self.proof_data)
    
    def to_calldata(self) -> bytes:
        """Serializar para llamada a contrato"""
        # Formato: proof || public_inputs
        inputs_bytes = b"".join(
            i.to_bytes(32, 'big') for i in self.public_inputs
        )
        return self.proof_data + inputs_bytes


class Prover(ABC):
    """Interfaz abstracta para provers"""
    
    @abstractmethod
    def generate_proof(
        self,
        proving_key: ProvingKey,
        witness: Witness
    ) -> Proof:
        pass
    
    @abstractmethod
    def get_mode(self) -> ProvingMode:
        pass


class ClientProver(Prover):
    """
    Prover del lado del cliente.
    
    Genera pruebas completamente en el dispositivo del usuario.
    Maxima privacidad pero puede ser lento.
    """
    
    def __init__(self, wasm_path: Optional[str] = None):
        self.wasm_path = wasm_path
        self.loaded_circuits: Dict[str, bool] = {}
        
    def load_circuit(self, circuit_id: str):
        """Cargar circuito WASM"""
        # En produccion: cargar archivo WASM real
        self.loaded_circuits[circuit_id] = True
    
    def generate_proof(
        self,
        proving_key: ProvingKey,
        witness: Witness
    ) -> Proof:
        """Generar prueba localmente"""
        start_time = time.time()
        
        if proving_key.circuit_id not in self.loaded_circuits:
            self.load_circuit(proving_key.circuit_id)
        
        # Simular generacion de prueba
        # En produccion: llamar a snarkjs/circom WASM
        proof_data = self._generate_groth16_proof(
            proving_key, witness
        )
        
        proving_time = (time.time() - start_time) * 1000
        
        return Proof(
            circuit_id=proving_key.circuit_id,
            proof_system=proving_key.proof_system,
            proof_data=proof_data,
            public_inputs=witness.public_inputs,
            proving_time_ms=proving_time,
            prover_mode=ProvingMode.CLIENT
        )
    
    def _generate_groth16_proof(
        self,
        proving_key: ProvingKey,
        witness: Witness
    ) -> bytes:
        """Generar prueba Groth16 (simulado)"""
        # En produccion: llamar a prover real
        # Groth16 proof: 3 puntos G1 + 1 punto G2 ≈ 256 bytes
        proof_hash = hashlib.sha256(
            proving_key.get_hash().encode() +
            json.dumps(witness.private_inputs).encode()
        ).digest()
        
        # Simular tamaño de prueba Groth16
        return proof_hash * 8  # ~256 bytes
    
    def get_mode(self) -> ProvingMode:
        return ProvingMode.CLIENT
    
    def estimate_time(self, circuit_id: str) -> float:
        """Estimar tiempo de proving en ms"""
        # Estimaciones tipicas para diferentes circuitos
        estimates = {
            "deposit": 5000,      # 5 segundos
            "transfer": 15000,    # 15 segundos
            "withdrawal": 8000,   # 8 segundos
        }
        return estimates.get(circuit_id, 10000)


class ServerProver(Prover):
    """
    Prover del lado del servidor.
    
    Rapido pero requiere enviar witness al servidor.
    El servidor ve los datos privados.
    """
    
    def __init__(self, server_url: str, api_key: str):
        self.server_url = server_url
        self.api_key = api_key
        self.request_count = 0
        
    def generate_proof(
        self,
        proving_key: ProvingKey,
        witness: Witness
    ) -> Proof:
        """Enviar witness al servidor para proving"""
        start_time = time.time()
        
        # WARNING: El servidor ve el witness completo
        if witness.is_sensitive():
            print("WARNING: Sending sensitive data to server!")
        
        # Simular request al servidor
        proof_data = self._request_proof(proving_key, witness)
        
        proving_time = (time.time() - start_time) * 1000
        self.request_count += 1
        
        return Proof(
            circuit_id=proving_key.circuit_id,
            proof_system=proving_key.proof_system,
            proof_data=proof_data,
            public_inputs=witness.public_inputs,
            proving_time_ms=proving_time,
            prover_mode=ProvingMode.SERVER
        )
    
    def _request_proof(
        self,
        proving_key: ProvingKey,
        witness: Witness
    ) -> bytes:
        """Request al servidor de proving"""
        # En produccion: HTTP POST a servidor
        # Simular respuesta
        return secrets.token_bytes(256)
    
    def get_mode(self) -> ProvingMode:
        return ProvingMode.SERVER


class HybridProver(Prover):
    """
    Prover hibrido.
    
    Divide el circuito en partes:
    - Partes sensibles: client-side
    - Partes computacionalmente pesadas: server-side
    
    Usa folding schemes o split proving.
    """
    
    def __init__(
        self,
        client_prover: ClientProver,
        server_prover: ServerProver
    ):
        self.client = client_prover
        self.server = server_prover
        
    def generate_proof(
        self,
        proving_key: ProvingKey,
        witness: Witness
    ) -> Proof:
        """
        Generar prueba en modo hibrido.
        
        1. Client genera commitment a datos sensibles
        2. Server genera prueba de parte publica
        3. Client combina ambas partes
        """
        start_time = time.time()
        
        # Separar witness en sensible y no-sensible
        sensitive_inputs = {}
        public_computation_inputs = {}
        
        for key, value in witness.private_inputs.items():
            if key in ["nullifier_secret", "private_key", "amount"]:
                sensitive_inputs[key] = value
            else:
                public_computation_inputs[key] = value
        
        # 1. Client-side: commitment a datos sensibles
        sensitive_commitment = self._client_commit(sensitive_inputs)
        
        # 2. Server-side: proving de parte publica
        # El servidor NO ve los datos sensibles, solo el commitment
        server_witness = Witness(
            circuit_id=witness.circuit_id + "_public",
            private_inputs={
                **public_computation_inputs,
                "sensitive_commitment": sensitive_commitment
            },
            public_inputs=witness.public_inputs
        )
        
        # Simulamos que el server hace parte del trabajo
        # pero sin ver datos sensibles
        
        proving_time = (time.time() - start_time) * 1000
        
        # Combinar pruebas
        combined_proof = self._combine_proofs(
            sensitive_commitment,
            witness
        )
        
        return Proof(
            circuit_id=proving_key.circuit_id,
            proof_system=proving_key.proof_system,
            proof_data=combined_proof,
            public_inputs=witness.public_inputs,
            proving_time_ms=proving_time,
            prover_mode=ProvingMode.HYBRID
        )
    
    def _client_commit(self, sensitive_inputs: Dict) -> str:
        """Generar commitment a datos sensibles"""
        data = json.dumps(sensitive_inputs, sort_keys=True)
        return hashlib.sha256(data.encode()).hexdigest()
    
    def _combine_proofs(
        self,
        commitment: str,
        witness: Witness
    ) -> bytes:
        """Combinar partes de prueba"""
        return hashlib.sha256(
            commitment.encode() +
            json.dumps(witness.public_inputs).encode()
        ).digest() * 8
    
    def get_mode(self) -> ProvingMode:
        return ProvingMode.HYBRID


class ProvingService:
    """
    Servicio completo de proving.
    
    Caracteristicas:
    - Seleccion automatica de modo
    - Cache de proving keys
    - Benchmarking
    - Rate limiting
    """
    
    def __init__(self):
        self.provers: Dict[ProvingMode, Prover] = {}
        self.proving_keys: Dict[str, ProvingKey] = {}
        self.verification_keys: Dict[str, VerificationKey] = {}
        self.cache: Dict[str, Proof] = {}  # Cache de pruebas
        self.benchmarks: Dict[str, List[float]] = {}
        
    def register_prover(self, prover: Prover):
        """Registrar prover disponible"""
        self.provers[prover.get_mode()] = prover
    
    def load_circuit(
        self,
        circuit_id: str,
        proving_key: ProvingKey,
        verification_key: VerificationKey
    ):
        """Cargar circuito con sus claves"""
        self.proving_keys[circuit_id] = proving_key
        self.verification_keys[circuit_id] = verification_key
    
    def generate_proof(
        self,
        circuit_id: str,
        witness: Witness,
        mode: Optional[ProvingMode] = None
    ) -> Proof:
        """
        Generar prueba con modo optimo.
        
        Si no se especifica modo, selecciona automaticamente.
        """
        if circuit_id not in self.proving_keys:
            raise ValueError(f"Circuit not loaded: {circuit_id}")
        
        # Seleccionar modo si no se especifica
        if mode is None:
            mode = self._select_optimal_mode(witness)
        
        if mode not in self.provers:
            raise ValueError(f"Prover not available for mode: {mode}")
        
        prover = self.provers[mode]
        proving_key = self.proving_keys[circuit_id]
        
        # Generar prueba
        proof = prover.generate_proof(proving_key, witness)
        
        # Actualizar benchmarks
        if circuit_id not in self.benchmarks:
            self.benchmarks[circuit_id] = []
        self.benchmarks[circuit_id].append(proof.proving_time_ms)
        
        return proof
    
    def _select_optimal_mode(self, witness: Witness) -> ProvingMode:
        """Seleccionar modo optimo basado en witness"""
        # Si tiene datos muy sensibles, preferir client
        if witness.is_sensitive():
            if ProvingMode.CLIENT in self.provers:
                return ProvingMode.CLIENT
            elif ProvingMode.HYBRID in self.provers:
                return ProvingMode.HYBRID
        
        # Si no hay datos sensibles, server es mas rapido
        if ProvingMode.SERVER in self.provers:
            return ProvingMode.SERVER
        
        # Default a client
        return ProvingMode.CLIENT
    
    def verify_proof(
        self,
        circuit_id: str,
        proof: Proof
    ) -> bool:
        """Verificar prueba"""
        if circuit_id not in self.verification_keys:
            raise ValueError(f"Verification key not found: {circuit_id}")
        
        vk = self.verification_keys[circuit_id]
        return vk.verify_proof(proof.proof_data, proof.public_inputs)
    
    def get_benchmarks(self) -> Dict[str, Dict]:
        """Obtener estadisticas de benchmarks"""
        stats = {}
        for circuit_id, times in self.benchmarks.items():
            if times:
                stats[circuit_id] = {
                    "count": len(times),
                    "avg_ms": sum(times) / len(times),
                    "min_ms": min(times),
                    "max_ms": max(times)
                }
        return stats


# Servicios de proving externos
@dataclass
class ExternalProvingService:
    """Servicio externo de proving"""
    name: str
    url: str
    supported_systems: List[ProofSystem]
    pricing_per_proof: float  # USD
    avg_latency_ms: float
    trusted: bool  # Si el servicio ve datos privados

PROVING_SERVICES = [
    ExternalProvingService(
        name="Bonsai (RISC Zero)",
        url="https://bonsai.risc0.dev",
        supported_systems=[ProofSystem.STARK],
        pricing_per_proof=0.10,
        avg_latency_ms=5000,
        trusted=True  # Execution encapsulada
    ),
    ExternalProvingService(
        name="Sindri",
        url="https://sindri.app",
        supported_systems=[ProofSystem.GROTH16, ProofSystem.PLONK],
        pricing_per_proof=0.05,
        avg_latency_ms=3000,
        trusted=False  # Ve witness
    ),
    ExternalProvingService(
        name="Succinct",
        url="https://succinct.xyz",
        supported_systems=[ProofSystem.PLONK, ProofSystem.HALO2],
        pricing_per_proof=0.08,
        avg_latency_ms=4000,
        trusted=False
    ),
]


# Demo
def proving_demo():
    print("=== Proving Infrastructure Demo ===\n")
    
    # Crear servicio de proving
    service = ProvingService()
    
    # Registrar provers
    client_prover = ClientProver()
    server_prover = ServerProver(
        server_url="https://prover.example.com",
        api_key="demo_key"
    )
    hybrid_prover = HybridProver(client_prover, server_prover)
    
    service.register_prover(client_prover)
    service.register_prover(server_prover)
    service.register_prover(hybrid_prover)
    
    # Cargar circuito
    proving_key = ProvingKey(
        circuit_id="withdrawal",
        proof_system=ProofSystem.GROTH16,
        key_data=secrets.token_bytes(1024),
        size_bytes=1024,
        created_at=datetime.now()
    )
    
    verification_key = VerificationKey(
        circuit_id="withdrawal",
        proof_system=ProofSystem.GROTH16,
        key_data=secrets.token_bytes(256),
        commitment="0xvk_commitment"
    )
    
    service.load_circuit("withdrawal", proving_key, verification_key)
    print("1. Circuito 'withdrawal' cargado")
    
    # Crear witness
    witness = Witness(
        circuit_id="withdrawal",
        private_inputs={
            "nullifier_secret": "0x" + secrets.token_hex(32),
            "amount": 1000000000000000000,
            "merkle_path": [1, 2, 3, 4, 5]
        },
        public_inputs=[
            int("0x1234", 16),  # root
            int("0xabcd", 16),  # nullifier
            int("0x9999", 16),  # recipient
        ]
    )
    print(f"\n2. Witness creado (sensible: {witness.is_sensitive()})")
    
    # Generar pruebas con diferentes modos
    modes = [ProvingMode.CLIENT, ProvingMode.SERVER, ProvingMode.HYBRID]
    
    print("\n3. Generando pruebas en diferentes modos:")
    for mode in modes:
        proof = service.generate_proof("withdrawal", witness, mode)
        print(f"\n   {mode.value.upper()}:")
        print(f"     Tiempo: {proof.proving_time_ms:.2f} ms")
        print(f"     Tamaño: {proof.size_bytes()} bytes")
        print(f"     Sistema: {proof.proof_system.value}")
    
    # Verificar
    proof = service.generate_proof("withdrawal", witness, ProvingMode.CLIENT)
    is_valid = service.verify_proof("withdrawal", proof)
    print(f"\n4. Verificacion: {'VALID' if is_valid else 'INVALID'}")
    
    # Benchmarks
    print("\n5. Benchmarks:")
    stats = service.get_benchmarks()
    for circuit, data in stats.items():
        print(f"   {circuit}: avg={data['avg_ms']:.2f}ms, count={data['count']}")
    
    # Servicios externos
    print("\n6. Servicios de proving externos:")
    for svc in PROVING_SERVICES:
        trust = "trusted" if svc.trusted else "untrusted"
        print(f"   {svc.name}: ${svc.pricing_per_proof}/proof, {svc.avg_latency_ms}ms ({trust})")

proving_demo()
\end{lstlisting}

\begin{table}[htbp]
\centering
\small
\begin{tabular}{lccccc}
\toprule
\textbf{Modo} & \textbf{Privacidad} & \textbf{Velocidad} & \textbf{Costo} & \textbf{Complejidad} \\
\midrule
Client-side & Máxima & Lenta & Gratis & Baja \\
Server-side & Baja & Rápida & \$/proof & Baja \\
Híbrido & Alta & Media & Parcial & Alta \\
MPC Delegated & Alta & Media & Alto & Muy alta \\
\bottomrule
\end{tabular}
\caption{Comparativa de modos de proving}
\end{table}

%===========================================
\section{Hardware Wallets y ZK}
\label{sec:hardware-zk}
%===========================================

\begin{definition}[Hardware Wallet ZK]
Dispositivo seguro que almacena claves privadas y potencialmente participa en operaciones ZK:
\begin{itemize}
    \item \textbf{Almacenamiento seguro:} Keys nunca salen del dispositivo
    \item \textbf{Firma segura:} Firmas generadas internamente
    \item \textbf{ZK limitado:} Capacidad de cómputo restringida
\end{itemize}
\end{definition}

\subsection{Limitaciones Actuales}

\begin{figure}[htbp]
\centering
\begin{tikzpicture}[
    hw/.style={rectangle, draw, rounded corners, minimum width=3cm, minimum height=1.5cm, align=center, font=\small, fill=gray!20},
    check/.style={font=\scriptsize, green!60!black},
    cross/.style={font=\scriptsize, red}
]
    % Hardware wallet
    \node[hw] (hw) at (0, 0) {\textbf{Hardware}\\Wallet};
    
    % Capacidades
    \node[check] at (4, 1.5) {\checkmark\ Almacenar seeds};
    \node[check] at (4, 1) {\checkmark\ Firmar transacciones};
    \node[check] at (4, 0.5) {\checkmark\ Almacenar notes (encriptadas)};
    \node[cross] at (4, -0.2) {$\times$\ Generar ZK proofs};
    \node[cross] at (4, -0.7) {$\times$\ Witness computation pesada};
    \node[cross] at (4, -1.2) {$\times$\ MSM/FFT grandes};
    
    % Specs tipicas
    \node[font=\tiny, align=left] at (0, -2.5) {
        Specs típicas:\\
        - CPU: ARM Cortex-M4\\
        - RAM: 256KB - 1MB\\
        - Storage: 2-4MB
    };
\end{tikzpicture}
\caption{Capacidades y limitaciones de hardware wallets actuales}
\end{figure}

\subsection{Estrategias de Integración}

\begin{lstlisting}[style=python, caption={Integración de hardware wallet con ZK}]
"""
Hardware Wallet ZK Integration
Estrategias para usar hardware wallets con sistemas ZK
"""

from dataclasses import dataclass, field
from typing import Optional, Dict, List, Any, Tuple
from enum import Enum
from abc import ABC, abstractmethod
import hashlib
import secrets

class HardwareCapability(Enum):
    """Capacidades de hardware wallet"""
    SIGN_TX = "sign_tx"           # Firmar transacciones
    SIGN_MESSAGE = "sign_message" # Firmar mensajes
    DERIVE_KEY = "derive_key"     # Derivar claves
    STORE_DATA = "store_data"     # Almacenar datos
    COMPUTE_HASH = "compute_hash" # Calcular hashes
    BASIC_CRYPTO = "basic_crypto" # Operaciones crypto basicas

@dataclass
class HardwareWalletSpec:
    """Especificaciones de hardware wallet"""
    model: str
    cpu: str
    ram_kb: int
    storage_kb: int
    capabilities: List[HardwareCapability]
    max_operation_ms: int  # Timeout maximo
    
    def can_generate_proof(self, circuit_constraints: int) -> bool:
        """Verificar si puede generar prueba para circuito"""
        # Estimacion: ~1KB RAM por 100 constraints
        required_ram = circuit_constraints * 10  # bytes
        return required_ram < self.ram_kb * 1024

# Dispositivos comunes
LEDGER_NANO_X = HardwareWalletSpec(
    model="Ledger Nano X",
    cpu="ST33J2M0",
    ram_kb=256,
    storage_kb=2048,
    capabilities=[
        HardwareCapability.SIGN_TX,
        HardwareCapability.SIGN_MESSAGE,
        HardwareCapability.DERIVE_KEY,
        HardwareCapability.COMPUTE_HASH
    ],
    max_operation_ms=30000
)

TREZOR_MODEL_T = HardwareWalletSpec(
    model="Trezor Model T",
    cpu="ARM Cortex-M4",
    ram_kb=256,
    storage_kb=4096,
    capabilities=[
        HardwareCapability.SIGN_TX,
        HardwareCapability.SIGN_MESSAGE,
        HardwareCapability.DERIVE_KEY,
        HardwareCapability.STORE_DATA
    ],
    max_operation_ms=60000
)


class HardwareWallet(ABC):
    """Interfaz abstracta para hardware wallet"""
    
    @abstractmethod
    def sign_transaction(self, tx_hash: bytes) -> bytes:
        pass
    
    @abstractmethod
    def derive_key(self, path: str) -> bytes:
        pass
    
    @abstractmethod
    def get_public_key(self, path: str) -> bytes:
        pass


class SimulatedHardwareWallet(HardwareWallet):
    """Hardware wallet simulado para desarrollo"""
    
    def __init__(self, spec: HardwareWalletSpec, seed: bytes):
        self.spec = spec
        self.seed = seed
        self.derived_keys: Dict[str, bytes] = {}
        
    def sign_transaction(self, tx_hash: bytes) -> bytes:
        """Firmar hash de transaccion"""
        # Simular firma ECDSA
        signature = hashlib.sha256(
            self.seed + tx_hash
        ).digest() + hashlib.sha256(
            tx_hash + self.seed
        ).digest()
        return signature[:64]  # 64 bytes signature
    
    def derive_key(self, path: str) -> bytes:
        """Derivar clave desde path BIP32"""
        if path not in self.derived_keys:
            key_material = hashlib.sha256(
                self.seed + path.encode()
            ).digest()
            self.derived_keys[path] = key_material
        return self.derived_keys[path]
    
    def get_public_key(self, path: str) -> bytes:
        """Obtener clave publica"""
        private = self.derive_key(path)
        # Simular derivacion de pubkey
        return hashlib.sha256(private).digest()


@dataclass
class ZKDelegationRequest:
    """Request para delegacion segura de proving"""
    nullifier_commitment: str     # H(nullifier_secret)
    amount_commitment: str        # H(amount, blinding)
    merkle_path_commitment: str   # H(path)
    public_inputs: List[int]
    
    # NO incluye secretos reales

@dataclass
class ZKDelegationProof:
    """Prueba de delegacion"""
    proof_data: bytes
    hw_signature: bytes  # Firma del hardware wallet
    delegation_commitment: str


class SecureZKDelegation:
    """
    Sistema de delegacion segura de proving.
    
    Permite que un servidor genere pruebas sin ver
    los secretos del usuario.
    
    Protocolo:
    1. HW genera commitments a secretos
    2. HW firma los commitments
    3. Servidor genera prueba sobre commitments
    4. HW verifica y autoriza
    """
    
    def __init__(
        self,
        hardware_wallet: HardwareWallet,
        server_url: str
    ):
        self.hw = hardware_wallet
        self.server_url = server_url
        
    def create_delegation_request(
        self,
        nullifier_secret: bytes,
        amount: int,
        merkle_path: List[int]
    ) -> ZKDelegationRequest:
        """
        Crear request de delegacion.
        
        IMPORTANTE: Los secretos nunca salen del cliente.
        Solo se envian commitments.
        """
        # Generar commitments
        nullifier_commitment = hashlib.sha256(nullifier_secret).hexdigest()
        
        blinding = secrets.token_bytes(32)
        amount_commitment = hashlib.sha256(
            amount.to_bytes(32, 'big') + blinding
        ).hexdigest()
        
        merkle_path_commitment = hashlib.sha256(
            b"".join(p.to_bytes(32, 'big') for p in merkle_path)
        ).hexdigest()
        
        # Calcular public inputs
        # (estos si se revelan)
        nullifier_hash = hashlib.sha256(
            nullifier_secret + b"nullifier"
        ).digest()
        
        public_inputs = [
            int.from_bytes(nullifier_hash[:16], 'big'),  # nullifier
            # merkle root, recipient, etc (simplificado)
        ]
        
        return ZKDelegationRequest(
            nullifier_commitment=nullifier_commitment,
            amount_commitment=amount_commitment,
            merkle_path_commitment=merkle_path_commitment,
            public_inputs=public_inputs
        )
    
    def authorize_delegation(
        self,
        request: ZKDelegationRequest
    ) -> bytes:
        """
        Hardware wallet autoriza la delegacion.
        
        Firma los commitments para que el servidor
        sepa que son validos.
        """
        # Crear mensaje a firmar
        message = (
            request.nullifier_commitment +
            request.amount_commitment +
            request.merkle_path_commitment
        ).encode()
        
        message_hash = hashlib.sha256(message).digest()
        
        # Firmar con hardware wallet
        signature = self.hw.sign_transaction(message_hash)
        
        return signature
    
    def request_proof(
        self,
        request: ZKDelegationRequest,
        hw_signature: bytes
    ) -> ZKDelegationProof:
        """
        Solicitar prueba al servidor.
        
        El servidor genera la prueba usando los commitments
        sin conocer los valores reales.
        """
        # En produccion: HTTP request al servidor
        
        # Simular respuesta
        proof_data = secrets.token_bytes(256)
        
        return ZKDelegationProof(
            proof_data=proof_data,
            hw_signature=hw_signature,
            delegation_commitment=request.nullifier_commitment
        )
    
    def verify_and_submit(
        self,
        proof: ZKDelegationProof,
        original_request: ZKDelegationRequest
    ) -> bool:
        """
        Verificar prueba y preparar para submit.
        
        El hardware wallet verifica que la prueba
        corresponde a los commitments originales.
        """
        # Verificar que la prueba usa nuestros commitments
        if proof.delegation_commitment != original_request.nullifier_commitment:
            return False
        
        # Verificar firma (ya la tenemos)
        # En produccion: verificar criptograficamente
        
        return True


class FutureHardwareZK:
    """
    Especulacion sobre futuro hardware ZK.
    
    Hardware especializado para ZK proving.
    """
    
    @staticmethod
    def projected_specs() -> Dict[str, Any]:
        return {
            "zk_accelerator": {
                "description": "ASIC/FPGA para MSM/NTT",
                "speedup": "100-1000x vs CPU",
                "timeline": "2025-2026"
            },
            "secure_element_zk": {
                "description": "Secure element con soporte ZK basico",
                "capabilities": ["small_circuits", "commitment_generation"],
                "timeline": "2024-2025"
            },
            "tee_proving": {
                "description": "TEE (SGX/TrustZone) para proving",
                "benefit": "Privacidad server-side",
                "timeline": "Available now (limited)"
            },
            "mpc_hardware": {
                "description": "Hardware optimizado para MPC",
                "use_case": "Distributed proving",
                "timeline": "2025+"
            }
        }


# Demo
def hardware_zk_demo():
    print("=== Hardware Wallet ZK Demo ===\n")
    
    # Mostrar specs de dispositivos
    print("1. Especificaciones de hardware wallets:")
    for spec in [LEDGER_NANO_X, TREZOR_MODEL_T]:
        print(f"\n   {spec.model}:")
        print(f"     CPU: {spec.cpu}")
        print(f"     RAM: {spec.ram_kb}KB")
        print(f"     Storage: {spec.storage_kb}KB")
        
        # Verificar capacidad ZK
        can_small = spec.can_generate_proof(1000)
        can_large = spec.can_generate_proof(100000)
        print(f"     ZK (1K constraints): {'Si' if can_small else 'No'}")
        print(f"     ZK (100K constraints): {'Si' if can_large else 'No'}")
    
    # Crear hardware wallet simulado
    print("\n2. Simulando delegacion segura de proving...")
    
    hw = SimulatedHardwareWallet(
        spec=LEDGER_NANO_X,
        seed=secrets.token_bytes(32)
    )
    
    delegation = SecureZKDelegation(
        hardware_wallet=hw,
        server_url="https://prover.example.com"
    )
    
    # Crear request (secretos se quedan en cliente)
    nullifier_secret = secrets.token_bytes(32)
    request = delegation.create_delegation_request(
        nullifier_secret=nullifier_secret,
        amount=10**18,
        merkle_path=[1, 2, 3, 4, 5]
    )
    
    print(f"   Nullifier commitment: {request.nullifier_commitment[:20]}...")
    print(f"   Amount commitment: {request.amount_commitment[:20]}...")
    print("   (Secretos nunca salen del cliente)")
    
    # Autorizar con HW
    print("\n3. Hardware wallet autoriza delegacion...")
    signature = delegation.authorize_delegation(request)
    print(f"   Firma: {signature[:20].hex()}...")
    
    # Request proof
    print("\n4. Servidor genera prueba (sin ver secretos)...")
    proof = delegation.request_proof(request, signature)
    print(f"   Proof size: {len(proof.proof_data)} bytes")
    
    # Verificar
    print("\n5. Verificando prueba...")
    is_valid = delegation.verify_and_submit(proof, request)
    print(f"   Valida: {'Si' if is_valid else 'No'}")
    
    # Futuro
    print("\n6. Futuro del hardware ZK:")
    future = FutureHardwareZK.projected_specs()
    for name, info in future.items():
        print(f"\n   {name}:")
        print(f"     {info['description']}")
        print(f"     Timeline: {info['timeline']}")

hardware_zk_demo()
\end{lstlisting}

%===========================================
\section*{Resumen del Bloque}
%===========================================

\begin{keyidea}
\textbf{Ecosistema de Privacidad - Componentes Esenciales:}

\textbf{1. Exploradores Privacy-Aware:}
\begin{itemize}
    \item Sin tracking (cookies, analytics, fingerprinting)
    \item Acceso via Tor (.onion services)
    \item Dummy queries para confusión estadística
    \item Self-hosting recomendado para máxima privacidad
    \item Índices locales sin dependencias externas
\end{itemize}

\textbf{2. Wallets ZK:}
\begin{itemize}
    \item Railgun: ERC-20 privacy, DeFi integration, shielded pools
    \item Umbra: Stealth payments simples, bajo overhead
    \item Zcash: Privacidad nativa (Zashi, YWallet, Nighthawk)
    \item Gestión segura de viewing keys y spending keys
    \item Derivación determinística compatible con hardware wallets
\end{itemize}

\textbf{3. RPC Privados:}
\begin{itemize}
    \item Flashbots Protect: No mempool pública, bundle submission
    \item MEV Blocker: OFA (Order Flow Auctions) con refunds
    \item Protección contra frontrunning, sandwich attacks, y liquidation sniping
    \item Private mempools con commit-reveal schemes
    \item Encrypted transaction propagation
\end{itemize}

\textbf{4. Relayer Networks:}
\begin{itemize}
    \item Descentralización para censorship resistance
    \item Competencia económica reduce fees
    \item Staking y slashing para accountability
    \item Rotación automática para unlinkability
    \item Múltiples jurisdicciones para resiliencia legal
\end{itemize}

\textbf{5. Note Management:}
\begin{itemize}
    \item Backup ANTES de depositar (regla de oro)
    \item Encriptación AES-256-GCM con key derivation segura
    \item Split backup con Shamir Secret Sharing (2-of-3 mínimo)
    \item Sincronización segura entre dispositivos
    \item Recovery testing periódico obligatorio
\end{itemize}

\textbf{6. Proving Infrastructure:}
\begin{itemize}
    \item Client-side (WASM/native): Máxima privacidad, witness nunca sale
    \item Server-side: Máxima velocidad, requiere confianza
    \item Híbrido: Balance óptimo con witness encryption
    \item GPU acceleration para proving intensivo
    \item Proof caching y batching para eficiencia
\end{itemize}

\textbf{7. Hardware Wallets:}
\begin{itemize}
    \item Limitaciones actuales: memoria insuficiente para ZK proving
    \item Delegación segura con commit-reveal posible
    \item Signing de transacciones blindadas
    \item Futuro: Hardware ZK especializado (FPGAs, ASICs)
    \item Integración con secure elements para key storage
\end{itemize}
\end{keyidea}

\begin{bancozk}
\textbf{Stack Completo de Banco ZK:}

\begin{center}
\begin{tikzpicture}[
    layer/.style={rectangle, draw, minimum width=12cm, minimum height=1cm, align=center, font=\small},
    note/.style={font=\tiny, align=left}
]
    % Stack completo
    \node[layer, fill=blue!20] (ui) at (0, 6) {\textbf{UI Layer:} Web App + Mobile (PWA) + Desktop};
    \node[layer, fill=cyan!20] (wallet) at (0, 5) {\textbf{Wallet Layer:} Note Manager + Backup System + Key Derivation};
    \node[layer, fill=green!20] (proving) at (0, 4) {\textbf{Proving Layer:} Client WASM + Hybrid Delegation + GPU Accel};
    \node[layer, fill=yellow!20] (privacy) at (0, 3) {\textbf{Privacy Layer:} Tor + Mixnet + Dummy Traffic};
    \node[layer, fill=orange!20] (network) at (0, 2) {\textbf{Network Layer:} Relayer Network + P2P Gossip};
    \node[layer, fill=red!20] (rpc) at (0, 1) {\textbf{RPC Layer:} Flashbots Protect + Private Mempool};
    \node[layer, fill=purple!20] (chain) at (0, 0) {\textbf{Chain Layer:} ZK Contracts + Merkle Trees + Nullifier Sets};
    
    % Notas laterales
    \node[note, right] at (6.5, 6) {React/Vue, Capacitor};
    \node[note, right] at (6.5, 5) {IndexedDB, Encrypted};
    \node[note, right] at (6.5, 4) {snarkjs, rapidsnark};
    \node[note, right] at (6.5, 3) {Arti, I2P optional};
    \node[note, right] at (6.5, 2) {libp2p, stake-weighted};
    \node[note, right] at (6.5, 1) {MEV protection};
    \node[note, right] at (6.5, 0) {Solidity, Circom};
\end{tikzpicture}
\end{center}

\textbf{Principios de Diseño:}
\begin{enumerate}
    \item \textbf{Privacy by default:} Toda operación es privada, transparencia es opt-in
    \item \textbf{Redundancia:} Múltiples relayers, backups distribuidos, failover automático
    \item \textbf{Verificabilidad:} Todo open source, builds reproducibles, auditable on-chain
    \item \textbf{Usabilidad:} Complejidad criptográfica oculta al usuario final
    \item \textbf{Descentralización:} Sin single points of failure, governance distribuida
    \item \textbf{Composabilidad:} Integración con DeFi existente vía adapters
    \item \textbf{Upgradeability:} Migración segura de circuitos y contratos
\end{enumerate}

\textbf{Métricas Objetivo:}
\begin{center}
\begin{tabular}{lcc}
\toprule
\textbf{Métrica} & \textbf{Objetivo} & \textbf{Crítico} \\
\midrule
Tiempo de depósito & $< 30$s & $< 60$s \\
Tiempo de retiro & $< 45$s & $< 90$s \\
Anonymity set & $> 10,000$ & $> 1,000$ \\
Uptime relayers & $> 99.9\%$ & $> 99\%$ \\
Recovery success & $100\%$ & $100\%$ \\
MEV leakage & $0\%$ & $< 1\%$ \\
False positive rate & $< 0.1\%$ & $< 1\%$ \\
Gas overhead vs EOA & $< 5x$ & $< 10x$ \\
\bottomrule
\end{tabular}
\end{center}

\textbf{Recomendaciones de Implementación:}
\begin{enumerate}
    \item Comenzar con client-side proving (privacidad máxima)
    \item Integrar Flashbots Protect desde día 1
    \item Implementar backup obligatorio antes de depósitos
    \item Usar Tor para todas las conexiones por defecto
    \item Auditar código antes de mainnet (mínimo 2 auditores)
    \item Programa de bug bounty activo desde testnet
    \item Monitoring de anonymity set en tiempo real
    \item Alertas automáticas para comportamiento anómalo
    \item Documentación completa para usuarios y desarrolladores
    \item Proceso de incident response documentado
\end{enumerate}

\textbf{Roadmap Sugerido:}
\begin{enumerate}
    \item \textbf{Fase 1 (MVP):} Depósitos/retiros básicos, single relayer, backup manual
    \item \textbf{Fase 2 (Beta):} Multi-relayer, backup automático, mobile app
    \item \textbf{Fase 3 (Mainnet):} DeFi integrations, hardware wallet support
    \item \textbf{Fase 4 (Mature):} Cross-chain bridges, institutional features
\end{enumerate}
\end{bancozk}

\begin{keyidea}
\textbf{Conclusión de la Serie Banco ZK:}

A lo largo de estos 20 bloques, hemos construido el conocimiento necesario para entender y desarrollar sistemas de privacidad en blockchain:

\begin{enumerate}
    \item \textbf{Fundamentos Matemáticos:} 
    \begin{itemize}
        \item Teoría de grupos y campos finitos
        \item Curvas elípticas (secp256k1, BN254, BLS12-381)
        \item Funciones hash criptográficas (Poseidon, Pedersen)
        \item Compromisos y esquemas de firma
    \end{itemize}
    
    \item \textbf{Zero-Knowledge Proofs:} 
    \begin{itemize}
        \item SNARKs: Groth16, PLONK, Marlin
        \item STARKs: Transparencia y post-quantum security
        \item Circuitos aritméticos y R1CS
        \item Sistemas de prueba interactivos y no interactivos
    \end{itemize}
    
    \item \textbf{Protocolos de Privacidad:} 
    \begin{itemize}
        \item Mixers y coin joins
        \item Commitment schemes y nullifiers
        \item Ring signatures y stealth addresses
        \item Confidential transactions
    \end{itemize}
    
    \item \textbf{Vectores de Ataque:} 
    \begin{itemize}
        \item Análisis de grafos y clustering
        \item Timing y traffic analysis
        \item Sybil attacks en anonymity sets
        \item Side-channel attacks en proving
    \end{itemize}
    
    \item \textbf{Infraestructura Práctica:} 
    \begin{itemize}
        \item Wallets y key management
        \item Relayer networks descentralizados
        \item Proving infrastructure escalable
        \item Integración con ecosistema existente
    \end{itemize}
\end{enumerate}

\textbf{El futuro de la privacidad en blockchain} depende de:
\begin{itemize}
    \item \textbf{Eficiencia:} Reducción de proving time de minutos a segundos
    \item \textbf{Hardware:} GPUs, FPGAs, y ASICs especializados para ZK
    \item \textbf{Estándares:} Interoperabilidad entre sistemas de privacidad
    \item \textbf{Educación:} Comprensión pública de la importancia de privacidad
    \item \textbf{Regulación:} Balance entre compliance y derechos fundamentales
    \item \textbf{Usabilidad:} Privacidad invisible para el usuario promedio
    \item \textbf{Investigación:} Nuevos sistemas de prueba más eficientes
\end{itemize}

\textbf{Desafíos Pendientes:}
\begin{itemize}
    \item Escalabilidad de anonymity sets a millones de usuarios
    \item Composabilidad completa con DeFi sin leakage
    \item Resistencia a computación cuántica (transición a STARKs)
    \item Compliance selectivo sin comprometer privacidad
    \item UX comparable a aplicaciones tradicionales
\end{itemize}

\vspace{0.5cm}
\begin{center}
\textit{``Privacy is not about having something to hide.}\\
\textit{Privacy is about having the right to choose what to reveal.''}\\
--- Anónimo
\end{center}

\vspace{0.5cm}
\begin{center}
\textit{``Arguing that you don't care about privacy because you have nothing to hide}\\
\textit{is no different than saying you don't care about free speech}\\
\textit{because you have nothing to say.''}\\
--- Edward Snowden
\end{center}
\end{keyidea}

\cleardoublepage

% ============================================
% CONCLUSIÓN GENERAL DEL VOLUMEN 2
% ============================================
\chapter*{Conclusión}
\addcontentsline{toc}{chapter}{Conclusión}

Este volumen ha cubierto los \textbf{Bloques 14 a 20} de la serie Banco ZK, enfocándose en aplicaciones avanzadas, seguridad, y el ecosistema de herramientas necesario para implementar sistemas de privacidad en blockchain.

\section*{Resumen por Bloque}

\subsection*{Bloque 14: Identidad y Credenciales Anónimas}
Exploramos cómo demostrar propiedades sobre uno mismo sin revelar información innecesaria:
\begin{itemize}
    \item Anonymous credentials y selective disclosure
    \item Pruebas de membresía en grupos (group signatures)
    \item Verificación de edad sin revelar fecha de nacimiento
    \item DIDs (Decentralized Identifiers) y VCs (Verifiable Credentials)
    \item Soulbound tokens con privacidad
    \item KYC preservando privacidad mediante ZK proofs
\end{itemize}

\subsection*{Bloque 15: Compliance vs Privacidad}
Analizamos el delicado balance entre requisitos regulatorios y privacidad del usuario:
\begin{itemize}
    \item Selective disclosure para auditorías
    \item Viewing keys y su gestión segura
    \item Travel Rule compliance con ZK
    \item Sanctions screening sin revelar identidad
    \item Arquitecturas de compliance privado
    \item Casos de estudio: Zcash, Tornado Cash, Railgun
\end{itemize}

\subsection*{Bloque 16: Metadata y Network Privacy}
Profundizamos en la protección contra análisis de metadatos:
\begin{itemize}
    \item Timing analysis y sus mitigaciones
    \item Traffic analysis en redes P2P
    \item Tor, I2P, y mixnets para privacidad de red
    \item Dummy transactions y cover traffic
    \item Fingerprinting de wallets y protecciones
    \item Análisis de grafos y técnicas de clustering
\end{itemize}

\subsection*{Bloque 17: Private Smart Contracts}
Cubrimos la ejecución privada de lógica arbitraria:
\begin{itemize}
    \item zkEVM y sus variantes (Type 1-4)
    \item Private state en smart contracts
    \item Aztec, Aleo, y Mina: comparativa
    \item Composabilidad entre contratos privados
    \item Oráculos privados y feeds de datos
    \item Private DeFi: AMMs, lending, derivatives
\end{itemize}

\subsection*{Bloque 18: Criptografía Avanzada}
Exploramos primitivas criptográficas de frontera:
\begin{itemize}
    \item Fully Homomorphic Encryption (FHE)
    \item Multi-Party Computation (MPC)
    \item Oblivious Transfer y Private Set Intersection
    \item Witness Encryption
    \item Post-quantum cryptography
    \item Combinación de técnicas para máxima seguridad
\end{itemize}

\subsection*{Bloque 19: Ataques y Defensa}
Analizamos vectores de ataque y estrategias defensivas:
\begin{itemize}
    \item Ataques al anonymity set (Sybil, poisoning)
    \item Deanonymization mediante análisis heurístico
    \item Side-channel attacks en proving
    \item Smart contract vulnerabilities específicas de privacidad
    \item Incident response y recovery
    \item Diseño de sistemas resilientes
\end{itemize}

\subsection*{Bloque 20: Herramientas y Ecosistema}
Presentamos el stack completo para implementación:
\begin{itemize}
    \item Exploradores privacy-aware
    \item Wallets con soporte ZK (Railgun, Umbra, Zcash)
    \item RPCs privados (Flashbots, MEV Blocker)
    \item Redes de relayers descentralizados
    \item Gestión de notas y backups
    \item Infraestructura de proving
    \item Integración con hardware wallets
\end{itemize}

\section*{Conexiones con Volumen 1}

El Volumen 1 (Bloques 1-13) estableció los fundamentos:
\begin{itemize}
    \item \textbf{Bloques 1-3:} Criptografía básica, curvas elípticas, hashing
    \item \textbf{Bloques 4-6:} Estructuras de datos (Merkle trees), compromisos
    \item \textbf{Bloques 7-9:} Firmas digitales, ring signatures, stealth addresses
    \item \textbf{Bloques 10-13:} Zero-knowledge proofs (SNARKs, STARKs, circuitos)
\end{itemize}

Este Volumen 2 construye sobre esos fundamentos para aplicaciones prácticas.

\section*{Implementación Recomendada}

Para implementar un sistema como Banco ZK, sugerimos:

\begin{enumerate}
    \item \textbf{Fase de Diseño (2-3 meses):}
    \begin{itemize}
        \item Definir modelo de amenazas
        \item Seleccionar primitivas criptográficas
        \item Diseñar arquitectura de circuitos
        \item Especificar interfaces de usuario
    \end{itemize}
    
    \item \textbf{Fase de Desarrollo (4-6 meses):}
    \begin{itemize}
        \item Implementar contratos core
        \item Desarrollar circuitos ZK
        \item Construir infraestructura de relayers
        \item Crear aplicaciones cliente
    \end{itemize}
    
    \item \textbf{Fase de Testing (2-3 meses):}
    \begin{itemize}
        \item Auditorías de seguridad (mínimo 2)
        \item Testnet público
        \item Bug bounty program
        \item Stress testing
    \end{itemize}
    
    \item \textbf{Fase de Lanzamiento:}
    \begin{itemize}
        \item Mainnet con límites iniciales
        \item Monitoreo 24/7
        \item Proceso de escalado gradual
        \item Documentación y soporte
    \end{itemize}
\end{enumerate}

\section*{Recursos Adicionales}

Para profundizar en los temas cubiertos:

\textbf{Documentación Técnica:}
\begin{itemize}
    \item Zcash Protocol Specification
    \item Tornado Cash documentation
    \item Railgun technical papers
    \item Aztec protocol documentation
\end{itemize}

\textbf{Papers Académicos:}
\begin{itemize}
    \item ``Zerocash: Decentralized Anonymous Payments'' (Sasson et al.)
    \item ``PLONK: Permutations over Lagrange-bases'' (Gabizon et al.)
    \item ``STARKs'' (Ben-Sasson et al.)
    \item ``Bulletproofs'' (Bünz et al.)
\end{itemize}

\textbf{Herramientas de Desarrollo:}
\begin{itemize}
    \item Circom + snarkjs
    \item Noir (Aztec)
    \item Leo (Aleo)
    \item Cairo (StarkWare)
\end{itemize}

\cleardoublepage

% ============================================
% GLOSARIO EXPANDIDO
% ============================================
\chapter*{Glosario}
\addcontentsline{toc}{chapter}{Glosario}

\section*{A}
\begin{description}
\item[AES] Advanced Encryption Standard. Algoritmo de cifrado simétrico estándar.
\item[AMM] Automated Market Maker. Protocolo DeFi para intercambio sin order book.
\item[Anonymity Set] Conjunto de usuarios indistinguibles entre sí para un observador.
\item[ASIC] Application-Specific Integrated Circuit. Hardware especializado.
\end{description}

\section*{B}
\begin{description}
\item[BFT] Byzantine Fault Tolerance. Tolerancia a fallos bizantinos en consenso.
\item[BLS] Boneh-Lynn-Shacham. Esquema de firmas con agregación eficiente.
\item[BN254] Barreto-Naehrig curve. Curva elíptica pairing-friendly de 254 bits.
\item[Bulletproofs] Sistema de range proofs sin trusted setup.
\end{description}

\section*{C}
\begin{description}
\item[CRS] Common Reference String. String de referencia para SNARKs con trusted setup.
\item[Commitment] Compromiso criptográfico que oculta un valor hasta su revelación.
\item[Cover Traffic] Tráfico ficticio para ocultar patrones reales.
\item[Circom] Lenguaje DSL para definir circuitos aritméticos ZK.
\end{description}

\section*{D}
\begin{description}
\item[DA] Data Availability. Disponibilidad de datos en sistemas L2.
\item[DID] Decentralized Identifier. Identificador descentralizado W3C.
\item[DKG] Distributed Key Generation. Generación distribuida de claves.
\item[DSL] Domain-Specific Language. Lenguaje especializado para un dominio.
\item[Dummy Query] Consulta ficticia para ocultar consultas reales.
\end{description}

\section*{E}
\begin{description}
\item[ECDLP] Elliptic Curve Discrete Log Problem. Problema base de seguridad ECC.
\item[ECDSA] Elliptic Curve Digital Signature Algorithm. Firma estándar en Bitcoin/Ethereum.
\item[EdDSA] Edwards-curve Digital Signature Algorithm. Firma en curvas Edwards.
\item[EIP] Ethereum Improvement Proposal. Propuesta de mejora de Ethereum.
\item[EVM] Ethereum Virtual Machine. Máquina virtual de Ethereum.
\end{description}

\section*{F}
\begin{description}
\item[FHE] Fully Homomorphic Encryption. Cifrado que permite cómputo sobre datos cifrados.
\item[FFT] Fast Fourier Transform. Algoritmo para evaluación rápida de polinomios.
\item[FPGA] Field-Programmable Gate Array. Hardware reconfigurable.
\item[FRI] Fast Reed-Solomon IOP. Protocolo de proximidad usado en STARKs.
\item[Frontrunning] Ataque que ejecuta transacciones antes de otras conocidas.
\end{description}

\section*{G}
\begin{description}
\item[Gas] Unidad de cómputo en Ethereum.
\item[Groth16] Sistema SNARK eficiente con trusted setup por circuito.
\item[Group Signature] Firma que prueba membresía sin revelar identidad.
\end{description}

\section*{H}
\begin{description}
\item[HMAC] Hash-based Message Authentication Code.
\item[HSM] Hardware Security Module. Módulo de seguridad hardware.
\item[Hybrid Proving] Combinación de proving local y delegado.
\end{description}

\section*{I}
\begin{description}
\item[IVC] Incrementally Verifiable Computation. Computación verificable incremental.
\item[I2P] Invisible Internet Project. Red de anonimato alternativa a Tor.
\end{description}

\section*{K}
\begin{description}
\item[KDF] Key Derivation Function. Función para derivar claves.
\item[KZG] Kate-Zaverucha-Goldberg. Esquema de commitment polinomial.
\item[KYC] Know Your Customer. Requisitos de identificación de clientes.
\end{description}

\section*{L}
\begin{description}
\item[L1] Layer 1. Blockchain base (ej: Ethereum mainnet).
\item[L2] Layer 2. Solución de escalabilidad sobre L1.
\item[Lookup Table] Tabla precalculada para optimizar circuitos.
\end{description}

\section*{M}
\begin{description}
\item[Merkle Tree] Árbol hash para compromisos eficientes de conjuntos.
\item[MEV] Maximal Extractable Value. Valor extraíble por reordenación de transacciones.
\item[Mixnet] Red de mezcla para anonimato de comunicaciones.
\item[MPC] Multi-Party Computation. Cómputo seguro entre múltiples partes.
\item[MSM] Multi-Scalar Multiplication. Operación crítica en SNARKs.
\end{description}

\section*{N}
\begin{description}
\item[NIZK] Non-Interactive Zero-Knowledge. ZK sin interacción.
\item[Note] Representación de valor en sistemas UTXO privados.
\item[Nullifier] Identificador único que previene doble gasto sin revelar origen.
\item[NTT] Number Theoretic Transform. FFT sobre campos finitos.
\end{description}

\section*{O}
\begin{description}
\item[OFA] Order Flow Auction. Subasta de flujo de órdenes para MEV.
\item[Onion Routing] Enrutamiento en capas (usado por Tor).
\item[OT] Oblivious Transfer. Transferencia inconsciente en MPC.
\item[Oracle] Servicio que proporciona datos externos a smart contracts.
\end{description}

\section*{P}
\begin{description}
\item[Pairing] Operación bilineal en curvas elípticas.
\item[Pedersen Hash] Hash algebraico eficiente en circuitos.
\item[PLONK] Permutations over Lagrange-bases. Sistema SNARK universal.
\item[Poseidon] Función hash optimizada para circuitos ZK.
\item[Prover] Entidad que genera pruebas ZK.
\item[PSI] Private Set Intersection. Intersección privada de conjuntos.
\item[PWA] Progressive Web App. Aplicación web instalable.
\end{description}

\section*{R}
\begin{description}
\item[R1CS] Rank-1 Constraint System. Sistema de restricciones para SNARKs.
\item[Range Proof] Prueba de que un valor está en un rango sin revelarlo.
\item[Relayer] Intermediario que envía transacciones por cuenta de usuarios.
\item[Ring Signature] Firma que oculta al firmante real entre un grupo.
\item[Rollup] Solución L2 que publica datos en L1.
\end{description}

\section*{S}
\begin{description}
\item[Sandwich Attack] Ataque MEV que rodea una transacción víctima.
\item[SBT] Soulbound Token. Token no transferible.
\item[Schnorr] Esquema de firma con agregación simple.
\item[Shamir] Shamir Secret Sharing. División de secretos con threshold.
\item[Shielded Pool] Pool de fondos privados en sistemas como Zcash.
\item[SNARK] Succinct Non-interactive ARgument of Knowledge.
\item[Spending Key] Clave que autoriza gastos en sistemas ZK.
\item[STARK] Scalable Transparent ARgument of Knowledge.
\item[Stealth Address] Dirección de un solo uso para privacidad.
\end{description}

\section*{T}
\begin{description}
\item[TEE] Trusted Execution Environment. Entorno de ejecución aislado.
\item[Tor] The Onion Router. Red de anonimato más popular.
\item[Trusted Setup] Ceremonia para generar CRS en algunos SNARKs.
\item[TWAP] Time-Weighted Average Price. Precio promedio ponderado por tiempo.
\end{description}

\section*{U}
\begin{description}
\item[Umbra] Protocolo de stealth addresses para Ethereum.
\item[UTXO] Unspent Transaction Output. Modelo de Bitcoin/Zcash.
\end{description}

\section*{V}
\begin{description}
\item[VC] Verifiable Credential. Credencial verificable W3C.
\item[Verifier] Entidad que verifica pruebas ZK.
\item[Viewing Key] Clave que permite ver transacciones sin gastar.
\item[VRF] Verifiable Random Function. Función aleatoria verificable.
\end{description}

\section*{W}
\begin{description}
\item[WASM] WebAssembly. Formato binario para ejecución en navegadores.
\item[Witness] Datos privados usados para generar pruebas ZK.
\item[Witness Encryption] Cifrado que requiere prueba ZK para descifrar.
\end{description}

\section*{Z}
\begin{description}
\item[Zcash] Criptomoneda pionera en privacidad con zk-SNARKs.
\item[zkEVM] Zero-Knowledge EVM. EVM compatible con pruebas ZK.
\item[zkML] Zero-Knowledge Machine Learning. ML con privacidad ZK.
\item[zkSNARK] Zero-Knowledge Succinct Non-interactive ARgument of Knowledge.
\item[zkSTARK] Zero-Knowledge STARK.
\item[zkVM] Zero-Knowledge Virtual Machine. VM genérica con ZK.
\item[ZKP] Zero-Knowledge Proof. Prueba de conocimiento cero.
\end{description}

\cleardoublepage

% ============================================
% ÍNDICE DE FIGURAS Y TABLAS (OPCIONAL)
% ============================================
\listoffigures
\addcontentsline{toc}{chapter}{Índice de Figuras}

\cleardoublepage

\listoftables
\addcontentsline{toc}{chapter}{Índice de Tablas}

\cleardoublepage

% ============================================
% REFERENCIAS Y RECURSOS
% ============================================
\chapter*{Referencias y Recursos}
\addcontentsline{toc}{chapter}{Referencias y Recursos}

\section*{Papers Fundamentales}

\begin{enumerate}
    \item Sasson, E. B., et al. (2014). \textit{Zerocash: Decentralized Anonymous Payments from Bitcoin.}
    \item Groth, J. (2016). \textit{On the Size of Pairing-based Non-interactive Arguments.}
    \item Gabizon, A., Williamson, Z., \& Ciobotaru, O. (2019). \textit{PLONK: Permutations over Lagrange-bases for Oecumenical Noninteractive arguments of Knowledge.}
    \item Ben-Sasson, E., et al. (2018). \textit{Scalable, transparent, and post-quantum secure computational integrity.}
    \item Bünz, B., et al. (2018). \textit{Bulletproofs: Short Proofs for Confidential Transactions and More.}
    \item Maller, M., et al. (2019). \textit{Sonic: Zero-Knowledge SNARKs from Linear-Size Universal and Updateable Structured Reference Strings.}
\end{enumerate}

\section*{Documentación de Protocolos}

\begin{itemize}
    \item \textbf{Zcash:} \url{https://zips.z.cash/}
    \item \textbf{Tornado Cash:} \url{https://docs.tornado.cash/}
    \item \textbf{Railgun:} \url{https://docs.railgun.org/}
    \item \textbf{Aztec:} \url{https://docs.aztec.network/}
    \item \textbf{Aleo:} \url{https://developer.aleo.org/}
\end{itemize}

\section*{Herramientas de Desarrollo}

\begin{itemize}
    \item \textbf{Circom:} \url{https://docs.circom.io/}
    \item \textbf{snarkjs:} \url{https://github.com/iden3/snarkjs}
    \item \textbf{Noir:} \url{https://noir-lang.org/}
    \item \textbf{Cairo:} \url{https://www.cairo-lang.org/}
    \item \textbf{Halo2:} \url{https://zcash.github.io/halo2/}
\end{itemize}

\section*{Comunidades y Foros}

\begin{itemize}
    \item \textbf{ZK Research:} \url{https://zkresear.ch/}
    \item \textbf{Ethereum Research:} \url{https://ethresear.ch/}
    \item \textbf{Zero Knowledge Podcast:} \url{https://zeroknowledge.fm/}
\end{itemize}

\cleardoublepage
\end{document}